\documentclass[a4paper,11pt]{article}

% =========================================================================
%  Electromagnetismo I — Tema 3: Conductores en equilibrio electrostático
%  Apuntes formalizados. Compilar con: pdflatex + biber + pdflatex ×2
%  (o latexmk -pdf main.tex). Todo el preámbulo está en preambulo.tex.
%  Los datos numéricos de las figuras «tipo Mathematica» se regeneran con
%  python3 scripts/generar_datos.py (ya se incluyen en figuras/datos/).
% =========================================================================
% =========================================================================
%  Electromagnetismo I — Tema 3: Conductores en equilibrio electrostático
%  Preámbulo. Tipografía, márgenes y espaciado del informe de referencia
%  (clase article, Latin Modern a 11 pt, márgenes de 2,5 cm), con los
%  entornos formales, la notación y el estilo gráfico de los temas 0, 1 y 2.
%  Paleta: azules sobrios con degradados suaves sobre grises neutros.
% =========================================================================

% --- IDIOMA Y CODIFICACIÓN ------------------------------------------------
\usepackage[utf8]{inputenc}
\usepackage[T1]{fontenc}
\usepackage[spanish,es-tabla,es-nodecimaldot,es-noshorthands]{babel}

% --- TIPOGRAFÍA: Latin Modern (Computer Modern escalable), como el informe --
\usepackage{lmodern}
\linespread{1.06}
\emergencystretch=2em
\usepackage[final]{microtype}
\setlength{\parskip}{0.25ex plus 0.15ex}

% --- MATEMÁTICAS -------------------------------------------------------------
\usepackage{amsmath, amssymb, mathtools}
\usepackage{bm}
\usepackage{mathrsfs}
\usepackage{esint}

% --- PÁGINA Y COLORES ---------------------------------------------------------
\usepackage[a4paper, margin=2.5cm, headheight=14pt, headsep=18pt, footskip=30pt]{geometry}
\usepackage{graphicx}
\usepackage[table]{xcolor}
% Paleta del Tema 3: azules sobrios (de la tinta al velo) y grises neutros
\definecolor{azulprof}{HTML}{0B2545}    % azul tinta (títulos, enunciados)
\definecolor{azulmedio}{HTML}{1F4E89}   % azul medio (números, enlaces)
\definecolor{azulacero}{HTML}{3A6EA5}   % azul acero (líneas de campo, acentos)
\definecolor{azulcielo}{HTML}{7FA7D3}   % azul cielo (equipotenciales, marcos)
\definecolor{azulclaro}{HTML}{BCD0EA}   % azul claro (marcos finos)
\definecolor{azulvelo}{HTML}{DCE7F5}    % azul velo (rellenos)
\definecolor{azulfondo}{HTML}{F1F5FB}   % fondo tenue de definiciones
\definecolor{grispizarra}{HTML}{5F6B7A} % gris de texto secundario
\definecolor{grismedio}{HTML}{8A94A3}
\definecolor{grisclaro}{HTML}{D7DDE5}   % rejillas
\definecolor{metal}{HTML}{C9D3E0}       % conductores (relleno)
\definecolor{metalosc}{HTML}{8EA0B8}    % conductores (sombra)
% Acentos mínimos para el signo de las cargas (únicos colores no azules)
\definecolor{cargapos}{HTML}{B03A48}    % carga positiva, vectores destacados
\definecolor{carganeg}{HTML}{1F4E89}    % carga negativa
\definecolor{verdeazul}{HTML}{2A7F8E}   % versores y normales
\colorlet{acento}{cargapos}

% --- TIKZ Y PGFPLOTS ----------------------------------------------------------
\usepackage{tikz}
\usepackage{tikz-3dplot}
\usetikzlibrary{babel, arrows.meta, calc, positioning, decorations.pathmorphing,
                decorations.markings, decorations.pathreplacing, patterns, 3d,
                intersections, angles, quotes, fadings, shapes.geometric, shadings,
                backgrounds, fit}
\usepackage{pgfplots}
\pgfplotsset{compat=1.17}
% Externalización: cada figura se compila una vez a PDF en figuras/ext/ y se
% reutiliza mientras su código no cambie (suma MD5). Requiere -shell-escape
% (ya incluido en latexmkrc; Overleaf lo permite por defecto).
\usepgfplotslibrary{external}
\tikzexternalize[prefix=figuras/ext/]
\tikzset{external/up to date check=md5}
\usepgfplotslibrary{colormaps, fillbetween}
\pgfplotsset{colormap={azulpot}{
  rgb255=(245,248,252) rgb255=(220,231,245) rgb255=(157,185,221)
  rgb255=(58,110,165) rgb255=(31,78,137) rgb255=(11,37,69)}}
\pgfplotsset{colormap={azulcampo}{
  rgb255=(11,37,69) rgb255=(31,78,137) rgb255=(58,134,181)
  rgb255=(127,195,217) rgb255=(207,234,241) rgb255=(255,255,255)}}

% --- TABLAS, LISTAS Y CAJAS ---------------------------------------------------
\usepackage{booktabs}
\usepackage{array}
\usepackage{multirow}
\usepackage{tabularx}
\usepackage{needspace}
\usepackage[section]{placeins}
\usepackage{enumitem}
\setlist[itemize]{leftmargin=1.6em, itemsep=2pt, topsep=4pt, label=\textcolor{azulacero}{\small$\blacktriangleright$}}
\setlist[enumerate]{leftmargin=1.9em, itemsep=2pt, topsep=4pt}
\usepackage[most]{tcolorbox}
\renewcommand{\arraystretch}{1.18}

% --- ENLACES Y REFERENCIAS CRUZADAS ----------------------------------------
% amsthm debe cargarse antes que cleveref para que los enunciados que comparten
% contador se citen con su propio nombre («teorema», «lema»...)
\usepackage{amsthm}
\usepackage[colorlinks=true, linkcolor=azulmedio, citecolor=azulmedio,
            urlcolor=azulmedio, plainpages=false, pdfusetitle=false,
            pdftitle={Tema 3. Conductores en equilibrio electrostático},
            pdfauthor={Daniel Martín Guerrero},
            pdfsubject={Electromagnetismo I}]{hyperref}
\usepackage[spanish,nameinlink]{cleveref}

% --- BIBLIOGRAFÍA -----------------------------------------------------------
\usepackage{csquotes}
\usepackage[backend=biber, style=numeric, sorting=none, giveninits=true,
            maxbibnames=6, minbibnames=3, doi=true, url=true,
            isbn=false]{biblatex}
\addbibresource{referencias.bib}
\DeclareFieldFormat[article,techreport,unpublished,book,inproceedings]{title}{#1\isdot}
\DeclareFieldFormat{url}{\url{#1}}
\renewcommand*{\bibfont}{\small}
\setlength{\bibitemsep}{3pt}

% =========================================================================
%        TÍTULOS DE SECCIÓN, ÍNDICE, CABECERAS Y PIES
% =========================================================================
\usepackage{titlesec}
\usepackage{titletoc}
% Filete con degradado bajo los títulos de sección (no se externaliza)
\newcommand{\filetedegradado}[1][\linewidth]{%
  \par\nointerlineskip\vspace{-0.35\baselineskip}%
  {\tikzset{external/export next=false}%
   \begin{tikzpicture}
     \shade[left color=azulmedio, right color=white, middle color=azulcielo!60]
       (0,0) rectangle (#1,0.75pt);
   \end{tikzpicture}}\par}
\titleformat{\section}[block]
  {\normalfont\Large\bfseries\color{azulprof}}
  {\textcolor{azulmedio}{\thesection.}}{0.7em}{}[\filetedegradado]
\titlespacing*{\section}{0pt}{4.2ex plus 1ex minus .4ex}{1.4ex plus .3ex}
\titleformat{\subsection}[block]
  {\normalfont\large\bfseries\color{azulprof}}
  {\textcolor{azulmedio}{\thesubsection.}}{0.65em}{}
\titlespacing*{\subsection}{0pt}{3.2ex plus 1ex minus .4ex}{1.5ex plus .3ex}
\titleformat{\subsubsection}[block]
  {\normalfont\normalsize\bfseries\color{azulprof}}
  {\textcolor{azulmedio}{\thesubsubsection.}}{0.6em}{}
\titlespacing*{\subsubsection}{0pt}{2.6ex plus .8ex minus .3ex}{1.1ex plus .2ex}
\setcounter{secnumdepth}{3}
\setcounter{tocdepth}{2}

% Índice en la línea del informe: secciones en negrita sin puntos,
% subsecciones con línea de puntos y número de página alineado
\titlecontents{section}[1.9em]{\addvspace{4.5pt}\bfseries\color{azulprof}}
  {\contentslabel[\textcolor{azulmedio}{\thecontentslabel.}]{1.9em}}{}
  {\hfill\textcolor{azulprof}{\contentspage}}
\titlecontents{subsection}[4.6em]{\addvspace{1pt}}
  {\contentslabel[\thecontentslabel.]{2.7em}}{}
  {\titlerule*[0.75pc]{\textcolor{grismedio}{.}}\contentspage}

% Cabeceras discretas en versalitas y número de página centrado (como el informe)
\usepackage{fancyhdr}
\pagestyle{fancy}
\fancyhf{}
\fancyhead[L]{\small\textsc{\color{grispizarra}Electromagnetismo I\enspace·\enspace Tema 3}}
\fancyhead[R]{\small\textsc{\color{grispizarra}\nouppercase{\leftmark}}}
\fancyfoot[C]{\small\color{azulprof}\thepage}
\renewcommand{\headrule}{{\color{azulclaro}\hrule height 0.4pt}\vspace{-0.4pt}}
\renewcommand{\sectionmark}[1]{\markboth{#1}{}}
\fancypagestyle{plain}{\fancyhf{}\fancyfoot[C]{\small\color{azulprof}\thepage}%
  \renewcommand{\headrule}{}}

\numberwithin{equation}{section}

% Pies de figura y tabla
\usepackage[font=small, labelfont={sc,color=azulmedio}, labelsep=period,
            margin=0.6cm, justification=justified]{caption}
\usepackage{subcaption}

% =========================================================================
%               NOTACIÓN Y MACROS DE USO FRECUENTE (temas 0, 1 y 2)
% =========================================================================
% --- Diferencial, derivadas parciales y números -----------------------------
\newcommand{\dif}{\mathrm{d}}
\newcommand{\pd}[2]{\frac{\partial #1}{\partial #2}}
\newcommand{\pdd}[2]{\frac{\partial^{2} #1}{\partial #2^{2}}}
\newcommand{\R}{\mathbb{R}}
\newcommand{\N}{\mathbb{N}}
\newcommand{\eul}{\mathrm{e}}

% --- Vectores, versores y operador nabla ---------------------------------
\newcommand{\vc}[1]{\vec{#1}}                       % vector con flecha
\newcommand{\uh}[1]{\hat{u}_{#1}}                   % versor û_#1
\newcommand{\ux}{\uh{x}}\newcommand{\uy}{\uh{y}}\newcommand{\uz}{\uh{z}}
\newcommand{\ur}{\uh{r}}\newcommand{\uth}{\uh{\theta}}\newcommand{\uph}{\uh{\varphi}}
\newcommand{\urho}{\uh{\rho}}\newcommand{\un}{\hat{n}}
\newcommand{\nab}{\vec{\nabla}}                     % ∇⃗
\newcommand{\Lap}{\nabla^{2}}                       % laplaciano
\newcommand{\norm}[1]{\lVert #1\rVert}
\newcommand{\abs}[1]{\lvert #1\rvert}
\newcommand{\rp}{\vc{r}\,'}                         % r⃗′ punto fuente
\newcommand{\dl}{\dif\vc{\ell}}                     % elemento de línea dirigido
\newcommand{\dS}{\dif\vc{S}}                        % elemento de área dirigido
\newcommand{\dsc}{\dif S}                           % elemento de área escalar
\newcommand{\dV}{\dif V}                            % elemento de volumen
\newcommand{\defeq}{\mathrel{\overset{\mathrm{def}}{=}}}
\newcommand{\bordeS}{\partial S}
\newcommand{\bordeV}{\partial V}
\DeclareMathOperator{\sen}{sen}
\DeclareMathOperator{\sgn}{sgn}
\DeclareMathOperator{\adj}{adj}
\renewcommand{\sin}{\sen}                            % «sen» en todo el texto
\newcommand{\delt}[2]{\delta_{#1#2}}                 % delta de Kronecker
\newcommand{\ord}{\mathcal{O}}                       % orden de Landau

% --- Notación de electrostática (Tema 1) -------------------------------------
\newcommand{\eo}{\varepsilon_{0}}                    % permitividad del vacío
\newcommand{\pot}{\phi}                              % potencial electrostático
\newcommand{\Qenc}{Q_{\mathrm{enc}}}                 % carga encerrada
\newcommand{\Ke}{\frac{1}{4\pi\eo}}                  % constante de Coulomb
\newcommand{\vE}{\vc{E}}
\newcommand{\vF}{\vc{F}}
\newcommand{\dSp}{\dif S\,'}                        % elemento de superficie fuente
\newcommand{\dVp}{\dif V\,'}                        % elemento de volumen fuente
\newcommand{\PhiE}{\Phi_{E}}                         % flujo eléctrico

% --- Notación específica del Tema 3 (conductores) ----------------------------
\newcommand{\Vc}{\mathcal{V}}                        % región ocupada por el conductor
\newcommand{\Sc}{\mathcal{S}}                        % superficie del conductor
\newcommand{\Erest}{\vE_{\mathrm{resto}}}            % campo del resto de la carga
\newcommand{\Edisc}{\vE_{\mathrm{disco}}}            % campo del elemento ΔS
\newcommand{\Pel}{\mathcal{P}}                       % presión electrostática
\newcommand{\Qint}{Q_{\mathrm{int}}}                 % carga de la cara interior
\newcommand{\Qext}{Q_{\mathrm{ext}}}                 % carga de la cara exterior
\newcommand{\Ecav}{\vE_{\mathrm{cav}}}               % campo en la cavidad
\newcommand{\Eext}{\vE_{\mathrm{ext}}}               % campo exterior
\newcommand{\Pm}{\mathsf{P}}                         % matriz de coeficientes de potencial
\newcommand{\Cm}{\mathsf{C}}                         % matriz de coef. de capacidad e influencia
\newcommand{\vQ}{\bm{Q}}                             % vector columna de cargas
\newcommand{\vphi}{\bm{\phi}}                        % vector columna de potenciales
\newcommand{\Erup}{E_{\mathrm{rup}}}                 % campo de ruptura dieléctrica
\newcommand{\DV}{\Delta V}                           % diferencia de potencial
\newcommand{\kc}{k}                                  % k = 1/(4π ε0)
\newcommand{\Tt}{^{\mathsf{T}}}                      % traspuesta

% =========================================================================
%        ENTORNOS FORMALES (enunciados numerados, estilo artículo)
% =========================================================================
% Enunciados en cursiva (resultados generales)
\newtheoremstyle{cienplano}{10pt}{10pt}{\itshape}{}%
  {\bfseries\color{azulprof}}{.}{0.6em}%
  {\thmname{#1}\thmnumber{ #2}\thmnote{\normalfont\upshape\ (#3)}}
% Definiciones, ejemplos y ejercicios (cuerpo en redonda)
\newtheoremstyle{ciendef}{10pt}{10pt}{\normalfont}{}%
  {\bfseries\color{azulprof}}{.}{0.6em}%
  {\thmname{#1}\thmnumber{ #2}\thmnote{\normalfont\upshape\ (#3)}}
% Observaciones y notas (encabezado en cursiva)
\newtheoremstyle{cienobs}{8pt}{8pt}{\normalfont}{}%
  {\itshape\color{azulmedio}}{.}{0.6em}%
  {\thmname{#1}\thmnumber{ #2}\thmnote{\normalfont\upshape\ (#3)}}

\theoremstyle{cienplano}
\newtheorem{proposicion}{Proposición}[section]
\newtheorem{teorema}[proposicion]{Teorema}
\newtheorem{lema}[proposicion]{Lema}
\newtheorem{corolario}[proposicion]{Corolario}
\theoremstyle{ciendef}
\newtheorem{notacion}[proposicion]{Notación}
\newtheorem{ejemplo}[proposicion]{Ejemplo}
\newtheorem{ejercicio}[proposicion]{Ejercicio}
\theoremstyle{cienobs}
\newtheorem{observacion}[proposicion]{Observación}

% Teoremas destacados: filete vertical azul con un velo en degradado
\tcolorboxenvironment{teorema}{enhanced, breakable, shield externalize,
  frame hidden, sharp corners, boxrule=0pt,
  interior style={left color=azulfondo, right color=white},
  borderline west={1.2pt}{0pt}{azulmedio},
  left=8pt, right=4pt, top=0pt, bottom=0pt, before skip=10pt, after skip=10pt}

% Definiciones: caja fina con degradado horizontal y filete lateral.
% Comparten contador con los enunciados (como en los temas anteriores).
\newcounter{definicion}
\renewcommand{\thedefinicion}{\theproposicion}
\newtcolorbox{definicionbox}[1]{%
  enhanced, breakable, sharp corners, shield externalize,
  frame hidden, boxrule=0pt,
  interior style={left color=azulfondo, right color=azulfondo!25!white},
  borderline west={2pt}{0pt}{azulmedio},
  borderline south={0.4pt}{0pt}{azulclaro},
  borderline north={0.4pt}{0pt}{azulclaro},
  left=10pt, right=9pt, top=6pt, bottom=6pt,
  before skip=12pt, after skip=12pt,
  title={Definición~\thedefinicion\ifx\relax#1\relax\else\ {\normalfont(#1)}\fi},
  coltitle=azulprof, fonttitle=\bfseries, detach title, before upper={\textbf{\color{azulprof}\tcbtitle.}\enspace}}
\newenvironment{definicion}[1][]
  {\par\stepcounter{proposicion}\refstepcounter{definicion}\begin{definicionbox}{#1}}
  {\end{definicionbox}}

% Enunciados de los problemas de los seminarios
\newtcolorbox{enunciadobox}[1]{%
  enhanced, breakable, sharp corners, shield externalize,
  frame hidden, boxrule=0pt,
  interior style={top color=azulvelo!55!white, bottom color=azulfondo!40!white},
  borderline west={2pt}{0pt}{azulprof},
  left=10pt, right=9pt, top=6pt, bottom=6pt, before skip=10pt, after skip=10pt,
  fontupper=\itshape,
  title={#1}, coltitle=azulprof, fonttitle=\upshape\bfseries, detach title,
  before upper={{\upshape\bfseries\color{azulprof}\tcbtitle.}\enspace}}
\newcounter{problema}
\renewcommand{\theproblema}{3.\arabic{problema}}
\newenvironment{problema}[1][]
  {\par\refstepcounter{problema}%
   \begin{enunciadobox}{Problema~\theproblema\ifx\relax#1\relax\else\ \normalfont\upshape(#1)\fi}}
  {\end{enunciadobox}}

% Cajas de «receta» y de comprobación (seminarios y guía de métodos)
\newtcolorbox{receta}[1][Receta]{%
  enhanced, breakable, sharp corners, shield externalize,
  colback=white, colframe=azulclaro, boxrule=0.5pt,
  borderline west={2pt}{0pt}{azulacero},
  left=9pt, right=8pt, top=5pt, bottom=5pt, before skip=10pt, after skip=10pt,
  title={#1}, fonttitle=\small\scshape, coltitle=azulprof,
  colbacktitle=azulvelo!70!white, toptitle=1pt, bottomtitle=1pt}

% Ecuaciones destacadas: marco fino con velo en degradado
\tcbset{resaltado/.style={enhanced, shield externalize, sharp corners, nobeforeafter,
  tcbox raise base, colframe=azulcielo, boxrule=0.45pt,
  interior style={left color=azulfondo, right color=white},
  left=4pt, right=4pt, top=2pt, bottom=2pt}}
\newcommand{\resalta}[1]{\tcboxmath[resaltado]{#1}}

% Resolución de ejercicios: misma maquinaria que «proof», con otro rótulo
\renewcommand{\qedsymbol}{\textcolor{azulmedio}{\ensuremath{\blacksquare}}}
\newenvironment{resolucion}[1][Resolución]
  {\begin{proof}[\itshape\color{azulmedio}#1]}{\end{proof}}
% Pasos rotulados dentro de las resoluciones
\newcommand{\paso}[1]{\par\smallskip\noindent\textbf{\color{azulprof}#1}\enspace}

\crefname{proposicion}{proposición}{proposiciones}
\Crefname{proposicion}{Proposición}{Proposiciones}
\crefname{teorema}{teorema}{teoremas}
\Crefname{teorema}{Teorema}{Teoremas}
\crefname{lema}{lema}{lemas}
\Crefname{lema}{Lema}{Lemas}
\crefname{corolario}{corolario}{corolarios}
\Crefname{corolario}{Corolario}{Corolarios}
\crefname{definicion}{definición}{definiciones}
\Crefname{definicion}{Definición}{Definiciones}
\crefname{notacion}{notación}{notaciones}
\Crefname{notacion}{Notación}{Notaciones}
\crefname{ejemplo}{ejemplo}{ejemplos}
\Crefname{ejemplo}{Ejemplo}{Ejemplos}
\crefname{ejercicio}{ejercicio}{ejercicios}
\Crefname{ejercicio}{Ejercicio}{Ejercicios}
\crefname{observacion}{observación}{observaciones}
\Crefname{observacion}{Observación}{Observaciones}
\crefname{problema}{problema}{problemas}
\Crefname{problema}{Problema}{Problemas}
\crefname{section}{sección}{secciones}
\Crefname{section}{Sección}{Secciones}
\crefname{subsection}{sección}{secciones}
\Crefname{subsection}{Sección}{Secciones}
\crefname{subsubsection}{sección}{secciones}
\Crefname{subsubsection}{Sección}{Secciones}
\crefname{table}{tabla}{tablas}
\Crefname{table}{Tabla}{Tablas}
\crefname{figure}{figura}{figuras}
\Crefname{figure}{Figura}{Figuras}
\crefname{equation}{ecuación}{ecuaciones}
\Crefname{equation}{Ecuación}{Ecuaciones}
\crefformat{equation}{ecuación~(#2#1#3)}
\Crefformat{equation}{Ecuación~(#2#1#3)}
\crefrangeformat{equation}{ecuaciones~(#3#1#4)--(#5#2#6)}
\crefmultiformat{equation}{ecuaciones~(#2#1#3)}{ y~(#2#1#3)}{, (#2#1#3)}{ y~(#2#1#3)}

% Referencia a un panel de figura: «figura 5\,(b)», enlazada en bloque
\newcommand{\figpanel}[2]{\hyperref[#1]{figura~\ref*{#1}\,\textup{(#2)}}}
\newcommand{\Figpanel}[2]{\hyperref[#1]{Figura~\ref*{#1}\,\textup{(#2)}}}

% =========================================================================
%        ESTILOS GRÁFICOS COMUNES (TikZ / pgfplots)
% =========================================================================
% =========================================================================
%  Estilos y macros gráficas comunes del Tema 3
%  - \mapcolor{t}: escala cromática azul (t en [0,1]); define el color mapc
%  - Estilos TikZ heredados de los temas 0, 1 y 2 (ejes, vectores, versores…)
%    con la paleta azul del documento
%  - Macros que dibujan los datos numéricos de figuras/datos/*.tex
%    (\LC, \EQ, \FL, \CP, \CN, \SG), generados por scripts/generar_datos.py
% =========================================================================

% --- Escala cromática azul (5 tramos: del velo a la tinta) -------------------
\definecolor{cm0}{HTML}{DCE7F5}
\definecolor{cm1}{HTML}{9DB9DD}
\definecolor{cm2}{HTML}{3A6EA5}
\definecolor{cm3}{HTML}{1F4E89}
\definecolor{cm4}{HTML}{0B2545}
\newcommand{\mapcolor}[1]{%
  \pgfmathsetmacro{\mt}{min(max(#1,0),0.9999)*4}%
  \pgfmathtruncatemacro{\mi}{floor(\mt)}%
  \pgfmathtruncatemacro{\mj}{\mi+1}%
  \pgfmathsetmacro{\mf}{100*(\mt-\mi)}%
  \colorlet{mapc}{cm\mj!\mf!cm\mi}}

% --- Estilos TikZ generales ------------------------------------------------
\tikzset{
  >=Stealth,
  eje/.style={->, line width=0.55pt, grispizarra},
  vect/.style={->, line width=1pt, azulprof},
  vectrojo/.style={->, line width=1pt, cargapos},
  vectverde/.style={->, line width=1pt, verdeazul},
  fino/.style={line width=0.4pt, grispizarra},
  punteado/.style={densely dotted, line width=0.5pt, grispizarra},
  etiq/.style={font=\small, text=black!85},
  cuadro/.style={draw=azulclaro, line width=0.4pt, fill=white, inner sep=2pt},
  % rótulo con fondo blanco semitransparente: evita que el texto se confunda con trazos
  fondo/.style={fill=white, fill opacity=0.85, text opacity=1, inner sep=1.2pt, rounded corners=1pt},
  orienta/.style={decoration={markings, mark=at position 0.55 with {\arrow{Stealth[length=5pt]}}}, postaction={decorate}},
  % conductores: relleno metálico en degradado y contorno en azul tinta
  conductor/.style={draw=azulprof, line width=0.6pt, top color=azulvelo!80!white,
                    bottom color=metal, shading angle=0},
  conductorplano/.style={draw=azulprof, line width=0.6pt, fill=metal!70!white},
  metalrayado/.style={draw=azulprof, line width=0.6pt, fill=azulvelo,
                      postaction={pattern=north east lines, pattern color=azulcielo!70}},
  % superficies gaussianas
  gauss/.style={densely dashed, line width=0.75pt, cargapos!85},
  gaussazul/.style={densely dashed, line width=0.7pt, azulacero},
  % líneas de campo y equipotenciales
  lcampo/.style={azulacero, line width=0.55pt, postaction={decorate},
                 decoration={markings, mark=at position 0.5 with {\arrow{Stealth[length=3.6pt]}}}},
  lcampo2/.style={azulacero, line width=0.55pt, postaction={decorate},
                 decoration={markings, mark=at position 0.3 with {\arrow{Stealth[length=3.6pt]}},
                             mark=at position 0.75 with {\arrow{Stealth[length=3.6pt]}}}},
  lequi/.style={azulcielo, line width=0.45pt, densely dashed},
  % tierra
  tierra/.pic={\draw[line width=0.6pt, azulprof] (0,0) -- (0,-0.18)
                 (-0.18,-0.18) -- (0.18,-0.18) (-0.11,-0.25) -- (0.11,-0.25)
                 (-0.045,-0.32) -- (0.045,-0.32);},
}

% Cargas puntuales y signos de carga superficial
\tikzset{
  qpos/.style={circle, fill=cargapos, draw=cargapos!70!black, line width=0.4pt, inner sep=0pt,
               minimum size=9pt, text=white, font=\scriptsize\bfseries},
  qneg/.style={circle, fill=carganeg, draw=azulprof, line width=0.4pt, inner sep=0pt,
               minimum size=9pt, text=white, font=\scriptsize\bfseries},
  signopos/.style={text=cargapos, font=\scriptsize\bfseries, inner sep=0pt},
  signoneg/.style={text=carganeg, font=\scriptsize\bfseries, inner sep=0pt},
}
\newcommand{\cargaP}[1]{\node[qpos] at (#1) {$+$};}
\newcommand{\cargaN}[1]{\node[qneg] at (#1) {$-$};}

% --- Estilo base de ejes pgfplots (como en los temas anteriores) -----------
\pgfplotsset{
  tema/.style={
    tick label style={font=\scriptsize}, label style={font=\small},
    axis lines=left, axis line style={grispizarra},
    grid=major, grid style={grisclaro, line width=0.35pt},
    every tick/.style={grispizarra},
    legend style={font=\scriptsize, draw=none, fill=white, fill opacity=0.85,
                  text opacity=1},
    legend cell align=left,
    /pgf/number format/use comma,
  },
  marco/.style={
    tick label style={font=\scriptsize}, label style={font=\small},
    axis line style={grispizarra!80, line width=0.5pt},
    every tick/.style={grispizarra}, tick align=inside,
    grid=major, grid style={grisclaro!70, line width=0.3pt},
    /pgf/number format/use comma,
    legend style={font=\scriptsize, draw=none, fill=white, fill opacity=0.85,
                  text opacity=1},
    legend cell align=left,
  },
}

% =========================================================================
%  Macros que interpretan los datos numéricos (figuras/datos/*.tex).
%  Cada figura puede redefinirlas localmente antes de \input.
% =========================================================================
\tikzset{
  numLC/.style={azulmedio, line width=0.45pt},
  numEQ/.style={azulcielo, line width=0.35pt},
  numFL/.style={-{Stealth[length=3.4pt, width=2.8pt]}, azulmedio, line width=0.45pt},
}
\newcommand{\LC}[1]{\draw[numLC] plot coordinates {#1};}
\newcommand{\EQ}[1]{\draw[numEQ] plot coordinates {#1};}
\newcommand{\FL}[3]{\draw[numFL] (#1,#2) -- ++(#3:0.012);}
\newcommand{\CP}[2]{\node[signopos] at (#1,#2) {$+$};}
\newcommand{\CN}[2]{\node[signoneg] at (#1,#2) {$-$};}
\newcommand{\SG}[5]{\mapcolor{#5}\draw[mapc, line width=1.6pt, line cap=round] (#1,#2) -- (#3,#4);}

% Marco fino con marcas (ventana [x0,x1] x [y0,y1]) para los paneles numéricos
\newcommand{\marcopanel}[4]{%
  \draw[azulclaro, line width=0.5pt] (#1,#3) rectangle (#2,#4);}

% =========================================================================
%  Primitivas 3D para tdplot (vista fija: \tdplotsetmaincoords{70}{125}),
%  heredadas de los temas 1 y 2.
% =========================================================================
\def\VA{125}
% \tdCirc[opciones]{radio}{z}   (mitad frontal continua, trasera a trazos)
\newcommand{\tdCirc}[3][]{%
  \draw[#1, domain=\VA-180:\VA, samples=36, variable=\tdt]
    plot ({#2*cos(\tdt)},{#2*sin(\tdt)},{#3});
  \draw[#1, densely dashed, domain=\VA:\VA+180, samples=36, variable=\tdt]
    plot ({#2*cos(\tdt)},{#2*sin(\tdt)},{#3});}
% \tdEsfera{radio}{color}  (esfera sombreada con ecuador)
\newcommand{\tdEsfera}[2]{%
  \begin{scope}[tdplot_screen_coords]
    \shade[ball color=azulvelo, opacity=0.55] (0,0) circle (#1);
    \draw[#2, line width=0.7pt] (0,0) circle (#1);
  \end{scope}
  \tdCirc[#2!60, line width=0.4pt]{#1}{0}}



% =========================================================================
%                     VARIABLES PARA LA PORTADA
% =========================================================================
\newcommand{\Universidad}{Universidad de Salamanca}
\newcommand{\Facultad}{Facultad de Ciencias}
\newcommand{\Titulacion}{Grado en Física}
\newcommand{\Asignatura}{Electromagnetismo I}
\newcommand{\NumTema}{Tema 3}
\newcommand{\TituloTema}{Conductores en equilibrio\\electrostático}
\newcommand{\SubtituloTema}{Propiedades del conductor perfecto, efecto punta, presión electrostática,
  cavidades y apantallamiento; coeficientes de potencial, de capacidad e influencia, y condensadores}
\newcommand{\tipotrabajo}{Apuntes formalizados con seminarios prácticos}
\newcommand{\miNombre}{Daniel Martín Guerrero}
\newcommand{\miProfesor}{Luis López Díaz}
\newcommand{\Afiliacion}{Grado en Física, Facultad de Ciencias, Universidad de Salamanca,
  Plaza de la Merced s/n, 37008 Salamanca, España}
\newcommand{\miUbicacion}{Salamanca, curso 2026/2027}
\newcommand{\ResumenPortada}{%
  Estudiamos la electrostática en presencia de conductores. A partir de la definición de conductor
  perfecto y de la condición de equilibrio deducimos sus propiedades: campo nulo en el interior,
  volumen equipotencial, carga confinada en la superficie y campo exterior normal de módulo
  $\sigma/\eo$. De ellas se siguen el efecto punta, la presión electrostática $\sigma^{2}/2\eo$, el
  apantallamiento de las cavidades y la capacidad de un conductor aislado. Para un sistema de $N$
  conductores, la linealidad de las ecuaciones de Laplace y Poisson define los coeficientes de potencial
  $p_{ij}$ y los de capacidad e influencia $c_{ij}$, cuya simetría, signo e inversión demostramos; la
  influencia total conduce al condensador. Los seminarios resuelven los siete problemas del tema.}

% =========================================================================
%                                DOCUMENTO
% =========================================================================
\begin{document}
% --- PORTADA ----------------------------------------------------------------
\pagenumbering{alph}
%%%%%%%%%%%%%%%%%%%%%%%%%%%%%%%%%%%%%%%%%%%%%%%%%%%%%%%%%%%%%%%%%%%%%%%%
% Portada con estructura de artículo científico: cabecera institucional,
% título y autoría alineados a la izquierda, resumen, palabras clave y
% clasificación, y un emblema calculado numéricamente (líneas de campo y
% equipotenciales de un conductor cargado con punta, figuras/datos/punta_*).
% Todos los bloques tienen coordenadas fijas en la página A4 (21 x 29,7 cm)
% para que ningún elemento se superponga a otro.
% Requiere en main.tex: \Universidad, \Facultad, \Titulacion, \Asignatura,
% \NumTema, \TituloTema, \SubtituloTema, \tipotrabajo, \miNombre,
% \miProfesor, \miUbicacion, \Afiliacion y el texto \ResumenPortada.
%%%%%%%%%%%%%%%%%%%%%%%%%%%%%%%%%%%%%%%%%%%%%%%%%%%%%%%%%%%%%%%%%%%%%%%%

\begin{titlepage}
\thispagestyle{empty}
\tikzexternaldisable
\begin{tikzpicture}[remember picture, overlay]
  \coordinate (O) at (current page.south west);

  % --- Bandas superior e inferior en degradado --------------------------
  \shade[left color=azulprof, right color=azulcielo!70]
    ($(O)+(0,29.25)$) rectangle ($(O)+(21,29.7)$);
  \shade[left color=azulcielo!45, right color=white]
    ($(O)+(0,29.15)$) rectangle ($(O)+(21,29.25)$);
  \shade[left color=azulprof, right color=azulcielo!70]
    ($(O)+(0,0)$) rectangle ($(O)+(21,0.3)$);
  % velo muy tenue en la cabecera
  \shade[top color=azulfondo, bottom color=white]
    ($(O)+(0,24.9)$) rectangle ($(O)+(21,29.15)$);

  % --- Cabecera institucional -------------------------------------------
  \node[anchor=north west, inner sep=0] at ($(O)+(2.5,28.3)$)
    {\includegraphics[height=1.2cm]{logo_usal.png}};
  \node[anchor=north east, inner sep=0] at ($(O)+(18.5,28.3)$)
    {\includegraphics[height=1.2cm]{logo_facultad.png}};
  \node[anchor=north, align=center, inner sep=0, text width=7.2cm] at ($(O)+(10.5,27.95)$)
    {{\color{azulprof}\scshape\textls[110]{\Universidad}}\\[3pt]
     {\color{grispizarra}\footnotesize\scshape\textls[50]{\Facultad\enspace·\enspace\Titulacion}}};
  \draw[azulclaro, line width=0.45pt] ($(O)+(2.5,26.55)$) -- ($(O)+(18.5,26.55)$);

  % --- Asignatura, título y subtítulo -----------------------------------
  \node[anchor=north west, align=left, inner sep=0, text width=16cm] at ($(O)+(2.5,25.95)$)
    {{\color{azulmedio}\scshape\textls[140]{\Asignatura\enspace·\enspace\NumTema}}};
  \node[anchor=north west, align=left, inner sep=0, text width=16cm,
        font=\bfseries\fontsize{29}{35}\selectfont, text=azulprof] at ($(O)+(2.5,25.05)$)
    {\TituloTema};
  \shade[left color=azulmedio, right color=white]
    ($(O)+(2.5,22.38)$) rectangle ($(O)+(10.5,22.44)$);
  \node[anchor=north west, align=left, inner sep=0, text width=15.6cm,
        font=\itshape\large, text=black!78] at ($(O)+(2.5,22.05)$)
    {\SubtituloTema};

  % --- Autoría y afiliación ----------------------------------------------
  \node[anchor=north west, align=left, inner sep=0, text width=7.6cm, font=\large] at ($(O)+(2.5,20.75)$)
    {{\scshape\small\color{azulmedio}\textls[80]{Autor}}\\[2pt]
     {\color{black!88}\miNombre\textsuperscript{\,a}}};
  \node[anchor=north west, align=left, inner sep=0, text width=7.6cm, font=\large] at ($(O)+(10.9,20.75)$)
    {{\scshape\small\color{azulmedio}\textls[80]{Profesor}}\\[2pt]
     {\color{black!88}\miProfesor\textsuperscript{\,a}}};
  \node[anchor=north west, align=left, inner sep=0, text width=16cm, font=\footnotesize] at ($(O)+(2.5,19.55)$)
    {\color{grispizarra}\textsuperscript{a}\,\Afiliacion};

  % --- Resumen (izquierda) y emblema (derecha) ---------------------------
  \node[anchor=north west, inner sep=0] at ($(O)+(2.5,18.6)$) {%
    \begin{tcolorbox}[enhanced, width=10.1cm, sharp corners, frame hidden, boxrule=0pt,
        interior style={left color=azulfondo, right color=white},
        borderline west={2pt}{0pt}{azulmedio},
        left=9pt, right=6pt, top=6pt, bottom=6pt]
      {\small\scshape\color{azulprof}\textls[90]{Resumen}}\par\smallskip
      {\footnotesize\ResumenPortada}
    \end{tcolorbox}};

  \begin{scope}[shift={($(O)+(15.7,15.4)$)}, x=1.06cm, y=1.06cm]
    \begin{scope}
      \clip[rounded corners=3pt] (-2.6,-2.0) rectangle (2.6,2.0);
      \fill[inner color=white, outer color=azulfondo] (-2.6,-2.0) rectangle (2.6,2.0);
      \tikzset{numLC/.style={azulacero!85, line width=0.42pt},
               numEQ/.style={azulcielo!80, line width=0.3pt, densely dashed},
               numFL/.style={-{Stealth[length=3pt, width=2.5pt]}, azulacero, line width=0.42pt}}
      % Conductor con punta - equipotenciales
% Generado por scripts/generar_datos.py (no editar a mano).
\EQ{(2.341,-2.040) (2.370,-2.012) (2.400,-1.982) (2.430,-1.951) (2.459,-1.920) (2.486,-1.890) (2.513,-1.860) (2.535,-1.835) (2.564,-1.800) (2.589,-1.770) (2.610,-1.743) (2.636,-1.710) (2.640,-1.705)}
\EQ{(-2.640,-1.913) (-2.610,-1.949) (-2.583,-1.980) (-2.556,-2.010) (-2.535,-2.033) (-2.529,-2.040)}
\EQ{(2.640,1.705) (2.613,1.740) (2.589,1.770) (2.565,1.799) (2.539,1.830) (2.513,1.860) (2.490,1.886) (2.460,1.919) (2.431,1.950) (2.402,1.980) (2.372,2.010) (2.341,2.040)}
\EQ{(-2.529,2.040) (-2.550,2.017) (-2.580,1.983) (-2.609,1.950) (-2.634,1.920) (-2.640,1.913)}
\EQ{(1.657,-2.040) (1.695,-2.015) (1.724,-1.995) (1.755,-1.973) (1.785,-1.952) (1.815,-1.930) (1.845,-1.907) (1.875,-1.884) (1.905,-1.860) (1.935,-1.835) (1.959,-1.815) (1.994,-1.785) (2.025,-1.758) (2.055,-1.730) (2.085,-1.702) (2.108,-1.680) (2.130,-1.658) (2.153,-1.635) (2.175,-1.612) (2.196,-1.590) (2.220,-1.564) (2.250,-1.531) (2.277,-1.500) (2.302,-1.470) (2.325,-1.442) (2.351,-1.410) (2.370,-1.385) (2.397,-1.350) (2.415,-1.325) (2.440,-1.290) (2.460,-1.260) (2.480,-1.230) (2.500,-1.200) (2.518,-1.170) (2.535,-1.143) (2.554,-1.110) (2.571,-1.080) (2.588,-1.050) (2.604,-1.020) (2.619,-0.990) (2.634,-0.960) (2.640,-0.947)}
\EQ{(-2.640,-1.202) (-2.625,-1.230) (-2.610,-1.258) (-2.595,-1.284) (-2.574,-1.320) (-2.556,-1.350) (-2.537,-1.380) (-2.520,-1.406) (-2.497,-1.440) (-2.477,-1.470) (-2.455,-1.500) (-2.433,-1.530) (-2.415,-1.554) (-2.387,-1.590) (-2.363,-1.620) (-2.340,-1.648) (-2.312,-1.680) (-2.286,-1.710) (-2.265,-1.733) (-2.244,-1.755) (-2.220,-1.780) (-2.190,-1.810) (-2.160,-1.839) (-2.138,-1.860) (-2.115,-1.881) (-2.089,-1.905) (-2.055,-1.934) (-2.025,-1.960) (-2.000,-1.980) (-1.965,-2.008) (-1.935,-2.031) (-1.923,-2.040)}
\EQ{(2.640,0.947) (2.626,0.975) (2.611,1.005) (2.596,1.035) (2.580,1.064) (2.565,1.091) (2.550,1.117) (2.528,1.155) (2.509,1.185) (2.490,1.215) (2.470,1.245) (2.450,1.275) (2.430,1.304) (2.408,1.335) (2.385,1.365) (2.363,1.395) (2.340,1.424) (2.315,1.455) (2.295,1.479) (2.265,1.514) (2.238,1.545) (2.210,1.575) (2.182,1.605) (2.160,1.628) (2.138,1.650) (2.115,1.673) (2.092,1.695) (2.070,1.716) (2.044,1.740) (2.011,1.770) (1.980,1.797) (1.950,1.823) (1.923,1.845) (1.890,1.872) (1.860,1.895) (1.830,1.918) (1.800,1.941) (1.770,1.963) (1.746,1.980) (1.710,2.005) (1.680,2.025) (1.657,2.040)}
\EQ{(-1.923,2.040) (-1.950,2.019) (-1.980,1.996) (-2.010,1.972) (-2.037,1.950) (-2.070,1.921) (-2.100,1.895) (-2.130,1.868) (-2.154,1.845) (-2.185,1.815) (-2.215,1.785) (-2.244,1.755) (-2.265,1.733) (-2.286,1.710) (-2.310,1.683) (-2.338,1.650) (-2.363,1.620) (-2.385,1.593) (-2.411,1.560) (-2.430,1.535) (-2.455,1.500) (-2.475,1.472) (-2.497,1.440) (-2.517,1.410) (-2.535,1.383) (-2.556,1.350) (-2.574,1.320) (-2.592,1.290) (-2.609,1.260) (-2.625,1.230) (-2.640,1.202)}
\EQ{(0.464,-2.040) (0.495,-2.032) (0.525,-2.025) (0.555,-2.017) (0.585,-2.008) (0.615,-2.000) (0.645,-1.991) (0.675,-1.982) (0.705,-1.972) (0.735,-1.962) (0.765,-1.952) (0.795,-1.941) (0.825,-1.930) (0.855,-1.919) (0.885,-1.908) (0.915,-1.896) (0.945,-1.884) (0.975,-1.871) (1.005,-1.858) (1.034,-1.845) (1.065,-1.831) (1.095,-1.817) (1.125,-1.803) (1.155,-1.788) (1.185,-1.772) (1.215,-1.757) (1.245,-1.741) (1.274,-1.725) (1.300,-1.710) (1.326,-1.695) (1.365,-1.672) (1.395,-1.654) (1.425,-1.635) (1.455,-1.616) (1.485,-1.596) (1.515,-1.576) (1.545,-1.555) (1.575,-1.534) (1.600,-1.515) (1.635,-1.489) (1.659,-1.470) (1.695,-1.442) (1.725,-1.417) (1.751,-1.395) (1.785,-1.366) (1.815,-1.339) (1.845,-1.311) (1.867,-1.290) (1.890,-1.267) (1.912,-1.245) (1.935,-1.222) (1.956,-1.200) (1.980,-1.174) (2.010,-1.140) (2.036,-1.110) (2.061,-1.080) (2.085,-1.050) (2.108,-1.020) (2.130,-0.991) (2.152,-0.960) (2.173,-0.930) (2.190,-0.905) (2.213,-0.870) (2.231,-0.840) (2.249,-0.810) (2.265,-0.782) (2.280,-0.755) (2.295,-0.726) (2.310,-0.697) (2.325,-0.666) (2.340,-0.633) (2.354,-0.600) (2.367,-0.570) (2.378,-0.540) (2.389,-0.510) (2.400,-0.480) (2.409,-0.450) (2.419,-0.420) (2.427,-0.390) (2.435,-0.360) (2.442,-0.330) (2.448,-0.300) (2.454,-0.270) (2.460,-0.240) (2.464,-0.210) (2.468,-0.180) (2.472,-0.150) (2.475,-0.120) (2.477,-0.090) (2.478,-0.060) (2.479,-0.030) (2.479,-0.000) (2.479,0.030) (2.478,0.060) (2.477,0.090) (2.475,0.120) (2.472,0.150) (2.468,0.180) (2.464,0.210) (2.460,0.240) (2.454,0.270) (2.448,0.300) (2.442,0.330) (2.435,0.360) (2.427,0.390) (2.419,0.420) (2.409,0.450) (2.400,0.479) (2.389,0.510) (2.378,0.540) (2.367,0.570) (2.355,0.598) (2.341,0.630) (2.328,0.660) (2.313,0.690) (2.298,0.720) (2.283,0.750) (2.266,0.780) (2.250,0.808) (2.231,0.840) (2.213,0.870) (2.193,0.900) (2.175,0.927) (2.152,0.960) (2.131,0.990) (2.108,1.020) (2.085,1.050) (2.061,1.080) (2.040,1.105) (2.010,1.140) (1.983,1.170) (1.956,1.200) (1.935,1.222) (1.912,1.245) (1.890,1.267) (1.867,1.290) (1.845,1.311) (1.819,1.335) (1.786,1.365) (1.755,1.392) (1.725,1.417) (1.697,1.440) (1.665,1.466) (1.640,1.485) (1.605,1.511) (1.580,1.530) (1.545,1.555) (1.516,1.575) (1.485,1.596) (1.455,1.616) (1.425,1.635) (1.395,1.654) (1.365,1.672) (1.335,1.690) (1.305,1.707) (1.275,1.724) (1.246,1.740) (1.218,1.755) (1.190,1.770) (1.160,1.785) (1.130,1.800) (1.099,1.815) (1.067,1.830) (1.035,1.845) (1.005,1.858) (0.975,1.871) (0.945,1.884) (0.915,1.896) (0.885,1.908) (0.855,1.919) (0.825,1.930) (0.795,1.941) (0.765,1.952) (0.735,1.962) (0.705,1.972) (0.675,1.982) (0.645,1.991) (0.615,2.000) (0.585,2.008) (0.555,2.017) (0.525,2.025) (0.495,2.032) (0.465,2.040) (0.464,2.040)}
\EQ{(-0.923,2.040) (-0.960,2.030) (-0.990,2.021) (-1.020,2.012) (-1.050,2.003) (-1.080,1.993) (-1.110,1.982) (-1.140,1.971) (-1.170,1.960) (-1.200,1.948) (-1.230,1.936) (-1.260,1.923) (-1.290,1.910) (-1.320,1.896) (-1.350,1.881) (-1.380,1.866) (-1.410,1.851) (-1.440,1.835) (-1.470,1.818) (-1.500,1.801) (-1.526,1.785) (-1.560,1.764) (-1.590,1.745) (-1.619,1.725) (-1.650,1.704) (-1.680,1.682) (-1.710,1.660) (-1.740,1.637) (-1.770,1.613) (-1.797,1.590) (-1.830,1.562) (-1.860,1.535) (-1.890,1.507) (-1.913,1.485) (-1.935,1.463) (-1.957,1.440) (-1.980,1.416) (-2.010,1.383) (-2.039,1.350) (-2.064,1.320) (-2.085,1.294) (-2.111,1.260) (-2.130,1.235) (-2.155,1.200) (-2.175,1.171) (-2.196,1.140) (-2.215,1.110) (-2.234,1.080) (-2.250,1.052) (-2.265,1.026) (-2.280,0.999) (-2.295,0.971) (-2.310,0.942) (-2.324,0.915) (-2.338,0.885) (-2.352,0.855) (-2.365,0.825) (-2.378,0.795) (-2.390,0.765) (-2.401,0.735) (-2.412,0.705) (-2.423,0.675) (-2.433,0.645) (-2.442,0.615) (-2.451,0.585) (-2.459,0.555) (-2.468,0.525) (-2.475,0.495) (-2.482,0.465) (-2.489,0.435) (-2.495,0.405) (-2.500,0.375) (-2.506,0.345) (-2.511,0.315) (-2.515,0.285) (-2.519,0.255) (-2.522,0.225) (-2.525,0.195) (-2.528,0.165) (-2.530,0.135) (-2.532,0.105) (-2.533,0.075) (-2.534,0.045) (-2.534,0.015) (-2.534,-0.030) (-2.533,-0.060) (-2.532,-0.090) (-2.531,-0.120) (-2.529,-0.150) (-2.527,-0.180) (-2.524,-0.210) (-2.520,-0.240) (-2.517,-0.270) (-2.513,-0.300) (-2.508,-0.330) (-2.503,-0.360) (-2.498,-0.390) (-2.492,-0.420) (-2.485,-0.450) (-2.479,-0.480) (-2.471,-0.510) (-2.464,-0.540) (-2.455,-0.570) (-2.447,-0.600) (-2.437,-0.630) (-2.428,-0.660) (-2.417,-0.690) (-2.407,-0.720) (-2.395,-0.750) (-2.384,-0.780) (-2.371,-0.810) (-2.358,-0.840) (-2.345,-0.870) (-2.331,-0.900) (-2.316,-0.930) (-2.301,-0.960) (-2.285,-0.990) (-2.269,-1.020) (-2.251,-1.050) (-2.235,-1.078) (-2.215,-1.110) (-2.196,-1.140) (-2.176,-1.170) (-2.155,-1.200) (-2.133,-1.230) (-2.115,-1.255) (-2.088,-1.290) (-2.064,-1.320) (-2.040,-1.348) (-2.012,-1.380) (-1.985,-1.410) (-1.957,-1.440) (-1.935,-1.463) (-1.913,-1.485) (-1.890,-1.507) (-1.865,-1.530) (-1.832,-1.560) (-1.800,-1.588) (-1.770,-1.613) (-1.742,-1.635) (-1.710,-1.660) (-1.683,-1.680) (-1.650,-1.704) (-1.620,-1.725) (-1.590,-1.745) (-1.560,-1.764) (-1.530,-1.783) (-1.501,-1.800) (-1.470,-1.818) (-1.440,-1.835) (-1.410,-1.851) (-1.380,-1.866) (-1.350,-1.881) (-1.320,-1.896) (-1.290,-1.910) (-1.260,-1.923) (-1.232,-1.935) (-1.200,-1.948) (-1.170,-1.960) (-1.140,-1.971) (-1.110,-1.982) (-1.080,-1.993) (-1.050,-2.003) (-1.020,-2.012) (-0.990,-2.021) (-0.960,-2.030) (-0.930,-2.038) (-0.923,-2.040)}
\EQ{(-0.510,-1.726) (-0.480,-1.728) (-0.450,-1.730) (-0.420,-1.731) (-0.390,-1.732) (-0.360,-1.733) (-0.330,-1.733) (-0.300,-1.733) (-0.270,-1.733) (-0.240,-1.732) (-0.210,-1.731) (-0.180,-1.729) (-0.150,-1.727) (-0.120,-1.725) (-0.090,-1.723) (-0.060,-1.720) (-0.030,-1.716) (-0.000,-1.713) (0.030,-1.709) (0.060,-1.705) (0.090,-1.700) (0.120,-1.695) (0.150,-1.690) (0.180,-1.684) (0.210,-1.678) (0.240,-1.672) (0.270,-1.666) (0.300,-1.659) (0.330,-1.652) (0.360,-1.644) (0.390,-1.636) (0.420,-1.628) (0.449,-1.620) (0.480,-1.611) (0.510,-1.602) (0.540,-1.593) (0.570,-1.583) (0.600,-1.573) (0.630,-1.562) (0.660,-1.552) (0.690,-1.541) (0.718,-1.530) (0.750,-1.518) (0.780,-1.506) (0.810,-1.493) (0.840,-1.481) (0.870,-1.468) (0.898,-1.455) (0.930,-1.441) (0.960,-1.426) (0.990,-1.412) (1.020,-1.397) (1.050,-1.382) (1.080,-1.366) (1.110,-1.350) (1.138,-1.335) (1.164,-1.320) (1.191,-1.305) (1.230,-1.282) (1.260,-1.264) (1.290,-1.245) (1.320,-1.226) (1.350,-1.206) (1.380,-1.186) (1.410,-1.165) (1.440,-1.144) (1.466,-1.125) (1.500,-1.099) (1.525,-1.080) (1.560,-1.052) (1.590,-1.028) (1.617,-1.005) (1.650,-0.976) (1.680,-0.949) (1.710,-0.921) (1.740,-0.891) (1.770,-0.861) (1.800,-0.829) (1.830,-0.796) (1.857,-0.765) (1.882,-0.735) (1.905,-0.706) (1.929,-0.675) (1.950,-0.646) (1.971,-0.615) (1.991,-0.585) (2.010,-0.555) (2.028,-0.525) (2.044,-0.495) (2.060,-0.465) (2.075,-0.435) (2.089,-0.405) (2.101,-0.375) (2.113,-0.345) (2.124,-0.315) (2.134,-0.285) (2.143,-0.255) (2.151,-0.225) (2.158,-0.195) (2.163,-0.165) (2.168,-0.135) (2.172,-0.105) (2.175,-0.075) (2.177,-0.045) (2.178,-0.015) (2.178,0.030) (2.176,0.060) (2.174,0.090) (2.170,0.120) (2.166,0.150) (2.161,0.180) (2.154,0.210) (2.147,0.240) (2.138,0.270) (2.129,0.300) (2.119,0.330) (2.107,0.360) (2.095,0.390) (2.082,0.420) (2.068,0.450) (2.052,0.480) (2.036,0.510) (2.019,0.540) (2.001,0.570) (1.981,0.600) (1.961,0.630) (1.940,0.660) (1.920,0.687) (1.894,0.720) (1.875,0.744) (1.845,0.779) (1.818,0.810) (1.790,0.840) (1.761,0.870) (1.731,0.900) (1.700,0.930) (1.668,0.960) (1.635,0.989) (1.605,1.015) (1.581,1.035) (1.545,1.064) (1.515,1.088) (1.486,1.110) (1.455,1.133) (1.425,1.155) (1.395,1.176) (1.365,1.196) (1.337,1.215) (1.305,1.235) (1.275,1.254) (1.245,1.273) (1.216,1.290) (1.185,1.308) (1.155,1.325) (1.125,1.342) (1.095,1.358) (1.065,1.374) (1.035,1.390) (1.005,1.405) (0.975,1.419) (0.945,1.434) (0.915,1.447) (0.885,1.461) (0.855,1.474) (0.825,1.487) (0.795,1.500) (0.765,1.512) (0.735,1.524) (0.705,1.535) (0.675,1.546) (0.645,1.557) (0.615,1.568) (0.585,1.578) (0.555,1.588) (0.525,1.597) (0.495,1.607) (0.465,1.615) (0.435,1.624) (0.405,1.632) (0.375,1.640) (0.345,1.648) (0.315,1.655) (0.285,1.662) (0.255,1.669) (0.225,1.675) (0.195,1.681) (0.165,1.687) (0.135,1.693) (0.105,1.698) (0.075,1.702) (0.045,1.707) (0.015,1.711) (-0.015,1.715) (-0.045,1.718) (-0.075,1.721) (-0.105,1.724) (-0.135,1.726) (-0.165,1.728) (-0.195,1.730) (-0.225,1.731) (-0.255,1.732) (-0.285,1.733) (-0.315,1.733) (-0.345,1.733) (-0.375,1.733) (-0.405,1.732) (-0.435,1.731) (-0.465,1.729) (-0.495,1.727) (-0.525,1.724) (-0.555,1.722) (-0.585,1.718) (-0.615,1.715) (-0.645,1.710) (-0.675,1.706) (-0.705,1.701) (-0.735,1.695) (-0.765,1.690) (-0.795,1.683) (-0.825,1.676) (-0.855,1.669) (-0.885,1.661) (-0.915,1.653) (-0.945,1.644) (-0.975,1.635) (-1.005,1.625) (-1.035,1.615) (-1.065,1.604) (-1.095,1.592) (-1.125,1.580) (-1.155,1.567) (-1.185,1.554) (-1.215,1.540) (-1.245,1.525) (-1.275,1.510) (-1.305,1.494) (-1.335,1.477) (-1.365,1.459) (-1.395,1.441) (-1.425,1.421) (-1.455,1.401) (-1.485,1.380) (-1.515,1.358) (-1.545,1.335) (-1.575,1.311) (-1.601,1.290) (-1.635,1.260) (-1.665,1.232) (-1.695,1.203) (-1.725,1.173) (-1.755,1.141) (-1.783,1.110) (-1.808,1.080) (-1.830,1.053) (-1.856,1.020) (-1.875,0.994) (-1.899,0.960) (-1.919,0.930) (-1.939,0.900) (-1.957,0.870) (-1.975,0.840) (-1.992,0.810) (-2.008,0.780) (-2.023,0.750) (-2.037,0.720) (-2.051,0.690) (-2.064,0.660) (-2.076,0.630) (-2.088,0.600) (-2.099,0.570) (-2.110,0.540) (-2.119,0.510) (-2.129,0.480) (-2.137,0.450) (-2.145,0.420) (-2.152,0.390) (-2.159,0.360) (-2.166,0.330) (-2.171,0.300) (-2.176,0.270) (-2.181,0.240) (-2.185,0.210) (-2.188,0.180) (-2.191,0.150) (-2.194,0.120) (-2.196,0.090) (-2.197,0.060) (-2.198,0.030) (-2.198,-0.000) (-2.198,-0.030) (-2.197,-0.060) (-2.196,-0.090) (-2.194,-0.120) (-2.191,-0.150) (-2.188,-0.180) (-2.185,-0.210) (-2.181,-0.240) (-2.176,-0.270) (-2.171,-0.300) (-2.166,-0.330) (-2.159,-0.360) (-2.152,-0.390) (-2.145,-0.420) (-2.137,-0.450) (-2.129,-0.480) (-2.119,-0.510) (-2.110,-0.540) (-2.099,-0.570) (-2.088,-0.600) (-2.076,-0.630) (-2.064,-0.660) (-2.051,-0.690) (-2.037,-0.720) (-2.023,-0.750) (-2.008,-0.780) (-1.992,-0.810) (-1.975,-0.840) (-1.957,-0.870) (-1.939,-0.900) (-1.920,-0.929) (-1.899,-0.960) (-1.878,-0.990) (-1.860,-1.014) (-1.832,-1.050) (-1.808,-1.080) (-1.785,-1.107) (-1.756,-1.140) (-1.728,-1.170) (-1.698,-1.200) (-1.667,-1.230) (-1.635,-1.260) (-1.605,-1.286) (-1.575,-1.311) (-1.546,-1.335) (-1.515,-1.358) (-1.485,-1.380) (-1.455,-1.401) (-1.425,-1.421) (-1.396,-1.440) (-1.365,-1.459) (-1.335,-1.477) (-1.305,-1.494) (-1.275,-1.510) (-1.245,-1.525) (-1.215,-1.540) (-1.185,-1.554) (-1.155,-1.567) (-1.125,-1.580) (-1.095,-1.592) (-1.065,-1.604) (-1.035,-1.615) (-1.005,-1.625) (-0.975,-1.635) (-0.945,-1.644) (-0.915,-1.653) (-0.885,-1.661) (-0.855,-1.669) (-0.825,-1.676) (-0.795,-1.683) (-0.765,-1.690) (-0.735,-1.695) (-0.705,-1.701) (-0.675,-1.706) (-0.645,-1.710) (-0.615,-1.715) (-0.585,-1.718) (-0.555,-1.722) (-0.525,-1.724) (-0.510,-1.726)}
\EQ{(-0.525,-1.395) (-0.495,-1.397) (-0.465,-1.398) (-0.435,-1.398) (-0.405,-1.398) (-0.375,-1.398) (-0.345,-1.397) (-0.315,-1.396) (-0.285,-1.394) (-0.255,-1.392) (-0.225,-1.389) (-0.195,-1.386) (-0.165,-1.383) (-0.135,-1.379) (-0.105,-1.375) (-0.075,-1.370) (-0.045,-1.366) (-0.015,-1.360) (0.015,-1.355) (0.045,-1.349) (0.075,-1.342) (0.105,-1.336) (0.135,-1.329) (0.165,-1.321) (0.195,-1.314) (0.225,-1.306) (0.255,-1.298) (0.285,-1.289) (0.315,-1.280) (0.345,-1.271) (0.375,-1.262) (0.405,-1.252) (0.435,-1.242) (0.465,-1.232) (0.495,-1.221) (0.525,-1.210) (0.555,-1.199) (0.585,-1.188) (0.615,-1.176) (0.645,-1.164) (0.675,-1.152) (0.704,-1.140) (0.735,-1.127) (0.765,-1.113) (0.795,-1.100) (0.825,-1.086) (0.855,-1.072) (0.885,-1.058) (0.915,-1.043) (0.945,-1.028) (0.975,-1.013) (1.005,-0.998) (1.035,-0.982) (1.065,-0.965) (1.095,-0.949) (1.125,-0.932) (1.154,-0.915) (1.185,-0.897) (1.215,-0.878) (1.245,-0.860) (1.275,-0.841) (1.305,-0.821) (1.335,-0.801) (1.365,-0.781) (1.395,-0.760) (1.425,-0.738) (1.449,-0.720) (1.485,-0.693) (1.515,-0.669) (1.544,-0.645) (1.575,-0.619) (1.605,-0.593) (1.630,-0.570) (1.662,-0.540) (1.692,-0.510) (1.720,-0.480) (1.747,-0.450) (1.770,-0.424) (1.797,-0.390) (1.819,-0.360) (1.840,-0.330) (1.859,-0.300) (1.875,-0.273) (1.890,-0.245) (1.905,-0.213) (1.919,-0.180) (1.929,-0.150) (1.938,-0.120) (1.945,-0.090) (1.949,-0.060) (1.952,-0.030) (1.953,-0.000) (1.952,0.030) (1.949,0.060) (1.945,0.090) (1.938,0.120) (1.929,0.150) (1.919,0.180) (1.907,0.210) (1.892,0.240) (1.877,0.270) (1.860,0.299) (1.840,0.330) (1.819,0.360) (1.800,0.386) (1.773,0.420) (1.747,0.450) (1.725,0.475) (1.695,0.507) (1.665,0.537) (1.635,0.565) (1.605,0.593) (1.580,0.615) (1.545,0.645) (1.515,0.669) (1.489,0.690) (1.455,0.716) (1.429,0.735) (1.395,0.760) (1.366,0.780) (1.335,0.801) (1.305,0.821) (1.276,0.840) (1.245,0.860) (1.215,0.878) (1.185,0.897) (1.155,0.914) (1.128,0.930) (1.102,0.945) (1.075,0.960) (1.047,0.975) (1.020,0.990) (0.991,1.005) (0.962,1.020) (0.932,1.035) (0.902,1.050) (0.870,1.065) (0.840,1.079) (0.810,1.093) (0.780,1.107) (0.750,1.120) (0.720,1.133) (0.690,1.146) (0.660,1.158) (0.631,1.170) (0.600,1.182) (0.570,1.194) (0.540,1.205) (0.510,1.216) (0.480,1.227) (0.450,1.237) (0.420,1.247) (0.390,1.257) (0.360,1.267) (0.330,1.276) (0.300,1.285) (0.270,1.293) (0.240,1.302) (0.210,1.310) (0.180,1.318) (0.150,1.325) (0.120,1.332) (0.090,1.339) (0.060,1.346) (0.030,1.352) (-0.000,1.358) (-0.030,1.363) (-0.060,1.368) (-0.090,1.373) (-0.120,1.377) (-0.150,1.381) (-0.180,1.385) (-0.210,1.388) (-0.240,1.391) (-0.270,1.393) (-0.300,1.395) (-0.330,1.397) (-0.360,1.398) (-0.390,1.398) (-0.420,1.398) (-0.450,1.398) (-0.480,1.397) (-0.510,1.396) (-0.540,1.394) (-0.570,1.392) (-0.600,1.389) (-0.630,1.386) (-0.660,1.382) (-0.690,1.378) (-0.720,1.373) (-0.750,1.368) (-0.780,1.362) (-0.810,1.355) (-0.840,1.348) (-0.870,1.340) (-0.900,1.331) (-0.930,1.322) (-0.960,1.312) (-0.990,1.302) (-1.020,1.291) (-1.050,1.279) (-1.080,1.266) (-1.110,1.252) (-1.140,1.238) (-1.170,1.223) (-1.200,1.207) (-1.230,1.190) (-1.260,1.172) (-1.286,1.155) (-1.320,1.132) (-1.350,1.111) (-1.380,1.089) (-1.410,1.065) (-1.440,1.040) (-1.470,1.013) (-1.495,0.990) (-1.525,0.960) (-1.554,0.930) (-1.575,0.906) (-1.605,0.871) (-1.630,0.840) (-1.650,0.813) (-1.674,0.780) (-1.694,0.750) (-1.713,0.720) (-1.731,0.690) (-1.748,0.660) (-1.764,0.630) (-1.779,0.600) (-1.794,0.570) (-1.807,0.540) (-1.819,0.510) (-1.831,0.480) (-1.842,0.450) (-1.852,0.420) (-1.861,0.390) (-1.869,0.360) (-1.877,0.330) (-1.884,0.300) (-1.891,0.270) (-1.896,0.240) (-1.901,0.210) (-1.905,0.180) (-1.909,0.150) (-1.912,0.120) (-1.914,0.090) (-1.916,0.060) (-1.917,0.030) (-1.917,-0.000) (-1.917,-0.030) (-1.916,-0.060) (-1.914,-0.090) (-1.912,-0.120) (-1.909,-0.150) (-1.905,-0.180) (-1.901,-0.210) (-1.896,-0.240) (-1.891,-0.270) (-1.884,-0.300) (-1.877,-0.330) (-1.869,-0.360) (-1.861,-0.390) (-1.852,-0.420) (-1.842,-0.450) (-1.831,-0.480) (-1.819,-0.510) (-1.807,-0.540) (-1.794,-0.570) (-1.779,-0.600) (-1.764,-0.630) (-1.748,-0.660) (-1.731,-0.690) (-1.713,-0.720) (-1.695,-0.749) (-1.674,-0.780) (-1.653,-0.810) (-1.630,-0.840) (-1.606,-0.870) (-1.580,-0.900) (-1.560,-0.923) (-1.539,-0.945) (-1.515,-0.970) (-1.485,-0.999) (-1.462,-1.020) (-1.428,-1.050) (-1.395,-1.077) (-1.365,-1.100) (-1.335,-1.122) (-1.309,-1.140) (-1.275,-1.162) (-1.245,-1.181) (-1.215,-1.198) (-1.185,-1.215) (-1.156,-1.230) (-1.125,-1.245) (-1.095,-1.259) (-1.065,-1.272) (-1.035,-1.285) (-1.005,-1.296) (-0.975,-1.307) (-0.945,-1.317) (-0.915,-1.327) (-0.885,-1.336) (-0.855,-1.344) (-0.825,-1.351) (-0.795,-1.358) (-0.765,-1.365) (-0.735,-1.370) (-0.705,-1.376) (-0.675,-1.380) (-0.645,-1.384) (-0.615,-1.388) (-0.585,-1.391) (-0.555,-1.393) (-0.525,-1.395)}
\EQ{(-0.660,-1.111) (-0.630,-1.113) (-0.600,-1.115) (-0.570,-1.117) (-0.540,-1.118) (-0.510,-1.118) (-0.480,-1.117) (-0.450,-1.116) (-0.420,-1.115) (-0.390,-1.112) (-0.360,-1.109) (-0.330,-1.106) (-0.300,-1.102) (-0.270,-1.097) (-0.240,-1.092) (-0.210,-1.086) (-0.180,-1.080) (-0.150,-1.073) (-0.120,-1.066) (-0.090,-1.059) (-0.060,-1.051) (-0.030,-1.043) (-0.000,-1.035) (0.030,-1.026) (0.060,-1.017) (0.090,-1.008) (0.120,-0.998) (0.150,-0.988) (0.180,-0.978) (0.210,-0.968) (0.240,-0.958) (0.270,-0.947) (0.300,-0.936) (0.330,-0.925) (0.360,-0.914) (0.390,-0.903) (0.420,-0.891) (0.450,-0.880) (0.480,-0.868) (0.510,-0.856) (0.540,-0.844) (0.570,-0.831) (0.600,-0.819) (0.630,-0.806) (0.660,-0.793) (0.690,-0.780) (0.720,-0.767) (0.750,-0.753) (0.780,-0.740) (0.810,-0.726) (0.840,-0.712) (0.870,-0.698) (0.900,-0.684) (0.930,-0.669) (0.960,-0.654) (0.990,-0.640) (1.020,-0.625) (1.050,-0.609) (1.080,-0.594) (1.110,-0.578) (1.140,-0.562) (1.170,-0.546) (1.200,-0.529) (1.230,-0.513) (1.260,-0.496) (1.287,-0.480) (1.320,-0.460) (1.350,-0.442) (1.380,-0.424) (1.410,-0.405) (1.440,-0.386) (1.470,-0.366) (1.500,-0.346) (1.530,-0.325) (1.560,-0.303) (1.590,-0.280) (1.620,-0.256) (1.650,-0.231) (1.674,-0.210) (1.705,-0.180) (1.733,-0.150) (1.755,-0.123) (1.777,-0.090) (1.792,-0.060) (1.801,-0.030) (1.805,-0.000) (1.801,0.030) (1.792,0.060) (1.777,0.090) (1.757,0.120) (1.733,0.150) (1.710,0.175) (1.680,0.204) (1.657,0.225) (1.621,0.255) (1.590,0.280) (1.564,0.300) (1.530,0.325) (1.501,0.345) (1.470,0.366) (1.440,0.386) (1.410,0.405) (1.380,0.424) (1.350,0.442) (1.320,0.460) (1.290,0.478) (1.261,0.495) (1.235,0.510) (1.208,0.525) (1.181,0.540) (1.153,0.555) (1.125,0.570) (1.097,0.585) (1.068,0.600) (1.039,0.615) (1.009,0.630) (0.979,0.645) (0.949,0.660) (0.918,0.675) (0.887,0.690) (0.855,0.705) (0.825,0.719) (0.795,0.733) (0.765,0.746) (0.735,0.760) (0.705,0.773) (0.675,0.786) (0.645,0.799) (0.615,0.812) (0.585,0.825) (0.555,0.837) (0.525,0.850) (0.495,0.862) (0.465,0.874) (0.436,0.885) (0.405,0.897) (0.375,0.909) (0.345,0.920) (0.315,0.931) (0.285,0.942) (0.255,0.953) (0.225,0.963) (0.195,0.973) (0.165,0.983) (0.135,0.993) (0.105,1.003) (0.075,1.012) (0.045,1.021) (0.015,1.030) (-0.015,1.039) (-0.045,1.047) (-0.075,1.055) (-0.105,1.063) (-0.135,1.070) (-0.165,1.077) (-0.195,1.083) (-0.225,1.089) (-0.255,1.094) (-0.285,1.099) (-0.315,1.104) (-0.345,1.108) (-0.375,1.111) (-0.405,1.113) (-0.435,1.115) (-0.465,1.117) (-0.495,1.118) (-0.525,1.118) (-0.555,1.117) (-0.585,1.116) (-0.615,1.114) (-0.645,1.112) (-0.675,1.109) (-0.705,1.105) (-0.735,1.101) (-0.765,1.095) (-0.795,1.089) (-0.825,1.083) (-0.855,1.075) (-0.885,1.067) (-0.915,1.057) (-0.945,1.047) (-0.975,1.036) (-1.005,1.023) (-1.035,1.010) (-1.065,0.996) (-1.095,0.980) (-1.125,0.964) (-1.155,0.946) (-1.185,0.927) (-1.215,0.906) (-1.244,0.885) (-1.275,0.861) (-1.300,0.840) (-1.333,0.810) (-1.364,0.780) (-1.393,0.750) (-1.419,0.720) (-1.440,0.695) (-1.467,0.660) (-1.485,0.636) (-1.509,0.600) (-1.528,0.570) (-1.545,0.540) (-1.560,0.513) (-1.575,0.483) (-1.590,0.450) (-1.603,0.420) (-1.615,0.390) (-1.625,0.360) (-1.635,0.330) (-1.644,0.300) (-1.652,0.270) (-1.659,0.240) (-1.665,0.210) (-1.670,0.180) (-1.674,0.150) (-1.677,0.120) (-1.680,0.090) (-1.682,0.060) (-1.683,0.030) (-1.683,-0.000) (-1.683,-0.030) (-1.682,-0.060) (-1.680,-0.090) (-1.677,-0.120) (-1.674,-0.150) (-1.670,-0.180) (-1.665,-0.210) (-1.659,-0.240) (-1.652,-0.270) (-1.644,-0.300) (-1.635,-0.330) (-1.625,-0.360) (-1.615,-0.390) (-1.603,-0.420) (-1.590,-0.450) (-1.576,-0.480) (-1.561,-0.510) (-1.545,-0.540) (-1.530,-0.566) (-1.509,-0.600) (-1.489,-0.630) (-1.470,-0.656) (-1.444,-0.690) (-1.425,-0.713) (-1.395,-0.747) (-1.365,-0.779) (-1.335,-0.808) (-1.305,-0.835) (-1.282,-0.855) (-1.245,-0.884) (-1.215,-0.906) (-1.185,-0.927) (-1.157,-0.945) (-1.125,-0.964) (-1.095,-0.980) (-1.065,-0.996) (-1.035,-1.010) (-1.005,-1.023) (-0.976,-1.035) (-0.945,-1.047) (-0.915,-1.057) (-0.885,-1.067) (-0.855,-1.075) (-0.825,-1.083) (-0.795,-1.089) (-0.765,-1.095) (-0.735,-1.101) (-0.705,-1.105) (-0.675,-1.109) (-0.660,-1.111)}

      % Conductor con punta - líneas de campo
% Generado por scripts/generar_datos.py (no editar a mano).
\LC{(-0.388,0.905) (-0.381,0.941) (-0.375,0.976) (-0.370,1.012) (-0.365,1.048) (-0.360,1.083) (-0.356,1.119) (-0.353,1.155) (-0.349,1.191) (-0.347,1.227) (-0.344,1.263) (-0.342,1.298) (-0.340,1.334) (-0.338,1.370) (-0.337,1.406) (-0.335,1.442) (-0.334,1.478) (-0.333,1.514) (-0.333,1.550) (-0.332,1.586) (-0.332,1.622) (-0.332,1.658) (-0.332,1.694) (-0.332,1.730) (-0.332,1.766) (-0.332,1.802) (-0.333,1.838) (-0.333,1.874) (-0.334,1.910) (-0.335,1.946) (-0.336,1.982) (-0.336,2.000)}
\LC{(-0.651,0.929) (-0.653,0.965) (-0.655,1.000) (-0.658,1.036) (-0.661,1.072) (-0.665,1.108) (-0.669,1.144) (-0.673,1.180) (-0.678,1.215) (-0.683,1.251) (-0.688,1.287) (-0.693,1.322) (-0.698,1.358) (-0.704,1.393) (-0.710,1.429) (-0.716,1.464) (-0.722,1.500) (-0.729,1.535) (-0.735,1.571) (-0.742,1.606) (-0.749,1.641) (-0.756,1.677) (-0.763,1.712) (-0.770,1.747) (-0.777,1.782) (-0.785,1.818) (-0.792,1.853) (-0.800,1.888) (-0.808,1.923) (-0.815,1.958) (-0.823,1.993) (-0.826,2.000)}
\LC{(-0.888,0.884) (-0.899,0.918) (-0.911,0.953) (-0.923,0.987) (-0.934,1.021) (-0.947,1.054) (-0.959,1.088) (-0.971,1.122) (-0.984,1.156) (-0.997,1.189) (-1.010,1.223) (-1.023,1.257) (-1.036,1.290) (-1.049,1.323) (-1.063,1.357) (-1.077,1.390) (-1.091,1.423) (-1.104,1.457) (-1.118,1.490) (-1.133,1.523) (-1.147,1.556) (-1.161,1.589) (-1.176,1.622) (-1.190,1.655) (-1.205,1.688) (-1.219,1.721) (-1.234,1.753) (-1.249,1.786) (-1.264,1.819) (-1.279,1.852) (-1.294,1.884) (-1.309,1.917) (-1.324,1.950) (-1.339,1.983) (-1.349,2.000)}
\LC{(-1.097,0.786) (-1.117,0.816) (-1.136,0.846) (-1.156,0.876) (-1.176,0.907) (-1.196,0.937) (-1.216,0.966) (-1.236,0.996) (-1.256,1.026) (-1.276,1.056) (-1.297,1.085) (-1.317,1.115) (-1.338,1.145) (-1.359,1.174) (-1.379,1.203) (-1.400,1.233) (-1.421,1.262) (-1.442,1.291) (-1.463,1.320) (-1.485,1.350) (-1.506,1.379) (-1.527,1.408) (-1.548,1.437) (-1.570,1.466) (-1.591,1.494) (-1.613,1.523) (-1.635,1.552) (-1.656,1.581) (-1.678,1.609) (-1.700,1.638) (-1.722,1.667) (-1.743,1.695) (-1.765,1.724) (-1.787,1.753) (-1.809,1.781) (-1.831,1.809) (-1.853,1.838) (-1.875,1.866) (-1.898,1.895) (-1.920,1.923) (-1.942,1.951) (-1.964,1.980) (-1.987,2.000)}
\LC{(-1.274,0.641) (-1.300,0.666) (-1.326,0.690) (-1.352,0.715) (-1.379,0.740) (-1.405,0.764) (-1.432,0.788) (-1.458,0.813) (-1.485,0.837) (-1.512,0.861) (-1.538,0.885) (-1.565,0.909) (-1.592,0.933) (-1.619,0.957) (-1.646,0.981) (-1.673,1.005) (-1.700,1.028) (-1.727,1.052) (-1.754,1.076) (-1.781,1.099) (-1.809,1.123) (-1.836,1.146) (-1.863,1.170) (-1.891,1.193) (-1.918,1.216) (-1.946,1.240) (-1.973,1.263) (-2.001,1.286) (-2.028,1.309) (-2.056,1.332) (-2.084,1.355) (-2.112,1.378) (-2.139,1.401) (-2.167,1.424) (-2.195,1.447) (-2.223,1.469) (-2.251,1.492) (-2.279,1.515) (-2.307,1.537) (-2.335,1.560) (-2.363,1.582) (-2.391,1.605) (-2.419,1.627) (-2.447,1.650) (-2.475,1.672) (-2.504,1.695) (-2.532,1.717) (-2.560,1.739) (-2.589,1.761) (-2.600,1.776)}
\LC{(-1.407,0.462) (-1.438,0.481) (-1.469,0.498) (-1.501,0.516) (-1.532,0.534) (-1.563,0.551) (-1.595,0.569) (-1.626,0.586) (-1.658,0.604) (-1.689,0.621) (-1.721,0.638) (-1.753,0.656) (-1.784,0.673) (-1.816,0.690) (-1.847,0.707) (-1.879,0.724) (-1.911,0.741) (-1.943,0.758) (-1.975,0.775) (-2.006,0.792) (-2.038,0.809) (-2.070,0.825) (-2.102,0.842) (-2.134,0.859) (-2.166,0.875) (-2.198,0.892) (-2.230,0.908) (-2.262,0.925) (-2.294,0.941) (-2.326,0.957) (-2.358,0.974) (-2.390,0.990) (-2.422,1.006) (-2.454,1.022) (-2.487,1.038) (-2.519,1.055) (-2.551,1.071) (-2.583,1.086) (-2.600,1.097)}
\LC{(-1.493,0.261) (-1.527,0.271) (-1.562,0.280) (-1.597,0.290) (-1.631,0.300) (-1.666,0.310) (-1.701,0.319) (-1.735,0.329) (-1.770,0.339) (-1.805,0.348) (-1.839,0.358) (-1.874,0.367) (-1.909,0.377) (-1.943,0.386) (-1.978,0.396) (-2.013,0.405) (-2.048,0.415) (-2.082,0.424) (-2.117,0.434) (-2.152,0.443) (-2.187,0.452) (-2.222,0.462) (-2.256,0.471) (-2.291,0.480) (-2.326,0.489) (-2.361,0.498) (-2.396,0.507) (-2.430,0.517) (-2.465,0.526) (-2.500,0.535) (-2.535,0.544) (-2.570,0.553) (-2.600,0.562)}
\LC{(-1.529,0.044) (-1.565,0.045) (-1.601,0.046) (-1.637,0.048) (-1.673,0.049) (-1.709,0.050) (-1.745,0.052) (-1.781,0.053) (-1.817,0.055) (-1.853,0.057) (-1.889,0.058) (-1.925,0.060) (-1.961,0.061) (-1.997,0.063) (-2.033,0.064) (-2.068,0.066) (-2.104,0.067) (-2.140,0.069) (-2.176,0.070) (-2.212,0.072) (-2.248,0.073) (-2.284,0.075) (-2.320,0.076) (-2.356,0.078) (-2.392,0.079) (-2.428,0.081) (-2.464,0.082) (-2.500,0.084) (-2.536,0.085) (-2.572,0.087) (-2.600,0.088)}
\LC{(-1.512,-0.181) (-1.548,-0.187) (-1.583,-0.194) (-1.618,-0.200) (-1.654,-0.207) (-1.689,-0.213) (-1.725,-0.220) (-1.760,-0.226) (-1.795,-0.233) (-1.831,-0.240) (-1.866,-0.246) (-1.902,-0.253) (-1.937,-0.259) (-1.972,-0.266) (-2.008,-0.272) (-2.043,-0.279) (-2.079,-0.285) (-2.114,-0.292) (-2.150,-0.298) (-2.185,-0.305) (-2.220,-0.311) (-2.256,-0.317) (-2.291,-0.324) (-2.327,-0.330) (-2.362,-0.336) (-2.398,-0.343) (-2.433,-0.349) (-2.468,-0.355) (-2.504,-0.361) (-2.539,-0.367) (-2.575,-0.374) (-2.600,-0.380)}
\LC{(-1.445,-0.389) (-1.478,-0.404) (-1.511,-0.418) (-1.543,-0.433) (-1.576,-0.448) (-1.609,-0.462) (-1.642,-0.477) (-1.675,-0.492) (-1.708,-0.506) (-1.741,-0.521) (-1.774,-0.535) (-1.807,-0.549) (-1.840,-0.564) (-1.873,-0.578) (-1.906,-0.592) (-1.939,-0.607) (-1.972,-0.621) (-2.005,-0.635) (-2.038,-0.649) (-2.071,-0.663) (-2.105,-0.677) (-2.138,-0.691) (-2.171,-0.705) (-2.204,-0.719) (-2.237,-0.733) (-2.271,-0.747) (-2.304,-0.760) (-2.337,-0.774) (-2.370,-0.788) (-2.404,-0.801) (-2.437,-0.815) (-2.471,-0.828) (-2.504,-0.842) (-2.537,-0.855) (-2.571,-0.869) (-2.600,-0.882)}
\LC{(-1.330,-0.576) (-1.358,-0.599) (-1.386,-0.621) (-1.415,-0.643) (-1.443,-0.665) (-1.472,-0.687) (-1.500,-0.709) (-1.529,-0.731) (-1.558,-0.753) (-1.586,-0.774) (-1.615,-0.796) (-1.644,-0.818) (-1.673,-0.839) (-1.702,-0.861) (-1.731,-0.882) (-1.760,-0.903) (-1.789,-0.925) (-1.818,-0.946) (-1.847,-0.967) (-1.876,-0.988) (-1.905,-1.009) (-1.934,-1.030) (-1.964,-1.051) (-1.993,-1.072) (-2.022,-1.093) (-2.052,-1.114) (-2.081,-1.134) (-2.111,-1.155) (-2.140,-1.176) (-2.170,-1.196) (-2.199,-1.217) (-2.229,-1.237) (-2.258,-1.258) (-2.288,-1.278) (-2.318,-1.299) (-2.347,-1.319) (-2.377,-1.339) (-2.407,-1.359) (-2.437,-1.380) (-2.467,-1.400) (-2.497,-1.420) (-2.526,-1.440) (-2.556,-1.460) (-2.586,-1.480) (-2.600,-1.493)}
\LC{(-1.170,-0.735) (-1.192,-0.763) (-1.214,-0.791) (-1.237,-0.820) (-1.260,-0.848) (-1.282,-0.876) (-1.305,-0.904) (-1.328,-0.931) (-1.351,-0.959) (-1.374,-0.987) (-1.397,-1.015) (-1.420,-1.042) (-1.443,-1.070) (-1.466,-1.097) (-1.489,-1.125) (-1.513,-1.152) (-1.536,-1.179) (-1.560,-1.206) (-1.583,-1.234) (-1.607,-1.261) (-1.631,-1.288) (-1.655,-1.315) (-1.678,-1.342) (-1.702,-1.369) (-1.726,-1.396) (-1.750,-1.423) (-1.774,-1.449) (-1.798,-1.476) (-1.823,-1.503) (-1.847,-1.529) (-1.871,-1.556) (-1.895,-1.583) (-1.920,-1.609) (-1.944,-1.636) (-1.968,-1.662) (-1.993,-1.689) (-2.017,-1.715) (-2.042,-1.741) (-2.066,-1.768) (-2.091,-1.794) (-2.116,-1.820) (-2.140,-1.846) (-2.165,-1.873) (-2.190,-1.899) (-2.215,-1.925) (-2.239,-1.951) (-2.264,-1.977) (-2.289,-2.000)}
\LC{(-0.973,-0.852) (-0.988,-0.885) (-1.002,-0.918) (-1.017,-0.951) (-1.032,-0.983) (-1.047,-1.016) (-1.063,-1.048) (-1.078,-1.081) (-1.094,-1.113) (-1.110,-1.146) (-1.126,-1.178) (-1.142,-1.210) (-1.158,-1.242) (-1.174,-1.274) (-1.191,-1.306) (-1.207,-1.338) (-1.224,-1.370) (-1.241,-1.402) (-1.258,-1.434) (-1.275,-1.466) (-1.292,-1.497) (-1.309,-1.529) (-1.326,-1.561) (-1.343,-1.592) (-1.360,-1.624) (-1.378,-1.655) (-1.395,-1.687) (-1.413,-1.718) (-1.431,-1.750) (-1.448,-1.781) (-1.466,-1.812) (-1.484,-1.844) (-1.501,-1.875) (-1.519,-1.906) (-1.537,-1.937) (-1.555,-1.969) (-1.573,-2.000) (-1.579,-2.000)}
\LC{(-0.745,-0.919) (-0.751,-0.954) (-0.756,-0.990) (-0.763,-1.025) (-0.769,-1.061) (-0.776,-1.096) (-0.783,-1.131) (-0.790,-1.166) (-0.798,-1.202) (-0.806,-1.237) (-0.814,-1.272) (-0.822,-1.307) (-0.831,-1.342) (-0.839,-1.377) (-0.848,-1.412) (-0.857,-1.447) (-0.866,-1.481) (-0.876,-1.516) (-0.885,-1.551) (-0.895,-1.586) (-0.904,-1.620) (-0.914,-1.655) (-0.924,-1.690) (-0.934,-1.724) (-0.944,-1.759) (-0.954,-1.793) (-0.964,-1.828) (-0.975,-1.862) (-0.985,-1.897) (-0.996,-1.931) (-1.006,-1.965) (-1.017,-2.000) (-1.021,-2.000)}
\LC{(-0.487,-0.923) (-0.485,-0.959) (-0.482,-0.995) (-0.480,-1.031) (-0.479,-1.067) (-0.477,-1.103) (-0.476,-1.139) (-0.476,-1.175) (-0.475,-1.211) (-0.475,-1.247) (-0.476,-1.283) (-0.476,-1.319) (-0.477,-1.355) (-0.478,-1.391) (-0.479,-1.427) (-0.481,-1.463) (-0.482,-1.499) (-0.484,-1.535) (-0.486,-1.571) (-0.488,-1.606) (-0.491,-1.642) (-0.493,-1.678) (-0.496,-1.714) (-0.499,-1.750) (-0.501,-1.786) (-0.504,-1.822) (-0.507,-1.858) (-0.511,-1.894) (-0.514,-1.929) (-0.517,-1.965) (-0.521,-2.000)}
\LC{(-0.215,-0.847) (-0.203,-0.881) (-0.192,-0.915) (-0.182,-0.950) (-0.172,-0.985) (-0.163,-1.019) (-0.155,-1.054) (-0.147,-1.090) (-0.139,-1.125) (-0.132,-1.160) (-0.125,-1.195) (-0.119,-1.231) (-0.112,-1.266) (-0.107,-1.302) (-0.101,-1.337) (-0.095,-1.373) (-0.090,-1.409) (-0.085,-1.444) (-0.080,-1.480) (-0.076,-1.516) (-0.071,-1.551) (-0.067,-1.587) (-0.063,-1.623) (-0.059,-1.659) (-0.055,-1.694) (-0.051,-1.730) (-0.047,-1.766) (-0.044,-1.802) (-0.040,-1.838) (-0.037,-1.873) (-0.033,-1.909) (-0.030,-1.945) (-0.027,-1.981) (-0.025,-2.000)}
\LC{(0.096,-0.718) (0.109,-0.752) (0.122,-0.785) (0.135,-0.819) (0.147,-0.853) (0.160,-0.887) (0.172,-0.920) (0.184,-0.954) (0.195,-0.989) (0.206,-1.023) (0.218,-1.057) (0.229,-1.091) (0.239,-1.126) (0.250,-1.160) (0.260,-1.194) (0.271,-1.229) (0.281,-1.263) (0.291,-1.298) (0.301,-1.333) (0.311,-1.367) (0.321,-1.402) (0.330,-1.437) (0.340,-1.471) (0.349,-1.506) (0.359,-1.541) (0.368,-1.576) (0.377,-1.610) (0.386,-1.645) (0.396,-1.680) (0.405,-1.715) (0.414,-1.750) (0.423,-1.785) (0.432,-1.819) (0.441,-1.854) (0.449,-1.889) (0.458,-1.924) (0.467,-1.959) (0.476,-1.994) (0.479,-2.000)}
\LC{(0.432,-0.579) (0.446,-0.612) (0.460,-0.645) (0.474,-0.678) (0.488,-0.712) (0.501,-0.745) (0.515,-0.778) (0.529,-0.812) (0.542,-0.845) (0.556,-0.878) (0.570,-0.912) (0.583,-0.945) (0.597,-0.978) (0.611,-1.011) (0.624,-1.045) (0.638,-1.078) (0.652,-1.111) (0.665,-1.145) (0.679,-1.178) (0.693,-1.211) (0.706,-1.245) (0.720,-1.278) (0.734,-1.311) (0.748,-1.344) (0.761,-1.378) (0.775,-1.411) (0.789,-1.444) (0.803,-1.477) (0.817,-1.511) (0.831,-1.544) (0.845,-1.577) (0.858,-1.610) (0.872,-1.643) (0.886,-1.677) (0.900,-1.710) (0.914,-1.743) (0.928,-1.776) (0.942,-1.809) (0.956,-1.842) (0.970,-1.876) (0.984,-1.909) (0.998,-1.942) (1.013,-1.975) (1.027,-2.000)}
\LC{(0.775,-0.437) (0.789,-0.470) (0.804,-0.503) (0.818,-0.536) (0.833,-0.569) (0.848,-0.601) (0.863,-0.634) (0.878,-0.667) (0.894,-0.699) (0.909,-0.732) (0.925,-0.764) (0.941,-0.796) (0.957,-0.828) (0.974,-0.860) (0.990,-0.892) (1.007,-0.924) (1.024,-0.956) (1.041,-0.988) (1.058,-1.019) (1.075,-1.051) (1.092,-1.083) (1.110,-1.114) (1.128,-1.145) (1.146,-1.177) (1.164,-1.208) (1.182,-1.239) (1.200,-1.270) (1.218,-1.301) (1.236,-1.332) (1.255,-1.363) (1.273,-1.394) (1.292,-1.425) (1.311,-1.455) (1.330,-1.486) (1.349,-1.517) (1.368,-1.547) (1.387,-1.578) (1.406,-1.608) (1.425,-1.639) (1.444,-1.669) (1.464,-1.699) (1.483,-1.730) (1.503,-1.760) (1.522,-1.790) (1.542,-1.820) (1.561,-1.851) (1.581,-1.881) (1.601,-1.911) (1.621,-1.941) (1.640,-1.971) (1.660,-2.000)}
\LC{(1.098,-0.303) (1.113,-0.336) (1.128,-0.368) (1.143,-0.401) (1.159,-0.433) (1.176,-0.465) (1.192,-0.497) (1.210,-0.529) (1.228,-0.560) (1.246,-0.591) (1.264,-0.622) (1.283,-0.652) (1.303,-0.683) (1.322,-0.713) (1.342,-0.743) (1.363,-0.773) (1.383,-0.802) (1.404,-0.831) (1.425,-0.860) (1.447,-0.889) (1.468,-0.918) (1.490,-0.947) (1.512,-0.975) (1.535,-1.003) (1.557,-1.031) (1.580,-1.059) (1.603,-1.087) (1.626,-1.114) (1.650,-1.142) (1.673,-1.169) (1.697,-1.196) (1.720,-1.224) (1.744,-1.250) (1.768,-1.277) (1.792,-1.304) (1.817,-1.331) (1.841,-1.357) (1.865,-1.384) (1.890,-1.410) (1.915,-1.436) (1.939,-1.462) (1.964,-1.488) (1.989,-1.514) (2.014,-1.540) (2.039,-1.566) (2.065,-1.592) (2.090,-1.617) (2.115,-1.643) (2.141,-1.668) (2.166,-1.694) (2.192,-1.719) (2.217,-1.744) (2.243,-1.770) (2.269,-1.795) (2.294,-1.820) (2.320,-1.845) (2.346,-1.870) (2.372,-1.895) (2.398,-1.920) (2.424,-1.945) (2.450,-1.970) (2.476,-1.994) (2.485,-2.000)}
\LC{(1.378,-0.187) (1.393,-0.219) (1.410,-0.251) (1.427,-0.283) (1.445,-0.314) (1.465,-0.344) (1.485,-0.374) (1.505,-0.404) (1.527,-0.432) (1.549,-0.461) (1.572,-0.488) (1.596,-0.516) (1.620,-0.543) (1.644,-0.569) (1.669,-0.595) (1.694,-0.620) (1.720,-0.646) (1.746,-0.670) (1.773,-0.695) (1.800,-0.719) (1.827,-0.743) (1.854,-0.766) (1.881,-0.789) (1.909,-0.812) (1.937,-0.835) (1.965,-0.858) (1.993,-0.880) (2.022,-0.902) (2.050,-0.924) (2.079,-0.945) (2.108,-0.967) (2.137,-0.988) (2.166,-1.009) (2.195,-1.030) (2.225,-1.051) (2.254,-1.072) (2.284,-1.092) (2.313,-1.113) (2.343,-1.133) (2.373,-1.153) (2.403,-1.173) (2.433,-1.193) (2.463,-1.213) (2.493,-1.233) (2.523,-1.252) (2.554,-1.272) (2.584,-1.291) (2.600,-1.304)}
\LC{(1.600,-0.095) (1.617,-0.126) (1.637,-0.156) (1.660,-0.184) (1.684,-0.211) (1.709,-0.237) (1.736,-0.261) (1.764,-0.284) (1.792,-0.306) (1.821,-0.327) (1.851,-0.347) (1.881,-0.366) (1.912,-0.385) (1.943,-0.403) (1.975,-0.421) (2.006,-0.438) (2.038,-0.455) (2.070,-0.471) (2.102,-0.487) (2.135,-0.503) (2.167,-0.519) (2.200,-0.534) (2.233,-0.549) (2.265,-0.563) (2.298,-0.578) (2.331,-0.592) (2.365,-0.606) (2.398,-0.620) (2.431,-0.634) (2.464,-0.647) (2.498,-0.661) (2.531,-0.674) (2.565,-0.687) (2.598,-0.700) (2.600,-0.704)}
\LC{(1.739,-0.037) (1.769,-0.057) (1.801,-0.073) (1.835,-0.085) (1.869,-0.096) (1.904,-0.106) (1.939,-0.115) (1.974,-0.123) (2.009,-0.131) (2.044,-0.138) (2.080,-0.145) (2.115,-0.152) (2.150,-0.159) (2.186,-0.165) (2.221,-0.171) (2.257,-0.177) (2.292,-0.182) (2.328,-0.188) (2.364,-0.193) (2.399,-0.199) (2.435,-0.204) (2.470,-0.209) (2.506,-0.214) (2.542,-0.219) (2.577,-0.224) (2.600,-0.227)}
\LC{(1.657,0.071) (1.677,0.101) (1.700,0.128) (1.726,0.154) (1.754,0.177) (1.782,0.198) (1.812,0.218) (1.843,0.237) (1.875,0.254) (1.907,0.271) (1.939,0.287) (1.971,0.302) (2.004,0.317) (2.037,0.331) (2.071,0.345) (2.104,0.358) (2.138,0.371) (2.171,0.384) (2.205,0.397) (2.239,0.409) (2.273,0.421) (2.307,0.432) (2.341,0.444) (2.375,0.455) (2.409,0.466) (2.444,0.477) (2.478,0.488) (2.512,0.499) (2.547,0.509) (2.581,0.520) (2.600,0.527)}
\LC{(1.466,0.150) (1.482,0.182) (1.499,0.214) (1.518,0.245) (1.538,0.275) (1.559,0.304) (1.580,0.333) (1.603,0.361) (1.627,0.388) (1.651,0.414) (1.677,0.440) (1.702,0.465) (1.729,0.490) (1.755,0.514) (1.783,0.537) (1.810,0.561) (1.838,0.583) (1.866,0.605) (1.895,0.627) (1.924,0.649) (1.953,0.670) (1.982,0.691) (2.012,0.711) (2.042,0.732) (2.072,0.752) (2.102,0.772) (2.132,0.791) (2.162,0.811) (2.193,0.830) (2.223,0.849) (2.254,0.867) (2.285,0.886) (2.316,0.905) (2.347,0.923) (2.378,0.941) (2.409,0.959) (2.440,0.977) (2.472,0.995) (2.503,1.012) (2.534,1.030) (2.566,1.047) (2.598,1.064) (2.600,1.070)}
\LC{(1.206,0.258) (1.221,0.291) (1.237,0.323) (1.253,0.355) (1.270,0.387) (1.287,0.419) (1.305,0.450) (1.323,0.481) (1.342,0.512) (1.361,0.542) (1.381,0.572) (1.402,0.602) (1.422,0.631) (1.444,0.660) (1.465,0.689) (1.487,0.717) (1.510,0.746) (1.532,0.774) (1.555,0.801) (1.578,0.829) (1.602,0.856) (1.626,0.883) (1.650,0.910) (1.674,0.937) (1.698,0.963) (1.723,0.989) (1.748,1.015) (1.773,1.041) (1.798,1.067) (1.824,1.092) (1.849,1.118) (1.875,1.143) (1.901,1.168) (1.927,1.193) (1.953,1.218) (1.979,1.242) (2.005,1.267) (2.032,1.291) (2.058,1.316) (2.085,1.340) (2.112,1.364) (2.138,1.388) (2.165,1.412) (2.192,1.436) (2.219,1.460) (2.247,1.483) (2.274,1.507) (2.301,1.530) (2.328,1.554) (2.356,1.577) (2.383,1.600) (2.411,1.623) (2.439,1.646) (2.466,1.669) (2.494,1.692) (2.522,1.715) (2.550,1.738) (2.578,1.761) (2.600,1.783)}
\LC{(0.902,0.384) (0.916,0.417) (0.930,0.450) (0.945,0.483) (0.960,0.516) (0.975,0.548) (0.991,0.581) (1.007,0.613) (1.023,0.645) (1.039,0.677) (1.056,0.709) (1.073,0.741) (1.090,0.773) (1.108,0.804) (1.126,0.835) (1.144,0.867) (1.162,0.898) (1.180,0.929) (1.199,0.959) (1.217,0.990) (1.236,1.021) (1.255,1.051) (1.275,1.082) (1.294,1.112) (1.314,1.142) (1.334,1.172) (1.353,1.202) (1.373,1.232) (1.394,1.262) (1.414,1.292) (1.434,1.321) (1.455,1.351) (1.476,1.380) (1.496,1.410) (1.517,1.439) (1.538,1.468) (1.559,1.498) (1.580,1.527) (1.602,1.556) (1.623,1.585) (1.644,1.614) (1.666,1.643) (1.687,1.672) (1.709,1.700) (1.731,1.729) (1.752,1.758) (1.774,1.786) (1.796,1.815) (1.818,1.844) (1.840,1.872) (1.862,1.901) (1.884,1.929) (1.906,1.957) (1.928,1.986) (1.943,2.000)}
\LC{(0.565,0.524) (0.579,0.557) (0.593,0.590) (0.607,0.623) (0.621,0.656) (0.635,0.689) (0.650,0.722) (0.664,0.756) (0.678,0.789) (0.693,0.822) (0.707,0.854) (0.722,0.887) (0.736,0.920) (0.751,0.953) (0.766,0.986) (0.781,1.019) (0.796,1.052) (0.811,1.084) (0.826,1.117) (0.841,1.150) (0.856,1.182) (0.871,1.215) (0.886,1.248) (0.901,1.280) (0.917,1.313) (0.932,1.345) (0.948,1.378) (0.963,1.410) (0.979,1.443) (0.995,1.475) (1.010,1.508) (1.026,1.540) (1.042,1.572) (1.058,1.605) (1.073,1.637) (1.089,1.669) (1.105,1.702) (1.121,1.734) (1.137,1.766) (1.153,1.798) (1.169,1.830) (1.186,1.863) (1.202,1.895) (1.218,1.927) (1.234,1.959) (1.250,1.991) (1.256,2.000)}
\LC{(0.223,0.666) (0.236,0.699) (0.249,0.733) (0.262,0.766) (0.275,0.800) (0.288,0.833) (0.301,0.867) (0.314,0.901) (0.326,0.934) (0.338,0.968) (0.351,1.002) (0.363,1.036) (0.375,1.070) (0.387,1.104) (0.399,1.138) (0.410,1.172) (0.422,1.206) (0.433,1.240) (0.445,1.274) (0.456,1.308) (0.468,1.343) (0.479,1.377) (0.490,1.411) (0.501,1.445) (0.513,1.479) (0.524,1.514) (0.535,1.548) (0.546,1.582) (0.557,1.616) (0.568,1.651) (0.579,1.685) (0.590,1.719) (0.601,1.754) (0.611,1.788) (0.622,1.822) (0.633,1.857) (0.644,1.891) (0.655,1.925) (0.665,1.960) (0.676,1.994) (0.680,2.000)}
\LC{(-0.107,0.803) (-0.095,0.836) (-0.082,0.870) (-0.071,0.904) (-0.060,0.938) (-0.049,0.973) (-0.039,1.007) (-0.029,1.042) (-0.019,1.077) (-0.010,1.111) (-0.002,1.146) (0.007,1.181) (0.015,1.216) (0.023,1.252) (0.031,1.287) (0.038,1.322) (0.045,1.357) (0.052,1.393) (0.059,1.428) (0.066,1.463) (0.073,1.499) (0.079,1.534) (0.085,1.569) (0.091,1.605) (0.098,1.640) (0.104,1.676) (0.109,1.711) (0.115,1.747) (0.121,1.783) (0.126,1.818) (0.132,1.854) (0.137,1.889) (0.143,1.925) (0.148,1.960) (0.153,1.996) (0.155,2.000)}

      % Conductor con punta - flechas
% Generado por scripts/generar_datos.py (no editar a mano).
\FL{-0.332}{1.586}{89.3}
\FL{-0.742}{1.606}{100.9}
\FL{-1.161}{1.589}{113.5}
\FL{-1.635}{1.552}{126.9}
\FL{-2.093}{1.363}{140.4}
\FL{-2.144}{0.864}{152.5}
\FL{-2.187}{0.452}{165.0}
\FL{-2.200}{0.071}{177.6}
\FL{-2.197}{-0.307}{-169.8}
\FL{-2.160}{-0.700}{-157.2}
\FL{-2.111}{-1.155}{-144.9}
\FL{-1.855}{-1.538}{-132.3}
\FL{-1.337}{-1.582}{-118.7}
\FL{-0.898}{-1.597}{-105.6}
\FL{-0.488}{-1.595}{-93.5}
\FL{-0.070}{-1.563}{-83.0}
\FL{0.352}{-1.518}{-74.8}
\FL{0.798}{-1.466}{-67.4}
\FL{1.292}{-1.425}{-58.7}
\FL{1.890}{-1.410}{-46.8}
\FL{2.079}{-0.945}{-36.8}
\FL{2.189}{-0.529}{-24.8}
\FL{2.269}{-0.179}{-9.2}
\FL{2.216}{0.401}{19.8}
\FL{2.122}{0.785}{32.9}
\FL{2.005}{1.267}{42.9}
\FL{1.503}{1.420}{54.6}
\FL{0.984}{1.454}{64.3}
\FL{0.516}{1.491}{72.0}
\FL{0.081}{1.546}{79.9}

      % Contorno de la gota
% Generado por scripts/generar_datos.py (no editar a mano).
\def\formagota{(-0.259,0.833) -- (-0.262,0.834) -- (-0.265,0.835) -- (-0.268,0.836) -- (-0.270,0.837) -- (-0.273,0.839) -- (-0.276,0.840) -- (-0.279,0.841) -- (-0.281,0.842) -- (-0.284,0.843) -- (-0.287,0.844) -- (-0.290,0.845) -- (-0.293,0.846) -- (-0.295,0.847) -- (-0.298,0.848) -- (-0.301,0.849) -- (-0.304,0.850) -- (-0.306,0.851) -- (-0.309,0.852) -- (-0.312,0.853) -- (-0.315,0.854) -- (-0.318,0.855) -- (-0.320,0.855) -- (-0.323,0.856) -- (-0.326,0.857) -- (-0.329,0.858) -- (-0.332,0.859) -- (-0.335,0.860) -- (-0.337,0.861) -- (-0.340,0.862) -- (-0.343,0.863) -- (-0.346,0.863) -- (-0.349,0.864) -- (-0.352,0.865) -- (-0.354,0.866) -- (-0.357,0.867) -- (-0.360,0.867) -- (-0.363,0.868) -- (-0.366,0.869) -- (-0.369,0.870) -- (-0.371,0.870) -- (-0.374,0.871) -- (-0.377,0.872) -- (-0.380,0.873) -- (-0.383,0.873) -- (-0.386,0.874) -- (-0.389,0.875) -- (-0.392,0.876) -- (-0.394,0.876) -- (-0.397,0.877) -- (-0.400,0.878) -- (-0.403,0.878) -- (-0.406,0.879) -- (-0.409,0.879) -- (-0.412,0.880) -- (-0.415,0.881) -- (-0.418,0.881) -- (-0.420,0.882) -- (-0.423,0.882) -- (-0.426,0.883) -- (-0.429,0.884) -- (-0.432,0.884) -- (-0.435,0.885) -- (-0.438,0.885) -- (-0.441,0.886) -- (-0.444,0.886) -- (-0.447,0.887) -- (-0.449,0.887) -- (-0.452,0.888) -- (-0.455,0.888) -- (-0.458,0.889) -- (-0.461,0.889) -- (-0.464,0.890) -- (-0.467,0.890) -- (-0.470,0.891) -- (-0.473,0.891) -- (-0.476,0.891) -- (-0.479,0.892) -- (-0.482,0.892) -- (-0.485,0.893) -- (-0.487,0.893) -- (-0.490,0.893) -- (-0.493,0.894) -- (-0.496,0.894) -- (-0.499,0.894) -- (-0.502,0.895) -- (-0.505,0.895) -- (-0.508,0.895) -- (-0.511,0.896) -- (-0.514,0.896) -- (-0.517,0.896) -- (-0.520,0.896) -- (-0.523,0.897) -- (-0.526,0.897) -- (-0.529,0.897) -- (-0.532,0.897) -- (-0.535,0.898) -- (-0.537,0.898) -- (-0.540,0.898) -- (-0.543,0.898) -- (-0.546,0.898) -- (-0.549,0.899) -- (-0.552,0.899) -- (-0.555,0.899) -- (-0.558,0.899) -- (-0.561,0.899) -- (-0.564,0.899) -- (-0.567,0.899) -- (-0.570,0.899) -- (-0.573,0.900) -- (-0.576,0.900) -- (-0.579,0.900) -- (-0.582,0.900) -- (-0.585,0.900) -- (-0.588,0.900) -- (-0.591,0.900) -- (-0.594,0.900) -- (-0.597,0.900) -- (-0.599,0.900) -- (-0.602,0.900) -- (-0.605,0.900) -- (-0.608,0.900) -- (-0.611,0.900) -- (-0.614,0.900) -- (-0.617,0.900) -- (-0.620,0.900) -- (-0.623,0.900) -- (-0.626,0.900) -- (-0.629,0.900) -- (-0.632,0.899) -- (-0.635,0.899) -- (-0.638,0.899) -- (-0.641,0.899) -- (-0.644,0.899) -- (-0.647,0.899) -- (-0.650,0.899) -- (-0.653,0.898) -- (-0.656,0.898) -- (-0.659,0.898) -- (-0.661,0.898) -- (-0.664,0.898) -- (-0.667,0.897) -- (-0.670,0.897) -- (-0.673,0.897) -- (-0.676,0.897) -- (-0.679,0.896) -- (-0.682,0.896) -- (-0.685,0.896) -- (-0.688,0.896) -- (-0.691,0.895) -- (-0.694,0.895) -- (-0.697,0.895) -- (-0.700,0.894) -- (-0.703,0.894) -- (-0.706,0.894) -- (-0.709,0.893) -- (-0.711,0.893) -- (-0.714,0.893) -- (-0.717,0.892) -- (-0.720,0.892) -- (-0.723,0.892) -- (-0.726,0.891) -- (-0.729,0.891) -- (-0.732,0.890) -- (-0.735,0.890) -- (-0.738,0.889) -- (-0.741,0.889) -- (-0.744,0.888) -- (-0.747,0.888) -- (-0.750,0.887) -- (-0.752,0.887) -- (-0.755,0.886) -- (-0.758,0.886) -- (-0.761,0.885) -- (-0.764,0.885) -- (-0.767,0.884) -- (-0.770,0.884) -- (-0.773,0.883) -- (-0.776,0.883) -- (-0.779,0.882) -- (-0.781,0.881) -- (-0.784,0.881) -- (-0.787,0.880) -- (-0.790,0.880) -- (-0.793,0.879) -- (-0.796,0.878) -- (-0.799,0.878) -- (-0.802,0.877) -- (-0.805,0.876) -- (-0.807,0.876) -- (-0.810,0.875) -- (-0.813,0.874) -- (-0.816,0.874) -- (-0.819,0.873) -- (-0.822,0.872) -- (-0.825,0.871) -- (-0.828,0.871) -- (-0.830,0.870) -- (-0.833,0.869) -- (-0.836,0.868) -- (-0.839,0.868) -- (-0.842,0.867) -- (-0.845,0.866) -- (-0.847,0.865) -- (-0.850,0.864) -- (-0.853,0.864) -- (-0.856,0.863) -- (-0.859,0.862) -- (-0.862,0.861) -- (-0.864,0.860) -- (-0.867,0.859) -- (-0.870,0.858) -- (-0.873,0.858) -- (-0.876,0.857) -- (-0.879,0.856) -- (-0.881,0.855) -- (-0.884,0.854) -- (-0.887,0.853) -- (-0.890,0.852) -- (-0.893,0.851) -- (-0.895,0.850) -- (-0.898,0.849) -- (-0.901,0.848) -- (-0.904,0.847) -- (-0.907,0.846) -- (-0.909,0.845) -- (-0.912,0.844) -- (-0.915,0.843) -- (-0.918,0.842) -- (-0.920,0.841) -- (-0.923,0.840) -- (-0.926,0.839) -- (-0.929,0.838) -- (-0.931,0.837) -- (-0.934,0.836) -- (-0.937,0.835) -- (-0.940,0.833) -- (-0.942,0.832) -- (-0.945,0.831) -- (-0.948,0.830) -- (-0.951,0.829) -- (-0.953,0.828) -- (-0.956,0.827) -- (-0.959,0.825) -- (-0.961,0.824) -- (-0.964,0.823) -- (-0.967,0.822) -- (-0.969,0.821) -- (-0.972,0.819) -- (-0.975,0.818) -- (-0.978,0.817) -- (-0.980,0.816) -- (-0.983,0.814) -- (-0.986,0.813) -- (-0.988,0.812) -- (-0.991,0.811) -- (-0.994,0.809) -- (-0.996,0.808) -- (-0.999,0.807) -- (-1.002,0.805) -- (-1.004,0.804) -- (-1.007,0.803) -- (-1.009,0.801) -- (-1.012,0.800) -- (-1.015,0.799) -- (-1.017,0.797) -- (-1.020,0.796) -- (-1.023,0.795) -- (-1.025,0.793) -- (-1.028,0.792) -- (-1.030,0.790) -- (-1.033,0.789) -- (-1.036,0.788) -- (-1.038,0.786) -- (-1.041,0.785) -- (-1.043,0.783) -- (-1.046,0.782) -- (-1.048,0.780) -- (-1.051,0.779) -- (-1.054,0.777) -- (-1.056,0.776) -- (-1.059,0.774) -- (-1.061,0.773) -- (-1.064,0.771) -- (-1.066,0.770) -- (-1.069,0.768) -- (-1.071,0.767) -- (-1.074,0.765) -- (-1.076,0.764) -- (-1.079,0.762) -- (-1.081,0.760) -- (-1.084,0.759) -- (-1.086,0.757) -- (-1.089,0.756) -- (-1.091,0.754) -- (-1.094,0.752) -- (-1.096,0.751) -- (-1.099,0.749) -- (-1.101,0.748) -- (-1.104,0.746) -- (-1.106,0.744) -- (-1.108,0.743) -- (-1.111,0.741) -- (-1.113,0.739) -- (-1.116,0.738) -- (-1.118,0.736) -- (-1.121,0.734) -- (-1.123,0.732) -- (-1.125,0.731) -- (-1.128,0.729) -- (-1.130,0.727) -- (-1.133,0.725) -- (-1.135,0.724) -- (-1.137,0.722) -- (-1.140,0.720) -- (-1.142,0.718) -- (-1.144,0.717) -- (-1.147,0.715) -- (-1.149,0.713) -- (-1.151,0.711) -- (-1.154,0.709) -- (-1.156,0.708) -- (-1.158,0.706) -- (-1.161,0.704) -- (-1.163,0.702) -- (-1.165,0.700) -- (-1.168,0.698) -- (-1.170,0.697) -- (-1.172,0.695) -- (-1.174,0.693) -- (-1.177,0.691) -- (-1.179,0.689) -- (-1.181,0.687) -- (-1.184,0.685) -- (-1.186,0.683) -- (-1.188,0.681) -- (-1.190,0.679) -- (-1.192,0.677) -- (-1.195,0.675) -- (-1.197,0.674) -- (-1.199,0.672) -- (-1.201,0.670) -- (-1.204,0.668) -- (-1.206,0.666) -- (-1.208,0.664) -- (-1.210,0.662) -- (-1.212,0.660) -- (-1.214,0.658) -- (-1.217,0.656) -- (-1.219,0.654) -- (-1.221,0.652) -- (-1.223,0.649) -- (-1.225,0.647) -- (-1.227,0.645) -- (-1.229,0.643) -- (-1.231,0.641) -- (-1.234,0.639) -- (-1.236,0.637) -- (-1.238,0.635) -- (-1.240,0.633) -- (-1.242,0.631) -- (-1.244,0.629) -- (-1.246,0.627) -- (-1.248,0.624) -- (-1.250,0.622) -- (-1.252,0.620) -- (-1.254,0.618) -- (-1.256,0.616) -- (-1.258,0.614) -- (-1.260,0.612) -- (-1.262,0.609) -- (-1.264,0.607) -- (-1.266,0.605) -- (-1.268,0.603) -- (-1.270,0.601) -- (-1.272,0.598) -- (-1.274,0.596) -- (-1.276,0.594) -- (-1.278,0.592) -- (-1.280,0.590) -- (-1.282,0.587) -- (-1.284,0.585) -- (-1.286,0.583) -- (-1.288,0.581) -- (-1.290,0.578) -- (-1.291,0.576) -- (-1.293,0.574) -- (-1.295,0.571) -- (-1.297,0.569) -- (-1.299,0.567) -- (-1.301,0.565) -- (-1.303,0.562) -- (-1.305,0.560) -- (-1.306,0.558) -- (-1.308,0.555) -- (-1.310,0.553) -- (-1.312,0.551) -- (-1.314,0.548) -- (-1.315,0.546) -- (-1.317,0.544) -- (-1.319,0.541) -- (-1.321,0.539) -- (-1.323,0.537) -- (-1.324,0.534) -- (-1.326,0.532) -- (-1.328,0.529) -- (-1.330,0.527) -- (-1.331,0.525) -- (-1.333,0.522) -- (-1.335,0.520) -- (-1.336,0.517) -- (-1.338,0.515) -- (-1.340,0.513) -- (-1.341,0.510) -- (-1.343,0.508) -- (-1.345,0.505) -- (-1.346,0.503) -- (-1.348,0.500) -- (-1.350,0.498) -- (-1.351,0.495) -- (-1.353,0.493) -- (-1.355,0.490) -- (-1.356,0.488) -- (-1.358,0.485) -- (-1.359,0.483) -- (-1.361,0.480) -- (-1.363,0.478) -- (-1.364,0.475) -- (-1.366,0.473) -- (-1.367,0.470) -- (-1.369,0.468) -- (-1.370,0.465) -- (-1.372,0.463) -- (-1.373,0.460) -- (-1.375,0.458) -- (-1.376,0.455) -- (-1.378,0.453) -- (-1.379,0.450) -- (-1.381,0.448) -- (-1.382,0.445) -- (-1.384,0.442) -- (-1.385,0.440) -- (-1.387,0.437) -- (-1.388,0.435) -- (-1.389,0.432) -- (-1.391,0.430) -- (-1.392,0.427) -- (-1.394,0.424) -- (-1.395,0.422) -- (-1.396,0.419) -- (-1.398,0.416) -- (-1.399,0.414) -- (-1.401,0.411) -- (-1.402,0.409) -- (-1.403,0.406) -- (-1.405,0.403) -- (-1.406,0.401) -- (-1.407,0.398) -- (-1.408,0.395) -- (-1.410,0.393) -- (-1.411,0.390) -- (-1.412,0.387) -- (-1.414,0.385) -- (-1.415,0.382) -- (-1.416,0.379) -- (-1.417,0.377) -- (-1.419,0.374) -- (-1.420,0.371) -- (-1.421,0.369) -- (-1.422,0.366) -- (-1.423,0.363) -- (-1.425,0.361) -- (-1.426,0.358) -- (-1.427,0.355) -- (-1.428,0.352) -- (-1.429,0.350) -- (-1.430,0.347) -- (-1.432,0.344) -- (-1.433,0.341) -- (-1.434,0.339) -- (-1.435,0.336) -- (-1.436,0.333) -- (-1.437,0.330) -- (-1.438,0.328) -- (-1.439,0.325) -- (-1.440,0.322) -- (-1.441,0.319) -- (-1.442,0.317) -- (-1.443,0.314) -- (-1.444,0.311) -- (-1.445,0.308) -- (-1.447,0.306) -- (-1.448,0.303) -- (-1.448,0.300) -- (-1.449,0.297) -- (-1.450,0.294) -- (-1.451,0.292) -- (-1.452,0.289) -- (-1.453,0.286) -- (-1.454,0.283) -- (-1.455,0.280) -- (-1.456,0.278) -- (-1.457,0.275) -- (-1.458,0.272) -- (-1.459,0.269) -- (-1.460,0.266) -- (-1.461,0.264) -- (-1.461,0.261) -- (-1.462,0.258) -- (-1.463,0.255) -- (-1.464,0.252) -- (-1.465,0.249) -- (-1.466,0.247) -- (-1.466,0.244) -- (-1.467,0.241) -- (-1.468,0.238) -- (-1.469,0.235) -- (-1.469,0.232) -- (-1.470,0.229) -- (-1.471,0.227) -- (-1.472,0.224) -- (-1.472,0.221) -- (-1.473,0.218) -- (-1.474,0.215) -- (-1.475,0.212) -- (-1.475,0.209) -- (-1.476,0.207) -- (-1.477,0.204) -- (-1.477,0.201) -- (-1.478,0.198) -- (-1.479,0.195) -- (-1.479,0.192) -- (-1.480,0.189) -- (-1.480,0.186) -- (-1.481,0.183) -- (-1.482,0.181) -- (-1.482,0.178) -- (-1.483,0.175) -- (-1.483,0.172) -- (-1.484,0.169) -- (-1.485,0.166) -- (-1.485,0.163) -- (-1.486,0.160) -- (-1.486,0.157) -- (-1.487,0.154) -- (-1.487,0.151) -- (-1.488,0.149) -- (-1.488,0.146) -- (-1.489,0.143) -- (-1.489,0.140) -- (-1.490,0.137) -- (-1.490,0.134) -- (-1.490,0.131) -- (-1.491,0.128) -- (-1.491,0.125) -- (-1.492,0.122) -- (-1.492,0.119) -- (-1.492,0.116) -- (-1.493,0.113) -- (-1.493,0.111) -- (-1.494,0.108) -- (-1.494,0.105) -- (-1.494,0.102) -- (-1.495,0.099) -- (-1.495,0.096) -- (-1.495,0.093) -- (-1.495,0.090) -- (-1.496,0.087) -- (-1.496,0.084) -- (-1.496,0.081) -- (-1.497,0.078) -- (-1.497,0.075) -- (-1.497,0.072) -- (-1.497,0.069) -- (-1.498,0.066) -- (-1.498,0.063) -- (-1.498,0.061) -- (-1.498,0.058) -- (-1.498,0.055) -- (-1.498,0.052) -- (-1.499,0.049) -- (-1.499,0.046) -- (-1.499,0.043) -- (-1.499,0.040) -- (-1.499,0.037) -- (-1.499,0.034) -- (-1.499,0.031) -- (-1.500,0.028) -- (-1.500,0.025) -- (-1.500,0.022) -- (-1.500,0.019) -- (-1.500,0.016) -- (-1.500,0.013) -- (-1.500,0.010) -- (-1.500,0.007) -- (-1.500,0.004) -- (-1.500,0.001) -- (-1.500,-0.001) -- (-1.500,-0.004) -- (-1.500,-0.007) -- (-1.500,-0.010) -- (-1.500,-0.013) -- (-1.500,-0.016) -- (-1.500,-0.019) -- (-1.500,-0.022) -- (-1.500,-0.025) -- (-1.500,-0.028) -- (-1.499,-0.031) -- (-1.499,-0.034) -- (-1.499,-0.037) -- (-1.499,-0.040) -- (-1.499,-0.043) -- (-1.499,-0.046) -- (-1.499,-0.049) -- (-1.498,-0.052) -- (-1.498,-0.055) -- (-1.498,-0.058) -- (-1.498,-0.061) -- (-1.498,-0.063) -- (-1.498,-0.066) -- (-1.497,-0.069) -- (-1.497,-0.072) -- (-1.497,-0.075) -- (-1.497,-0.078) -- (-1.496,-0.081) -- (-1.496,-0.084) -- (-1.496,-0.087) -- (-1.495,-0.090) -- (-1.495,-0.093) -- (-1.495,-0.096) -- (-1.495,-0.099) -- (-1.494,-0.102) -- (-1.494,-0.105) -- (-1.494,-0.108) -- (-1.493,-0.111) -- (-1.493,-0.113) -- (-1.492,-0.116) -- (-1.492,-0.119) -- (-1.492,-0.122) -- (-1.491,-0.125) -- (-1.491,-0.128) -- (-1.490,-0.131) -- (-1.490,-0.134) -- (-1.490,-0.137) -- (-1.489,-0.140) -- (-1.489,-0.143) -- (-1.488,-0.146) -- (-1.488,-0.149) -- (-1.487,-0.151) -- (-1.487,-0.154) -- (-1.486,-0.157) -- (-1.486,-0.160) -- (-1.485,-0.163) -- (-1.485,-0.166) -- (-1.484,-0.169) -- (-1.483,-0.172) -- (-1.483,-0.175) -- (-1.482,-0.178) -- (-1.482,-0.181) -- (-1.481,-0.183) -- (-1.480,-0.186) -- (-1.480,-0.189) -- (-1.479,-0.192) -- (-1.479,-0.195) -- (-1.478,-0.198) -- (-1.477,-0.201) -- (-1.477,-0.204) -- (-1.476,-0.207) -- (-1.475,-0.209) -- (-1.475,-0.212) -- (-1.474,-0.215) -- (-1.473,-0.218) -- (-1.472,-0.221) -- (-1.472,-0.224) -- (-1.471,-0.227) -- (-1.470,-0.229) -- (-1.469,-0.232) -- (-1.469,-0.235) -- (-1.468,-0.238) -- (-1.467,-0.241) -- (-1.466,-0.244) -- (-1.466,-0.247) -- (-1.465,-0.249) -- (-1.464,-0.252) -- (-1.463,-0.255) -- (-1.462,-0.258) -- (-1.461,-0.261) -- (-1.461,-0.264) -- (-1.460,-0.266) -- (-1.459,-0.269) -- (-1.458,-0.272) -- (-1.457,-0.275) -- (-1.456,-0.278) -- (-1.455,-0.280) -- (-1.454,-0.283) -- (-1.453,-0.286) -- (-1.452,-0.289) -- (-1.451,-0.292) -- (-1.450,-0.294) -- (-1.449,-0.297) -- (-1.448,-0.300) -- (-1.448,-0.303) -- (-1.447,-0.306) -- (-1.445,-0.308) -- (-1.444,-0.311) -- (-1.443,-0.314) -- (-1.442,-0.317) -- (-1.441,-0.319) -- (-1.440,-0.322) -- (-1.439,-0.325) -- (-1.438,-0.328) -- (-1.437,-0.330) -- (-1.436,-0.333) -- (-1.435,-0.336) -- (-1.434,-0.339) -- (-1.433,-0.341) -- (-1.432,-0.344) -- (-1.430,-0.347) -- (-1.429,-0.350) -- (-1.428,-0.352) -- (-1.427,-0.355) -- (-1.426,-0.358) -- (-1.425,-0.361) -- (-1.423,-0.363) -- (-1.422,-0.366) -- (-1.421,-0.369) -- (-1.420,-0.371) -- (-1.419,-0.374) -- (-1.417,-0.377) -- (-1.416,-0.379) -- (-1.415,-0.382) -- (-1.414,-0.385) -- (-1.412,-0.387) -- (-1.411,-0.390) -- (-1.410,-0.393) -- (-1.408,-0.395) -- (-1.407,-0.398) -- (-1.406,-0.401) -- (-1.405,-0.403) -- (-1.403,-0.406) -- (-1.402,-0.409) -- (-1.401,-0.411) -- (-1.399,-0.414) -- (-1.398,-0.416) -- (-1.396,-0.419) -- (-1.395,-0.422) -- (-1.394,-0.424) -- (-1.392,-0.427) -- (-1.391,-0.430) -- (-1.389,-0.432) -- (-1.388,-0.435) -- (-1.387,-0.437) -- (-1.385,-0.440) -- (-1.384,-0.442) -- (-1.382,-0.445) -- (-1.381,-0.448) -- (-1.379,-0.450) -- (-1.378,-0.453) -- (-1.376,-0.455) -- (-1.375,-0.458) -- (-1.373,-0.460) -- (-1.372,-0.463) -- (-1.370,-0.465) -- (-1.369,-0.468) -- (-1.367,-0.470) -- (-1.366,-0.473) -- (-1.364,-0.475) -- (-1.363,-0.478) -- (-1.361,-0.480) -- (-1.359,-0.483) -- (-1.358,-0.485) -- (-1.356,-0.488) -- (-1.355,-0.490) -- (-1.353,-0.493) -- (-1.351,-0.495) -- (-1.350,-0.498) -- (-1.348,-0.500) -- (-1.346,-0.503) -- (-1.345,-0.505) -- (-1.343,-0.508) -- (-1.341,-0.510) -- (-1.340,-0.513) -- (-1.338,-0.515) -- (-1.336,-0.517) -- (-1.335,-0.520) -- (-1.333,-0.522) -- (-1.331,-0.525) -- (-1.330,-0.527) -- (-1.328,-0.529) -- (-1.326,-0.532) -- (-1.324,-0.534) -- (-1.323,-0.537) -- (-1.321,-0.539) -- (-1.319,-0.541) -- (-1.317,-0.544) -- (-1.315,-0.546) -- (-1.314,-0.548) -- (-1.312,-0.551) -- (-1.310,-0.553) -- (-1.308,-0.555) -- (-1.306,-0.558) -- (-1.305,-0.560) -- (-1.303,-0.562) -- (-1.301,-0.565) -- (-1.299,-0.567) -- (-1.297,-0.569) -- (-1.295,-0.571) -- (-1.293,-0.574) -- (-1.291,-0.576) -- (-1.290,-0.578) -- (-1.288,-0.581) -- (-1.286,-0.583) -- (-1.284,-0.585) -- (-1.282,-0.587) -- (-1.280,-0.590) -- (-1.278,-0.592) -- (-1.276,-0.594) -- (-1.274,-0.596) -- (-1.272,-0.598) -- (-1.270,-0.601) -- (-1.268,-0.603) -- (-1.266,-0.605) -- (-1.264,-0.607) -- (-1.262,-0.609) -- (-1.260,-0.612) -- (-1.258,-0.614) -- (-1.256,-0.616) -- (-1.254,-0.618) -- (-1.252,-0.620) -- (-1.250,-0.622) -- (-1.248,-0.624) -- (-1.246,-0.627) -- (-1.244,-0.629) -- (-1.242,-0.631) -- (-1.240,-0.633) -- (-1.238,-0.635) -- (-1.236,-0.637) -- (-1.234,-0.639) -- (-1.231,-0.641) -- (-1.229,-0.643) -- (-1.227,-0.645) -- (-1.225,-0.647) -- (-1.223,-0.649) -- (-1.221,-0.652) -- (-1.219,-0.654) -- (-1.217,-0.656) -- (-1.214,-0.658) -- (-1.212,-0.660) -- (-1.210,-0.662) -- (-1.208,-0.664) -- (-1.206,-0.666) -- (-1.204,-0.668) -- (-1.201,-0.670) -- (-1.199,-0.672) -- (-1.197,-0.674) -- (-1.195,-0.675) -- (-1.192,-0.677) -- (-1.190,-0.679) -- (-1.188,-0.681) -- (-1.186,-0.683) -- (-1.184,-0.685) -- (-1.181,-0.687) -- (-1.179,-0.689) -- (-1.177,-0.691) -- (-1.174,-0.693) -- (-1.172,-0.695) -- (-1.170,-0.697) -- (-1.168,-0.698) -- (-1.165,-0.700) -- (-1.163,-0.702) -- (-1.161,-0.704) -- (-1.158,-0.706) -- (-1.156,-0.708) -- (-1.154,-0.709) -- (-1.151,-0.711) -- (-1.149,-0.713) -- (-1.147,-0.715) -- (-1.144,-0.717) -- (-1.142,-0.718) -- (-1.140,-0.720) -- (-1.137,-0.722) -- (-1.135,-0.724) -- (-1.133,-0.725) -- (-1.130,-0.727) -- (-1.128,-0.729) -- (-1.125,-0.731) -- (-1.123,-0.732) -- (-1.121,-0.734) -- (-1.118,-0.736) -- (-1.116,-0.738) -- (-1.113,-0.739) -- (-1.111,-0.741) -- (-1.108,-0.743) -- (-1.106,-0.744) -- (-1.104,-0.746) -- (-1.101,-0.748) -- (-1.099,-0.749) -- (-1.096,-0.751) -- (-1.094,-0.752) -- (-1.091,-0.754) -- (-1.089,-0.756) -- (-1.086,-0.757) -- (-1.084,-0.759) -- (-1.081,-0.760) -- (-1.079,-0.762) -- (-1.076,-0.764) -- (-1.074,-0.765) -- (-1.071,-0.767) -- (-1.069,-0.768) -- (-1.066,-0.770) -- (-1.064,-0.771) -- (-1.061,-0.773) -- (-1.059,-0.774) -- (-1.056,-0.776) -- (-1.054,-0.777) -- (-1.051,-0.779) -- (-1.048,-0.780) -- (-1.046,-0.782) -- (-1.043,-0.783) -- (-1.041,-0.785) -- (-1.038,-0.786) -- (-1.036,-0.788) -- (-1.033,-0.789) -- (-1.030,-0.790) -- (-1.028,-0.792) -- (-1.025,-0.793) -- (-1.023,-0.795) -- (-1.020,-0.796) -- (-1.017,-0.797) -- (-1.015,-0.799) -- (-1.012,-0.800) -- (-1.009,-0.801) -- (-1.007,-0.803) -- (-1.004,-0.804) -- (-1.002,-0.805) -- (-0.999,-0.807) -- (-0.996,-0.808) -- (-0.994,-0.809) -- (-0.991,-0.811) -- (-0.988,-0.812) -- (-0.986,-0.813) -- (-0.983,-0.814) -- (-0.980,-0.816) -- (-0.978,-0.817) -- (-0.975,-0.818) -- (-0.972,-0.819) -- (-0.969,-0.821) -- (-0.967,-0.822) -- (-0.964,-0.823) -- (-0.961,-0.824) -- (-0.959,-0.825) -- (-0.956,-0.827) -- (-0.953,-0.828) -- (-0.951,-0.829) -- (-0.948,-0.830) -- (-0.945,-0.831) -- (-0.942,-0.832) -- (-0.940,-0.833) -- (-0.937,-0.835) -- (-0.934,-0.836) -- (-0.931,-0.837) -- (-0.929,-0.838) -- (-0.926,-0.839) -- (-0.923,-0.840) -- (-0.920,-0.841) -- (-0.918,-0.842) -- (-0.915,-0.843) -- (-0.912,-0.844) -- (-0.909,-0.845) -- (-0.907,-0.846) -- (-0.904,-0.847) -- (-0.901,-0.848) -- (-0.898,-0.849) -- (-0.895,-0.850) -- (-0.893,-0.851) -- (-0.890,-0.852) -- (-0.887,-0.853) -- (-0.884,-0.854) -- (-0.881,-0.855) -- (-0.879,-0.856) -- (-0.876,-0.857) -- (-0.873,-0.858) -- (-0.870,-0.858) -- (-0.867,-0.859) -- (-0.864,-0.860) -- (-0.862,-0.861) -- (-0.859,-0.862) -- (-0.856,-0.863) -- (-0.853,-0.864) -- (-0.850,-0.864) -- (-0.847,-0.865) -- (-0.845,-0.866) -- (-0.842,-0.867) -- (-0.839,-0.868) -- (-0.836,-0.868) -- (-0.833,-0.869) -- (-0.830,-0.870) -- (-0.828,-0.871) -- (-0.825,-0.871) -- (-0.822,-0.872) -- (-0.819,-0.873) -- (-0.816,-0.874) -- (-0.813,-0.874) -- (-0.810,-0.875) -- (-0.807,-0.876) -- (-0.805,-0.876) -- (-0.802,-0.877) -- (-0.799,-0.878) -- (-0.796,-0.878) -- (-0.793,-0.879) -- (-0.790,-0.880) -- (-0.787,-0.880) -- (-0.784,-0.881) -- (-0.781,-0.881) -- (-0.779,-0.882) -- (-0.776,-0.883) -- (-0.773,-0.883) -- (-0.770,-0.884) -- (-0.767,-0.884) -- (-0.764,-0.885) -- (-0.761,-0.885) -- (-0.758,-0.886) -- (-0.755,-0.886) -- (-0.752,-0.887) -- (-0.750,-0.887) -- (-0.747,-0.888) -- (-0.744,-0.888) -- (-0.741,-0.889) -- (-0.738,-0.889) -- (-0.735,-0.890) -- (-0.732,-0.890) -- (-0.729,-0.891) -- (-0.726,-0.891) -- (-0.723,-0.892) -- (-0.720,-0.892) -- (-0.717,-0.892) -- (-0.714,-0.893) -- (-0.711,-0.893) -- (-0.709,-0.893) -- (-0.706,-0.894) -- (-0.703,-0.894) -- (-0.700,-0.894) -- (-0.697,-0.895) -- (-0.694,-0.895) -- (-0.691,-0.895) -- (-0.688,-0.896) -- (-0.685,-0.896) -- (-0.682,-0.896) -- (-0.679,-0.896) -- (-0.676,-0.897) -- (-0.673,-0.897) -- (-0.670,-0.897) -- (-0.667,-0.897) -- (-0.664,-0.898) -- (-0.661,-0.898) -- (-0.659,-0.898) -- (-0.656,-0.898) -- (-0.653,-0.898) -- (-0.650,-0.899) -- (-0.647,-0.899) -- (-0.644,-0.899) -- (-0.641,-0.899) -- (-0.638,-0.899) -- (-0.635,-0.899) -- (-0.632,-0.899) -- (-0.629,-0.900) -- (-0.626,-0.900) -- (-0.623,-0.900) -- (-0.620,-0.900) -- (-0.617,-0.900) -- (-0.614,-0.900) -- (-0.611,-0.900) -- (-0.608,-0.900) -- (-0.605,-0.900) -- (-0.602,-0.900) -- (-0.599,-0.900) -- (-0.597,-0.900) -- (-0.594,-0.900) -- (-0.591,-0.900) -- (-0.588,-0.900) -- (-0.585,-0.900) -- (-0.582,-0.900) -- (-0.579,-0.900) -- (-0.576,-0.900) -- (-0.573,-0.900) -- (-0.570,-0.899) -- (-0.567,-0.899) -- (-0.564,-0.899) -- (-0.561,-0.899) -- (-0.558,-0.899) -- (-0.555,-0.899) -- (-0.552,-0.899) -- (-0.549,-0.899) -- (-0.546,-0.898) -- (-0.543,-0.898) -- (-0.540,-0.898) -- (-0.537,-0.898) -- (-0.535,-0.898) -- (-0.532,-0.897) -- (-0.529,-0.897) -- (-0.526,-0.897) -- (-0.523,-0.897) -- (-0.520,-0.896) -- (-0.517,-0.896) -- (-0.514,-0.896) -- (-0.511,-0.896) -- (-0.508,-0.895) -- (-0.505,-0.895) -- (-0.502,-0.895) -- (-0.499,-0.894) -- (-0.496,-0.894) -- (-0.493,-0.894) -- (-0.490,-0.893) -- (-0.487,-0.893) -- (-0.485,-0.893) -- (-0.482,-0.892) -- (-0.479,-0.892) -- (-0.476,-0.891) -- (-0.473,-0.891) -- (-0.470,-0.891) -- (-0.467,-0.890) -- (-0.464,-0.890) -- (-0.461,-0.889) -- (-0.458,-0.889) -- (-0.455,-0.888) -- (-0.452,-0.888) -- (-0.449,-0.887) -- (-0.447,-0.887) -- (-0.444,-0.886) -- (-0.441,-0.886) -- (-0.438,-0.885) -- (-0.435,-0.885) -- (-0.432,-0.884) -- (-0.429,-0.884) -- (-0.426,-0.883) -- (-0.423,-0.882) -- (-0.420,-0.882) -- (-0.418,-0.881) -- (-0.415,-0.881) -- (-0.412,-0.880) -- (-0.409,-0.879) -- (-0.406,-0.879) -- (-0.403,-0.878) -- (-0.400,-0.878) -- (-0.397,-0.877) -- (-0.394,-0.876) -- (-0.392,-0.876) -- (-0.389,-0.875) -- (-0.386,-0.874) -- (-0.383,-0.873) -- (-0.380,-0.873) -- (-0.377,-0.872) -- (-0.374,-0.871) -- (-0.371,-0.870) -- (-0.369,-0.870) -- (-0.366,-0.869) -- (-0.363,-0.868) -- (-0.360,-0.867) -- (-0.357,-0.867) -- (-0.354,-0.866) -- (-0.352,-0.865) -- (-0.349,-0.864) -- (-0.346,-0.863) -- (-0.343,-0.863) -- (-0.340,-0.862) -- (-0.337,-0.861) -- (-0.335,-0.860) -- (-0.332,-0.859) -- (-0.329,-0.858) -- (-0.326,-0.857) -- (-0.323,-0.856) -- (-0.320,-0.855) -- (-0.318,-0.855) -- (-0.315,-0.854) -- (-0.312,-0.853) -- (-0.309,-0.852) -- (-0.306,-0.851) -- (-0.304,-0.850) -- (-0.301,-0.849) -- (-0.298,-0.848) -- (-0.295,-0.847) -- (-0.293,-0.846) -- (-0.290,-0.845) -- (-0.287,-0.844) -- (-0.284,-0.843) -- (-0.281,-0.842) -- (-0.279,-0.841) -- (-0.276,-0.840) -- (-0.273,-0.839) -- (-0.270,-0.837) -- (-0.268,-0.836) -- (-0.265,-0.835) -- (-0.262,-0.834) -- (-0.259,-0.833) -- (-0.256,-0.832) -- (-0.251,-0.830) -- (-0.246,-0.827) -- (-0.239,-0.825) -- (-0.233,-0.822) -- (-0.227,-0.820) -- (-0.220,-0.817) -- (-0.214,-0.814) -- (-0.208,-0.812) -- (-0.201,-0.809) -- (-0.195,-0.806) -- (-0.189,-0.804) -- (-0.182,-0.801) -- (-0.176,-0.798) -- (-0.170,-0.796) -- (-0.163,-0.793) -- (-0.157,-0.791) -- (-0.151,-0.788) -- (-0.144,-0.785) -- (-0.138,-0.783) -- (-0.132,-0.780) -- (-0.125,-0.777) -- (-0.119,-0.775) -- (-0.113,-0.772) -- (-0.106,-0.770) -- (-0.100,-0.767) -- (-0.093,-0.764) -- (-0.087,-0.762) -- (-0.081,-0.759) -- (-0.074,-0.756) -- (-0.068,-0.754) -- (-0.062,-0.751) -- (-0.055,-0.749) -- (-0.049,-0.746) -- (-0.043,-0.743) -- (-0.036,-0.741) -- (-0.030,-0.738) -- (-0.024,-0.735) -- (-0.017,-0.733) -- (-0.011,-0.730) -- (-0.005,-0.727) -- (0.002,-0.725) -- (0.008,-0.722) -- (0.014,-0.720) -- (0.021,-0.717) -- (0.027,-0.714) -- (0.033,-0.712) -- (0.040,-0.709) -- (0.046,-0.706) -- (0.052,-0.704) -- (0.059,-0.701) -- (0.065,-0.699) -- (0.071,-0.696) -- (0.078,-0.693) -- (0.084,-0.691) -- (0.091,-0.688) -- (0.097,-0.685) -- (0.103,-0.683) -- (0.110,-0.680) -- (0.116,-0.677) -- (0.122,-0.675) -- (0.129,-0.672) -- (0.135,-0.670) -- (0.141,-0.667) -- (0.148,-0.664) -- (0.154,-0.662) -- (0.160,-0.659) -- (0.167,-0.656) -- (0.173,-0.654) -- (0.179,-0.651) -- (0.186,-0.649) -- (0.192,-0.646) -- (0.198,-0.643) -- (0.205,-0.641) -- (0.211,-0.638) -- (0.217,-0.635) -- (0.224,-0.633) -- (0.230,-0.630) -- (0.236,-0.627) -- (0.243,-0.625) -- (0.249,-0.622) -- (0.256,-0.620) -- (0.262,-0.617) -- (0.268,-0.614) -- (0.275,-0.612) -- (0.281,-0.609) -- (0.287,-0.606) -- (0.294,-0.604) -- (0.300,-0.601) -- (0.306,-0.599) -- (0.313,-0.596) -- (0.319,-0.593) -- (0.325,-0.591) -- (0.332,-0.588) -- (0.338,-0.585) -- (0.344,-0.583) -- (0.351,-0.580) -- (0.357,-0.577) -- (0.363,-0.575) -- (0.370,-0.572) -- (0.376,-0.570) -- (0.382,-0.567) -- (0.389,-0.564) -- (0.395,-0.562) -- (0.401,-0.559) -- (0.408,-0.556) -- (0.414,-0.554) -- (0.421,-0.551) -- (0.427,-0.549) -- (0.433,-0.546) -- (0.440,-0.543) -- (0.446,-0.541) -- (0.452,-0.538) -- (0.459,-0.535) -- (0.465,-0.533) -- (0.471,-0.530) -- (0.478,-0.528) -- (0.484,-0.525) -- (0.490,-0.522) -- (0.497,-0.520) -- (0.503,-0.517) -- (0.509,-0.514) -- (0.516,-0.512) -- (0.522,-0.509) -- (0.528,-0.506) -- (0.535,-0.504) -- (0.541,-0.501) -- (0.547,-0.499) -- (0.554,-0.496) -- (0.560,-0.493) -- (0.566,-0.491) -- (0.573,-0.488) -- (0.579,-0.485) -- (0.586,-0.483) -- (0.592,-0.480) -- (0.598,-0.478) -- (0.605,-0.475) -- (0.611,-0.472) -- (0.617,-0.470) -- (0.624,-0.467) -- (0.630,-0.464) -- (0.636,-0.462) -- (0.643,-0.459) -- (0.649,-0.456) -- (0.655,-0.454) -- (0.662,-0.451) -- (0.668,-0.449) -- (0.674,-0.446) -- (0.681,-0.443) -- (0.687,-0.441) -- (0.693,-0.438) -- (0.700,-0.435) -- (0.706,-0.433) -- (0.712,-0.430) -- (0.719,-0.428) -- (0.725,-0.425) -- (0.731,-0.422) -- (0.738,-0.420) -- (0.744,-0.417) -- (0.751,-0.414) -- (0.757,-0.412) -- (0.763,-0.409) -- (0.770,-0.406) -- (0.776,-0.404) -- (0.782,-0.401) -- (0.789,-0.399) -- (0.795,-0.396) -- (0.801,-0.393) -- (0.808,-0.391) -- (0.814,-0.388) -- (0.820,-0.385) -- (0.827,-0.383) -- (0.833,-0.380) -- (0.839,-0.378) -- (0.846,-0.375) -- (0.852,-0.372) -- (0.858,-0.370) -- (0.865,-0.367) -- (0.871,-0.364) -- (0.877,-0.362) -- (0.884,-0.359) -- (0.890,-0.356) -- (0.896,-0.354) -- (0.903,-0.351) -- (0.909,-0.349) -- (0.916,-0.346) -- (0.922,-0.343) -- (0.928,-0.341) -- (0.935,-0.338) -- (0.941,-0.335) -- (0.947,-0.333) -- (0.954,-0.330) -- (0.960,-0.328) -- (0.966,-0.325) -- (0.973,-0.322) -- (0.979,-0.320) -- (0.985,-0.317) -- (0.992,-0.314) -- (0.998,-0.312) -- (1.004,-0.309) -- (1.011,-0.307) -- (1.017,-0.304) -- (1.023,-0.301) -- (1.030,-0.299) -- (1.036,-0.296) -- (1.042,-0.293) -- (1.049,-0.291) -- (1.055,-0.288) -- (1.061,-0.285) -- (1.068,-0.283) -- (1.074,-0.280) -- (1.081,-0.278) -- (1.087,-0.275) -- (1.093,-0.272) -- (1.100,-0.270) -- (1.106,-0.267) -- (1.112,-0.264) -- (1.119,-0.262) -- (1.125,-0.259) -- (1.131,-0.257) -- (1.138,-0.254) -- (1.144,-0.251) -- (1.150,-0.249) -- (1.157,-0.246) -- (1.163,-0.243) -- (1.169,-0.241) -- (1.176,-0.238) -- (1.182,-0.235) -- (1.188,-0.233) -- (1.195,-0.230) -- (1.201,-0.228) -- (1.207,-0.225) -- (1.214,-0.222) -- (1.220,-0.220) -- (1.226,-0.217) -- (1.233,-0.214) -- (1.239,-0.212) -- (1.245,-0.209) -- (1.252,-0.207) -- (1.258,-0.204) -- (1.265,-0.201) -- (1.271,-0.199) -- (1.277,-0.196) -- (1.284,-0.193) -- (1.290,-0.191) -- (1.296,-0.188) -- (1.303,-0.185) -- (1.309,-0.183) -- (1.315,-0.180) -- (1.322,-0.178) -- (1.328,-0.175) -- (1.334,-0.172) -- (1.341,-0.170) -- (1.347,-0.167) -- (1.353,-0.164) -- (1.360,-0.162) -- (1.366,-0.159) -- (1.372,-0.157) -- (1.379,-0.154) -- (1.385,-0.151) -- (1.391,-0.149) -- (1.398,-0.146) -- (1.404,-0.143) -- (1.410,-0.141) -- (1.417,-0.138) -- (1.423,-0.135) -- (1.430,-0.133) -- (1.436,-0.130) -- (1.442,-0.128) -- (1.449,-0.125) -- (1.455,-0.122) -- (1.461,-0.120) -- (1.468,-0.117) -- (1.474,-0.114) -- (1.480,-0.112) -- (1.487,-0.109) -- (1.493,-0.107) -- (1.499,-0.104) -- (1.506,-0.101) -- (1.512,-0.099) -- (1.518,-0.096) -- (1.525,-0.093) -- (1.531,-0.091) -- (1.537,-0.088) -- (1.544,-0.086) -- (1.550,-0.083) -- (1.556,-0.080) -- (1.563,-0.078) -- (1.569,-0.075) -- (1.575,-0.072) -- (1.582,-0.070) -- (1.588,-0.067) -- (1.595,-0.064) -- (1.601,-0.062) -- (1.607,-0.059) -- (1.614,-0.057) -- (1.620,-0.054) -- (1.626,-0.051) -- (1.633,-0.049) -- (1.639,-0.046) -- (1.645,-0.043) -- (1.652,-0.041) -- (1.658,-0.038) -- (1.664,-0.036) -- (1.671,-0.033) -- (1.677,-0.030) -- (1.683,-0.028) -- (1.690,-0.025) -- (1.696,-0.022) -- (1.702,-0.020) -- (1.709,-0.017) -- (1.715,-0.014) -- (1.721,-0.012) -- (1.728,-0.009) -- (1.734,-0.007) -- (1.740,-0.004) -- (1.744,-0.001) -- (1.744,0.001) -- (1.740,0.004) -- (1.734,0.007) -- (1.728,0.009) -- (1.721,0.012) -- (1.715,0.014) -- (1.709,0.017) -- (1.702,0.020) -- (1.696,0.022) -- (1.690,0.025) -- (1.683,0.028) -- (1.677,0.030) -- (1.671,0.033) -- (1.664,0.036) -- (1.658,0.038) -- (1.652,0.041) -- (1.645,0.043) -- (1.639,0.046) -- (1.633,0.049) -- (1.626,0.051) -- (1.620,0.054) -- (1.614,0.057) -- (1.607,0.059) -- (1.601,0.062) -- (1.595,0.064) -- (1.588,0.067) -- (1.582,0.070) -- (1.575,0.072) -- (1.569,0.075) -- (1.563,0.078) -- (1.556,0.080) -- (1.550,0.083) -- (1.544,0.086) -- (1.537,0.088) -- (1.531,0.091) -- (1.525,0.093) -- (1.518,0.096) -- (1.512,0.099) -- (1.506,0.101) -- (1.499,0.104) -- (1.493,0.107) -- (1.487,0.109) -- (1.480,0.112) -- (1.474,0.114) -- (1.468,0.117) -- (1.461,0.120) -- (1.455,0.122) -- (1.449,0.125) -- (1.442,0.128) -- (1.436,0.130) -- (1.430,0.133) -- (1.423,0.135) -- (1.417,0.138) -- (1.410,0.141) -- (1.404,0.143) -- (1.398,0.146) -- (1.391,0.149) -- (1.385,0.151) -- (1.379,0.154) -- (1.372,0.157) -- (1.366,0.159) -- (1.360,0.162) -- (1.353,0.164) -- (1.347,0.167) -- (1.341,0.170) -- (1.334,0.172) -- (1.328,0.175) -- (1.322,0.178) -- (1.315,0.180) -- (1.309,0.183) -- (1.303,0.185) -- (1.296,0.188) -- (1.290,0.191) -- (1.284,0.193) -- (1.277,0.196) -- (1.271,0.199) -- (1.265,0.201) -- (1.258,0.204) -- (1.252,0.207) -- (1.245,0.209) -- (1.239,0.212) -- (1.233,0.214) -- (1.226,0.217) -- (1.220,0.220) -- (1.214,0.222) -- (1.207,0.225) -- (1.201,0.228) -- (1.195,0.230) -- (1.188,0.233) -- (1.182,0.235) -- (1.176,0.238) -- (1.169,0.241) -- (1.163,0.243) -- (1.157,0.246) -- (1.150,0.249) -- (1.144,0.251) -- (1.138,0.254) -- (1.131,0.257) -- (1.125,0.259) -- (1.119,0.262) -- (1.112,0.264) -- (1.106,0.267) -- (1.100,0.270) -- (1.093,0.272) -- (1.087,0.275) -- (1.081,0.278) -- (1.074,0.280) -- (1.068,0.283) -- (1.061,0.285) -- (1.055,0.288) -- (1.049,0.291) -- (1.042,0.293) -- (1.036,0.296) -- (1.030,0.299) -- (1.023,0.301) -- (1.017,0.304) -- (1.011,0.307) -- (1.004,0.309) -- (0.998,0.312) -- (0.992,0.314) -- (0.985,0.317) -- (0.979,0.320) -- (0.973,0.322) -- (0.966,0.325) -- (0.960,0.328) -- (0.954,0.330) -- (0.947,0.333) -- (0.941,0.335) -- (0.935,0.338) -- (0.928,0.341) -- (0.922,0.343) -- (0.916,0.346) -- (0.909,0.349) -- (0.903,0.351) -- (0.896,0.354) -- (0.890,0.356) -- (0.884,0.359) -- (0.877,0.362) -- (0.871,0.364) -- (0.865,0.367) -- (0.858,0.370) -- (0.852,0.372) -- (0.846,0.375) -- (0.839,0.378) -- (0.833,0.380) -- (0.827,0.383) -- (0.820,0.385) -- (0.814,0.388) -- (0.808,0.391) -- (0.801,0.393) -- (0.795,0.396) -- (0.789,0.399) -- (0.782,0.401) -- (0.776,0.404) -- (0.770,0.406) -- (0.763,0.409) -- (0.757,0.412) -- (0.751,0.414) -- (0.744,0.417) -- (0.738,0.420) -- (0.731,0.422) -- (0.725,0.425) -- (0.719,0.428) -- (0.712,0.430) -- (0.706,0.433) -- (0.700,0.435) -- (0.693,0.438) -- (0.687,0.441) -- (0.681,0.443) -- (0.674,0.446) -- (0.668,0.449) -- (0.662,0.451) -- (0.655,0.454) -- (0.649,0.456) -- (0.643,0.459) -- (0.636,0.462) -- (0.630,0.464) -- (0.624,0.467) -- (0.617,0.470) -- (0.611,0.472) -- (0.605,0.475) -- (0.598,0.478) -- (0.592,0.480) -- (0.586,0.483) -- (0.579,0.485) -- (0.573,0.488) -- (0.566,0.491) -- (0.560,0.493) -- (0.554,0.496) -- (0.547,0.499) -- (0.541,0.501) -- (0.535,0.504) -- (0.528,0.506) -- (0.522,0.509) -- (0.516,0.512) -- (0.509,0.514) -- (0.503,0.517) -- (0.497,0.520) -- (0.490,0.522) -- (0.484,0.525) -- (0.478,0.528) -- (0.471,0.530) -- (0.465,0.533) -- (0.459,0.535) -- (0.452,0.538) -- (0.446,0.541) -- (0.440,0.543) -- (0.433,0.546) -- (0.427,0.549) -- (0.421,0.551) -- (0.414,0.554) -- (0.408,0.556) -- (0.401,0.559) -- (0.395,0.562) -- (0.389,0.564) -- (0.382,0.567) -- (0.376,0.570) -- (0.370,0.572) -- (0.363,0.575) -- (0.357,0.577) -- (0.351,0.580) -- (0.344,0.583) -- (0.338,0.585) -- (0.332,0.588) -- (0.325,0.591) -- (0.319,0.593) -- (0.313,0.596) -- (0.306,0.599) -- (0.300,0.601) -- (0.294,0.604) -- (0.287,0.606) -- (0.281,0.609) -- (0.275,0.612) -- (0.268,0.614) -- (0.262,0.617) -- (0.256,0.620) -- (0.249,0.622) -- (0.243,0.625) -- (0.236,0.627) -- (0.230,0.630) -- (0.224,0.633) -- (0.217,0.635) -- (0.211,0.638) -- (0.205,0.641) -- (0.198,0.643) -- (0.192,0.646) -- (0.186,0.649) -- (0.179,0.651) -- (0.173,0.654) -- (0.167,0.656) -- (0.160,0.659) -- (0.154,0.662) -- (0.148,0.664) -- (0.141,0.667) -- (0.135,0.670) -- (0.129,0.672) -- (0.122,0.675) -- (0.116,0.677) -- (0.110,0.680) -- (0.103,0.683) -- (0.097,0.685) -- (0.091,0.688) -- (0.084,0.691) -- (0.078,0.693) -- (0.071,0.696) -- (0.065,0.699) -- (0.059,0.701) -- (0.052,0.704) -- (0.046,0.706) -- (0.040,0.709) -- (0.033,0.712) -- (0.027,0.714) -- (0.021,0.717) -- (0.014,0.720) -- (0.008,0.722) -- (0.002,0.725) -- (-0.005,0.727) -- (-0.011,0.730) -- (-0.017,0.733) -- (-0.024,0.735) -- (-0.030,0.738) -- (-0.036,0.741) -- (-0.043,0.743) -- (-0.049,0.746) -- (-0.055,0.749) -- (-0.062,0.751) -- (-0.068,0.754) -- (-0.074,0.756) -- (-0.081,0.759) -- (-0.087,0.762) -- (-0.093,0.764) -- (-0.100,0.767) -- (-0.106,0.770) -- (-0.113,0.772) -- (-0.119,0.775) -- (-0.125,0.777) -- (-0.132,0.780) -- (-0.138,0.783) -- (-0.144,0.785) -- (-0.151,0.788) -- (-0.157,0.791) -- (-0.163,0.793) -- (-0.170,0.796) -- (-0.176,0.798) -- (-0.182,0.801) -- (-0.189,0.804) -- (-0.195,0.806) -- (-0.201,0.809) -- (-0.208,0.812) -- (-0.214,0.814) -- (-0.220,0.817) -- (-0.227,0.820) -- (-0.233,0.822) -- (-0.239,0.825) -- (-0.246,0.827) -- (-0.251,0.830) -- (-0.256,0.832) -- cycle}

      \fill[top color=azulvelo!70!white, bottom color=metal] \formagota;
      % Conductor con punta - densidad superficial
% Generado por scripts/generar_datos.py (no editar a mano).
\SG{-0.259}{0.833}{-0.262}{0.834}{0.183}
\SG{-0.262}{0.834}{-0.265}{0.835}{0.181}
\SG{-0.265}{0.835}{-0.268}{0.836}{0.179}
\SG{-0.268}{0.836}{-0.270}{0.837}{0.179}
\SG{-0.270}{0.837}{-0.273}{0.839}{0.181}
\SG{-0.273}{0.839}{-0.276}{0.840}{0.183}
\SG{-0.276}{0.840}{-0.279}{0.841}{0.185}
\SG{-0.279}{0.841}{-0.281}{0.842}{0.187}
\SG{-0.281}{0.842}{-0.284}{0.843}{0.190}
\SG{-0.284}{0.843}{-0.287}{0.844}{0.194}
\SG{-0.287}{0.844}{-0.290}{0.845}{0.197}
\SG{-0.290}{0.845}{-0.293}{0.846}{0.201}
\SG{-0.293}{0.846}{-0.295}{0.847}{0.203}
\SG{-0.295}{0.847}{-0.298}{0.848}{0.206}
\SG{-0.298}{0.848}{-0.301}{0.849}{0.203}
\SG{-0.301}{0.849}{-0.304}{0.850}{0.201}
\SG{-0.304}{0.850}{-0.306}{0.851}{0.200}
\SG{-0.306}{0.851}{-0.309}{0.852}{0.199}
\SG{-0.309}{0.852}{-0.312}{0.853}{0.199}
\SG{-0.312}{0.853}{-0.315}{0.854}{0.201}
\SG{-0.315}{0.854}{-0.318}{0.855}{0.204}
\SG{-0.318}{0.855}{-0.320}{0.855}{0.206}
\SG{-0.320}{0.855}{-0.323}{0.856}{0.209}
\SG{-0.323}{0.856}{-0.326}{0.857}{0.212}
\SG{-0.326}{0.857}{-0.329}{0.858}{0.216}
\SG{-0.329}{0.858}{-0.332}{0.859}{0.221}
\SG{-0.332}{0.859}{-0.335}{0.860}{0.225}
\SG{-0.335}{0.860}{-0.337}{0.861}{0.229}
\SG{-0.337}{0.861}{-0.340}{0.862}{0.233}
\SG{-0.340}{0.862}{-0.343}{0.863}{0.234}
\SG{-0.343}{0.863}{-0.346}{0.863}{0.232}
\SG{-0.346}{0.863}{-0.349}{0.864}{0.231}
\SG{-0.349}{0.864}{-0.352}{0.865}{0.231}
\SG{-0.352}{0.865}{-0.354}{0.866}{0.231}
\SG{-0.354}{0.866}{-0.357}{0.867}{0.232}
\SG{-0.357}{0.867}{-0.360}{0.867}{0.236}
\SG{-0.360}{0.867}{-0.363}{0.868}{0.239}
\SG{-0.363}{0.868}{-0.366}{0.869}{0.242}
\SG{-0.366}{0.869}{-0.369}{0.870}{0.246}
\SG{-0.369}{0.870}{-0.371}{0.870}{0.250}
\SG{-0.371}{0.870}{-0.374}{0.871}{0.257}
\SG{-0.374}{0.871}{-0.377}{0.872}{0.262}
\SG{-0.377}{0.872}{-0.380}{0.873}{0.266}
\SG{-0.380}{0.873}{-0.383}{0.873}{0.271}
\SG{-0.383}{0.873}{-0.386}{0.874}{0.273}
\SG{-0.386}{0.874}{-0.389}{0.875}{0.269}
\SG{-0.389}{0.875}{-0.392}{0.876}{0.266}
\SG{-0.392}{0.876}{-0.394}{0.876}{0.262}
\SG{-0.394}{0.876}{-0.397}{0.877}{0.259}
\SG{-0.397}{0.877}{-0.400}{0.878}{0.257}
\SG{-0.400}{0.878}{-0.403}{0.878}{0.253}
\SG{-0.403}{0.878}{-0.406}{0.879}{0.251}
\SG{-0.406}{0.879}{-0.409}{0.879}{0.248}
\SG{-0.409}{0.879}{-0.412}{0.880}{0.246}
\SG{-0.412}{0.880}{-0.415}{0.881}{0.244}
\SG{-0.415}{0.881}{-0.418}{0.881}{0.244}
\SG{-0.418}{0.881}{-0.420}{0.882}{0.245}
\SG{-0.420}{0.882}{-0.423}{0.882}{0.245}
\SG{-0.423}{0.882}{-0.426}{0.883}{0.246}
\SG{-0.426}{0.883}{-0.429}{0.884}{0.247}
\SG{-0.429}{0.884}{-0.432}{0.884}{0.252}
\SG{-0.432}{0.884}{-0.435}{0.885}{0.257}
\SG{-0.435}{0.885}{-0.438}{0.885}{0.262}
\SG{-0.438}{0.885}{-0.441}{0.886}{0.266}
\SG{-0.441}{0.886}{-0.444}{0.886}{0.271}
\SG{-0.444}{0.886}{-0.447}{0.887}{0.279}
\SG{-0.447}{0.887}{-0.449}{0.887}{0.286}
\SG{-0.449}{0.887}{-0.452}{0.888}{0.292}
\SG{-0.452}{0.888}{-0.455}{0.888}{0.299}
\SG{-0.455}{0.888}{-0.458}{0.889}{0.303}
\SG{-0.458}{0.889}{-0.461}{0.889}{0.302}
\SG{-0.461}{0.889}{-0.464}{0.890}{0.301}
\SG{-0.464}{0.890}{-0.467}{0.890}{0.299}
\SG{-0.467}{0.890}{-0.470}{0.891}{0.298}
\SG{-0.470}{0.891}{-0.473}{0.891}{0.297}
\SG{-0.473}{0.891}{-0.476}{0.891}{0.295}
\SG{-0.476}{0.891}{-0.479}{0.892}{0.293}
\SG{-0.479}{0.892}{-0.482}{0.892}{0.290}
\SG{-0.482}{0.892}{-0.485}{0.893}{0.289}
\SG{-0.485}{0.893}{-0.487}{0.893}{0.287}
\SG{-0.487}{0.893}{-0.490}{0.893}{0.285}
\SG{-0.490}{0.893}{-0.493}{0.894}{0.283}
\SG{-0.493}{0.894}{-0.496}{0.894}{0.281}
\SG{-0.496}{0.894}{-0.499}{0.894}{0.280}
\SG{-0.499}{0.894}{-0.502}{0.895}{0.278}
\SG{-0.502}{0.895}{-0.505}{0.895}{0.277}
\SG{-0.505}{0.895}{-0.508}{0.895}{0.276}
\SG{-0.508}{0.895}{-0.511}{0.896}{0.275}
\SG{-0.511}{0.896}{-0.514}{0.896}{0.273}
\SG{-0.514}{0.896}{-0.517}{0.896}{0.272}
\SG{-0.517}{0.896}{-0.520}{0.896}{0.271}
\SG{-0.520}{0.896}{-0.523}{0.897}{0.271}
\SG{-0.523}{0.897}{-0.526}{0.897}{0.270}
\SG{-0.526}{0.897}{-0.529}{0.897}{0.269}
\SG{-0.529}{0.897}{-0.532}{0.897}{0.268}
\SG{-0.532}{0.897}{-0.535}{0.898}{0.268}
\SG{-0.535}{0.898}{-0.537}{0.898}{0.267}
\SG{-0.537}{0.898}{-0.540}{0.898}{0.267}
\SG{-0.540}{0.898}{-0.543}{0.898}{0.266}
\SG{-0.543}{0.898}{-0.546}{0.898}{0.266}
\SG{-0.546}{0.898}{-0.549}{0.899}{0.266}
\SG{-0.549}{0.899}{-0.552}{0.899}{0.266}
\SG{-0.552}{0.899}{-0.555}{0.899}{0.265}
\SG{-0.555}{0.899}{-0.558}{0.899}{0.265}
\SG{-0.558}{0.899}{-0.561}{0.899}{0.265}
\SG{-0.561}{0.899}{-0.564}{0.899}{0.265}
\SG{-0.564}{0.899}{-0.567}{0.899}{0.265}
\SG{-0.567}{0.899}{-0.570}{0.899}{0.265}
\SG{-0.570}{0.899}{-0.573}{0.900}{0.265}
\SG{-0.573}{0.900}{-0.576}{0.900}{0.265}
\SG{-0.576}{0.900}{-0.579}{0.900}{0.266}
\SG{-0.579}{0.900}{-0.582}{0.900}{0.266}
\SG{-0.582}{0.900}{-0.585}{0.900}{0.266}
\SG{-0.585}{0.900}{-0.588}{0.900}{0.267}
\SG{-0.588}{0.900}{-0.591}{0.900}{0.267}
\SG{-0.591}{0.900}{-0.594}{0.900}{0.268}
\SG{-0.594}{0.900}{-0.597}{0.900}{0.268}
\SG{-0.597}{0.900}{-0.599}{0.900}{0.269}
\SG{-0.599}{0.900}{-0.602}{0.900}{0.269}
\SG{-0.602}{0.900}{-0.605}{0.900}{0.270}
\SG{-0.605}{0.900}{-0.608}{0.900}{0.270}
\SG{-0.608}{0.900}{-0.611}{0.900}{0.271}
\SG{-0.611}{0.900}{-0.614}{0.900}{0.272}
\SG{-0.614}{0.900}{-0.617}{0.900}{0.273}
\SG{-0.617}{0.900}{-0.620}{0.900}{0.273}
\SG{-0.620}{0.900}{-0.623}{0.900}{0.274}
\SG{-0.623}{0.900}{-0.626}{0.900}{0.275}
\SG{-0.626}{0.900}{-0.629}{0.900}{0.276}
\SG{-0.629}{0.900}{-0.632}{0.899}{0.277}
\SG{-0.632}{0.899}{-0.635}{0.899}{0.278}
\SG{-0.635}{0.899}{-0.638}{0.899}{0.280}
\SG{-0.638}{0.899}{-0.641}{0.899}{0.281}
\SG{-0.641}{0.899}{-0.644}{0.899}{0.282}
\SG{-0.644}{0.899}{-0.647}{0.899}{0.283}
\SG{-0.647}{0.899}{-0.650}{0.899}{0.284}
\SG{-0.650}{0.899}{-0.653}{0.898}{0.286}
\SG{-0.653}{0.898}{-0.656}{0.898}{0.287}
\SG{-0.656}{0.898}{-0.659}{0.898}{0.289}
\SG{-0.659}{0.898}{-0.661}{0.898}{0.290}
\SG{-0.661}{0.898}{-0.664}{0.898}{0.292}
\SG{-0.664}{0.898}{-0.667}{0.897}{0.294}
\SG{-0.667}{0.897}{-0.670}{0.897}{0.295}
\SG{-0.670}{0.897}{-0.673}{0.897}{0.297}
\SG{-0.673}{0.897}{-0.676}{0.897}{0.299}
\SG{-0.676}{0.897}{-0.679}{0.896}{0.301}
\SG{-0.679}{0.896}{-0.682}{0.896}{0.303}
\SG{-0.682}{0.896}{-0.685}{0.896}{0.305}
\SG{-0.685}{0.896}{-0.688}{0.896}{0.307}
\SG{-0.688}{0.896}{-0.691}{0.895}{0.310}
\SG{-0.691}{0.895}{-0.694}{0.895}{0.312}
\SG{-0.694}{0.895}{-0.697}{0.895}{0.315}
\SG{-0.697}{0.895}{-0.700}{0.894}{0.317}
\SG{-0.700}{0.894}{-0.703}{0.894}{0.320}
\SG{-0.703}{0.894}{-0.706}{0.894}{0.322}
\SG{-0.706}{0.894}{-0.709}{0.893}{0.325}
\SG{-0.709}{0.893}{-0.711}{0.893}{0.328}
\SG{-0.711}{0.893}{-0.714}{0.893}{0.331}
\SG{-0.714}{0.893}{-0.717}{0.892}{0.334}
\SG{-0.717}{0.892}{-0.720}{0.892}{0.338}
\SG{-0.720}{0.892}{-0.723}{0.892}{0.341}
\SG{-0.723}{0.892}{-0.726}{0.891}{0.344}
\SG{-0.726}{0.891}{-0.729}{0.891}{0.348}
\SG{-0.729}{0.891}{-0.732}{0.890}{0.351}
\SG{-0.732}{0.890}{-0.735}{0.890}{0.354}
\SG{-0.735}{0.890}{-0.738}{0.889}{0.357}
\SG{-0.738}{0.889}{-0.741}{0.889}{0.359}
\SG{-0.741}{0.889}{-0.744}{0.888}{0.362}
\SG{-0.744}{0.888}{-0.747}{0.888}{0.365}
\SG{-0.747}{0.888}{-0.750}{0.887}{0.364}
\SG{-0.750}{0.887}{-0.752}{0.887}{0.359}
\SG{-0.752}{0.887}{-0.755}{0.886}{0.353}
\SG{-0.755}{0.886}{-0.758}{0.886}{0.348}
\SG{-0.758}{0.886}{-0.761}{0.885}{0.342}
\SG{-0.761}{0.885}{-0.764}{0.885}{0.338}
\SG{-0.764}{0.885}{-0.767}{0.884}{0.334}
\SG{-0.767}{0.884}{-0.770}{0.884}{0.331}
\SG{-0.770}{0.884}{-0.773}{0.883}{0.328}
\SG{-0.773}{0.883}{-0.776}{0.883}{0.324}
\SG{-0.776}{0.883}{-0.779}{0.882}{0.323}
\SG{-0.779}{0.882}{-0.781}{0.881}{0.324}
\SG{-0.781}{0.881}{-0.784}{0.881}{0.325}
\SG{-0.784}{0.881}{-0.787}{0.880}{0.326}
\SG{-0.787}{0.880}{-0.790}{0.880}{0.327}
\SG{-0.790}{0.880}{-0.793}{0.879}{0.330}
\SG{-0.793}{0.879}{-0.796}{0.878}{0.334}
\SG{-0.796}{0.878}{-0.799}{0.878}{0.338}
\SG{-0.799}{0.878}{-0.802}{0.877}{0.342}
\SG{-0.802}{0.877}{-0.805}{0.876}{0.346}
\SG{-0.805}{0.876}{-0.807}{0.876}{0.351}
\SG{-0.807}{0.876}{-0.810}{0.875}{0.355}
\SG{-0.810}{0.875}{-0.813}{0.874}{0.360}
\SG{-0.813}{0.874}{-0.816}{0.874}{0.365}
\SG{-0.816}{0.874}{-0.819}{0.873}{0.370}
\SG{-0.819}{0.873}{-0.822}{0.872}{0.373}
\SG{-0.822}{0.872}{-0.825}{0.871}{0.370}
\SG{-0.825}{0.871}{-0.828}{0.871}{0.367}
\SG{-0.828}{0.871}{-0.830}{0.870}{0.363}
\SG{-0.830}{0.870}{-0.833}{0.869}{0.359}
\SG{-0.833}{0.869}{-0.836}{0.868}{0.355}
\SG{-0.836}{0.868}{-0.839}{0.868}{0.353}
\SG{-0.839}{0.868}{-0.842}{0.867}{0.352}
\SG{-0.842}{0.867}{-0.845}{0.866}{0.350}
\SG{-0.845}{0.866}{-0.847}{0.865}{0.349}
\SG{-0.847}{0.865}{-0.850}{0.864}{0.348}
\SG{-0.850}{0.864}{-0.853}{0.864}{0.350}
\SG{-0.853}{0.864}{-0.856}{0.863}{0.352}
\SG{-0.856}{0.863}{-0.859}{0.862}{0.355}
\SG{-0.859}{0.862}{-0.862}{0.861}{0.358}
\SG{-0.862}{0.861}{-0.864}{0.860}{0.360}
\SG{-0.864}{0.860}{-0.867}{0.859}{0.358}
\SG{-0.867}{0.859}{-0.870}{0.858}{0.356}
\SG{-0.870}{0.858}{-0.873}{0.858}{0.354}
\SG{-0.873}{0.858}{-0.876}{0.857}{0.351}
\SG{-0.876}{0.857}{-0.879}{0.856}{0.349}
\SG{-0.879}{0.856}{-0.881}{0.855}{0.348}
\SG{-0.881}{0.855}{-0.884}{0.854}{0.347}
\SG{-0.884}{0.854}{-0.887}{0.853}{0.346}
\SG{-0.887}{0.853}{-0.890}{0.852}{0.345}
\SG{-0.890}{0.852}{-0.893}{0.851}{0.345}
\SG{-0.893}{0.851}{-0.895}{0.850}{0.347}
\SG{-0.895}{0.850}{-0.898}{0.849}{0.350}
\SG{-0.898}{0.849}{-0.901}{0.848}{0.353}
\SG{-0.901}{0.848}{-0.904}{0.847}{0.357}
\SG{-0.904}{0.847}{-0.907}{0.846}{0.362}
\SG{-0.907}{0.846}{-0.909}{0.845}{0.363}
\SG{-0.909}{0.845}{-0.912}{0.844}{0.362}
\SG{-0.912}{0.844}{-0.915}{0.843}{0.361}
\SG{-0.915}{0.843}{-0.918}{0.842}{0.360}
\SG{-0.918}{0.842}{-0.920}{0.841}{0.358}
\SG{-0.920}{0.841}{-0.923}{0.840}{0.356}
\SG{-0.923}{0.840}{-0.926}{0.839}{0.356}
\SG{-0.926}{0.839}{-0.929}{0.838}{0.355}
\SG{-0.929}{0.838}{-0.931}{0.837}{0.355}
\SG{-0.931}{0.837}{-0.934}{0.836}{0.355}
\SG{-0.934}{0.836}{-0.937}{0.835}{0.354}
\SG{-0.937}{0.835}{-0.940}{0.833}{0.353}
\SG{-0.940}{0.833}{-0.942}{0.832}{0.351}
\SG{-0.942}{0.832}{-0.945}{0.831}{0.350}
\SG{-0.945}{0.831}{-0.948}{0.830}{0.348}
\SG{-0.948}{0.830}{-0.951}{0.829}{0.345}
\SG{-0.951}{0.829}{-0.953}{0.828}{0.345}
\SG{-0.953}{0.828}{-0.956}{0.827}{0.346}
\SG{-0.956}{0.827}{-0.959}{0.825}{0.346}
\SG{-0.959}{0.825}{-0.961}{0.824}{0.346}
\SG{-0.961}{0.824}{-0.964}{0.823}{0.346}
\SG{-0.964}{0.823}{-0.967}{0.822}{0.348}
\SG{-0.967}{0.822}{-0.969}{0.821}{0.352}
\SG{-0.969}{0.821}{-0.972}{0.819}{0.357}
\SG{-0.972}{0.819}{-0.975}{0.818}{0.362}
\SG{-0.975}{0.818}{-0.978}{0.817}{0.369}
\SG{-0.978}{0.817}{-0.980}{0.816}{0.375}
\SG{-0.980}{0.816}{-0.983}{0.814}{0.377}
\SG{-0.983}{0.814}{-0.986}{0.813}{0.378}
\SG{-0.986}{0.813}{-0.988}{0.812}{0.379}
\SG{-0.988}{0.812}{-0.991}{0.811}{0.380}
\SG{-0.991}{0.811}{-0.994}{0.809}{0.380}
\SG{-0.994}{0.809}{-0.996}{0.808}{0.381}
\SG{-0.996}{0.808}{-0.999}{0.807}{0.383}
\SG{-0.999}{0.807}{-1.002}{0.805}{0.385}
\SG{-1.002}{0.805}{-1.004}{0.804}{0.388}
\SG{-1.004}{0.804}{-1.007}{0.803}{0.391}
\SG{-1.007}{0.803}{-1.009}{0.801}{0.394}
\SG{-1.009}{0.801}{-1.012}{0.800}{0.393}
\SG{-1.012}{0.800}{-1.015}{0.799}{0.392}
\SG{-1.015}{0.799}{-1.017}{0.797}{0.391}
\SG{-1.017}{0.797}{-1.020}{0.796}{0.390}
\SG{-1.020}{0.796}{-1.023}{0.795}{0.389}
\SG{-1.023}{0.795}{-1.025}{0.793}{0.386}
\SG{-1.025}{0.793}{-1.028}{0.792}{0.383}
\SG{-1.028}{0.792}{-1.030}{0.790}{0.380}
\SG{-1.030}{0.790}{-1.033}{0.789}{0.377}
\SG{-1.033}{0.789}{-1.036}{0.788}{0.373}
\SG{-1.036}{0.788}{-1.038}{0.786}{0.368}
\SG{-1.038}{0.786}{-1.041}{0.785}{0.367}
\SG{-1.041}{0.785}{-1.043}{0.783}{0.366}
\SG{-1.043}{0.783}{-1.046}{0.782}{0.365}
\SG{-1.046}{0.782}{-1.048}{0.780}{0.363}
\SG{-1.048}{0.780}{-1.051}{0.779}{0.363}
\SG{-1.051}{0.779}{-1.054}{0.777}{0.363}
\SG{-1.054}{0.777}{-1.056}{0.776}{0.363}
\SG{-1.056}{0.776}{-1.059}{0.774}{0.363}
\SG{-1.059}{0.774}{-1.061}{0.773}{0.363}
\SG{-1.061}{0.773}{-1.064}{0.771}{0.363}
\SG{-1.064}{0.771}{-1.066}{0.770}{0.363}
\SG{-1.066}{0.770}{-1.069}{0.768}{0.365}
\SG{-1.069}{0.768}{-1.071}{0.767}{0.367}
\SG{-1.071}{0.767}{-1.074}{0.765}{0.369}
\SG{-1.074}{0.765}{-1.076}{0.764}{0.372}
\SG{-1.076}{0.764}{-1.079}{0.762}{0.377}
\SG{-1.079}{0.762}{-1.081}{0.760}{0.382}
\SG{-1.081}{0.760}{-1.084}{0.759}{0.383}
\SG{-1.084}{0.759}{-1.086}{0.757}{0.385}
\SG{-1.086}{0.757}{-1.089}{0.756}{0.386}
\SG{-1.089}{0.756}{-1.091}{0.754}{0.388}
\SG{-1.091}{0.754}{-1.094}{0.752}{0.390}
\SG{-1.094}{0.752}{-1.096}{0.751}{0.391}
\SG{-1.096}{0.751}{-1.099}{0.749}{0.391}
\SG{-1.099}{0.749}{-1.101}{0.748}{0.391}
\SG{-1.101}{0.748}{-1.104}{0.746}{0.390}
\SG{-1.104}{0.746}{-1.106}{0.744}{0.388}
\SG{-1.106}{0.744}{-1.108}{0.743}{0.386}
\SG{-1.108}{0.743}{-1.111}{0.741}{0.384}
\SG{-1.111}{0.741}{-1.113}{0.739}{0.382}
\SG{-1.113}{0.739}{-1.116}{0.738}{0.380}
\SG{-1.116}{0.738}{-1.118}{0.736}{0.377}
\SG{-1.118}{0.736}{-1.121}{0.734}{0.374}
\SG{-1.121}{0.734}{-1.123}{0.732}{0.371}
\SG{-1.123}{0.732}{-1.125}{0.731}{0.368}
\SG{-1.125}{0.731}{-1.128}{0.729}{0.367}
\SG{-1.128}{0.729}{-1.130}{0.727}{0.367}
\SG{-1.130}{0.727}{-1.133}{0.725}{0.368}
\SG{-1.133}{0.725}{-1.135}{0.724}{0.368}
\SG{-1.135}{0.724}{-1.137}{0.722}{0.369}
\SG{-1.137}{0.722}{-1.140}{0.720}{0.372}
\SG{-1.140}{0.720}{-1.142}{0.718}{0.373}
\SG{-1.142}{0.718}{-1.144}{0.717}{0.375}
\SG{-1.144}{0.717}{-1.147}{0.715}{0.377}
\SG{-1.147}{0.715}{-1.149}{0.713}{0.379}
\SG{-1.149}{0.713}{-1.151}{0.711}{0.381}
\SG{-1.151}{0.711}{-1.154}{0.709}{0.384}
\SG{-1.154}{0.709}{-1.156}{0.708}{0.386}
\SG{-1.156}{0.708}{-1.158}{0.706}{0.388}
\SG{-1.158}{0.706}{-1.161}{0.704}{0.390}
\SG{-1.161}{0.704}{-1.163}{0.702}{0.391}
\SG{-1.163}{0.702}{-1.165}{0.700}{0.391}
\SG{-1.165}{0.700}{-1.168}{0.698}{0.392}
\SG{-1.168}{0.698}{-1.170}{0.697}{0.392}
\SG{-1.170}{0.697}{-1.172}{0.695}{0.391}
\SG{-1.172}{0.695}{-1.174}{0.693}{0.391}
\SG{-1.174}{0.693}{-1.177}{0.691}{0.390}
\SG{-1.177}{0.691}{-1.179}{0.689}{0.388}
\SG{-1.179}{0.689}{-1.181}{0.687}{0.387}
\SG{-1.181}{0.687}{-1.184}{0.685}{0.386}
\SG{-1.184}{0.685}{-1.186}{0.683}{0.385}
\SG{-1.186}{0.683}{-1.188}{0.681}{0.384}
\SG{-1.188}{0.681}{-1.190}{0.679}{0.382}
\SG{-1.190}{0.679}{-1.192}{0.677}{0.380}
\SG{-1.192}{0.677}{-1.195}{0.675}{0.379}
\SG{-1.195}{0.675}{-1.197}{0.674}{0.377}
\SG{-1.197}{0.674}{-1.199}{0.672}{0.376}
\SG{-1.199}{0.672}{-1.201}{0.670}{0.375}
\SG{-1.201}{0.670}{-1.204}{0.668}{0.374}
\SG{-1.204}{0.668}{-1.206}{0.666}{0.373}
\SG{-1.206}{0.666}{-1.208}{0.664}{0.372}
\SG{-1.208}{0.664}{-1.210}{0.662}{0.371}
\SG{-1.210}{0.662}{-1.212}{0.660}{0.370}
\SG{-1.212}{0.660}{-1.214}{0.658}{0.370}
\SG{-1.214}{0.658}{-1.217}{0.656}{0.369}
\SG{-1.217}{0.656}{-1.219}{0.654}{0.369}
\SG{-1.219}{0.654}{-1.221}{0.652}{0.368}
\SG{-1.221}{0.652}{-1.223}{0.649}{0.368}
\SG{-1.223}{0.649}{-1.225}{0.647}{0.367}
\SG{-1.225}{0.647}{-1.227}{0.645}{0.367}
\SG{-1.227}{0.645}{-1.229}{0.643}{0.367}
\SG{-1.229}{0.643}{-1.231}{0.641}{0.367}
\SG{-1.231}{0.641}{-1.234}{0.639}{0.367}
\SG{-1.234}{0.639}{-1.236}{0.637}{0.367}
\SG{-1.236}{0.637}{-1.238}{0.635}{0.367}
\SG{-1.238}{0.635}{-1.240}{0.633}{0.367}
\SG{-1.240}{0.633}{-1.242}{0.631}{0.367}
\SG{-1.242}{0.631}{-1.244}{0.629}{0.368}
\SG{-1.244}{0.629}{-1.246}{0.627}{0.368}
\SG{-1.246}{0.627}{-1.248}{0.624}{0.368}
\SG{-1.248}{0.624}{-1.250}{0.622}{0.369}
\SG{-1.250}{0.622}{-1.252}{0.620}{0.370}
\SG{-1.252}{0.620}{-1.254}{0.618}{0.370}
\SG{-1.254}{0.618}{-1.256}{0.616}{0.371}
\SG{-1.256}{0.616}{-1.258}{0.614}{0.372}
\SG{-1.258}{0.614}{-1.260}{0.612}{0.373}
\SG{-1.260}{0.612}{-1.262}{0.609}{0.374}
\SG{-1.262}{0.609}{-1.264}{0.607}{0.375}
\SG{-1.264}{0.607}{-1.266}{0.605}{0.376}
\SG{-1.266}{0.605}{-1.268}{0.603}{0.377}
\SG{-1.268}{0.603}{-1.270}{0.601}{0.379}
\SG{-1.270}{0.601}{-1.272}{0.598}{0.380}
\SG{-1.272}{0.598}{-1.274}{0.596}{0.381}
\SG{-1.274}{0.596}{-1.276}{0.594}{0.382}
\SG{-1.276}{0.594}{-1.278}{0.592}{0.384}
\SG{-1.278}{0.592}{-1.280}{0.590}{0.386}
\SG{-1.280}{0.590}{-1.282}{0.587}{0.388}
\SG{-1.282}{0.587}{-1.284}{0.585}{0.390}
\SG{-1.284}{0.585}{-1.286}{0.583}{0.392}
\SG{-1.286}{0.583}{-1.288}{0.581}{0.393}
\SG{-1.288}{0.581}{-1.290}{0.578}{0.394}
\SG{-1.290}{0.578}{-1.291}{0.576}{0.396}
\SG{-1.291}{0.576}{-1.293}{0.574}{0.398}
\SG{-1.293}{0.574}{-1.295}{0.571}{0.399}
\SG{-1.295}{0.571}{-1.297}{0.569}{0.400}
\SG{-1.297}{0.569}{-1.299}{0.567}{0.401}
\SG{-1.299}{0.567}{-1.301}{0.565}{0.402}
\SG{-1.301}{0.565}{-1.303}{0.562}{0.402}
\SG{-1.303}{0.562}{-1.305}{0.560}{0.402}
\SG{-1.305}{0.560}{-1.306}{0.558}{0.402}
\SG{-1.306}{0.558}{-1.308}{0.555}{0.401}
\SG{-1.308}{0.555}{-1.310}{0.553}{0.400}
\SG{-1.310}{0.553}{-1.312}{0.551}{0.398}
\SG{-1.312}{0.551}{-1.314}{0.548}{0.396}
\SG{-1.314}{0.548}{-1.315}{0.546}{0.393}
\SG{-1.315}{0.546}{-1.317}{0.544}{0.391}
\SG{-1.317}{0.544}{-1.319}{0.541}{0.389}
\SG{-1.319}{0.541}{-1.321}{0.539}{0.388}
\SG{-1.321}{0.539}{-1.323}{0.537}{0.387}
\SG{-1.323}{0.537}{-1.324}{0.534}{0.385}
\SG{-1.324}{0.534}{-1.326}{0.532}{0.384}
\SG{-1.326}{0.532}{-1.328}{0.529}{0.383}
\SG{-1.328}{0.529}{-1.330}{0.527}{0.383}
\SG{-1.330}{0.527}{-1.331}{0.525}{0.383}
\SG{-1.331}{0.525}{-1.333}{0.522}{0.384}
\SG{-1.333}{0.522}{-1.335}{0.520}{0.385}
\SG{-1.335}{0.520}{-1.336}{0.517}{0.389}
\SG{-1.336}{0.517}{-1.338}{0.515}{0.392}
\SG{-1.338}{0.515}{-1.340}{0.513}{0.396}
\SG{-1.340}{0.513}{-1.341}{0.510}{0.399}
\SG{-1.341}{0.510}{-1.343}{0.508}{0.401}
\SG{-1.343}{0.508}{-1.345}{0.505}{0.403}
\SG{-1.345}{0.505}{-1.346}{0.503}{0.406}
\SG{-1.346}{0.503}{-1.348}{0.500}{0.408}
\SG{-1.348}{0.500}{-1.350}{0.498}{0.410}
\SG{-1.350}{0.498}{-1.351}{0.495}{0.411}
\SG{-1.351}{0.495}{-1.353}{0.493}{0.412}
\SG{-1.353}{0.493}{-1.355}{0.490}{0.411}
\SG{-1.355}{0.490}{-1.356}{0.488}{0.410}
\SG{-1.356}{0.488}{-1.358}{0.485}{0.408}
\SG{-1.358}{0.485}{-1.359}{0.483}{0.407}
\SG{-1.359}{0.483}{-1.361}{0.480}{0.406}
\SG{-1.361}{0.480}{-1.363}{0.478}{0.405}
\SG{-1.363}{0.478}{-1.364}{0.475}{0.402}
\SG{-1.364}{0.475}{-1.366}{0.473}{0.398}
\SG{-1.366}{0.473}{-1.367}{0.470}{0.394}
\SG{-1.367}{0.470}{-1.369}{0.468}{0.392}
\SG{-1.369}{0.468}{-1.370}{0.465}{0.390}
\SG{-1.370}{0.465}{-1.372}{0.463}{0.388}
\SG{-1.372}{0.463}{-1.373}{0.460}{0.388}
\SG{-1.373}{0.460}{-1.375}{0.458}{0.388}
\SG{-1.375}{0.458}{-1.376}{0.455}{0.388}
\SG{-1.376}{0.455}{-1.378}{0.453}{0.388}
\SG{-1.378}{0.453}{-1.379}{0.450}{0.389}
\SG{-1.379}{0.450}{-1.381}{0.448}{0.389}
\SG{-1.381}{0.448}{-1.382}{0.445}{0.389}
\SG{-1.382}{0.445}{-1.384}{0.442}{0.391}
\SG{-1.384}{0.442}{-1.385}{0.440}{0.393}
\SG{-1.385}{0.440}{-1.387}{0.437}{0.394}
\SG{-1.387}{0.437}{-1.388}{0.435}{0.395}
\SG{-1.388}{0.435}{-1.389}{0.432}{0.398}
\SG{-1.389}{0.432}{-1.391}{0.430}{0.403}
\SG{-1.391}{0.430}{-1.392}{0.427}{0.407}
\SG{-1.392}{0.427}{-1.394}{0.424}{0.411}
\SG{-1.394}{0.424}{-1.395}{0.422}{0.414}
\SG{-1.395}{0.422}{-1.396}{0.419}{0.417}
\SG{-1.396}{0.419}{-1.398}{0.416}{0.420}
\SG{-1.398}{0.416}{-1.399}{0.414}{0.422}
\SG{-1.399}{0.414}{-1.401}{0.411}{0.423}
\SG{-1.401}{0.411}{-1.402}{0.409}{0.424}
\SG{-1.402}{0.409}{-1.403}{0.406}{0.426}
\SG{-1.403}{0.406}{-1.405}{0.403}{0.426}
\SG{-1.405}{0.403}{-1.406}{0.401}{0.423}
\SG{-1.406}{0.401}{-1.407}{0.398}{0.420}
\SG{-1.407}{0.398}{-1.408}{0.395}{0.418}
\SG{-1.408}{0.395}{-1.410}{0.393}{0.416}
\SG{-1.410}{0.393}{-1.411}{0.390}{0.414}
\SG{-1.411}{0.390}{-1.412}{0.387}{0.415}
\SG{-1.412}{0.387}{-1.414}{0.385}{0.415}
\SG{-1.414}{0.385}{-1.415}{0.382}{0.414}
\SG{-1.415}{0.382}{-1.416}{0.379}{0.413}
\SG{-1.416}{0.379}{-1.417}{0.377}{0.412}
\SG{-1.417}{0.377}{-1.419}{0.374}{0.409}
\SG{-1.419}{0.374}{-1.420}{0.371}{0.403}
\SG{-1.420}{0.371}{-1.421}{0.369}{0.397}
\SG{-1.421}{0.369}{-1.422}{0.366}{0.392}
\SG{-1.422}{0.366}{-1.423}{0.363}{0.388}
\SG{-1.423}{0.363}{-1.425}{0.361}{0.385}
\SG{-1.425}{0.361}{-1.426}{0.358}{0.384}
\SG{-1.426}{0.358}{-1.427}{0.355}{0.385}
\SG{-1.427}{0.355}{-1.428}{0.352}{0.385}
\SG{-1.428}{0.352}{-1.429}{0.350}{0.385}
\SG{-1.429}{0.350}{-1.430}{0.347}{0.385}
\SG{-1.430}{0.347}{-1.432}{0.344}{0.387}
\SG{-1.432}{0.344}{-1.433}{0.341}{0.389}
\SG{-1.433}{0.341}{-1.434}{0.339}{0.392}
\SG{-1.434}{0.339}{-1.435}{0.336}{0.394}
\SG{-1.435}{0.336}{-1.436}{0.333}{0.396}
\SG{-1.436}{0.333}{-1.437}{0.330}{0.397}
\SG{-1.437}{0.330}{-1.438}{0.328}{0.397}
\SG{-1.438}{0.328}{-1.439}{0.325}{0.398}
\SG{-1.439}{0.325}{-1.440}{0.322}{0.399}
\SG{-1.440}{0.322}{-1.441}{0.319}{0.400}
\SG{-1.441}{0.319}{-1.442}{0.317}{0.401}
\SG{-1.442}{0.317}{-1.443}{0.314}{0.403}
\SG{-1.443}{0.314}{-1.444}{0.311}{0.405}
\SG{-1.444}{0.311}{-1.445}{0.308}{0.407}
\SG{-1.445}{0.308}{-1.447}{0.306}{0.409}
\SG{-1.447}{0.306}{-1.448}{0.303}{0.410}
\SG{-1.448}{0.303}{-1.448}{0.300}{0.407}
\SG{-1.448}{0.300}{-1.449}{0.297}{0.403}
\SG{-1.449}{0.297}{-1.450}{0.294}{0.399}
\SG{-1.450}{0.294}{-1.451}{0.292}{0.397}
\SG{-1.451}{0.292}{-1.452}{0.289}{0.394}
\SG{-1.452}{0.289}{-1.453}{0.286}{0.394}
\SG{-1.453}{0.286}{-1.454}{0.283}{0.395}
\SG{-1.454}{0.283}{-1.455}{0.280}{0.396}
\SG{-1.455}{0.280}{-1.456}{0.278}{0.397}
\SG{-1.456}{0.278}{-1.457}{0.275}{0.398}
\SG{-1.457}{0.275}{-1.458}{0.272}{0.400}
\SG{-1.458}{0.272}{-1.459}{0.269}{0.403}
\SG{-1.459}{0.269}{-1.460}{0.266}{0.406}
\SG{-1.460}{0.266}{-1.461}{0.264}{0.409}
\SG{-1.461}{0.264}{-1.461}{0.261}{0.411}
\SG{-1.461}{0.261}{-1.462}{0.258}{0.412}
\SG{-1.462}{0.258}{-1.463}{0.255}{0.409}
\SG{-1.463}{0.255}{-1.464}{0.252}{0.407}
\SG{-1.464}{0.252}{-1.465}{0.249}{0.405}
\SG{-1.465}{0.249}{-1.466}{0.247}{0.403}
\SG{-1.466}{0.247}{-1.466}{0.244}{0.402}
\SG{-1.466}{0.244}{-1.467}{0.241}{0.404}
\SG{-1.467}{0.241}{-1.468}{0.238}{0.406}
\SG{-1.468}{0.238}{-1.469}{0.235}{0.408}
\SG{-1.469}{0.235}{-1.469}{0.232}{0.410}
\SG{-1.469}{0.232}{-1.470}{0.229}{0.412}
\SG{-1.470}{0.229}{-1.471}{0.227}{0.418}
\SG{-1.471}{0.227}{-1.472}{0.224}{0.422}
\SG{-1.472}{0.224}{-1.472}{0.221}{0.426}
\SG{-1.472}{0.221}{-1.473}{0.218}{0.429}
\SG{-1.473}{0.218}{-1.474}{0.215}{0.432}
\SG{-1.474}{0.215}{-1.475}{0.212}{0.428}
\SG{-1.475}{0.212}{-1.475}{0.209}{0.423}
\SG{-1.475}{0.209}{-1.476}{0.207}{0.419}
\SG{-1.476}{0.207}{-1.477}{0.204}{0.414}
\SG{-1.477}{0.204}{-1.477}{0.201}{0.411}
\SG{-1.477}{0.201}{-1.478}{0.198}{0.406}
\SG{-1.478}{0.198}{-1.479}{0.195}{0.402}
\SG{-1.479}{0.195}{-1.479}{0.192}{0.399}
\SG{-1.479}{0.192}{-1.480}{0.189}{0.395}
\SG{-1.480}{0.189}{-1.480}{0.186}{0.392}
\SG{-1.480}{0.186}{-1.481}{0.183}{0.391}
\SG{-1.481}{0.183}{-1.482}{0.181}{0.390}
\SG{-1.482}{0.181}{-1.482}{0.178}{0.389}
\SG{-1.482}{0.178}{-1.483}{0.175}{0.389}
\SG{-1.483}{0.175}{-1.483}{0.172}{0.389}
\SG{-1.483}{0.172}{-1.484}{0.169}{0.392}
\SG{-1.484}{0.169}{-1.485}{0.166}{0.396}
\SG{-1.485}{0.166}{-1.485}{0.163}{0.400}
\SG{-1.485}{0.163}{-1.486}{0.160}{0.403}
\SG{-1.486}{0.160}{-1.486}{0.157}{0.407}
\SG{-1.486}{0.157}{-1.487}{0.154}{0.413}
\SG{-1.487}{0.154}{-1.487}{0.151}{0.419}
\SG{-1.487}{0.151}{-1.488}{0.149}{0.425}
\SG{-1.488}{0.149}{-1.488}{0.146}{0.430}
\SG{-1.488}{0.146}{-1.489}{0.143}{0.436}
\SG{-1.489}{0.143}{-1.489}{0.140}{0.434}
\SG{-1.489}{0.140}{-1.490}{0.137}{0.432}
\SG{-1.490}{0.137}{-1.490}{0.134}{0.430}
\SG{-1.490}{0.134}{-1.490}{0.131}{0.428}
\SG{-1.490}{0.131}{-1.491}{0.128}{0.426}
\SG{-1.491}{0.128}{-1.491}{0.125}{0.423}
\SG{-1.491}{0.125}{-1.492}{0.122}{0.420}
\SG{-1.492}{0.122}{-1.492}{0.119}{0.417}
\SG{-1.492}{0.119}{-1.492}{0.116}{0.414}
\SG{-1.492}{0.116}{-1.493}{0.113}{0.412}
\SG{-1.493}{0.113}{-1.493}{0.111}{0.409}
\SG{-1.493}{0.111}{-1.494}{0.108}{0.407}
\SG{-1.494}{0.108}{-1.494}{0.105}{0.404}
\SG{-1.494}{0.105}{-1.494}{0.102}{0.402}
\SG{-1.494}{0.102}{-1.495}{0.099}{0.399}
\SG{-1.495}{0.099}{-1.495}{0.096}{0.397}
\SG{-1.495}{0.096}{-1.495}{0.093}{0.395}
\SG{-1.495}{0.093}{-1.495}{0.090}{0.393}
\SG{-1.495}{0.090}{-1.496}{0.087}{0.391}
\SG{-1.496}{0.087}{-1.496}{0.084}{0.389}
\SG{-1.496}{0.084}{-1.496}{0.081}{0.388}
\SG{-1.496}{0.081}{-1.497}{0.078}{0.386}
\SG{-1.497}{0.078}{-1.497}{0.075}{0.384}
\SG{-1.497}{0.075}{-1.497}{0.072}{0.383}
\SG{-1.497}{0.072}{-1.497}{0.069}{0.381}
\SG{-1.497}{0.069}{-1.498}{0.066}{0.380}
\SG{-1.498}{0.066}{-1.498}{0.063}{0.379}
\SG{-1.498}{0.063}{-1.498}{0.061}{0.378}
\SG{-1.498}{0.061}{-1.498}{0.058}{0.377}
\SG{-1.498}{0.058}{-1.498}{0.055}{0.375}
\SG{-1.498}{0.055}{-1.498}{0.052}{0.375}
\SG{-1.498}{0.052}{-1.499}{0.049}{0.374}
\SG{-1.499}{0.049}{-1.499}{0.046}{0.373}
\SG{-1.499}{0.046}{-1.499}{0.043}{0.372}
\SG{-1.499}{0.043}{-1.499}{0.040}{0.371}
\SG{-1.499}{0.040}{-1.499}{0.037}{0.370}
\SG{-1.499}{0.037}{-1.499}{0.034}{0.370}
\SG{-1.499}{0.034}{-1.499}{0.031}{0.369}
\SG{-1.499}{0.031}{-1.500}{0.028}{0.369}
\SG{-1.500}{0.028}{-1.500}{0.025}{0.368}
\SG{-1.500}{0.025}{-1.500}{0.022}{0.368}
\SG{-1.500}{0.022}{-1.500}{0.019}{0.367}
\SG{-1.500}{0.019}{-1.500}{0.016}{0.367}
\SG{-1.500}{0.016}{-1.500}{0.013}{0.367}
\SG{-1.500}{0.013}{-1.500}{0.010}{0.366}
\SG{-1.500}{0.010}{-1.500}{0.007}{0.366}
\SG{-1.500}{0.007}{-1.500}{0.004}{0.366}
\SG{-1.500}{0.004}{-1.500}{0.001}{0.366}
\SG{-1.500}{0.001}{-1.500}{-0.001}{0.366}
\SG{-1.500}{-0.001}{-1.500}{-0.004}{0.366}
\SG{-1.500}{-0.004}{-1.500}{-0.007}{0.366}
\SG{-1.500}{-0.007}{-1.500}{-0.010}{0.366}
\SG{-1.500}{-0.010}{-1.500}{-0.013}{0.366}
\SG{-1.500}{-0.013}{-1.500}{-0.016}{0.366}
\SG{-1.500}{-0.016}{-1.500}{-0.019}{0.367}
\SG{-1.500}{-0.019}{-1.500}{-0.022}{0.367}
\SG{-1.500}{-0.022}{-1.500}{-0.025}{0.367}
\SG{-1.500}{-0.025}{-1.500}{-0.028}{0.368}
\SG{-1.500}{-0.028}{-1.499}{-0.031}{0.368}
\SG{-1.499}{-0.031}{-1.499}{-0.034}{0.369}
\SG{-1.499}{-0.034}{-1.499}{-0.037}{0.369}
\SG{-1.499}{-0.037}{-1.499}{-0.040}{0.370}
\SG{-1.499}{-0.040}{-1.499}{-0.043}{0.370}
\SG{-1.499}{-0.043}{-1.499}{-0.046}{0.371}
\SG{-1.499}{-0.046}{-1.499}{-0.049}{0.372}
\SG{-1.499}{-0.049}{-1.498}{-0.052}{0.373}
\SG{-1.498}{-0.052}{-1.498}{-0.055}{0.374}
\SG{-1.498}{-0.055}{-1.498}{-0.058}{0.374}
\SG{-1.498}{-0.058}{-1.498}{-0.061}{0.375}
\SG{-1.498}{-0.061}{-1.498}{-0.063}{0.376}
\SG{-1.498}{-0.063}{-1.498}{-0.066}{0.378}
\SG{-1.498}{-0.066}{-1.497}{-0.069}{0.379}
\SG{-1.497}{-0.069}{-1.497}{-0.072}{0.380}
\SG{-1.497}{-0.072}{-1.497}{-0.075}{0.381}
\SG{-1.497}{-0.075}{-1.497}{-0.078}{0.383}
\SG{-1.497}{-0.078}{-1.496}{-0.081}{0.384}
\SG{-1.496}{-0.081}{-1.496}{-0.084}{0.386}
\SG{-1.496}{-0.084}{-1.496}{-0.087}{0.388}
\SG{-1.496}{-0.087}{-1.495}{-0.090}{0.389}
\SG{-1.495}{-0.090}{-1.495}{-0.093}{0.391}
\SG{-1.495}{-0.093}{-1.495}{-0.096}{0.393}
\SG{-1.495}{-0.096}{-1.495}{-0.099}{0.395}
\SG{-1.495}{-0.099}{-1.494}{-0.102}{0.397}
\SG{-1.494}{-0.102}{-1.494}{-0.105}{0.399}
\SG{-1.494}{-0.105}{-1.494}{-0.108}{0.402}
\SG{-1.494}{-0.108}{-1.493}{-0.111}{0.404}
\SG{-1.493}{-0.111}{-1.493}{-0.113}{0.407}
\SG{-1.493}{-0.113}{-1.492}{-0.116}{0.409}
\SG{-1.492}{-0.116}{-1.492}{-0.119}{0.412}
\SG{-1.492}{-0.119}{-1.492}{-0.122}{0.414}
\SG{-1.492}{-0.122}{-1.491}{-0.125}{0.417}
\SG{-1.491}{-0.125}{-1.491}{-0.128}{0.420}
\SG{-1.491}{-0.128}{-1.490}{-0.131}{0.423}
\SG{-1.490}{-0.131}{-1.490}{-0.134}{0.426}
\SG{-1.490}{-0.134}{-1.490}{-0.137}{0.428}
\SG{-1.490}{-0.137}{-1.489}{-0.140}{0.430}
\SG{-1.489}{-0.140}{-1.489}{-0.143}{0.432}
\SG{-1.489}{-0.143}{-1.488}{-0.146}{0.435}
\SG{-1.488}{-0.146}{-1.488}{-0.149}{0.436}
\SG{-1.488}{-0.149}{-1.487}{-0.151}{0.430}
\SG{-1.487}{-0.151}{-1.487}{-0.154}{0.425}
\SG{-1.487}{-0.154}{-1.486}{-0.157}{0.419}
\SG{-1.486}{-0.157}{-1.486}{-0.160}{0.412}
\SG{-1.486}{-0.160}{-1.485}{-0.163}{0.407}
\SG{-1.485}{-0.163}{-1.485}{-0.166}{0.403}
\SG{-1.485}{-0.166}{-1.484}{-0.169}{0.399}
\SG{-1.484}{-0.169}{-1.483}{-0.172}{0.396}
\SG{-1.483}{-0.172}{-1.483}{-0.175}{0.392}
\SG{-1.483}{-0.175}{-1.482}{-0.178}{0.389}
\SG{-1.482}{-0.178}{-1.482}{-0.181}{0.389}
\SG{-1.482}{-0.181}{-1.481}{-0.183}{0.389}
\SG{-1.481}{-0.183}{-1.480}{-0.186}{0.390}
\SG{-1.480}{-0.186}{-1.480}{-0.189}{0.391}
\SG{-1.480}{-0.189}{-1.479}{-0.192}{0.392}
\SG{-1.479}{-0.192}{-1.479}{-0.195}{0.395}
\SG{-1.479}{-0.195}{-1.478}{-0.198}{0.399}
\SG{-1.478}{-0.198}{-1.477}{-0.201}{0.402}
\SG{-1.477}{-0.201}{-1.477}{-0.204}{0.406}
\SG{-1.477}{-0.204}{-1.476}{-0.207}{0.411}
\SG{-1.476}{-0.207}{-1.475}{-0.209}{0.414}
\SG{-1.475}{-0.209}{-1.475}{-0.212}{0.419}
\SG{-1.475}{-0.212}{-1.474}{-0.215}{0.423}
\SG{-1.474}{-0.215}{-1.473}{-0.218}{0.428}
\SG{-1.473}{-0.218}{-1.472}{-0.221}{0.432}
\SG{-1.472}{-0.221}{-1.472}{-0.224}{0.429}
\SG{-1.472}{-0.224}{-1.471}{-0.227}{0.426}
\SG{-1.471}{-0.227}{-1.470}{-0.229}{0.422}
\SG{-1.470}{-0.229}{-1.469}{-0.232}{0.418}
\SG{-1.469}{-0.232}{-1.469}{-0.235}{0.412}
\SG{-1.469}{-0.235}{-1.468}{-0.238}{0.410}
\SG{-1.468}{-0.238}{-1.467}{-0.241}{0.408}
\SG{-1.467}{-0.241}{-1.466}{-0.244}{0.406}
\SG{-1.466}{-0.244}{-1.466}{-0.247}{0.404}
\SG{-1.466}{-0.247}{-1.465}{-0.249}{0.402}
\SG{-1.465}{-0.249}{-1.464}{-0.252}{0.403}
\SG{-1.464}{-0.252}{-1.463}{-0.255}{0.405}
\SG{-1.463}{-0.255}{-1.462}{-0.258}{0.407}
\SG{-1.462}{-0.258}{-1.461}{-0.261}{0.410}
\SG{-1.461}{-0.261}{-1.461}{-0.264}{0.412}
\SG{-1.461}{-0.264}{-1.460}{-0.266}{0.411}
\SG{-1.460}{-0.266}{-1.459}{-0.269}{0.409}
\SG{-1.459}{-0.269}{-1.458}{-0.272}{0.406}
\SG{-1.458}{-0.272}{-1.457}{-0.275}{0.403}
\SG{-1.457}{-0.275}{-1.456}{-0.278}{0.400}
\SG{-1.456}{-0.278}{-1.455}{-0.280}{0.398}
\SG{-1.455}{-0.280}{-1.454}{-0.283}{0.397}
\SG{-1.454}{-0.283}{-1.453}{-0.286}{0.396}
\SG{-1.453}{-0.286}{-1.452}{-0.289}{0.395}
\SG{-1.452}{-0.289}{-1.451}{-0.292}{0.394}
\SG{-1.451}{-0.292}{-1.450}{-0.294}{0.394}
\SG{-1.450}{-0.294}{-1.449}{-0.297}{0.397}
\SG{-1.449}{-0.297}{-1.448}{-0.300}{0.399}
\SG{-1.448}{-0.300}{-1.448}{-0.303}{0.403}
\SG{-1.448}{-0.303}{-1.447}{-0.306}{0.407}
\SG{-1.447}{-0.306}{-1.445}{-0.308}{0.410}
\SG{-1.445}{-0.308}{-1.444}{-0.311}{0.409}
\SG{-1.444}{-0.311}{-1.443}{-0.314}{0.407}
\SG{-1.443}{-0.314}{-1.442}{-0.317}{0.405}
\SG{-1.442}{-0.317}{-1.441}{-0.319}{0.403}
\SG{-1.441}{-0.319}{-1.440}{-0.322}{0.401}
\SG{-1.440}{-0.322}{-1.439}{-0.325}{0.400}
\SG{-1.439}{-0.325}{-1.438}{-0.328}{0.399}
\SG{-1.438}{-0.328}{-1.437}{-0.330}{0.398}
\SG{-1.437}{-0.330}{-1.436}{-0.333}{0.397}
\SG{-1.436}{-0.333}{-1.435}{-0.336}{0.397}
\SG{-1.435}{-0.336}{-1.434}{-0.339}{0.396}
\SG{-1.434}{-0.339}{-1.433}{-0.341}{0.394}
\SG{-1.433}{-0.341}{-1.432}{-0.344}{0.392}
\SG{-1.432}{-0.344}{-1.430}{-0.347}{0.389}
\SG{-1.430}{-0.347}{-1.429}{-0.350}{0.386}
\SG{-1.429}{-0.350}{-1.428}{-0.352}{0.385}
\SG{-1.428}{-0.352}{-1.427}{-0.355}{0.385}
\SG{-1.427}{-0.355}{-1.426}{-0.358}{0.385}
\SG{-1.426}{-0.358}{-1.425}{-0.361}{0.385}
\SG{-1.425}{-0.361}{-1.423}{-0.363}{0.384}
\SG{-1.423}{-0.363}{-1.422}{-0.366}{0.385}
\SG{-1.422}{-0.366}{-1.421}{-0.369}{0.388}
\SG{-1.421}{-0.369}{-1.420}{-0.371}{0.392}
\SG{-1.420}{-0.371}{-1.419}{-0.374}{0.397}
\SG{-1.419}{-0.374}{-1.417}{-0.377}{0.403}
\SG{-1.417}{-0.377}{-1.416}{-0.379}{0.410}
\SG{-1.416}{-0.379}{-1.415}{-0.382}{0.412}
\SG{-1.415}{-0.382}{-1.414}{-0.385}{0.413}
\SG{-1.414}{-0.385}{-1.412}{-0.387}{0.414}
\SG{-1.412}{-0.387}{-1.411}{-0.390}{0.415}
\SG{-1.411}{-0.390}{-1.410}{-0.393}{0.415}
\SG{-1.410}{-0.393}{-1.408}{-0.395}{0.414}
\SG{-1.408}{-0.395}{-1.407}{-0.398}{0.416}
\SG{-1.407}{-0.398}{-1.406}{-0.401}{0.418}
\SG{-1.406}{-0.401}{-1.405}{-0.403}{0.420}
\SG{-1.405}{-0.403}{-1.403}{-0.406}{0.423}
\SG{-1.403}{-0.406}{-1.402}{-0.409}{0.426}
\SG{-1.402}{-0.409}{-1.401}{-0.411}{0.425}
\SG{-1.401}{-0.411}{-1.399}{-0.414}{0.424}
\SG{-1.399}{-0.414}{-1.398}{-0.416}{0.423}
\SG{-1.398}{-0.416}{-1.396}{-0.419}{0.422}
\SG{-1.396}{-0.419}{-1.395}{-0.422}{0.420}
\SG{-1.395}{-0.422}{-1.394}{-0.424}{0.417}
\SG{-1.394}{-0.424}{-1.392}{-0.427}{0.414}
\SG{-1.392}{-0.427}{-1.391}{-0.430}{0.411}
\SG{-1.391}{-0.430}{-1.389}{-0.432}{0.407}
\SG{-1.389}{-0.432}{-1.388}{-0.435}{0.403}
\SG{-1.388}{-0.435}{-1.387}{-0.437}{0.398}
\SG{-1.387}{-0.437}{-1.385}{-0.440}{0.395}
\SG{-1.385}{-0.440}{-1.384}{-0.442}{0.394}
\SG{-1.384}{-0.442}{-1.382}{-0.445}{0.393}
\SG{-1.382}{-0.445}{-1.381}{-0.448}{0.391}
\SG{-1.381}{-0.448}{-1.379}{-0.450}{0.389}
\SG{-1.379}{-0.450}{-1.378}{-0.453}{0.389}
\SG{-1.378}{-0.453}{-1.376}{-0.455}{0.389}
\SG{-1.376}{-0.455}{-1.375}{-0.458}{0.388}
\SG{-1.375}{-0.458}{-1.373}{-0.460}{0.388}
\SG{-1.373}{-0.460}{-1.372}{-0.463}{0.388}
\SG{-1.372}{-0.463}{-1.370}{-0.465}{0.388}
\SG{-1.370}{-0.465}{-1.369}{-0.468}{0.388}
\SG{-1.369}{-0.468}{-1.367}{-0.470}{0.390}
\SG{-1.367}{-0.470}{-1.366}{-0.473}{0.392}
\SG{-1.366}{-0.473}{-1.364}{-0.475}{0.394}
\SG{-1.364}{-0.475}{-1.363}{-0.478}{0.398}
\SG{-1.363}{-0.478}{-1.361}{-0.480}{0.402}
\SG{-1.361}{-0.480}{-1.359}{-0.483}{0.405}
\SG{-1.359}{-0.483}{-1.358}{-0.485}{0.406}
\SG{-1.358}{-0.485}{-1.356}{-0.488}{0.407}
\SG{-1.356}{-0.488}{-1.355}{-0.490}{0.408}
\SG{-1.355}{-0.490}{-1.353}{-0.493}{0.410}
\SG{-1.353}{-0.493}{-1.351}{-0.495}{0.412}
\SG{-1.351}{-0.495}{-1.350}{-0.498}{0.412}
\SG{-1.350}{-0.498}{-1.348}{-0.500}{0.411}
\SG{-1.348}{-0.500}{-1.346}{-0.503}{0.410}
\SG{-1.346}{-0.503}{-1.345}{-0.505}{0.408}
\SG{-1.345}{-0.505}{-1.343}{-0.508}{0.406}
\SG{-1.343}{-0.508}{-1.341}{-0.510}{0.403}
\SG{-1.341}{-0.510}{-1.340}{-0.513}{0.401}
\SG{-1.340}{-0.513}{-1.338}{-0.515}{0.399}
\SG{-1.338}{-0.515}{-1.336}{-0.517}{0.396}
\SG{-1.336}{-0.517}{-1.335}{-0.520}{0.392}
\SG{-1.335}{-0.520}{-1.333}{-0.522}{0.389}
\SG{-1.333}{-0.522}{-1.331}{-0.525}{0.385}
\SG{-1.331}{-0.525}{-1.330}{-0.527}{0.384}
\SG{-1.330}{-0.527}{-1.328}{-0.529}{0.383}
\SG{-1.328}{-0.529}{-1.326}{-0.532}{0.383}
\SG{-1.326}{-0.532}{-1.324}{-0.534}{0.383}
\SG{-1.324}{-0.534}{-1.323}{-0.537}{0.384}
\SG{-1.323}{-0.537}{-1.321}{-0.539}{0.385}
\SG{-1.321}{-0.539}{-1.319}{-0.541}{0.387}
\SG{-1.319}{-0.541}{-1.317}{-0.544}{0.388}
\SG{-1.317}{-0.544}{-1.315}{-0.546}{0.389}
\SG{-1.315}{-0.546}{-1.314}{-0.548}{0.391}
\SG{-1.314}{-0.548}{-1.312}{-0.551}{0.393}
\SG{-1.312}{-0.551}{-1.310}{-0.553}{0.396}
\SG{-1.310}{-0.553}{-1.308}{-0.555}{0.398}
\SG{-1.308}{-0.555}{-1.306}{-0.558}{0.400}
\SG{-1.306}{-0.558}{-1.305}{-0.560}{0.401}
\SG{-1.305}{-0.560}{-1.303}{-0.562}{0.402}
\SG{-1.303}{-0.562}{-1.301}{-0.565}{0.402}
\SG{-1.301}{-0.565}{-1.299}{-0.567}{0.402}
\SG{-1.299}{-0.567}{-1.297}{-0.569}{0.402}
\SG{-1.297}{-0.569}{-1.295}{-0.571}{0.401}
\SG{-1.295}{-0.571}{-1.293}{-0.574}{0.400}
\SG{-1.293}{-0.574}{-1.291}{-0.576}{0.399}
\SG{-1.291}{-0.576}{-1.290}{-0.578}{0.398}
\SG{-1.290}{-0.578}{-1.288}{-0.581}{0.396}
\SG{-1.288}{-0.581}{-1.286}{-0.583}{0.394}
\SG{-1.286}{-0.583}{-1.284}{-0.585}{0.393}
\SG{-1.284}{-0.585}{-1.282}{-0.587}{0.392}
\SG{-1.282}{-0.587}{-1.280}{-0.590}{0.390}
\SG{-1.280}{-0.590}{-1.278}{-0.592}{0.388}
\SG{-1.278}{-0.592}{-1.276}{-0.594}{0.386}
\SG{-1.276}{-0.594}{-1.274}{-0.596}{0.384}
\SG{-1.274}{-0.596}{-1.272}{-0.598}{0.382}
\SG{-1.272}{-0.598}{-1.270}{-0.601}{0.381}
\SG{-1.270}{-0.601}{-1.268}{-0.603}{0.380}
\SG{-1.268}{-0.603}{-1.266}{-0.605}{0.379}
\SG{-1.266}{-0.605}{-1.264}{-0.607}{0.377}
\SG{-1.264}{-0.607}{-1.262}{-0.609}{0.376}
\SG{-1.262}{-0.609}{-1.260}{-0.612}{0.375}
\SG{-1.260}{-0.612}{-1.258}{-0.614}{0.374}
\SG{-1.258}{-0.614}{-1.256}{-0.616}{0.373}
\SG{-1.256}{-0.616}{-1.254}{-0.618}{0.372}
\SG{-1.254}{-0.618}{-1.252}{-0.620}{0.371}
\SG{-1.252}{-0.620}{-1.250}{-0.622}{0.370}
\SG{-1.250}{-0.622}{-1.248}{-0.624}{0.370}
\SG{-1.248}{-0.624}{-1.246}{-0.627}{0.369}
\SG{-1.246}{-0.627}{-1.244}{-0.629}{0.368}
\SG{-1.244}{-0.629}{-1.242}{-0.631}{0.368}
\SG{-1.242}{-0.631}{-1.240}{-0.633}{0.368}
\SG{-1.240}{-0.633}{-1.238}{-0.635}{0.367}
\SG{-1.238}{-0.635}{-1.236}{-0.637}{0.367}
\SG{-1.236}{-0.637}{-1.234}{-0.639}{0.367}
\SG{-1.234}{-0.639}{-1.231}{-0.641}{0.367}
\SG{-1.231}{-0.641}{-1.229}{-0.643}{0.367}
\SG{-1.229}{-0.643}{-1.227}{-0.645}{0.367}
\SG{-1.227}{-0.645}{-1.225}{-0.647}{0.367}
\SG{-1.225}{-0.647}{-1.223}{-0.649}{0.367}
\SG{-1.223}{-0.649}{-1.221}{-0.652}{0.367}
\SG{-1.221}{-0.652}{-1.219}{-0.654}{0.368}
\SG{-1.219}{-0.654}{-1.217}{-0.656}{0.368}
\SG{-1.217}{-0.656}{-1.214}{-0.658}{0.369}
\SG{-1.214}{-0.658}{-1.212}{-0.660}{0.369}
\SG{-1.212}{-0.660}{-1.210}{-0.662}{0.370}
\SG{-1.210}{-0.662}{-1.208}{-0.664}{0.370}
\SG{-1.208}{-0.664}{-1.206}{-0.666}{0.371}
\SG{-1.206}{-0.666}{-1.204}{-0.668}{0.372}
\SG{-1.204}{-0.668}{-1.201}{-0.670}{0.373}
\SG{-1.201}{-0.670}{-1.199}{-0.672}{0.374}
\SG{-1.199}{-0.672}{-1.197}{-0.674}{0.375}
\SG{-1.197}{-0.674}{-1.195}{-0.675}{0.376}
\SG{-1.195}{-0.675}{-1.192}{-0.677}{0.377}
\SG{-1.192}{-0.677}{-1.190}{-0.679}{0.379}
\SG{-1.190}{-0.679}{-1.188}{-0.681}{0.380}
\SG{-1.188}{-0.681}{-1.186}{-0.683}{0.382}
\SG{-1.186}{-0.683}{-1.184}{-0.685}{0.384}
\SG{-1.184}{-0.685}{-1.181}{-0.687}{0.385}
\SG{-1.181}{-0.687}{-1.179}{-0.689}{0.385}
\SG{-1.179}{-0.689}{-1.177}{-0.691}{0.387}
\SG{-1.177}{-0.691}{-1.174}{-0.693}{0.388}
\SG{-1.174}{-0.693}{-1.172}{-0.695}{0.390}
\SG{-1.172}{-0.695}{-1.170}{-0.697}{0.391}
\SG{-1.170}{-0.697}{-1.168}{-0.698}{0.391}
\SG{-1.168}{-0.698}{-1.165}{-0.700}{0.391}
\SG{-1.165}{-0.700}{-1.163}{-0.702}{0.392}
\SG{-1.163}{-0.702}{-1.161}{-0.704}{0.391}
\SG{-1.161}{-0.704}{-1.158}{-0.706}{0.391}
\SG{-1.158}{-0.706}{-1.156}{-0.708}{0.390}
\SG{-1.156}{-0.708}{-1.154}{-0.709}{0.388}
\SG{-1.154}{-0.709}{-1.151}{-0.711}{0.386}
\SG{-1.151}{-0.711}{-1.149}{-0.713}{0.384}
\SG{-1.149}{-0.713}{-1.147}{-0.715}{0.381}
\SG{-1.147}{-0.715}{-1.144}{-0.717}{0.379}
\SG{-1.144}{-0.717}{-1.142}{-0.718}{0.377}
\SG{-1.142}{-0.718}{-1.140}{-0.720}{0.375}
\SG{-1.140}{-0.720}{-1.137}{-0.722}{0.373}
\SG{-1.137}{-0.722}{-1.135}{-0.724}{0.372}
\SG{-1.135}{-0.724}{-1.133}{-0.725}{0.369}
\SG{-1.133}{-0.725}{-1.130}{-0.727}{0.368}
\SG{-1.130}{-0.727}{-1.128}{-0.729}{0.368}
\SG{-1.128}{-0.729}{-1.125}{-0.731}{0.367}
\SG{-1.125}{-0.731}{-1.123}{-0.732}{0.367}
\SG{-1.123}{-0.732}{-1.121}{-0.734}{0.368}
\SG{-1.121}{-0.734}{-1.118}{-0.736}{0.371}
\SG{-1.118}{-0.736}{-1.116}{-0.738}{0.374}
\SG{-1.116}{-0.738}{-1.113}{-0.739}{0.377}
\SG{-1.113}{-0.739}{-1.111}{-0.741}{0.380}
\SG{-1.111}{-0.741}{-1.108}{-0.743}{0.382}
\SG{-1.108}{-0.743}{-1.106}{-0.744}{0.384}
\SG{-1.106}{-0.744}{-1.104}{-0.746}{0.386}
\SG{-1.104}{-0.746}{-1.101}{-0.748}{0.388}
\SG{-1.101}{-0.748}{-1.099}{-0.749}{0.390}
\SG{-1.099}{-0.749}{-1.096}{-0.751}{0.391}
\SG{-1.096}{-0.751}{-1.094}{-0.752}{0.391}
\SG{-1.094}{-0.752}{-1.091}{-0.754}{0.391}
\SG{-1.091}{-0.754}{-1.089}{-0.756}{0.390}
\SG{-1.089}{-0.756}{-1.086}{-0.757}{0.388}
\SG{-1.086}{-0.757}{-1.084}{-0.759}{0.386}
\SG{-1.084}{-0.759}{-1.081}{-0.760}{0.385}
\SG{-1.081}{-0.760}{-1.079}{-0.762}{0.383}
\SG{-1.079}{-0.762}{-1.076}{-0.764}{0.381}
\SG{-1.076}{-0.764}{-1.074}{-0.765}{0.376}
\SG{-1.074}{-0.765}{-1.071}{-0.767}{0.372}
\SG{-1.071}{-0.767}{-1.069}{-0.768}{0.369}
\SG{-1.069}{-0.768}{-1.066}{-0.770}{0.367}
\SG{-1.066}{-0.770}{-1.064}{-0.771}{0.365}
\SG{-1.064}{-0.771}{-1.061}{-0.773}{0.363}
\SG{-1.061}{-0.773}{-1.059}{-0.774}{0.363}
\SG{-1.059}{-0.774}{-1.056}{-0.776}{0.363}
\SG{-1.056}{-0.776}{-1.054}{-0.777}{0.363}
\SG{-1.054}{-0.777}{-1.051}{-0.779}{0.362}
\SG{-1.051}{-0.779}{-1.048}{-0.780}{0.363}
\SG{-1.048}{-0.780}{-1.046}{-0.782}{0.363}
\SG{-1.046}{-0.782}{-1.043}{-0.783}{0.363}
\SG{-1.043}{-0.783}{-1.041}{-0.785}{0.365}
\SG{-1.041}{-0.785}{-1.038}{-0.786}{0.366}
\SG{-1.038}{-0.786}{-1.036}{-0.788}{0.367}
\SG{-1.036}{-0.788}{-1.033}{-0.789}{0.368}
\SG{-1.033}{-0.789}{-1.030}{-0.790}{0.373}
\SG{-1.030}{-0.790}{-1.028}{-0.792}{0.377}
\SG{-1.028}{-0.792}{-1.025}{-0.793}{0.380}
\SG{-1.025}{-0.793}{-1.023}{-0.795}{0.383}
\SG{-1.023}{-0.795}{-1.020}{-0.796}{0.386}
\SG{-1.020}{-0.796}{-1.017}{-0.797}{0.389}
\SG{-1.017}{-0.797}{-1.015}{-0.799}{0.390}
\SG{-1.015}{-0.799}{-1.012}{-0.800}{0.391}
\SG{-1.012}{-0.800}{-1.009}{-0.801}{0.392}
\SG{-1.009}{-0.801}{-1.007}{-0.803}{0.393}
\SG{-1.007}{-0.803}{-1.004}{-0.804}{0.394}
\SG{-1.004}{-0.804}{-1.002}{-0.805}{0.391}
\SG{-1.002}{-0.805}{-0.999}{-0.807}{0.388}
\SG{-0.999}{-0.807}{-0.996}{-0.808}{0.385}
\SG{-0.996}{-0.808}{-0.994}{-0.809}{0.383}
\SG{-0.994}{-0.809}{-0.991}{-0.811}{0.381}
\SG{-0.991}{-0.811}{-0.988}{-0.812}{0.380}
\SG{-0.988}{-0.812}{-0.986}{-0.813}{0.380}
\SG{-0.986}{-0.813}{-0.983}{-0.814}{0.379}
\SG{-0.983}{-0.814}{-0.980}{-0.816}{0.378}
\SG{-0.980}{-0.816}{-0.978}{-0.817}{0.377}
\SG{-0.978}{-0.817}{-0.975}{-0.818}{0.375}
\SG{-0.975}{-0.818}{-0.972}{-0.819}{0.368}
\SG{-0.972}{-0.819}{-0.969}{-0.821}{0.362}
\SG{-0.969}{-0.821}{-0.967}{-0.822}{0.356}
\SG{-0.967}{-0.822}{-0.964}{-0.823}{0.352}
\SG{-0.964}{-0.823}{-0.961}{-0.824}{0.348}
\SG{-0.961}{-0.824}{-0.959}{-0.825}{0.346}
\SG{-0.959}{-0.825}{-0.956}{-0.827}{0.346}
\SG{-0.956}{-0.827}{-0.953}{-0.828}{0.346}
\SG{-0.953}{-0.828}{-0.951}{-0.829}{0.346}
\SG{-0.951}{-0.829}{-0.948}{-0.830}{0.345}
\SG{-0.948}{-0.830}{-0.945}{-0.831}{0.345}
\SG{-0.945}{-0.831}{-0.942}{-0.832}{0.348}
\SG{-0.942}{-0.832}{-0.940}{-0.833}{0.350}
\SG{-0.940}{-0.833}{-0.937}{-0.835}{0.352}
\SG{-0.937}{-0.835}{-0.934}{-0.836}{0.353}
\SG{-0.934}{-0.836}{-0.931}{-0.837}{0.354}
\SG{-0.931}{-0.837}{-0.929}{-0.838}{0.355}
\SG{-0.929}{-0.838}{-0.926}{-0.839}{0.355}
\SG{-0.926}{-0.839}{-0.923}{-0.840}{0.355}
\SG{-0.923}{-0.840}{-0.920}{-0.841}{0.356}
\SG{-0.920}{-0.841}{-0.918}{-0.842}{0.356}
\SG{-0.918}{-0.842}{-0.915}{-0.843}{0.358}
\SG{-0.915}{-0.843}{-0.912}{-0.844}{0.360}
\SG{-0.912}{-0.844}{-0.909}{-0.845}{0.361}
\SG{-0.909}{-0.845}{-0.907}{-0.846}{0.363}
\SG{-0.907}{-0.846}{-0.904}{-0.847}{0.363}
\SG{-0.904}{-0.847}{-0.901}{-0.848}{0.362}
\SG{-0.901}{-0.848}{-0.898}{-0.849}{0.357}
\SG{-0.898}{-0.849}{-0.895}{-0.850}{0.353}
\SG{-0.895}{-0.850}{-0.893}{-0.851}{0.350}
\SG{-0.893}{-0.851}{-0.890}{-0.852}{0.347}
\SG{-0.890}{-0.852}{-0.887}{-0.853}{0.345}
\SG{-0.887}{-0.853}{-0.884}{-0.854}{0.345}
\SG{-0.884}{-0.854}{-0.881}{-0.855}{0.346}
\SG{-0.881}{-0.855}{-0.879}{-0.856}{0.347}
\SG{-0.879}{-0.856}{-0.876}{-0.857}{0.348}
\SG{-0.876}{-0.857}{-0.873}{-0.858}{0.349}
\SG{-0.873}{-0.858}{-0.870}{-0.858}{0.352}
\SG{-0.870}{-0.858}{-0.867}{-0.859}{0.354}
\SG{-0.867}{-0.859}{-0.864}{-0.860}{0.356}
\SG{-0.864}{-0.860}{-0.862}{-0.861}{0.358}
\SG{-0.862}{-0.861}{-0.859}{-0.862}{0.360}
\SG{-0.859}{-0.862}{-0.856}{-0.863}{0.357}
\SG{-0.856}{-0.863}{-0.853}{-0.864}{0.354}
\SG{-0.853}{-0.864}{-0.850}{-0.864}{0.352}
\SG{-0.850}{-0.864}{-0.847}{-0.865}{0.350}
\SG{-0.847}{-0.865}{-0.845}{-0.866}{0.348}
\SG{-0.845}{-0.866}{-0.842}{-0.867}{0.349}
\SG{-0.842}{-0.867}{-0.839}{-0.868}{0.350}
\SG{-0.839}{-0.868}{-0.836}{-0.868}{0.352}
\SG{-0.836}{-0.868}{-0.833}{-0.869}{0.353}
\SG{-0.833}{-0.869}{-0.830}{-0.870}{0.355}
\SG{-0.830}{-0.870}{-0.828}{-0.871}{0.359}
\SG{-0.828}{-0.871}{-0.825}{-0.871}{0.364}
\SG{-0.825}{-0.871}{-0.822}{-0.872}{0.367}
\SG{-0.822}{-0.872}{-0.819}{-0.873}{0.370}
\SG{-0.819}{-0.873}{-0.816}{-0.874}{0.373}
\SG{-0.816}{-0.874}{-0.813}{-0.874}{0.370}
\SG{-0.813}{-0.874}{-0.810}{-0.875}{0.365}
\SG{-0.810}{-0.875}{-0.807}{-0.876}{0.360}
\SG{-0.807}{-0.876}{-0.805}{-0.876}{0.355}
\SG{-0.805}{-0.876}{-0.802}{-0.877}{0.351}
\SG{-0.802}{-0.877}{-0.799}{-0.878}{0.346}
\SG{-0.799}{-0.878}{-0.796}{-0.878}{0.342}
\SG{-0.796}{-0.878}{-0.793}{-0.879}{0.338}
\SG{-0.793}{-0.879}{-0.790}{-0.880}{0.334}
\SG{-0.790}{-0.880}{-0.787}{-0.880}{0.330}
\SG{-0.787}{-0.880}{-0.784}{-0.881}{0.327}
\SG{-0.784}{-0.881}{-0.781}{-0.881}{0.326}
\SG{-0.781}{-0.881}{-0.779}{-0.882}{0.325}
\SG{-0.779}{-0.882}{-0.776}{-0.883}{0.324}
\SG{-0.776}{-0.883}{-0.773}{-0.883}{0.323}
\SG{-0.773}{-0.883}{-0.770}{-0.884}{0.325}
\SG{-0.770}{-0.884}{-0.767}{-0.884}{0.328}
\SG{-0.767}{-0.884}{-0.764}{-0.885}{0.331}
\SG{-0.764}{-0.885}{-0.761}{-0.885}{0.335}
\SG{-0.761}{-0.885}{-0.758}{-0.886}{0.338}
\SG{-0.758}{-0.886}{-0.755}{-0.886}{0.342}
\SG{-0.755}{-0.886}{-0.752}{-0.887}{0.348}
\SG{-0.752}{-0.887}{-0.750}{-0.887}{0.354}
\SG{-0.750}{-0.887}{-0.747}{-0.888}{0.359}
\SG{-0.747}{-0.888}{-0.744}{-0.888}{0.364}
\SG{-0.744}{-0.888}{-0.741}{-0.889}{0.365}
\SG{-0.741}{-0.889}{-0.738}{-0.889}{0.362}
\SG{-0.738}{-0.889}{-0.735}{-0.890}{0.359}
\SG{-0.735}{-0.890}{-0.732}{-0.890}{0.357}
\SG{-0.732}{-0.890}{-0.729}{-0.891}{0.354}
\SG{-0.729}{-0.891}{-0.726}{-0.891}{0.351}
\SG{-0.726}{-0.891}{-0.723}{-0.892}{0.348}
\SG{-0.723}{-0.892}{-0.720}{-0.892}{0.344}
\SG{-0.720}{-0.892}{-0.717}{-0.892}{0.341}
\SG{-0.717}{-0.892}{-0.714}{-0.893}{0.338}
\SG{-0.714}{-0.893}{-0.711}{-0.893}{0.334}
\SG{-0.711}{-0.893}{-0.709}{-0.893}{0.331}
\SG{-0.709}{-0.893}{-0.706}{-0.894}{0.328}
\SG{-0.706}{-0.894}{-0.703}{-0.894}{0.325}
\SG{-0.703}{-0.894}{-0.700}{-0.894}{0.322}
\SG{-0.700}{-0.894}{-0.697}{-0.895}{0.320}
\SG{-0.697}{-0.895}{-0.694}{-0.895}{0.317}
\SG{-0.694}{-0.895}{-0.691}{-0.895}{0.315}
\SG{-0.691}{-0.895}{-0.688}{-0.896}{0.312}
\SG{-0.688}{-0.896}{-0.685}{-0.896}{0.310}
\SG{-0.685}{-0.896}{-0.682}{-0.896}{0.307}
\SG{-0.682}{-0.896}{-0.679}{-0.896}{0.305}
\SG{-0.679}{-0.896}{-0.676}{-0.897}{0.303}
\SG{-0.676}{-0.897}{-0.673}{-0.897}{0.301}
\SG{-0.673}{-0.897}{-0.670}{-0.897}{0.299}
\SG{-0.670}{-0.897}{-0.667}{-0.897}{0.297}
\SG{-0.667}{-0.897}{-0.664}{-0.898}{0.295}
\SG{-0.664}{-0.898}{-0.661}{-0.898}{0.294}
\SG{-0.661}{-0.898}{-0.659}{-0.898}{0.292}
\SG{-0.659}{-0.898}{-0.656}{-0.898}{0.290}
\SG{-0.656}{-0.898}{-0.653}{-0.898}{0.289}
\SG{-0.653}{-0.898}{-0.650}{-0.899}{0.287}
\SG{-0.650}{-0.899}{-0.647}{-0.899}{0.286}
\SG{-0.647}{-0.899}{-0.644}{-0.899}{0.284}
\SG{-0.644}{-0.899}{-0.641}{-0.899}{0.283}
\SG{-0.641}{-0.899}{-0.638}{-0.899}{0.282}
\SG{-0.638}{-0.899}{-0.635}{-0.899}{0.281}
\SG{-0.635}{-0.899}{-0.632}{-0.899}{0.280}
\SG{-0.632}{-0.899}{-0.629}{-0.900}{0.278}
\SG{-0.629}{-0.900}{-0.626}{-0.900}{0.277}
\SG{-0.626}{-0.900}{-0.623}{-0.900}{0.276}
\SG{-0.623}{-0.900}{-0.620}{-0.900}{0.275}
\SG{-0.620}{-0.900}{-0.617}{-0.900}{0.274}
\SG{-0.617}{-0.900}{-0.614}{-0.900}{0.273}
\SG{-0.614}{-0.900}{-0.611}{-0.900}{0.273}
\SG{-0.611}{-0.900}{-0.608}{-0.900}{0.272}
\SG{-0.608}{-0.900}{-0.605}{-0.900}{0.271}
\SG{-0.605}{-0.900}{-0.602}{-0.900}{0.270}
\SG{-0.602}{-0.900}{-0.599}{-0.900}{0.270}
\SG{-0.599}{-0.900}{-0.597}{-0.900}{0.269}
\SG{-0.597}{-0.900}{-0.594}{-0.900}{0.268}
\SG{-0.594}{-0.900}{-0.591}{-0.900}{0.268}
\SG{-0.591}{-0.900}{-0.588}{-0.900}{0.267}
\SG{-0.588}{-0.900}{-0.585}{-0.900}{0.267}
\SG{-0.585}{-0.900}{-0.582}{-0.900}{0.267}
\SG{-0.582}{-0.900}{-0.579}{-0.900}{0.266}
\SG{-0.579}{-0.900}{-0.576}{-0.900}{0.266}
\SG{-0.576}{-0.900}{-0.573}{-0.900}{0.266}
\SG{-0.573}{-0.900}{-0.570}{-0.899}{0.265}
\SG{-0.570}{-0.899}{-0.567}{-0.899}{0.265}
\SG{-0.567}{-0.899}{-0.564}{-0.899}{0.265}
\SG{-0.564}{-0.899}{-0.561}{-0.899}{0.265}
\SG{-0.561}{-0.899}{-0.558}{-0.899}{0.265}
\SG{-0.558}{-0.899}{-0.555}{-0.899}{0.265}
\SG{-0.555}{-0.899}{-0.552}{-0.899}{0.265}
\SG{-0.552}{-0.899}{-0.549}{-0.899}{0.265}
\SG{-0.549}{-0.899}{-0.546}{-0.898}{0.265}
\SG{-0.546}{-0.898}{-0.543}{-0.898}{0.266}
\SG{-0.543}{-0.898}{-0.540}{-0.898}{0.266}
\SG{-0.540}{-0.898}{-0.537}{-0.898}{0.266}
\SG{-0.537}{-0.898}{-0.535}{-0.898}{0.267}
\SG{-0.535}{-0.898}{-0.532}{-0.897}{0.267}
\SG{-0.532}{-0.897}{-0.529}{-0.897}{0.268}
\SG{-0.529}{-0.897}{-0.526}{-0.897}{0.268}
\SG{-0.526}{-0.897}{-0.523}{-0.897}{0.269}
\SG{-0.523}{-0.897}{-0.520}{-0.896}{0.270}
\SG{-0.520}{-0.896}{-0.517}{-0.896}{0.270}
\SG{-0.517}{-0.896}{-0.514}{-0.896}{0.271}
\SG{-0.514}{-0.896}{-0.511}{-0.896}{0.272}
\SG{-0.511}{-0.896}{-0.508}{-0.895}{0.273}
\SG{-0.508}{-0.895}{-0.505}{-0.895}{0.274}
\SG{-0.505}{-0.895}{-0.502}{-0.895}{0.276}
\SG{-0.502}{-0.895}{-0.499}{-0.894}{0.277}
\SG{-0.499}{-0.894}{-0.496}{-0.894}{0.278}
\SG{-0.496}{-0.894}{-0.493}{-0.894}{0.280}
\SG{-0.493}{-0.894}{-0.490}{-0.893}{0.281}
\SG{-0.490}{-0.893}{-0.487}{-0.893}{0.283}
\SG{-0.487}{-0.893}{-0.485}{-0.893}{0.285}
\SG{-0.485}{-0.893}{-0.482}{-0.892}{0.287}
\SG{-0.482}{-0.892}{-0.479}{-0.892}{0.288}
\SG{-0.479}{-0.892}{-0.476}{-0.891}{0.290}
\SG{-0.476}{-0.891}{-0.473}{-0.891}{0.292}
\SG{-0.473}{-0.891}{-0.470}{-0.891}{0.295}
\SG{-0.470}{-0.891}{-0.467}{-0.890}{0.297}
\SG{-0.467}{-0.890}{-0.464}{-0.890}{0.298}
\SG{-0.464}{-0.890}{-0.461}{-0.889}{0.299}
\SG{-0.461}{-0.889}{-0.458}{-0.889}{0.300}
\SG{-0.458}{-0.889}{-0.455}{-0.888}{0.302}
\SG{-0.455}{-0.888}{-0.452}{-0.888}{0.303}
\SG{-0.452}{-0.888}{-0.449}{-0.887}{0.298}
\SG{-0.449}{-0.887}{-0.447}{-0.887}{0.292}
\SG{-0.447}{-0.887}{-0.444}{-0.886}{0.285}
\SG{-0.444}{-0.886}{-0.441}{-0.886}{0.278}
\SG{-0.441}{-0.886}{-0.438}{-0.885}{0.271}
\SG{-0.438}{-0.885}{-0.435}{-0.885}{0.266}
\SG{-0.435}{-0.885}{-0.432}{-0.884}{0.262}
\SG{-0.432}{-0.884}{-0.429}{-0.884}{0.257}
\SG{-0.429}{-0.884}{-0.426}{-0.883}{0.252}
\SG{-0.426}{-0.883}{-0.423}{-0.882}{0.247}
\SG{-0.423}{-0.882}{-0.420}{-0.882}{0.246}
\SG{-0.420}{-0.882}{-0.418}{-0.881}{0.245}
\SG{-0.418}{-0.881}{-0.415}{-0.881}{0.245}
\SG{-0.415}{-0.881}{-0.412}{-0.880}{0.244}
\SG{-0.412}{-0.880}{-0.409}{-0.879}{0.244}
\SG{-0.409}{-0.879}{-0.406}{-0.879}{0.246}
\SG{-0.406}{-0.879}{-0.403}{-0.878}{0.248}
\SG{-0.403}{-0.878}{-0.400}{-0.878}{0.251}
\SG{-0.400}{-0.878}{-0.397}{-0.877}{0.253}
\SG{-0.397}{-0.877}{-0.394}{-0.876}{0.256}
\SG{-0.394}{-0.876}{-0.392}{-0.876}{0.259}
\SG{-0.392}{-0.876}{-0.389}{-0.875}{0.262}
\SG{-0.389}{-0.875}{-0.386}{-0.874}{0.266}
\SG{-0.386}{-0.874}{-0.383}{-0.873}{0.269}
\SG{-0.383}{-0.873}{-0.380}{-0.873}{0.273}
\SG{-0.380}{-0.873}{-0.377}{-0.872}{0.271}
\SG{-0.377}{-0.872}{-0.374}{-0.871}{0.266}
\SG{-0.374}{-0.871}{-0.371}{-0.870}{0.262}
\SG{-0.371}{-0.870}{-0.369}{-0.870}{0.256}
\SG{-0.369}{-0.870}{-0.366}{-0.869}{0.250}
\SG{-0.366}{-0.869}{-0.363}{-0.868}{0.246}
\SG{-0.363}{-0.868}{-0.360}{-0.867}{0.242}
\SG{-0.360}{-0.867}{-0.357}{-0.867}{0.239}
\SG{-0.357}{-0.867}{-0.354}{-0.866}{0.236}
\SG{-0.354}{-0.866}{-0.352}{-0.865}{0.232}
\SG{-0.352}{-0.865}{-0.349}{-0.864}{0.231}
\SG{-0.349}{-0.864}{-0.346}{-0.863}{0.231}
\SG{-0.346}{-0.863}{-0.343}{-0.863}{0.231}
\SG{-0.343}{-0.863}{-0.340}{-0.862}{0.232}
\SG{-0.340}{-0.862}{-0.337}{-0.861}{0.234}
\SG{-0.337}{-0.861}{-0.335}{-0.860}{0.233}
\SG{-0.335}{-0.860}{-0.332}{-0.859}{0.229}
\SG{-0.332}{-0.859}{-0.329}{-0.858}{0.225}
\SG{-0.329}{-0.858}{-0.326}{-0.857}{0.221}
\SG{-0.326}{-0.857}{-0.323}{-0.856}{0.216}
\SG{-0.323}{-0.856}{-0.320}{-0.855}{0.212}
\SG{-0.320}{-0.855}{-0.318}{-0.855}{0.209}
\SG{-0.318}{-0.855}{-0.315}{-0.854}{0.206}
\SG{-0.315}{-0.854}{-0.312}{-0.853}{0.203}
\SG{-0.312}{-0.853}{-0.309}{-0.852}{0.201}
\SG{-0.309}{-0.852}{-0.306}{-0.851}{0.199}
\SG{-0.306}{-0.851}{-0.304}{-0.850}{0.199}
\SG{-0.304}{-0.850}{-0.301}{-0.849}{0.200}
\SG{-0.301}{-0.849}{-0.298}{-0.848}{0.201}
\SG{-0.298}{-0.848}{-0.295}{-0.847}{0.203}
\SG{-0.295}{-0.847}{-0.293}{-0.846}{0.206}
\SG{-0.293}{-0.846}{-0.290}{-0.845}{0.203}
\SG{-0.290}{-0.845}{-0.287}{-0.844}{0.200}
\SG{-0.287}{-0.844}{-0.284}{-0.843}{0.197}
\SG{-0.284}{-0.843}{-0.281}{-0.842}{0.194}
\SG{-0.281}{-0.842}{-0.279}{-0.841}{0.190}
\SG{-0.279}{-0.841}{-0.276}{-0.840}{0.187}
\SG{-0.276}{-0.840}{-0.273}{-0.839}{0.185}
\SG{-0.273}{-0.839}{-0.270}{-0.837}{0.183}
\SG{-0.270}{-0.837}{-0.268}{-0.836}{0.181}
\SG{-0.268}{-0.836}{-0.265}{-0.835}{0.179}
\SG{-0.265}{-0.835}{-0.262}{-0.834}{0.179}
\SG{-0.262}{-0.834}{-0.259}{-0.833}{0.181}
\SG{-0.259}{-0.833}{-0.256}{-0.832}{0.183}
\SG{-0.256}{-0.832}{-0.251}{-0.830}{0.187}
\SG{-0.251}{-0.830}{-0.246}{-0.827}{0.193}
\SG{-0.246}{-0.827}{-0.239}{-0.825}{0.189}
\SG{-0.239}{-0.825}{-0.233}{-0.822}{0.183}
\SG{-0.233}{-0.822}{-0.227}{-0.820}{0.175}
\SG{-0.227}{-0.820}{-0.220}{-0.817}{0.171}
\SG{-0.220}{-0.817}{-0.214}{-0.814}{0.167}
\SG{-0.214}{-0.814}{-0.208}{-0.812}{0.158}
\SG{-0.208}{-0.812}{-0.201}{-0.809}{0.147}
\SG{-0.201}{-0.809}{-0.195}{-0.806}{0.138}
\SG{-0.195}{-0.806}{-0.189}{-0.804}{0.132}
\SG{-0.189}{-0.804}{-0.182}{-0.801}{0.127}
\SG{-0.182}{-0.801}{-0.176}{-0.798}{0.119}
\SG{-0.176}{-0.798}{-0.170}{-0.796}{0.110}
\SG{-0.170}{-0.796}{-0.163}{-0.793}{0.106}
\SG{-0.163}{-0.793}{-0.157}{-0.791}{0.102}
\SG{-0.157}{-0.791}{-0.151}{-0.788}{0.104}
\SG{-0.151}{-0.788}{-0.144}{-0.785}{0.112}
\SG{-0.144}{-0.785}{-0.138}{-0.783}{0.118}
\SG{-0.138}{-0.783}{-0.132}{-0.780}{0.114}
\SG{-0.132}{-0.780}{-0.125}{-0.777}{0.108}
\SG{-0.125}{-0.777}{-0.119}{-0.775}{0.106}
\SG{-0.119}{-0.775}{-0.113}{-0.772}{0.104}
\SG{-0.113}{-0.772}{-0.106}{-0.770}{0.099}
\SG{-0.106}{-0.770}{-0.100}{-0.767}{0.092}
\SG{-0.100}{-0.767}{-0.093}{-0.764}{0.085}
\SG{-0.093}{-0.764}{-0.087}{-0.762}{0.081}
\SG{-0.087}{-0.762}{-0.081}{-0.759}{0.079}
\SG{-0.081}{-0.759}{-0.074}{-0.756}{0.085}
\SG{-0.074}{-0.756}{-0.068}{-0.754}{0.095}
\SG{-0.068}{-0.754}{-0.062}{-0.751}{0.099}
\SG{-0.062}{-0.751}{-0.055}{-0.749}{0.097}
\SG{-0.055}{-0.749}{-0.049}{-0.746}{0.091}
\SG{-0.049}{-0.746}{-0.043}{-0.743}{0.089}
\SG{-0.043}{-0.743}{-0.036}{-0.741}{0.089}
\SG{-0.036}{-0.741}{-0.030}{-0.738}{0.084}
\SG{-0.030}{-0.738}{-0.024}{-0.735}{0.076}
\SG{-0.024}{-0.735}{-0.017}{-0.733}{0.068}
\SG{-0.017}{-0.733}{-0.011}{-0.730}{0.064}
\SG{-0.011}{-0.730}{-0.005}{-0.727}{0.061}
\SG{-0.005}{-0.727}{0.002}{-0.725}{0.055}
\SG{0.002}{-0.725}{0.008}{-0.722}{0.048}
\SG{0.008}{-0.722}{0.014}{-0.720}{0.045}
\SG{0.014}{-0.720}{0.021}{-0.717}{0.042}
\SG{0.021}{-0.717}{0.027}{-0.714}{0.043}
\SG{0.027}{-0.714}{0.033}{-0.712}{0.051}
\SG{0.033}{-0.712}{0.040}{-0.709}{0.062}
\SG{0.040}{-0.709}{0.046}{-0.706}{0.060}
\SG{0.046}{-0.706}{0.052}{-0.704}{0.057}
\SG{0.052}{-0.704}{0.059}{-0.701}{0.054}
\SG{0.059}{-0.701}{0.065}{-0.699}{0.053}
\SG{0.065}{-0.699}{0.071}{-0.696}{0.052}
\SG{0.071}{-0.696}{0.078}{-0.693}{0.046}
\SG{0.078}{-0.693}{0.084}{-0.691}{0.039}
\SG{0.084}{-0.691}{0.091}{-0.688}{0.037}
\SG{0.091}{-0.688}{0.097}{-0.685}{0.035}
\SG{0.097}{-0.685}{0.103}{-0.683}{0.038}
\SG{0.103}{-0.683}{0.110}{-0.680}{0.047}
\SG{0.110}{-0.680}{0.116}{-0.677}{0.058}
\SG{0.116}{-0.677}{0.122}{-0.675}{0.056}
\SG{0.122}{-0.675}{0.129}{-0.672}{0.053}
\SG{0.129}{-0.672}{0.135}{-0.670}{0.050}
\SG{0.135}{-0.670}{0.141}{-0.667}{0.050}
\SG{0.141}{-0.667}{0.148}{-0.664}{0.048}
\SG{0.148}{-0.664}{0.154}{-0.662}{0.042}
\SG{0.154}{-0.662}{0.160}{-0.659}{0.034}
\SG{0.160}{-0.659}{0.167}{-0.656}{0.030}
\SG{0.167}{-0.656}{0.173}{-0.654}{0.028}
\SG{0.173}{-0.654}{0.179}{-0.651}{0.024}
\SG{0.179}{-0.651}{0.186}{-0.649}{0.018}
\SG{0.186}{-0.649}{0.192}{-0.646}{0.013}
\SG{0.192}{-0.646}{0.198}{-0.643}{0.012}
\SG{0.198}{-0.643}{0.205}{-0.641}{0.011}
\SG{0.205}{-0.641}{0.211}{-0.638}{0.017}
\SG{0.211}{-0.638}{0.217}{-0.635}{0.028}
\SG{0.217}{-0.635}{0.224}{-0.633}{0.032}
\SG{0.224}{-0.633}{0.230}{-0.630}{0.030}
\SG{0.230}{-0.630}{0.236}{-0.627}{0.027}
\SG{0.236}{-0.627}{0.243}{-0.625}{0.026}
\SG{0.243}{-0.625}{0.249}{-0.622}{0.026}
\SG{0.249}{-0.622}{0.256}{-0.620}{0.022}
\SG{0.256}{-0.620}{0.262}{-0.617}{0.016}
\SG{0.262}{-0.617}{0.268}{-0.614}{0.013}
\SG{0.268}{-0.614}{0.275}{-0.612}{0.011}
\SG{0.275}{-0.612}{0.281}{-0.609}{0.011}
\SG{0.281}{-0.609}{0.287}{-0.606}{0.019}
\SG{0.287}{-0.606}{0.294}{-0.604}{0.031}
\SG{0.294}{-0.604}{0.300}{-0.601}{0.033}
\SG{0.300}{-0.601}{0.306}{-0.599}{0.031}
\SG{0.306}{-0.599}{0.313}{-0.596}{0.027}
\SG{0.313}{-0.596}{0.319}{-0.593}{0.027}
\SG{0.319}{-0.593}{0.325}{-0.591}{0.027}
\SG{0.325}{-0.591}{0.332}{-0.588}{0.022}
\SG{0.332}{-0.588}{0.338}{-0.585}{0.016}
\SG{0.338}{-0.585}{0.344}{-0.583}{0.012}
\SG{0.344}{-0.583}{0.351}{-0.580}{0.011}
\SG{0.351}{-0.580}{0.357}{-0.577}{0.014}
\SG{0.357}{-0.577}{0.363}{-0.575}{0.024}
\SG{0.363}{-0.575}{0.370}{-0.572}{0.038}
\SG{0.370}{-0.572}{0.376}{-0.570}{0.037}
\SG{0.376}{-0.570}{0.382}{-0.567}{0.035}
\SG{0.382}{-0.567}{0.389}{-0.564}{0.032}
\SG{0.389}{-0.564}{0.395}{-0.562}{0.032}
\SG{0.395}{-0.562}{0.401}{-0.559}{0.032}
\SG{0.401}{-0.559}{0.408}{-0.556}{0.027}
\SG{0.408}{-0.556}{0.414}{-0.554}{0.020}
\SG{0.414}{-0.554}{0.421}{-0.551}{0.015}
\SG{0.421}{-0.551}{0.427}{-0.549}{0.014}
\SG{0.427}{-0.549}{0.433}{-0.546}{0.011}
\SG{0.433}{-0.546}{0.440}{-0.543}{0.006}
\SG{0.440}{-0.543}{0.446}{-0.541}{0.001}
\SG{0.446}{-0.541}{0.452}{-0.538}{0.000}
\SG{0.452}{-0.538}{0.459}{-0.535}{0.000}
\SG{0.459}{-0.535}{0.465}{-0.533}{0.006}
\SG{0.465}{-0.533}{0.471}{-0.530}{0.017}
\SG{0.471}{-0.530}{0.478}{-0.528}{0.024}
\SG{0.478}{-0.528}{0.484}{-0.525}{0.023}
\SG{0.484}{-0.525}{0.490}{-0.522}{0.020}
\SG{0.490}{-0.522}{0.497}{-0.520}{0.020}
\SG{0.497}{-0.520}{0.503}{-0.517}{0.021}
\SG{0.503}{-0.517}{0.509}{-0.514}{0.019}
\SG{0.509}{-0.514}{0.516}{-0.512}{0.014}
\SG{0.516}{-0.512}{0.522}{-0.509}{0.010}
\SG{0.522}{-0.509}{0.528}{-0.506}{0.009}
\SG{0.528}{-0.506}{0.535}{-0.504}{0.009}
\SG{0.535}{-0.504}{0.541}{-0.501}{0.018}
\SG{0.541}{-0.501}{0.547}{-0.499}{0.030}
\SG{0.547}{-0.499}{0.554}{-0.496}{0.036}
\SG{0.554}{-0.496}{0.560}{-0.493}{0.036}
\SG{0.560}{-0.493}{0.566}{-0.491}{0.032}
\SG{0.566}{-0.491}{0.573}{-0.488}{0.032}
\SG{0.573}{-0.488}{0.579}{-0.485}{0.034}
\SG{0.579}{-0.485}{0.586}{-0.483}{0.030}
\SG{0.586}{-0.483}{0.592}{-0.480}{0.025}
\SG{0.592}{-0.480}{0.598}{-0.478}{0.019}
\SG{0.598}{-0.478}{0.605}{-0.475}{0.017}
\SG{0.605}{-0.475}{0.611}{-0.472}{0.016}
\SG{0.611}{-0.472}{0.617}{-0.470}{0.012}
\SG{0.617}{-0.470}{0.624}{-0.467}{0.007}
\SG{0.624}{-0.467}{0.630}{-0.464}{0.006}
\SG{0.630}{-0.464}{0.636}{-0.462}{0.006}
\SG{0.636}{-0.462}{0.643}{-0.459}{0.009}
\SG{0.643}{-0.459}{0.649}{-0.456}{0.019}
\SG{0.649}{-0.456}{0.655}{-0.454}{0.031}
\SG{0.655}{-0.454}{0.662}{-0.451}{0.031}
\SG{0.662}{-0.451}{0.668}{-0.449}{0.030}
\SG{0.668}{-0.449}{0.674}{-0.446}{0.028}
\SG{0.674}{-0.446}{0.681}{-0.443}{0.029}
\SG{0.681}{-0.443}{0.687}{-0.441}{0.029}
\SG{0.687}{-0.441}{0.693}{-0.438}{0.026}
\SG{0.693}{-0.438}{0.700}{-0.435}{0.020}
\SG{0.700}{-0.435}{0.706}{-0.433}{0.020}
\SG{0.706}{-0.433}{0.712}{-0.430}{0.020}
\SG{0.712}{-0.430}{0.719}{-0.428}{0.026}
\SG{0.719}{-0.428}{0.725}{-0.425}{0.037}
\SG{0.725}{-0.425}{0.731}{-0.422}{0.048}
\SG{0.731}{-0.422}{0.738}{-0.420}{0.048}
\SG{0.738}{-0.420}{0.744}{-0.417}{0.046}
\SG{0.744}{-0.417}{0.751}{-0.414}{0.045}
\SG{0.751}{-0.414}{0.757}{-0.412}{0.047}
\SG{0.757}{-0.412}{0.763}{-0.409}{0.046}
\SG{0.763}{-0.409}{0.770}{-0.406}{0.042}
\SG{0.770}{-0.406}{0.776}{-0.404}{0.036}
\SG{0.776}{-0.404}{0.782}{-0.401}{0.033}
\SG{0.782}{-0.401}{0.789}{-0.399}{0.033}
\SG{0.789}{-0.399}{0.795}{-0.396}{0.030}
\SG{0.795}{-0.396}{0.801}{-0.393}{0.026}
\SG{0.801}{-0.393}{0.808}{-0.391}{0.024}
\SG{0.808}{-0.391}{0.814}{-0.388}{0.024}
\SG{0.814}{-0.388}{0.820}{-0.385}{0.025}
\SG{0.820}{-0.385}{0.827}{-0.383}{0.034}
\SG{0.827}{-0.383}{0.833}{-0.380}{0.046}
\SG{0.833}{-0.380}{0.839}{-0.378}{0.050}
\SG{0.839}{-0.378}{0.846}{-0.375}{0.050}
\SG{0.846}{-0.375}{0.852}{-0.372}{0.048}
\SG{0.852}{-0.372}{0.858}{-0.370}{0.049}
\SG{0.858}{-0.370}{0.865}{-0.367}{0.051}
\SG{0.865}{-0.367}{0.871}{-0.364}{0.048}
\SG{0.871}{-0.364}{0.877}{-0.362}{0.044}
\SG{0.877}{-0.362}{0.884}{-0.359}{0.042}
\SG{0.884}{-0.359}{0.890}{-0.356}{0.042}
\SG{0.890}{-0.356}{0.896}{-0.354}{0.045}
\SG{0.896}{-0.354}{0.903}{-0.351}{0.055}
\SG{0.903}{-0.351}{0.909}{-0.349}{0.069}
\SG{0.909}{-0.349}{0.916}{-0.346}{0.071}
\SG{0.916}{-0.346}{0.922}{-0.343}{0.070}
\SG{0.922}{-0.343}{0.928}{-0.341}{0.068}
\SG{0.928}{-0.341}{0.935}{-0.338}{0.070}
\SG{0.935}{-0.338}{0.941}{-0.335}{0.071}
\SG{0.941}{-0.335}{0.947}{-0.333}{0.068}
\SG{0.947}{-0.333}{0.954}{-0.330}{0.064}
\SG{0.954}{-0.330}{0.960}{-0.328}{0.062}
\SG{0.960}{-0.328}{0.966}{-0.325}{0.063}
\SG{0.966}{-0.325}{0.973}{-0.322}{0.068}
\SG{0.973}{-0.322}{0.979}{-0.320}{0.080}
\SG{0.979}{-0.320}{0.985}{-0.317}{0.095}
\SG{0.985}{-0.317}{0.992}{-0.314}{0.095}
\SG{0.992}{-0.314}{0.998}{-0.312}{0.094}
\SG{0.998}{-0.312}{1.004}{-0.309}{0.094}
\SG{1.004}{-0.309}{1.011}{-0.307}{0.096}
\SG{1.011}{-0.307}{1.017}{-0.304}{0.097}
\SG{1.017}{-0.304}{1.023}{-0.301}{0.094}
\SG{1.023}{-0.301}{1.030}{-0.299}{0.087}
\SG{1.030}{-0.299}{1.036}{-0.296}{0.086}
\SG{1.036}{-0.296}{1.042}{-0.293}{0.086}
\SG{1.042}{-0.293}{1.049}{-0.291}{0.085}
\SG{1.049}{-0.291}{1.055}{-0.288}{0.082}
\SG{1.055}{-0.288}{1.061}{-0.285}{0.079}
\SG{1.061}{-0.285}{1.068}{-0.283}{0.080}
\SG{1.068}{-0.283}{1.074}{-0.280}{0.082}
\SG{1.074}{-0.280}{1.081}{-0.278}{0.091}
\SG{1.081}{-0.278}{1.087}{-0.275}{0.104}
\SG{1.087}{-0.275}{1.093}{-0.272}{0.112}
\SG{1.093}{-0.272}{1.100}{-0.270}{0.112}
\SG{1.100}{-0.270}{1.106}{-0.267}{0.112}
\SG{1.106}{-0.267}{1.112}{-0.264}{0.113}
\SG{1.112}{-0.264}{1.119}{-0.262}{0.116}
\SG{1.119}{-0.262}{1.125}{-0.259}{0.116}
\SG{1.125}{-0.259}{1.131}{-0.257}{0.113}
\SG{1.131}{-0.257}{1.138}{-0.254}{0.111}
\SG{1.138}{-0.254}{1.144}{-0.251}{0.112}
\SG{1.144}{-0.251}{1.150}{-0.249}{0.115}
\SG{1.150}{-0.249}{1.157}{-0.246}{0.126}
\SG{1.157}{-0.246}{1.163}{-0.243}{0.141}
\SG{1.163}{-0.243}{1.169}{-0.241}{0.147}
\SG{1.169}{-0.241}{1.176}{-0.238}{0.149}
\SG{1.176}{-0.238}{1.182}{-0.235}{0.148}
\SG{1.182}{-0.235}{1.188}{-0.233}{0.151}
\SG{1.188}{-0.233}{1.195}{-0.230}{0.155}
\SG{1.195}{-0.230}{1.201}{-0.228}{0.152}
\SG{1.201}{-0.228}{1.207}{-0.225}{0.149}
\SG{1.207}{-0.225}{1.214}{-0.222}{0.146}
\SG{1.214}{-0.222}{1.220}{-0.220}{0.146}
\SG{1.220}{-0.220}{1.226}{-0.217}{0.147}
\SG{1.226}{-0.217}{1.233}{-0.214}{0.146}
\SG{1.233}{-0.214}{1.239}{-0.212}{0.143}
\SG{1.239}{-0.212}{1.245}{-0.209}{0.145}
\SG{1.245}{-0.209}{1.252}{-0.207}{0.147}
\SG{1.252}{-0.207}{1.258}{-0.204}{0.154}
\SG{1.258}{-0.204}{1.265}{-0.201}{0.166}
\SG{1.265}{-0.201}{1.271}{-0.199}{0.180}
\SG{1.271}{-0.199}{1.277}{-0.196}{0.183}
\SG{1.277}{-0.196}{1.284}{-0.193}{0.184}
\SG{1.284}{-0.193}{1.290}{-0.191}{0.185}
\SG{1.290}{-0.191}{1.296}{-0.188}{0.189}
\SG{1.296}{-0.188}{1.303}{-0.185}{0.191}
\SG{1.303}{-0.185}{1.309}{-0.183}{0.190}
\SG{1.309}{-0.183}{1.315}{-0.180}{0.189}
\SG{1.315}{-0.180}{1.322}{-0.178}{0.191}
\SG{1.322}{-0.178}{1.328}{-0.175}{0.194}
\SG{1.328}{-0.175}{1.334}{-0.172}{0.204}
\SG{1.334}{-0.172}{1.341}{-0.170}{0.219}
\SG{1.341}{-0.170}{1.347}{-0.167}{0.232}
\SG{1.347}{-0.167}{1.353}{-0.164}{0.235}
\SG{1.353}{-0.164}{1.360}{-0.162}{0.235}
\SG{1.360}{-0.162}{1.366}{-0.159}{0.239}
\SG{1.366}{-0.159}{1.372}{-0.157}{0.244}
\SG{1.372}{-0.157}{1.379}{-0.154}{0.246}
\SG{1.379}{-0.154}{1.385}{-0.151}{0.245}
\SG{1.385}{-0.151}{1.391}{-0.149}{0.243}
\SG{1.391}{-0.149}{1.398}{-0.146}{0.244}
\SG{1.398}{-0.146}{1.404}{-0.143}{0.247}
\SG{1.404}{-0.143}{1.410}{-0.141}{0.248}
\SG{1.410}{-0.141}{1.417}{-0.138}{0.248}
\SG{1.417}{-0.138}{1.423}{-0.135}{0.250}
\SG{1.423}{-0.135}{1.430}{-0.133}{0.254}
\SG{1.430}{-0.133}{1.436}{-0.130}{0.260}
\SG{1.436}{-0.130}{1.442}{-0.128}{0.273}
\SG{1.442}{-0.128}{1.449}{-0.125}{0.290}
\SG{1.449}{-0.125}{1.455}{-0.122}{0.298}
\SG{1.455}{-0.122}{1.461}{-0.120}{0.302}
\SG{1.461}{-0.120}{1.468}{-0.117}{0.305}
\SG{1.468}{-0.117}{1.474}{-0.114}{0.311}
\SG{1.474}{-0.114}{1.480}{-0.112}{0.317}
\SG{1.480}{-0.112}{1.487}{-0.109}{0.319}
\SG{1.487}{-0.109}{1.493}{-0.107}{0.320}
\SG{1.493}{-0.107}{1.499}{-0.104}{0.324}
\SG{1.499}{-0.104}{1.506}{-0.101}{0.330}
\SG{1.506}{-0.101}{1.512}{-0.099}{0.339}
\SG{1.512}{-0.099}{1.518}{-0.096}{0.356}
\SG{1.518}{-0.096}{1.525}{-0.093}{0.376}
\SG{1.525}{-0.093}{1.531}{-0.091}{0.383}
\SG{1.531}{-0.091}{1.537}{-0.088}{0.388}
\SG{1.537}{-0.088}{1.544}{-0.086}{0.393}
\SG{1.544}{-0.086}{1.550}{-0.083}{0.401}
\SG{1.550}{-0.083}{1.556}{-0.080}{0.410}
\SG{1.556}{-0.080}{1.563}{-0.078}{0.414}
\SG{1.563}{-0.078}{1.569}{-0.075}{0.417}
\SG{1.569}{-0.075}{1.575}{-0.072}{0.424}
\SG{1.575}{-0.072}{1.582}{-0.070}{0.433}
\SG{1.582}{-0.070}{1.588}{-0.067}{0.448}
\SG{1.588}{-0.067}{1.595}{-0.064}{0.470}
\SG{1.595}{-0.064}{1.601}{-0.062}{0.493}
\SG{1.601}{-0.062}{1.607}{-0.059}{0.503}
\SG{1.607}{-0.059}{1.614}{-0.057}{0.511}
\SG{1.614}{-0.057}{1.620}{-0.054}{0.522}
\SG{1.620}{-0.054}{1.626}{-0.051}{0.536}
\SG{1.626}{-0.051}{1.633}{-0.049}{0.548}
\SG{1.633}{-0.049}{1.639}{-0.046}{0.556}
\SG{1.639}{-0.046}{1.645}{-0.043}{0.562}
\SG{1.645}{-0.043}{1.652}{-0.041}{0.575}
\SG{1.652}{-0.041}{1.658}{-0.038}{0.589}
\SG{1.658}{-0.038}{1.664}{-0.036}{0.603}
\SG{1.664}{-0.036}{1.671}{-0.033}{0.615}
\SG{1.671}{-0.033}{1.677}{-0.030}{0.629}
\SG{1.677}{-0.030}{1.683}{-0.028}{0.648}
\SG{1.683}{-0.028}{1.690}{-0.025}{0.670}
\SG{1.690}{-0.025}{1.696}{-0.022}{0.701}
\SG{1.696}{-0.022}{1.702}{-0.020}{0.737}
\SG{1.702}{-0.020}{1.709}{-0.017}{0.764}
\SG{1.709}{-0.017}{1.715}{-0.014}{0.783}
\SG{1.715}{-0.014}{1.721}{-0.012}{0.803}
\SG{1.721}{-0.012}{1.728}{-0.009}{0.824}
\SG{1.728}{-0.009}{1.734}{-0.007}{0.850}
\SG{1.734}{-0.007}{1.740}{-0.004}{0.830}
\SG{1.740}{-0.004}{1.744}{-0.001}{0.869}
\SG{1.744}{-0.001}{1.744}{0.001}{1.000}
\SG{1.744}{0.001}{1.740}{0.004}{0.854}
\SG{1.740}{0.004}{1.734}{0.007}{0.792}
\SG{1.734}{0.007}{1.728}{0.009}{0.830}
\SG{1.728}{0.009}{1.721}{0.012}{0.850}
\SG{1.721}{0.012}{1.715}{0.014}{0.824}
\SG{1.715}{0.014}{1.709}{0.017}{0.803}
\SG{1.709}{0.017}{1.702}{0.020}{0.783}
\SG{1.702}{0.020}{1.696}{0.022}{0.764}
\SG{1.696}{0.022}{1.690}{0.025}{0.737}
\SG{1.690}{0.025}{1.683}{0.028}{0.701}
\SG{1.683}{0.028}{1.677}{0.030}{0.670}
\SG{1.677}{0.030}{1.671}{0.033}{0.648}
\SG{1.671}{0.033}{1.664}{0.036}{0.629}
\SG{1.664}{0.036}{1.658}{0.038}{0.615}
\SG{1.658}{0.038}{1.652}{0.041}{0.603}
\SG{1.652}{0.041}{1.645}{0.043}{0.589}
\SG{1.645}{0.043}{1.639}{0.046}{0.575}
\SG{1.639}{0.046}{1.633}{0.049}{0.562}
\SG{1.633}{0.049}{1.626}{0.051}{0.556}
\SG{1.626}{0.051}{1.620}{0.054}{0.548}
\SG{1.620}{0.054}{1.614}{0.057}{0.536}
\SG{1.614}{0.057}{1.607}{0.059}{0.522}
\SG{1.607}{0.059}{1.601}{0.062}{0.511}
\SG{1.601}{0.062}{1.595}{0.064}{0.503}
\SG{1.595}{0.064}{1.588}{0.067}{0.493}
\SG{1.588}{0.067}{1.582}{0.070}{0.470}
\SG{1.582}{0.070}{1.575}{0.072}{0.448}
\SG{1.575}{0.072}{1.569}{0.075}{0.433}
\SG{1.569}{0.075}{1.563}{0.078}{0.424}
\SG{1.563}{0.078}{1.556}{0.080}{0.417}
\SG{1.556}{0.080}{1.550}{0.083}{0.414}
\SG{1.550}{0.083}{1.544}{0.086}{0.410}
\SG{1.544}{0.086}{1.537}{0.088}{0.401}
\SG{1.537}{0.088}{1.531}{0.091}{0.393}
\SG{1.531}{0.091}{1.525}{0.093}{0.388}
\SG{1.525}{0.093}{1.518}{0.096}{0.383}
\SG{1.518}{0.096}{1.512}{0.099}{0.376}
\SG{1.512}{0.099}{1.506}{0.101}{0.356}
\SG{1.506}{0.101}{1.499}{0.104}{0.339}
\SG{1.499}{0.104}{1.493}{0.107}{0.330}
\SG{1.493}{0.107}{1.487}{0.109}{0.324}
\SG{1.487}{0.109}{1.480}{0.112}{0.320}
\SG{1.480}{0.112}{1.474}{0.114}{0.319}
\SG{1.474}{0.114}{1.468}{0.117}{0.317}
\SG{1.468}{0.117}{1.461}{0.120}{0.311}
\SG{1.461}{0.120}{1.455}{0.122}{0.305}
\SG{1.455}{0.122}{1.449}{0.125}{0.302}
\SG{1.449}{0.125}{1.442}{0.128}{0.298}
\SG{1.442}{0.128}{1.436}{0.130}{0.290}
\SG{1.436}{0.130}{1.430}{0.133}{0.273}
\SG{1.430}{0.133}{1.423}{0.135}{0.260}
\SG{1.423}{0.135}{1.417}{0.138}{0.254}
\SG{1.417}{0.138}{1.410}{0.141}{0.250}
\SG{1.410}{0.141}{1.404}{0.143}{0.248}
\SG{1.404}{0.143}{1.398}{0.146}{0.248}
\SG{1.398}{0.146}{1.391}{0.149}{0.247}
\SG{1.391}{0.149}{1.385}{0.151}{0.244}
\SG{1.385}{0.151}{1.379}{0.154}{0.243}
\SG{1.379}{0.154}{1.372}{0.157}{0.245}
\SG{1.372}{0.157}{1.366}{0.159}{0.246}
\SG{1.366}{0.159}{1.360}{0.162}{0.244}
\SG{1.360}{0.162}{1.353}{0.164}{0.239}
\SG{1.353}{0.164}{1.347}{0.167}{0.235}
\SG{1.347}{0.167}{1.341}{0.170}{0.235}
\SG{1.341}{0.170}{1.334}{0.172}{0.232}
\SG{1.334}{0.172}{1.328}{0.175}{0.219}
\SG{1.328}{0.175}{1.322}{0.178}{0.204}
\SG{1.322}{0.178}{1.315}{0.180}{0.194}
\SG{1.315}{0.180}{1.309}{0.183}{0.191}
\SG{1.309}{0.183}{1.303}{0.185}{0.189}
\SG{1.303}{0.185}{1.296}{0.188}{0.190}
\SG{1.296}{0.188}{1.290}{0.191}{0.191}
\SG{1.290}{0.191}{1.284}{0.193}{0.189}
\SG{1.284}{0.193}{1.277}{0.196}{0.185}
\SG{1.277}{0.196}{1.271}{0.199}{0.184}
\SG{1.271}{0.199}{1.265}{0.201}{0.183}
\SG{1.265}{0.201}{1.258}{0.204}{0.180}
\SG{1.258}{0.204}{1.252}{0.207}{0.166}
\SG{1.252}{0.207}{1.245}{0.209}{0.154}
\SG{1.245}{0.209}{1.239}{0.212}{0.147}
\SG{1.239}{0.212}{1.233}{0.214}{0.145}
\SG{1.233}{0.214}{1.226}{0.217}{0.143}
\SG{1.226}{0.217}{1.220}{0.220}{0.146}
\SG{1.220}{0.220}{1.214}{0.222}{0.147}
\SG{1.214}{0.222}{1.207}{0.225}{0.146}
\SG{1.207}{0.225}{1.201}{0.228}{0.146}
\SG{1.201}{0.228}{1.195}{0.230}{0.149}
\SG{1.195}{0.230}{1.188}{0.233}{0.152}
\SG{1.188}{0.233}{1.182}{0.235}{0.155}
\SG{1.182}{0.235}{1.176}{0.238}{0.151}
\SG{1.176}{0.238}{1.169}{0.241}{0.148}
\SG{1.169}{0.241}{1.163}{0.243}{0.149}
\SG{1.163}{0.243}{1.157}{0.246}{0.147}
\SG{1.157}{0.246}{1.150}{0.249}{0.141}
\SG{1.150}{0.249}{1.144}{0.251}{0.126}
\SG{1.144}{0.251}{1.138}{0.254}{0.115}
\SG{1.138}{0.254}{1.131}{0.257}{0.112}
\SG{1.131}{0.257}{1.125}{0.259}{0.111}
\SG{1.125}{0.259}{1.119}{0.262}{0.113}
\SG{1.119}{0.262}{1.112}{0.264}{0.116}
\SG{1.112}{0.264}{1.106}{0.267}{0.116}
\SG{1.106}{0.267}{1.100}{0.270}{0.113}
\SG{1.100}{0.270}{1.093}{0.272}{0.112}
\SG{1.093}{0.272}{1.087}{0.275}{0.112}
\SG{1.087}{0.275}{1.081}{0.278}{0.112}
\SG{1.081}{0.278}{1.074}{0.280}{0.104}
\SG{1.074}{0.280}{1.068}{0.283}{0.091}
\SG{1.068}{0.283}{1.061}{0.285}{0.082}
\SG{1.061}{0.285}{1.055}{0.288}{0.080}
\SG{1.055}{0.288}{1.049}{0.291}{0.079}
\SG{1.049}{0.291}{1.042}{0.293}{0.082}
\SG{1.042}{0.293}{1.036}{0.296}{0.085}
\SG{1.036}{0.296}{1.030}{0.299}{0.086}
\SG{1.030}{0.299}{1.023}{0.301}{0.086}
\SG{1.023}{0.301}{1.017}{0.304}{0.087}
\SG{1.017}{0.304}{1.011}{0.307}{0.094}
\SG{1.011}{0.307}{1.004}{0.309}{0.097}
\SG{1.004}{0.309}{0.998}{0.312}{0.096}
\SG{0.998}{0.312}{0.992}{0.314}{0.094}
\SG{0.992}{0.314}{0.985}{0.317}{0.094}
\SG{0.985}{0.317}{0.979}{0.320}{0.095}
\SG{0.979}{0.320}{0.973}{0.322}{0.095}
\SG{0.973}{0.322}{0.966}{0.325}{0.080}
\SG{0.966}{0.325}{0.960}{0.328}{0.068}
\SG{0.960}{0.328}{0.954}{0.330}{0.063}
\SG{0.954}{0.330}{0.947}{0.333}{0.062}
\SG{0.947}{0.333}{0.941}{0.335}{0.064}
\SG{0.941}{0.335}{0.935}{0.338}{0.068}
\SG{0.935}{0.338}{0.928}{0.341}{0.071}
\SG{0.928}{0.341}{0.922}{0.343}{0.070}
\SG{0.922}{0.343}{0.916}{0.346}{0.068}
\SG{0.916}{0.346}{0.909}{0.349}{0.070}
\SG{0.909}{0.349}{0.903}{0.351}{0.071}
\SG{0.903}{0.351}{0.896}{0.354}{0.069}
\SG{0.896}{0.354}{0.890}{0.356}{0.055}
\SG{0.890}{0.356}{0.884}{0.359}{0.045}
\SG{0.884}{0.359}{0.877}{0.362}{0.042}
\SG{0.877}{0.362}{0.871}{0.364}{0.042}
\SG{0.871}{0.364}{0.865}{0.367}{0.044}
\SG{0.865}{0.367}{0.858}{0.370}{0.048}
\SG{0.858}{0.370}{0.852}{0.372}{0.051}
\SG{0.852}{0.372}{0.846}{0.375}{0.049}
\SG{0.846}{0.375}{0.839}{0.378}{0.048}
\SG{0.839}{0.378}{0.833}{0.380}{0.050}
\SG{0.833}{0.380}{0.827}{0.383}{0.050}
\SG{0.827}{0.383}{0.820}{0.385}{0.046}
\SG{0.820}{0.385}{0.814}{0.388}{0.034}
\SG{0.814}{0.388}{0.808}{0.391}{0.025}
\SG{0.808}{0.391}{0.801}{0.393}{0.024}
\SG{0.801}{0.393}{0.795}{0.396}{0.024}
\SG{0.795}{0.396}{0.789}{0.399}{0.026}
\SG{0.789}{0.399}{0.782}{0.401}{0.030}
\SG{0.782}{0.401}{0.776}{0.404}{0.033}
\SG{0.776}{0.404}{0.770}{0.406}{0.033}
\SG{0.770}{0.406}{0.763}{0.409}{0.036}
\SG{0.763}{0.409}{0.757}{0.412}{0.042}
\SG{0.757}{0.412}{0.751}{0.414}{0.046}
\SG{0.751}{0.414}{0.744}{0.417}{0.047}
\SG{0.744}{0.417}{0.738}{0.420}{0.045}
\SG{0.738}{0.420}{0.731}{0.422}{0.046}
\SG{0.731}{0.422}{0.725}{0.425}{0.048}
\SG{0.725}{0.425}{0.719}{0.428}{0.048}
\SG{0.719}{0.428}{0.712}{0.430}{0.037}
\SG{0.712}{0.430}{0.706}{0.433}{0.026}
\SG{0.706}{0.433}{0.700}{0.435}{0.020}
\SG{0.700}{0.435}{0.693}{0.438}{0.020}
\SG{0.693}{0.438}{0.687}{0.441}{0.020}
\SG{0.687}{0.441}{0.681}{0.443}{0.026}
\SG{0.681}{0.443}{0.674}{0.446}{0.029}
\SG{0.674}{0.446}{0.668}{0.449}{0.029}
\SG{0.668}{0.449}{0.662}{0.451}{0.028}
\SG{0.662}{0.451}{0.655}{0.454}{0.030}
\SG{0.655}{0.454}{0.649}{0.456}{0.031}
\SG{0.649}{0.456}{0.643}{0.459}{0.031}
\SG{0.643}{0.459}{0.636}{0.462}{0.019}
\SG{0.636}{0.462}{0.630}{0.464}{0.009}
\SG{0.630}{0.464}{0.624}{0.467}{0.006}
\SG{0.624}{0.467}{0.617}{0.470}{0.006}
\SG{0.617}{0.470}{0.611}{0.472}{0.007}
\SG{0.611}{0.472}{0.605}{0.475}{0.012}
\SG{0.605}{0.475}{0.598}{0.478}{0.016}
\SG{0.598}{0.478}{0.592}{0.480}{0.017}
\SG{0.592}{0.480}{0.586}{0.483}{0.019}
\SG{0.586}{0.483}{0.579}{0.485}{0.025}
\SG{0.579}{0.485}{0.573}{0.488}{0.030}
\SG{0.573}{0.488}{0.566}{0.491}{0.034}
\SG{0.566}{0.491}{0.560}{0.493}{0.032}
\SG{0.560}{0.493}{0.554}{0.496}{0.032}
\SG{0.554}{0.496}{0.547}{0.499}{0.036}
\SG{0.547}{0.499}{0.541}{0.501}{0.036}
\SG{0.541}{0.501}{0.535}{0.504}{0.030}
\SG{0.535}{0.504}{0.528}{0.506}{0.018}
\SG{0.528}{0.506}{0.522}{0.509}{0.009}
\SG{0.522}{0.509}{0.516}{0.512}{0.009}
\SG{0.516}{0.512}{0.509}{0.514}{0.010}
\SG{0.509}{0.514}{0.503}{0.517}{0.014}
\SG{0.503}{0.517}{0.497}{0.520}{0.019}
\SG{0.497}{0.520}{0.490}{0.522}{0.021}
\SG{0.490}{0.522}{0.484}{0.525}{0.020}
\SG{0.484}{0.525}{0.478}{0.528}{0.020}
\SG{0.478}{0.528}{0.471}{0.530}{0.023}
\SG{0.471}{0.530}{0.465}{0.533}{0.024}
\SG{0.465}{0.533}{0.459}{0.535}{0.017}
\SG{0.459}{0.535}{0.452}{0.538}{0.006}
\SG{0.452}{0.538}{0.446}{0.541}{0.000}
\SG{0.446}{0.541}{0.440}{0.543}{0.000}
\SG{0.440}{0.543}{0.433}{0.546}{0.001}
\SG{0.433}{0.546}{0.427}{0.549}{0.006}
\SG{0.427}{0.549}{0.421}{0.551}{0.011}
\SG{0.421}{0.551}{0.414}{0.554}{0.014}
\SG{0.414}{0.554}{0.408}{0.556}{0.015}
\SG{0.408}{0.556}{0.401}{0.559}{0.020}
\SG{0.401}{0.559}{0.395}{0.562}{0.027}
\SG{0.395}{0.562}{0.389}{0.564}{0.032}
\SG{0.389}{0.564}{0.382}{0.567}{0.032}
\SG{0.382}{0.567}{0.376}{0.570}{0.032}
\SG{0.376}{0.570}{0.370}{0.572}{0.035}
\SG{0.370}{0.572}{0.363}{0.575}{0.037}
\SG{0.363}{0.575}{0.357}{0.577}{0.038}
\SG{0.357}{0.577}{0.351}{0.580}{0.024}
\SG{0.351}{0.580}{0.344}{0.583}{0.014}
\SG{0.344}{0.583}{0.338}{0.585}{0.011}
\SG{0.338}{0.585}{0.332}{0.588}{0.012}
\SG{0.332}{0.588}{0.325}{0.591}{0.016}
\SG{0.325}{0.591}{0.319}{0.593}{0.022}
\SG{0.319}{0.593}{0.313}{0.596}{0.027}
\SG{0.313}{0.596}{0.306}{0.599}{0.027}
\SG{0.306}{0.599}{0.300}{0.601}{0.027}
\SG{0.300}{0.601}{0.294}{0.604}{0.031}
\SG{0.294}{0.604}{0.287}{0.606}{0.033}
\SG{0.287}{0.606}{0.281}{0.609}{0.031}
\SG{0.281}{0.609}{0.275}{0.612}{0.019}
\SG{0.275}{0.612}{0.268}{0.614}{0.011}
\SG{0.268}{0.614}{0.262}{0.617}{0.011}
\SG{0.262}{0.617}{0.256}{0.620}{0.013}
\SG{0.256}{0.620}{0.249}{0.622}{0.016}
\SG{0.249}{0.622}{0.243}{0.625}{0.022}
\SG{0.243}{0.625}{0.236}{0.627}{0.026}
\SG{0.236}{0.627}{0.230}{0.630}{0.026}
\SG{0.230}{0.630}{0.224}{0.633}{0.027}
\SG{0.224}{0.633}{0.217}{0.635}{0.030}
\SG{0.217}{0.635}{0.211}{0.638}{0.032}
\SG{0.211}{0.638}{0.205}{0.641}{0.028}
\SG{0.205}{0.641}{0.198}{0.643}{0.017}
\SG{0.198}{0.643}{0.192}{0.646}{0.011}
\SG{0.192}{0.646}{0.186}{0.649}{0.012}
\SG{0.186}{0.649}{0.179}{0.651}{0.013}
\SG{0.179}{0.651}{0.173}{0.654}{0.018}
\SG{0.173}{0.654}{0.167}{0.656}{0.024}
\SG{0.167}{0.656}{0.160}{0.659}{0.028}
\SG{0.160}{0.659}{0.154}{0.662}{0.030}
\SG{0.154}{0.662}{0.148}{0.664}{0.034}
\SG{0.148}{0.664}{0.141}{0.667}{0.042}
\SG{0.141}{0.667}{0.135}{0.670}{0.048}
\SG{0.135}{0.670}{0.129}{0.672}{0.050}
\SG{0.129}{0.672}{0.122}{0.675}{0.050}
\SG{0.122}{0.675}{0.116}{0.677}{0.053}
\SG{0.116}{0.677}{0.110}{0.680}{0.056}
\SG{0.110}{0.680}{0.103}{0.683}{0.058}
\SG{0.103}{0.683}{0.097}{0.685}{0.047}
\SG{0.097}{0.685}{0.091}{0.688}{0.038}
\SG{0.091}{0.688}{0.084}{0.691}{0.035}
\SG{0.084}{0.691}{0.078}{0.693}{0.037}
\SG{0.078}{0.693}{0.071}{0.696}{0.039}
\SG{0.071}{0.696}{0.065}{0.699}{0.046}
\SG{0.065}{0.699}{0.059}{0.701}{0.052}
\SG{0.059}{0.701}{0.052}{0.704}{0.053}
\SG{0.052}{0.704}{0.046}{0.706}{0.054}
\SG{0.046}{0.706}{0.040}{0.709}{0.057}
\SG{0.040}{0.709}{0.033}{0.712}{0.060}
\SG{0.033}{0.712}{0.027}{0.714}{0.062}
\SG{0.027}{0.714}{0.021}{0.717}{0.051}
\SG{0.021}{0.717}{0.014}{0.720}{0.043}
\SG{0.014}{0.720}{0.008}{0.722}{0.042}
\SG{0.008}{0.722}{0.002}{0.725}{0.045}
\SG{0.002}{0.725}{-0.005}{0.727}{0.048}
\SG{-0.005}{0.727}{-0.011}{0.730}{0.055}
\SG{-0.011}{0.730}{-0.017}{0.733}{0.061}
\SG{-0.017}{0.733}{-0.024}{0.735}{0.064}
\SG{-0.024}{0.735}{-0.030}{0.738}{0.068}
\SG{-0.030}{0.738}{-0.036}{0.741}{0.076}
\SG{-0.036}{0.741}{-0.043}{0.743}{0.084}
\SG{-0.043}{0.743}{-0.049}{0.746}{0.089}
\SG{-0.049}{0.746}{-0.055}{0.749}{0.089}
\SG{-0.055}{0.749}{-0.062}{0.751}{0.091}
\SG{-0.062}{0.751}{-0.068}{0.754}{0.097}
\SG{-0.068}{0.754}{-0.074}{0.756}{0.099}
\SG{-0.074}{0.756}{-0.081}{0.759}{0.095}
\SG{-0.081}{0.759}{-0.087}{0.762}{0.085}
\SG{-0.087}{0.762}{-0.093}{0.764}{0.079}
\SG{-0.093}{0.764}{-0.100}{0.767}{0.081}
\SG{-0.100}{0.767}{-0.106}{0.770}{0.085}
\SG{-0.106}{0.770}{-0.113}{0.772}{0.092}
\SG{-0.113}{0.772}{-0.119}{0.775}{0.099}
\SG{-0.119}{0.775}{-0.125}{0.777}{0.104}
\SG{-0.125}{0.777}{-0.132}{0.780}{0.106}
\SG{-0.132}{0.780}{-0.138}{0.783}{0.108}
\SG{-0.138}{0.783}{-0.144}{0.785}{0.114}
\SG{-0.144}{0.785}{-0.151}{0.788}{0.118}
\SG{-0.151}{0.788}{-0.157}{0.791}{0.112}
\SG{-0.157}{0.791}{-0.163}{0.793}{0.104}
\SG{-0.163}{0.793}{-0.170}{0.796}{0.102}
\SG{-0.170}{0.796}{-0.176}{0.798}{0.106}
\SG{-0.176}{0.798}{-0.182}{0.801}{0.110}
\SG{-0.182}{0.801}{-0.189}{0.804}{0.119}
\SG{-0.189}{0.804}{-0.195}{0.806}{0.127}
\SG{-0.195}{0.806}{-0.201}{0.809}{0.132}
\SG{-0.201}{0.809}{-0.208}{0.812}{0.138}
\SG{-0.208}{0.812}{-0.214}{0.814}{0.147}
\SG{-0.214}{0.814}{-0.220}{0.817}{0.158}
\SG{-0.220}{0.817}{-0.227}{0.820}{0.167}
\SG{-0.227}{0.820}{-0.233}{0.822}{0.171}
\SG{-0.233}{0.822}{-0.239}{0.825}{0.175}
\SG{-0.239}{0.825}{-0.246}{0.827}{0.183}
\SG{-0.246}{0.827}{-0.251}{0.830}{0.189}
\SG{-0.251}{0.830}{-0.256}{0.832}{0.193}
\SG{-0.256}{0.832}{-0.259}{0.833}{0.187}

    \end{scope}
    \draw[azulclaro, line width=0.45pt, rounded corners=3pt] (-2.6,-2.0) rectangle (2.6,2.0);
    \node[anchor=north, align=center, text width=5.3cm, inner sep=0, font=\scriptsize,
          text=grispizarra] at (0,-2.2)
      {Líneas de campo y equipotenciales de un conductor cargado con punta;
       el color del contorno indica la densidad superficial $\sigma$.};
  \end{scope}

  % --- Palabras clave y clasificación -----------------------------------
  \node[anchor=north west, align=justify, inner sep=0, text width=16cm, font=\small] at ($(O)+(2.5,11.0)$)
    {{\scshape\color{azulmedio}\textls[60]{Palabras clave}}\enspace
     conductor perfecto · equilibrio electrostático · efecto punta · presión electrostática ·
     jaula de Faraday · coeficientes de potencial · coeficientes de capacidad e influencia ·
     influencia total · condensador\\[4pt]
     {\scshape\color{azulmedio}\textls[60]{Clasificación PACS}}\enspace
     41.20.Cv (Electrostática; ecuaciones de Poisson y Laplace, problemas de contorno)};

  % --- Expresiones centrales del tema ---------------------------------------
  \draw[azulclaro, line width=0.45pt] ($(O)+(2.5,8.55)$) -- ($(O)+(18.5,8.55)$);
  \node[anchor=north, align=center, inner sep=0, text width=16cm] at ($(O)+(10.5,8.1)$)
    {\color{azulprof}
     $\displaystyle\vE\big|_{\Sc^{+}}=\frac{\sigma}{\eo}\,\un\qquad
      \Pel=\frac{\sigma^{2}}{2\eo}\qquad
      \phi_i=\sum_{j=1}^{N}p_{ij}\,Q_j\qquad
      Q_i=\sum_{j=1}^{N}c_{ij}\,\phi_j\qquad
      C=\frac{1}{p_{11}+p_{22}-2p_{12}}$};
  \draw[azulclaro, line width=0.45pt] ($(O)+(2.5,6.25)$) -- ($(O)+(18.5,6.25)$);

  % --- Tipo de documento, lugar y curso ------------------------------------
  \node[anchor=south, align=center, inner sep=0, text width=16cm] at ($(O)+(10.5,2.2)$)
    {{\color{azulprof}\scshape\large\textls[100]{\tipotrabajo}}\\[5pt]
     {\color{grispizarra}\scshape\small\textls[80]{\miUbicacion}}};
\end{tikzpicture}
\tikzexternalenable
\end{titlepage}

\clearpage
\pagenumbering{arabic}
\setcounter{page}{1}

% --- ÍNDICE -----------------------------------------------------------------
{\hypersetup{linkcolor=azulprof}\tableofcontents}
\thispagestyle{plain}
\clearpage

% =========================================================================
\section{Introducción: el conductor perfecto}\label{sec:intro}
% =========================================================================
Hasta ahora todo nuestro estudio se ha desarrollado en el vacío: dadas las cargas, la ley de Coulomb, la ley de Gauss y el potencial
permiten calcular el campo~\cite{TEMA1}, y el desarrollo multipolar permite aproximarlo a grandes distancias~\cite{TEMA2}. En este tema
damos el primer paso hacia la materia. Los materiales se dividen, a efectos electrostáticos, en \emph{conductores} y \emph{dieléctricos};
aquí estudiamos los primeros. La novedad es profunda: en un conductor la distribución de carga \emph{no es un dato}, sino una incógnita que
se ajusta sola hasta alcanzar el equilibrio. Veremos que esa autorregulación impone condiciones tan fuertes que, paradójicamente, simplifica
el problema: basta conocer la geometría y unas pocas cargas o potenciales para determinarlo todo.

Seguimos la estructura de los apuntes manuscritos y de teoría de la asignatura~\cite{APUNTES-T3}, cuyas figuras~\cite{FIGURAS-T3} hemos
recalculado y redibujado; para ampliar, pueden consultarse los textos clásicos~\cite{GRIFFITHS,WANGSNESS,VICTORIANO,PURCELL,FEYNMAN}.

% -------------------------------------------------------------------------
\subsection{Conductor perfecto y equilibrio electrostático}\label{sec:conductor-perfecto}
% -------------------------------------------------------------------------
De forma genérica, un material conductor es un cuerpo en el que las cargas eléctricas pueden moverse libremente bajo la acción de un campo
eléctrico. El ejemplo más común son los metales (cobre, plata, oro, aluminio\ldots): en el cobre hay aproximadamente un electrón libre por
átomo. Pero no son los únicos: en el agua salada, por ejemplo, lo que se desplaza son los iones.

\begin{definicion}[Conductor perfecto]\label{def:conductor}
Un \emph{conductor perfecto} es un cuerpo material, que ocupa una región cerrada y acotada $\Vc\subset\R^{3}$ de frontera $\Sc=\partial\Vc$,
que cumple:
\begin{enumerate}[label=(\roman*)]
  \item las cargas libres se mueven por su interior sin oposición alguna (se desprecia la agitación térmica);
  \item admite un número ilimitado de cargas libres: dispone de una reserva inagotable de carga de ambos signos.
\end{enumerate}
\end{definicion}

\begin{observacion}[Una idealización excelente]\label{obs:idealizacion}
Ningún conductor real es perfecto: los metales tienen una resistividad no nula y una densidad de electrones libres finita
($n\simeq8{,}5\cdot10^{28}\,\mathrm{m^{-3}}$ en el cobre). Sin embargo, la aproximación funciona extraordinariamente bien en electrostática,
donde solo interesa el estado final de equilibrio. Las cargas que hay que redistribuir para anular un campo externo típico son una fracción
insignificante de las disponibles, y la redistribución es prácticamente instantánea: en un conductor óhmico de conductividad $\sigma_c$ cualquier
exceso de carga en el interior decae como $\eul^{-t/\tau}$, con $\tau=\eo/\sigma_c$, que para el cobre ($\sigma_c\simeq5{,}96\cdot10^{7}\,\mathrm{S/m}$)
vale $\tau\simeq1{,}5\cdot10^{-19}\,\mathrm{s}$~\cite{GRIFFITHS}.
\end{observacion}

\begin{definicion}[Equilibrio electrostático]\label{def:equilibrio}
Un conductor está en \emph{equilibrio electrostático} cuando sus cargas libres están en reposo a escala macroscópica, es decir, elemento de
volumen a elemento de volumen: no hay ningún flujo neto de carga en su interior ni sobre su superficie.
\end{definicion}

\begin{proposicion}[Campo nulo en el interior]\label{prop:campo-nulo}
En un conductor perfecto en equilibrio electrostático el campo eléctrico es nulo en todo punto interior:
\begin{equation}\label{eq:campo-nulo}
  \resalta{\vE(\vc{r})=\vc{0}\qquad\forall\,\vc{r}\in\Vc\setminus\Sc .}
\end{equation}
\end{proposicion}
\begin{proof}
Supongamos que existe un punto interior $\vc{r}_0$ con $\vE(\vc{r}_0)\neq\vc{0}$. Las cargas libres situadas en un entorno de $\vc{r}_0$ (que existen
por la hipótesis~(ii) de la \cref{def:conductor}) experimentan la fuerza $\dif\vF=\dif q\,\vE$, que por la hipótesis~(i) no encuentra
oposición: se pondrían en movimiento, lo que contradice la \cref{def:equilibrio}. Luego $\vE=\vc{0}$ en todo el interior.
\end{proof}

La demostración dice \emph{qué} ocurre en el equilibrio, pero no \emph{cómo} se llega a él. Supongamos que introducimos un conductor
descargado en un campo externo uniforme $\vE_0$ (\cref{fig:placa}). En los primeros instantes el campo penetra en el material y empuja a los
electrones libres en sentido opuesto a $\vE_0$: se acumulan en la cara por la que entran las líneas, que queda con carga negativa, mientras
que la cara opuesta, que pierde electrones, queda con carga positiva. Esta \emph{carga inducida} crea, en el interior, un campo $\vE_{\mathrm{ind}}$
dirigido de la carga positiva a la negativa, es decir, \emph{opuesto} a $\vE_0$. La redistribución continúa mientras quede campo neto en el
interior, y se detiene exactamente cuando
\[
  \vE=\vE_0+\vE_{\mathrm{ind}}=\vc{0}\qquad\text{en el interior.}
\]
La clave es, pues, que el campo creado por la redistribución de carga inducida por el campo externo se opone a dicho campo y acaba
cancelándolo. Fuera del conductor, en cambio, el campo inducido \emph{no} cancela al externo: lo deforma, de modo que las líneas de campo
terminan en las cargas negativas y vuelven a nacer en las positivas.

% =========================================================================
%  Figura: placa conductora en un campo externo uniforme.
%  (a) instante inicial: el campo penetra y desplaza a los electrones;
%  (b) equilibrio: la carga inducida crea E_ind = -E_0 en el interior.
% =========================================================================
\begin{figure}[t]
\centering
\begin{tikzpicture}[x=1cm, y=1cm, font=\small]
  % ---------------------------------------------------------------- (a)
  \begin{scope}
    \fill[conductorplano, fill=metal!45!white, draw=azulprof] (-0.8,-1.7) rectangle (0.8,1.7);
    \foreach \y in {-1.3,-0.65,0,0.65,1.3}{
      \draw[azulacero, line width=0.6pt] (-2.5,\y) -- (-0.8,\y);
      \draw[azulacero!55, line width=0.6pt, densely dashed] (-0.8,\y) -- (0.8,\y);
      \draw[azulacero, line width=0.6pt, -{Stealth[length=4pt]}] (0.8,\y) -- (2.5,\y);
      \draw[azulacero, line width=0.6pt, -{Stealth[length=4pt]}] (-2.5,\y) -- (-1.4,\y);}
    \foreach \y in {-0.98,-0.33,0.33,0.98}{
      \draw[-{Stealth[length=3.5pt]}, carganeg, line width=0.7pt] (0.3,\y) -- (-0.2,\y);
      \node[font=\scriptsize, text=carganeg, anchor=west, inner sep=1pt] at (0.32,\y) {$e^{-}$};}
    \node[above, text=azulprof] at (-1.6,1.75) {$\vE_0$};
    \node[text=azulprof, font=\footnotesize] at (0,-2.05) {conductor};
    \node[font=\small] at (0,-2.65) {(a) Instante inicial};
  \end{scope}
  % ---------------------------------------------------------------- (b)
  \begin{scope}[xshift=6.6cm]
    \fill[metal!45!white, draw=azulprof, line width=0.6pt] (-0.8,-1.7) rectangle (0.8,1.7);
    \foreach \y in {-1.3,-0.65,0,0.65,1.3}{
      \draw[azulacero, line width=0.6pt, -{Stealth[length=4pt]}] (-2.5,\y) -- (-1.4,\y);
      \draw[azulacero, line width=0.6pt] (-1.4,\y) -- (-0.87,\y);
      \draw[azulacero, line width=0.6pt, -{Stealth[length=4pt]}] (0.87,\y) -- (2.5,\y);}
    \foreach \y in {-1.5,-1.1,...,1.55}{
      \node[signoneg] at (-0.66,\y) {$-$};
      \node[signopos] at (0.66,\y) {$+$};}
    \draw[-{Stealth[length=3.5pt]}, cargapos, line width=0.8pt] (0.32,0.5) -- (-0.32,0.5);
    \node[font=\scriptsize, text=cargapos] at (0,0.74) {$\vE_{\mathrm{ind}}$};
    \draw[-{Stealth[length=3.5pt]}, azulacero, line width=0.8pt] (-0.32,-0.5) -- (0.32,-0.5);
    \node[font=\scriptsize, text=azulacero] at (0,-0.74) {$\vE_{0}$};
    \node[font=\footnotesize, text=azulprof, fill=white, inner sep=1.5pt, rounded corners=1pt] at (0,0.02) {$\vE=\vc{0}$};
    \node[above, text=azulprof] at (-1.6,1.75) {$\vE_0$};
    \node[text=carganeg, font=\footnotesize] at (-0.95,-2.05) {$\sigma<0$};
    \node[text=cargapos, font=\footnotesize] at (0.95,-2.05) {$\sigma>0$};
    \node[font=\small] at (0,-2.65) {(b) Equilibrio electrostático};
  \end{scope}
\end{tikzpicture}
\caption{Conductor neutro en un campo externo uniforme $\vE_0$. (a) Al principio el campo penetra en el metal (trazos) y arrastra a los
electrones libres en sentido opuesto a $\vE_0$. (b) En el equilibrio, la carga inducida en las caras crea en el interior un campo
$\vE_{\mathrm{ind}}=-\vE_0$ que anula al externo: las líneas terminan en la cara negativa y renacen en la positiva.}
\label{fig:placa}
\end{figure}


La \cref{fig:uniforme} muestra el resultado de ese proceso para dos conductores neutros: una esfera, cuya solución es analítica, y un
conductor de forma irregular, que hemos resuelto numéricamente. En ambos casos el campo exterior se deforma para llegar
\emph{perpendicularmente} a la superficie, la carga inducida se concentra en las zonas más curvadas y el interior queda libre de campo.
Las secciones siguientes demuestran cada uno de estos rasgos.

% =========================================================================
%  Figura «tipo Mathematica»: conductores neutros en un campo uniforme.
%  (a) esfera (solución analítica, plano meridiano): líneas de campo como
%      curvas de nivel de la función de Stokes Psi = (E0/2) y^2 (1 + 2a^3/r^3);
%  (b) conductor irregular (ecuación de Laplace 2D por diferencias finitas).
%  Datos: figuras/datos/esfera_uniforme_* e irregular_uniforme_*.
% =========================================================================
\begin{figure}[t]
\centering
\begin{tikzpicture}[font=\small]
  % ---------------------------------------------------------------- (a)
  \begin{scope}[x=1.04cm, y=1.04cm]
    \fill[azulfondo!60!white] (-2.6,-2.6) rectangle (2.6,2.6);
    \begin{scope}
      \clip (-2.6,-2.6) rectangle (2.6,2.6);
      \tikzset{numEQ/.style={azulcielo, line width=0.35pt, densely dashed}}
      % Esfera en campo uniforme - equipotenciales
% Generado por scripts/generar_datos.py (no editar a mano).
\EQ{(2.454,-2.600) (2.455,-2.567) (2.456,-2.535) (2.457,-2.502) (2.458,-2.470) (2.459,-2.438) (2.460,-2.405) (2.462,-2.373) (2.463,-2.340) (2.464,-2.308) (2.465,-2.275) (2.467,-2.243) (2.468,-2.210) (2.469,-2.178) (2.471,-2.145) (2.472,-2.112) (2.473,-2.080) (2.475,-2.047) (2.476,-2.015) (2.478,-1.982) (2.479,-1.950) (2.480,-1.917) (2.482,-1.885) (2.484,-1.853) (2.485,-1.820) (2.487,-1.788) (2.488,-1.755) (2.490,-1.723) (2.491,-1.690) (2.493,-1.657) (2.495,-1.625) (2.496,-1.593) (2.498,-1.560) (2.499,-1.528) (2.501,-1.495) (2.503,-1.462) (2.504,-1.430) (2.506,-1.397) (2.508,-1.365) (2.509,-1.333) (2.511,-1.300) (2.513,-1.268) (2.514,-1.235) (2.516,-1.202) (2.518,-1.170) (2.519,-1.137) (2.521,-1.105) (2.522,-1.073) (2.524,-1.040) (2.526,-1.007) (2.527,-0.975) (2.529,-0.942) (2.530,-0.910) (2.532,-0.877) (2.533,-0.845) (2.534,-0.812) (2.536,-0.780) (2.537,-0.747) (2.538,-0.715) (2.540,-0.682) (2.541,-0.650) (2.542,-0.617) (2.543,-0.585) (2.544,-0.552) (2.545,-0.520) (2.546,-0.487) (2.547,-0.455) (2.548,-0.422) (2.549,-0.390) (2.549,-0.357) (2.550,-0.325) (2.551,-0.292) (2.551,-0.260) (2.552,-0.228) (2.552,-0.195) (2.553,-0.163) (2.553,-0.130) (2.553,-0.097) (2.553,-0.065) (2.553,-0.032) (2.553,0.000) (2.553,0.033) (2.553,0.065) (2.553,0.098) (2.553,0.130) (2.553,0.163) (2.552,0.195) (2.552,0.228) (2.551,0.260) (2.551,0.292) (2.550,0.325) (2.549,0.358) (2.549,0.390) (2.548,0.423) (2.547,0.455) (2.546,0.488) (2.545,0.520) (2.544,0.553) (2.543,0.585) (2.542,0.618) (2.541,0.650) (2.540,0.683) (2.538,0.715) (2.537,0.748) (2.536,0.780) (2.534,0.812) (2.533,0.845) (2.532,0.878) (2.530,0.910) (2.529,0.943) (2.527,0.975) (2.526,1.008) (2.524,1.040) (2.522,1.073) (2.521,1.105) (2.519,1.138) (2.518,1.170) (2.516,1.203) (2.514,1.235) (2.513,1.268) (2.511,1.300) (2.509,1.333) (2.508,1.365) (2.506,1.398) (2.504,1.430) (2.503,1.462) (2.501,1.495) (2.499,1.528) (2.498,1.560) (2.496,1.593) (2.495,1.625) (2.493,1.658) (2.491,1.690) (2.490,1.723) (2.488,1.755) (2.487,1.788) (2.485,1.820) (2.484,1.853) (2.482,1.885) (2.480,1.917) (2.479,1.950) (2.478,1.983) (2.476,2.015) (2.475,2.047) (2.473,2.080) (2.472,2.113) (2.471,2.145) (2.469,2.178) (2.468,2.210) (2.467,2.243) (2.465,2.275) (2.464,2.308) (2.463,2.340) (2.462,2.373) (2.460,2.405) (2.459,2.438) (2.458,2.470) (2.457,2.501) (2.456,2.535) (2.455,2.568) (2.454,2.600)}
\EQ{(2.056,-2.600) (2.058,-2.567) (2.059,-2.535) (2.060,-2.502) (2.062,-2.470) (2.063,-2.438) (2.065,-2.405) (2.066,-2.373) (2.068,-2.340) (2.069,-2.308) (2.071,-2.275) (2.073,-2.243) (2.074,-2.210) (2.076,-2.178) (2.078,-2.145) (2.080,-2.112) (2.082,-2.080) (2.084,-2.047) (2.086,-2.015) (2.087,-1.982) (2.090,-1.950) (2.092,-1.917) (2.094,-1.885) (2.096,-1.853) (2.098,-1.820) (2.100,-1.788) (2.102,-1.755) (2.105,-1.723) (2.107,-1.690) (2.109,-1.657) (2.112,-1.625) (2.114,-1.593) (2.116,-1.560) (2.119,-1.528) (2.121,-1.495) (2.124,-1.462) (2.126,-1.430) (2.129,-1.397) (2.131,-1.365) (2.134,-1.333) (2.137,-1.300) (2.139,-1.268) (2.142,-1.235) (2.144,-1.202) (2.147,-1.170) (2.149,-1.137) (2.152,-1.105) (2.155,-1.073) (2.157,-1.040) (2.160,-1.007) (2.162,-0.975) (2.165,-0.943) (2.167,-0.910) (2.169,-0.877) (2.172,-0.845) (2.174,-0.812) (2.176,-0.780) (2.178,-0.747) (2.180,-0.715) (2.183,-0.682) (2.185,-0.650) (2.186,-0.617) (2.188,-0.585) (2.190,-0.552) (2.192,-0.520) (2.193,-0.487) (2.195,-0.455) (2.196,-0.422) (2.198,-0.390) (2.199,-0.357) (2.200,-0.325) (2.201,-0.292) (2.202,-0.260) (2.203,-0.228) (2.204,-0.197) (2.204,-0.163) (2.205,-0.130) (2.205,-0.097) (2.205,-0.065) (2.206,-0.032) (2.206,0.000) (2.206,0.033) (2.205,0.065) (2.205,0.098) (2.205,0.130) (2.204,0.163) (2.204,0.195) (2.203,0.228) (2.202,0.260) (2.201,0.292) (2.200,0.325) (2.199,0.358) (2.198,0.390) (2.196,0.423) (2.195,0.455) (2.193,0.488) (2.192,0.520) (2.190,0.553) (2.188,0.585) (2.186,0.618) (2.185,0.650) (2.183,0.683) (2.180,0.715) (2.178,0.748) (2.176,0.780) (2.174,0.812) (2.172,0.845) (2.169,0.878) (2.167,0.910) (2.165,0.943) (2.162,0.975) (2.160,1.008) (2.157,1.040) (2.155,1.073) (2.152,1.105) (2.149,1.138) (2.147,1.170) (2.144,1.203) (2.142,1.235) (2.139,1.268) (2.137,1.300) (2.134,1.333) (2.131,1.365) (2.129,1.398) (2.126,1.430) (2.124,1.462) (2.121,1.495) (2.119,1.526) (2.116,1.560) (2.114,1.593) (2.112,1.625) (2.109,1.658) (2.107,1.690) (2.105,1.723) (2.102,1.755) (2.100,1.788) (2.098,1.820) (2.096,1.853) (2.094,1.885) (2.092,1.917) (2.090,1.950) (2.087,1.983) (2.086,2.015) (2.084,2.047) (2.082,2.080) (2.080,2.113) (2.078,2.145) (2.076,2.178) (2.074,2.210) (2.073,2.243) (2.071,2.275) (2.069,2.308) (2.068,2.340) (2.066,2.373) (2.065,2.405) (2.063,2.438) (2.062,2.470) (2.061,2.502) (2.059,2.535) (2.058,2.568) (2.056,2.600)}
\EQ{(1.657,-2.600) (1.658,-2.567) (1.660,-2.535) (1.661,-2.502) (1.663,-2.470) (1.665,-2.438) (1.667,-2.405) (1.668,-2.373) (1.670,-2.340) (1.672,-2.308) (1.674,-2.275) (1.676,-2.243) (1.679,-2.210) (1.681,-2.178) (1.683,-2.145) (1.685,-2.112) (1.688,-2.080) (1.690,-2.047) (1.693,-2.015) (1.696,-1.982) (1.698,-1.950) (1.701,-1.917) (1.704,-1.885) (1.707,-1.853) (1.710,-1.820) (1.713,-1.788) (1.716,-1.755) (1.719,-1.723) (1.723,-1.691) (1.726,-1.657) (1.729,-1.625) (1.733,-1.593) (1.737,-1.560) (1.740,-1.528) (1.744,-1.495) (1.748,-1.462) (1.752,-1.430) (1.755,-1.397) (1.759,-1.365) (1.763,-1.333) (1.767,-1.300) (1.771,-1.268) (1.776,-1.235) (1.780,-1.202) (1.784,-1.170) (1.788,-1.137) (1.792,-1.105) (1.796,-1.073) (1.800,-1.040) (1.804,-1.007) (1.809,-0.975) (1.813,-0.942) (1.817,-0.910) (1.821,-0.877) (1.824,-0.845) (1.828,-0.812) (1.832,-0.780) (1.836,-0.747) (1.839,-0.715) (1.843,-0.682) (1.846,-0.652) (1.849,-0.617) (1.853,-0.586) (1.856,-0.552) (1.859,-0.520) (1.861,-0.487) (1.864,-0.455) (1.866,-0.422) (1.869,-0.390) (1.871,-0.357) (1.873,-0.325) (1.875,-0.292) (1.876,-0.260) (1.878,-0.228) (1.879,-0.195) (1.880,-0.163) (1.881,-0.130) (1.881,-0.097) (1.882,-0.065) (1.882,-0.032) (1.882,0.000) (1.882,0.033) (1.882,0.065) (1.881,0.098) (1.881,0.130) (1.880,0.163) (1.879,0.195) (1.878,0.228) (1.876,0.260) (1.875,0.292) (1.873,0.325) (1.871,0.358) (1.869,0.390) (1.866,0.423) (1.864,0.455) (1.861,0.488) (1.859,0.520) (1.856,0.553) (1.853,0.585) (1.849,0.618) (1.846,0.650) (1.843,0.683) (1.840,0.713) (1.836,0.748) (1.832,0.780) (1.828,0.812) (1.824,0.845) (1.821,0.878) (1.817,0.910) (1.813,0.943) (1.809,0.975) (1.804,1.008) (1.800,1.039) (1.796,1.073) (1.792,1.105) (1.788,1.138) (1.784,1.170) (1.780,1.203) (1.776,1.235) (1.771,1.268) (1.767,1.300) (1.763,1.333) (1.759,1.365) (1.755,1.398) (1.752,1.430) (1.748,1.462) (1.744,1.495) (1.740,1.528) (1.737,1.560) (1.733,1.593) (1.729,1.625) (1.726,1.658) (1.723,1.690) (1.719,1.723) (1.716,1.755) (1.713,1.788) (1.710,1.820) (1.707,1.853) (1.704,1.885) (1.701,1.917) (1.698,1.950) (1.696,1.983) (1.693,2.015) (1.690,2.047) (1.688,2.080) (1.685,2.112) (1.683,2.145) (1.681,2.178) (1.679,2.210) (1.676,2.243) (1.674,2.275) (1.672,2.308) (1.670,2.340) (1.668,2.373) (1.667,2.405) (1.665,2.438) (1.663,2.470) (1.661,2.502) (1.660,2.535) (1.658,2.568) (1.657,2.600)}
\EQ{(1.252,-2.600) (1.254,-2.567) (1.255,-2.535) (1.257,-2.502) (1.259,-2.470) (1.261,-2.438) (1.263,-2.405) (1.265,-2.373) (1.267,-2.340) (1.269,-2.308) (1.272,-2.275) (1.274,-2.243) (1.277,-2.210) (1.279,-2.178) (1.282,-2.145) (1.285,-2.112) (1.288,-2.080) (1.291,-2.047) (1.294,-2.015) (1.298,-1.982) (1.301,-1.950) (1.305,-1.917) (1.308,-1.885) (1.312,-1.853) (1.316,-1.820) (1.320,-1.788) (1.325,-1.755) (1.329,-1.723) (1.334,-1.690) (1.338,-1.657) (1.343,-1.625) (1.348,-1.593) (1.354,-1.560) (1.359,-1.527) (1.365,-1.495) (1.370,-1.462) (1.376,-1.430) (1.382,-1.397) (1.388,-1.365) (1.394,-1.333) (1.401,-1.300) (1.407,-1.268) (1.414,-1.235) (1.420,-1.202) (1.427,-1.170) (1.434,-1.137) (1.441,-1.105) (1.448,-1.073) (1.454,-1.040) (1.461,-1.007) (1.468,-0.975) (1.475,-0.942) (1.482,-0.910) (1.489,-0.878) (1.495,-0.846) (1.502,-0.813) (1.508,-0.781) (1.514,-0.747) (1.521,-0.715) (1.527,-0.682) (1.532,-0.650) (1.538,-0.617) (1.543,-0.585) (1.548,-0.552) (1.553,-0.520) (1.558,-0.487) (1.563,-0.455) (1.567,-0.424) (1.571,-0.390) (1.574,-0.357) (1.578,-0.325) (1.581,-0.292) (1.583,-0.260) (1.586,-0.228) (1.588,-0.195) (1.590,-0.163) (1.591,-0.130) (1.592,-0.097) (1.593,-0.065) (1.594,-0.032) (1.594,0.000) (1.594,0.033) (1.593,0.065) (1.592,0.098) (1.591,0.130) (1.590,0.163) (1.588,0.195) (1.586,0.228) (1.583,0.260) (1.581,0.292) (1.578,0.325) (1.574,0.358) (1.571,0.390) (1.567,0.423) (1.563,0.455) (1.558,0.488) (1.554,0.519) (1.548,0.553) (1.543,0.585) (1.538,0.618) (1.532,0.650) (1.527,0.683) (1.521,0.712) (1.515,0.747) (1.508,0.780) (1.502,0.812) (1.495,0.845) (1.489,0.878) (1.482,0.909) (1.476,0.940) (1.469,0.971) (1.463,1.001) (1.456,1.033) (1.450,1.063) (1.443,1.094) (1.437,1.125) (1.430,1.156) (1.424,1.187) (1.417,1.219) (1.411,1.251) (1.405,1.281) (1.398,1.313) (1.392,1.346) (1.386,1.378) (1.380,1.411) (1.374,1.443) (1.368,1.476) (1.362,1.508) (1.357,1.541) (1.352,1.573) (1.346,1.606) (1.341,1.638) (1.336,1.671) (1.332,1.703) (1.327,1.736) (1.323,1.768) (1.319,1.800) (1.315,1.833) (1.311,1.866) (1.307,1.898) (1.303,1.930) (1.300,1.963) (1.296,1.996) (1.293,2.028) (1.290,2.061) (1.287,2.090) (1.284,2.126) (1.281,2.158) (1.278,2.191) (1.276,2.223) (1.273,2.256) (1.271,2.288) (1.269,2.321) (1.266,2.353) (1.264,2.386) (1.262,2.418) (1.260,2.451) (1.258,2.483) (1.257,2.516) (1.255,2.548) (1.253,2.581) (1.252,2.600)}
\EQ{(0.841,-2.600) (0.843,-2.567) (0.844,-2.535) (0.846,-2.502) (0.848,-2.470) (0.849,-2.438) (0.851,-2.405) (0.853,-2.373) (0.855,-2.340) (0.857,-2.308) (0.860,-2.275) (0.862,-2.243) (0.864,-2.212) (0.867,-2.178) (0.870,-2.145) (0.873,-2.112) (0.876,-2.080) (0.879,-2.047) (0.883,-2.015) (0.887,-1.982) (0.890,-1.950) (0.894,-1.917) (0.899,-1.885) (0.903,-1.853) (0.908,-1.820) (0.913,-1.788) (0.918,-1.755) (0.924,-1.723) (0.930,-1.690) (0.936,-1.657) (0.942,-1.625) (0.949,-1.593) (0.956,-1.563) (0.964,-1.527) (0.971,-1.495) (0.980,-1.462) (0.988,-1.431) (0.997,-1.397) (1.006,-1.365) (1.016,-1.333) (1.026,-1.300) (1.036,-1.268) (1.047,-1.236) (1.058,-1.202) (1.069,-1.170) (1.079,-1.140) (1.089,-1.111) (1.101,-1.079) (1.112,-1.049) (1.124,-1.014) (1.136,-0.982) (1.148,-0.949) (1.159,-0.916) (1.170,-0.886) (1.182,-0.851) (1.194,-0.819) (1.205,-0.786) (1.215,-0.754) (1.226,-0.722) (1.235,-0.693) (1.244,-0.663) (1.254,-0.630) (1.263,-0.598) (1.272,-0.566) (1.280,-0.533) (1.288,-0.500) (1.296,-0.468) (1.303,-0.435) (1.309,-0.403) (1.315,-0.370) (1.321,-0.338) (1.326,-0.307) (1.331,-0.273) (1.335,-0.240) (1.339,-0.208) (1.342,-0.175) (1.344,-0.143) (1.346,-0.111) (1.348,-0.078) (1.349,-0.045) (1.349,-0.013) (1.349,0.020) (1.349,0.052) (1.348,0.085) (1.346,0.117) (1.344,0.150) (1.341,0.182) (1.338,0.215) (1.334,0.247) (1.330,0.280) (1.325,0.312) (1.320,0.345) (1.314,0.377) (1.308,0.409) (1.301,0.442) (1.294,0.475) (1.287,0.505) (1.278,0.540) (1.270,0.572) (1.261,0.605) (1.252,0.637) (1.242,0.670) (1.232,0.702) (1.222,0.734) (1.211,0.767) (1.200,0.800) (1.190,0.831) (1.178,0.865) (1.166,0.897) (1.155,0.930) (1.144,0.959) (1.133,0.988) (1.122,1.021) (1.110,1.053) (1.099,1.085) (1.087,1.118) (1.075,1.151) (1.064,1.183) (1.053,1.216) (1.042,1.248) (1.032,1.281) (1.022,1.313) (1.012,1.346) (1.003,1.378) (0.993,1.411) (0.985,1.443) (0.976,1.476) (0.969,1.507) (0.961,1.541) (0.953,1.573) (0.946,1.606) (0.940,1.638) (0.933,1.671) (0.927,1.703) (0.921,1.736) (0.917,1.765) (0.911,1.800) (0.906,1.833) (0.901,1.866) (0.897,1.898) (0.893,1.930) (0.889,1.963) (0.885,1.996) (0.882,2.028) (0.878,2.061) (0.875,2.093) (0.872,2.126) (0.869,2.158) (0.866,2.191) (0.864,2.223) (0.861,2.256) (0.859,2.288) (0.857,2.321) (0.854,2.353) (0.852,2.385) (0.851,2.418) (0.849,2.451) (0.847,2.483) (0.845,2.516) (0.844,2.548) (0.842,2.581) (0.841,2.600)}
\EQ{(0.423,-2.600) (0.424,-2.567) (0.425,-2.535) (0.426,-2.502) (0.427,-2.470) (0.428,-2.438) (0.429,-2.405) (0.431,-2.373) (0.432,-2.340) (0.433,-2.308) (0.435,-2.275) (0.437,-2.243) (0.438,-2.210) (0.440,-2.178) (0.442,-2.146) (0.444,-2.112) (0.446,-2.080) (0.449,-2.047) (0.451,-2.015) (0.454,-1.982) (0.457,-1.950) (0.460,-1.917) (0.463,-1.885) (0.467,-1.853) (0.471,-1.820) (0.475,-1.788) (0.480,-1.755) (0.485,-1.723) (0.490,-1.690) (0.496,-1.657) (0.502,-1.625) (0.509,-1.593) (0.516,-1.560) (0.525,-1.528) (0.533,-1.496) (0.543,-1.462) (0.553,-1.433) (0.563,-1.404) (0.575,-1.371) (0.588,-1.339) (0.602,-1.306) (0.618,-1.274) (0.634,-1.242) (0.650,-1.212) (0.666,-1.183) (0.681,-1.157) (0.697,-1.131) (0.714,-1.105) (0.730,-1.079) (0.747,-1.053) (0.764,-1.027) (0.782,-1.001) (0.799,-0.975) (0.816,-0.949) (0.833,-0.923) (0.850,-0.897) (0.867,-0.871) (0.884,-0.845) (0.900,-0.819) (0.915,-0.793) (0.931,-0.767) (0.949,-0.735) (0.967,-0.702) (0.982,-0.674) (0.997,-0.643) (1.012,-0.611) (1.027,-0.579) (1.040,-0.548) (1.051,-0.520) (1.064,-0.487) (1.075,-0.455) (1.086,-0.425) (1.096,-0.390) (1.105,-0.358) (1.113,-0.325) (1.121,-0.292) (1.128,-0.260) (1.134,-0.228) (1.139,-0.195) (1.143,-0.163) (1.146,-0.130) (1.149,-0.097) (1.151,-0.065) (1.152,-0.032) (1.153,0.000) (1.152,0.033) (1.151,0.065) (1.149,0.098) (1.146,0.130) (1.143,0.163) (1.139,0.195) (1.134,0.228) (1.128,0.260) (1.121,0.292) (1.113,0.325) (1.105,0.358) (1.096,0.390) (1.086,0.423) (1.075,0.455) (1.064,0.488) (1.053,0.516) (1.041,0.546) (1.027,0.579) (1.014,0.607) (1.001,0.635) (0.987,0.663) (0.970,0.696) (0.956,0.723) (0.938,0.754) (0.923,0.780) (0.908,0.806) (0.892,0.832) (0.875,0.858) (0.859,0.884) (0.842,0.910) (0.825,0.936) (0.808,0.962) (0.790,0.988) (0.774,1.013) (0.756,1.040) (0.739,1.066) (0.722,1.092) (0.705,1.118) (0.689,1.144) (0.674,1.170) (0.659,1.196) (0.644,1.224) (0.627,1.255) (0.611,1.287) (0.598,1.316) (0.585,1.346) (0.572,1.378) (0.560,1.411) (0.549,1.443) (0.540,1.474) (0.530,1.508) (0.521,1.540) (0.514,1.572) (0.506,1.606) (0.499,1.638) (0.493,1.671) (0.488,1.703) (0.483,1.736) (0.478,1.768) (0.473,1.800) (0.469,1.833) (0.465,1.866) (0.462,1.898) (0.459,1.930) (0.456,1.963) (0.453,1.996) (0.450,2.028) (0.448,2.061) (0.445,2.093) (0.443,2.126) (0.441,2.158) (0.439,2.191) (0.438,2.223) (0.436,2.256) (0.434,2.288) (0.433,2.321) (0.432,2.353) (0.430,2.385) (0.429,2.417) (0.428,2.451) (0.427,2.483) (0.426,2.516) (0.425,2.548) (0.424,2.581) (0.423,2.600)}
\EQ{(0.000,-2.600) (0.000,-2.567) (0.000,-2.535) (0.000,-2.502) (0.000,-2.470) (0.000,-2.438) (0.000,-2.405) (0.000,-2.373) (0.000,-2.340) (0.000,-2.307) (0.000,-2.275) (0.000,-2.243) (0.000,-2.210) (0.000,-2.178) (0.000,-2.145) (0.000,-2.113) (0.000,-2.080) (0.000,-2.047) (0.000,-2.015) (0.000,-1.982) (0.000,-1.950) (0.000,-1.918) (0.000,-1.885) (0.000,-1.853) (0.000,-1.820) (0.000,-1.788) (0.000,-1.755) (0.000,-1.723) (0.000,-1.690) (0.000,-1.657) (0.000,-1.625) (0.000,-1.593) (0.000,-1.560) (0.000,-1.528) (0.000,-1.495) (0.000,-1.462) (0.000,-1.430) (0.000,-1.398) (0.000,-1.365) (0.000,-1.333) (0.000,-1.300) (0.000,-1.268) (0.000,-1.235) (0.000,-1.202) (0.000,-1.170) (0.000,-1.138) (0.000,-1.105) (0.000,-1.073) (0.000,-1.040) (0.000,-1.008) (0.000,-1.001)}
\EQ{(0.000,1.001) (0.000,1.034) (0.000,1.066) (0.000,1.099) (0.000,1.131) (0.000,1.164) (0.000,1.196) (0.000,1.229) (0.000,1.261) (0.000,1.294) (0.000,1.326) (0.000,1.359) (0.000,1.391) (0.000,1.424) (0.000,1.456) (0.000,1.489) (0.000,1.521) (0.000,1.554) (0.000,1.586) (0.000,1.619) (0.000,1.651) (0.000,1.683) (0.000,1.716) (0.000,1.749) (0.000,1.781) (0.000,1.813) (0.000,1.846) (0.000,1.879) (0.000,1.911) (0.000,1.944) (0.000,1.976) (0.000,2.009) (0.000,2.041) (0.000,2.074) (0.000,2.106) (0.000,2.139) (0.000,2.171) (0.000,2.204) (0.000,2.236) (0.000,2.268) (0.000,2.301) (0.000,2.334) (0.000,2.366) (0.000,2.399) (0.000,2.431) (0.000,2.464) (0.000,2.496) (0.000,2.529) (0.000,2.561) (0.000,2.594) (0.000,2.600)}
\EQ{(-0.423,-2.600) (-0.424,-2.567) (-0.425,-2.535) (-0.426,-2.502) (-0.427,-2.470) (-0.428,-2.438) (-0.429,-2.405) (-0.431,-2.373) (-0.432,-2.340) (-0.433,-2.308) (-0.435,-2.275) (-0.437,-2.243) (-0.438,-2.210) (-0.440,-2.178) (-0.442,-2.146) (-0.444,-2.112) (-0.446,-2.080) (-0.449,-2.047) (-0.451,-2.015) (-0.454,-1.982) (-0.457,-1.950) (-0.460,-1.918) (-0.463,-1.885) (-0.467,-1.853) (-0.471,-1.820) (-0.475,-1.788) (-0.480,-1.755) (-0.485,-1.723) (-0.490,-1.690) (-0.496,-1.658) (-0.502,-1.625) (-0.509,-1.593) (-0.516,-1.560) (-0.525,-1.528) (-0.533,-1.496) (-0.543,-1.462) (-0.552,-1.433) (-0.563,-1.404) (-0.575,-1.371) (-0.588,-1.339) (-0.602,-1.306) (-0.617,-1.274) (-0.634,-1.241) (-0.650,-1.212) (-0.666,-1.183) (-0.681,-1.157) (-0.697,-1.131) (-0.714,-1.105) (-0.730,-1.079) (-0.747,-1.053) (-0.764,-1.027) (-0.782,-1.001) (-0.799,-0.975) (-0.816,-0.949) (-0.833,-0.923) (-0.850,-0.897) (-0.867,-0.871) (-0.884,-0.845) (-0.900,-0.819) (-0.915,-0.793) (-0.931,-0.767) (-0.949,-0.735) (-0.967,-0.702) (-0.982,-0.674) (-0.997,-0.643) (-1.012,-0.611) (-1.027,-0.579) (-1.040,-0.548) (-1.051,-0.520) (-1.064,-0.487) (-1.075,-0.455) (-1.085,-0.425) (-1.096,-0.390) (-1.105,-0.358) (-1.113,-0.325) (-1.121,-0.292) (-1.128,-0.260) (-1.134,-0.228) (-1.139,-0.195) (-1.143,-0.163) (-1.146,-0.130) (-1.149,-0.097) (-1.151,-0.065) (-1.152,-0.032) (-1.153,0.000) (-1.152,0.033) (-1.151,0.065) (-1.149,0.098) (-1.146,0.130) (-1.143,0.163) (-1.139,0.195) (-1.134,0.228) (-1.128,0.260) (-1.121,0.292) (-1.113,0.325) (-1.105,0.358) (-1.096,0.390) (-1.086,0.423) (-1.075,0.455) (-1.064,0.488) (-1.053,0.516) (-1.041,0.546) (-1.027,0.579) (-1.014,0.607) (-1.001,0.635) (-0.987,0.663) (-0.970,0.696) (-0.956,0.723) (-0.938,0.754) (-0.923,0.780) (-0.908,0.806) (-0.892,0.832) (-0.875,0.858) (-0.859,0.884) (-0.842,0.910) (-0.825,0.936) (-0.808,0.962) (-0.790,0.988) (-0.773,1.013) (-0.756,1.040) (-0.739,1.066) (-0.722,1.092) (-0.705,1.118) (-0.689,1.144) (-0.674,1.170) (-0.659,1.196) (-0.643,1.224) (-0.627,1.255) (-0.611,1.287) (-0.598,1.316) (-0.585,1.346) (-0.572,1.378) (-0.560,1.411) (-0.549,1.443) (-0.539,1.474) (-0.530,1.508) (-0.521,1.541) (-0.514,1.572) (-0.506,1.606) (-0.499,1.638) (-0.493,1.671) (-0.488,1.703) (-0.483,1.736) (-0.478,1.768) (-0.473,1.800) (-0.469,1.833) (-0.465,1.866) (-0.462,1.898) (-0.459,1.930) (-0.456,1.963) (-0.453,1.996) (-0.450,2.028) (-0.448,2.061) (-0.445,2.093) (-0.443,2.126) (-0.441,2.158) (-0.439,2.191) (-0.438,2.223) (-0.436,2.256) (-0.434,2.288) (-0.433,2.321) (-0.432,2.353) (-0.430,2.385) (-0.429,2.417) (-0.428,2.451) (-0.427,2.483) (-0.426,2.516) (-0.425,2.548) (-0.424,2.581) (-0.423,2.600)}
\EQ{(-0.841,-2.600) (-0.843,-2.567) (-0.844,-2.535) (-0.846,-2.502) (-0.848,-2.470) (-0.849,-2.438) (-0.851,-2.405) (-0.853,-2.373) (-0.855,-2.340) (-0.857,-2.308) (-0.860,-2.275) (-0.862,-2.243) (-0.865,-2.212) (-0.867,-2.178) (-0.870,-2.145) (-0.873,-2.112) (-0.876,-2.080) (-0.879,-2.047) (-0.883,-2.015) (-0.887,-1.982) (-0.890,-1.950) (-0.894,-1.917) (-0.899,-1.885) (-0.903,-1.853) (-0.908,-1.820) (-0.913,-1.788) (-0.918,-1.755) (-0.924,-1.723) (-0.929,-1.690) (-0.936,-1.657) (-0.942,-1.625) (-0.949,-1.593) (-0.956,-1.563) (-0.964,-1.527) (-0.971,-1.495) (-0.980,-1.462) (-0.988,-1.431) (-0.997,-1.397) (-1.006,-1.365) (-1.016,-1.333) (-1.026,-1.300) (-1.036,-1.268) (-1.046,-1.236) (-1.058,-1.202) (-1.069,-1.170) (-1.079,-1.140) (-1.089,-1.111) (-1.101,-1.079) (-1.111,-1.049) (-1.124,-1.014) (-1.136,-0.982) (-1.148,-0.949) (-1.159,-0.916) (-1.170,-0.886) (-1.182,-0.851) (-1.194,-0.819) (-1.205,-0.786) (-1.215,-0.754) (-1.226,-0.722) (-1.235,-0.693) (-1.244,-0.663) (-1.254,-0.630) (-1.263,-0.598) (-1.272,-0.566) (-1.280,-0.533) (-1.288,-0.500) (-1.296,-0.468) (-1.303,-0.435) (-1.309,-0.403) (-1.315,-0.370) (-1.321,-0.338) (-1.326,-0.307) (-1.331,-0.273) (-1.335,-0.240) (-1.339,-0.208) (-1.342,-0.175) (-1.344,-0.143) (-1.346,-0.111) (-1.348,-0.078) (-1.349,-0.045) (-1.349,-0.013) (-1.349,0.020) (-1.349,0.052) (-1.348,0.085) (-1.346,0.117) (-1.344,0.150) (-1.341,0.182) (-1.338,0.215) (-1.334,0.247) (-1.330,0.280) (-1.325,0.312) (-1.320,0.345) (-1.314,0.377) (-1.308,0.409) (-1.301,0.442) (-1.294,0.475) (-1.287,0.505) (-1.278,0.540) (-1.270,0.572) (-1.261,0.605) (-1.252,0.637) (-1.242,0.670) (-1.232,0.702) (-1.222,0.734) (-1.211,0.767) (-1.200,0.800) (-1.189,0.831) (-1.178,0.865) (-1.166,0.897) (-1.155,0.930) (-1.144,0.959) (-1.133,0.988) (-1.122,1.021) (-1.110,1.053) (-1.099,1.085) (-1.087,1.118) (-1.075,1.151) (-1.064,1.183) (-1.053,1.216) (-1.042,1.248) (-1.032,1.281) (-1.022,1.313) (-1.012,1.346) (-1.003,1.378) (-0.993,1.411) (-0.985,1.443) (-0.976,1.476) (-0.968,1.507) (-0.961,1.541) (-0.953,1.573) (-0.946,1.606) (-0.940,1.638) (-0.933,1.671) (-0.927,1.703) (-0.921,1.736) (-0.916,1.765) (-0.911,1.800) (-0.906,1.833) (-0.901,1.866) (-0.897,1.898) (-0.893,1.930) (-0.889,1.963) (-0.885,1.996) (-0.882,2.028) (-0.878,2.061) (-0.875,2.093) (-0.872,2.126) (-0.869,2.158) (-0.866,2.191) (-0.864,2.223) (-0.861,2.256) (-0.859,2.288) (-0.857,2.321) (-0.854,2.353) (-0.852,2.385) (-0.851,2.418) (-0.849,2.450) (-0.847,2.483) (-0.845,2.516) (-0.844,2.548) (-0.842,2.581) (-0.841,2.600)}
\EQ{(-1.252,-2.600) (-1.254,-2.567) (-1.255,-2.535) (-1.257,-2.502) (-1.259,-2.470) (-1.261,-2.438) (-1.263,-2.405) (-1.265,-2.373) (-1.267,-2.340) (-1.269,-2.308) (-1.272,-2.275) (-1.274,-2.243) (-1.277,-2.210) (-1.279,-2.178) (-1.282,-2.145) (-1.285,-2.112) (-1.288,-2.080) (-1.291,-2.047) (-1.294,-2.015) (-1.298,-1.983) (-1.301,-1.950) (-1.305,-1.917) (-1.308,-1.885) (-1.312,-1.853) (-1.316,-1.820) (-1.320,-1.788) (-1.325,-1.755) (-1.329,-1.723) (-1.334,-1.690) (-1.338,-1.657) (-1.343,-1.625) (-1.348,-1.593) (-1.354,-1.560) (-1.359,-1.528) (-1.365,-1.495) (-1.370,-1.462) (-1.376,-1.430) (-1.382,-1.397) (-1.388,-1.365) (-1.394,-1.333) (-1.401,-1.300) (-1.407,-1.268) (-1.414,-1.235) (-1.420,-1.202) (-1.427,-1.170) (-1.434,-1.138) (-1.441,-1.105) (-1.448,-1.073) (-1.454,-1.040) (-1.461,-1.007) (-1.468,-0.975) (-1.475,-0.942) (-1.482,-0.910) (-1.488,-0.878) (-1.495,-0.846) (-1.502,-0.813) (-1.508,-0.781) (-1.514,-0.747) (-1.521,-0.715) (-1.527,-0.682) (-1.532,-0.650) (-1.538,-0.617) (-1.543,-0.585) (-1.548,-0.552) (-1.553,-0.520) (-1.558,-0.487) (-1.563,-0.455) (-1.567,-0.424) (-1.571,-0.390) (-1.574,-0.357) (-1.578,-0.325) (-1.581,-0.292) (-1.583,-0.260) (-1.586,-0.228) (-1.588,-0.195) (-1.590,-0.163) (-1.591,-0.130) (-1.592,-0.097) (-1.593,-0.065) (-1.594,-0.032) (-1.594,0.000) (-1.594,0.033) (-1.593,0.065) (-1.592,0.098) (-1.591,0.130) (-1.590,0.163) (-1.588,0.195) (-1.586,0.228) (-1.583,0.260) (-1.581,0.292) (-1.578,0.325) (-1.574,0.358) (-1.571,0.390) (-1.567,0.423) (-1.563,0.455) (-1.558,0.488) (-1.554,0.519) (-1.548,0.553) (-1.543,0.585) (-1.538,0.618) (-1.532,0.650) (-1.527,0.683) (-1.521,0.712) (-1.514,0.747) (-1.508,0.780) (-1.502,0.812) (-1.495,0.845) (-1.489,0.878) (-1.482,0.909) (-1.476,0.940) (-1.469,0.971) (-1.463,1.001) (-1.456,1.033) (-1.450,1.063) (-1.443,1.094) (-1.437,1.125) (-1.430,1.156) (-1.423,1.187) (-1.417,1.219) (-1.411,1.251) (-1.405,1.281) (-1.398,1.313) (-1.392,1.346) (-1.386,1.378) (-1.380,1.411) (-1.374,1.443) (-1.368,1.476) (-1.362,1.508) (-1.357,1.541) (-1.352,1.573) (-1.346,1.606) (-1.341,1.638) (-1.336,1.671) (-1.332,1.703) (-1.327,1.736) (-1.323,1.768) (-1.319,1.800) (-1.315,1.833) (-1.311,1.866) (-1.307,1.898) (-1.303,1.930) (-1.300,1.963) (-1.296,1.996) (-1.293,2.028) (-1.290,2.061) (-1.287,2.090) (-1.284,2.126) (-1.281,2.158) (-1.278,2.191) (-1.276,2.223) (-1.273,2.256) (-1.271,2.288) (-1.269,2.321) (-1.266,2.353) (-1.264,2.385) (-1.262,2.418) (-1.260,2.451) (-1.258,2.483) (-1.257,2.516) (-1.255,2.548) (-1.253,2.581) (-1.252,2.600)}
\EQ{(-1.657,-2.600) (-1.658,-2.567) (-1.660,-2.535) (-1.661,-2.502) (-1.663,-2.470) (-1.665,-2.438) (-1.667,-2.405) (-1.668,-2.373) (-1.670,-2.340) (-1.672,-2.308) (-1.674,-2.275) (-1.676,-2.243) (-1.679,-2.210) (-1.681,-2.178) (-1.683,-2.145) (-1.685,-2.112) (-1.688,-2.080) (-1.690,-2.047) (-1.693,-2.015) (-1.696,-1.982) (-1.698,-1.950) (-1.701,-1.917) (-1.704,-1.885) (-1.707,-1.853) (-1.710,-1.820) (-1.713,-1.788) (-1.716,-1.755) (-1.719,-1.723) (-1.723,-1.691) (-1.726,-1.657) (-1.729,-1.625) (-1.733,-1.592) (-1.737,-1.560) (-1.740,-1.528) (-1.744,-1.495) (-1.748,-1.462) (-1.752,-1.430) (-1.755,-1.397) (-1.759,-1.365) (-1.763,-1.333) (-1.767,-1.300) (-1.771,-1.268) (-1.776,-1.235) (-1.780,-1.202) (-1.784,-1.170) (-1.788,-1.137) (-1.792,-1.105) (-1.796,-1.073) (-1.800,-1.040) (-1.804,-1.007) (-1.809,-0.975) (-1.813,-0.942) (-1.817,-0.910) (-1.821,-0.877) (-1.824,-0.845) (-1.828,-0.812) (-1.832,-0.780) (-1.836,-0.747) (-1.839,-0.715) (-1.843,-0.682) (-1.846,-0.652) (-1.849,-0.617) (-1.853,-0.586) (-1.856,-0.552) (-1.859,-0.520) (-1.861,-0.487) (-1.864,-0.455) (-1.866,-0.422) (-1.869,-0.390) (-1.871,-0.357) (-1.873,-0.325) (-1.875,-0.292) (-1.876,-0.260) (-1.878,-0.228) (-1.879,-0.195) (-1.880,-0.163) (-1.881,-0.130) (-1.881,-0.097) (-1.882,-0.065) (-1.882,-0.032) (-1.882,0.000) (-1.882,0.033) (-1.882,0.065) (-1.881,0.098) (-1.881,0.130) (-1.880,0.163) (-1.879,0.195) (-1.878,0.228) (-1.876,0.260) (-1.875,0.292) (-1.873,0.325) (-1.871,0.358) (-1.869,0.390) (-1.866,0.423) (-1.864,0.455) (-1.861,0.488) (-1.859,0.520) (-1.856,0.553) (-1.853,0.585) (-1.849,0.618) (-1.846,0.650) (-1.843,0.683) (-1.840,0.713) (-1.836,0.748) (-1.832,0.780) (-1.828,0.812) (-1.824,0.845) (-1.821,0.878) (-1.817,0.910) (-1.813,0.943) (-1.809,0.975) (-1.804,1.008) (-1.801,1.039) (-1.796,1.073) (-1.792,1.105) (-1.788,1.138) (-1.784,1.170) (-1.780,1.203) (-1.776,1.235) (-1.771,1.268) (-1.767,1.300) (-1.763,1.333) (-1.759,1.365) (-1.755,1.398) (-1.752,1.430) (-1.748,1.462) (-1.744,1.495) (-1.740,1.528) (-1.737,1.560) (-1.733,1.593) (-1.729,1.625) (-1.726,1.658) (-1.723,1.690) (-1.719,1.723) (-1.716,1.755) (-1.713,1.788) (-1.710,1.820) (-1.707,1.853) (-1.704,1.885) (-1.701,1.917) (-1.698,1.950) (-1.696,1.983) (-1.693,2.015) (-1.690,2.047) (-1.688,2.080) (-1.685,2.113) (-1.683,2.145) (-1.681,2.178) (-1.679,2.210) (-1.676,2.243) (-1.674,2.275) (-1.672,2.308) (-1.670,2.340) (-1.668,2.373) (-1.667,2.405) (-1.665,2.438) (-1.663,2.470) (-1.661,2.502) (-1.660,2.535) (-1.658,2.568) (-1.657,2.600)}
\EQ{(-2.056,-2.600) (-2.058,-2.567) (-2.059,-2.535) (-2.060,-2.502) (-2.062,-2.470) (-2.063,-2.438) (-2.065,-2.405) (-2.066,-2.372) (-2.068,-2.340) (-2.069,-2.308) (-2.071,-2.275) (-2.073,-2.243) (-2.074,-2.210) (-2.076,-2.178) (-2.078,-2.145) (-2.080,-2.112) (-2.082,-2.080) (-2.084,-2.047) (-2.086,-2.015) (-2.087,-1.982) (-2.090,-1.950) (-2.092,-1.917) (-2.094,-1.885) (-2.096,-1.853) (-2.098,-1.820) (-2.100,-1.788) (-2.102,-1.755) (-2.105,-1.723) (-2.107,-1.690) (-2.109,-1.657) (-2.112,-1.625) (-2.114,-1.593) (-2.116,-1.560) (-2.119,-1.528) (-2.121,-1.495) (-2.124,-1.462) (-2.126,-1.430) (-2.129,-1.397) (-2.131,-1.365) (-2.134,-1.333) (-2.137,-1.300) (-2.139,-1.268) (-2.142,-1.235) (-2.144,-1.202) (-2.147,-1.170) (-2.149,-1.137) (-2.152,-1.105) (-2.155,-1.073) (-2.157,-1.040) (-2.160,-1.007) (-2.162,-0.975) (-2.164,-0.943) (-2.167,-0.910) (-2.169,-0.878) (-2.172,-0.845) (-2.174,-0.812) (-2.176,-0.780) (-2.178,-0.747) (-2.180,-0.715) (-2.183,-0.682) (-2.185,-0.650) (-2.186,-0.617) (-2.188,-0.585) (-2.190,-0.552) (-2.192,-0.520) (-2.193,-0.487) (-2.195,-0.455) (-2.196,-0.422) (-2.198,-0.390) (-2.199,-0.357) (-2.200,-0.325) (-2.201,-0.292) (-2.202,-0.260) (-2.203,-0.228) (-2.204,-0.197) (-2.204,-0.163) (-2.205,-0.130) (-2.205,-0.097) (-2.205,-0.065) (-2.206,-0.032) (-2.206,0.000) (-2.206,0.033) (-2.205,0.065) (-2.205,0.098) (-2.205,0.130) (-2.204,0.163) (-2.204,0.195) (-2.203,0.228) (-2.202,0.260) (-2.201,0.292) (-2.200,0.325) (-2.199,0.358) (-2.198,0.390) (-2.196,0.423) (-2.195,0.455) (-2.193,0.488) (-2.192,0.520) (-2.190,0.553) (-2.188,0.585) (-2.186,0.618) (-2.185,0.650) (-2.183,0.683) (-2.180,0.715) (-2.178,0.748) (-2.176,0.780) (-2.174,0.812) (-2.172,0.845) (-2.169,0.878) (-2.167,0.910) (-2.165,0.943) (-2.162,0.975) (-2.160,1.008) (-2.157,1.040) (-2.155,1.073) (-2.152,1.105) (-2.149,1.138) (-2.147,1.170) (-2.144,1.203) (-2.142,1.235) (-2.139,1.268) (-2.137,1.300) (-2.134,1.333) (-2.131,1.365) (-2.129,1.398) (-2.126,1.430) (-2.124,1.462) (-2.121,1.495) (-2.119,1.526) (-2.116,1.560) (-2.114,1.593) (-2.112,1.625) (-2.109,1.658) (-2.107,1.690) (-2.105,1.723) (-2.102,1.755) (-2.100,1.788) (-2.098,1.820) (-2.096,1.853) (-2.094,1.885) (-2.092,1.917) (-2.090,1.950) (-2.087,1.983) (-2.086,2.015) (-2.084,2.047) (-2.082,2.080) (-2.080,2.113) (-2.078,2.145) (-2.076,2.178) (-2.074,2.210) (-2.073,2.243) (-2.071,2.275) (-2.069,2.308) (-2.068,2.340) (-2.066,2.372) (-2.065,2.405) (-2.063,2.438) (-2.062,2.470) (-2.061,2.502) (-2.059,2.535) (-2.058,2.568) (-2.056,2.600)}
\EQ{(-2.454,-2.600) (-2.455,-2.567) (-2.456,-2.535) (-2.457,-2.502) (-2.458,-2.470) (-2.459,-2.438) (-2.460,-2.405) (-2.462,-2.373) (-2.463,-2.340) (-2.464,-2.308) (-2.465,-2.275) (-2.467,-2.243) (-2.468,-2.210) (-2.469,-2.178) (-2.471,-2.145) (-2.472,-2.112) (-2.473,-2.080) (-2.475,-2.047) (-2.476,-2.015) (-2.478,-1.982) (-2.479,-1.950) (-2.480,-1.917) (-2.482,-1.885) (-2.484,-1.853) (-2.485,-1.820) (-2.487,-1.788) (-2.488,-1.755) (-2.490,-1.723) (-2.491,-1.690) (-2.493,-1.657) (-2.495,-1.625) (-2.496,-1.593) (-2.498,-1.560) (-2.499,-1.527) (-2.501,-1.495) (-2.503,-1.462) (-2.504,-1.430) (-2.506,-1.397) (-2.508,-1.365) (-2.509,-1.333) (-2.511,-1.300) (-2.513,-1.268) (-2.514,-1.235) (-2.516,-1.202) (-2.518,-1.170) (-2.519,-1.137) (-2.521,-1.105) (-2.522,-1.073) (-2.524,-1.040) (-2.526,-1.007) (-2.527,-0.975) (-2.529,-0.942) (-2.530,-0.910) (-2.532,-0.877) (-2.533,-0.845) (-2.534,-0.812) (-2.536,-0.780) (-2.537,-0.747) (-2.538,-0.715) (-2.540,-0.682) (-2.541,-0.650) (-2.542,-0.617) (-2.543,-0.585) (-2.544,-0.552) (-2.545,-0.520) (-2.546,-0.487) (-2.547,-0.455) (-2.548,-0.422) (-2.549,-0.390) (-2.549,-0.357) (-2.550,-0.325) (-2.551,-0.292) (-2.551,-0.260) (-2.552,-0.228) (-2.552,-0.195) (-2.553,-0.163) (-2.553,-0.130) (-2.553,-0.097) (-2.553,-0.065) (-2.553,-0.032) (-2.553,0.000) (-2.553,0.033) (-2.553,0.065) (-2.553,0.098) (-2.553,0.130) (-2.553,0.163) (-2.552,0.195) (-2.552,0.228) (-2.551,0.260) (-2.551,0.292) (-2.550,0.325) (-2.549,0.358) (-2.549,0.390) (-2.548,0.423) (-2.547,0.455) (-2.546,0.488) (-2.545,0.520) (-2.544,0.553) (-2.543,0.585) (-2.542,0.618) (-2.541,0.650) (-2.540,0.683) (-2.538,0.715) (-2.537,0.748) (-2.536,0.780) (-2.534,0.812) (-2.533,0.845) (-2.532,0.878) (-2.530,0.910) (-2.529,0.943) (-2.527,0.975) (-2.526,1.008) (-2.524,1.040) (-2.522,1.073) (-2.521,1.105) (-2.519,1.138) (-2.518,1.170) (-2.516,1.203) (-2.514,1.235) (-2.513,1.268) (-2.511,1.300) (-2.509,1.333) (-2.508,1.365) (-2.506,1.398) (-2.504,1.430) (-2.503,1.462) (-2.501,1.495) (-2.499,1.528) (-2.498,1.560) (-2.496,1.593) (-2.495,1.625) (-2.493,1.658) (-2.491,1.690) (-2.490,1.723) (-2.488,1.755) (-2.487,1.788) (-2.485,1.820) (-2.484,1.853) (-2.482,1.885) (-2.480,1.917) (-2.479,1.950) (-2.478,1.983) (-2.476,2.015) (-2.475,2.047) (-2.473,2.080) (-2.472,2.113) (-2.471,2.145) (-2.469,2.178) (-2.468,2.210) (-2.467,2.243) (-2.465,2.275) (-2.464,2.308) (-2.463,2.340) (-2.462,2.373) (-2.460,2.405) (-2.459,2.438) (-2.458,2.470) (-2.457,2.501) (-2.456,2.535) (-2.455,2.568) (-2.454,2.600)}

      % Esfera en campo uniforme - líneas de campo
% Generado por scripts/generar_datos.py (no editar a mano).
\LC{(1.000,-0.105) (1.034,-0.108) (1.066,-0.111) (1.099,-0.114) (1.129,-0.117) (1.164,-0.120) (1.196,-0.123) (1.229,-0.125) (1.261,-0.128) (1.294,-0.130) (1.326,-0.132) (1.359,-0.135) (1.391,-0.137) (1.424,-0.139) (1.456,-0.141) (1.489,-0.142) (1.521,-0.144) (1.554,-0.146) (1.586,-0.147) (1.619,-0.149) (1.651,-0.150) (1.683,-0.151) (1.716,-0.153) (1.749,-0.154) (1.781,-0.155) (1.813,-0.156) (1.846,-0.157) (1.879,-0.158) (1.911,-0.159) (1.944,-0.160) (1.976,-0.161) (2.009,-0.161) (2.041,-0.162) (2.074,-0.163) (2.106,-0.163) (2.139,-0.164) (2.171,-0.165) (2.204,-0.165) (2.236,-0.166) (2.268,-0.166) (2.301,-0.167) (2.334,-0.167) (2.366,-0.168) (2.399,-0.168) (2.431,-0.169) (2.464,-0.169) (2.496,-0.169) (2.529,-0.170) (2.561,-0.170) (2.594,-0.171) (2.600,-0.171)}
\LC{(-2.600,-0.171) (-2.567,-0.170) (-2.535,-0.170) (-2.502,-0.170) (-2.470,-0.169) (-2.438,-0.169) (-2.405,-0.168) (-2.373,-0.168) (-2.340,-0.167) (-2.308,-0.167) (-2.275,-0.166) (-2.242,-0.166) (-2.210,-0.165) (-2.178,-0.165) (-2.145,-0.164) (-2.112,-0.164) (-2.080,-0.163) (-2.047,-0.162) (-2.015,-0.161) (-1.982,-0.161) (-1.950,-0.160) (-1.917,-0.159) (-1.885,-0.158) (-1.853,-0.157) (-1.820,-0.156) (-1.788,-0.155) (-1.755,-0.154) (-1.723,-0.153) (-1.690,-0.152) (-1.657,-0.150) (-1.625,-0.149) (-1.593,-0.147) (-1.560,-0.146) (-1.527,-0.144) (-1.495,-0.143) (-1.462,-0.141) (-1.430,-0.139) (-1.397,-0.137) (-1.365,-0.135) (-1.333,-0.133) (-1.300,-0.131) (-1.268,-0.128) (-1.235,-0.126) (-1.202,-0.123) (-1.170,-0.121) (-1.137,-0.118) (-1.105,-0.115) (-1.073,-0.112) (-1.040,-0.109) (-1.007,-0.105) (-1.000,-0.105)}
\LC{(1.000,0.105) (1.034,0.108) (1.066,0.111) (1.099,0.114) (1.129,0.117) (1.164,0.120) (1.196,0.123) (1.229,0.125) (1.261,0.128) (1.294,0.130) (1.326,0.132) (1.359,0.135) (1.391,0.137) (1.424,0.139) (1.456,0.141) (1.489,0.142) (1.521,0.144) (1.554,0.146) (1.586,0.147) (1.619,0.149) (1.651,0.150) (1.683,0.151) (1.716,0.153) (1.749,0.154) (1.781,0.155) (1.814,0.156) (1.846,0.157) (1.879,0.158) (1.911,0.159) (1.944,0.160) (1.976,0.161) (2.009,0.161) (2.041,0.162) (2.074,0.163) (2.106,0.163) (2.139,0.164) (2.171,0.165) (2.204,0.165) (2.236,0.166) (2.268,0.166) (2.301,0.167) (2.334,0.167) (2.366,0.168) (2.399,0.168) (2.431,0.169) (2.464,0.169) (2.496,0.169) (2.529,0.170) (2.561,0.170) (2.594,0.171) (2.600,0.171)}
\LC{(-2.600,0.171) (-2.567,0.170) (-2.535,0.170) (-2.502,0.170) (-2.470,0.169) (-2.438,0.169) (-2.405,0.168) (-2.373,0.168) (-2.340,0.167) (-2.308,0.167) (-2.275,0.166) (-2.242,0.166) (-2.210,0.165) (-2.178,0.165) (-2.145,0.164) (-2.112,0.164) (-2.080,0.163) (-2.047,0.162) (-2.015,0.161) (-1.982,0.161) (-1.950,0.160) (-1.917,0.159) (-1.885,0.158) (-1.853,0.157) (-1.820,0.156) (-1.788,0.155) (-1.755,0.154) (-1.723,0.153) (-1.690,0.152) (-1.657,0.150) (-1.625,0.149) (-1.593,0.147) (-1.560,0.146) (-1.528,0.144) (-1.495,0.143) (-1.462,0.141) (-1.430,0.139) (-1.397,0.137) (-1.365,0.135) (-1.333,0.133) (-1.300,0.131) (-1.268,0.128) (-1.235,0.126) (-1.202,0.123) (-1.170,0.121) (-1.137,0.118) (-1.105,0.115) (-1.073,0.112) (-1.040,0.109) (-1.007,0.105) (-1.000,0.105)}
\LC{(0.964,-0.278) (0.995,-0.287) (1.027,-0.295) (1.060,-0.304) (1.092,-0.312) (1.125,-0.320) (1.157,-0.327) (1.190,-0.334) (1.222,-0.341) (1.255,-0.348) (1.287,-0.354) (1.320,-0.360) (1.352,-0.365) (1.385,-0.370) (1.417,-0.375) (1.450,-0.380) (1.482,-0.385) (1.515,-0.389) (1.547,-0.393) (1.579,-0.397) (1.612,-0.400) (1.645,-0.404) (1.677,-0.407) (1.710,-0.410) (1.742,-0.413) (1.773,-0.416) (1.807,-0.419) (1.840,-0.421) (1.872,-0.424) (1.905,-0.426) (1.937,-0.428) (1.970,-0.430) (2.002,-0.432) (2.034,-0.434) (2.067,-0.436) (2.100,-0.438) (2.132,-0.439) (2.165,-0.441) (2.197,-0.442) (2.230,-0.444) (2.262,-0.445) (2.295,-0.446) (2.327,-0.448) (2.360,-0.449) (2.392,-0.450) (2.425,-0.451) (2.457,-0.452) (2.490,-0.453) (2.522,-0.454) (2.555,-0.455) (2.587,-0.456) (2.600,-0.456)}
\LC{(-2.600,-0.456) (-2.567,-0.455) (-2.535,-0.454) (-2.502,-0.453) (-2.470,-0.452) (-2.438,-0.451) (-2.405,-0.450) (-2.373,-0.449) (-2.340,-0.448) (-2.308,-0.447) (-2.275,-0.446) (-2.243,-0.444) (-2.210,-0.443) (-2.178,-0.441) (-2.145,-0.440) (-2.112,-0.438) (-2.080,-0.437) (-2.047,-0.435) (-2.015,-0.433) (-1.982,-0.431) (-1.950,-0.429) (-1.917,-0.427) (-1.885,-0.425) (-1.854,-0.422) (-1.820,-0.420) (-1.788,-0.417) (-1.755,-0.414) (-1.723,-0.412) (-1.690,-0.409) (-1.657,-0.405) (-1.625,-0.402) (-1.593,-0.398) (-1.560,-0.395) (-1.528,-0.391) (-1.495,-0.386) (-1.462,-0.382) (-1.430,-0.377) (-1.397,-0.372) (-1.365,-0.367) (-1.333,-0.362) (-1.300,-0.356) (-1.267,-0.350) (-1.235,-0.344) (-1.202,-0.337) (-1.170,-0.330) (-1.137,-0.323) (-1.105,-0.315) (-1.073,-0.307) (-1.040,-0.299) (-1.007,-0.290) (-0.975,-0.281) (-0.964,-0.278)}
\LC{(0.964,0.278) (0.995,0.287) (1.027,0.295) (1.060,0.304) (1.092,0.312) (1.125,0.320) (1.157,0.327) (1.190,0.334) (1.222,0.341) (1.255,0.348) (1.287,0.354) (1.320,0.360) (1.352,0.365) (1.385,0.370) (1.417,0.375) (1.450,0.380) (1.482,0.385) (1.515,0.389) (1.547,0.393) (1.579,0.397) (1.612,0.400) (1.645,0.404) (1.677,0.407) (1.710,0.410) (1.742,0.413) (1.773,0.416) (1.807,0.419) (1.840,0.421) (1.872,0.424) (1.905,0.426) (1.937,0.428) (1.970,0.430) (2.002,0.432) (2.034,0.434) (2.067,0.436) (2.100,0.438) (2.132,0.439) (2.165,0.441) (2.197,0.442) (2.230,0.444) (2.262,0.445) (2.295,0.446) (2.327,0.448) (2.360,0.449) (2.392,0.450) (2.425,0.451) (2.457,0.452) (2.490,0.453) (2.522,0.454) (2.555,0.455) (2.587,0.456) (2.600,0.456)}
\LC{(-2.600,0.456) (-2.567,0.455) (-2.535,0.454) (-2.502,0.453) (-2.470,0.452) (-2.438,0.451) (-2.405,0.450) (-2.373,0.449) (-2.340,0.448) (-2.308,0.447) (-2.275,0.446) (-2.243,0.444) (-2.210,0.443) (-2.178,0.441) (-2.145,0.440) (-2.112,0.438) (-2.080,0.437) (-2.047,0.435) (-2.015,0.433) (-1.982,0.431) (-1.950,0.429) (-1.917,0.427) (-1.885,0.425) (-1.854,0.423) (-1.820,0.420) (-1.788,0.417) (-1.755,0.414) (-1.723,0.412) (-1.690,0.409) (-1.657,0.405) (-1.625,0.402) (-1.593,0.398) (-1.560,0.395) (-1.528,0.391) (-1.495,0.386) (-1.462,0.382) (-1.430,0.377) (-1.397,0.372) (-1.365,0.367) (-1.333,0.362) (-1.300,0.356) (-1.267,0.350) (-1.235,0.344) (-1.202,0.337) (-1.170,0.330) (-1.137,0.323) (-1.105,0.315) (-1.073,0.307) (-1.040,0.299) (-1.007,0.290) (-0.975,0.281) (-0.964,0.278)}
\LC{(0.894,-0.451) (0.923,-0.465) (0.956,-0.481) (0.985,-0.494) (1.014,-0.506) (1.047,-0.520) (1.079,-0.532) (1.112,-0.544) (1.144,-0.555) (1.175,-0.566) (1.209,-0.576) (1.239,-0.585) (1.274,-0.595) (1.307,-0.603) (1.337,-0.611) (1.372,-0.619) (1.404,-0.626) (1.437,-0.633) (1.469,-0.640) (1.502,-0.646) (1.534,-0.652) (1.567,-0.657) (1.599,-0.663) (1.632,-0.668) (1.664,-0.672) (1.696,-0.677) (1.729,-0.681) (1.762,-0.685) (1.794,-0.689) (1.827,-0.693) (1.859,-0.696) (1.892,-0.699) (1.924,-0.703) (1.957,-0.706) (1.989,-0.708) (2.022,-0.711) (2.054,-0.714) (2.087,-0.716) (2.119,-0.718) (2.151,-0.721) (2.184,-0.723) (2.216,-0.725) (2.249,-0.727) (2.282,-0.729) (2.314,-0.731) (2.346,-0.732) (2.379,-0.734) (2.412,-0.735) (2.444,-0.737) (2.477,-0.738) (2.509,-0.740) (2.540,-0.741) (2.574,-0.742) (2.600,-0.743)}
\LC{(-2.600,-0.743) (-2.567,-0.742) (-2.535,-0.741) (-2.502,-0.739) (-2.470,-0.738) (-2.438,-0.737) (-2.405,-0.735) (-2.373,-0.734) (-2.340,-0.732) (-2.308,-0.730) (-2.275,-0.728) (-2.243,-0.726) (-2.210,-0.725) (-2.178,-0.722) (-2.145,-0.720) (-2.112,-0.718) (-2.080,-0.716) (-2.047,-0.713) (-2.015,-0.711) (-1.982,-0.708) (-1.950,-0.705) (-1.919,-0.702) (-1.885,-0.699) (-1.854,-0.696) (-1.820,-0.692) (-1.788,-0.688) (-1.755,-0.684) (-1.723,-0.680) (-1.691,-0.676) (-1.657,-0.671) (-1.625,-0.667) (-1.593,-0.662) (-1.561,-0.656) (-1.528,-0.651) (-1.495,-0.645) (-1.462,-0.639) (-1.430,-0.632) (-1.397,-0.625) (-1.365,-0.618) (-1.333,-0.610) (-1.300,-0.602) (-1.268,-0.593) (-1.235,-0.584) (-1.202,-0.574) (-1.170,-0.564) (-1.137,-0.553) (-1.105,-0.542) (-1.073,-0.530) (-1.040,-0.517) (-1.007,-0.504) (-0.975,-0.489) (-0.943,-0.474) (-0.915,-0.462) (-0.894,-0.451)}
\LC{(0.894,0.451) (0.923,0.465) (0.956,0.481) (0.985,0.494) (1.014,0.506) (1.047,0.520) (1.079,0.532) (1.112,0.544) (1.144,0.555) (1.175,0.566) (1.209,0.576) (1.239,0.585) (1.274,0.595) (1.307,0.603) (1.337,0.611) (1.372,0.619) (1.404,0.626) (1.437,0.633) (1.469,0.640) (1.502,0.646) (1.534,0.652) (1.567,0.657) (1.599,0.663) (1.632,0.668) (1.664,0.672) (1.696,0.677) (1.729,0.681) (1.762,0.685) (1.794,0.689) (1.827,0.693) (1.859,0.696) (1.892,0.699) (1.924,0.703) (1.957,0.706) (1.989,0.708) (2.022,0.711) (2.054,0.714) (2.087,0.716) (2.119,0.718) (2.151,0.721) (2.184,0.723) (2.217,0.725) (2.249,0.727) (2.282,0.729) (2.314,0.731) (2.347,0.732) (2.379,0.734) (2.412,0.735) (2.444,0.737) (2.476,0.738) (2.509,0.740) (2.540,0.741) (2.574,0.742) (2.600,0.743)}
\LC{(-2.600,0.743) (-2.567,0.742) (-2.535,0.741) (-2.503,0.739) (-2.470,0.738) (-2.438,0.737) (-2.405,0.735) (-2.373,0.734) (-2.340,0.732) (-2.308,0.730) (-2.275,0.728) (-2.243,0.726) (-2.210,0.725) (-2.178,0.722) (-2.145,0.720) (-2.112,0.718) (-2.080,0.716) (-2.047,0.713) (-2.015,0.711) (-1.983,0.708) (-1.950,0.705) (-1.919,0.702) (-1.885,0.699) (-1.854,0.696) (-1.820,0.692) (-1.788,0.688) (-1.755,0.684) (-1.723,0.680) (-1.691,0.676) (-1.658,0.671) (-1.625,0.667) (-1.592,0.662) (-1.561,0.657) (-1.528,0.651) (-1.495,0.645) (-1.462,0.639) (-1.430,0.632) (-1.397,0.625) (-1.365,0.618) (-1.333,0.610) (-1.300,0.602) (-1.268,0.593) (-1.235,0.584) (-1.202,0.574) (-1.170,0.564) (-1.137,0.553) (-1.105,0.542) (-1.073,0.530) (-1.040,0.517) (-1.007,0.504) (-0.975,0.489) (-0.943,0.475) (-0.915,0.462) (-0.894,0.451)}
\LC{(0.785,-0.626) (0.812,-0.647) (0.839,-0.666) (0.865,-0.684) (0.891,-0.701) (0.917,-0.717) (0.943,-0.732) (0.971,-0.747) (1.001,-0.763) (1.034,-0.779) (1.065,-0.793) (1.095,-0.806) (1.125,-0.818) (1.157,-0.830) (1.190,-0.842) (1.219,-0.851) (1.248,-0.861) (1.281,-0.871) (1.313,-0.880) (1.346,-0.889) (1.378,-0.897) (1.411,-0.905) (1.443,-0.912) (1.476,-0.919) (1.508,-0.926) (1.541,-0.932) (1.573,-0.938) (1.606,-0.944) (1.637,-0.949) (1.671,-0.954) (1.703,-0.959) (1.736,-0.964) (1.768,-0.968) (1.800,-0.972) (1.833,-0.976) (1.866,-0.980) (1.898,-0.984) (1.930,-0.987) (1.963,-0.991) (1.996,-0.994) (2.028,-0.997) (2.061,-1.000) (2.093,-1.003) (2.126,-1.005) (2.158,-1.008) (2.191,-1.010) (2.223,-1.013) (2.256,-1.015) (2.288,-1.017) (2.321,-1.019) (2.353,-1.021) (2.385,-1.023) (2.418,-1.025) (2.451,-1.027) (2.483,-1.028) (2.516,-1.030) (2.548,-1.031) (2.581,-1.033) (2.600,-1.034)}
\LC{(-2.600,-1.034) (-2.567,-1.032) (-2.535,-1.031) (-2.502,-1.029) (-2.470,-1.028) (-2.438,-1.026) (-2.405,-1.024) (-2.373,-1.022) (-2.342,-1.020) (-2.308,-1.018) (-2.275,-1.016) (-2.243,-1.014) (-2.210,-1.012) (-2.178,-1.009) (-2.145,-1.007) (-2.112,-1.004) (-2.080,-1.002) (-2.047,-0.999) (-2.015,-0.996) (-1.982,-0.993) (-1.950,-0.989) (-1.917,-0.986) (-1.885,-0.983) (-1.853,-0.979) (-1.821,-0.975) (-1.788,-0.971) (-1.755,-0.966) (-1.723,-0.962) (-1.690,-0.957) (-1.657,-0.952) (-1.625,-0.947) (-1.593,-0.941) (-1.562,-0.936) (-1.528,-0.930) (-1.495,-0.923) (-1.463,-0.916) (-1.434,-0.910) (-1.404,-0.903) (-1.371,-0.895) (-1.339,-0.887) (-1.306,-0.878) (-1.274,-0.869) (-1.242,-0.859) (-1.209,-0.848) (-1.180,-0.839) (-1.150,-0.828) (-1.118,-0.815) (-1.085,-0.802) (-1.053,-0.788) (-1.022,-0.773) (-0.994,-0.760) (-0.962,-0.743) (-0.935,-0.728) (-0.903,-0.709) (-0.877,-0.693) (-0.851,-0.675) (-0.825,-0.657) (-0.799,-0.637) (-0.785,-0.626)}
\LC{(0.785,0.626) (0.812,0.647) (0.839,0.666) (0.865,0.684) (0.891,0.701) (0.917,0.717) (0.943,0.732) (0.971,0.748) (1.001,0.763) (1.034,0.779) (1.065,0.793) (1.095,0.806) (1.125,0.818) (1.157,0.830) (1.190,0.842) (1.219,0.852) (1.248,0.861) (1.281,0.871) (1.313,0.880) (1.346,0.889) (1.378,0.897) (1.411,0.905) (1.443,0.912) (1.476,0.919) (1.508,0.926) (1.541,0.932) (1.573,0.938) (1.606,0.944) (1.637,0.949) (1.671,0.954) (1.703,0.959) (1.736,0.964) (1.768,0.968) (1.800,0.972) (1.833,0.976) (1.866,0.980) (1.898,0.984) (1.930,0.987) (1.963,0.991) (1.996,0.994) (2.028,0.997) (2.061,1.000) (2.093,1.003) (2.126,1.005) (2.158,1.008) (2.191,1.010) (2.223,1.013) (2.256,1.015) (2.288,1.017) (2.321,1.019) (2.353,1.021) (2.385,1.023) (2.418,1.025) (2.451,1.027) (2.483,1.028) (2.516,1.030) (2.548,1.031) (2.581,1.033) (2.600,1.034)}
\LC{(-2.600,1.034) (-2.568,1.032) (-2.535,1.031) (-2.502,1.029) (-2.470,1.028) (-2.438,1.026) (-2.405,1.024) (-2.373,1.022) (-2.342,1.021) (-2.308,1.018) (-2.275,1.016) (-2.243,1.014) (-2.210,1.012) (-2.178,1.009) (-2.145,1.007) (-2.112,1.004) (-2.080,1.002) (-2.047,0.999) (-2.015,0.996) (-1.982,0.993) (-1.950,0.989) (-1.917,0.986) (-1.885,0.983) (-1.853,0.979) (-1.821,0.975) (-1.788,0.971) (-1.755,0.966) (-1.723,0.962) (-1.690,0.957) (-1.657,0.952) (-1.625,0.947) (-1.593,0.941) (-1.562,0.936) (-1.528,0.930) (-1.495,0.923) (-1.463,0.917) (-1.434,0.910) (-1.404,0.903) (-1.371,0.895) (-1.339,0.887) (-1.306,0.878) (-1.274,0.869) (-1.242,0.859) (-1.209,0.848) (-1.180,0.839) (-1.150,0.828) (-1.118,0.815) (-1.085,0.802) (-1.053,0.788) (-1.022,0.774) (-0.994,0.760) (-0.962,0.743) (-0.935,0.728) (-0.903,0.709) (-0.877,0.693) (-0.852,0.675) (-0.825,0.657) (-0.799,0.637) (-0.785,0.626)}
\LC{(0.606,-0.798) (0.627,-0.825) (0.650,-0.851) (0.675,-0.877) (0.702,-0.903) (0.726,-0.923) (0.750,-0.942) (0.777,-0.962) (0.806,-0.982) (0.832,-0.998) (0.858,-1.013) (0.891,-1.031) (0.920,-1.046) (0.947,-1.059) (0.975,-1.072) (1.007,-1.085) (1.039,-1.099) (1.073,-1.111) (1.105,-1.122) (1.138,-1.133) (1.170,-1.143) (1.203,-1.153) (1.235,-1.162) (1.264,-1.170) (1.300,-1.179) (1.333,-1.187) (1.365,-1.194) (1.398,-1.201) (1.430,-1.208) (1.462,-1.214) (1.495,-1.221) (1.528,-1.226) (1.560,-1.232) (1.593,-1.237) (1.625,-1.242) (1.658,-1.247) (1.690,-1.252) (1.723,-1.256) (1.755,-1.261) (1.788,-1.265) (1.820,-1.269) (1.853,-1.272) (1.885,-1.276) (1.917,-1.279) (1.950,-1.283) (1.983,-1.286) (2.015,-1.289) (2.047,-1.292) (2.080,-1.295) (2.113,-1.297) (2.144,-1.300) (2.178,-1.303) (2.210,-1.305) (2.243,-1.307) (2.275,-1.309) (2.308,-1.312) (2.340,-1.314) (2.373,-1.316) (2.405,-1.318) (2.438,-1.319) (2.470,-1.321) (2.502,-1.323) (2.535,-1.325) (2.568,-1.326) (2.600,-1.328)}
\LC{(-2.600,-1.328) (-2.567,-1.326) (-2.535,-1.325) (-2.502,-1.323) (-2.470,-1.321) (-2.439,-1.319) (-2.405,-1.318) (-2.373,-1.316) (-2.340,-1.314) (-2.308,-1.312) (-2.275,-1.309) (-2.243,-1.307) (-2.210,-1.305) (-2.178,-1.303) (-2.145,-1.300) (-2.112,-1.297) (-2.080,-1.295) (-2.047,-1.292) (-2.015,-1.289) (-1.982,-1.286) (-1.950,-1.283) (-1.917,-1.279) (-1.885,-1.276) (-1.853,-1.272) (-1.820,-1.269) (-1.788,-1.265) (-1.757,-1.261) (-1.723,-1.256) (-1.690,-1.252) (-1.657,-1.247) (-1.625,-1.242) (-1.592,-1.237) (-1.560,-1.232) (-1.528,-1.226) (-1.495,-1.221) (-1.462,-1.214) (-1.430,-1.208) (-1.397,-1.201) (-1.365,-1.194) (-1.333,-1.187) (-1.300,-1.179) (-1.268,-1.171) (-1.235,-1.162) (-1.202,-1.153) (-1.172,-1.144) (-1.137,-1.133) (-1.105,-1.122) (-1.074,-1.111) (-1.040,-1.099) (-1.007,-1.086) (-0.976,-1.073) (-0.947,-1.059) (-0.920,-1.046) (-0.890,-1.031) (-0.859,-1.014) (-0.832,-0.998) (-0.806,-0.982) (-0.780,-0.964) (-0.754,-0.945) (-0.728,-0.925) (-0.703,-0.903) (-0.676,-0.878) (-0.650,-0.851) (-0.627,-0.825) (-0.606,-0.799) (-0.606,-0.798)}
\LC{(0.606,0.798) (0.627,0.826) (0.650,0.851) (0.675,0.878) (0.702,0.903) (0.726,0.923) (0.750,0.943) (0.777,0.962) (0.806,0.982) (0.832,0.998) (0.858,1.013) (0.891,1.031) (0.920,1.047) (0.947,1.060) (0.975,1.072) (1.007,1.086) (1.039,1.099) (1.073,1.111) (1.105,1.122) (1.138,1.133) (1.170,1.143) (1.203,1.153) (1.235,1.162) (1.264,1.170) (1.300,1.179) (1.333,1.187) (1.365,1.194) (1.398,1.201) (1.430,1.208) (1.462,1.214) (1.495,1.221) (1.528,1.226) (1.560,1.232) (1.593,1.237) (1.625,1.242) (1.657,1.247) (1.690,1.252) (1.723,1.256) (1.755,1.261) (1.788,1.265) (1.820,1.269) (1.853,1.272) (1.885,1.276) (1.917,1.279) (1.950,1.283) (1.983,1.286) (2.015,1.289) (2.047,1.292) (2.080,1.295) (2.113,1.297) (2.144,1.300) (2.178,1.303) (2.210,1.305) (2.243,1.307) (2.275,1.309) (2.308,1.312) (2.340,1.314) (2.373,1.316) (2.405,1.318) (2.438,1.319) (2.470,1.321) (2.502,1.323) (2.535,1.325) (2.568,1.326) (2.600,1.328)}
\LC{(-2.600,1.328) (-2.567,1.326) (-2.535,1.325) (-2.502,1.323) (-2.470,1.321) (-2.439,1.320) (-2.405,1.318) (-2.373,1.316) (-2.340,1.314) (-2.308,1.312) (-2.275,1.309) (-2.243,1.307) (-2.210,1.305) (-2.178,1.303) (-2.145,1.300) (-2.112,1.297) (-2.080,1.295) (-2.047,1.292) (-2.015,1.289) (-1.982,1.286) (-1.950,1.283) (-1.917,1.279) (-1.885,1.276) (-1.853,1.272) (-1.820,1.269) (-1.788,1.265) (-1.757,1.261) (-1.723,1.256) (-1.690,1.252) (-1.657,1.247) (-1.625,1.242) (-1.593,1.237) (-1.560,1.232) (-1.528,1.226) (-1.495,1.221) (-1.462,1.214) (-1.430,1.208) (-1.397,1.201) (-1.365,1.194) (-1.333,1.187) (-1.300,1.179) (-1.268,1.171) (-1.235,1.162) (-1.202,1.153) (-1.172,1.144) (-1.137,1.133) (-1.105,1.122) (-1.074,1.112) (-1.040,1.099) (-1.007,1.086) (-0.976,1.073) (-0.947,1.060) (-0.920,1.047) (-0.890,1.031) (-0.859,1.014) (-0.832,0.998) (-0.806,0.982) (-0.780,0.964) (-0.754,0.945) (-0.728,0.925) (-0.703,0.904) (-0.676,0.878) (-0.650,0.852) (-0.627,0.826) (-0.606,0.799) (-0.606,0.798)}
\LC{(0.245,-0.975) (0.256,-1.007) (0.272,-1.040) (0.290,-1.066) (0.310,-1.092) (0.332,-1.115) (0.356,-1.137) (0.379,-1.157) (0.403,-1.176) (0.429,-1.194) (0.455,-1.212) (0.481,-1.228) (0.507,-1.243) (0.540,-1.261) (0.572,-1.278) (0.605,-1.293) (0.633,-1.307) (0.663,-1.319) (0.694,-1.333) (0.727,-1.345) (0.761,-1.358) (0.793,-1.370) (0.826,-1.381) (0.857,-1.391) (0.891,-1.402) (0.920,-1.411) (0.956,-1.421) (0.988,-1.430) (1.021,-1.438) (1.053,-1.446) (1.086,-1.454) (1.118,-1.462) (1.150,-1.469) (1.181,-1.476) (1.212,-1.482) (1.242,-1.488) (1.274,-1.494) (1.306,-1.500) (1.339,-1.506) (1.372,-1.512) (1.404,-1.517) (1.437,-1.522) (1.469,-1.527) (1.502,-1.532) (1.534,-1.537) (1.567,-1.541) (1.599,-1.545) (1.632,-1.549) (1.664,-1.553) (1.696,-1.557) (1.729,-1.561) (1.762,-1.565) (1.794,-1.568) (1.827,-1.571) (1.859,-1.575) (1.892,-1.578) (1.924,-1.581) (1.957,-1.584) (1.989,-1.586) (2.022,-1.589) (2.054,-1.592) (2.087,-1.594) (2.119,-1.597) (2.151,-1.599) (2.184,-1.601) (2.217,-1.604) (2.247,-1.606) (2.282,-1.608) (2.314,-1.610) (2.347,-1.612) (2.379,-1.613) (2.412,-1.615) (2.444,-1.617) (2.477,-1.619) (2.509,-1.620) (2.542,-1.622) (2.574,-1.623) (2.600,-1.625)}
\LC{(-2.600,-1.625) (-2.567,-1.623) (-2.535,-1.622) (-2.502,-1.620) (-2.470,-1.618) (-2.438,-1.617) (-2.405,-1.615) (-2.373,-1.613) (-2.340,-1.611) (-2.308,-1.609) (-2.275,-1.607) (-2.243,-1.605) (-2.210,-1.603) (-2.178,-1.601) (-2.145,-1.599) (-2.112,-1.596) (-2.080,-1.594) (-2.047,-1.591) (-2.015,-1.589) (-1.984,-1.586) (-1.950,-1.583) (-1.917,-1.580) (-1.885,-1.577) (-1.853,-1.574) (-1.820,-1.571) (-1.788,-1.567) (-1.755,-1.564) (-1.723,-1.560) (-1.690,-1.557) (-1.657,-1.553) (-1.625,-1.549) (-1.592,-1.544) (-1.563,-1.540) (-1.528,-1.536) (-1.495,-1.531) (-1.462,-1.526) (-1.430,-1.521) (-1.397,-1.516) (-1.365,-1.511) (-1.333,-1.505) (-1.300,-1.499) (-1.268,-1.493) (-1.235,-1.487) (-1.202,-1.480) (-1.170,-1.473) (-1.137,-1.466) (-1.105,-1.459) (-1.073,-1.451) (-1.040,-1.443) (-1.007,-1.435) (-0.975,-1.426) (-0.943,-1.417) (-0.910,-1.407) (-0.877,-1.398) (-0.845,-1.387) (-0.812,-1.376) (-0.780,-1.365) (-0.747,-1.353) (-0.715,-1.341) (-0.682,-1.328) (-0.650,-1.314) (-0.619,-1.300) (-0.591,-1.287) (-0.559,-1.271) (-0.528,-1.254) (-0.500,-1.239) (-0.474,-1.224) (-0.449,-1.207) (-0.422,-1.190) (-0.397,-1.171) (-0.371,-1.150) (-0.345,-1.127) (-0.322,-1.105) (-0.300,-1.079) (-0.281,-1.053) (-0.265,-1.027) (-0.251,-0.994) (-0.245,-0.975)}
\LC{(0.245,0.975) (0.256,1.008) (0.272,1.040) (0.290,1.066) (0.310,1.092) (0.332,1.115) (0.356,1.138) (0.379,1.157) (0.403,1.176) (0.429,1.194) (0.455,1.212) (0.481,1.228) (0.507,1.243) (0.540,1.261) (0.572,1.278) (0.605,1.293) (0.633,1.307) (0.663,1.319) (0.694,1.333) (0.727,1.346) (0.761,1.358) (0.793,1.370) (0.826,1.381) (0.857,1.391) (0.891,1.402) (0.920,1.411) (0.956,1.421) (0.988,1.430) (1.021,1.438) (1.053,1.446) (1.086,1.454) (1.118,1.462) (1.150,1.469) (1.181,1.476) (1.212,1.482) (1.242,1.488) (1.274,1.494) (1.307,1.500) (1.339,1.506) (1.372,1.512) (1.404,1.517) (1.437,1.522) (1.469,1.527) (1.502,1.532) (1.534,1.537) (1.567,1.541) (1.599,1.545) (1.632,1.549) (1.664,1.553) (1.696,1.557) (1.729,1.561) (1.762,1.565) (1.794,1.568) (1.827,1.571) (1.859,1.575) (1.892,1.578) (1.924,1.581) (1.957,1.584) (1.989,1.586) (2.022,1.589) (2.054,1.592) (2.087,1.594) (2.119,1.597) (2.151,1.599) (2.184,1.601) (2.216,1.604) (2.247,1.606) (2.282,1.608) (2.314,1.610) (2.347,1.612) (2.379,1.613) (2.412,1.615) (2.444,1.617) (2.477,1.619) (2.509,1.620) (2.542,1.622) (2.574,1.623) (2.600,1.625)}
\LC{(-2.600,1.625) (-2.567,1.623) (-2.535,1.622) (-2.502,1.620) (-2.470,1.618) (-2.438,1.617) (-2.405,1.615) (-2.373,1.613) (-2.340,1.611) (-2.308,1.609) (-2.275,1.607) (-2.243,1.605) (-2.210,1.603) (-2.178,1.601) (-2.145,1.599) (-2.112,1.596) (-2.080,1.594) (-2.047,1.591) (-2.015,1.589) (-1.984,1.586) (-1.950,1.583) (-1.917,1.580) (-1.885,1.577) (-1.853,1.574) (-1.820,1.571) (-1.788,1.567) (-1.755,1.564) (-1.723,1.560) (-1.690,1.557) (-1.657,1.553) (-1.625,1.549) (-1.593,1.544) (-1.563,1.541) (-1.528,1.536) (-1.495,1.531) (-1.462,1.526) (-1.430,1.521) (-1.397,1.516) (-1.365,1.511) (-1.333,1.505) (-1.300,1.499) (-1.268,1.493) (-1.235,1.487) (-1.202,1.480) (-1.170,1.473) (-1.137,1.466) (-1.105,1.459) (-1.073,1.451) (-1.040,1.443) (-1.007,1.435) (-0.975,1.426) (-0.943,1.417) (-0.910,1.407) (-0.877,1.398) (-0.845,1.387) (-0.812,1.376) (-0.780,1.365) (-0.747,1.353) (-0.715,1.341) (-0.682,1.328) (-0.650,1.314) (-0.619,1.300) (-0.591,1.287) (-0.559,1.271) (-0.528,1.255) (-0.500,1.239) (-0.474,1.224) (-0.449,1.207) (-0.422,1.190) (-0.397,1.171) (-0.371,1.151) (-0.345,1.127) (-0.322,1.105) (-0.300,1.079) (-0.281,1.053) (-0.265,1.027) (-0.251,0.995) (-0.245,0.975)}
\LC{(-2.600,-1.924) (-2.567,-1.923) (-2.535,-1.921) (-2.502,-1.920) (-2.470,-1.918) (-2.438,-1.917) (-2.405,-1.915) (-2.373,-1.914) (-2.340,-1.912) (-2.308,-1.910) (-2.275,-1.908) (-2.243,-1.907) (-2.210,-1.905) (-2.178,-1.903) (-2.145,-1.901) (-2.112,-1.899) (-2.080,-1.897) (-2.047,-1.895) (-2.015,-1.892) (-1.982,-1.890) (-1.950,-1.888) (-1.917,-1.885) (-1.885,-1.883) (-1.853,-1.880) (-1.820,-1.878) (-1.788,-1.875) (-1.755,-1.872) (-1.723,-1.869) (-1.690,-1.867) (-1.658,-1.864) (-1.625,-1.860) (-1.593,-1.857) (-1.560,-1.854) (-1.528,-1.851) (-1.495,-1.847) (-1.462,-1.844) (-1.430,-1.840) (-1.397,-1.836) (-1.365,-1.832) (-1.333,-1.828) (-1.300,-1.824) (-1.268,-1.820) (-1.235,-1.816) (-1.202,-1.812) (-1.170,-1.807) (-1.137,-1.803) (-1.105,-1.798) (-1.073,-1.793) (-1.040,-1.788) (-1.007,-1.783) (-0.975,-1.778) (-0.942,-1.773) (-0.911,-1.768) (-0.877,-1.762) (-0.845,-1.757) (-0.812,-1.752) (-0.780,-1.746) (-0.747,-1.740) (-0.715,-1.735) (-0.682,-1.729) (-0.650,-1.723) (-0.617,-1.718) (-0.585,-1.712) (-0.552,-1.706) (-0.520,-1.701) (-0.487,-1.695) (-0.458,-1.690) (-0.422,-1.684) (-0.390,-1.679) (-0.357,-1.674) (-0.325,-1.670) (-0.292,-1.665) (-0.260,-1.661) (-0.228,-1.657) (-0.195,-1.654) (-0.163,-1.651) (-0.130,-1.649) (-0.097,-1.647) (-0.065,-1.646) (-0.032,-1.645) (0.000,-1.644) (0.033,-1.645) (0.065,-1.646) (0.098,-1.647) (0.130,-1.649) (0.163,-1.651) (0.195,-1.654) (0.228,-1.657) (0.260,-1.661) (0.292,-1.665) (0.325,-1.670) (0.358,-1.674) (0.390,-1.679) (0.423,-1.684) (0.455,-1.690) (0.488,-1.695) (0.520,-1.701) (0.553,-1.706) (0.585,-1.712) (0.618,-1.718) (0.650,-1.723) (0.682,-1.729) (0.715,-1.735) (0.748,-1.740) (0.780,-1.746) (0.812,-1.752) (0.845,-1.757) (0.878,-1.762) (0.910,-1.768) (0.943,-1.773) (0.975,-1.778) (1.008,-1.783) (1.040,-1.788) (1.073,-1.793) (1.105,-1.798) (1.138,-1.803) (1.169,-1.807) (1.203,-1.812) (1.235,-1.816) (1.266,-1.820) (1.300,-1.824) (1.333,-1.828) (1.365,-1.832) (1.398,-1.836) (1.430,-1.840) (1.462,-1.844) (1.495,-1.847) (1.528,-1.851) (1.560,-1.854) (1.593,-1.857) (1.625,-1.860) (1.658,-1.864) (1.690,-1.867) (1.723,-1.869) (1.755,-1.872) (1.788,-1.875) (1.820,-1.878) (1.853,-1.880) (1.885,-1.883) (1.917,-1.885) (1.950,-1.888) (1.983,-1.890) (2.015,-1.892) (2.047,-1.895) (2.080,-1.897) (2.113,-1.899) (2.145,-1.901) (2.178,-1.903) (2.210,-1.905) (2.243,-1.907) (2.275,-1.908) (2.308,-1.910) (2.340,-1.912) (2.373,-1.914) (2.405,-1.915) (2.438,-1.917) (2.470,-1.918) (2.502,-1.920) (2.535,-1.921) (2.568,-1.923) (2.600,-1.924)}
\LC{(-2.600,1.924) (-2.567,1.923) (-2.535,1.921) (-2.502,1.920) (-2.470,1.918) (-2.438,1.917) (-2.405,1.915) (-2.373,1.914) (-2.340,1.912) (-2.308,1.910) (-2.275,1.908) (-2.243,1.907) (-2.210,1.905) (-2.178,1.903) (-2.145,1.901) (-2.112,1.899) (-2.080,1.897) (-2.047,1.895) (-2.015,1.892) (-1.982,1.890) (-1.950,1.888) (-1.917,1.885) (-1.885,1.883) (-1.853,1.880) (-1.820,1.878) (-1.788,1.875) (-1.755,1.872) (-1.723,1.869) (-1.690,1.867) (-1.657,1.864) (-1.625,1.860) (-1.593,1.857) (-1.560,1.854) (-1.528,1.851) (-1.495,1.847) (-1.462,1.844) (-1.430,1.840) (-1.397,1.836) (-1.365,1.832) (-1.333,1.828) (-1.300,1.824) (-1.268,1.820) (-1.235,1.816) (-1.202,1.812) (-1.170,1.807) (-1.137,1.803) (-1.105,1.798) (-1.073,1.793) (-1.040,1.788) (-1.007,1.783) (-0.975,1.778) (-0.942,1.773) (-0.911,1.768) (-0.877,1.762) (-0.845,1.757) (-0.812,1.752) (-0.780,1.746) (-0.747,1.740) (-0.715,1.735) (-0.682,1.729) (-0.650,1.723) (-0.617,1.718) (-0.585,1.712) (-0.552,1.706) (-0.520,1.701) (-0.487,1.695) (-0.458,1.690) (-0.422,1.684) (-0.390,1.679) (-0.357,1.674) (-0.325,1.670) (-0.292,1.665) (-0.260,1.661) (-0.228,1.658) (-0.195,1.654) (-0.163,1.651) (-0.130,1.649) (-0.097,1.647) (-0.065,1.646) (-0.032,1.645) (0.000,1.644) (0.033,1.645) (0.065,1.646) (0.098,1.647) (0.130,1.649) (0.163,1.651) (0.195,1.654) (0.228,1.657) (0.260,1.661) (0.292,1.665) (0.325,1.670) (0.358,1.674) (0.390,1.679) (0.423,1.684) (0.455,1.690) (0.488,1.695) (0.520,1.701) (0.553,1.706) (0.585,1.712) (0.618,1.718) (0.650,1.723) (0.682,1.729) (0.715,1.735) (0.748,1.740) (0.780,1.746) (0.812,1.752) (0.845,1.757) (0.878,1.762) (0.910,1.768) (0.943,1.773) (0.975,1.778) (1.008,1.783) (1.040,1.788) (1.073,1.793) (1.105,1.798) (1.138,1.803) (1.169,1.807) (1.203,1.812) (1.235,1.816) (1.266,1.820) (1.300,1.824) (1.333,1.828) (1.365,1.832) (1.398,1.836) (1.430,1.840) (1.462,1.844) (1.495,1.847) (1.528,1.851) (1.560,1.854) (1.593,1.857) (1.625,1.860) (1.658,1.864) (1.690,1.867) (1.723,1.869) (1.755,1.872) (1.788,1.875) (1.820,1.878) (1.853,1.880) (1.885,1.883) (1.917,1.885) (1.950,1.888) (1.983,1.890) (2.015,1.892) (2.047,1.895) (2.080,1.897) (2.113,1.899) (2.145,1.901) (2.178,1.903) (2.210,1.905) (2.243,1.907) (2.275,1.908) (2.308,1.910) (2.340,1.912) (2.373,1.914) (2.405,1.915) (2.438,1.917) (2.470,1.918) (2.502,1.920) (2.535,1.921) (2.568,1.923) (2.600,1.924)}
\LC{(-2.600,-2.225) (-2.567,-2.224) (-2.535,-2.223) (-2.502,-2.221) (-2.470,-2.220) (-2.438,-2.219) (-2.405,-2.218) (-2.373,-2.216) (-2.340,-2.215) (-2.308,-2.213) (-2.275,-2.212) (-2.243,-2.210) (-2.210,-2.209) (-2.178,-2.207) (-2.145,-2.206) (-2.112,-2.204) (-2.080,-2.202) (-2.048,-2.200) (-2.015,-2.199) (-1.982,-2.197) (-1.950,-2.195) (-1.917,-2.193) (-1.885,-2.191) (-1.853,-2.189) (-1.820,-2.187) (-1.788,-2.185) (-1.755,-2.183) (-1.723,-2.181) (-1.690,-2.178) (-1.657,-2.176) (-1.625,-2.174) (-1.593,-2.172) (-1.560,-2.169) (-1.528,-2.167) (-1.495,-2.164) (-1.462,-2.162) (-1.430,-2.159) (-1.397,-2.156) (-1.365,-2.154) (-1.333,-2.151) (-1.300,-2.148) (-1.268,-2.146) (-1.235,-2.143) (-1.202,-2.140) (-1.170,-2.137) (-1.137,-2.134) (-1.105,-2.131) (-1.073,-2.128) (-1.040,-2.125) (-1.007,-2.122) (-0.975,-2.119) (-0.942,-2.116) (-0.910,-2.113) (-0.877,-2.110) (-0.845,-2.107) (-0.812,-2.104) (-0.780,-2.101) (-0.747,-2.098) (-0.715,-2.095) (-0.682,-2.092) (-0.650,-2.089) (-0.618,-2.087) (-0.585,-2.084) (-0.552,-2.081) (-0.520,-2.078) (-0.487,-2.076) (-0.455,-2.074) (-0.422,-2.071) (-0.390,-2.069) (-0.357,-2.067) (-0.325,-2.065) (-0.292,-2.063) (-0.260,-2.062) (-0.228,-2.060) (-0.195,-2.059) (-0.163,-2.058) (-0.130,-2.057) (-0.097,-2.056) (-0.065,-2.056) (-0.032,-2.056) (0.000,-2.056) (0.033,-2.056) (0.065,-2.056) (0.098,-2.056) (0.130,-2.057) (0.163,-2.058) (0.195,-2.059) (0.228,-2.060) (0.260,-2.062) (0.292,-2.063) (0.325,-2.065) (0.357,-2.067) (0.390,-2.069) (0.423,-2.071) (0.454,-2.074) (0.488,-2.076) (0.520,-2.078) (0.553,-2.081) (0.585,-2.084) (0.618,-2.086) (0.650,-2.089) (0.683,-2.092) (0.715,-2.095) (0.748,-2.098) (0.780,-2.101) (0.812,-2.104) (0.845,-2.107) (0.878,-2.110) (0.910,-2.113) (0.943,-2.116) (0.973,-2.119) (1.008,-2.122) (1.040,-2.125) (1.073,-2.128) (1.105,-2.131) (1.138,-2.134) (1.170,-2.137) (1.203,-2.140) (1.235,-2.143) (1.268,-2.146) (1.300,-2.148) (1.333,-2.151) (1.365,-2.154) (1.398,-2.156) (1.430,-2.159) (1.462,-2.162) (1.495,-2.164) (1.528,-2.167) (1.560,-2.169) (1.593,-2.172) (1.625,-2.174) (1.658,-2.176) (1.690,-2.178) (1.723,-2.181) (1.755,-2.183) (1.788,-2.185) (1.820,-2.187) (1.853,-2.189) (1.885,-2.191) (1.917,-2.193) (1.950,-2.195) (1.983,-2.197) (2.015,-2.199) (2.048,-2.200) (2.080,-2.202) (2.113,-2.204) (2.145,-2.206) (2.178,-2.207) (2.210,-2.209) (2.243,-2.210) (2.275,-2.212) (2.308,-2.213) (2.340,-2.215) (2.373,-2.216) (2.405,-2.218) (2.438,-2.219) (2.470,-2.220) (2.502,-2.221) (2.535,-2.223) (2.568,-2.224) (2.600,-2.225)}
\LC{(-2.600,2.225) (-2.567,2.224) (-2.535,2.223) (-2.502,2.221) (-2.470,2.220) (-2.438,2.219) (-2.405,2.218) (-2.372,2.216) (-2.340,2.215) (-2.308,2.213) (-2.275,2.212) (-2.243,2.210) (-2.210,2.209) (-2.178,2.207) (-2.145,2.206) (-2.112,2.204) (-2.080,2.202) (-2.047,2.200) (-2.015,2.199) (-1.982,2.197) (-1.950,2.195) (-1.917,2.193) (-1.885,2.191) (-1.853,2.189) (-1.820,2.187) (-1.788,2.185) (-1.755,2.183) (-1.723,2.181) (-1.690,2.178) (-1.658,2.176) (-1.625,2.174) (-1.593,2.172) (-1.560,2.169) (-1.527,2.167) (-1.495,2.164) (-1.462,2.162) (-1.430,2.159) (-1.397,2.156) (-1.365,2.154) (-1.333,2.151) (-1.300,2.148) (-1.268,2.146) (-1.235,2.143) (-1.202,2.140) (-1.170,2.137) (-1.137,2.134) (-1.105,2.131) (-1.073,2.128) (-1.040,2.125) (-1.007,2.122) (-0.975,2.119) (-0.942,2.116) (-0.910,2.113) (-0.877,2.110) (-0.845,2.107) (-0.812,2.104) (-0.780,2.101) (-0.747,2.098) (-0.715,2.095) (-0.682,2.092) (-0.650,2.089) (-0.618,2.087) (-0.585,2.084) (-0.552,2.081) (-0.520,2.078) (-0.487,2.076) (-0.455,2.074) (-0.422,2.071) (-0.390,2.069) (-0.357,2.067) (-0.325,2.065) (-0.292,2.063) (-0.260,2.062) (-0.228,2.060) (-0.195,2.059) (-0.163,2.058) (-0.130,2.057) (-0.097,2.056) (-0.065,2.056) (-0.032,2.056) (0.000,2.056) (0.033,2.056) (0.065,2.056) (0.098,2.056) (0.130,2.057) (0.163,2.058) (0.195,2.059) (0.228,2.060) (0.260,2.062) (0.292,2.063) (0.325,2.065) (0.357,2.067) (0.390,2.069) (0.423,2.071) (0.454,2.074) (0.488,2.076) (0.520,2.078) (0.553,2.081) (0.585,2.084) (0.618,2.086) (0.650,2.089) (0.683,2.092) (0.715,2.095) (0.748,2.098) (0.780,2.101) (0.812,2.104) (0.845,2.107) (0.878,2.110) (0.910,2.113) (0.943,2.116) (0.973,2.119) (1.008,2.122) (1.040,2.125) (1.073,2.128) (1.105,2.131) (1.138,2.134) (1.170,2.137) (1.203,2.140) (1.235,2.143) (1.268,2.146) (1.300,2.148) (1.333,2.151) (1.365,2.154) (1.398,2.156) (1.430,2.159) (1.462,2.162) (1.495,2.164) (1.528,2.167) (1.560,2.169) (1.593,2.172) (1.625,2.174) (1.658,2.176) (1.690,2.178) (1.723,2.181) (1.755,2.183) (1.788,2.185) (1.820,2.187) (1.853,2.189) (1.885,2.191) (1.917,2.193) (1.950,2.195) (1.983,2.197) (2.015,2.199) (2.047,2.200) (2.080,2.202) (2.113,2.204) (2.145,2.206) (2.178,2.207) (2.210,2.209) (2.243,2.210) (2.275,2.212) (2.308,2.213) (2.340,2.215) (2.372,2.216) (2.405,2.218) (2.438,2.219) (2.470,2.220) (2.502,2.221) (2.535,2.223) (2.568,2.224) (2.600,2.225)}
\LC{(-2.600,-2.528) (-2.567,-2.527) (-2.535,-2.525) (-2.502,-2.524) (-2.470,-2.523) (-2.438,-2.522) (-2.405,-2.521) (-2.373,-2.520) (-2.340,-2.519) (-2.308,-2.518) (-2.275,-2.516) (-2.243,-2.515) (-2.210,-2.514) (-2.178,-2.513) (-2.145,-2.511) (-2.112,-2.510) (-2.080,-2.509) (-2.047,-2.507) (-2.015,-2.506) (-1.982,-2.504) (-1.950,-2.503) (-1.918,-2.501) (-1.885,-2.500) (-1.853,-2.498) (-1.820,-2.497) (-1.788,-2.495) (-1.755,-2.494) (-1.723,-2.492) (-1.690,-2.490) (-1.657,-2.489) (-1.625,-2.487) (-1.592,-2.485) (-1.560,-2.483) (-1.528,-2.482) (-1.495,-2.480) (-1.462,-2.478) (-1.430,-2.476) (-1.397,-2.474) (-1.365,-2.473) (-1.333,-2.471) (-1.300,-2.469) (-1.268,-2.467) (-1.235,-2.465) (-1.203,-2.463) (-1.170,-2.461) (-1.137,-2.459) (-1.105,-2.457) (-1.073,-2.456) (-1.040,-2.454) (-1.007,-2.452) (-0.975,-2.450) (-0.942,-2.448) (-0.910,-2.446) (-0.877,-2.444) (-0.845,-2.442) (-0.812,-2.441) (-0.780,-2.439) (-0.747,-2.437) (-0.715,-2.435) (-0.682,-2.434) (-0.650,-2.432) (-0.617,-2.430) (-0.585,-2.429) (-0.552,-2.427) (-0.520,-2.426) (-0.487,-2.425) (-0.455,-2.423) (-0.422,-2.422) (-0.390,-2.421) (-0.357,-2.420) (-0.325,-2.419) (-0.292,-2.418) (-0.260,-2.417) (-0.228,-2.417) (-0.195,-2.416) (-0.163,-2.415) (-0.130,-2.415) (-0.097,-2.415) (-0.065,-2.414) (-0.032,-2.414) (0.000,-2.414) (0.033,-2.414) (0.065,-2.414) (0.098,-2.415) (0.130,-2.415) (0.163,-2.415) (0.195,-2.416) (0.228,-2.417) (0.260,-2.417) (0.292,-2.418) (0.325,-2.419) (0.358,-2.420) (0.390,-2.421) (0.423,-2.422) (0.455,-2.423) (0.488,-2.425) (0.520,-2.426) (0.553,-2.427) (0.585,-2.429) (0.618,-2.430) (0.650,-2.432) (0.683,-2.434) (0.715,-2.435) (0.748,-2.437) (0.780,-2.439) (0.812,-2.441) (0.845,-2.442) (0.878,-2.444) (0.910,-2.446) (0.943,-2.448) (0.975,-2.450) (1.008,-2.452) (1.040,-2.454) (1.073,-2.456) (1.105,-2.457) (1.138,-2.459) (1.170,-2.461) (1.203,-2.463) (1.235,-2.465) (1.268,-2.467) (1.300,-2.469) (1.333,-2.471) (1.365,-2.473) (1.398,-2.474) (1.430,-2.476) (1.462,-2.478) (1.495,-2.480) (1.528,-2.482) (1.560,-2.483) (1.593,-2.485) (1.625,-2.487) (1.657,-2.489) (1.690,-2.490) (1.723,-2.492) (1.755,-2.494) (1.788,-2.495) (1.820,-2.497) (1.853,-2.498) (1.885,-2.500) (1.918,-2.501) (1.950,-2.503) (1.983,-2.504) (2.015,-2.506) (2.047,-2.507) (2.080,-2.509) (2.113,-2.510) (2.145,-2.511) (2.178,-2.513) (2.210,-2.514) (2.243,-2.515) (2.275,-2.516) (2.308,-2.518) (2.340,-2.519) (2.373,-2.520) (2.405,-2.521) (2.438,-2.522) (2.470,-2.523) (2.502,-2.524) (2.535,-2.525) (2.568,-2.527) (2.600,-2.528)}
\LC{(-2.600,2.528) (-2.567,2.527) (-2.535,2.525) (-2.502,2.524) (-2.470,2.523) (-2.438,2.522) (-2.405,2.521) (-2.373,2.520) (-2.340,2.519) (-2.308,2.518) (-2.275,2.516) (-2.243,2.515) (-2.210,2.514) (-2.178,2.513) (-2.145,2.511) (-2.112,2.510) (-2.080,2.509) (-2.047,2.507) (-2.015,2.506) (-1.982,2.504) (-1.950,2.503) (-1.918,2.501) (-1.885,2.500) (-1.853,2.498) (-1.820,2.497) (-1.788,2.495) (-1.755,2.494) (-1.723,2.492) (-1.690,2.490) (-1.657,2.489) (-1.625,2.487) (-1.592,2.485) (-1.560,2.483) (-1.528,2.482) (-1.495,2.480) (-1.462,2.478) (-1.430,2.476) (-1.397,2.474) (-1.365,2.473) (-1.333,2.471) (-1.300,2.469) (-1.268,2.467) (-1.235,2.465) (-1.202,2.463) (-1.170,2.461) (-1.137,2.459) (-1.105,2.457) (-1.073,2.456) (-1.040,2.454) (-1.007,2.452) (-0.975,2.450) (-0.942,2.448) (-0.910,2.446) (-0.877,2.444) (-0.845,2.442) (-0.812,2.441) (-0.780,2.439) (-0.747,2.437) (-0.715,2.435) (-0.682,2.434) (-0.650,2.432) (-0.617,2.430) (-0.585,2.429) (-0.552,2.427) (-0.520,2.426) (-0.487,2.425) (-0.455,2.423) (-0.422,2.422) (-0.390,2.421) (-0.357,2.420) (-0.325,2.419) (-0.292,2.418) (-0.260,2.417) (-0.228,2.417) (-0.195,2.416) (-0.163,2.415) (-0.130,2.415) (-0.097,2.415) (-0.065,2.414) (-0.032,2.414) (0.000,2.414) (0.033,2.414) (0.065,2.414) (0.098,2.415) (0.130,2.415) (0.163,2.415) (0.195,2.416) (0.228,2.417) (0.260,2.417) (0.292,2.418) (0.325,2.419) (0.358,2.420) (0.390,2.421) (0.423,2.422) (0.455,2.423) (0.488,2.425) (0.520,2.426) (0.553,2.427) (0.585,2.429) (0.618,2.430) (0.650,2.432) (0.683,2.434) (0.715,2.435) (0.748,2.437) (0.780,2.439) (0.812,2.441) (0.845,2.442) (0.878,2.444) (0.910,2.446) (0.943,2.448) (0.975,2.450) (1.008,2.452) (1.040,2.454) (1.073,2.456) (1.105,2.457) (1.138,2.459) (1.170,2.461) (1.203,2.463) (1.235,2.465) (1.268,2.467) (1.300,2.469) (1.333,2.471) (1.365,2.473) (1.398,2.474) (1.430,2.476) (1.462,2.478) (1.495,2.480) (1.528,2.482) (1.560,2.483) (1.593,2.485) (1.625,2.487) (1.657,2.489) (1.690,2.490) (1.723,2.492) (1.755,2.494) (1.788,2.495) (1.820,2.497) (1.853,2.498) (1.885,2.500) (1.918,2.501) (1.950,2.503) (1.983,2.504) (2.015,2.506) (2.047,2.507) (2.080,2.509) (2.113,2.510) (2.145,2.511) (2.178,2.513) (2.210,2.514) (2.243,2.515) (2.275,2.516) (2.308,2.518) (2.340,2.519) (2.373,2.520) (2.405,2.521) (2.438,2.522) (2.470,2.523) (2.502,2.524) (2.535,2.525) (2.568,2.527) (2.600,2.528)}
\LC{(1.000,0.105) (1.034,0.108) (1.066,0.111) (1.099,0.114) (1.129,0.117) (1.164,0.120) (1.196,0.123) (1.229,0.125) (1.261,0.128) (1.294,0.130) (1.326,0.132) (1.359,0.135) (1.391,0.137) (1.424,0.139) (1.456,0.141) (1.489,0.142) (1.521,0.144) (1.554,0.146) (1.586,0.147) (1.619,0.149) (1.651,0.150) (1.683,0.151) (1.716,0.153) (1.749,0.154) (1.781,0.155) (1.813,0.156) (1.846,0.157) (1.879,0.158) (1.911,0.159) (1.944,0.160) (1.976,0.161) (2.009,0.161) (2.041,0.162) (2.074,0.163) (2.106,0.163) (2.139,0.164) (2.171,0.165) (2.204,0.165) (2.236,0.166) (2.268,0.166) (2.301,0.167) (2.334,0.167) (2.366,0.168) (2.399,0.168) (2.431,0.169) (2.464,0.169) (2.496,0.169) (2.529,0.170) (2.561,0.170) (2.594,0.171) (2.600,0.171)}
\LC{(-2.600,0.171) (-2.567,0.170) (-2.535,0.170) (-2.502,0.170) (-2.470,0.169) (-2.438,0.169) (-2.405,0.168) (-2.373,0.168) (-2.340,0.167) (-2.308,0.167) (-2.275,0.166) (-2.242,0.166) (-2.210,0.165) (-2.178,0.165) (-2.145,0.164) (-2.112,0.164) (-2.080,0.163) (-2.047,0.162) (-2.015,0.161) (-1.982,0.161) (-1.950,0.160) (-1.917,0.159) (-1.885,0.158) (-1.853,0.157) (-1.820,0.156) (-1.788,0.155) (-1.755,0.154) (-1.723,0.153) (-1.690,0.152) (-1.657,0.150) (-1.625,0.149) (-1.593,0.147) (-1.560,0.146) (-1.527,0.144) (-1.495,0.143) (-1.462,0.141) (-1.430,0.139) (-1.397,0.137) (-1.365,0.135) (-1.333,0.133) (-1.300,0.131) (-1.268,0.128) (-1.235,0.126) (-1.202,0.123) (-1.170,0.121) (-1.137,0.118) (-1.105,0.115) (-1.073,0.112) (-1.040,0.109) (-1.007,0.105) (-1.000,0.105)}
\LC{(1.000,-0.105) (1.034,-0.108) (1.066,-0.111) (1.099,-0.114) (1.129,-0.117) (1.164,-0.120) (1.196,-0.123) (1.229,-0.125) (1.261,-0.128) (1.294,-0.130) (1.326,-0.132) (1.359,-0.135) (1.391,-0.137) (1.424,-0.139) (1.456,-0.141) (1.489,-0.142) (1.521,-0.144) (1.554,-0.146) (1.586,-0.147) (1.619,-0.149) (1.651,-0.150) (1.683,-0.151) (1.716,-0.153) (1.749,-0.154) (1.781,-0.155) (1.814,-0.156) (1.846,-0.157) (1.879,-0.158) (1.911,-0.159) (1.944,-0.160) (1.976,-0.161) (2.009,-0.161) (2.041,-0.162) (2.074,-0.163) (2.106,-0.163) (2.139,-0.164) (2.171,-0.165) (2.204,-0.165) (2.236,-0.166) (2.268,-0.166) (2.301,-0.167) (2.334,-0.167) (2.366,-0.168) (2.399,-0.168) (2.431,-0.169) (2.464,-0.169) (2.496,-0.169) (2.529,-0.170) (2.561,-0.170) (2.594,-0.171) (2.600,-0.171)}
\LC{(-2.600,-0.171) (-2.567,-0.170) (-2.535,-0.170) (-2.502,-0.170) (-2.470,-0.169) (-2.438,-0.169) (-2.405,-0.168) (-2.373,-0.168) (-2.340,-0.167) (-2.308,-0.167) (-2.275,-0.166) (-2.242,-0.166) (-2.210,-0.165) (-2.178,-0.165) (-2.145,-0.164) (-2.112,-0.164) (-2.080,-0.163) (-2.047,-0.162) (-2.015,-0.161) (-1.982,-0.161) (-1.950,-0.160) (-1.917,-0.159) (-1.885,-0.158) (-1.853,-0.157) (-1.820,-0.156) (-1.788,-0.155) (-1.755,-0.154) (-1.723,-0.153) (-1.690,-0.152) (-1.657,-0.150) (-1.625,-0.149) (-1.593,-0.147) (-1.560,-0.146) (-1.528,-0.144) (-1.495,-0.143) (-1.462,-0.141) (-1.430,-0.139) (-1.397,-0.137) (-1.365,-0.135) (-1.333,-0.133) (-1.300,-0.131) (-1.268,-0.128) (-1.235,-0.126) (-1.202,-0.123) (-1.170,-0.121) (-1.137,-0.118) (-1.105,-0.115) (-1.073,-0.112) (-1.040,-0.109) (-1.007,-0.105) (-1.000,-0.105)}
\LC{(0.964,0.278) (0.995,0.287) (1.027,0.295) (1.060,0.304) (1.092,0.312) (1.125,0.320) (1.157,0.327) (1.190,0.334) (1.222,0.341) (1.255,0.348) (1.287,0.354) (1.320,0.360) (1.352,0.365) (1.385,0.370) (1.417,0.375) (1.450,0.380) (1.482,0.385) (1.515,0.389) (1.547,0.393) (1.579,0.397) (1.612,0.400) (1.645,0.404) (1.677,0.407) (1.710,0.410) (1.742,0.413) (1.773,0.416) (1.807,0.419) (1.840,0.421) (1.872,0.424) (1.905,0.426) (1.937,0.428) (1.970,0.430) (2.002,0.432) (2.034,0.434) (2.067,0.436) (2.100,0.438) (2.132,0.439) (2.165,0.441) (2.197,0.442) (2.230,0.444) (2.262,0.445) (2.295,0.446) (2.327,0.448) (2.360,0.449) (2.392,0.450) (2.425,0.451) (2.457,0.452) (2.490,0.453) (2.522,0.454) (2.555,0.455) (2.587,0.456) (2.600,0.456)}
\LC{(-2.600,0.456) (-2.567,0.455) (-2.535,0.454) (-2.502,0.453) (-2.470,0.452) (-2.438,0.451) (-2.405,0.450) (-2.373,0.449) (-2.340,0.448) (-2.308,0.447) (-2.275,0.446) (-2.243,0.444) (-2.210,0.443) (-2.178,0.441) (-2.145,0.440) (-2.112,0.438) (-2.080,0.437) (-2.047,0.435) (-2.015,0.433) (-1.982,0.431) (-1.950,0.429) (-1.917,0.427) (-1.885,0.425) (-1.854,0.422) (-1.820,0.420) (-1.788,0.417) (-1.755,0.414) (-1.723,0.412) (-1.690,0.409) (-1.657,0.405) (-1.625,0.402) (-1.593,0.398) (-1.560,0.395) (-1.528,0.391) (-1.495,0.386) (-1.462,0.382) (-1.430,0.377) (-1.397,0.372) (-1.365,0.367) (-1.333,0.362) (-1.300,0.356) (-1.267,0.350) (-1.235,0.344) (-1.202,0.337) (-1.170,0.330) (-1.137,0.323) (-1.105,0.315) (-1.073,0.307) (-1.040,0.299) (-1.007,0.290) (-0.975,0.281) (-0.964,0.278)}
\LC{(0.964,-0.278) (0.995,-0.287) (1.027,-0.295) (1.060,-0.304) (1.092,-0.312) (1.125,-0.320) (1.157,-0.327) (1.190,-0.334) (1.222,-0.341) (1.255,-0.348) (1.287,-0.354) (1.320,-0.360) (1.352,-0.365) (1.385,-0.370) (1.417,-0.375) (1.450,-0.380) (1.482,-0.385) (1.515,-0.389) (1.547,-0.393) (1.579,-0.397) (1.612,-0.400) (1.645,-0.404) (1.677,-0.407) (1.710,-0.410) (1.742,-0.413) (1.773,-0.416) (1.807,-0.419) (1.840,-0.421) (1.872,-0.424) (1.905,-0.426) (1.937,-0.428) (1.970,-0.430) (2.002,-0.432) (2.034,-0.434) (2.067,-0.436) (2.100,-0.438) (2.132,-0.439) (2.165,-0.441) (2.197,-0.442) (2.230,-0.444) (2.262,-0.445) (2.295,-0.446) (2.327,-0.448) (2.360,-0.449) (2.392,-0.450) (2.425,-0.451) (2.457,-0.452) (2.490,-0.453) (2.522,-0.454) (2.555,-0.455) (2.587,-0.456) (2.600,-0.456)}
\LC{(-2.600,-0.456) (-2.567,-0.455) (-2.535,-0.454) (-2.502,-0.453) (-2.470,-0.452) (-2.438,-0.451) (-2.405,-0.450) (-2.373,-0.449) (-2.340,-0.448) (-2.308,-0.447) (-2.275,-0.446) (-2.243,-0.444) (-2.210,-0.443) (-2.178,-0.441) (-2.145,-0.440) (-2.112,-0.438) (-2.080,-0.437) (-2.047,-0.435) (-2.015,-0.433) (-1.982,-0.431) (-1.950,-0.429) (-1.917,-0.427) (-1.885,-0.425) (-1.854,-0.423) (-1.820,-0.420) (-1.788,-0.417) (-1.755,-0.414) (-1.723,-0.412) (-1.690,-0.409) (-1.657,-0.405) (-1.625,-0.402) (-1.593,-0.398) (-1.560,-0.395) (-1.528,-0.391) (-1.495,-0.386) (-1.462,-0.382) (-1.430,-0.377) (-1.397,-0.372) (-1.365,-0.367) (-1.333,-0.362) (-1.300,-0.356) (-1.267,-0.350) (-1.235,-0.344) (-1.202,-0.337) (-1.170,-0.330) (-1.137,-0.323) (-1.105,-0.315) (-1.073,-0.307) (-1.040,-0.299) (-1.007,-0.290) (-0.975,-0.281) (-0.964,-0.278)}
\LC{(0.894,0.451) (0.923,0.465) (0.956,0.481) (0.985,0.494) (1.014,0.506) (1.047,0.520) (1.079,0.532) (1.112,0.544) (1.144,0.555) (1.175,0.566) (1.209,0.576) (1.239,0.585) (1.274,0.595) (1.307,0.603) (1.337,0.611) (1.372,0.619) (1.404,0.626) (1.437,0.633) (1.469,0.640) (1.502,0.646) (1.534,0.652) (1.567,0.657) (1.599,0.663) (1.632,0.668) (1.664,0.672) (1.696,0.677) (1.729,0.681) (1.762,0.685) (1.794,0.689) (1.827,0.693) (1.859,0.696) (1.892,0.699) (1.924,0.703) (1.957,0.706) (1.989,0.708) (2.022,0.711) (2.054,0.714) (2.087,0.716) (2.119,0.718) (2.151,0.721) (2.184,0.723) (2.216,0.725) (2.249,0.727) (2.282,0.729) (2.314,0.731) (2.346,0.732) (2.379,0.734) (2.412,0.735) (2.444,0.737) (2.477,0.738) (2.509,0.740) (2.540,0.741) (2.574,0.742) (2.600,0.743)}
\LC{(-2.600,0.743) (-2.567,0.742) (-2.535,0.741) (-2.502,0.739) (-2.470,0.738) (-2.438,0.737) (-2.405,0.735) (-2.373,0.734) (-2.340,0.732) (-2.308,0.730) (-2.275,0.728) (-2.243,0.726) (-2.210,0.725) (-2.178,0.722) (-2.145,0.720) (-2.112,0.718) (-2.080,0.716) (-2.047,0.713) (-2.015,0.711) (-1.982,0.708) (-1.950,0.705) (-1.919,0.702) (-1.885,0.699) (-1.854,0.696) (-1.820,0.692) (-1.788,0.688) (-1.755,0.684) (-1.723,0.680) (-1.691,0.676) (-1.657,0.671) (-1.625,0.667) (-1.593,0.662) (-1.561,0.656) (-1.528,0.651) (-1.495,0.645) (-1.462,0.639) (-1.430,0.632) (-1.397,0.625) (-1.365,0.618) (-1.333,0.610) (-1.300,0.602) (-1.268,0.593) (-1.235,0.584) (-1.202,0.574) (-1.170,0.564) (-1.137,0.553) (-1.105,0.542) (-1.073,0.530) (-1.040,0.517) (-1.007,0.504) (-0.975,0.489) (-0.943,0.474) (-0.915,0.462) (-0.894,0.451)}
\LC{(0.894,-0.451) (0.923,-0.465) (0.956,-0.481) (0.985,-0.494) (1.014,-0.506) (1.047,-0.520) (1.079,-0.532) (1.112,-0.544) (1.144,-0.555) (1.175,-0.566) (1.209,-0.576) (1.239,-0.585) (1.274,-0.595) (1.307,-0.603) (1.337,-0.611) (1.372,-0.619) (1.404,-0.626) (1.437,-0.633) (1.469,-0.640) (1.502,-0.646) (1.534,-0.652) (1.567,-0.657) (1.599,-0.663) (1.632,-0.668) (1.664,-0.672) (1.696,-0.677) (1.729,-0.681) (1.762,-0.685) (1.794,-0.689) (1.827,-0.693) (1.859,-0.696) (1.892,-0.699) (1.924,-0.703) (1.957,-0.706) (1.989,-0.708) (2.022,-0.711) (2.054,-0.714) (2.087,-0.716) (2.119,-0.718) (2.151,-0.721) (2.184,-0.723) (2.217,-0.725) (2.249,-0.727) (2.282,-0.729) (2.314,-0.731) (2.347,-0.732) (2.379,-0.734) (2.412,-0.735) (2.444,-0.737) (2.476,-0.738) (2.509,-0.740) (2.540,-0.741) (2.574,-0.742) (2.600,-0.743)}
\LC{(-2.600,-0.743) (-2.567,-0.742) (-2.535,-0.741) (-2.503,-0.739) (-2.470,-0.738) (-2.438,-0.737) (-2.405,-0.735) (-2.373,-0.734) (-2.340,-0.732) (-2.308,-0.730) (-2.275,-0.728) (-2.243,-0.726) (-2.210,-0.725) (-2.178,-0.722) (-2.145,-0.720) (-2.112,-0.718) (-2.080,-0.716) (-2.047,-0.713) (-2.015,-0.711) (-1.983,-0.708) (-1.950,-0.705) (-1.919,-0.702) (-1.885,-0.699) (-1.854,-0.696) (-1.820,-0.692) (-1.788,-0.688) (-1.755,-0.684) (-1.723,-0.680) (-1.691,-0.676) (-1.658,-0.671) (-1.625,-0.667) (-1.592,-0.662) (-1.561,-0.657) (-1.528,-0.651) (-1.495,-0.645) (-1.462,-0.639) (-1.430,-0.632) (-1.397,-0.625) (-1.365,-0.618) (-1.333,-0.610) (-1.300,-0.602) (-1.268,-0.593) (-1.235,-0.584) (-1.202,-0.574) (-1.170,-0.564) (-1.137,-0.553) (-1.105,-0.542) (-1.073,-0.530) (-1.040,-0.517) (-1.007,-0.504) (-0.975,-0.489) (-0.943,-0.475) (-0.915,-0.462) (-0.894,-0.451)}
\LC{(0.785,0.626) (0.812,0.647) (0.839,0.666) (0.865,0.684) (0.891,0.701) (0.917,0.717) (0.943,0.732) (0.971,0.747) (1.001,0.763) (1.034,0.779) (1.065,0.793) (1.095,0.806) (1.125,0.818) (1.157,0.830) (1.190,0.842) (1.219,0.851) (1.248,0.861) (1.281,0.871) (1.313,0.880) (1.346,0.889) (1.378,0.897) (1.411,0.905) (1.443,0.912) (1.476,0.919) (1.508,0.926) (1.541,0.932) (1.573,0.938) (1.606,0.944) (1.637,0.949) (1.671,0.954) (1.703,0.959) (1.736,0.964) (1.768,0.968) (1.800,0.972) (1.833,0.976) (1.866,0.980) (1.898,0.984) (1.930,0.987) (1.963,0.991) (1.996,0.994) (2.028,0.997) (2.061,1.000) (2.093,1.003) (2.126,1.005) (2.158,1.008) (2.191,1.010) (2.223,1.013) (2.256,1.015) (2.288,1.017) (2.321,1.019) (2.353,1.021) (2.385,1.023) (2.418,1.025) (2.451,1.027) (2.483,1.028) (2.516,1.030) (2.548,1.031) (2.581,1.033) (2.600,1.034)}
\LC{(-2.600,1.034) (-2.567,1.032) (-2.535,1.031) (-2.502,1.029) (-2.470,1.028) (-2.438,1.026) (-2.405,1.024) (-2.373,1.022) (-2.342,1.020) (-2.308,1.018) (-2.275,1.016) (-2.243,1.014) (-2.210,1.012) (-2.178,1.009) (-2.145,1.007) (-2.112,1.004) (-2.080,1.002) (-2.047,0.999) (-2.015,0.996) (-1.982,0.993) (-1.950,0.989) (-1.917,0.986) (-1.885,0.983) (-1.853,0.979) (-1.821,0.975) (-1.788,0.971) (-1.755,0.966) (-1.723,0.962) (-1.690,0.957) (-1.657,0.952) (-1.625,0.947) (-1.593,0.941) (-1.562,0.936) (-1.528,0.930) (-1.495,0.923) (-1.463,0.916) (-1.434,0.910) (-1.404,0.903) (-1.371,0.895) (-1.339,0.887) (-1.306,0.878) (-1.274,0.869) (-1.242,0.859) (-1.209,0.848) (-1.180,0.839) (-1.150,0.828) (-1.118,0.815) (-1.085,0.802) (-1.053,0.788) (-1.022,0.773) (-0.994,0.760) (-0.962,0.743) (-0.935,0.728) (-0.903,0.709) (-0.877,0.693) (-0.851,0.675) (-0.825,0.657) (-0.799,0.637) (-0.785,0.626)}
\LC{(0.785,-0.626) (0.812,-0.647) (0.839,-0.666) (0.865,-0.684) (0.891,-0.701) (0.917,-0.717) (0.943,-0.732) (0.971,-0.748) (1.001,-0.763) (1.034,-0.779) (1.065,-0.793) (1.095,-0.806) (1.125,-0.818) (1.157,-0.830) (1.190,-0.842) (1.219,-0.852) (1.248,-0.861) (1.281,-0.871) (1.313,-0.880) (1.346,-0.889) (1.378,-0.897) (1.411,-0.905) (1.443,-0.912) (1.476,-0.919) (1.508,-0.926) (1.541,-0.932) (1.573,-0.938) (1.606,-0.944) (1.637,-0.949) (1.671,-0.954) (1.703,-0.959) (1.736,-0.964) (1.768,-0.968) (1.800,-0.972) (1.833,-0.976) (1.866,-0.980) (1.898,-0.984) (1.930,-0.987) (1.963,-0.991) (1.996,-0.994) (2.028,-0.997) (2.061,-1.000) (2.093,-1.003) (2.126,-1.005) (2.158,-1.008) (2.191,-1.010) (2.223,-1.013) (2.256,-1.015) (2.288,-1.017) (2.321,-1.019) (2.353,-1.021) (2.385,-1.023) (2.418,-1.025) (2.451,-1.027) (2.483,-1.028) (2.516,-1.030) (2.548,-1.031) (2.581,-1.033) (2.600,-1.034)}
\LC{(-2.600,-1.034) (-2.568,-1.032) (-2.535,-1.031) (-2.502,-1.029) (-2.470,-1.028) (-2.438,-1.026) (-2.405,-1.024) (-2.373,-1.022) (-2.342,-1.021) (-2.308,-1.018) (-2.275,-1.016) (-2.243,-1.014) (-2.210,-1.012) (-2.178,-1.009) (-2.145,-1.007) (-2.112,-1.004) (-2.080,-1.002) (-2.047,-0.999) (-2.015,-0.996) (-1.982,-0.993) (-1.950,-0.989) (-1.917,-0.986) (-1.885,-0.983) (-1.853,-0.979) (-1.821,-0.975) (-1.788,-0.971) (-1.755,-0.966) (-1.723,-0.962) (-1.690,-0.957) (-1.657,-0.952) (-1.625,-0.947) (-1.593,-0.941) (-1.562,-0.936) (-1.528,-0.930) (-1.495,-0.923) (-1.463,-0.917) (-1.434,-0.910) (-1.404,-0.903) (-1.371,-0.895) (-1.339,-0.887) (-1.306,-0.878) (-1.274,-0.869) (-1.242,-0.859) (-1.209,-0.848) (-1.180,-0.839) (-1.150,-0.828) (-1.118,-0.815) (-1.085,-0.802) (-1.053,-0.788) (-1.022,-0.774) (-0.994,-0.760) (-0.962,-0.743) (-0.935,-0.728) (-0.903,-0.709) (-0.877,-0.693) (-0.852,-0.675) (-0.825,-0.657) (-0.799,-0.637) (-0.785,-0.626)}
\LC{(0.606,0.798) (0.627,0.825) (0.650,0.851) (0.675,0.877) (0.702,0.903) (0.726,0.923) (0.750,0.942) (0.777,0.962) (0.806,0.982) (0.832,0.998) (0.858,1.013) (0.891,1.031) (0.920,1.046) (0.947,1.059) (0.975,1.072) (1.007,1.085) (1.039,1.099) (1.073,1.111) (1.105,1.122) (1.138,1.133) (1.170,1.143) (1.203,1.153) (1.235,1.162) (1.264,1.170) (1.300,1.179) (1.333,1.187) (1.365,1.194) (1.398,1.201) (1.430,1.208) (1.462,1.214) (1.495,1.221) (1.528,1.226) (1.560,1.232) (1.593,1.237) (1.625,1.242) (1.658,1.247) (1.690,1.252) (1.723,1.256) (1.755,1.261) (1.788,1.265) (1.820,1.269) (1.853,1.272) (1.885,1.276) (1.917,1.279) (1.950,1.283) (1.983,1.286) (2.015,1.289) (2.047,1.292) (2.080,1.295) (2.113,1.297) (2.144,1.300) (2.178,1.303) (2.210,1.305) (2.243,1.307) (2.275,1.309) (2.308,1.312) (2.340,1.314) (2.373,1.316) (2.405,1.318) (2.438,1.319) (2.470,1.321) (2.502,1.323) (2.535,1.325) (2.568,1.326) (2.600,1.328)}
\LC{(-2.600,1.328) (-2.567,1.326) (-2.535,1.325) (-2.502,1.323) (-2.470,1.321) (-2.439,1.319) (-2.405,1.318) (-2.373,1.316) (-2.340,1.314) (-2.308,1.312) (-2.275,1.309) (-2.243,1.307) (-2.210,1.305) (-2.178,1.303) (-2.145,1.300) (-2.112,1.297) (-2.080,1.295) (-2.047,1.292) (-2.015,1.289) (-1.982,1.286) (-1.950,1.283) (-1.917,1.279) (-1.885,1.276) (-1.853,1.272) (-1.820,1.269) (-1.788,1.265) (-1.757,1.261) (-1.723,1.256) (-1.690,1.252) (-1.657,1.247) (-1.625,1.242) (-1.592,1.237) (-1.560,1.232) (-1.528,1.226) (-1.495,1.221) (-1.462,1.214) (-1.430,1.208) (-1.397,1.201) (-1.365,1.194) (-1.333,1.187) (-1.300,1.179) (-1.268,1.171) (-1.235,1.162) (-1.202,1.153) (-1.172,1.144) (-1.137,1.133) (-1.105,1.122) (-1.074,1.111) (-1.040,1.099) (-1.007,1.086) (-0.976,1.073) (-0.947,1.059) (-0.920,1.046) (-0.890,1.031) (-0.859,1.014) (-0.832,0.998) (-0.806,0.982) (-0.780,0.964) (-0.754,0.945) (-0.728,0.925) (-0.703,0.903) (-0.676,0.878) (-0.650,0.851) (-0.627,0.825) (-0.606,0.799) (-0.606,0.798)}
\LC{(0.606,-0.798) (0.627,-0.826) (0.650,-0.851) (0.675,-0.878) (0.702,-0.903) (0.726,-0.923) (0.750,-0.943) (0.777,-0.962) (0.806,-0.982) (0.832,-0.998) (0.858,-1.013) (0.891,-1.031) (0.920,-1.047) (0.947,-1.060) (0.975,-1.072) (1.007,-1.086) (1.039,-1.099) (1.073,-1.111) (1.105,-1.122) (1.138,-1.133) (1.170,-1.143) (1.203,-1.153) (1.235,-1.162) (1.264,-1.170) (1.300,-1.179) (1.333,-1.187) (1.365,-1.194) (1.398,-1.201) (1.430,-1.208) (1.462,-1.214) (1.495,-1.221) (1.528,-1.226) (1.560,-1.232) (1.593,-1.237) (1.625,-1.242) (1.657,-1.247) (1.690,-1.252) (1.723,-1.256) (1.755,-1.261) (1.788,-1.265) (1.820,-1.269) (1.853,-1.272) (1.885,-1.276) (1.917,-1.279) (1.950,-1.283) (1.983,-1.286) (2.015,-1.289) (2.047,-1.292) (2.080,-1.295) (2.113,-1.297) (2.144,-1.300) (2.178,-1.303) (2.210,-1.305) (2.243,-1.307) (2.275,-1.309) (2.308,-1.312) (2.340,-1.314) (2.373,-1.316) (2.405,-1.318) (2.438,-1.319) (2.470,-1.321) (2.502,-1.323) (2.535,-1.325) (2.568,-1.326) (2.600,-1.328)}
\LC{(-2.600,-1.328) (-2.567,-1.326) (-2.535,-1.325) (-2.502,-1.323) (-2.470,-1.321) (-2.439,-1.320) (-2.405,-1.318) (-2.373,-1.316) (-2.340,-1.314) (-2.308,-1.312) (-2.275,-1.309) (-2.243,-1.307) (-2.210,-1.305) (-2.178,-1.303) (-2.145,-1.300) (-2.112,-1.297) (-2.080,-1.295) (-2.047,-1.292) (-2.015,-1.289) (-1.982,-1.286) (-1.950,-1.283) (-1.917,-1.279) (-1.885,-1.276) (-1.853,-1.272) (-1.820,-1.269) (-1.788,-1.265) (-1.757,-1.261) (-1.723,-1.256) (-1.690,-1.252) (-1.657,-1.247) (-1.625,-1.242) (-1.593,-1.237) (-1.560,-1.232) (-1.528,-1.226) (-1.495,-1.221) (-1.462,-1.214) (-1.430,-1.208) (-1.397,-1.201) (-1.365,-1.194) (-1.333,-1.187) (-1.300,-1.179) (-1.268,-1.171) (-1.235,-1.162) (-1.202,-1.153) (-1.172,-1.144) (-1.137,-1.133) (-1.105,-1.122) (-1.074,-1.112) (-1.040,-1.099) (-1.007,-1.086) (-0.976,-1.073) (-0.947,-1.060) (-0.920,-1.047) (-0.890,-1.031) (-0.859,-1.014) (-0.832,-0.998) (-0.806,-0.982) (-0.780,-0.964) (-0.754,-0.945) (-0.728,-0.925) (-0.703,-0.904) (-0.676,-0.878) (-0.650,-0.852) (-0.627,-0.826) (-0.606,-0.799) (-0.606,-0.798)}
\LC{(0.245,0.975) (0.256,1.007) (0.272,1.040) (0.290,1.066) (0.310,1.092) (0.332,1.115) (0.356,1.137) (0.379,1.157) (0.403,1.176) (0.429,1.194) (0.455,1.212) (0.481,1.228) (0.507,1.243) (0.540,1.261) (0.572,1.278) (0.605,1.293) (0.633,1.307) (0.663,1.319) (0.694,1.333) (0.727,1.345) (0.761,1.358) (0.793,1.370) (0.826,1.381) (0.857,1.391) (0.891,1.402) (0.920,1.411) (0.956,1.421) (0.988,1.430) (1.021,1.438) (1.053,1.446) (1.086,1.454) (1.118,1.462) (1.150,1.469) (1.181,1.476) (1.212,1.482) (1.242,1.488) (1.274,1.494) (1.306,1.500) (1.339,1.506) (1.372,1.512) (1.404,1.517) (1.437,1.522) (1.469,1.527) (1.502,1.532) (1.534,1.537) (1.567,1.541) (1.599,1.545) (1.632,1.549) (1.664,1.553) (1.696,1.557) (1.729,1.561) (1.762,1.565) (1.794,1.568) (1.827,1.571) (1.859,1.575) (1.892,1.578) (1.924,1.581) (1.957,1.584) (1.989,1.586) (2.022,1.589) (2.054,1.592) (2.087,1.594) (2.119,1.597) (2.151,1.599) (2.184,1.601) (2.217,1.604) (2.247,1.606) (2.282,1.608) (2.314,1.610) (2.347,1.612) (2.379,1.613) (2.412,1.615) (2.444,1.617) (2.477,1.619) (2.509,1.620) (2.542,1.622) (2.574,1.623) (2.600,1.625)}
\LC{(-2.600,1.625) (-2.567,1.623) (-2.535,1.622) (-2.502,1.620) (-2.470,1.618) (-2.438,1.617) (-2.405,1.615) (-2.373,1.613) (-2.340,1.611) (-2.308,1.609) (-2.275,1.607) (-2.243,1.605) (-2.210,1.603) (-2.178,1.601) (-2.145,1.599) (-2.112,1.596) (-2.080,1.594) (-2.047,1.591) (-2.015,1.589) (-1.984,1.586) (-1.950,1.583) (-1.917,1.580) (-1.885,1.577) (-1.853,1.574) (-1.820,1.571) (-1.788,1.567) (-1.755,1.564) (-1.723,1.560) (-1.690,1.557) (-1.657,1.553) (-1.625,1.549) (-1.592,1.544) (-1.563,1.540) (-1.528,1.536) (-1.495,1.531) (-1.462,1.526) (-1.430,1.521) (-1.397,1.516) (-1.365,1.511) (-1.333,1.505) (-1.300,1.499) (-1.268,1.493) (-1.235,1.487) (-1.202,1.480) (-1.170,1.473) (-1.137,1.466) (-1.105,1.459) (-1.073,1.451) (-1.040,1.443) (-1.007,1.435) (-0.975,1.426) (-0.943,1.417) (-0.910,1.407) (-0.877,1.398) (-0.845,1.387) (-0.812,1.376) (-0.780,1.365) (-0.747,1.353) (-0.715,1.341) (-0.682,1.328) (-0.650,1.314) (-0.619,1.300) (-0.591,1.287) (-0.559,1.271) (-0.528,1.254) (-0.500,1.239) (-0.474,1.224) (-0.449,1.207) (-0.422,1.190) (-0.397,1.171) (-0.371,1.150) (-0.345,1.127) (-0.322,1.105) (-0.300,1.079) (-0.281,1.053) (-0.265,1.027) (-0.251,0.994) (-0.245,0.975)}
\LC{(0.245,-0.975) (0.256,-1.008) (0.272,-1.040) (0.290,-1.066) (0.310,-1.092) (0.332,-1.115) (0.356,-1.138) (0.379,-1.157) (0.403,-1.176) (0.429,-1.194) (0.455,-1.212) (0.481,-1.228) (0.507,-1.243) (0.540,-1.261) (0.572,-1.278) (0.605,-1.293) (0.633,-1.307) (0.663,-1.319) (0.694,-1.333) (0.727,-1.346) (0.761,-1.358) (0.793,-1.370) (0.826,-1.381) (0.857,-1.391) (0.891,-1.402) (0.920,-1.411) (0.956,-1.421) (0.988,-1.430) (1.021,-1.438) (1.053,-1.446) (1.086,-1.454) (1.118,-1.462) (1.150,-1.469) (1.181,-1.476) (1.212,-1.482) (1.242,-1.488) (1.274,-1.494) (1.307,-1.500) (1.339,-1.506) (1.372,-1.512) (1.404,-1.517) (1.437,-1.522) (1.469,-1.527) (1.502,-1.532) (1.534,-1.537) (1.567,-1.541) (1.599,-1.545) (1.632,-1.549) (1.664,-1.553) (1.696,-1.557) (1.729,-1.561) (1.762,-1.565) (1.794,-1.568) (1.827,-1.571) (1.859,-1.575) (1.892,-1.578) (1.924,-1.581) (1.957,-1.584) (1.989,-1.586) (2.022,-1.589) (2.054,-1.592) (2.087,-1.594) (2.119,-1.597) (2.151,-1.599) (2.184,-1.601) (2.216,-1.604) (2.247,-1.606) (2.282,-1.608) (2.314,-1.610) (2.347,-1.612) (2.379,-1.613) (2.412,-1.615) (2.444,-1.617) (2.477,-1.619) (2.509,-1.620) (2.542,-1.622) (2.574,-1.623) (2.600,-1.625)}
\LC{(-2.600,-1.625) (-2.567,-1.623) (-2.535,-1.622) (-2.502,-1.620) (-2.470,-1.618) (-2.438,-1.617) (-2.405,-1.615) (-2.373,-1.613) (-2.340,-1.611) (-2.308,-1.609) (-2.275,-1.607) (-2.243,-1.605) (-2.210,-1.603) (-2.178,-1.601) (-2.145,-1.599) (-2.112,-1.596) (-2.080,-1.594) (-2.047,-1.591) (-2.015,-1.589) (-1.984,-1.586) (-1.950,-1.583) (-1.917,-1.580) (-1.885,-1.577) (-1.853,-1.574) (-1.820,-1.571) (-1.788,-1.567) (-1.755,-1.564) (-1.723,-1.560) (-1.690,-1.557) (-1.657,-1.553) (-1.625,-1.549) (-1.593,-1.544) (-1.563,-1.541) (-1.528,-1.536) (-1.495,-1.531) (-1.462,-1.526) (-1.430,-1.521) (-1.397,-1.516) (-1.365,-1.511) (-1.333,-1.505) (-1.300,-1.499) (-1.268,-1.493) (-1.235,-1.487) (-1.202,-1.480) (-1.170,-1.473) (-1.137,-1.466) (-1.105,-1.459) (-1.073,-1.451) (-1.040,-1.443) (-1.007,-1.435) (-0.975,-1.426) (-0.943,-1.417) (-0.910,-1.407) (-0.877,-1.398) (-0.845,-1.387) (-0.812,-1.376) (-0.780,-1.365) (-0.747,-1.353) (-0.715,-1.341) (-0.682,-1.328) (-0.650,-1.314) (-0.619,-1.300) (-0.591,-1.287) (-0.559,-1.271) (-0.528,-1.255) (-0.500,-1.239) (-0.474,-1.224) (-0.449,-1.207) (-0.422,-1.190) (-0.397,-1.171) (-0.371,-1.151) (-0.345,-1.127) (-0.322,-1.105) (-0.300,-1.079) (-0.281,-1.053) (-0.265,-1.027) (-0.251,-0.995) (-0.245,-0.975)}
\LC{(-2.600,1.924) (-2.567,1.923) (-2.535,1.921) (-2.502,1.920) (-2.470,1.918) (-2.438,1.917) (-2.405,1.915) (-2.373,1.914) (-2.340,1.912) (-2.308,1.910) (-2.275,1.908) (-2.243,1.907) (-2.210,1.905) (-2.178,1.903) (-2.145,1.901) (-2.112,1.899) (-2.080,1.897) (-2.047,1.895) (-2.015,1.892) (-1.982,1.890) (-1.950,1.888) (-1.917,1.885) (-1.885,1.883) (-1.853,1.880) (-1.820,1.878) (-1.788,1.875) (-1.755,1.872) (-1.723,1.869) (-1.690,1.867) (-1.658,1.864) (-1.625,1.860) (-1.593,1.857) (-1.560,1.854) (-1.528,1.851) (-1.495,1.847) (-1.462,1.844) (-1.430,1.840) (-1.397,1.836) (-1.365,1.832) (-1.333,1.828) (-1.300,1.824) (-1.268,1.820) (-1.235,1.816) (-1.202,1.812) (-1.170,1.807) (-1.137,1.803) (-1.105,1.798) (-1.073,1.793) (-1.040,1.788) (-1.007,1.783) (-0.975,1.778) (-0.942,1.773) (-0.911,1.768) (-0.877,1.762) (-0.845,1.757) (-0.812,1.752) (-0.780,1.746) (-0.747,1.740) (-0.715,1.735) (-0.682,1.729) (-0.650,1.723) (-0.617,1.718) (-0.585,1.712) (-0.552,1.706) (-0.520,1.701) (-0.487,1.695) (-0.458,1.690) (-0.422,1.684) (-0.390,1.679) (-0.357,1.674) (-0.325,1.670) (-0.292,1.665) (-0.260,1.661) (-0.228,1.657) (-0.195,1.654) (-0.163,1.651) (-0.130,1.649) (-0.097,1.647) (-0.065,1.646) (-0.032,1.645) (0.000,1.644) (0.033,1.645) (0.065,1.646) (0.098,1.647) (0.130,1.649) (0.163,1.651) (0.195,1.654) (0.228,1.657) (0.260,1.661) (0.292,1.665) (0.325,1.670) (0.358,1.674) (0.390,1.679) (0.423,1.684) (0.455,1.690) (0.488,1.695) (0.520,1.701) (0.553,1.706) (0.585,1.712) (0.618,1.718) (0.650,1.723) (0.682,1.729) (0.715,1.735) (0.748,1.740) (0.780,1.746) (0.812,1.752) (0.845,1.757) (0.878,1.762) (0.910,1.768) (0.943,1.773) (0.975,1.778) (1.008,1.783) (1.040,1.788) (1.073,1.793) (1.105,1.798) (1.138,1.803) (1.169,1.807) (1.203,1.812) (1.235,1.816) (1.266,1.820) (1.300,1.824) (1.333,1.828) (1.365,1.832) (1.398,1.836) (1.430,1.840) (1.462,1.844) (1.495,1.847) (1.528,1.851) (1.560,1.854) (1.593,1.857) (1.625,1.860) (1.658,1.864) (1.690,1.867) (1.723,1.869) (1.755,1.872) (1.788,1.875) (1.820,1.878) (1.853,1.880) (1.885,1.883) (1.917,1.885) (1.950,1.888) (1.983,1.890) (2.015,1.892) (2.047,1.895) (2.080,1.897) (2.113,1.899) (2.145,1.901) (2.178,1.903) (2.210,1.905) (2.243,1.907) (2.275,1.908) (2.308,1.910) (2.340,1.912) (2.373,1.914) (2.405,1.915) (2.438,1.917) (2.470,1.918) (2.502,1.920) (2.535,1.921) (2.568,1.923) (2.600,1.924)}
\LC{(-2.600,-1.924) (-2.567,-1.923) (-2.535,-1.921) (-2.502,-1.920) (-2.470,-1.918) (-2.438,-1.917) (-2.405,-1.915) (-2.373,-1.914) (-2.340,-1.912) (-2.308,-1.910) (-2.275,-1.908) (-2.243,-1.907) (-2.210,-1.905) (-2.178,-1.903) (-2.145,-1.901) (-2.112,-1.899) (-2.080,-1.897) (-2.047,-1.895) (-2.015,-1.892) (-1.982,-1.890) (-1.950,-1.888) (-1.917,-1.885) (-1.885,-1.883) (-1.853,-1.880) (-1.820,-1.878) (-1.788,-1.875) (-1.755,-1.872) (-1.723,-1.869) (-1.690,-1.867) (-1.657,-1.864) (-1.625,-1.860) (-1.593,-1.857) (-1.560,-1.854) (-1.528,-1.851) (-1.495,-1.847) (-1.462,-1.844) (-1.430,-1.840) (-1.397,-1.836) (-1.365,-1.832) (-1.333,-1.828) (-1.300,-1.824) (-1.268,-1.820) (-1.235,-1.816) (-1.202,-1.812) (-1.170,-1.807) (-1.137,-1.803) (-1.105,-1.798) (-1.073,-1.793) (-1.040,-1.788) (-1.007,-1.783) (-0.975,-1.778) (-0.942,-1.773) (-0.911,-1.768) (-0.877,-1.762) (-0.845,-1.757) (-0.812,-1.752) (-0.780,-1.746) (-0.747,-1.740) (-0.715,-1.735) (-0.682,-1.729) (-0.650,-1.723) (-0.617,-1.718) (-0.585,-1.712) (-0.552,-1.706) (-0.520,-1.701) (-0.487,-1.695) (-0.458,-1.690) (-0.422,-1.684) (-0.390,-1.679) (-0.357,-1.674) (-0.325,-1.670) (-0.292,-1.665) (-0.260,-1.661) (-0.228,-1.658) (-0.195,-1.654) (-0.163,-1.651) (-0.130,-1.649) (-0.097,-1.647) (-0.065,-1.646) (-0.032,-1.645) (0.000,-1.644) (0.033,-1.645) (0.065,-1.646) (0.098,-1.647) (0.130,-1.649) (0.163,-1.651) (0.195,-1.654) (0.228,-1.657) (0.260,-1.661) (0.292,-1.665) (0.325,-1.670) (0.358,-1.674) (0.390,-1.679) (0.423,-1.684) (0.455,-1.690) (0.488,-1.695) (0.520,-1.701) (0.553,-1.706) (0.585,-1.712) (0.618,-1.718) (0.650,-1.723) (0.682,-1.729) (0.715,-1.735) (0.748,-1.740) (0.780,-1.746) (0.812,-1.752) (0.845,-1.757) (0.878,-1.762) (0.910,-1.768) (0.943,-1.773) (0.975,-1.778) (1.008,-1.783) (1.040,-1.788) (1.073,-1.793) (1.105,-1.798) (1.138,-1.803) (1.169,-1.807) (1.203,-1.812) (1.235,-1.816) (1.266,-1.820) (1.300,-1.824) (1.333,-1.828) (1.365,-1.832) (1.398,-1.836) (1.430,-1.840) (1.462,-1.844) (1.495,-1.847) (1.528,-1.851) (1.560,-1.854) (1.593,-1.857) (1.625,-1.860) (1.658,-1.864) (1.690,-1.867) (1.723,-1.869) (1.755,-1.872) (1.788,-1.875) (1.820,-1.878) (1.853,-1.880) (1.885,-1.883) (1.917,-1.885) (1.950,-1.888) (1.983,-1.890) (2.015,-1.892) (2.047,-1.895) (2.080,-1.897) (2.113,-1.899) (2.145,-1.901) (2.178,-1.903) (2.210,-1.905) (2.243,-1.907) (2.275,-1.908) (2.308,-1.910) (2.340,-1.912) (2.373,-1.914) (2.405,-1.915) (2.438,-1.917) (2.470,-1.918) (2.502,-1.920) (2.535,-1.921) (2.568,-1.923) (2.600,-1.924)}
\LC{(-2.600,2.225) (-2.567,2.224) (-2.535,2.223) (-2.502,2.221) (-2.470,2.220) (-2.438,2.219) (-2.405,2.218) (-2.373,2.216) (-2.340,2.215) (-2.308,2.213) (-2.275,2.212) (-2.243,2.210) (-2.210,2.209) (-2.178,2.207) (-2.145,2.206) (-2.112,2.204) (-2.080,2.202) (-2.048,2.200) (-2.015,2.199) (-1.982,2.197) (-1.950,2.195) (-1.917,2.193) (-1.885,2.191) (-1.853,2.189) (-1.820,2.187) (-1.788,2.185) (-1.755,2.183) (-1.723,2.181) (-1.690,2.178) (-1.657,2.176) (-1.625,2.174) (-1.593,2.172) (-1.560,2.169) (-1.528,2.167) (-1.495,2.164) (-1.462,2.162) (-1.430,2.159) (-1.397,2.156) (-1.365,2.154) (-1.333,2.151) (-1.300,2.148) (-1.268,2.146) (-1.235,2.143) (-1.202,2.140) (-1.170,2.137) (-1.137,2.134) (-1.105,2.131) (-1.073,2.128) (-1.040,2.125) (-1.007,2.122) (-0.975,2.119) (-0.942,2.116) (-0.910,2.113) (-0.877,2.110) (-0.845,2.107) (-0.812,2.104) (-0.780,2.101) (-0.747,2.098) (-0.715,2.095) (-0.682,2.092) (-0.650,2.089) (-0.618,2.087) (-0.585,2.084) (-0.552,2.081) (-0.520,2.078) (-0.487,2.076) (-0.455,2.074) (-0.422,2.071) (-0.390,2.069) (-0.357,2.067) (-0.325,2.065) (-0.292,2.063) (-0.260,2.062) (-0.228,2.060) (-0.195,2.059) (-0.163,2.058) (-0.130,2.057) (-0.097,2.056) (-0.065,2.056) (-0.032,2.056) (0.000,2.056) (0.033,2.056) (0.065,2.056) (0.098,2.056) (0.130,2.057) (0.163,2.058) (0.195,2.059) (0.228,2.060) (0.260,2.062) (0.292,2.063) (0.325,2.065) (0.357,2.067) (0.390,2.069) (0.423,2.071) (0.454,2.074) (0.488,2.076) (0.520,2.078) (0.553,2.081) (0.585,2.084) (0.618,2.086) (0.650,2.089) (0.683,2.092) (0.715,2.095) (0.748,2.098) (0.780,2.101) (0.812,2.104) (0.845,2.107) (0.878,2.110) (0.910,2.113) (0.943,2.116) (0.973,2.119) (1.008,2.122) (1.040,2.125) (1.073,2.128) (1.105,2.131) (1.138,2.134) (1.170,2.137) (1.203,2.140) (1.235,2.143) (1.268,2.146) (1.300,2.148) (1.333,2.151) (1.365,2.154) (1.398,2.156) (1.430,2.159) (1.462,2.162) (1.495,2.164) (1.528,2.167) (1.560,2.169) (1.593,2.172) (1.625,2.174) (1.658,2.176) (1.690,2.178) (1.723,2.181) (1.755,2.183) (1.788,2.185) (1.820,2.187) (1.853,2.189) (1.885,2.191) (1.917,2.193) (1.950,2.195) (1.983,2.197) (2.015,2.199) (2.048,2.200) (2.080,2.202) (2.113,2.204) (2.145,2.206) (2.178,2.207) (2.210,2.209) (2.243,2.210) (2.275,2.212) (2.308,2.213) (2.340,2.215) (2.373,2.216) (2.405,2.218) (2.438,2.219) (2.470,2.220) (2.502,2.221) (2.535,2.223) (2.568,2.224) (2.600,2.225)}
\LC{(-2.600,-2.225) (-2.567,-2.224) (-2.535,-2.223) (-2.502,-2.221) (-2.470,-2.220) (-2.438,-2.219) (-2.405,-2.218) (-2.372,-2.216) (-2.340,-2.215) (-2.308,-2.213) (-2.275,-2.212) (-2.243,-2.210) (-2.210,-2.209) (-2.178,-2.207) (-2.145,-2.206) (-2.112,-2.204) (-2.080,-2.202) (-2.047,-2.200) (-2.015,-2.199) (-1.982,-2.197) (-1.950,-2.195) (-1.917,-2.193) (-1.885,-2.191) (-1.853,-2.189) (-1.820,-2.187) (-1.788,-2.185) (-1.755,-2.183) (-1.723,-2.181) (-1.690,-2.178) (-1.658,-2.176) (-1.625,-2.174) (-1.593,-2.172) (-1.560,-2.169) (-1.527,-2.167) (-1.495,-2.164) (-1.462,-2.162) (-1.430,-2.159) (-1.397,-2.156) (-1.365,-2.154) (-1.333,-2.151) (-1.300,-2.148) (-1.268,-2.146) (-1.235,-2.143) (-1.202,-2.140) (-1.170,-2.137) (-1.137,-2.134) (-1.105,-2.131) (-1.073,-2.128) (-1.040,-2.125) (-1.007,-2.122) (-0.975,-2.119) (-0.942,-2.116) (-0.910,-2.113) (-0.877,-2.110) (-0.845,-2.107) (-0.812,-2.104) (-0.780,-2.101) (-0.747,-2.098) (-0.715,-2.095) (-0.682,-2.092) (-0.650,-2.089) (-0.618,-2.087) (-0.585,-2.084) (-0.552,-2.081) (-0.520,-2.078) (-0.487,-2.076) (-0.455,-2.074) (-0.422,-2.071) (-0.390,-2.069) (-0.357,-2.067) (-0.325,-2.065) (-0.292,-2.063) (-0.260,-2.062) (-0.228,-2.060) (-0.195,-2.059) (-0.163,-2.058) (-0.130,-2.057) (-0.097,-2.056) (-0.065,-2.056) (-0.032,-2.056) (0.000,-2.056) (0.033,-2.056) (0.065,-2.056) (0.098,-2.056) (0.130,-2.057) (0.163,-2.058) (0.195,-2.059) (0.228,-2.060) (0.260,-2.062) (0.292,-2.063) (0.325,-2.065) (0.357,-2.067) (0.390,-2.069) (0.423,-2.071) (0.454,-2.074) (0.488,-2.076) (0.520,-2.078) (0.553,-2.081) (0.585,-2.084) (0.618,-2.086) (0.650,-2.089) (0.683,-2.092) (0.715,-2.095) (0.748,-2.098) (0.780,-2.101) (0.812,-2.104) (0.845,-2.107) (0.878,-2.110) (0.910,-2.113) (0.943,-2.116) (0.973,-2.119) (1.008,-2.122) (1.040,-2.125) (1.073,-2.128) (1.105,-2.131) (1.138,-2.134) (1.170,-2.137) (1.203,-2.140) (1.235,-2.143) (1.268,-2.146) (1.300,-2.148) (1.333,-2.151) (1.365,-2.154) (1.398,-2.156) (1.430,-2.159) (1.462,-2.162) (1.495,-2.164) (1.528,-2.167) (1.560,-2.169) (1.593,-2.172) (1.625,-2.174) (1.658,-2.176) (1.690,-2.178) (1.723,-2.181) (1.755,-2.183) (1.788,-2.185) (1.820,-2.187) (1.853,-2.189) (1.885,-2.191) (1.917,-2.193) (1.950,-2.195) (1.983,-2.197) (2.015,-2.199) (2.047,-2.200) (2.080,-2.202) (2.113,-2.204) (2.145,-2.206) (2.178,-2.207) (2.210,-2.209) (2.243,-2.210) (2.275,-2.212) (2.308,-2.213) (2.340,-2.215) (2.372,-2.216) (2.405,-2.218) (2.438,-2.219) (2.470,-2.220) (2.502,-2.221) (2.535,-2.223) (2.568,-2.224) (2.600,-2.225)}
\LC{(-2.600,2.528) (-2.567,2.527) (-2.535,2.525) (-2.502,2.524) (-2.470,2.523) (-2.438,2.522) (-2.405,2.521) (-2.373,2.520) (-2.340,2.519) (-2.308,2.518) (-2.275,2.516) (-2.243,2.515) (-2.210,2.514) (-2.178,2.513) (-2.145,2.511) (-2.112,2.510) (-2.080,2.509) (-2.047,2.507) (-2.015,2.506) (-1.982,2.504) (-1.950,2.503) (-1.918,2.501) (-1.885,2.500) (-1.853,2.498) (-1.820,2.497) (-1.788,2.495) (-1.755,2.494) (-1.723,2.492) (-1.690,2.490) (-1.657,2.489) (-1.625,2.487) (-1.592,2.485) (-1.560,2.483) (-1.528,2.482) (-1.495,2.480) (-1.462,2.478) (-1.430,2.476) (-1.397,2.474) (-1.365,2.473) (-1.333,2.471) (-1.300,2.469) (-1.268,2.467) (-1.235,2.465) (-1.203,2.463) (-1.170,2.461) (-1.137,2.459) (-1.105,2.457) (-1.073,2.456) (-1.040,2.454) (-1.007,2.452) (-0.975,2.450) (-0.942,2.448) (-0.910,2.446) (-0.877,2.444) (-0.845,2.442) (-0.812,2.441) (-0.780,2.439) (-0.747,2.437) (-0.715,2.435) (-0.682,2.434) (-0.650,2.432) (-0.617,2.430) (-0.585,2.429) (-0.552,2.427) (-0.520,2.426) (-0.487,2.425) (-0.455,2.423) (-0.422,2.422) (-0.390,2.421) (-0.357,2.420) (-0.325,2.419) (-0.292,2.418) (-0.260,2.417) (-0.228,2.417) (-0.195,2.416) (-0.163,2.415) (-0.130,2.415) (-0.097,2.415) (-0.065,2.414) (-0.032,2.414) (0.000,2.414) (0.033,2.414) (0.065,2.414) (0.098,2.415) (0.130,2.415) (0.163,2.415) (0.195,2.416) (0.228,2.417) (0.260,2.417) (0.292,2.418) (0.325,2.419) (0.358,2.420) (0.390,2.421) (0.423,2.422) (0.455,2.423) (0.488,2.425) (0.520,2.426) (0.553,2.427) (0.585,2.429) (0.618,2.430) (0.650,2.432) (0.683,2.434) (0.715,2.435) (0.748,2.437) (0.780,2.439) (0.812,2.441) (0.845,2.442) (0.878,2.444) (0.910,2.446) (0.943,2.448) (0.975,2.450) (1.008,2.452) (1.040,2.454) (1.073,2.456) (1.105,2.457) (1.138,2.459) (1.170,2.461) (1.203,2.463) (1.235,2.465) (1.268,2.467) (1.300,2.469) (1.333,2.471) (1.365,2.473) (1.398,2.474) (1.430,2.476) (1.462,2.478) (1.495,2.480) (1.528,2.482) (1.560,2.483) (1.593,2.485) (1.625,2.487) (1.657,2.489) (1.690,2.490) (1.723,2.492) (1.755,2.494) (1.788,2.495) (1.820,2.497) (1.853,2.498) (1.885,2.500) (1.918,2.501) (1.950,2.503) (1.983,2.504) (2.015,2.506) (2.047,2.507) (2.080,2.509) (2.113,2.510) (2.145,2.511) (2.178,2.513) (2.210,2.514) (2.243,2.515) (2.275,2.516) (2.308,2.518) (2.340,2.519) (2.373,2.520) (2.405,2.521) (2.438,2.522) (2.470,2.523) (2.502,2.524) (2.535,2.525) (2.568,2.527) (2.600,2.528)}
\LC{(-2.600,-2.528) (-2.567,-2.527) (-2.535,-2.525) (-2.502,-2.524) (-2.470,-2.523) (-2.438,-2.522) (-2.405,-2.521) (-2.373,-2.520) (-2.340,-2.519) (-2.308,-2.518) (-2.275,-2.516) (-2.243,-2.515) (-2.210,-2.514) (-2.178,-2.513) (-2.145,-2.511) (-2.112,-2.510) (-2.080,-2.509) (-2.047,-2.507) (-2.015,-2.506) (-1.982,-2.504) (-1.950,-2.503) (-1.918,-2.501) (-1.885,-2.500) (-1.853,-2.498) (-1.820,-2.497) (-1.788,-2.495) (-1.755,-2.494) (-1.723,-2.492) (-1.690,-2.490) (-1.657,-2.489) (-1.625,-2.487) (-1.592,-2.485) (-1.560,-2.483) (-1.528,-2.482) (-1.495,-2.480) (-1.462,-2.478) (-1.430,-2.476) (-1.397,-2.474) (-1.365,-2.473) (-1.333,-2.471) (-1.300,-2.469) (-1.268,-2.467) (-1.235,-2.465) (-1.202,-2.463) (-1.170,-2.461) (-1.137,-2.459) (-1.105,-2.457) (-1.073,-2.456) (-1.040,-2.454) (-1.007,-2.452) (-0.975,-2.450) (-0.942,-2.448) (-0.910,-2.446) (-0.877,-2.444) (-0.845,-2.442) (-0.812,-2.441) (-0.780,-2.439) (-0.747,-2.437) (-0.715,-2.435) (-0.682,-2.434) (-0.650,-2.432) (-0.617,-2.430) (-0.585,-2.429) (-0.552,-2.427) (-0.520,-2.426) (-0.487,-2.425) (-0.455,-2.423) (-0.422,-2.422) (-0.390,-2.421) (-0.357,-2.420) (-0.325,-2.419) (-0.292,-2.418) (-0.260,-2.417) (-0.228,-2.417) (-0.195,-2.416) (-0.163,-2.415) (-0.130,-2.415) (-0.097,-2.415) (-0.065,-2.414) (-0.032,-2.414) (0.000,-2.414) (0.033,-2.414) (0.065,-2.414) (0.098,-2.415) (0.130,-2.415) (0.163,-2.415) (0.195,-2.416) (0.228,-2.417) (0.260,-2.417) (0.292,-2.418) (0.325,-2.419) (0.358,-2.420) (0.390,-2.421) (0.423,-2.422) (0.455,-2.423) (0.488,-2.425) (0.520,-2.426) (0.553,-2.427) (0.585,-2.429) (0.618,-2.430) (0.650,-2.432) (0.683,-2.434) (0.715,-2.435) (0.748,-2.437) (0.780,-2.439) (0.812,-2.441) (0.845,-2.442) (0.878,-2.444) (0.910,-2.446) (0.943,-2.448) (0.975,-2.450) (1.008,-2.452) (1.040,-2.454) (1.073,-2.456) (1.105,-2.457) (1.138,-2.459) (1.170,-2.461) (1.203,-2.463) (1.235,-2.465) (1.268,-2.467) (1.300,-2.469) (1.333,-2.471) (1.365,-2.473) (1.398,-2.474) (1.430,-2.476) (1.462,-2.478) (1.495,-2.480) (1.528,-2.482) (1.560,-2.483) (1.593,-2.485) (1.625,-2.487) (1.657,-2.489) (1.690,-2.490) (1.723,-2.492) (1.755,-2.494) (1.788,-2.495) (1.820,-2.497) (1.853,-2.498) (1.885,-2.500) (1.918,-2.501) (1.950,-2.503) (1.983,-2.504) (2.015,-2.506) (2.047,-2.507) (2.080,-2.509) (2.113,-2.510) (2.145,-2.511) (2.178,-2.513) (2.210,-2.514) (2.243,-2.515) (2.275,-2.516) (2.308,-2.518) (2.340,-2.519) (2.373,-2.520) (2.405,-2.521) (2.438,-2.522) (2.470,-2.523) (2.502,-2.524) (2.535,-2.525) (2.568,-2.527) (2.600,-2.528)}
\LC{(-2.600,0.000) (-1.000,0.000)}
\LC{(1.000,0.000) (2.600,0.000)}

      % Esfera en campo uniforme - flechas
% Generado por scripts/generar_datos.py (no editar a mano).
\FL{1.898}{-0.158}{-1.6}
\FL{-1.898}{-0.158}{1.6}
\FL{1.898}{0.158}{1.6}
\FL{-1.898}{0.158}{-1.6}
\FL{1.898}{-0.426}{-4.0}
\FL{-1.898}{-0.426}{4.0}
\FL{1.898}{0.426}{4.0}
\FL{-1.898}{0.426}{-4.0}
\FL{1.898}{-0.700}{-5.6}
\FL{-1.898}{-0.700}{5.6}
\FL{1.898}{0.700}{5.6}
\FL{-1.898}{0.700}{-5.6}
\FL{1.898}{-0.984}{-6.3}
\FL{-1.898}{-0.984}{6.3}
\FL{1.898}{0.984}{6.3}
\FL{-1.898}{0.984}{-6.3}
\FL{1.898}{-1.277}{-6.1}
\FL{-1.898}{-1.277}{6.1}
\FL{1.898}{1.277}{6.1}
\FL{-1.898}{1.277}{-6.1}
\FL{1.898}{-1.578}{-5.3}
\FL{-1.898}{-1.578}{5.3}
\FL{1.898}{1.578}{5.3}
\FL{-1.898}{1.578}{-5.3}
\FL{-1.898}{-1.884}{4.4}
\FL{1.898}{-1.884}{-4.4}
\FL{-1.898}{1.884}{-4.4}
\FL{1.898}{1.884}{4.4}
\FL{-1.898}{-2.192}{3.4}
\FL{1.898}{-2.192}{-3.4}
\FL{-1.898}{2.192}{-3.4}
\FL{1.898}{2.192}{3.4}
\FL{-1.898}{-2.500}{2.7}
\FL{1.898}{-2.500}{-2.7}
\FL{-1.898}{2.500}{-2.7}
\FL{1.898}{2.500}{2.7}
\FL{1.898}{0.158}{1.6}
\FL{-1.898}{0.158}{-1.6}
\FL{1.898}{-0.158}{-1.6}
\FL{-1.898}{-0.158}{1.6}
\FL{1.898}{0.426}{4.0}
\FL{-1.898}{0.426}{-4.0}
\FL{1.898}{-0.426}{-4.0}
\FL{-1.898}{-0.426}{4.0}
\FL{1.898}{0.700}{5.6}
\FL{-1.898}{0.700}{-5.6}
\FL{1.898}{-0.700}{-5.6}
\FL{-1.898}{-0.700}{5.6}
\FL{1.898}{0.984}{6.3}
\FL{-1.898}{0.984}{-6.3}
\FL{1.898}{-0.984}{-6.3}
\FL{-1.898}{-0.984}{6.3}
\FL{1.898}{1.277}{6.1}
\FL{-1.898}{1.277}{-6.1}
\FL{1.898}{-1.277}{-6.1}
\FL{-1.898}{-1.277}{6.1}
\FL{1.898}{1.578}{5.3}
\FL{-1.898}{1.578}{-5.3}
\FL{1.898}{-1.578}{-5.3}
\FL{-1.898}{-1.578}{5.3}
\FL{-1.898}{1.884}{-4.4}
\FL{1.898}{1.884}{4.4}
\FL{-1.898}{-1.884}{4.4}
\FL{1.898}{-1.884}{-4.4}
\FL{-1.898}{2.192}{-3.4}
\FL{1.898}{2.192}{3.4}
\FL{-1.898}{-2.192}{3.4}
\FL{1.898}{-2.192}{-3.4}
\FL{-1.898}{2.500}{-2.7}
\FL{1.898}{2.500}{2.7}
\FL{-1.898}{-2.500}{2.7}
\FL{1.898}{-2.500}{-2.7}
\FL{-2.600}{0.000}{0.0}
\FL{2.600}{0.000}{0.0}

    \end{scope}
    \shade[ball color=azulvelo, draw=azulprof] (0,0) circle (1);
    \draw[azulprof, line width=0.6pt] (0,0) circle (1);
    % Esfera en campo uniforme - cargas inducidas
% Generado por scripts/generar_datos.py (no editar a mano).
\CP{1.123}{0.128}
\CP{1.064}{0.379}
\CP{0.939}{0.629}
\CP{0.709}{0.880}
\CP{0.060}{1.128}
\CP{0.712}{-0.877}
\CP{0.941}{-0.626}
\CP{1.066}{-0.376}
\CP{1.123}{-0.124}
\CN{-0.519}{1.004}
\CN{-0.845}{0.750}
\CN{-1.013}{0.500}
\CN{-1.103}{0.247}
\CN{-1.130}{-0.004}
\CN{-1.101}{-0.253}
\CN{-1.010}{-0.507}
\CN{-0.841}{-0.755}
\CN{-0.513}{-1.007}

    \node[font=\footnotesize, text=azulprof] at (0,0) {$\vE=\vc{0}$};
    \draw[azulclaro, line width=0.5pt] (-2.6,-2.6) rectangle (2.6,2.6);
    \node[anchor=north] at (0,-2.75) {(a) Esfera conductora};
  \end{scope}
  % ---------------------------------------------------------------- (b)
  \begin{scope}[xshift=6.45cm, yshift=-0.35cm, x=1.17cm, y=1.17cm]
    \fill[azulfondo!60!white] (-2.3,-1.9) rectangle (2.3,1.9);
    \begin{scope}
      \clip (-2.3,-1.9) rectangle (2.3,1.9);
      \tikzset{numEQ/.style={azulcielo, line width=0.35pt, densely dashed}}
      % Conductor neutro en campo uniforme - equipotenciales
% Generado por scripts/generar_datos.py (no editar a mano).
\EQ{(2.165,-1.940) (2.168,-1.900) (2.171,-1.860) (2.174,-1.820) (2.177,-1.780) (2.180,-1.740) (2.183,-1.700) (2.186,-1.660) (2.189,-1.620) (2.192,-1.580) (2.195,-1.540) (2.199,-1.500) (2.202,-1.460) (2.205,-1.420) (2.209,-1.380) (2.212,-1.340) (2.216,-1.300) (2.219,-1.260) (2.223,-1.220) (2.226,-1.180) (2.230,-1.140) (2.234,-1.100) (2.237,-1.060) (2.240,-1.029) (2.245,-0.980) (2.248,-0.940) (2.252,-0.900) (2.256,-0.860) (2.259,-0.820) (2.263,-0.780) (2.266,-0.740) (2.270,-0.700) (2.273,-0.660) (2.277,-0.620) (2.280,-0.584) (2.284,-0.540) (2.287,-0.500) (2.290,-0.460) (2.293,-0.420) (2.296,-0.380) (2.299,-0.340) (2.302,-0.300) (2.305,-0.260) (2.308,-0.220) (2.310,-0.180) (2.313,-0.140) (2.315,-0.100) (2.317,-0.060) (2.319,-0.020) (2.321,0.020) (2.322,0.060) (2.324,0.100) (2.325,0.140) (2.326,0.180) (2.327,0.220) (2.328,0.260) (2.328,0.300) (2.328,0.340) (2.329,0.380) (2.329,0.420) (2.328,0.460) (2.328,0.500) (2.327,0.540) (2.326,0.580) (2.325,0.620) (2.324,0.660) (2.323,0.700) (2.321,0.740) (2.319,0.780) (2.318,0.820) (2.316,0.860) (2.313,0.900) (2.311,0.940) (2.308,0.980) (2.306,1.020) (2.303,1.060) (2.300,1.100) (2.297,1.140) (2.294,1.180) (2.291,1.220) (2.288,1.260) (2.284,1.300) (2.281,1.340) (2.277,1.380) (2.274,1.420) (2.270,1.460) (2.266,1.500) (2.263,1.540) (2.260,1.570) (2.255,1.620) (2.252,1.660) (2.248,1.700) (2.244,1.740) (2.240,1.780) (2.237,1.820) (2.233,1.860) (2.229,1.900) (2.225,1.940)}
\EQ{(1.870,-1.940) (1.873,-1.900) (1.876,-1.860) (1.879,-1.820) (1.883,-1.780) (1.886,-1.740) (1.890,-1.700) (1.893,-1.660) (1.897,-1.620) (1.900,-1.586) (1.904,-1.540) (1.908,-1.500) (1.912,-1.460) (1.916,-1.420) (1.920,-1.386) (1.925,-1.340) (1.929,-1.300) (1.933,-1.260) (1.938,-1.220) (1.942,-1.180) (1.947,-1.140) (1.951,-1.100) (1.956,-1.060) (1.960,-1.026) (1.965,-0.980) (1.970,-0.940) (1.975,-0.900) (1.980,-0.860) (1.984,-0.820) (1.989,-0.780) (1.994,-0.740) (1.999,-0.700) (2.003,-0.660) (2.008,-0.620) (2.013,-0.580) (2.017,-0.540) (2.022,-0.500) (2.026,-0.460) (2.031,-0.420) (2.035,-0.380) (2.039,-0.340) (2.043,-0.300) (2.047,-0.260) (2.051,-0.220) (2.054,-0.180) (2.058,-0.140) (2.060,-0.109) (2.064,-0.060) (2.066,-0.020) (2.069,0.020) (2.071,0.060) (2.073,0.100) (2.075,0.140) (2.077,0.180) (2.078,0.220) (2.079,0.260) (2.080,0.297) (2.081,0.340) (2.081,0.380) (2.081,0.420) (2.080,0.460) (2.080,0.500) (2.079,0.540) (2.078,0.580) (2.076,0.620) (2.075,0.660) (2.073,0.700) (2.070,0.740) (2.068,0.780) (2.065,0.820) (2.062,0.860) (2.059,0.900) (2.056,0.940) (2.052,0.980) (2.048,1.020) (2.045,1.060) (2.040,1.100) (2.036,1.140) (2.032,1.180) (2.028,1.220) (2.023,1.260) (2.018,1.300) (2.014,1.340) (2.009,1.380) (2.004,1.420) (2.000,1.455) (1.994,1.500) (1.990,1.540) (1.985,1.580) (1.980,1.618) (1.975,1.660) (1.970,1.700) (1.965,1.740) (1.960,1.780) (1.956,1.820) (1.951,1.860) (1.946,1.900) (1.942,1.940)}
\EQ{(1.569,-1.940) (1.572,-1.900) (1.576,-1.860) (1.580,-1.820) (1.583,-1.780) (1.587,-1.740) (1.591,-1.700) (1.595,-1.660) (1.600,-1.620) (1.604,-1.580) (1.608,-1.540) (1.613,-1.500) (1.618,-1.460) (1.622,-1.420) (1.627,-1.380) (1.632,-1.340) (1.638,-1.300) (1.643,-1.260) (1.648,-1.220) (1.654,-1.180) (1.659,-1.140) (1.665,-1.100) (1.671,-1.060) (1.677,-1.020) (1.683,-0.980) (1.689,-0.940) (1.695,-0.900) (1.700,-0.868) (1.708,-0.820) (1.714,-0.780) (1.720,-0.741) (1.727,-0.700) (1.733,-0.660) (1.739,-0.620) (1.746,-0.580) (1.752,-0.540) (1.758,-0.500) (1.764,-0.460) (1.770,-0.420) (1.776,-0.380) (1.782,-0.340) (1.788,-0.300) (1.793,-0.260) (1.799,-0.220) (1.804,-0.180) (1.809,-0.140) (1.813,-0.100) (1.818,-0.060) (1.822,-0.020) (1.826,0.020) (1.829,0.060) (1.832,0.100) (1.835,0.140) (1.837,0.180) (1.839,0.220) (1.841,0.260) (1.842,0.300) (1.843,0.340) (1.843,0.380) (1.843,0.420) (1.843,0.460) (1.842,0.500) (1.841,0.540) (1.839,0.580) (1.837,0.620) (1.834,0.660) (1.831,0.700) (1.828,0.740) (1.825,0.780) (1.821,0.820) (1.816,0.860) (1.812,0.900) (1.807,0.940) (1.801,0.980) (1.796,1.020) (1.790,1.060) (1.785,1.100) (1.780,1.131) (1.772,1.180) (1.766,1.220) (1.760,1.259) (1.753,1.300) (1.747,1.340) (1.740,1.380) (1.734,1.420) (1.727,1.460) (1.721,1.500) (1.714,1.540) (1.708,1.580) (1.701,1.620) (1.695,1.660) (1.689,1.700) (1.683,1.740) (1.676,1.780) (1.670,1.820) (1.665,1.860) (1.660,1.892) (1.653,1.940)}
\EQ{(1.262,-1.940) (1.266,-1.900) (1.269,-1.860) (1.273,-1.820) (1.277,-1.780) (1.281,-1.740) (1.286,-1.700) (1.290,-1.660) (1.295,-1.620) (1.300,-1.580) (1.305,-1.540) (1.310,-1.500) (1.315,-1.460) (1.320,-1.424) (1.326,-1.380) (1.332,-1.340) (1.338,-1.300) (1.344,-1.260) (1.351,-1.220) (1.357,-1.180) (1.364,-1.140) (1.371,-1.100) (1.378,-1.060) (1.386,-1.020) (1.393,-0.980) (1.400,-0.945) (1.409,-0.900) (1.417,-0.860) (1.425,-0.820) (1.433,-0.780) (1.440,-0.747) (1.450,-0.700) (1.458,-0.660) (1.467,-0.620) (1.476,-0.580) (1.484,-0.540) (1.493,-0.500) (1.500,-0.467) (1.510,-0.420) (1.518,-0.380) (1.527,-0.340) (1.535,-0.300) (1.543,-0.260) (1.551,-0.220) (1.559,-0.180) (1.566,-0.140) (1.573,-0.100) (1.580,-0.060) (1.586,-0.020) (1.592,0.020) (1.597,0.060) (1.602,0.100) (1.606,0.140) (1.610,0.180) (1.614,0.220) (1.616,0.260) (1.618,0.300) (1.620,0.340) (1.620,0.380) (1.621,0.420) (1.620,0.454) (1.619,0.500) (1.617,0.540) (1.614,0.580) (1.611,0.620) (1.607,0.660) (1.602,0.700) (1.597,0.740) (1.591,0.780) (1.585,0.820) (1.580,0.852) (1.572,0.900) (1.564,0.940) (1.556,0.980) (1.548,1.020) (1.540,1.058) (1.531,1.100) (1.522,1.140) (1.513,1.180) (1.504,1.220) (1.495,1.260) (1.486,1.300) (1.477,1.340) (1.468,1.380) (1.460,1.415) (1.450,1.460) (1.441,1.500) (1.433,1.540) (1.424,1.580) (1.416,1.620) (1.408,1.660) (1.400,1.699) (1.392,1.740) (1.384,1.780) (1.377,1.820) (1.370,1.860) (1.363,1.900) (1.356,1.940)}
\EQ{(0.948,-1.940) (0.951,-1.900) (0.954,-1.860) (0.958,-1.820) (0.962,-1.780) (0.966,-1.740) (0.970,-1.700) (0.975,-1.660) (0.980,-1.620) (0.985,-1.580) (0.990,-1.540) (0.995,-1.500) (1.000,-1.466) (1.007,-1.420) (1.013,-1.380) (1.019,-1.340) (1.026,-1.300) (1.033,-1.260) (1.040,-1.222) (1.048,-1.180) (1.056,-1.140) (1.064,-1.100) (1.072,-1.060) (1.080,-1.026) (1.090,-0.980) (1.100,-0.940) (1.109,-0.900) (1.119,-0.860) (1.130,-0.820) (1.140,-0.781) (1.151,-0.740) (1.160,-0.707) (1.173,-0.660) (1.185,-0.620) (1.196,-0.580) (1.208,-0.540) (1.220,-0.501) (1.232,-0.460) (1.244,-0.420) (1.256,-0.380) (1.269,-0.340) (1.280,-0.302) (1.293,-0.260) (1.304,-0.220) (1.316,-0.180) (1.327,-0.140) (1.338,-0.100) (1.349,-0.060) (1.359,-0.020) (1.368,0.020) (1.377,0.060) (1.385,0.100) (1.392,0.140) (1.399,0.180) (1.404,0.220) (1.409,0.260) (1.413,0.300) (1.415,0.340) (1.417,0.380) (1.417,0.420) (1.416,0.460) (1.414,0.500) (1.411,0.540) (1.407,0.580) (1.402,0.620) (1.395,0.660) (1.388,0.700) (1.380,0.736) (1.370,0.780) (1.360,0.819) (1.349,0.860) (1.340,0.891) (1.331,0.920) (1.320,0.957) (1.307,1.000) (1.294,1.040) (1.280,1.080) (1.267,1.120) (1.254,1.160) (1.241,1.200) (1.227,1.240) (1.214,1.280) (1.201,1.320) (1.189,1.360) (1.176,1.400) (1.164,1.440) (1.153,1.480) (1.142,1.520) (1.131,1.560) (1.120,1.600) (1.110,1.640) (1.100,1.680) (1.091,1.720) (1.082,1.760) (1.073,1.800) (1.065,1.840) (1.057,1.880) (1.049,1.920) (1.046,1.940)}
\EQ{(0.624,-1.940) (0.627,-1.900) (0.630,-1.860) (0.633,-1.820) (0.636,-1.780) (0.639,-1.740) (0.643,-1.700) (0.647,-1.660) (0.651,-1.620) (0.655,-1.580) (0.660,-1.540) (0.665,-1.500) (0.670,-1.460) (0.676,-1.420) (0.680,-1.390) (0.685,-1.360) (0.691,-1.320) (0.698,-1.280) (0.705,-1.240) (0.713,-1.200) (0.720,-1.166) (0.730,-1.120) (0.739,-1.080) (0.749,-1.040) (0.759,-1.000) (0.769,-0.960) (0.780,-0.922) (0.792,-0.880) (0.804,-0.840) (0.817,-0.800) (0.830,-0.760) (0.840,-0.731) (0.851,-0.700) (0.865,-0.660) (0.880,-0.620) (0.895,-0.580) (0.911,-0.540) (0.927,-0.500) (0.940,-0.469) (0.952,-0.440) (0.969,-0.400) (0.986,-0.360) (1.000,-0.329) (1.013,-0.300) (1.030,-0.260) (1.048,-0.220) (1.066,-0.180) (1.080,-0.148) (1.092,-0.120) (1.109,-0.080) (1.126,-0.040) (1.140,-0.006) (1.158,0.040) (1.172,0.080) (1.186,0.120) (1.198,0.160) (1.210,0.200) (1.219,0.240) (1.227,0.280) (1.233,0.320) (1.238,0.360) (1.240,0.400) (1.240,0.440) (1.238,0.480) (1.234,0.520) (1.227,0.560) (1.220,0.594) (1.208,0.640) (1.195,0.680) (1.181,0.720) (1.165,0.760) (1.147,0.800) (1.129,0.840) (1.109,0.880) (1.089,0.920) (1.068,0.960) (1.047,1.000) (1.026,1.040) (1.005,1.080) (0.984,1.120) (0.964,1.160) (0.944,1.200) (0.925,1.240) (0.907,1.280) (0.890,1.320) (0.873,1.360) (0.860,1.394) (0.843,1.440) (0.828,1.480) (0.815,1.520) (0.803,1.560) (0.791,1.600) (0.780,1.639) (0.769,1.680) (0.760,1.717) (0.750,1.760) (0.741,1.800) (0.733,1.840) (0.725,1.880) (0.718,1.920) (0.715,1.940)}
\EQ{(0.292,-1.940) (0.293,-1.900) (0.295,-1.860) (0.296,-1.820) (0.298,-1.780) (0.300,-1.740) (0.302,-1.700) (0.304,-1.660) (0.306,-1.620) (0.309,-1.580) (0.312,-1.540) (0.315,-1.500) (0.318,-1.460) (0.322,-1.420) (0.326,-1.380) (0.330,-1.340) (0.335,-1.300) (0.340,-1.263) (0.346,-1.220) (0.353,-1.180) (0.360,-1.141) (0.368,-1.100) (0.377,-1.060) (0.387,-1.020) (0.397,-0.980) (0.409,-0.940) (0.420,-0.903) (0.434,-0.860) (0.448,-0.820) (0.460,-0.789) (0.471,-0.760) (0.487,-0.720) (0.500,-0.689) (0.512,-0.660) (0.530,-0.620) (0.549,-0.580) (0.568,-0.540) (0.588,-0.500) (0.609,-0.460) (0.631,-0.420) (0.653,-0.380) (0.677,-0.340) (0.700,-0.302) (0.720,-0.270) (0.739,-0.240) (0.760,-0.209) (0.779,-0.180) (0.800,-0.150) (0.820,-0.121) (0.840,-0.092) (0.860,-0.063) (0.880,-0.035) (0.900,-0.006) (0.918,0.020) (0.940,0.052) (0.958,0.080) (0.980,0.114) (0.996,0.140) (1.019,0.180) (1.039,0.220) (1.058,0.260) (1.073,0.300) (1.085,0.340) (1.094,0.380) (1.098,0.420) (1.098,0.460) (1.092,0.500) (1.081,0.540) (1.065,0.580) (1.044,0.620) (1.020,0.659) (1.000,0.686) (0.980,0.712) (0.960,0.736) (0.940,0.759) (0.902,0.800) (0.880,0.823) (0.844,0.860) (0.820,0.885) (0.786,0.920) (0.760,0.947) (0.729,0.980) (0.700,1.013) (0.680,1.037) (0.661,1.060) (0.640,1.087) (0.620,1.114) (0.602,1.140) (0.580,1.173) (0.564,1.200) (0.541,1.240) (0.521,1.280) (0.502,1.320) (0.486,1.360) (0.471,1.400) (0.460,1.432) (0.445,1.480) (0.434,1.520) (0.423,1.560) (0.414,1.600) (0.405,1.640) (0.398,1.680) (0.390,1.720) (0.384,1.760) (0.378,1.800) (0.372,1.840) (0.367,1.880) (0.362,1.920) (0.359,1.940)}
\EQ{(-0.383,-1.940) (-0.385,-1.900) (-0.388,-1.860) (-0.390,-1.820) (-0.393,-1.780) (-0.396,-1.740) (-0.399,-1.700) (-0.402,-1.660) (-0.406,-1.620) (-0.409,-1.580) (-0.413,-1.540) (-0.418,-1.500) (-0.422,-1.460) (-0.427,-1.420) (-0.433,-1.380) (-0.439,-1.340) (-0.445,-1.300) (-0.453,-1.260) (-0.460,-1.222) (-0.469,-1.180) (-0.478,-1.140) (-0.488,-1.100) (-0.499,-1.060) (-0.511,-1.020) (-0.524,-0.980) (-0.538,-0.940) (-0.553,-0.900) (-0.568,-0.860) (-0.580,-0.831) (-0.593,-0.800) (-0.610,-0.760) (-0.628,-0.720) (-0.640,-0.692) (-0.655,-0.660) (-0.673,-0.620) (-0.692,-0.580) (-0.711,-0.540) (-0.731,-0.500) (-0.752,-0.460) (-0.773,-0.420) (-0.795,-0.380) (-0.818,-0.340) (-0.840,-0.301) (-0.860,-0.268) (-0.877,-0.240) (-0.900,-0.203) (-0.920,-0.172) (-0.940,-0.141) (-0.960,-0.110) (-0.979,-0.080) (-1.000,-0.047) (-1.016,-0.020) (-1.040,0.020) (-1.060,0.056) (-1.080,0.095) (-1.100,0.139) (-1.115,0.180) (-1.127,0.220) (-1.135,0.260) (-1.139,0.300) (-1.139,0.340) (-1.135,0.380) (-1.126,0.420) (-1.113,0.460) (-1.100,0.492) (-1.086,0.520) (-1.063,0.560) (-1.040,0.596) (-1.020,0.624) (-1.000,0.651) (-0.980,0.675) (-0.960,0.699) (-0.940,0.722) (-0.920,0.745) (-0.888,0.780) (-0.860,0.811) (-0.833,0.840) (-0.800,0.876) (-0.779,0.900) (-0.744,0.940) (-0.720,0.968) (-0.700,0.993) (-0.680,1.019) (-0.660,1.046) (-0.640,1.074) (-0.622,1.100) (-0.600,1.136) (-0.580,1.171) (-0.564,1.200) (-0.545,1.240) (-0.527,1.280) (-0.511,1.320) (-0.500,1.349) (-0.489,1.380) (-0.480,1.410) (-0.471,1.440) (-0.460,1.480) (-0.451,1.520) (-0.442,1.560) (-0.434,1.600) (-0.427,1.640) (-0.420,1.680) (-0.414,1.720) (-0.408,1.760) (-0.403,1.800) (-0.398,1.840) (-0.394,1.880) (-0.390,1.920) (-0.388,1.940)}
\EQ{(-0.713,-1.940) (-0.717,-1.900) (-0.720,-1.862) (-0.724,-1.820) (-0.728,-1.780) (-0.732,-1.740) (-0.737,-1.700) (-0.741,-1.660) (-0.746,-1.620) (-0.752,-1.580) (-0.757,-1.540) (-0.763,-1.500) (-0.769,-1.460) (-0.776,-1.420) (-0.783,-1.380) (-0.790,-1.340) (-0.798,-1.300) (-0.806,-1.260) (-0.815,-1.220) (-0.824,-1.180) (-0.833,-1.140) (-0.843,-1.100) (-0.854,-1.060) (-0.864,-1.020) (-0.876,-0.980) (-0.887,-0.940) (-0.900,-0.900) (-0.912,-0.860) (-0.925,-0.820) (-0.938,-0.780) (-0.952,-0.740) (-0.966,-0.700) (-0.980,-0.662) (-0.995,-0.620) (-1.010,-0.580) (-1.026,-0.540) (-1.040,-0.503) (-1.057,-0.460) (-1.073,-0.420) (-1.089,-0.380) (-1.105,-0.340) (-1.120,-0.304) (-1.138,-0.260) (-1.154,-0.220) (-1.170,-0.180) (-1.185,-0.140) (-1.200,-0.100) (-1.214,-0.060) (-1.227,-0.020) (-1.240,0.020) (-1.251,0.060) (-1.260,0.096) (-1.269,0.140) (-1.276,0.180) (-1.282,0.220) (-1.285,0.260) (-1.286,0.300) (-1.286,0.340) (-1.283,0.380) (-1.280,0.411) (-1.272,0.460) (-1.264,0.500) (-1.254,0.540) (-1.241,0.580) (-1.228,0.620) (-1.213,0.660) (-1.200,0.691) (-1.187,0.720) (-1.169,0.760) (-1.150,0.800) (-1.131,0.840) (-1.110,0.880) (-1.090,0.920) (-1.069,0.960) (-1.049,1.000) (-1.029,1.040) (-1.009,1.080) (-0.989,1.120) (-0.970,1.160) (-0.952,1.200) (-0.940,1.228) (-0.926,1.260) (-0.910,1.300) (-0.895,1.340) (-0.880,1.380) (-0.866,1.420) (-0.853,1.460) (-0.840,1.500) (-0.829,1.540) (-0.820,1.572) (-0.807,1.620) (-0.800,1.651) (-0.793,1.680) (-0.784,1.720) (-0.776,1.760) (-0.768,1.800) (-0.761,1.840) (-0.754,1.880) (-0.748,1.920) (-0.745,1.940)}
\EQ{(-1.034,-1.940) (-1.038,-1.900) (-1.042,-1.860) (-1.046,-1.820) (-1.050,-1.780) (-1.055,-1.740) (-1.060,-1.700) (-1.065,-1.660) (-1.070,-1.620) (-1.076,-1.580) (-1.082,-1.540) (-1.088,-1.500) (-1.094,-1.460) (-1.100,-1.422) (-1.107,-1.380) (-1.114,-1.340) (-1.120,-1.308) (-1.129,-1.260) (-1.137,-1.220) (-1.145,-1.180) (-1.153,-1.140) (-1.160,-1.108) (-1.171,-1.060) (-1.180,-1.020) (-1.189,-0.980) (-1.199,-0.940) (-1.209,-0.900) (-1.219,-0.860) (-1.229,-0.820) (-1.239,-0.780) (-1.250,-0.740) (-1.260,-0.703) (-1.272,-0.660) (-1.280,-0.630) (-1.288,-0.600) (-1.299,-0.560) (-1.311,-0.520) (-1.320,-0.487) (-1.333,-0.440) (-1.344,-0.400) (-1.355,-0.360) (-1.366,-0.320) (-1.376,-0.280) (-1.387,-0.240) (-1.397,-0.200) (-1.406,-0.160) (-1.415,-0.120) (-1.424,-0.080) (-1.432,-0.040) (-1.439,-0.000) (-1.446,0.040) (-1.452,0.080) (-1.457,0.120) (-1.460,0.153) (-1.464,0.200) (-1.466,0.240) (-1.467,0.280) (-1.467,0.320) (-1.466,0.360) (-1.463,0.400) (-1.460,0.440) (-1.456,0.480) (-1.450,0.520) (-1.444,0.560) (-1.437,0.600) (-1.428,0.640) (-1.420,0.677) (-1.409,0.720) (-1.400,0.756) (-1.388,0.800) (-1.376,0.840) (-1.364,0.880) (-1.351,0.920) (-1.340,0.956) (-1.326,1.000) (-1.313,1.040) (-1.300,1.078) (-1.286,1.120) (-1.273,1.160) (-1.260,1.200) (-1.248,1.240) (-1.235,1.280) (-1.223,1.320) (-1.211,1.360) (-1.200,1.397) (-1.188,1.440) (-1.180,1.469) (-1.172,1.500) (-1.161,1.540) (-1.151,1.580) (-1.142,1.620) (-1.133,1.660) (-1.124,1.700) (-1.115,1.740) (-1.107,1.780) (-1.100,1.818) (-1.092,1.860) (-1.085,1.900) (-1.080,1.930) (-1.078,1.940)}
\EQ{(-1.345,-1.940) (-1.349,-1.900) (-1.353,-1.860) (-1.357,-1.820) (-1.361,-1.780) (-1.366,-1.740) (-1.371,-1.700) (-1.375,-1.660) (-1.380,-1.624) (-1.386,-1.580) (-1.391,-1.540) (-1.396,-1.500) (-1.402,-1.460) (-1.408,-1.420) (-1.414,-1.380) (-1.420,-1.340) (-1.426,-1.300) (-1.433,-1.260) (-1.439,-1.220) (-1.446,-1.180) (-1.453,-1.140) (-1.460,-1.101) (-1.467,-1.060) (-1.475,-1.020) (-1.482,-0.980) (-1.490,-0.940) (-1.498,-0.900) (-1.505,-0.860) (-1.513,-0.820) (-1.520,-0.786) (-1.529,-0.740) (-1.537,-0.700) (-1.545,-0.660) (-1.553,-0.620) (-1.560,-0.586) (-1.569,-0.540) (-1.577,-0.500) (-1.585,-0.460) (-1.593,-0.420) (-1.600,-0.380) (-1.607,-0.340) (-1.614,-0.300) (-1.620,-0.268) (-1.628,-0.220) (-1.634,-0.180) (-1.640,-0.140) (-1.645,-0.100) (-1.651,-0.060) (-1.655,-0.020) (-1.659,0.020) (-1.663,0.060) (-1.666,0.100) (-1.669,0.140) (-1.671,0.180) (-1.672,0.220) (-1.673,0.260) (-1.673,0.300) (-1.672,0.340) (-1.671,0.380) (-1.669,0.420) (-1.667,0.460) (-1.664,0.500) (-1.660,0.540) (-1.656,0.580) (-1.651,0.620) (-1.646,0.660) (-1.640,0.698) (-1.633,0.740) (-1.627,0.780) (-1.620,0.817) (-1.612,0.860) (-1.604,0.900) (-1.596,0.940) (-1.587,0.980) (-1.580,1.014) (-1.570,1.060) (-1.561,1.100) (-1.552,1.140) (-1.543,1.180) (-1.534,1.220) (-1.525,1.260) (-1.516,1.300) (-1.507,1.340) (-1.500,1.372) (-1.489,1.420) (-1.481,1.460) (-1.472,1.500) (-1.464,1.540) (-1.456,1.580) (-1.448,1.620) (-1.440,1.660) (-1.433,1.700) (-1.426,1.740) (-1.420,1.772) (-1.412,1.820) (-1.405,1.860) (-1.400,1.892) (-1.392,1.940)}
\EQ{(-1.648,-1.940) (-1.652,-1.900) (-1.655,-1.860) (-1.659,-1.820) (-1.663,-1.780) (-1.667,-1.740) (-1.672,-1.700) (-1.676,-1.660) (-1.680,-1.624) (-1.685,-1.580) (-1.690,-1.540) (-1.694,-1.500) (-1.699,-1.460) (-1.704,-1.420) (-1.709,-1.380) (-1.714,-1.340) (-1.720,-1.300) (-1.725,-1.260) (-1.730,-1.220) (-1.736,-1.180) (-1.742,-1.140) (-1.747,-1.100) (-1.753,-1.060) (-1.759,-1.020) (-1.765,-0.980) (-1.770,-0.940) (-1.776,-0.900) (-1.782,-0.860) (-1.788,-0.820) (-1.794,-0.780) (-1.800,-0.740) (-1.806,-0.700) (-1.812,-0.660) (-1.818,-0.620) (-1.823,-0.580) (-1.829,-0.540) (-1.835,-0.500) (-1.840,-0.461) (-1.845,-0.420) (-1.851,-0.380) (-1.856,-0.340) (-1.860,-0.303) (-1.865,-0.260) (-1.869,-0.220) (-1.873,-0.180) (-1.877,-0.140) (-1.880,-0.109) (-1.884,-0.060) (-1.887,-0.020) (-1.890,0.020) (-1.892,0.060) (-1.894,0.100) (-1.895,0.140) (-1.896,0.180) (-1.897,0.220) (-1.898,0.260) (-1.898,0.300) (-1.897,0.340) (-1.896,0.380) (-1.895,0.420) (-1.893,0.460) (-1.891,0.500) (-1.889,0.540) (-1.886,0.580) (-1.883,0.620) (-1.880,0.652) (-1.875,0.700) (-1.871,0.740) (-1.867,0.780) (-1.862,0.820) (-1.857,0.860) (-1.852,0.900) (-1.846,0.940) (-1.840,0.980) (-1.835,1.020) (-1.829,1.060) (-1.822,1.100) (-1.816,1.140) (-1.810,1.180) (-1.803,1.220) (-1.797,1.260) (-1.790,1.300) (-1.784,1.340) (-1.777,1.380) (-1.771,1.420) (-1.764,1.460) (-1.758,1.500) (-1.752,1.540) (-1.745,1.580) (-1.740,1.614) (-1.733,1.660) (-1.727,1.700) (-1.721,1.740) (-1.715,1.780) (-1.709,1.820) (-1.704,1.860) (-1.698,1.900) (-1.693,1.940)}
\EQ{(-1.945,-1.940) (-1.949,-1.900) (-1.952,-1.860) (-1.955,-1.820) (-1.959,-1.780) (-1.962,-1.740) (-1.966,-1.700) (-1.970,-1.660) (-1.974,-1.620) (-1.977,-1.580) (-1.981,-1.540) (-1.985,-1.500) (-1.989,-1.460) (-1.993,-1.420) (-1.998,-1.380) (-2.002,-1.340) (-2.006,-1.300) (-2.010,-1.260) (-2.015,-1.220) (-2.019,-1.180) (-2.024,-1.140) (-2.028,-1.100) (-2.033,-1.060) (-2.037,-1.020) (-2.042,-0.980) (-2.046,-0.940) (-2.051,-0.900) (-2.055,-0.860) (-2.059,-0.820) (-2.064,-0.780) (-2.068,-0.740) (-2.073,-0.700) (-2.077,-0.660) (-2.081,-0.620) (-2.085,-0.580) (-2.089,-0.540) (-2.093,-0.500) (-2.097,-0.460) (-2.101,-0.420) (-2.105,-0.380) (-2.108,-0.340) (-2.111,-0.300) (-2.115,-0.260) (-2.118,-0.220) (-2.120,-0.185) (-2.123,-0.140) (-2.125,-0.100) (-2.127,-0.060) (-2.129,-0.020) (-2.131,0.020) (-2.133,0.060) (-2.134,0.100) (-2.135,0.140) (-2.135,0.180) (-2.136,0.220) (-2.136,0.260) (-2.136,0.300) (-2.136,0.340) (-2.135,0.380) (-2.134,0.420) (-2.133,0.460) (-2.132,0.500) (-2.130,0.540) (-2.128,0.580) (-2.126,0.620) (-2.123,0.660) (-2.121,0.700) (-2.118,0.740) (-2.115,0.780) (-2.111,0.820) (-2.108,0.860) (-2.104,0.900) (-2.100,0.940) (-2.096,0.980) (-2.092,1.020) (-2.088,1.060) (-2.084,1.100) (-2.080,1.134) (-2.075,1.180) (-2.070,1.220) (-2.065,1.260) (-2.061,1.300) (-2.056,1.340) (-2.051,1.380) (-2.046,1.420) (-2.041,1.460) (-2.036,1.500) (-2.032,1.540) (-2.027,1.580) (-2.022,1.620) (-2.017,1.660) (-2.012,1.700) (-2.008,1.740) (-2.003,1.780) (-1.998,1.820) (-1.994,1.860) (-1.989,1.900) (-1.985,1.940)}
\EQ{(-2.238,-1.940) (-2.241,-1.900) (-2.244,-1.860) (-2.247,-1.820) (-2.250,-1.780) (-2.253,-1.740) (-2.256,-1.700) (-2.259,-1.660) (-2.262,-1.620) (-2.265,-1.580) (-2.269,-1.540) (-2.272,-1.500) (-2.275,-1.460) (-2.279,-1.420) (-2.282,-1.380) (-2.285,-1.340) (-2.289,-1.300) (-2.292,-1.260) (-2.296,-1.220) (-2.299,-1.180) (-2.303,-1.140) (-2.306,-1.100) (-2.310,-1.060) (-2.313,-1.020) (-2.316,-0.980) (-2.320,-0.940) (-2.323,-0.900) (-2.327,-0.860) (-2.330,-0.820) (-2.333,-0.780) (-2.337,-0.740) (-2.340,-0.700) (-2.340,-0.699)}
\EQ{(-2.340,1.188) (-2.337,1.220) (-2.334,1.260) (-2.330,1.300) (-2.327,1.340) (-2.323,1.380) (-2.320,1.415) (-2.316,1.460) (-2.312,1.500) (-2.308,1.540) (-2.305,1.580) (-2.301,1.620) (-2.297,1.660) (-2.294,1.700) (-2.290,1.740) (-2.286,1.780) (-2.283,1.820) (-2.279,1.860) (-2.275,1.900) (-2.272,1.940)}

      % Conductor neutro en campo uniforme - líneas
% Generado por scripts/generar_datos.py (no editar a mano).
\LC{(-2.300,-1.850) (-2.264,-1.847) (-2.228,-1.845) (-2.192,-1.842) (-2.156,-1.839) (-2.121,-1.837) (-2.085,-1.834) (-2.049,-1.831) (-2.013,-1.828) (-1.977,-1.825) (-1.941,-1.822) (-1.905,-1.818) (-1.869,-1.815) (-1.834,-1.812) (-1.798,-1.809) (-1.762,-1.805) (-1.726,-1.802) (-1.690,-1.798) (-1.654,-1.795) (-1.619,-1.791) (-1.583,-1.787) (-1.547,-1.784) (-1.511,-1.780) (-1.475,-1.776) (-1.440,-1.772) (-1.404,-1.768) (-1.368,-1.764) (-1.332,-1.760) (-1.296,-1.756) (-1.261,-1.752) (-1.225,-1.747) (-1.189,-1.743) (-1.153,-1.739) (-1.118,-1.735) (-1.082,-1.730) (-1.046,-1.726) (-1.011,-1.722) (-0.975,-1.717) (-0.939,-1.713) (-0.903,-1.709) (-0.868,-1.704) (-0.832,-1.700) (-0.796,-1.696) (-0.760,-1.692) (-0.725,-1.687) (-0.689,-1.683) (-0.653,-1.679) (-0.617,-1.675) (-0.581,-1.671) (-0.546,-1.668) (-0.510,-1.664) (-0.474,-1.661) (-0.438,-1.657) (-0.402,-1.654) (-0.366,-1.651) (-0.331,-1.648) (-0.295,-1.646) (-0.259,-1.644) (-0.223,-1.642) (-0.187,-1.640) (-0.151,-1.638) (-0.115,-1.637) (-0.079,-1.636) (-0.043,-1.635) (-0.007,-1.635) (0.029,-1.635) (0.065,-1.635) (0.101,-1.636) (0.137,-1.637) (0.173,-1.638) (0.209,-1.639) (0.245,-1.641) (0.281,-1.643) (0.317,-1.645) (0.353,-1.647) (0.389,-1.650) (0.425,-1.652) (0.461,-1.655) (0.496,-1.658) (0.532,-1.661) (0.568,-1.664) (0.604,-1.668) (0.640,-1.671) (0.676,-1.675) (0.711,-1.678) (0.747,-1.682) (0.783,-1.686) (0.819,-1.689) (0.855,-1.693) (0.891,-1.697) (0.926,-1.701) (0.962,-1.705) (0.998,-1.709) (1.034,-1.713) (1.069,-1.717) (1.105,-1.720) (1.141,-1.724) (1.177,-1.728) (1.213,-1.732) (1.248,-1.736) (1.284,-1.740) (1.320,-1.743) (1.356,-1.747) (1.392,-1.751) (1.427,-1.755) (1.463,-1.758) (1.499,-1.762) (1.535,-1.765) (1.571,-1.769) (1.607,-1.772) (1.642,-1.776) (1.678,-1.779) (1.714,-1.782) (1.750,-1.786) (1.786,-1.789) (1.822,-1.792) (1.858,-1.795) (1.893,-1.798) (1.929,-1.801) (1.965,-1.804) (2.001,-1.807) (2.037,-1.810) (2.073,-1.812) (2.109,-1.815) (2.145,-1.818) (2.181,-1.820) (2.216,-1.823) (2.252,-1.825) (2.288,-1.828) (2.300,-1.829)}
\LC{(-2.300,-1.645) (-2.264,-1.642) (-2.228,-1.639) (-2.192,-1.636) (-2.156,-1.633) (-2.121,-1.630) (-2.085,-1.627) (-2.049,-1.624) (-2.013,-1.621) (-1.977,-1.617) (-1.941,-1.614) (-1.906,-1.610) (-1.870,-1.607) (-1.834,-1.603) (-1.798,-1.599) (-1.762,-1.595) (-1.727,-1.591) (-1.691,-1.587) (-1.655,-1.583) (-1.619,-1.579) (-1.583,-1.575) (-1.548,-1.570) (-1.512,-1.566) (-1.476,-1.561) (-1.441,-1.557) (-1.405,-1.552) (-1.369,-1.547) (-1.334,-1.542) (-1.298,-1.538) (-1.262,-1.532) (-1.227,-1.527) (-1.191,-1.522) (-1.155,-1.517) (-1.120,-1.512) (-1.084,-1.506) (-1.049,-1.501) (-1.013,-1.495) (-0.977,-1.490) (-0.942,-1.484) (-0.906,-1.479) (-0.871,-1.473) (-0.835,-1.467) (-0.800,-1.462) (-0.764,-1.456) (-0.729,-1.450) (-0.693,-1.445) (-0.657,-1.439) (-0.622,-1.433) (-0.586,-1.428) (-0.551,-1.422) (-0.515,-1.417) (-0.480,-1.412) (-0.444,-1.407) (-0.408,-1.402) (-0.373,-1.397) (-0.337,-1.393) (-0.301,-1.389) (-0.265,-1.385) (-0.229,-1.382) (-0.193,-1.379) (-0.158,-1.377) (-0.122,-1.375) (-0.086,-1.373) (-0.050,-1.372) (-0.014,-1.372) (0.022,-1.372) (0.058,-1.372) (0.094,-1.374) (0.130,-1.375) (0.166,-1.377) (0.202,-1.380) (0.238,-1.382) (0.274,-1.386) (0.310,-1.389) (0.346,-1.393) (0.381,-1.397) (0.417,-1.401) (0.453,-1.405) (0.489,-1.410) (0.524,-1.414) (0.560,-1.419) (0.596,-1.424) (0.631,-1.429) (0.667,-1.434) (0.703,-1.439) (0.738,-1.444) (0.774,-1.449) (0.810,-1.454) (0.845,-1.459) (0.881,-1.464) (0.916,-1.469) (0.952,-1.474) (0.988,-1.479) (1.023,-1.484) (1.059,-1.489) (1.095,-1.494) (1.130,-1.499) (1.166,-1.504) (1.202,-1.509) (1.237,-1.513) (1.273,-1.518) (1.309,-1.523) (1.344,-1.527) (1.380,-1.532) (1.416,-1.536) (1.452,-1.541) (1.487,-1.545) (1.523,-1.549) (1.559,-1.553) (1.595,-1.557) (1.630,-1.561) (1.666,-1.565) (1.702,-1.569) (1.738,-1.573) (1.774,-1.576) (1.809,-1.580) (1.845,-1.584) (1.881,-1.587) (1.917,-1.590) (1.953,-1.594) (1.989,-1.597) (2.025,-1.600) (2.060,-1.603) (2.096,-1.606) (2.132,-1.609) (2.168,-1.612) (2.204,-1.615) (2.240,-1.618) (2.276,-1.621) (2.300,-1.623)}
\LC{(-2.300,-1.440) (-2.264,-1.437) (-2.228,-1.434) (-2.192,-1.431) (-2.157,-1.428) (-2.121,-1.424) (-2.085,-1.421) (-2.049,-1.417) (-2.013,-1.414) (-1.977,-1.410) (-1.942,-1.406) (-1.906,-1.402) (-1.870,-1.398) (-1.834,-1.394) (-1.799,-1.390) (-1.763,-1.386) (-1.727,-1.381) (-1.691,-1.377) (-1.656,-1.372) (-1.620,-1.367) (-1.584,-1.362) (-1.549,-1.357) (-1.513,-1.352) (-1.477,-1.347) (-1.442,-1.341) (-1.406,-1.336) (-1.371,-1.330) (-1.335,-1.324) (-1.300,-1.318) (-1.264,-1.312) (-1.229,-1.306) (-1.193,-1.300) (-1.158,-1.294) (-1.122,-1.287) (-1.087,-1.280) (-1.052,-1.273) (-1.016,-1.267) (-0.981,-1.259) (-0.946,-1.252) (-0.911,-1.245) (-0.875,-1.237) (-0.840,-1.230) (-0.805,-1.222) (-0.770,-1.214) (-0.735,-1.207) (-0.700,-1.199) (-0.665,-1.191) (-0.629,-1.182) (-0.594,-1.174) (-0.559,-1.166) (-0.524,-1.157) (-0.489,-1.149) (-0.454,-1.141) (-0.419,-1.132) (-0.384,-1.124) (-0.349,-1.116) (-0.314,-1.108) (-0.279,-1.100) (-0.244,-1.093) (-0.208,-1.086) (-0.173,-1.080) (-0.137,-1.076) (-0.101,-1.072) (-0.065,-1.070) (-0.029,-1.069) (0.007,-1.070) (0.043,-1.072) (0.078,-1.076) (0.114,-1.080) (0.150,-1.085) (0.185,-1.091) (0.221,-1.097) (0.256,-1.103) (0.292,-1.110) (0.327,-1.116) (0.362,-1.123) (0.398,-1.130) (0.433,-1.137) (0.468,-1.145) (0.503,-1.152) (0.539,-1.159) (0.574,-1.166) (0.609,-1.174) (0.645,-1.181) (0.680,-1.188) (0.715,-1.195) (0.750,-1.202) (0.786,-1.209) (0.821,-1.216) (0.856,-1.223) (0.892,-1.230) (0.927,-1.236) (0.963,-1.243) (0.998,-1.249) (1.033,-1.256) (1.069,-1.262) (1.104,-1.268) (1.140,-1.274) (1.175,-1.280) (1.211,-1.286) (1.246,-1.292) (1.282,-1.298) (1.317,-1.303) (1.353,-1.309) (1.389,-1.314) (1.424,-1.319) (1.460,-1.324) (1.495,-1.330) (1.531,-1.334) (1.567,-1.339) (1.603,-1.344) (1.638,-1.348) (1.674,-1.353) (1.710,-1.357) (1.745,-1.362) (1.781,-1.366) (1.817,-1.370) (1.853,-1.374) (1.889,-1.378) (1.924,-1.381) (1.960,-1.385) (1.996,-1.389) (2.032,-1.392) (2.068,-1.396) (2.103,-1.399) (2.139,-1.402) (2.175,-1.405) (2.211,-1.408) (2.247,-1.411) (2.283,-1.414) (2.300,-1.416)}
\LC{(-2.300,-1.235) (-2.264,-1.232) (-2.228,-1.229) (-2.192,-1.225) (-2.157,-1.222) (-2.121,-1.218) (-2.085,-1.215) (-2.049,-1.211) (-2.013,-1.207) (-1.978,-1.203) (-1.942,-1.199) (-1.906,-1.195) (-1.870,-1.190) (-1.835,-1.186) (-1.799,-1.181) (-1.763,-1.176) (-1.728,-1.172) (-1.692,-1.166) (-1.656,-1.161) (-1.621,-1.156) (-1.585,-1.150) (-1.550,-1.145) (-1.514,-1.139) (-1.479,-1.133) (-1.443,-1.126) (-1.408,-1.120) (-1.372,-1.113) (-1.337,-1.107) (-1.302,-1.100) (-1.266,-1.093) (-1.231,-1.085) (-1.196,-1.078) (-1.161,-1.070) (-1.126,-1.062) (-1.091,-1.054) (-1.056,-1.045) (-1.021,-1.037) (-0.986,-1.028) (-0.951,-1.019) (-0.916,-1.009) (-0.882,-1.000) (-0.847,-0.990) (-0.812,-0.980) (-0.778,-0.969) (-0.744,-0.959) (-0.709,-0.948) (-0.675,-0.936) (-0.641,-0.925) (-0.607,-0.913) (-0.573,-0.900) (-0.540,-0.887) (-0.506,-0.874) (-0.473,-0.860) (-0.440,-0.846) (-0.407,-0.830) (-0.375,-0.814) (-0.344,-0.796) (-0.313,-0.778) (-0.284,-0.757) (-0.256,-0.734) (-0.229,-0.710)}
\LC{(-2.300,-1.030) (-2.264,-1.027) (-2.228,-1.023) (-2.192,-1.020) (-2.157,-1.017) (-2.121,-1.013) (-2.085,-1.009) (-2.049,-1.005) (-2.013,-1.001) (-1.978,-0.997) (-1.942,-0.993) (-1.906,-0.988) (-1.871,-0.983) (-1.835,-0.979) (-1.799,-0.974) (-1.764,-0.968) (-1.728,-0.963) (-1.692,-0.958) (-1.657,-0.952) (-1.621,-0.946) (-1.586,-0.940) (-1.550,-0.934) (-1.515,-0.927) (-1.480,-0.920) (-1.444,-0.913) (-1.409,-0.906) (-1.374,-0.899) (-1.339,-0.891) (-1.304,-0.883) (-1.269,-0.874) (-1.234,-0.866) (-1.199,-0.857) (-1.164,-0.848) (-1.129,-0.838) (-1.095,-0.829) (-1.060,-0.818) (-1.026,-0.808) (-0.991,-0.797) (-0.957,-0.786) (-0.923,-0.774) (-0.889,-0.763) (-0.855,-0.750) (-0.821,-0.738) (-0.788,-0.725) (-0.754,-0.711) (-0.721,-0.697) (-0.688,-0.683) (-0.655,-0.668) (-0.623,-0.653) (-0.590,-0.638) (-0.558,-0.622) (-0.526,-0.605) (-0.494,-0.589) (-0.462,-0.572) (-0.430,-0.555) (-0.399,-0.538) (-0.367,-0.520)}
\LC{(-2.300,-0.825) (-2.264,-0.822) (-2.228,-0.819) (-2.192,-0.815) (-2.157,-0.812) (-2.121,-0.808) (-2.085,-0.804) (-2.049,-0.800) (-2.013,-0.796) (-1.978,-0.792) (-1.942,-0.787) (-1.906,-0.783) (-1.871,-0.778) (-1.835,-0.773) (-1.799,-0.768) (-1.764,-0.763) (-1.728,-0.757) (-1.693,-0.751) (-1.657,-0.745) (-1.622,-0.739) (-1.586,-0.732) (-1.551,-0.725) (-1.516,-0.718) (-1.480,-0.711) (-1.445,-0.703) (-1.410,-0.695) (-1.375,-0.687) (-1.340,-0.678) (-1.305,-0.669) (-1.271,-0.660) (-1.236,-0.650) (-1.201,-0.640) (-1.167,-0.630) (-1.133,-0.619) (-1.098,-0.607) (-1.064,-0.596) (-1.031,-0.583) (-0.997,-0.571) (-0.963,-0.557) (-0.930,-0.544) (-0.897,-0.529) (-0.864,-0.515) (-0.832,-0.499) (-0.799,-0.484) (-0.767,-0.468) (-0.735,-0.451) (-0.703,-0.434) (-0.672,-0.416) (-0.641,-0.398) (-0.610,-0.380) (-0.579,-0.362) (-0.548,-0.343) (-0.518,-0.323) (-0.489,-0.302)}
\LC{(-2.300,-0.620) (-2.264,-0.617) (-2.228,-0.614) (-2.192,-0.611) (-2.157,-0.607) (-2.121,-0.604) (-2.085,-0.600) (-2.049,-0.596) (-2.013,-0.593) (-1.978,-0.588) (-1.942,-0.584) (-1.906,-0.580) (-1.870,-0.575) (-1.835,-0.570) (-1.799,-0.565) (-1.764,-0.560) (-1.728,-0.554) (-1.692,-0.548) (-1.657,-0.542) (-1.622,-0.536) (-1.586,-0.529) (-1.551,-0.522) (-1.516,-0.515) (-1.480,-0.507) (-1.445,-0.499) (-1.410,-0.490) (-1.375,-0.482) (-1.341,-0.472) (-1.306,-0.463) (-1.271,-0.452) (-1.237,-0.442) (-1.203,-0.430) (-1.169,-0.419) (-1.135,-0.406) (-1.101,-0.394) (-1.068,-0.380) (-1.035,-0.366) (-1.002,-0.351) (-0.970,-0.335) (-0.938,-0.319) (-0.906,-0.302) (-0.875,-0.284) (-0.844,-0.266) (-0.813,-0.247) (-0.783,-0.226) (-0.754,-0.205) (-0.725,-0.184) (-0.697,-0.161) (-0.670,-0.138) (-0.643,-0.114) (-0.635,-0.105)}
\LC{(-2.300,-0.415) (-2.264,-0.412) (-2.228,-0.410) (-2.192,-0.407) (-2.156,-0.404) (-2.121,-0.401) (-2.085,-0.398) (-2.049,-0.394) (-2.013,-0.391) (-1.977,-0.387) (-1.941,-0.383) (-1.906,-0.379) (-1.870,-0.375) (-1.834,-0.370) (-1.799,-0.365) (-1.763,-0.360) (-1.727,-0.355) (-1.692,-0.350) (-1.656,-0.344) (-1.621,-0.338) (-1.585,-0.332) (-1.550,-0.325) (-1.515,-0.318) (-1.479,-0.310) (-1.444,-0.302) (-1.409,-0.294) (-1.374,-0.285) (-1.340,-0.275) (-1.305,-0.265) (-1.271,-0.255) (-1.236,-0.244) (-1.202,-0.232) (-1.169,-0.219) (-1.135,-0.206) (-1.102,-0.191) (-1.070,-0.176) (-1.037,-0.160) (-1.006,-0.143) (-0.975,-0.125) (-0.944,-0.106) (-0.915,-0.085) (-0.886,-0.063) (-0.859,-0.040) (-0.833,-0.015) (-0.807,0.011) (-0.791,0.028)}
\LC{(-2.300,-0.210) (-2.264,-0.208) (-2.228,-0.206) (-2.192,-0.204) (-2.156,-0.201) (-2.120,-0.199) (-2.084,-0.196) (-2.049,-0.194) (-2.013,-0.191) (-1.977,-0.188) (-1.941,-0.185) (-1.905,-0.181) (-1.869,-0.178) (-1.833,-0.174) (-1.798,-0.170) (-1.762,-0.166) (-1.726,-0.162) (-1.690,-0.157) (-1.655,-0.152) (-1.619,-0.147) (-1.584,-0.142) (-1.548,-0.136) (-1.513,-0.130) (-1.477,-0.123) (-1.442,-0.116) (-1.407,-0.108) (-1.372,-0.100) (-1.337,-0.092) (-1.302,-0.082) (-1.267,-0.072) (-1.233,-0.061) (-1.199,-0.050) (-1.165,-0.037) (-1.132,-0.023) (-1.099,-0.009) (-1.067,0.007) (-1.035,0.025) (-1.005,0.044) (-0.975,0.064) (-0.947,0.087) (-0.920,0.111) (-0.912,0.119)}
\LC{(-2.300,-0.005) (-2.264,-0.004) (-2.228,-0.002) (-2.192,-0.001) (-2.156,0.000) (-2.120,0.002) (-2.084,0.003) (-2.048,0.005) (-2.012,0.007) (-1.976,0.009) (-1.940,0.011) (-1.904,0.013) (-1.868,0.015) (-1.833,0.018) (-1.797,0.020) (-1.761,0.023) (-1.725,0.026) (-1.689,0.029) (-1.653,0.032) (-1.617,0.036) (-1.582,0.039) (-1.546,0.043) (-1.510,0.048) (-1.474,0.053) (-1.439,0.058) (-1.403,0.063) (-1.368,0.069) (-1.332,0.076) (-1.297,0.083) (-1.262,0.091) (-1.227,0.100) (-1.192,0.109) (-1.158,0.120) (-1.124,0.132) (-1.090,0.145) (-1.058,0.160) (-1.026,0.177) (-0.997,0.198) (-0.988,0.206)}
\LC{(-2.300,0.200) (-2.264,0.200) (-2.228,0.201) (-2.192,0.201) (-2.156,0.201) (-2.120,0.202) (-2.084,0.202) (-2.048,0.203) (-2.012,0.203) (-1.976,0.204) (-1.940,0.204) (-1.904,0.205) (-1.868,0.205) (-1.832,0.206) (-1.796,0.207) (-1.760,0.208) (-1.724,0.209) (-1.688,0.210) (-1.652,0.211) (-1.616,0.212) (-1.580,0.213) (-1.544,0.215) (-1.508,0.216) (-1.472,0.218) (-1.436,0.220) (-1.400,0.222) (-1.364,0.224) (-1.329,0.227) (-1.293,0.230) (-1.257,0.234) (-1.221,0.238) (-1.185,0.242) (-1.150,0.248) (-1.114,0.254) (-1.079,0.262) (-1.044,0.271) (-1.021,0.276)}
\LC{(-2.300,0.405) (-2.264,0.404) (-2.228,0.404) (-2.192,0.403) (-2.156,0.402) (-2.120,0.401) (-2.084,0.400) (-2.048,0.400) (-2.012,0.399) (-1.976,0.398) (-1.940,0.397) (-1.904,0.396) (-1.868,0.394) (-1.832,0.393) (-1.796,0.392) (-1.760,0.391) (-1.724,0.390) (-1.688,0.388) (-1.652,0.387) (-1.616,0.385) (-1.580,0.384) (-1.544,0.382) (-1.508,0.380) (-1.472,0.378) (-1.437,0.376) (-1.401,0.374) (-1.365,0.372) (-1.329,0.370) (-1.293,0.367) (-1.257,0.365) (-1.221,0.362) (-1.185,0.359) (-1.149,0.356) (-1.113,0.352) (-1.078,0.349) (-1.042,0.342) (-1.031,0.337)}
\LC{(-2.300,0.610) (-2.264,0.608) (-2.228,0.607) (-2.192,0.605) (-2.156,0.603) (-2.120,0.601) (-2.084,0.599) (-2.048,0.597) (-2.012,0.595) (-1.976,0.592) (-1.941,0.590) (-1.905,0.587) (-1.869,0.585) (-1.833,0.582) (-1.797,0.579) (-1.761,0.575) (-1.725,0.572) (-1.689,0.568) (-1.654,0.565) (-1.618,0.561) (-1.582,0.556) (-1.546,0.552) (-1.511,0.547) (-1.475,0.542) (-1.440,0.536) (-1.404,0.530) (-1.369,0.524) (-1.333,0.517) (-1.298,0.510) (-1.263,0.502) (-1.228,0.493) (-1.193,0.484) (-1.159,0.473) (-1.125,0.462) (-1.091,0.449) (-1.058,0.435) (-1.027,0.417) (-1.017,0.410)}
\LC{(-2.300,0.815) (-2.264,0.813) (-2.228,0.810) (-2.192,0.808) (-2.156,0.805) (-2.120,0.802) (-2.085,0.799) (-2.049,0.796) (-2.013,0.792) (-1.977,0.789) (-1.941,0.785) (-1.905,0.782) (-1.870,0.778) (-1.834,0.773) (-1.798,0.769) (-1.762,0.764) (-1.727,0.759) (-1.691,0.754) (-1.656,0.749) (-1.620,0.743) (-1.585,0.737) (-1.549,0.730) (-1.514,0.723) (-1.479,0.716) (-1.444,0.708) (-1.409,0.699) (-1.374,0.690) (-1.339,0.681) (-1.305,0.670) (-1.270,0.659) (-1.236,0.647) (-1.203,0.635) (-1.169,0.621) (-1.137,0.606) (-1.105,0.590) (-1.073,0.572) (-1.043,0.553) (-1.013,0.532) (-0.985,0.509) (-0.959,0.485) (-0.950,0.477)}
\LC{(-2.300,1.020) (-2.264,1.017) (-2.228,1.014) (-2.192,1.011) (-2.157,1.007) (-2.121,1.004) (-2.085,1.000) (-2.049,0.996) (-2.013,0.992) (-1.978,0.988) (-1.942,0.984) (-1.906,0.979) (-1.870,0.975) (-1.835,0.969) (-1.799,0.964) (-1.764,0.959) (-1.728,0.953) (-1.693,0.947) (-1.657,0.940) (-1.622,0.933) (-1.587,0.926) (-1.551,0.919) (-1.516,0.911) (-1.481,0.902) (-1.446,0.893) (-1.412,0.884) (-1.377,0.874) (-1.343,0.863) (-1.309,0.852) (-1.275,0.840) (-1.241,0.827) (-1.208,0.813) (-1.175,0.799) (-1.142,0.783) (-1.110,0.767) (-1.079,0.749) (-1.048,0.730) (-1.018,0.710) (-0.990,0.688) (-0.962,0.665) (-0.936,0.641) (-0.911,0.614) (-0.889,0.586) (-0.869,0.556) (-0.856,0.536)}
\LC{(-2.300,1.225) (-2.264,1.222) (-2.228,1.218) (-2.192,1.215) (-2.157,1.211) (-2.121,1.207) (-2.085,1.203) (-2.049,1.199) (-2.014,1.195) (-1.978,1.190) (-1.942,1.185) (-1.907,1.180) (-1.871,1.175) (-1.835,1.170) (-1.800,1.164) (-1.764,1.158) (-1.729,1.152) (-1.693,1.146) (-1.658,1.139) (-1.623,1.132) (-1.587,1.125) (-1.552,1.117) (-1.517,1.109) (-1.482,1.100) (-1.447,1.091) (-1.413,1.082) (-1.378,1.072) (-1.344,1.062) (-1.309,1.051) (-1.275,1.039) (-1.241,1.027) (-1.208,1.014) (-1.174,1.000) (-1.141,0.986) (-1.109,0.971) (-1.076,0.955) (-1.045,0.938) (-1.013,0.921) (-0.983,0.902) (-0.952,0.882) (-0.923,0.861) (-0.895,0.839) (-0.868,0.815) (-0.841,0.790) (-0.817,0.764) (-0.793,0.737) (-0.772,0.708) (-0.752,0.678) (-0.735,0.647) (-0.721,0.613) (-0.718,0.602)}
\LC{(-2.300,1.430) (-2.264,1.427) (-2.228,1.423) (-2.193,1.419) (-2.157,1.416) (-2.121,1.412) (-2.085,1.408) (-2.049,1.403) (-2.014,1.399) (-1.978,1.394) (-1.942,1.390) (-1.907,1.385) (-1.871,1.379) (-1.835,1.374) (-1.800,1.369) (-1.764,1.363) (-1.729,1.357) (-1.693,1.350) (-1.658,1.344) (-1.623,1.337) (-1.587,1.330) (-1.552,1.323) (-1.517,1.315) (-1.482,1.307) (-1.447,1.299) (-1.412,1.290) (-1.377,1.281) (-1.342,1.272) (-1.308,1.262) (-1.273,1.252) (-1.239,1.241) (-1.205,1.230) (-1.170,1.218) (-1.137,1.206) (-1.103,1.193) (-1.069,1.180) (-1.036,1.166) (-1.003,1.151) (-0.971,1.136) (-0.938,1.120) (-0.906,1.104) (-0.875,1.086) (-0.844,1.068) (-0.813,1.049) (-0.783,1.029) (-0.754,1.008) (-0.725,0.987) (-0.697,0.964) (-0.670,0.940) (-0.644,0.915) (-0.620,0.889) (-0.596,0.862) (-0.574,0.833) (-0.554,0.803) (-0.536,0.772) (-0.520,0.740) (-0.506,0.707) (-0.500,0.684)}
\LC{(-2.300,1.635) (-2.264,1.632) (-2.228,1.628) (-2.193,1.624) (-2.157,1.621) (-2.121,1.617) (-2.085,1.613) (-2.049,1.609) (-2.014,1.604) (-1.978,1.600) (-1.942,1.595) (-1.907,1.591) (-1.871,1.586) (-1.835,1.581) (-1.800,1.576) (-1.764,1.570) (-1.728,1.565) (-1.693,1.559) (-1.657,1.553) (-1.622,1.547) (-1.587,1.540) (-1.551,1.534) (-1.516,1.527) (-1.480,1.520) (-1.445,1.512) (-1.410,1.505) (-1.375,1.497) (-1.340,1.489) (-1.305,1.480) (-1.270,1.472) (-1.235,1.463) (-1.200,1.453) (-1.166,1.444) (-1.131,1.434) (-1.096,1.424) (-1.062,1.413) (-1.028,1.402) (-0.993,1.391) (-0.959,1.380) (-0.925,1.368) (-0.892,1.355) (-0.858,1.343) (-0.824,1.330) (-0.791,1.316) (-0.758,1.302) (-0.725,1.288) (-0.692,1.274) (-0.659,1.259) (-0.626,1.244) (-0.594,1.228) (-0.562,1.212) (-0.529,1.196) (-0.498,1.179) (-0.466,1.162) (-0.434,1.145) (-0.403,1.127) (-0.371,1.110) (-0.340,1.092) (-0.309,1.073) (-0.278,1.055) (-0.247,1.037) (-0.216,1.019) (-0.184,1.001) (-0.153,0.984) (-0.120,0.968) (-0.087,0.954) (-0.053,0.944) (-0.018,0.938) (0.018,0.938) (0.054,0.945) (0.087,0.957) (0.120,0.973) (0.152,0.990) (0.182,1.009) (0.213,1.028) (0.243,1.047) (0.273,1.067) (0.304,1.086) (0.334,1.106) (0.364,1.125) (0.395,1.144) (0.426,1.163) (0.457,1.181) (0.488,1.199) (0.519,1.217) (0.551,1.234) (0.582,1.251) (0.614,1.268) (0.646,1.284) (0.679,1.300) (0.711,1.315) (0.744,1.330) (0.777,1.345) (0.810,1.359) (0.844,1.372) (0.877,1.385) (0.911,1.398) (0.945,1.411) (0.978,1.423) (1.013,1.434) (1.047,1.446) (1.081,1.456) (1.116,1.467) (1.150,1.477) (1.185,1.487) (1.219,1.496) (1.254,1.505) (1.289,1.514) (1.324,1.523) (1.359,1.531) (1.394,1.539) (1.429,1.547) (1.465,1.554) (1.500,1.561) (1.535,1.568) (1.571,1.575) (1.606,1.581) (1.642,1.587) (1.677,1.593) (1.713,1.599) (1.748,1.605) (1.784,1.610) (1.819,1.615) (1.855,1.620) (1.891,1.625) (1.926,1.630) (1.962,1.634) (1.998,1.638) (2.034,1.643) (2.069,1.647) (2.105,1.651) (2.141,1.654) (2.177,1.658) (2.213,1.662) (2.248,1.665) (2.284,1.668) (2.300,1.671)}
\LC{(-2.300,1.840) (-2.264,1.837) (-2.228,1.833) (-2.192,1.830) (-2.157,1.826) (-2.121,1.823) (-2.085,1.819) (-2.049,1.815) (-2.013,1.811) (-1.978,1.807) (-1.942,1.803) (-1.906,1.798) (-1.870,1.794) (-1.835,1.789) (-1.799,1.784) (-1.763,1.780) (-1.728,1.775) (-1.692,1.769) (-1.657,1.764) (-1.621,1.759) (-1.585,1.753) (-1.550,1.747) (-1.514,1.741) (-1.479,1.735) (-1.443,1.729) (-1.408,1.722) (-1.373,1.716) (-1.337,1.709) (-1.302,1.702) (-1.267,1.695) (-1.231,1.688) (-1.196,1.680) (-1.161,1.673) (-1.126,1.665) (-1.091,1.657) (-1.056,1.649) (-1.021,1.640) (-0.986,1.632) (-0.951,1.623) (-0.916,1.615) (-0.881,1.606) (-0.846,1.597) (-0.811,1.588) (-0.776,1.578) (-0.742,1.569) (-0.707,1.560) (-0.672,1.551) (-0.637,1.541) (-0.603,1.532) (-0.568,1.523) (-0.533,1.514) (-0.498,1.505) (-0.463,1.496) (-0.428,1.487) (-0.393,1.479) (-0.358,1.471) (-0.323,1.464) (-0.287,1.457) (-0.252,1.451) (-0.216,1.445) (-0.181,1.441) (-0.145,1.437) (-0.109,1.434) (-0.073,1.432) (-0.037,1.431) (-0.001,1.431) (0.035,1.433) (0.071,1.435) (0.107,1.439) (0.142,1.443) (0.178,1.448) (0.213,1.455) (0.249,1.461) (0.284,1.469) (0.319,1.477) (0.354,1.486) (0.389,1.495) (0.424,1.504) (0.458,1.514) (0.493,1.523) (0.528,1.533) (0.562,1.543) (0.597,1.554) (0.631,1.564) (0.666,1.574) (0.700,1.584) (0.735,1.594) (0.770,1.604) (0.804,1.614) (0.839,1.623) (0.874,1.633) (0.908,1.642) (0.943,1.651) (0.978,1.660) (1.013,1.669) (1.048,1.678) (1.083,1.686) (1.118,1.695) (1.153,1.703) (1.188,1.711) (1.223,1.718) (1.258,1.726) (1.294,1.733) (1.329,1.740) (1.364,1.747) (1.400,1.754) (1.435,1.761) (1.470,1.767) (1.506,1.773) (1.541,1.779) (1.577,1.785) (1.612,1.791) (1.648,1.796) (1.684,1.802) (1.719,1.807) (1.755,1.812) (1.791,1.817) (1.826,1.822) (1.862,1.826) (1.898,1.831) (1.933,1.835) (1.969,1.839) (2.005,1.844) (2.041,1.848) (2.076,1.851) (2.112,1.855) (2.148,1.859) (2.184,1.862) (2.220,1.866) (2.256,1.869) (2.291,1.872) (2.300,1.873)}
\LC{(0.050,-0.713) (0.077,-0.737) (0.107,-0.757) (0.138,-0.775) (0.170,-0.792) (0.203,-0.807) (0.236,-0.821) (0.269,-0.834) (0.303,-0.846) (0.337,-0.859) (0.371,-0.870) (0.405,-0.882) (0.439,-0.893) (0.474,-0.904) (0.508,-0.915) (0.542,-0.925) (0.577,-0.936) (0.612,-0.946) (0.646,-0.956) (0.681,-0.965) (0.716,-0.975) (0.750,-0.984) (0.785,-0.993) (0.820,-1.002) (0.855,-1.011) (0.890,-1.020) (0.925,-1.028) (0.960,-1.036) (0.995,-1.044) (1.030,-1.052) (1.065,-1.060) (1.100,-1.068) (1.136,-1.075) (1.171,-1.082) (1.206,-1.089) (1.242,-1.096) (1.277,-1.103) (1.312,-1.109) (1.348,-1.116) (1.383,-1.122) (1.419,-1.128) (1.454,-1.134) (1.490,-1.140) (1.525,-1.145) (1.561,-1.151) (1.597,-1.156) (1.632,-1.161) (1.668,-1.166) (1.703,-1.171) (1.739,-1.176) (1.775,-1.180) (1.811,-1.185) (1.846,-1.189) (1.882,-1.193) (1.918,-1.198) (1.954,-1.202) (1.989,-1.205) (2.025,-1.209) (2.061,-1.213) (2.097,-1.217) (2.133,-1.220) (2.169,-1.223) (2.204,-1.227) (2.240,-1.230) (2.276,-1.233) (2.300,-1.235)}
\LC{(0.193,-0.477) (0.224,-0.495) (0.257,-0.511) (0.288,-0.528) (0.320,-0.544) (0.353,-0.561) (0.385,-0.577) (0.417,-0.593) (0.450,-0.608) (0.482,-0.624) (0.515,-0.638) (0.548,-0.653) (0.581,-0.667) (0.614,-0.681) (0.648,-0.694) (0.681,-0.707) (0.715,-0.720) (0.749,-0.732) (0.783,-0.744) (0.817,-0.756) (0.851,-0.768) (0.885,-0.779) (0.919,-0.789) (0.954,-0.800) (0.988,-0.810) (1.023,-0.820) (1.058,-0.829) (1.093,-0.839) (1.127,-0.848) (1.162,-0.856) (1.197,-0.865) (1.232,-0.873) (1.267,-0.881) (1.303,-0.889) (1.338,-0.896) (1.373,-0.904) (1.408,-0.911) (1.444,-0.918) (1.479,-0.924) (1.515,-0.931) (1.550,-0.937) (1.585,-0.943) (1.621,-0.949) (1.657,-0.954) (1.692,-0.960) (1.728,-0.965) (1.763,-0.970) (1.799,-0.975) (1.835,-0.980) (1.870,-0.984) (1.906,-0.989) (1.942,-0.993) (1.978,-0.998) (2.013,-1.002) (2.049,-1.006) (2.085,-1.009) (2.121,-1.013) (2.157,-1.017) (2.192,-1.020) (2.228,-1.024) (2.264,-1.027) (2.300,-1.030)}
\LC{(0.341,-0.231) (0.372,-0.249) (0.401,-0.270) (0.430,-0.291) (0.460,-0.312) (0.490,-0.332) (0.520,-0.352) (0.550,-0.371) (0.581,-0.389) (0.612,-0.407) (0.644,-0.425) (0.675,-0.442) (0.707,-0.458) (0.740,-0.474) (0.772,-0.490) (0.805,-0.505) (0.838,-0.519) (0.871,-0.533) (0.904,-0.547) (0.938,-0.560) (0.972,-0.573) (1.005,-0.585) (1.040,-0.596) (1.074,-0.608) (1.108,-0.619) (1.142,-0.629) (1.177,-0.639) (1.212,-0.649) (1.246,-0.659) (1.281,-0.668) (1.316,-0.676) (1.351,-0.685) (1.386,-0.693) (1.421,-0.701) (1.457,-0.708) (1.492,-0.715) (1.527,-0.722) (1.562,-0.729) (1.598,-0.736) (1.633,-0.742) (1.669,-0.748) (1.704,-0.754) (1.740,-0.759) (1.776,-0.765) (1.811,-0.770) (1.847,-0.775) (1.882,-0.780) (1.918,-0.784) (1.954,-0.789) (1.990,-0.793) (2.025,-0.797) (2.061,-0.801) (2.097,-0.805) (2.133,-0.809) (2.169,-0.813) (2.204,-0.816) (2.240,-0.820) (2.276,-0.823) (2.300,-0.825)}
\LC{(0.525,-0.039) (0.550,-0.064) (0.576,-0.089) (0.603,-0.113) (0.630,-0.136) (0.658,-0.159) (0.687,-0.181) (0.716,-0.202) (0.746,-0.222) (0.776,-0.242) (0.807,-0.260) (0.838,-0.278) (0.870,-0.296) (0.902,-0.312) (0.934,-0.328) (0.967,-0.343) (1.000,-0.358) (1.033,-0.371) (1.066,-0.385) (1.100,-0.398) (1.134,-0.410) (1.168,-0.421) (1.202,-0.433) (1.237,-0.443) (1.271,-0.453) (1.306,-0.463) (1.340,-0.473) (1.375,-0.481) (1.410,-0.490) (1.445,-0.498) (1.480,-0.506) (1.516,-0.514) (1.551,-0.521) (1.586,-0.528) (1.622,-0.534) (1.657,-0.541) (1.693,-0.547) (1.728,-0.553) (1.764,-0.558) (1.799,-0.564) (1.835,-0.569) (1.871,-0.574) (1.906,-0.578) (1.942,-0.583) (1.978,-0.587) (2.013,-0.592) (2.049,-0.596) (2.085,-0.600) (2.121,-0.603) (2.157,-0.607) (2.192,-0.610) (2.228,-0.614) (2.264,-0.617) (2.300,-0.620)}
\LC{(0.694,0.096) (0.717,0.068) (0.741,0.042) (0.766,0.016) (0.793,-0.008) (0.821,-0.031) (0.849,-0.052) (0.879,-0.073) (0.909,-0.093) (0.940,-0.111) (0.971,-0.129) (1.003,-0.145) (1.036,-0.161) (1.068,-0.176) (1.102,-0.190) (1.135,-0.204) (1.169,-0.216) (1.203,-0.229) (1.237,-0.240) (1.271,-0.251) (1.306,-0.261) (1.340,-0.271) (1.375,-0.280) (1.410,-0.289) (1.445,-0.297) (1.480,-0.305) (1.515,-0.313) (1.550,-0.320) (1.586,-0.327) (1.621,-0.334) (1.657,-0.340) (1.692,-0.346) (1.728,-0.352) (1.763,-0.357) (1.799,-0.362) (1.835,-0.367) (1.870,-0.372) (1.906,-0.376) (1.942,-0.381) (1.977,-0.385) (2.013,-0.389) (2.049,-0.393) (2.085,-0.396) (2.121,-0.400) (2.157,-0.403) (2.192,-0.406) (2.228,-0.409) (2.264,-0.412) (2.300,-0.415)}
\LC{(0.814,0.196) (0.840,0.170) (0.866,0.146) (0.894,0.123) (0.922,0.101) (0.952,0.080) (0.982,0.061) (1.014,0.043) (1.046,0.026) (1.078,0.011) (1.111,-0.004) (1.144,-0.018) (1.178,-0.030) (1.212,-0.042) (1.246,-0.054) (1.280,-0.064) (1.315,-0.074) (1.350,-0.084) (1.385,-0.092) (1.420,-0.101) (1.455,-0.108) (1.490,-0.116) (1.525,-0.123) (1.561,-0.129) (1.596,-0.135) (1.632,-0.141) (1.667,-0.147) (1.703,-0.152) (1.739,-0.157) (1.774,-0.162) (1.810,-0.166) (1.846,-0.170) (1.882,-0.175) (1.917,-0.178) (1.953,-0.182) (1.989,-0.186) (2.025,-0.189) (2.061,-0.192) (2.097,-0.195) (2.132,-0.198) (2.168,-0.201) (2.204,-0.203) (2.240,-0.206) (2.276,-0.208) (2.300,-0.210)}
\LC{(0.911,0.282) (0.938,0.258) (0.967,0.236) (0.997,0.217) (1.028,0.199) (1.061,0.183) (1.094,0.169) (1.127,0.156) (1.161,0.144) (1.195,0.132) (1.230,0.122) (1.264,0.113) (1.299,0.104) (1.334,0.096) (1.370,0.088) (1.405,0.081) (1.440,0.075) (1.476,0.069) (1.511,0.063) (1.547,0.058) (1.583,0.053) (1.618,0.048) (1.654,0.044) (1.690,0.039) (1.726,0.036) (1.761,0.032) (1.797,0.028) (1.833,0.025) (1.869,0.022) (1.905,0.019) (1.941,0.016) (1.977,0.014) (2.012,0.011) (2.048,0.009) (2.084,0.006) (2.120,0.004) (2.156,0.002) (2.192,0.000) (2.228,-0.002) (2.264,-0.003) (2.300,-0.005)}
\LC{(0.973,0.353) (1.005,0.337) (1.038,0.323) (1.072,0.311) (1.107,0.300) (1.141,0.291) (1.176,0.283) (1.212,0.276) (1.247,0.269) (1.283,0.263) (1.318,0.258) (1.354,0.253) (1.390,0.249) (1.425,0.245) (1.461,0.241) (1.497,0.238) (1.533,0.235) (1.569,0.232) (1.605,0.229) (1.641,0.227) (1.677,0.224) (1.712,0.222) (1.748,0.220) (1.784,0.218) (1.820,0.216) (1.856,0.215) (1.892,0.213) (1.928,0.212) (1.964,0.210) (2.000,0.209) (2.036,0.207) (2.072,0.206) (2.108,0.205) (2.144,0.204) (2.180,0.203) (2.216,0.202) (2.252,0.201) (2.288,0.200) (2.300,0.200)}
\LC{(0.994,0.428) (1.029,0.419) (1.064,0.415) (1.100,0.412) (1.136,0.410) (1.172,0.408) (1.208,0.407) (1.244,0.406) (1.280,0.406) (1.316,0.405) (1.352,0.405) (1.388,0.404) (1.424,0.404) (1.460,0.404) (1.496,0.404) (1.532,0.404) (1.568,0.404) (1.604,0.404) (1.640,0.404) (1.676,0.404) (1.712,0.404) (1.748,0.404) (1.784,0.404) (1.820,0.404) (1.856,0.404) (1.892,0.404) (1.928,0.404) (1.964,0.404) (2.000,0.404) (2.036,0.404) (2.072,0.405) (2.108,0.405) (2.144,0.405) (2.180,0.405) (2.216,0.405) (2.252,0.405) (2.288,0.405) (2.300,0.405)}
\LC{(0.991,0.477) (1.024,0.492) (1.059,0.502) (1.094,0.511) (1.129,0.519) (1.164,0.526) (1.200,0.532) (1.235,0.538) (1.271,0.543) (1.306,0.548) (1.342,0.553) (1.378,0.557) (1.414,0.561) (1.449,0.565) (1.485,0.568) (1.521,0.571) (1.557,0.574) (1.593,0.577) (1.629,0.580) (1.665,0.582) (1.701,0.585) (1.736,0.587) (1.772,0.589) (1.808,0.591) (1.844,0.593) (1.880,0.595) (1.916,0.596) (1.952,0.598) (1.988,0.599) (2.024,0.601) (2.060,0.602) (2.096,0.604) (2.132,0.605) (2.168,0.606) (2.204,0.607) (2.240,0.608) (2.276,0.609) (2.300,0.610)}
\LC{(0.956,0.539) (0.983,0.563) (1.012,0.585) (1.042,0.603) (1.074,0.620) (1.107,0.636) (1.140,0.649) (1.174,0.662) (1.208,0.673) (1.242,0.684) (1.277,0.693) (1.312,0.702) (1.347,0.711) (1.382,0.718) (1.417,0.726) (1.453,0.732) (1.488,0.739) (1.524,0.744) (1.559,0.750) (1.595,0.755) (1.630,0.760) (1.666,0.764) (1.702,0.769) (1.738,0.773) (1.773,0.777) (1.809,0.780) (1.845,0.784) (1.881,0.787) (1.917,0.790) (1.953,0.793) (1.989,0.796) (2.025,0.798) (2.060,0.801) (2.096,0.803) (2.132,0.806) (2.168,0.808) (2.204,0.810) (2.240,0.812) (2.276,0.814) (2.300,0.815)}
\LC{(0.873,0.589) (0.893,0.619) (0.915,0.648) (0.939,0.674) (0.965,0.699) (0.993,0.722) (1.022,0.743) (1.052,0.762) (1.083,0.780) (1.115,0.797) (1.148,0.813) (1.181,0.827) (1.214,0.840) (1.248,0.853) (1.282,0.864) (1.317,0.875) (1.351,0.885) (1.386,0.895) (1.421,0.904) (1.456,0.912) (1.491,0.920) (1.526,0.927) (1.561,0.934) (1.597,0.941) (1.632,0.947) (1.668,0.953) (1.703,0.959) (1.739,0.964) (1.775,0.969) (1.810,0.974) (1.846,0.978) (1.882,0.982) (1.917,0.987) (1.953,0.990) (1.989,0.994) (2.025,0.998) (2.061,1.001) (2.097,1.004) (2.132,1.007) (2.168,1.010) (2.204,1.013) (2.240,1.016) (2.276,1.018) (2.300,1.020)}
\LC{(0.752,0.644) (0.764,0.678) (0.780,0.710) (0.799,0.741) (0.821,0.770) (0.844,0.798) (0.868,0.824) (0.895,0.848) (0.922,0.871) (0.951,0.893) (0.981,0.914) (1.011,0.933) (1.042,0.951) (1.074,0.968) (1.106,0.984) (1.139,0.999) (1.172,1.013) (1.205,1.027) (1.239,1.039) (1.273,1.051) (1.307,1.062) (1.342,1.073) (1.376,1.083) (1.411,1.092) (1.446,1.101) (1.481,1.110) (1.516,1.118) (1.551,1.126) (1.586,1.133) (1.622,1.140) (1.657,1.146) (1.692,1.153) (1.728,1.159) (1.764,1.164) (1.799,1.170) (1.835,1.175) (1.870,1.180) (1.906,1.185) (1.942,1.189) (1.978,1.193) (2.013,1.198) (2.049,1.202) (2.085,1.205) (2.121,1.209) (2.157,1.212) (2.192,1.216) (2.228,1.219) (2.264,1.222) (2.300,1.225)}
\LC{(0.559,0.710) (0.569,0.745) (0.582,0.778) (0.598,0.811) (0.616,0.842) (0.636,0.872) (0.658,0.900) (0.682,0.927) (0.707,0.954) (0.733,0.978) (0.760,1.002) (0.788,1.024) (0.817,1.046) (0.847,1.066) (0.877,1.086) (0.908,1.104) (0.939,1.122) (0.971,1.138) (1.004,1.154) (1.036,1.169) (1.069,1.184) (1.102,1.197) (1.136,1.211) (1.170,1.223) (1.204,1.235) (1.238,1.246) (1.272,1.257) (1.307,1.267) (1.341,1.277) (1.376,1.287) (1.411,1.295) (1.446,1.304) (1.481,1.312) (1.516,1.320) (1.551,1.328) (1.587,1.335) (1.622,1.342) (1.657,1.348) (1.693,1.354) (1.728,1.360) (1.764,1.366) (1.799,1.372) (1.835,1.377) (1.871,1.382) (1.906,1.387) (1.942,1.392) (1.978,1.396) (2.014,1.401) (2.049,1.405) (2.085,1.409) (2.121,1.413) (2.157,1.416) (2.192,1.420) (2.228,1.424) (2.264,1.427) (2.300,1.430)}

      % Conductor neutro en campo uniforme - flechas
% Generado por scripts/generar_datos.py (no editar a mano).
\FL{-1.463}{-1.775}{6.2}
\FL{1.475}{-1.759}{-5.7}
\FL{-1.464}{-1.560}{7.3}
\FL{1.476}{-1.543}{-6.8}
\FL{-1.466}{-1.345}{8.6}
\FL{1.472}{-1.326}{-8.0}
\FL{-1.906}{-1.195}{6.9}
\FL{-0.562}{-0.896}{20.8}
\FL{-1.930}{-0.991}{7.1}
\FL{-0.688}{-0.683}{23.7}
\FL{-1.954}{-0.789}{7.1}
\FL{-0.778}{-0.473}{26.9}
\FL{-1.978}{-0.588}{6.7}
\FL{-0.885}{-0.290}{29.8}
\FL{-2.001}{-0.389}{5.9}
\FL{-1.016}{-0.149}{28.7}
\FL{-2.037}{-0.193}{4.5}
\FL{-1.121}{-0.019}{23.9}
\FL{-2.048}{0.005}{2.8}
\FL{-1.204}{0.106}{15.7}
\FL{-2.060}{0.202}{0.7}
\FL{-1.245}{0.235}{6.2}
\FL{-2.060}{0.400}{-1.4}
\FL{-1.257}{0.365}{-4.2}
\FL{-2.060}{0.598}{-3.4}
\FL{-1.228}{0.493}{-14.5}
\FL{-2.037}{0.795}{-5.2}
\FL{-1.158}{0.616}{-24.2}
\FL{-2.013}{0.992}{-6.5}
\FL{-1.048}{0.730}{-32.8}
\FL{-1.978}{1.190}{-7.4}
\FL{-0.904}{0.846}{-38.5}
\FL{-1.930}{1.388}{-7.9}
\FL{-0.707}{0.972}{-39.7}
\FL{-1.422}{1.507}{-12.2}
\FL{1.441}{1.549}{11.9}
\FL{-1.455}{1.731}{-10.1}
\FL{1.470}{1.767}{10.1}
\FL{0.428}{-0.889}{-18.1}
\FL{1.894}{-1.195}{-6.6}
\FL{0.548}{-0.653}{-23.4}
\FL{1.918}{-0.990}{-7.0}
\FL{0.665}{-0.436}{-28.1}
\FL{1.930}{-0.786}{-7.2}
\FL{0.797}{-0.254}{-31.1}
\FL{1.966}{-0.586}{-6.9}
\FL{0.930}{-0.105}{-30.7}
\FL{2.001}{-0.387}{-6.3}
\FL{1.046}{0.026}{-26.7}
\FL{2.025}{-0.189}{-5.2}
\FL{1.138}{0.152}{-20.0}
\FL{2.048}{0.009}{-3.7}
\FL{1.212}{0.276}{-11.0}
\FL{2.060}{0.207}{-1.9}
\FL{1.232}{0.407}{-1.3}
\FL{2.072}{0.405}{0.1}
\FL{1.223}{0.536}{9.1}
\FL{2.072}{0.603}{2.1}
\FL{1.174}{0.662}{19.4}
\FL{2.060}{0.801}{3.9}
\FL{1.083}{0.780}{28.7}
\FL{2.025}{0.998}{5.5}
\FL{0.941}{0.886}{36.8}
\FL{2.001}{1.196}{6.5}
\FL{0.760}{1.002}{39.8}
\FL{1.954}{1.393}{7.3}

    \end{scope}
    % Contorno del conductor irregular
% Generado por scripts/generar_datos.py (no editar a mano).
\def\formairregular{(-0.534,0.681) -- (-0.416,0.719) -- (-0.297,0.748) -- (-0.178,0.769) -- (-0.059,0.782) -- (0.059,0.786) -- (0.178,0.782) -- (0.297,0.770) -- (0.416,0.749) -- (0.534,0.721) -- (0.640,0.691) -- (0.732,0.661) -- (0.811,0.631) -- (0.877,0.599) -- (0.929,0.568) -- (0.968,0.535) -- (0.994,0.502) -- (1.006,0.468) -- (1.006,0.431) -- (0.994,0.390) -- (0.969,0.345) -- (0.931,0.297) -- (0.881,0.245) -- (0.819,0.190) -- (0.744,0.131) -- (0.656,0.069) -- (0.575,0.004) -- (0.500,-0.063) -- (0.431,-0.133) -- (0.369,-0.205) -- (0.312,-0.279) -- (0.262,-0.355) -- (0.219,-0.434) -- (0.181,-0.516) -- (0.144,-0.586) -- (0.106,-0.647) -- (0.069,-0.697) -- (0.031,-0.736) -- (-0.006,-0.765) -- (-0.044,-0.783) -- (-0.081,-0.791) -- (-0.119,-0.789) -- (-0.156,-0.777) -- (-0.194,-0.755) -- (-0.231,-0.724) -- (-0.269,-0.683) -- (-0.306,-0.633) -- (-0.344,-0.574) -- (-0.381,-0.504) -- (-0.419,-0.426) -- (-0.460,-0.350) -- (-0.505,-0.279) -- (-0.555,-0.211) -- (-0.608,-0.147) -- (-0.665,-0.086) -- (-0.726,-0.029) -- (-0.791,0.025) -- (-0.859,0.075) -- (-0.917,0.123) -- (-0.964,0.170) -- (-1.000,0.216) -- (-1.025,0.259) -- (-1.039,0.302) -- (-1.042,0.342) -- (-1.034,0.381) -- (-1.016,0.419) -- (-0.984,0.456) -- (-0.941,0.494) -- (-0.884,0.531) -- (-0.816,0.569) -- (-0.734,0.606) -- (-0.641,0.644) -- cycle}

    \fill[top color=azulvelo!80!white, bottom color=metal, draw=azulprof, line width=0.6pt] \formairregular;
    % Conductor neutro en campo uniforme - cargas
% Generado por scripts/generar_datos.py (no editar a mano).
\CP{0.223}{-0.629}
\CP{0.442}{-0.257}
\CP{0.699}{-0.012}
\CP{0.903}{0.145}
\CP{1.034}{0.283}
\CP{1.096}{0.459}
\CP{0.999}{0.626}
\CP{0.820}{0.724}
\CP{0.492}{0.823}
\CN{-0.744}{0.699}
\CN{-0.969}{0.583}
\CN{-1.107}{0.436}
\CN{-1.119}{0.256}
\CN{-1.016}{0.095}
\CN{-0.848}{-0.044}
\CN{-0.635}{-0.255}
\CN{-0.448}{-0.570}
\CN{-0.273}{0.844}

    \node[font=\footnotesize, text=azulprof] at (-0.15,0.18) {$\vE=\vc{0}$};
    \draw[azulclaro, line width=0.5pt] (-2.3,-1.9) rectangle (2.3,1.9);
    \node[anchor=north] at (0,-2.0) {(b) Conductor irregular};
  \end{scope}
  % flecha del campo aplicado
  \draw[-{Stealth[length=5pt]}, azulprof, line width=0.9pt] (4.95,2.95) -- (6.25,2.95)
    node[midway, above, font=\small] {$\vE_0$};
\end{tikzpicture}
\caption{Conductores neutros en un campo uniforme $\vE_0$ (de izquierda a derecha). Líneas de campo (continuas) y equipotenciales (trazos).
(a) Esfera: la solución exacta, $\phi=-E_0\bigl(r-a^{3}/r^{2}\bigr)\cos\theta$, induce $\sigma=3\eo E_0\cos\theta$. (b) Conductor de forma
arbitraria, calculado numéricamente resolviendo $\Lap\phi=0$ con la condición de carga neta nula. En ambos casos las líneas llegan
perpendiculares a la superficie y el campo interior es nulo.}
\label{fig:uniforme}
\end{figure}


% -------------------------------------------------------------------------
\subsection{Notación y convenios}\label{sec:notacion}
% -------------------------------------------------------------------------
\begin{notacion}\label{not:tema3}
A lo largo del tema usamos la notación de los temas anteriores~\cite{TEMA0,TEMA1,TEMA2}, y además:
\begin{itemize}
  \item $\Vc$ es la región (cerrada) ocupada por un conductor y $\Sc=\partial\Vc$ su superficie; $\un$ es la normal unitaria \emph{exterior} al
        conductor, es decir, la que apunta del metal al vacío. Si el conductor tiene cavidades, $\Sc$ incluye sus caras interiores, y en ellas
        $\un$ apunta hacia la cavidad.
  \item $\sigma(\vc{r})$ es la densidad superficial de carga sobre $\Sc$, y $\Qint$, $\Qext$ las cargas de las caras interior y exterior de un
        conductor hueco.
  \item En un sistema de $N$ conductores, $Q_i$ y $\phi_i$ son la carga total y el potencial del conductor $i$; los agrupamos en los vectores
        columna $\vQ=(Q_1,\dots,Q_N)\Tt$ y $\vphi=(\phi_1,\dots,\phi_N)\Tt$. Las matrices de coeficientes de potencial y de capacidad e influencia
        son $\Pm=(p_{ij})$ y $\Cm=(c_{ij})$.
  \item Mientras no se diga lo contrario, el origen de potenciales está en el infinito, $\phi(\infty)=0$, y escribimos
        $\kc\equiv1/(4\pi\eo)\simeq8{,}988\cdot10^{9}\,\mathrm{N\,m^{2}\,C^{-2}}$~\cite{CODATA}.
  \item «A tierra» significa $\phi=0$ (la Tierra es un conductor tan grande que su potencial no cambia al intercambiar carga con el sistema);
        «aislado» significa que la carga del conductor se conserva.
\end{itemize}
\end{notacion}

% =========================================================================
\section{Propiedades del conductor en equilibrio}\label{sec:propiedades}
% =========================================================================
La \cref{prop:campo-nulo} es la semilla de todo el tema. Combinada con las leyes de la electrostática del vacío (la ley de Gauss en
forma diferencial e integral, el carácter conservativo del campo y las condiciones de contorno~\cite{TEMA1}), da lugar a un conjunto de
propiedades que caracterizan por completo al conductor en equilibrio.

% -------------------------------------------------------------------------
\subsection{Propiedades básicas}\label{sec:basicas}
% -------------------------------------------------------------------------
\begin{teorema}[Propiedades del conductor en equilibrio electrostático]\label{teo:propiedades}
Sea un conductor perfecto $\Vc$, de superficie $\Sc$, en equilibrio electrostático. Entonces:
\begin{enumerate}[label=\textup{(\roman*)}]
  \item\label{it:p1} el campo eléctrico es nulo en su interior, $\vE=\vc{0}$;
  \item\label{it:p2} el conductor es un \emph{volumen equipotencial}: $\phi(\vc{r})=\phi_{\Vc}=\mathrm{cte}$ para todo $\vc{r}\in\Vc$, y su
        superficie $\Sc$ es una superficie equipotencial;
  \item\label{it:p3} la densidad volúmica de carga es nula en el interior, $\rho(\vc{r})=0$;
  \item\label{it:p4} si el conductor está cargado, toda su carga se sitúa en su superficie: $Q=\oiint_{\Sc}\sigma\,\dsc$;
  \item\label{it:p5} el campo en el exterior e inmediatamente próximo a la superficie es normal a ella y vale
        \begin{equation}\label{eq:campo-superficie}
          \resalta{\vE\big|_{\Sc^{+}}=\frac{\sigma}{\eo}\,\un .}
        \end{equation}
\end{enumerate}
\end{teorema}
\begin{proof}
\ref{it:p1} Es la \cref{prop:campo-nulo}.

\ref{it:p2} Sean $A$ y $B$ dos puntos cualesquiera del conductor y $\Gamma$ una curva que los une \emph{por dentro} del material (que existe, pues
$\Vc$ es conexo; véase la \cref{fig:propiedades}). Como el campo electrostático es conservativo, $\vE=-\nab\phi$, y la diferencia de potencial
no depende del camino~\cite{TEMA0}:
\[
  \phi(B)-\phi(A)=-\int_{\Gamma}\vE\cdot\dl=-\int_{\Gamma}\vc{0}\cdot\dl=0 .
\]
Por continuidad del potencial, la igualdad se extiende a los puntos de $\Sc$.

\ref{it:p3} Por la ley de Gauss en forma diferencial, en todo punto interior
\[
  \nab\cdot\vE=\frac{\rho}{\eo}\quad\text{y}\quad\vE=\vc{0}\quad\Longrightarrow\quad\rho(\vc{r})=\eo\,\nab\cdot\vc{0}=0 .
\]
Equivalentemente, en forma integral: para cualquier superficie gaussiana $S_G$ contenida en el interior del conductor,
$\oiint_{S_G}\vE\cdot\dS=0=\Qenc/\eo$, luego $\Qenc=0$; como $S_G$ es arbitraria, no hay carga en ningún volumen interior.

\ref{it:p4} Si la carga no puede estar en el interior, el único lugar que le queda es la superficie, donde $\vE$ puede ser discontinuo y la ley de
Gauss diferencial deja de aplicarse.

\ref{it:p5} Aplicamos las condiciones de contorno del campo electrostático~\cite[§6]{TEMA1} en un punto de $\Sc$, con el medio~1 el conductor, el
medio~2 el vacío y $\un$ la normal del 1 al 2 (la exterior):
\[
  \vE_{2t}-\vE_{1t}=\vc{0},\qquad E_{2n}-E_{1n}=\frac{\sigma}{\eo} .
\]
Como dentro $\vE_1=\vc{0}$, queda $\vE_{2t}=\vc{0}$ y $E_{2n}=\sigma/\eo$: el campo exterior no tiene componente tangencial y su componente
normal es $\sigma/\eo$, que es \cref{eq:campo-superficie}.
\end{proof}

% =========================================================================
%  Figura: propiedades básicas. (a) Conductor cargado: campo normal a la
%  superficie, de módulo sigma/eps0, y camino interior entre A y B
%  (volumen equipotencial). (b) Cilindro de Gauss en la superficie.
%  Contorno en polares: r(t) = 1.45 + 0.30 cos 2t + 0.18 sen 3t; la normal
%  exterior se calcula con la derivada r'(t) en cada punto.
% =========================================================================
\begin{figure}[t]
\centering
\begin{tikzpicture}[font=\small]
  % ---------------------------------------------------------------- (a)
  \begin{scope}
    \fill[top color=azulvelo!80!white, bottom color=metal, draw=azulprof, line width=0.6pt]
      plot[domain=0:360, samples=120, smooth cycle, variable=\t]
      ({(1.45+0.30*cos(2*\t)+0.18*sin(3*\t))*cos(\t)},{(1.45+0.30*cos(2*\t)+0.18*sin(3*\t))*sin(\t)});
    \foreach \t in {0,20,...,340}{
      \pgfmathsetmacro{\r}{1.45+0.30*cos(2*\t)+0.18*sin(3*\t)}
      \pgfmathsetmacro{\rp}{-0.60*sin(2*\t)+0.54*cos(3*\t)}
      \pgfmathsetmacro{\tx}{\rp*cos(\t)-\r*sin(\t)}
      \pgfmathsetmacro{\ty}{\rp*sin(\t)+\r*cos(\t)}
      \pgfmathsetmacro{\nn}{sqrt(\tx*\tx+\ty*\ty)}
      \pgfmathsetmacro{\nx}{\ty/\nn}\pgfmathsetmacro{\ny}{-\tx/\nn}
      \draw[-{Stealth[length=4pt]}, azulacero, line width=0.6pt]
        ({\r*cos(\t)},{\r*sin(\t)}) -- ++({0.5*\nx},{0.5*\ny});
      \node[signopos, font=\tiny\bfseries] at ({\r*cos(\t)-0.13*\nx},{\r*sin(\t)-0.13*\ny}) {$+$};}
    % camino interior entre A y B
    \draw[cargapos!80, line width=0.7pt, densely dashed, -{Stealth[length=4pt]}]
      (-0.85,0.05) .. controls (-0.3,1.0) and (0.4,-0.5) .. (0.75,0.45);
    \fill[azulprof] (-0.85,0.05) circle (1.3pt) node[left, font=\footnotesize] {$A$};
    \fill[azulprof] (0.75,0.45) circle (1.3pt) node[right, font=\footnotesize] {$B$};
    \node[font=\footnotesize, text=azulprof, align=center] at (0.1,-0.75) {$\vE=\vc{0}$,\ $\phi=\mathrm{cte}$};
    \node[font=\footnotesize, text=azulacero] at (2.35,1.25) {$\vE=\dfrac{\sigma}{\eo}\,\un$};
    \node[font=\footnotesize, align=center] at (0,-2.55)
      {$\displaystyle\phi_B-\phi_A=-\int_A^B\vE\cdot\dl=0$};
    \node at (0,-3.25) {(a)};
  \end{scope}
  % ---------------------------------------------------------------- (b)
  \begin{scope}[xshift=6.6cm, yshift=-0.2cm]
    \fill[top color=azulvelo!80!white, bottom color=metal] (-2.2,-1.6) rectangle (2.2,0);
    \draw[azulprof, line width=0.7pt] (-2.2,0) -- (2.2,0);
    \foreach \x in {-1.9,-1.5,...,1.95}{\node[signopos, font=\tiny\bfseries] at (\x,-0.12) {$+$};}
    % cilindro de Gauss
    \draw[gaussazul, fill=white, fill opacity=0.35] (-0.75,-0.75) rectangle (0.75,0.75);
    \draw[gaussazul] (0,0.75) ellipse (0.75 and 0.13);
    \draw[gaussazul] (0,-0.75) ellipse (0.75 and 0.13);
    \foreach \x in {-0.45,0,0.45}{\draw[-{Stealth[length=4pt]}, azulacero, line width=0.7pt] (\x,0.05) -- (\x,1.55);}
    \draw[-{Stealth[length=4pt]}, verdeazul, line width=0.8pt] (1.15,0.75) -- (1.15,1.35) node[right, font=\footnotesize] {$\un$};
    \node[font=\footnotesize, text=azulacero] at (0.0,1.85) {$\vE=E\,\un$};
    \node[font=\footnotesize] at (1.3,0.42) {$A$};
    \node[font=\footnotesize, text=azulprof] at (-1.45,-1.05) {$\vE=\vc{0}$};
    \node[font=\footnotesize, text=grispizarra] at (1.7,-1.25) {conductor};
    \node[font=\footnotesize, text=grispizarra] at (-1.65,0.55) {vacío};
    \node[font=\footnotesize, align=center] at (0,-2.35)
      {$\displaystyle\oiint\vE\cdot\dS=EA=\frac{\sigma A}{\eo}$};
    \node at (0,-3.05) {(b)};
  \end{scope}
\end{tikzpicture}
\caption{Propiedades básicas del conductor en equilibrio. (a) Conductor cargado: el interior es un volumen equipotencial sin campo (camino
de $A$ a $B$ por dentro del metal), la carga está en la superficie y el campo exterior es normal a ella, de módulo $\sigma/\eo$. (b) Cálculo del
campo superficial con un cilindro de Gauss: solo la base exterior contribuye al flujo.}
\label{fig:propiedades}
\end{figure}


\begin{observacion}[Comprobación directa con un cilindro de Gauss]\label{obs:pastillero}
El resultado \ref{it:p5} se obtiene también con un pequeño cilindro de bases $A$, una dentro y otra fuera del conductor, perpendicular a la
superficie (\figpanel{fig:propiedades}{b}). Dentro $\vE=\vc{0}$; en la pared lateral, que puede hacerse tan corta como se quiera, el flujo es
despreciable; y en la base exterior $\vE=E\,\un$ es normal. Así $\PhiE=E\,A=\Qenc/\eo=\sigma A/\eo$, es decir, $E=\sigma/\eo$.
\end{observacion}

\begin{corolario}[Líneas de campo y conductores]\label{cor:lineas}
\begin{enumerate}[label=\textup{(\alph*)}]
  \item Las líneas de campo nacen o mueren perpendicularmente en la superficie de los conductores: salen de las zonas con $\sigma>0$ y llegan a
        las zonas con $\sigma<0$.
  \item Ninguna línea de campo puede salir de un conductor y volver a él.
  \item El potencial decrece estrictamente a lo largo de una línea de campo, en el sentido de $\vE$.
\end{enumerate}
\end{corolario}
\begin{proof}
(a) es consecuencia directa de \cref{eq:campo-superficie}. (c) Si $\Gamma$ es un arco de línea de campo recorrido en el sentido de $\vE$, el vector
$\dl$ es paralelo a $\vE$ y $\phi(B)-\phi(A)=-\int_{\Gamma}\vE\cdot\dl=-\int_{\Gamma}\norm{\vE}\,\dif\ell<0$. (b) Si una línea saliese del conductor en
$A$ y volviese a él en $B$, por (c) sería $\phi(B)<\phi(A)$, lo que contradice que el conductor es equipotencial.
\end{proof}

\begin{observacion}[Por qué la carga va a la superficie: el teorema de Thomson]\label{obs:thomson}
Hay una forma intuitiva de entender \ref{it:p4}: las cargas libres del mismo signo se repelen con fuerzas que decaen como $1/R^{2}$, y cada
una intenta alejarse lo más posible de las demás; el sitio más alejado al que pueden llegar sin abandonar el conductor es su superficie. El
\emph{teorema de Thomson} (lord Kelvin) da a esta idea una forma precisa: entre todas las distribuciones de una carga dada sobre un conductor,
la de equilibrio es la que hace mínima la energía electrostática del sistema~\cite{JACKSON}. Lo demostraremos cuando estudiemos la energía.
\end{observacion}

% -------------------------------------------------------------------------
\subsection{Efecto punta}\label{sec:punta}
% -------------------------------------------------------------------------
La densidad superficial de carga de un conductor en equilibrio no tiene por qué ser constante: aunque todo el conductor está al mismo
potencial, la carga se acumula de forma preferente en las zonas de la superficie con \emph{menor radio de curvatura}. Es el llamado
\emph{efecto punta}. Podemos demostrarlo de manera sencilla en una geometría favorable.

\begin{proposicion}[Esferas lejanas unidas por un hilo]\label{prop:punta}
Sean dos esferas conductoras de radios $R_1$ y $R_2$, unidas por un hilo conductor muy fino y separadas una distancia mucho mayor que sus
radios, con carga total $Q=Q_1+Q_2$. Entonces, en el equilibrio,
\begin{equation}\label{eq:punta}
  \frac{Q_1}{R_1}=\frac{Q_2}{R_2},\qquad\resalta{\sigma_1R_1=\sigma_2R_2,}\qquad E_1R_1=E_2R_2 ,
\end{equation}
donde $\sigma_i=Q_i/(4\pi R_i^{2})$ y $E_i=\sigma_i/\eo$ son la densidad de carga y el campo en la superficie de cada esfera. En particular, la
esfera más pequeña tiene mayor densidad de carga y mayor campo en su superficie.
\end{proposicion}
\begin{proof}
Al estar muy alejadas, la influencia de una esfera sobre la otra es despreciable: la carga de cada una se reparte uniformemente sobre su
superficie y su potencial es el de una esfera aislada, $\phi_i=\kc Q_i/R_i$ (la carga del hilo también se desprecia). Como están unidas por un
conductor, forman un único volumen equipotencial (\cref{teo:propiedades}\,\ref{it:p2}), luego
\[
  \frac{1}{4\pi\eo}\,\frac{Q_1}{R_1}=\frac{1}{4\pi\eo}\,\frac{Q_2}{R_2}
  \;\Longrightarrow\;\frac{\sigma_1\,4\pi R_1^{2}}{R_1}=\frac{\sigma_2\,4\pi R_2^{2}}{R_2}
  \;\Longrightarrow\;\sigma_1R_1=\sigma_2R_2 .
\]
Multiplicando por $1/\eo$ se obtiene la relación entre los campos. Si $R_1>R_2$, entonces $\sigma_2=\sigma_1R_1/R_2>\sigma_1$.
\end{proof}

% =========================================================================
%  Figura: efecto punta. (a) Dos esferas lejanas unidas por un hilo;
%  (b) y (c) descomposición de la repulsión de una carga vecina, situada a la
%  misma distancia de arco s = 0,9 sobre superficies de radios R = 3 y R = 0,8.
%  Con P en el origen y N en el arco, F ∝ (P - N)/|P - N|:
%    R = 3:   N = (-0.887,-0.134),  F = (1.187, 0.179)
%    R = 0,8: N = (-0.722,-0.456),  F = (1.014, 0.641)
% =========================================================================
\begin{figure}[t]
\centering
\begin{tikzpicture}[font=\small]
  % ---------------------------------------------------------------- (a)
  \begin{scope}
    \shade[ball color=azulvelo] (0,0) circle (1.0);
    \draw[azulprof, line width=0.6pt] (0,0) circle (1.0);
    \shade[ball color=azulvelo] (4.1,0.0) circle (0.35);
    \draw[azulprof, line width=0.6pt] (4.1,0) circle (0.35);
    \draw[azulprof, line width=0.8pt] (0.97,0.25) .. controls (2.0,0.9) and (3.0,0.7) .. (3.77,0.12);
    \draw[-{Stealth[length=4pt]}, azulprof, line width=0.5pt] (0,0) -- (-0.62,0.78) node[midway, right, font=\footnotesize] {$R_1$};
    \draw[-{Stealth[length=3.5pt]}, azulprof, line width=0.5pt] (4.1,0) -- (4.36,0.23);
    \node[font=\footnotesize] at (4.1,-0.6) {$R_2$};
    \node[font=\footnotesize] at (2.3,1.05) {hilo};
    \node[font=\footnotesize, text=azulprof] at (0,-1.3) {$Q_1,\ \sigma_1$};
    \node[font=\footnotesize, text=azulprof] at (4.1,-1.0) {$Q_2,\ \sigma_2$};
    \node[font=\footnotesize, align=center] at (2.05,-1.95)
      {$\displaystyle\frac{Q_1}{4\pi\eo R_1}=\frac{Q_2}{4\pi\eo R_2}\;\Longrightarrow\;\sigma_1R_1=\sigma_2R_2$};
    \node at (2.05,-2.75) {(a)};
  \end{scope}
  % ---------------------------------------------------------------- (b) R = 3
  \begin{scope}[xshift=7.4cm, yshift=0.2cm]
    \draw[azulprof, line width=0.8pt] (0,-3) ++(118:3) arc (118:62:3);
    \fill[cargapos] (0,0) circle (1.6pt) node[below=2pt, font=\footnotesize] {$P$};
    \fill[azulprof] (-0.887,-0.134) circle (1.6pt) node[below, font=\footnotesize] {$N$};
    \draw[punteado] (-0.887,-0.134) -- (0,0);
    \draw[-{Stealth[length=4pt]}, azulprof, line width=0.9pt] (0,0) -- (1.187,0.179) node[right, font=\footnotesize] {$\vF$};
    \draw[-{Stealth[length=4pt]}, verdeazul, line width=0.8pt] (0,0) -- (1.187,0) node[below right=-1pt, font=\footnotesize] {$\vF_t$};
    \draw[-{Stealth[length=4pt]}, cargapos, line width=0.8pt] (0,0) -- (0,0.179);
    \node[font=\footnotesize, text=cargapos, anchor=south] at (0,0.2) {$\vF_n$};
    \node[font=\footnotesize, text=grispizarra] at (0,-1.15) {radio $R$ grande};
    \node at (0,-2.95) {(b)};
  \end{scope}
  % ---------------------------------------------------------------- (c) R = 0,8
  \begin{scope}[xshift=11.4cm, yshift=0.2cm]
    \draw[azulprof, line width=0.8pt] (0,-0.8) ++(175:0.8) arc (175:5:0.8);
    \fill[cargapos] (0,0) circle (1.6pt) node[above left=-1pt, font=\footnotesize] {$P$};
    \fill[azulprof] (-0.722,-0.456) circle (1.6pt) node[left, font=\footnotesize] {$N$};
    \draw[punteado] (-0.722,-0.456) -- (0,0);
    \draw[-{Stealth[length=4pt]}, azulprof, line width=0.9pt] (0,0) -- (1.014,0.641) node[right, font=\footnotesize] {$\vF$};
    \draw[-{Stealth[length=4pt]}, verdeazul, line width=0.8pt] (0,0) -- (1.014,0) node[below, font=\footnotesize] {$\vF_t$};
    \draw[-{Stealth[length=4pt]}, cargapos, line width=0.8pt] (0,0) -- (0,0.641) node[above, font=\footnotesize] {$\vF_n$};
    \node[font=\footnotesize, text=grispizarra] at (0,-1.15) {radio $R$ pequeño};
    \node at (0,-2.95) {(c)};
  \end{scope}
\end{tikzpicture}
\caption{Efecto punta. (a) Dos esferas lejanas unidas por un hilo conductor están al mismo potencial, de modo que la pequeña tiene mayor
densidad de carga. (b,\,c) Repulsión $\vF$ que una carga vecina $N$, situada a la misma distancia sobre la superficie, ejerce sobre una carga $P$:
su componente tangencial $\vF_t$, que tiende a separar las cargas a lo largo de la superficie, pesa menos cuanto menor es el radio de curvatura.}
\label{fig:esferas-hilo}
\end{figure}


\begin{observacion}[Interpretación: la componente tangencial de la repulsión]\label{obs:curvatura}
La \cref{prop:punta} no es una mera curiosidad de dos esferas: una superficie arbitraria se comporta, localmente, como una esfera de radio igual
a su radio de curvatura. Hay una forma mecánica de verlo (\figpanel{fig:esferas-hilo}{b,\,c}). Sobre una carga de la superficie, la repulsión
$\vF$ del resto de cargas se descompone en una componente normal $\vF_n$, que la empuja a salir del conductor y que el material equilibra, y
una componente tangencial $\vF_t$, que la desplaza a lo largo de la superficie. En una zona muy curvada (radio de curvatura pequeño) las cargas
vecinas quedan «por debajo» de la tangente, y su repulsión es casi normal: $\vF_t$ disminuye con el radio de curvatura. Por eso las cargas pueden
apiñarse en las puntas sin que la componente tangencial las disperse.
\end{observacion}

La \cref{fig:punta} muestra el efecto en un conductor con forma de pera, resuelto numéricamente. Las líneas de campo, trazadas de forma que
entre dos consecutivas salga el mismo flujo, se agolpan en la punta; y la densidad de carga calculada sobre el contorno crece monótonamente al
acercarse a ella.

% =========================================================================
%  Figura «tipo Mathematica»: efecto punta en un conductor cargado con forma
%  de pera (sección de un conductor cilíndrico; Laplace 2D, caja a tierra).
%  (a) Líneas de campo con flujo constante entre dos consecutivas,
%      equipotenciales y contorno coloreado según sigma (escala logarítmica).
%  (b) sigma/sigma_min a lo largo del contorno superior frente a x.
%  Datos: figuras/datos/punta_*.
% =========================================================================
\begin{figure}[t]
\centering
\begin{tikzpicture}[font=\small]
  % ---------------------------------------------------------------- (a)
  \begin{scope}[x=1.12cm, y=1.12cm]
    \fill[azulfondo!60!white] (-2.6,-2.0) rectangle (2.6,2.0);
    \begin{scope}
      \clip (-2.6,-2.0) rectangle (2.6,2.0);
      \tikzset{numEQ/.style={azulcielo, line width=0.35pt, densely dashed}}
      % Conductor con punta - equipotenciales
% Generado por scripts/generar_datos.py (no editar a mano).
\EQ{(2.341,-2.040) (2.370,-2.012) (2.400,-1.982) (2.430,-1.951) (2.459,-1.920) (2.486,-1.890) (2.513,-1.860) (2.535,-1.835) (2.564,-1.800) (2.589,-1.770) (2.610,-1.743) (2.636,-1.710) (2.640,-1.705)}
\EQ{(-2.640,-1.913) (-2.610,-1.949) (-2.583,-1.980) (-2.556,-2.010) (-2.535,-2.033) (-2.529,-2.040)}
\EQ{(2.640,1.705) (2.613,1.740) (2.589,1.770) (2.565,1.799) (2.539,1.830) (2.513,1.860) (2.490,1.886) (2.460,1.919) (2.431,1.950) (2.402,1.980) (2.372,2.010) (2.341,2.040)}
\EQ{(-2.529,2.040) (-2.550,2.017) (-2.580,1.983) (-2.609,1.950) (-2.634,1.920) (-2.640,1.913)}
\EQ{(1.657,-2.040) (1.695,-2.015) (1.724,-1.995) (1.755,-1.973) (1.785,-1.952) (1.815,-1.930) (1.845,-1.907) (1.875,-1.884) (1.905,-1.860) (1.935,-1.835) (1.959,-1.815) (1.994,-1.785) (2.025,-1.758) (2.055,-1.730) (2.085,-1.702) (2.108,-1.680) (2.130,-1.658) (2.153,-1.635) (2.175,-1.612) (2.196,-1.590) (2.220,-1.564) (2.250,-1.531) (2.277,-1.500) (2.302,-1.470) (2.325,-1.442) (2.351,-1.410) (2.370,-1.385) (2.397,-1.350) (2.415,-1.325) (2.440,-1.290) (2.460,-1.260) (2.480,-1.230) (2.500,-1.200) (2.518,-1.170) (2.535,-1.143) (2.554,-1.110) (2.571,-1.080) (2.588,-1.050) (2.604,-1.020) (2.619,-0.990) (2.634,-0.960) (2.640,-0.947)}
\EQ{(-2.640,-1.202) (-2.625,-1.230) (-2.610,-1.258) (-2.595,-1.284) (-2.574,-1.320) (-2.556,-1.350) (-2.537,-1.380) (-2.520,-1.406) (-2.497,-1.440) (-2.477,-1.470) (-2.455,-1.500) (-2.433,-1.530) (-2.415,-1.554) (-2.387,-1.590) (-2.363,-1.620) (-2.340,-1.648) (-2.312,-1.680) (-2.286,-1.710) (-2.265,-1.733) (-2.244,-1.755) (-2.220,-1.780) (-2.190,-1.810) (-2.160,-1.839) (-2.138,-1.860) (-2.115,-1.881) (-2.089,-1.905) (-2.055,-1.934) (-2.025,-1.960) (-2.000,-1.980) (-1.965,-2.008) (-1.935,-2.031) (-1.923,-2.040)}
\EQ{(2.640,0.947) (2.626,0.975) (2.611,1.005) (2.596,1.035) (2.580,1.064) (2.565,1.091) (2.550,1.117) (2.528,1.155) (2.509,1.185) (2.490,1.215) (2.470,1.245) (2.450,1.275) (2.430,1.304) (2.408,1.335) (2.385,1.365) (2.363,1.395) (2.340,1.424) (2.315,1.455) (2.295,1.479) (2.265,1.514) (2.238,1.545) (2.210,1.575) (2.182,1.605) (2.160,1.628) (2.138,1.650) (2.115,1.673) (2.092,1.695) (2.070,1.716) (2.044,1.740) (2.011,1.770) (1.980,1.797) (1.950,1.823) (1.923,1.845) (1.890,1.872) (1.860,1.895) (1.830,1.918) (1.800,1.941) (1.770,1.963) (1.746,1.980) (1.710,2.005) (1.680,2.025) (1.657,2.040)}
\EQ{(-1.923,2.040) (-1.950,2.019) (-1.980,1.996) (-2.010,1.972) (-2.037,1.950) (-2.070,1.921) (-2.100,1.895) (-2.130,1.868) (-2.154,1.845) (-2.185,1.815) (-2.215,1.785) (-2.244,1.755) (-2.265,1.733) (-2.286,1.710) (-2.310,1.683) (-2.338,1.650) (-2.363,1.620) (-2.385,1.593) (-2.411,1.560) (-2.430,1.535) (-2.455,1.500) (-2.475,1.472) (-2.497,1.440) (-2.517,1.410) (-2.535,1.383) (-2.556,1.350) (-2.574,1.320) (-2.592,1.290) (-2.609,1.260) (-2.625,1.230) (-2.640,1.202)}
\EQ{(0.464,-2.040) (0.495,-2.032) (0.525,-2.025) (0.555,-2.017) (0.585,-2.008) (0.615,-2.000) (0.645,-1.991) (0.675,-1.982) (0.705,-1.972) (0.735,-1.962) (0.765,-1.952) (0.795,-1.941) (0.825,-1.930) (0.855,-1.919) (0.885,-1.908) (0.915,-1.896) (0.945,-1.884) (0.975,-1.871) (1.005,-1.858) (1.034,-1.845) (1.065,-1.831) (1.095,-1.817) (1.125,-1.803) (1.155,-1.788) (1.185,-1.772) (1.215,-1.757) (1.245,-1.741) (1.274,-1.725) (1.300,-1.710) (1.326,-1.695) (1.365,-1.672) (1.395,-1.654) (1.425,-1.635) (1.455,-1.616) (1.485,-1.596) (1.515,-1.576) (1.545,-1.555) (1.575,-1.534) (1.600,-1.515) (1.635,-1.489) (1.659,-1.470) (1.695,-1.442) (1.725,-1.417) (1.751,-1.395) (1.785,-1.366) (1.815,-1.339) (1.845,-1.311) (1.867,-1.290) (1.890,-1.267) (1.912,-1.245) (1.935,-1.222) (1.956,-1.200) (1.980,-1.174) (2.010,-1.140) (2.036,-1.110) (2.061,-1.080) (2.085,-1.050) (2.108,-1.020) (2.130,-0.991) (2.152,-0.960) (2.173,-0.930) (2.190,-0.905) (2.213,-0.870) (2.231,-0.840) (2.249,-0.810) (2.265,-0.782) (2.280,-0.755) (2.295,-0.726) (2.310,-0.697) (2.325,-0.666) (2.340,-0.633) (2.354,-0.600) (2.367,-0.570) (2.378,-0.540) (2.389,-0.510) (2.400,-0.480) (2.409,-0.450) (2.419,-0.420) (2.427,-0.390) (2.435,-0.360) (2.442,-0.330) (2.448,-0.300) (2.454,-0.270) (2.460,-0.240) (2.464,-0.210) (2.468,-0.180) (2.472,-0.150) (2.475,-0.120) (2.477,-0.090) (2.478,-0.060) (2.479,-0.030) (2.479,-0.000) (2.479,0.030) (2.478,0.060) (2.477,0.090) (2.475,0.120) (2.472,0.150) (2.468,0.180) (2.464,0.210) (2.460,0.240) (2.454,0.270) (2.448,0.300) (2.442,0.330) (2.435,0.360) (2.427,0.390) (2.419,0.420) (2.409,0.450) (2.400,0.479) (2.389,0.510) (2.378,0.540) (2.367,0.570) (2.355,0.598) (2.341,0.630) (2.328,0.660) (2.313,0.690) (2.298,0.720) (2.283,0.750) (2.266,0.780) (2.250,0.808) (2.231,0.840) (2.213,0.870) (2.193,0.900) (2.175,0.927) (2.152,0.960) (2.131,0.990) (2.108,1.020) (2.085,1.050) (2.061,1.080) (2.040,1.105) (2.010,1.140) (1.983,1.170) (1.956,1.200) (1.935,1.222) (1.912,1.245) (1.890,1.267) (1.867,1.290) (1.845,1.311) (1.819,1.335) (1.786,1.365) (1.755,1.392) (1.725,1.417) (1.697,1.440) (1.665,1.466) (1.640,1.485) (1.605,1.511) (1.580,1.530) (1.545,1.555) (1.516,1.575) (1.485,1.596) (1.455,1.616) (1.425,1.635) (1.395,1.654) (1.365,1.672) (1.335,1.690) (1.305,1.707) (1.275,1.724) (1.246,1.740) (1.218,1.755) (1.190,1.770) (1.160,1.785) (1.130,1.800) (1.099,1.815) (1.067,1.830) (1.035,1.845) (1.005,1.858) (0.975,1.871) (0.945,1.884) (0.915,1.896) (0.885,1.908) (0.855,1.919) (0.825,1.930) (0.795,1.941) (0.765,1.952) (0.735,1.962) (0.705,1.972) (0.675,1.982) (0.645,1.991) (0.615,2.000) (0.585,2.008) (0.555,2.017) (0.525,2.025) (0.495,2.032) (0.465,2.040) (0.464,2.040)}
\EQ{(-0.923,2.040) (-0.960,2.030) (-0.990,2.021) (-1.020,2.012) (-1.050,2.003) (-1.080,1.993) (-1.110,1.982) (-1.140,1.971) (-1.170,1.960) (-1.200,1.948) (-1.230,1.936) (-1.260,1.923) (-1.290,1.910) (-1.320,1.896) (-1.350,1.881) (-1.380,1.866) (-1.410,1.851) (-1.440,1.835) (-1.470,1.818) (-1.500,1.801) (-1.526,1.785) (-1.560,1.764) (-1.590,1.745) (-1.619,1.725) (-1.650,1.704) (-1.680,1.682) (-1.710,1.660) (-1.740,1.637) (-1.770,1.613) (-1.797,1.590) (-1.830,1.562) (-1.860,1.535) (-1.890,1.507) (-1.913,1.485) (-1.935,1.463) (-1.957,1.440) (-1.980,1.416) (-2.010,1.383) (-2.039,1.350) (-2.064,1.320) (-2.085,1.294) (-2.111,1.260) (-2.130,1.235) (-2.155,1.200) (-2.175,1.171) (-2.196,1.140) (-2.215,1.110) (-2.234,1.080) (-2.250,1.052) (-2.265,1.026) (-2.280,0.999) (-2.295,0.971) (-2.310,0.942) (-2.324,0.915) (-2.338,0.885) (-2.352,0.855) (-2.365,0.825) (-2.378,0.795) (-2.390,0.765) (-2.401,0.735) (-2.412,0.705) (-2.423,0.675) (-2.433,0.645) (-2.442,0.615) (-2.451,0.585) (-2.459,0.555) (-2.468,0.525) (-2.475,0.495) (-2.482,0.465) (-2.489,0.435) (-2.495,0.405) (-2.500,0.375) (-2.506,0.345) (-2.511,0.315) (-2.515,0.285) (-2.519,0.255) (-2.522,0.225) (-2.525,0.195) (-2.528,0.165) (-2.530,0.135) (-2.532,0.105) (-2.533,0.075) (-2.534,0.045) (-2.534,0.015) (-2.534,-0.030) (-2.533,-0.060) (-2.532,-0.090) (-2.531,-0.120) (-2.529,-0.150) (-2.527,-0.180) (-2.524,-0.210) (-2.520,-0.240) (-2.517,-0.270) (-2.513,-0.300) (-2.508,-0.330) (-2.503,-0.360) (-2.498,-0.390) (-2.492,-0.420) (-2.485,-0.450) (-2.479,-0.480) (-2.471,-0.510) (-2.464,-0.540) (-2.455,-0.570) (-2.447,-0.600) (-2.437,-0.630) (-2.428,-0.660) (-2.417,-0.690) (-2.407,-0.720) (-2.395,-0.750) (-2.384,-0.780) (-2.371,-0.810) (-2.358,-0.840) (-2.345,-0.870) (-2.331,-0.900) (-2.316,-0.930) (-2.301,-0.960) (-2.285,-0.990) (-2.269,-1.020) (-2.251,-1.050) (-2.235,-1.078) (-2.215,-1.110) (-2.196,-1.140) (-2.176,-1.170) (-2.155,-1.200) (-2.133,-1.230) (-2.115,-1.255) (-2.088,-1.290) (-2.064,-1.320) (-2.040,-1.348) (-2.012,-1.380) (-1.985,-1.410) (-1.957,-1.440) (-1.935,-1.463) (-1.913,-1.485) (-1.890,-1.507) (-1.865,-1.530) (-1.832,-1.560) (-1.800,-1.588) (-1.770,-1.613) (-1.742,-1.635) (-1.710,-1.660) (-1.683,-1.680) (-1.650,-1.704) (-1.620,-1.725) (-1.590,-1.745) (-1.560,-1.764) (-1.530,-1.783) (-1.501,-1.800) (-1.470,-1.818) (-1.440,-1.835) (-1.410,-1.851) (-1.380,-1.866) (-1.350,-1.881) (-1.320,-1.896) (-1.290,-1.910) (-1.260,-1.923) (-1.232,-1.935) (-1.200,-1.948) (-1.170,-1.960) (-1.140,-1.971) (-1.110,-1.982) (-1.080,-1.993) (-1.050,-2.003) (-1.020,-2.012) (-0.990,-2.021) (-0.960,-2.030) (-0.930,-2.038) (-0.923,-2.040)}
\EQ{(-0.510,-1.726) (-0.480,-1.728) (-0.450,-1.730) (-0.420,-1.731) (-0.390,-1.732) (-0.360,-1.733) (-0.330,-1.733) (-0.300,-1.733) (-0.270,-1.733) (-0.240,-1.732) (-0.210,-1.731) (-0.180,-1.729) (-0.150,-1.727) (-0.120,-1.725) (-0.090,-1.723) (-0.060,-1.720) (-0.030,-1.716) (-0.000,-1.713) (0.030,-1.709) (0.060,-1.705) (0.090,-1.700) (0.120,-1.695) (0.150,-1.690) (0.180,-1.684) (0.210,-1.678) (0.240,-1.672) (0.270,-1.666) (0.300,-1.659) (0.330,-1.652) (0.360,-1.644) (0.390,-1.636) (0.420,-1.628) (0.449,-1.620) (0.480,-1.611) (0.510,-1.602) (0.540,-1.593) (0.570,-1.583) (0.600,-1.573) (0.630,-1.562) (0.660,-1.552) (0.690,-1.541) (0.718,-1.530) (0.750,-1.518) (0.780,-1.506) (0.810,-1.493) (0.840,-1.481) (0.870,-1.468) (0.898,-1.455) (0.930,-1.441) (0.960,-1.426) (0.990,-1.412) (1.020,-1.397) (1.050,-1.382) (1.080,-1.366) (1.110,-1.350) (1.138,-1.335) (1.164,-1.320) (1.191,-1.305) (1.230,-1.282) (1.260,-1.264) (1.290,-1.245) (1.320,-1.226) (1.350,-1.206) (1.380,-1.186) (1.410,-1.165) (1.440,-1.144) (1.466,-1.125) (1.500,-1.099) (1.525,-1.080) (1.560,-1.052) (1.590,-1.028) (1.617,-1.005) (1.650,-0.976) (1.680,-0.949) (1.710,-0.921) (1.740,-0.891) (1.770,-0.861) (1.800,-0.829) (1.830,-0.796) (1.857,-0.765) (1.882,-0.735) (1.905,-0.706) (1.929,-0.675) (1.950,-0.646) (1.971,-0.615) (1.991,-0.585) (2.010,-0.555) (2.028,-0.525) (2.044,-0.495) (2.060,-0.465) (2.075,-0.435) (2.089,-0.405) (2.101,-0.375) (2.113,-0.345) (2.124,-0.315) (2.134,-0.285) (2.143,-0.255) (2.151,-0.225) (2.158,-0.195) (2.163,-0.165) (2.168,-0.135) (2.172,-0.105) (2.175,-0.075) (2.177,-0.045) (2.178,-0.015) (2.178,0.030) (2.176,0.060) (2.174,0.090) (2.170,0.120) (2.166,0.150) (2.161,0.180) (2.154,0.210) (2.147,0.240) (2.138,0.270) (2.129,0.300) (2.119,0.330) (2.107,0.360) (2.095,0.390) (2.082,0.420) (2.068,0.450) (2.052,0.480) (2.036,0.510) (2.019,0.540) (2.001,0.570) (1.981,0.600) (1.961,0.630) (1.940,0.660) (1.920,0.687) (1.894,0.720) (1.875,0.744) (1.845,0.779) (1.818,0.810) (1.790,0.840) (1.761,0.870) (1.731,0.900) (1.700,0.930) (1.668,0.960) (1.635,0.989) (1.605,1.015) (1.581,1.035) (1.545,1.064) (1.515,1.088) (1.486,1.110) (1.455,1.133) (1.425,1.155) (1.395,1.176) (1.365,1.196) (1.337,1.215) (1.305,1.235) (1.275,1.254) (1.245,1.273) (1.216,1.290) (1.185,1.308) (1.155,1.325) (1.125,1.342) (1.095,1.358) (1.065,1.374) (1.035,1.390) (1.005,1.405) (0.975,1.419) (0.945,1.434) (0.915,1.447) (0.885,1.461) (0.855,1.474) (0.825,1.487) (0.795,1.500) (0.765,1.512) (0.735,1.524) (0.705,1.535) (0.675,1.546) (0.645,1.557) (0.615,1.568) (0.585,1.578) (0.555,1.588) (0.525,1.597) (0.495,1.607) (0.465,1.615) (0.435,1.624) (0.405,1.632) (0.375,1.640) (0.345,1.648) (0.315,1.655) (0.285,1.662) (0.255,1.669) (0.225,1.675) (0.195,1.681) (0.165,1.687) (0.135,1.693) (0.105,1.698) (0.075,1.702) (0.045,1.707) (0.015,1.711) (-0.015,1.715) (-0.045,1.718) (-0.075,1.721) (-0.105,1.724) (-0.135,1.726) (-0.165,1.728) (-0.195,1.730) (-0.225,1.731) (-0.255,1.732) (-0.285,1.733) (-0.315,1.733) (-0.345,1.733) (-0.375,1.733) (-0.405,1.732) (-0.435,1.731) (-0.465,1.729) (-0.495,1.727) (-0.525,1.724) (-0.555,1.722) (-0.585,1.718) (-0.615,1.715) (-0.645,1.710) (-0.675,1.706) (-0.705,1.701) (-0.735,1.695) (-0.765,1.690) (-0.795,1.683) (-0.825,1.676) (-0.855,1.669) (-0.885,1.661) (-0.915,1.653) (-0.945,1.644) (-0.975,1.635) (-1.005,1.625) (-1.035,1.615) (-1.065,1.604) (-1.095,1.592) (-1.125,1.580) (-1.155,1.567) (-1.185,1.554) (-1.215,1.540) (-1.245,1.525) (-1.275,1.510) (-1.305,1.494) (-1.335,1.477) (-1.365,1.459) (-1.395,1.441) (-1.425,1.421) (-1.455,1.401) (-1.485,1.380) (-1.515,1.358) (-1.545,1.335) (-1.575,1.311) (-1.601,1.290) (-1.635,1.260) (-1.665,1.232) (-1.695,1.203) (-1.725,1.173) (-1.755,1.141) (-1.783,1.110) (-1.808,1.080) (-1.830,1.053) (-1.856,1.020) (-1.875,0.994) (-1.899,0.960) (-1.919,0.930) (-1.939,0.900) (-1.957,0.870) (-1.975,0.840) (-1.992,0.810) (-2.008,0.780) (-2.023,0.750) (-2.037,0.720) (-2.051,0.690) (-2.064,0.660) (-2.076,0.630) (-2.088,0.600) (-2.099,0.570) (-2.110,0.540) (-2.119,0.510) (-2.129,0.480) (-2.137,0.450) (-2.145,0.420) (-2.152,0.390) (-2.159,0.360) (-2.166,0.330) (-2.171,0.300) (-2.176,0.270) (-2.181,0.240) (-2.185,0.210) (-2.188,0.180) (-2.191,0.150) (-2.194,0.120) (-2.196,0.090) (-2.197,0.060) (-2.198,0.030) (-2.198,-0.000) (-2.198,-0.030) (-2.197,-0.060) (-2.196,-0.090) (-2.194,-0.120) (-2.191,-0.150) (-2.188,-0.180) (-2.185,-0.210) (-2.181,-0.240) (-2.176,-0.270) (-2.171,-0.300) (-2.166,-0.330) (-2.159,-0.360) (-2.152,-0.390) (-2.145,-0.420) (-2.137,-0.450) (-2.129,-0.480) (-2.119,-0.510) (-2.110,-0.540) (-2.099,-0.570) (-2.088,-0.600) (-2.076,-0.630) (-2.064,-0.660) (-2.051,-0.690) (-2.037,-0.720) (-2.023,-0.750) (-2.008,-0.780) (-1.992,-0.810) (-1.975,-0.840) (-1.957,-0.870) (-1.939,-0.900) (-1.920,-0.929) (-1.899,-0.960) (-1.878,-0.990) (-1.860,-1.014) (-1.832,-1.050) (-1.808,-1.080) (-1.785,-1.107) (-1.756,-1.140) (-1.728,-1.170) (-1.698,-1.200) (-1.667,-1.230) (-1.635,-1.260) (-1.605,-1.286) (-1.575,-1.311) (-1.546,-1.335) (-1.515,-1.358) (-1.485,-1.380) (-1.455,-1.401) (-1.425,-1.421) (-1.396,-1.440) (-1.365,-1.459) (-1.335,-1.477) (-1.305,-1.494) (-1.275,-1.510) (-1.245,-1.525) (-1.215,-1.540) (-1.185,-1.554) (-1.155,-1.567) (-1.125,-1.580) (-1.095,-1.592) (-1.065,-1.604) (-1.035,-1.615) (-1.005,-1.625) (-0.975,-1.635) (-0.945,-1.644) (-0.915,-1.653) (-0.885,-1.661) (-0.855,-1.669) (-0.825,-1.676) (-0.795,-1.683) (-0.765,-1.690) (-0.735,-1.695) (-0.705,-1.701) (-0.675,-1.706) (-0.645,-1.710) (-0.615,-1.715) (-0.585,-1.718) (-0.555,-1.722) (-0.525,-1.724) (-0.510,-1.726)}
\EQ{(-0.525,-1.395) (-0.495,-1.397) (-0.465,-1.398) (-0.435,-1.398) (-0.405,-1.398) (-0.375,-1.398) (-0.345,-1.397) (-0.315,-1.396) (-0.285,-1.394) (-0.255,-1.392) (-0.225,-1.389) (-0.195,-1.386) (-0.165,-1.383) (-0.135,-1.379) (-0.105,-1.375) (-0.075,-1.370) (-0.045,-1.366) (-0.015,-1.360) (0.015,-1.355) (0.045,-1.349) (0.075,-1.342) (0.105,-1.336) (0.135,-1.329) (0.165,-1.321) (0.195,-1.314) (0.225,-1.306) (0.255,-1.298) (0.285,-1.289) (0.315,-1.280) (0.345,-1.271) (0.375,-1.262) (0.405,-1.252) (0.435,-1.242) (0.465,-1.232) (0.495,-1.221) (0.525,-1.210) (0.555,-1.199) (0.585,-1.188) (0.615,-1.176) (0.645,-1.164) (0.675,-1.152) (0.704,-1.140) (0.735,-1.127) (0.765,-1.113) (0.795,-1.100) (0.825,-1.086) (0.855,-1.072) (0.885,-1.058) (0.915,-1.043) (0.945,-1.028) (0.975,-1.013) (1.005,-0.998) (1.035,-0.982) (1.065,-0.965) (1.095,-0.949) (1.125,-0.932) (1.154,-0.915) (1.185,-0.897) (1.215,-0.878) (1.245,-0.860) (1.275,-0.841) (1.305,-0.821) (1.335,-0.801) (1.365,-0.781) (1.395,-0.760) (1.425,-0.738) (1.449,-0.720) (1.485,-0.693) (1.515,-0.669) (1.544,-0.645) (1.575,-0.619) (1.605,-0.593) (1.630,-0.570) (1.662,-0.540) (1.692,-0.510) (1.720,-0.480) (1.747,-0.450) (1.770,-0.424) (1.797,-0.390) (1.819,-0.360) (1.840,-0.330) (1.859,-0.300) (1.875,-0.273) (1.890,-0.245) (1.905,-0.213) (1.919,-0.180) (1.929,-0.150) (1.938,-0.120) (1.945,-0.090) (1.949,-0.060) (1.952,-0.030) (1.953,-0.000) (1.952,0.030) (1.949,0.060) (1.945,0.090) (1.938,0.120) (1.929,0.150) (1.919,0.180) (1.907,0.210) (1.892,0.240) (1.877,0.270) (1.860,0.299) (1.840,0.330) (1.819,0.360) (1.800,0.386) (1.773,0.420) (1.747,0.450) (1.725,0.475) (1.695,0.507) (1.665,0.537) (1.635,0.565) (1.605,0.593) (1.580,0.615) (1.545,0.645) (1.515,0.669) (1.489,0.690) (1.455,0.716) (1.429,0.735) (1.395,0.760) (1.366,0.780) (1.335,0.801) (1.305,0.821) (1.276,0.840) (1.245,0.860) (1.215,0.878) (1.185,0.897) (1.155,0.914) (1.128,0.930) (1.102,0.945) (1.075,0.960) (1.047,0.975) (1.020,0.990) (0.991,1.005) (0.962,1.020) (0.932,1.035) (0.902,1.050) (0.870,1.065) (0.840,1.079) (0.810,1.093) (0.780,1.107) (0.750,1.120) (0.720,1.133) (0.690,1.146) (0.660,1.158) (0.631,1.170) (0.600,1.182) (0.570,1.194) (0.540,1.205) (0.510,1.216) (0.480,1.227) (0.450,1.237) (0.420,1.247) (0.390,1.257) (0.360,1.267) (0.330,1.276) (0.300,1.285) (0.270,1.293) (0.240,1.302) (0.210,1.310) (0.180,1.318) (0.150,1.325) (0.120,1.332) (0.090,1.339) (0.060,1.346) (0.030,1.352) (-0.000,1.358) (-0.030,1.363) (-0.060,1.368) (-0.090,1.373) (-0.120,1.377) (-0.150,1.381) (-0.180,1.385) (-0.210,1.388) (-0.240,1.391) (-0.270,1.393) (-0.300,1.395) (-0.330,1.397) (-0.360,1.398) (-0.390,1.398) (-0.420,1.398) (-0.450,1.398) (-0.480,1.397) (-0.510,1.396) (-0.540,1.394) (-0.570,1.392) (-0.600,1.389) (-0.630,1.386) (-0.660,1.382) (-0.690,1.378) (-0.720,1.373) (-0.750,1.368) (-0.780,1.362) (-0.810,1.355) (-0.840,1.348) (-0.870,1.340) (-0.900,1.331) (-0.930,1.322) (-0.960,1.312) (-0.990,1.302) (-1.020,1.291) (-1.050,1.279) (-1.080,1.266) (-1.110,1.252) (-1.140,1.238) (-1.170,1.223) (-1.200,1.207) (-1.230,1.190) (-1.260,1.172) (-1.286,1.155) (-1.320,1.132) (-1.350,1.111) (-1.380,1.089) (-1.410,1.065) (-1.440,1.040) (-1.470,1.013) (-1.495,0.990) (-1.525,0.960) (-1.554,0.930) (-1.575,0.906) (-1.605,0.871) (-1.630,0.840) (-1.650,0.813) (-1.674,0.780) (-1.694,0.750) (-1.713,0.720) (-1.731,0.690) (-1.748,0.660) (-1.764,0.630) (-1.779,0.600) (-1.794,0.570) (-1.807,0.540) (-1.819,0.510) (-1.831,0.480) (-1.842,0.450) (-1.852,0.420) (-1.861,0.390) (-1.869,0.360) (-1.877,0.330) (-1.884,0.300) (-1.891,0.270) (-1.896,0.240) (-1.901,0.210) (-1.905,0.180) (-1.909,0.150) (-1.912,0.120) (-1.914,0.090) (-1.916,0.060) (-1.917,0.030) (-1.917,-0.000) (-1.917,-0.030) (-1.916,-0.060) (-1.914,-0.090) (-1.912,-0.120) (-1.909,-0.150) (-1.905,-0.180) (-1.901,-0.210) (-1.896,-0.240) (-1.891,-0.270) (-1.884,-0.300) (-1.877,-0.330) (-1.869,-0.360) (-1.861,-0.390) (-1.852,-0.420) (-1.842,-0.450) (-1.831,-0.480) (-1.819,-0.510) (-1.807,-0.540) (-1.794,-0.570) (-1.779,-0.600) (-1.764,-0.630) (-1.748,-0.660) (-1.731,-0.690) (-1.713,-0.720) (-1.695,-0.749) (-1.674,-0.780) (-1.653,-0.810) (-1.630,-0.840) (-1.606,-0.870) (-1.580,-0.900) (-1.560,-0.923) (-1.539,-0.945) (-1.515,-0.970) (-1.485,-0.999) (-1.462,-1.020) (-1.428,-1.050) (-1.395,-1.077) (-1.365,-1.100) (-1.335,-1.122) (-1.309,-1.140) (-1.275,-1.162) (-1.245,-1.181) (-1.215,-1.198) (-1.185,-1.215) (-1.156,-1.230) (-1.125,-1.245) (-1.095,-1.259) (-1.065,-1.272) (-1.035,-1.285) (-1.005,-1.296) (-0.975,-1.307) (-0.945,-1.317) (-0.915,-1.327) (-0.885,-1.336) (-0.855,-1.344) (-0.825,-1.351) (-0.795,-1.358) (-0.765,-1.365) (-0.735,-1.370) (-0.705,-1.376) (-0.675,-1.380) (-0.645,-1.384) (-0.615,-1.388) (-0.585,-1.391) (-0.555,-1.393) (-0.525,-1.395)}
\EQ{(-0.660,-1.111) (-0.630,-1.113) (-0.600,-1.115) (-0.570,-1.117) (-0.540,-1.118) (-0.510,-1.118) (-0.480,-1.117) (-0.450,-1.116) (-0.420,-1.115) (-0.390,-1.112) (-0.360,-1.109) (-0.330,-1.106) (-0.300,-1.102) (-0.270,-1.097) (-0.240,-1.092) (-0.210,-1.086) (-0.180,-1.080) (-0.150,-1.073) (-0.120,-1.066) (-0.090,-1.059) (-0.060,-1.051) (-0.030,-1.043) (-0.000,-1.035) (0.030,-1.026) (0.060,-1.017) (0.090,-1.008) (0.120,-0.998) (0.150,-0.988) (0.180,-0.978) (0.210,-0.968) (0.240,-0.958) (0.270,-0.947) (0.300,-0.936) (0.330,-0.925) (0.360,-0.914) (0.390,-0.903) (0.420,-0.891) (0.450,-0.880) (0.480,-0.868) (0.510,-0.856) (0.540,-0.844) (0.570,-0.831) (0.600,-0.819) (0.630,-0.806) (0.660,-0.793) (0.690,-0.780) (0.720,-0.767) (0.750,-0.753) (0.780,-0.740) (0.810,-0.726) (0.840,-0.712) (0.870,-0.698) (0.900,-0.684) (0.930,-0.669) (0.960,-0.654) (0.990,-0.640) (1.020,-0.625) (1.050,-0.609) (1.080,-0.594) (1.110,-0.578) (1.140,-0.562) (1.170,-0.546) (1.200,-0.529) (1.230,-0.513) (1.260,-0.496) (1.287,-0.480) (1.320,-0.460) (1.350,-0.442) (1.380,-0.424) (1.410,-0.405) (1.440,-0.386) (1.470,-0.366) (1.500,-0.346) (1.530,-0.325) (1.560,-0.303) (1.590,-0.280) (1.620,-0.256) (1.650,-0.231) (1.674,-0.210) (1.705,-0.180) (1.733,-0.150) (1.755,-0.123) (1.777,-0.090) (1.792,-0.060) (1.801,-0.030) (1.805,-0.000) (1.801,0.030) (1.792,0.060) (1.777,0.090) (1.757,0.120) (1.733,0.150) (1.710,0.175) (1.680,0.204) (1.657,0.225) (1.621,0.255) (1.590,0.280) (1.564,0.300) (1.530,0.325) (1.501,0.345) (1.470,0.366) (1.440,0.386) (1.410,0.405) (1.380,0.424) (1.350,0.442) (1.320,0.460) (1.290,0.478) (1.261,0.495) (1.235,0.510) (1.208,0.525) (1.181,0.540) (1.153,0.555) (1.125,0.570) (1.097,0.585) (1.068,0.600) (1.039,0.615) (1.009,0.630) (0.979,0.645) (0.949,0.660) (0.918,0.675) (0.887,0.690) (0.855,0.705) (0.825,0.719) (0.795,0.733) (0.765,0.746) (0.735,0.760) (0.705,0.773) (0.675,0.786) (0.645,0.799) (0.615,0.812) (0.585,0.825) (0.555,0.837) (0.525,0.850) (0.495,0.862) (0.465,0.874) (0.436,0.885) (0.405,0.897) (0.375,0.909) (0.345,0.920) (0.315,0.931) (0.285,0.942) (0.255,0.953) (0.225,0.963) (0.195,0.973) (0.165,0.983) (0.135,0.993) (0.105,1.003) (0.075,1.012) (0.045,1.021) (0.015,1.030) (-0.015,1.039) (-0.045,1.047) (-0.075,1.055) (-0.105,1.063) (-0.135,1.070) (-0.165,1.077) (-0.195,1.083) (-0.225,1.089) (-0.255,1.094) (-0.285,1.099) (-0.315,1.104) (-0.345,1.108) (-0.375,1.111) (-0.405,1.113) (-0.435,1.115) (-0.465,1.117) (-0.495,1.118) (-0.525,1.118) (-0.555,1.117) (-0.585,1.116) (-0.615,1.114) (-0.645,1.112) (-0.675,1.109) (-0.705,1.105) (-0.735,1.101) (-0.765,1.095) (-0.795,1.089) (-0.825,1.083) (-0.855,1.075) (-0.885,1.067) (-0.915,1.057) (-0.945,1.047) (-0.975,1.036) (-1.005,1.023) (-1.035,1.010) (-1.065,0.996) (-1.095,0.980) (-1.125,0.964) (-1.155,0.946) (-1.185,0.927) (-1.215,0.906) (-1.244,0.885) (-1.275,0.861) (-1.300,0.840) (-1.333,0.810) (-1.364,0.780) (-1.393,0.750) (-1.419,0.720) (-1.440,0.695) (-1.467,0.660) (-1.485,0.636) (-1.509,0.600) (-1.528,0.570) (-1.545,0.540) (-1.560,0.513) (-1.575,0.483) (-1.590,0.450) (-1.603,0.420) (-1.615,0.390) (-1.625,0.360) (-1.635,0.330) (-1.644,0.300) (-1.652,0.270) (-1.659,0.240) (-1.665,0.210) (-1.670,0.180) (-1.674,0.150) (-1.677,0.120) (-1.680,0.090) (-1.682,0.060) (-1.683,0.030) (-1.683,-0.000) (-1.683,-0.030) (-1.682,-0.060) (-1.680,-0.090) (-1.677,-0.120) (-1.674,-0.150) (-1.670,-0.180) (-1.665,-0.210) (-1.659,-0.240) (-1.652,-0.270) (-1.644,-0.300) (-1.635,-0.330) (-1.625,-0.360) (-1.615,-0.390) (-1.603,-0.420) (-1.590,-0.450) (-1.576,-0.480) (-1.561,-0.510) (-1.545,-0.540) (-1.530,-0.566) (-1.509,-0.600) (-1.489,-0.630) (-1.470,-0.656) (-1.444,-0.690) (-1.425,-0.713) (-1.395,-0.747) (-1.365,-0.779) (-1.335,-0.808) (-1.305,-0.835) (-1.282,-0.855) (-1.245,-0.884) (-1.215,-0.906) (-1.185,-0.927) (-1.157,-0.945) (-1.125,-0.964) (-1.095,-0.980) (-1.065,-0.996) (-1.035,-1.010) (-1.005,-1.023) (-0.976,-1.035) (-0.945,-1.047) (-0.915,-1.057) (-0.885,-1.067) (-0.855,-1.075) (-0.825,-1.083) (-0.795,-1.089) (-0.765,-1.095) (-0.735,-1.101) (-0.705,-1.105) (-0.675,-1.109) (-0.660,-1.111)}

      % Conductor con punta - líneas de campo
% Generado por scripts/generar_datos.py (no editar a mano).
\LC{(-0.388,0.905) (-0.381,0.941) (-0.375,0.976) (-0.370,1.012) (-0.365,1.048) (-0.360,1.083) (-0.356,1.119) (-0.353,1.155) (-0.349,1.191) (-0.347,1.227) (-0.344,1.263) (-0.342,1.298) (-0.340,1.334) (-0.338,1.370) (-0.337,1.406) (-0.335,1.442) (-0.334,1.478) (-0.333,1.514) (-0.333,1.550) (-0.332,1.586) (-0.332,1.622) (-0.332,1.658) (-0.332,1.694) (-0.332,1.730) (-0.332,1.766) (-0.332,1.802) (-0.333,1.838) (-0.333,1.874) (-0.334,1.910) (-0.335,1.946) (-0.336,1.982) (-0.336,2.000)}
\LC{(-0.651,0.929) (-0.653,0.965) (-0.655,1.000) (-0.658,1.036) (-0.661,1.072) (-0.665,1.108) (-0.669,1.144) (-0.673,1.180) (-0.678,1.215) (-0.683,1.251) (-0.688,1.287) (-0.693,1.322) (-0.698,1.358) (-0.704,1.393) (-0.710,1.429) (-0.716,1.464) (-0.722,1.500) (-0.729,1.535) (-0.735,1.571) (-0.742,1.606) (-0.749,1.641) (-0.756,1.677) (-0.763,1.712) (-0.770,1.747) (-0.777,1.782) (-0.785,1.818) (-0.792,1.853) (-0.800,1.888) (-0.808,1.923) (-0.815,1.958) (-0.823,1.993) (-0.826,2.000)}
\LC{(-0.888,0.884) (-0.899,0.918) (-0.911,0.953) (-0.923,0.987) (-0.934,1.021) (-0.947,1.054) (-0.959,1.088) (-0.971,1.122) (-0.984,1.156) (-0.997,1.189) (-1.010,1.223) (-1.023,1.257) (-1.036,1.290) (-1.049,1.323) (-1.063,1.357) (-1.077,1.390) (-1.091,1.423) (-1.104,1.457) (-1.118,1.490) (-1.133,1.523) (-1.147,1.556) (-1.161,1.589) (-1.176,1.622) (-1.190,1.655) (-1.205,1.688) (-1.219,1.721) (-1.234,1.753) (-1.249,1.786) (-1.264,1.819) (-1.279,1.852) (-1.294,1.884) (-1.309,1.917) (-1.324,1.950) (-1.339,1.983) (-1.349,2.000)}
\LC{(-1.097,0.786) (-1.117,0.816) (-1.136,0.846) (-1.156,0.876) (-1.176,0.907) (-1.196,0.937) (-1.216,0.966) (-1.236,0.996) (-1.256,1.026) (-1.276,1.056) (-1.297,1.085) (-1.317,1.115) (-1.338,1.145) (-1.359,1.174) (-1.379,1.203) (-1.400,1.233) (-1.421,1.262) (-1.442,1.291) (-1.463,1.320) (-1.485,1.350) (-1.506,1.379) (-1.527,1.408) (-1.548,1.437) (-1.570,1.466) (-1.591,1.494) (-1.613,1.523) (-1.635,1.552) (-1.656,1.581) (-1.678,1.609) (-1.700,1.638) (-1.722,1.667) (-1.743,1.695) (-1.765,1.724) (-1.787,1.753) (-1.809,1.781) (-1.831,1.809) (-1.853,1.838) (-1.875,1.866) (-1.898,1.895) (-1.920,1.923) (-1.942,1.951) (-1.964,1.980) (-1.987,2.000)}
\LC{(-1.274,0.641) (-1.300,0.666) (-1.326,0.690) (-1.352,0.715) (-1.379,0.740) (-1.405,0.764) (-1.432,0.788) (-1.458,0.813) (-1.485,0.837) (-1.512,0.861) (-1.538,0.885) (-1.565,0.909) (-1.592,0.933) (-1.619,0.957) (-1.646,0.981) (-1.673,1.005) (-1.700,1.028) (-1.727,1.052) (-1.754,1.076) (-1.781,1.099) (-1.809,1.123) (-1.836,1.146) (-1.863,1.170) (-1.891,1.193) (-1.918,1.216) (-1.946,1.240) (-1.973,1.263) (-2.001,1.286) (-2.028,1.309) (-2.056,1.332) (-2.084,1.355) (-2.112,1.378) (-2.139,1.401) (-2.167,1.424) (-2.195,1.447) (-2.223,1.469) (-2.251,1.492) (-2.279,1.515) (-2.307,1.537) (-2.335,1.560) (-2.363,1.582) (-2.391,1.605) (-2.419,1.627) (-2.447,1.650) (-2.475,1.672) (-2.504,1.695) (-2.532,1.717) (-2.560,1.739) (-2.589,1.761) (-2.600,1.776)}
\LC{(-1.407,0.462) (-1.438,0.481) (-1.469,0.498) (-1.501,0.516) (-1.532,0.534) (-1.563,0.551) (-1.595,0.569) (-1.626,0.586) (-1.658,0.604) (-1.689,0.621) (-1.721,0.638) (-1.753,0.656) (-1.784,0.673) (-1.816,0.690) (-1.847,0.707) (-1.879,0.724) (-1.911,0.741) (-1.943,0.758) (-1.975,0.775) (-2.006,0.792) (-2.038,0.809) (-2.070,0.825) (-2.102,0.842) (-2.134,0.859) (-2.166,0.875) (-2.198,0.892) (-2.230,0.908) (-2.262,0.925) (-2.294,0.941) (-2.326,0.957) (-2.358,0.974) (-2.390,0.990) (-2.422,1.006) (-2.454,1.022) (-2.487,1.038) (-2.519,1.055) (-2.551,1.071) (-2.583,1.086) (-2.600,1.097)}
\LC{(-1.493,0.261) (-1.527,0.271) (-1.562,0.280) (-1.597,0.290) (-1.631,0.300) (-1.666,0.310) (-1.701,0.319) (-1.735,0.329) (-1.770,0.339) (-1.805,0.348) (-1.839,0.358) (-1.874,0.367) (-1.909,0.377) (-1.943,0.386) (-1.978,0.396) (-2.013,0.405) (-2.048,0.415) (-2.082,0.424) (-2.117,0.434) (-2.152,0.443) (-2.187,0.452) (-2.222,0.462) (-2.256,0.471) (-2.291,0.480) (-2.326,0.489) (-2.361,0.498) (-2.396,0.507) (-2.430,0.517) (-2.465,0.526) (-2.500,0.535) (-2.535,0.544) (-2.570,0.553) (-2.600,0.562)}
\LC{(-1.529,0.044) (-1.565,0.045) (-1.601,0.046) (-1.637,0.048) (-1.673,0.049) (-1.709,0.050) (-1.745,0.052) (-1.781,0.053) (-1.817,0.055) (-1.853,0.057) (-1.889,0.058) (-1.925,0.060) (-1.961,0.061) (-1.997,0.063) (-2.033,0.064) (-2.068,0.066) (-2.104,0.067) (-2.140,0.069) (-2.176,0.070) (-2.212,0.072) (-2.248,0.073) (-2.284,0.075) (-2.320,0.076) (-2.356,0.078) (-2.392,0.079) (-2.428,0.081) (-2.464,0.082) (-2.500,0.084) (-2.536,0.085) (-2.572,0.087) (-2.600,0.088)}
\LC{(-1.512,-0.181) (-1.548,-0.187) (-1.583,-0.194) (-1.618,-0.200) (-1.654,-0.207) (-1.689,-0.213) (-1.725,-0.220) (-1.760,-0.226) (-1.795,-0.233) (-1.831,-0.240) (-1.866,-0.246) (-1.902,-0.253) (-1.937,-0.259) (-1.972,-0.266) (-2.008,-0.272) (-2.043,-0.279) (-2.079,-0.285) (-2.114,-0.292) (-2.150,-0.298) (-2.185,-0.305) (-2.220,-0.311) (-2.256,-0.317) (-2.291,-0.324) (-2.327,-0.330) (-2.362,-0.336) (-2.398,-0.343) (-2.433,-0.349) (-2.468,-0.355) (-2.504,-0.361) (-2.539,-0.367) (-2.575,-0.374) (-2.600,-0.380)}
\LC{(-1.445,-0.389) (-1.478,-0.404) (-1.511,-0.418) (-1.543,-0.433) (-1.576,-0.448) (-1.609,-0.462) (-1.642,-0.477) (-1.675,-0.492) (-1.708,-0.506) (-1.741,-0.521) (-1.774,-0.535) (-1.807,-0.549) (-1.840,-0.564) (-1.873,-0.578) (-1.906,-0.592) (-1.939,-0.607) (-1.972,-0.621) (-2.005,-0.635) (-2.038,-0.649) (-2.071,-0.663) (-2.105,-0.677) (-2.138,-0.691) (-2.171,-0.705) (-2.204,-0.719) (-2.237,-0.733) (-2.271,-0.747) (-2.304,-0.760) (-2.337,-0.774) (-2.370,-0.788) (-2.404,-0.801) (-2.437,-0.815) (-2.471,-0.828) (-2.504,-0.842) (-2.537,-0.855) (-2.571,-0.869) (-2.600,-0.882)}
\LC{(-1.330,-0.576) (-1.358,-0.599) (-1.386,-0.621) (-1.415,-0.643) (-1.443,-0.665) (-1.472,-0.687) (-1.500,-0.709) (-1.529,-0.731) (-1.558,-0.753) (-1.586,-0.774) (-1.615,-0.796) (-1.644,-0.818) (-1.673,-0.839) (-1.702,-0.861) (-1.731,-0.882) (-1.760,-0.903) (-1.789,-0.925) (-1.818,-0.946) (-1.847,-0.967) (-1.876,-0.988) (-1.905,-1.009) (-1.934,-1.030) (-1.964,-1.051) (-1.993,-1.072) (-2.022,-1.093) (-2.052,-1.114) (-2.081,-1.134) (-2.111,-1.155) (-2.140,-1.176) (-2.170,-1.196) (-2.199,-1.217) (-2.229,-1.237) (-2.258,-1.258) (-2.288,-1.278) (-2.318,-1.299) (-2.347,-1.319) (-2.377,-1.339) (-2.407,-1.359) (-2.437,-1.380) (-2.467,-1.400) (-2.497,-1.420) (-2.526,-1.440) (-2.556,-1.460) (-2.586,-1.480) (-2.600,-1.493)}
\LC{(-1.170,-0.735) (-1.192,-0.763) (-1.214,-0.791) (-1.237,-0.820) (-1.260,-0.848) (-1.282,-0.876) (-1.305,-0.904) (-1.328,-0.931) (-1.351,-0.959) (-1.374,-0.987) (-1.397,-1.015) (-1.420,-1.042) (-1.443,-1.070) (-1.466,-1.097) (-1.489,-1.125) (-1.513,-1.152) (-1.536,-1.179) (-1.560,-1.206) (-1.583,-1.234) (-1.607,-1.261) (-1.631,-1.288) (-1.655,-1.315) (-1.678,-1.342) (-1.702,-1.369) (-1.726,-1.396) (-1.750,-1.423) (-1.774,-1.449) (-1.798,-1.476) (-1.823,-1.503) (-1.847,-1.529) (-1.871,-1.556) (-1.895,-1.583) (-1.920,-1.609) (-1.944,-1.636) (-1.968,-1.662) (-1.993,-1.689) (-2.017,-1.715) (-2.042,-1.741) (-2.066,-1.768) (-2.091,-1.794) (-2.116,-1.820) (-2.140,-1.846) (-2.165,-1.873) (-2.190,-1.899) (-2.215,-1.925) (-2.239,-1.951) (-2.264,-1.977) (-2.289,-2.000)}
\LC{(-0.973,-0.852) (-0.988,-0.885) (-1.002,-0.918) (-1.017,-0.951) (-1.032,-0.983) (-1.047,-1.016) (-1.063,-1.048) (-1.078,-1.081) (-1.094,-1.113) (-1.110,-1.146) (-1.126,-1.178) (-1.142,-1.210) (-1.158,-1.242) (-1.174,-1.274) (-1.191,-1.306) (-1.207,-1.338) (-1.224,-1.370) (-1.241,-1.402) (-1.258,-1.434) (-1.275,-1.466) (-1.292,-1.497) (-1.309,-1.529) (-1.326,-1.561) (-1.343,-1.592) (-1.360,-1.624) (-1.378,-1.655) (-1.395,-1.687) (-1.413,-1.718) (-1.431,-1.750) (-1.448,-1.781) (-1.466,-1.812) (-1.484,-1.844) (-1.501,-1.875) (-1.519,-1.906) (-1.537,-1.937) (-1.555,-1.969) (-1.573,-2.000) (-1.579,-2.000)}
\LC{(-0.745,-0.919) (-0.751,-0.954) (-0.756,-0.990) (-0.763,-1.025) (-0.769,-1.061) (-0.776,-1.096) (-0.783,-1.131) (-0.790,-1.166) (-0.798,-1.202) (-0.806,-1.237) (-0.814,-1.272) (-0.822,-1.307) (-0.831,-1.342) (-0.839,-1.377) (-0.848,-1.412) (-0.857,-1.447) (-0.866,-1.481) (-0.876,-1.516) (-0.885,-1.551) (-0.895,-1.586) (-0.904,-1.620) (-0.914,-1.655) (-0.924,-1.690) (-0.934,-1.724) (-0.944,-1.759) (-0.954,-1.793) (-0.964,-1.828) (-0.975,-1.862) (-0.985,-1.897) (-0.996,-1.931) (-1.006,-1.965) (-1.017,-2.000) (-1.021,-2.000)}
\LC{(-0.487,-0.923) (-0.485,-0.959) (-0.482,-0.995) (-0.480,-1.031) (-0.479,-1.067) (-0.477,-1.103) (-0.476,-1.139) (-0.476,-1.175) (-0.475,-1.211) (-0.475,-1.247) (-0.476,-1.283) (-0.476,-1.319) (-0.477,-1.355) (-0.478,-1.391) (-0.479,-1.427) (-0.481,-1.463) (-0.482,-1.499) (-0.484,-1.535) (-0.486,-1.571) (-0.488,-1.606) (-0.491,-1.642) (-0.493,-1.678) (-0.496,-1.714) (-0.499,-1.750) (-0.501,-1.786) (-0.504,-1.822) (-0.507,-1.858) (-0.511,-1.894) (-0.514,-1.929) (-0.517,-1.965) (-0.521,-2.000)}
\LC{(-0.215,-0.847) (-0.203,-0.881) (-0.192,-0.915) (-0.182,-0.950) (-0.172,-0.985) (-0.163,-1.019) (-0.155,-1.054) (-0.147,-1.090) (-0.139,-1.125) (-0.132,-1.160) (-0.125,-1.195) (-0.119,-1.231) (-0.112,-1.266) (-0.107,-1.302) (-0.101,-1.337) (-0.095,-1.373) (-0.090,-1.409) (-0.085,-1.444) (-0.080,-1.480) (-0.076,-1.516) (-0.071,-1.551) (-0.067,-1.587) (-0.063,-1.623) (-0.059,-1.659) (-0.055,-1.694) (-0.051,-1.730) (-0.047,-1.766) (-0.044,-1.802) (-0.040,-1.838) (-0.037,-1.873) (-0.033,-1.909) (-0.030,-1.945) (-0.027,-1.981) (-0.025,-2.000)}
\LC{(0.096,-0.718) (0.109,-0.752) (0.122,-0.785) (0.135,-0.819) (0.147,-0.853) (0.160,-0.887) (0.172,-0.920) (0.184,-0.954) (0.195,-0.989) (0.206,-1.023) (0.218,-1.057) (0.229,-1.091) (0.239,-1.126) (0.250,-1.160) (0.260,-1.194) (0.271,-1.229) (0.281,-1.263) (0.291,-1.298) (0.301,-1.333) (0.311,-1.367) (0.321,-1.402) (0.330,-1.437) (0.340,-1.471) (0.349,-1.506) (0.359,-1.541) (0.368,-1.576) (0.377,-1.610) (0.386,-1.645) (0.396,-1.680) (0.405,-1.715) (0.414,-1.750) (0.423,-1.785) (0.432,-1.819) (0.441,-1.854) (0.449,-1.889) (0.458,-1.924) (0.467,-1.959) (0.476,-1.994) (0.479,-2.000)}
\LC{(0.432,-0.579) (0.446,-0.612) (0.460,-0.645) (0.474,-0.678) (0.488,-0.712) (0.501,-0.745) (0.515,-0.778) (0.529,-0.812) (0.542,-0.845) (0.556,-0.878) (0.570,-0.912) (0.583,-0.945) (0.597,-0.978) (0.611,-1.011) (0.624,-1.045) (0.638,-1.078) (0.652,-1.111) (0.665,-1.145) (0.679,-1.178) (0.693,-1.211) (0.706,-1.245) (0.720,-1.278) (0.734,-1.311) (0.748,-1.344) (0.761,-1.378) (0.775,-1.411) (0.789,-1.444) (0.803,-1.477) (0.817,-1.511) (0.831,-1.544) (0.845,-1.577) (0.858,-1.610) (0.872,-1.643) (0.886,-1.677) (0.900,-1.710) (0.914,-1.743) (0.928,-1.776) (0.942,-1.809) (0.956,-1.842) (0.970,-1.876) (0.984,-1.909) (0.998,-1.942) (1.013,-1.975) (1.027,-2.000)}
\LC{(0.775,-0.437) (0.789,-0.470) (0.804,-0.503) (0.818,-0.536) (0.833,-0.569) (0.848,-0.601) (0.863,-0.634) (0.878,-0.667) (0.894,-0.699) (0.909,-0.732) (0.925,-0.764) (0.941,-0.796) (0.957,-0.828) (0.974,-0.860) (0.990,-0.892) (1.007,-0.924) (1.024,-0.956) (1.041,-0.988) (1.058,-1.019) (1.075,-1.051) (1.092,-1.083) (1.110,-1.114) (1.128,-1.145) (1.146,-1.177) (1.164,-1.208) (1.182,-1.239) (1.200,-1.270) (1.218,-1.301) (1.236,-1.332) (1.255,-1.363) (1.273,-1.394) (1.292,-1.425) (1.311,-1.455) (1.330,-1.486) (1.349,-1.517) (1.368,-1.547) (1.387,-1.578) (1.406,-1.608) (1.425,-1.639) (1.444,-1.669) (1.464,-1.699) (1.483,-1.730) (1.503,-1.760) (1.522,-1.790) (1.542,-1.820) (1.561,-1.851) (1.581,-1.881) (1.601,-1.911) (1.621,-1.941) (1.640,-1.971) (1.660,-2.000)}
\LC{(1.098,-0.303) (1.113,-0.336) (1.128,-0.368) (1.143,-0.401) (1.159,-0.433) (1.176,-0.465) (1.192,-0.497) (1.210,-0.529) (1.228,-0.560) (1.246,-0.591) (1.264,-0.622) (1.283,-0.652) (1.303,-0.683) (1.322,-0.713) (1.342,-0.743) (1.363,-0.773) (1.383,-0.802) (1.404,-0.831) (1.425,-0.860) (1.447,-0.889) (1.468,-0.918) (1.490,-0.947) (1.512,-0.975) (1.535,-1.003) (1.557,-1.031) (1.580,-1.059) (1.603,-1.087) (1.626,-1.114) (1.650,-1.142) (1.673,-1.169) (1.697,-1.196) (1.720,-1.224) (1.744,-1.250) (1.768,-1.277) (1.792,-1.304) (1.817,-1.331) (1.841,-1.357) (1.865,-1.384) (1.890,-1.410) (1.915,-1.436) (1.939,-1.462) (1.964,-1.488) (1.989,-1.514) (2.014,-1.540) (2.039,-1.566) (2.065,-1.592) (2.090,-1.617) (2.115,-1.643) (2.141,-1.668) (2.166,-1.694) (2.192,-1.719) (2.217,-1.744) (2.243,-1.770) (2.269,-1.795) (2.294,-1.820) (2.320,-1.845) (2.346,-1.870) (2.372,-1.895) (2.398,-1.920) (2.424,-1.945) (2.450,-1.970) (2.476,-1.994) (2.485,-2.000)}
\LC{(1.378,-0.187) (1.393,-0.219) (1.410,-0.251) (1.427,-0.283) (1.445,-0.314) (1.465,-0.344) (1.485,-0.374) (1.505,-0.404) (1.527,-0.432) (1.549,-0.461) (1.572,-0.488) (1.596,-0.516) (1.620,-0.543) (1.644,-0.569) (1.669,-0.595) (1.694,-0.620) (1.720,-0.646) (1.746,-0.670) (1.773,-0.695) (1.800,-0.719) (1.827,-0.743) (1.854,-0.766) (1.881,-0.789) (1.909,-0.812) (1.937,-0.835) (1.965,-0.858) (1.993,-0.880) (2.022,-0.902) (2.050,-0.924) (2.079,-0.945) (2.108,-0.967) (2.137,-0.988) (2.166,-1.009) (2.195,-1.030) (2.225,-1.051) (2.254,-1.072) (2.284,-1.092) (2.313,-1.113) (2.343,-1.133) (2.373,-1.153) (2.403,-1.173) (2.433,-1.193) (2.463,-1.213) (2.493,-1.233) (2.523,-1.252) (2.554,-1.272) (2.584,-1.291) (2.600,-1.304)}
\LC{(1.600,-0.095) (1.617,-0.126) (1.637,-0.156) (1.660,-0.184) (1.684,-0.211) (1.709,-0.237) (1.736,-0.261) (1.764,-0.284) (1.792,-0.306) (1.821,-0.327) (1.851,-0.347) (1.881,-0.366) (1.912,-0.385) (1.943,-0.403) (1.975,-0.421) (2.006,-0.438) (2.038,-0.455) (2.070,-0.471) (2.102,-0.487) (2.135,-0.503) (2.167,-0.519) (2.200,-0.534) (2.233,-0.549) (2.265,-0.563) (2.298,-0.578) (2.331,-0.592) (2.365,-0.606) (2.398,-0.620) (2.431,-0.634) (2.464,-0.647) (2.498,-0.661) (2.531,-0.674) (2.565,-0.687) (2.598,-0.700) (2.600,-0.704)}
\LC{(1.739,-0.037) (1.769,-0.057) (1.801,-0.073) (1.835,-0.085) (1.869,-0.096) (1.904,-0.106) (1.939,-0.115) (1.974,-0.123) (2.009,-0.131) (2.044,-0.138) (2.080,-0.145) (2.115,-0.152) (2.150,-0.159) (2.186,-0.165) (2.221,-0.171) (2.257,-0.177) (2.292,-0.182) (2.328,-0.188) (2.364,-0.193) (2.399,-0.199) (2.435,-0.204) (2.470,-0.209) (2.506,-0.214) (2.542,-0.219) (2.577,-0.224) (2.600,-0.227)}
\LC{(1.657,0.071) (1.677,0.101) (1.700,0.128) (1.726,0.154) (1.754,0.177) (1.782,0.198) (1.812,0.218) (1.843,0.237) (1.875,0.254) (1.907,0.271) (1.939,0.287) (1.971,0.302) (2.004,0.317) (2.037,0.331) (2.071,0.345) (2.104,0.358) (2.138,0.371) (2.171,0.384) (2.205,0.397) (2.239,0.409) (2.273,0.421) (2.307,0.432) (2.341,0.444) (2.375,0.455) (2.409,0.466) (2.444,0.477) (2.478,0.488) (2.512,0.499) (2.547,0.509) (2.581,0.520) (2.600,0.527)}
\LC{(1.466,0.150) (1.482,0.182) (1.499,0.214) (1.518,0.245) (1.538,0.275) (1.559,0.304) (1.580,0.333) (1.603,0.361) (1.627,0.388) (1.651,0.414) (1.677,0.440) (1.702,0.465) (1.729,0.490) (1.755,0.514) (1.783,0.537) (1.810,0.561) (1.838,0.583) (1.866,0.605) (1.895,0.627) (1.924,0.649) (1.953,0.670) (1.982,0.691) (2.012,0.711) (2.042,0.732) (2.072,0.752) (2.102,0.772) (2.132,0.791) (2.162,0.811) (2.193,0.830) (2.223,0.849) (2.254,0.867) (2.285,0.886) (2.316,0.905) (2.347,0.923) (2.378,0.941) (2.409,0.959) (2.440,0.977) (2.472,0.995) (2.503,1.012) (2.534,1.030) (2.566,1.047) (2.598,1.064) (2.600,1.070)}
\LC{(1.206,0.258) (1.221,0.291) (1.237,0.323) (1.253,0.355) (1.270,0.387) (1.287,0.419) (1.305,0.450) (1.323,0.481) (1.342,0.512) (1.361,0.542) (1.381,0.572) (1.402,0.602) (1.422,0.631) (1.444,0.660) (1.465,0.689) (1.487,0.717) (1.510,0.746) (1.532,0.774) (1.555,0.801) (1.578,0.829) (1.602,0.856) (1.626,0.883) (1.650,0.910) (1.674,0.937) (1.698,0.963) (1.723,0.989) (1.748,1.015) (1.773,1.041) (1.798,1.067) (1.824,1.092) (1.849,1.118) (1.875,1.143) (1.901,1.168) (1.927,1.193) (1.953,1.218) (1.979,1.242) (2.005,1.267) (2.032,1.291) (2.058,1.316) (2.085,1.340) (2.112,1.364) (2.138,1.388) (2.165,1.412) (2.192,1.436) (2.219,1.460) (2.247,1.483) (2.274,1.507) (2.301,1.530) (2.328,1.554) (2.356,1.577) (2.383,1.600) (2.411,1.623) (2.439,1.646) (2.466,1.669) (2.494,1.692) (2.522,1.715) (2.550,1.738) (2.578,1.761) (2.600,1.783)}
\LC{(0.902,0.384) (0.916,0.417) (0.930,0.450) (0.945,0.483) (0.960,0.516) (0.975,0.548) (0.991,0.581) (1.007,0.613) (1.023,0.645) (1.039,0.677) (1.056,0.709) (1.073,0.741) (1.090,0.773) (1.108,0.804) (1.126,0.835) (1.144,0.867) (1.162,0.898) (1.180,0.929) (1.199,0.959) (1.217,0.990) (1.236,1.021) (1.255,1.051) (1.275,1.082) (1.294,1.112) (1.314,1.142) (1.334,1.172) (1.353,1.202) (1.373,1.232) (1.394,1.262) (1.414,1.292) (1.434,1.321) (1.455,1.351) (1.476,1.380) (1.496,1.410) (1.517,1.439) (1.538,1.468) (1.559,1.498) (1.580,1.527) (1.602,1.556) (1.623,1.585) (1.644,1.614) (1.666,1.643) (1.687,1.672) (1.709,1.700) (1.731,1.729) (1.752,1.758) (1.774,1.786) (1.796,1.815) (1.818,1.844) (1.840,1.872) (1.862,1.901) (1.884,1.929) (1.906,1.957) (1.928,1.986) (1.943,2.000)}
\LC{(0.565,0.524) (0.579,0.557) (0.593,0.590) (0.607,0.623) (0.621,0.656) (0.635,0.689) (0.650,0.722) (0.664,0.756) (0.678,0.789) (0.693,0.822) (0.707,0.854) (0.722,0.887) (0.736,0.920) (0.751,0.953) (0.766,0.986) (0.781,1.019) (0.796,1.052) (0.811,1.084) (0.826,1.117) (0.841,1.150) (0.856,1.182) (0.871,1.215) (0.886,1.248) (0.901,1.280) (0.917,1.313) (0.932,1.345) (0.948,1.378) (0.963,1.410) (0.979,1.443) (0.995,1.475) (1.010,1.508) (1.026,1.540) (1.042,1.572) (1.058,1.605) (1.073,1.637) (1.089,1.669) (1.105,1.702) (1.121,1.734) (1.137,1.766) (1.153,1.798) (1.169,1.830) (1.186,1.863) (1.202,1.895) (1.218,1.927) (1.234,1.959) (1.250,1.991) (1.256,2.000)}
\LC{(0.223,0.666) (0.236,0.699) (0.249,0.733) (0.262,0.766) (0.275,0.800) (0.288,0.833) (0.301,0.867) (0.314,0.901) (0.326,0.934) (0.338,0.968) (0.351,1.002) (0.363,1.036) (0.375,1.070) (0.387,1.104) (0.399,1.138) (0.410,1.172) (0.422,1.206) (0.433,1.240) (0.445,1.274) (0.456,1.308) (0.468,1.343) (0.479,1.377) (0.490,1.411) (0.501,1.445) (0.513,1.479) (0.524,1.514) (0.535,1.548) (0.546,1.582) (0.557,1.616) (0.568,1.651) (0.579,1.685) (0.590,1.719) (0.601,1.754) (0.611,1.788) (0.622,1.822) (0.633,1.857) (0.644,1.891) (0.655,1.925) (0.665,1.960) (0.676,1.994) (0.680,2.000)}
\LC{(-0.107,0.803) (-0.095,0.836) (-0.082,0.870) (-0.071,0.904) (-0.060,0.938) (-0.049,0.973) (-0.039,1.007) (-0.029,1.042) (-0.019,1.077) (-0.010,1.111) (-0.002,1.146) (0.007,1.181) (0.015,1.216) (0.023,1.252) (0.031,1.287) (0.038,1.322) (0.045,1.357) (0.052,1.393) (0.059,1.428) (0.066,1.463) (0.073,1.499) (0.079,1.534) (0.085,1.569) (0.091,1.605) (0.098,1.640) (0.104,1.676) (0.109,1.711) (0.115,1.747) (0.121,1.783) (0.126,1.818) (0.132,1.854) (0.137,1.889) (0.143,1.925) (0.148,1.960) (0.153,1.996) (0.155,2.000)}

      % Conductor con punta - flechas
% Generado por scripts/generar_datos.py (no editar a mano).
\FL{-0.332}{1.586}{89.3}
\FL{-0.742}{1.606}{100.9}
\FL{-1.161}{1.589}{113.5}
\FL{-1.635}{1.552}{126.9}
\FL{-2.093}{1.363}{140.4}
\FL{-2.144}{0.864}{152.5}
\FL{-2.187}{0.452}{165.0}
\FL{-2.200}{0.071}{177.6}
\FL{-2.197}{-0.307}{-169.8}
\FL{-2.160}{-0.700}{-157.2}
\FL{-2.111}{-1.155}{-144.9}
\FL{-1.855}{-1.538}{-132.3}
\FL{-1.337}{-1.582}{-118.7}
\FL{-0.898}{-1.597}{-105.6}
\FL{-0.488}{-1.595}{-93.5}
\FL{-0.070}{-1.563}{-83.0}
\FL{0.352}{-1.518}{-74.8}
\FL{0.798}{-1.466}{-67.4}
\FL{1.292}{-1.425}{-58.7}
\FL{1.890}{-1.410}{-46.8}
\FL{2.079}{-0.945}{-36.8}
\FL{2.189}{-0.529}{-24.8}
\FL{2.269}{-0.179}{-9.2}
\FL{2.216}{0.401}{19.8}
\FL{2.122}{0.785}{32.9}
\FL{2.005}{1.267}{42.9}
\FL{1.503}{1.420}{54.6}
\FL{0.984}{1.454}{64.3}
\FL{0.516}{1.491}{72.0}
\FL{0.081}{1.546}{79.9}

    \end{scope}
    % Contorno de la gota
% Generado por scripts/generar_datos.py (no editar a mano).
\def\formagota{(-0.259,0.833) -- (-0.262,0.834) -- (-0.265,0.835) -- (-0.268,0.836) -- (-0.270,0.837) -- (-0.273,0.839) -- (-0.276,0.840) -- (-0.279,0.841) -- (-0.281,0.842) -- (-0.284,0.843) -- (-0.287,0.844) -- (-0.290,0.845) -- (-0.293,0.846) -- (-0.295,0.847) -- (-0.298,0.848) -- (-0.301,0.849) -- (-0.304,0.850) -- (-0.306,0.851) -- (-0.309,0.852) -- (-0.312,0.853) -- (-0.315,0.854) -- (-0.318,0.855) -- (-0.320,0.855) -- (-0.323,0.856) -- (-0.326,0.857) -- (-0.329,0.858) -- (-0.332,0.859) -- (-0.335,0.860) -- (-0.337,0.861) -- (-0.340,0.862) -- (-0.343,0.863) -- (-0.346,0.863) -- (-0.349,0.864) -- (-0.352,0.865) -- (-0.354,0.866) -- (-0.357,0.867) -- (-0.360,0.867) -- (-0.363,0.868) -- (-0.366,0.869) -- (-0.369,0.870) -- (-0.371,0.870) -- (-0.374,0.871) -- (-0.377,0.872) -- (-0.380,0.873) -- (-0.383,0.873) -- (-0.386,0.874) -- (-0.389,0.875) -- (-0.392,0.876) -- (-0.394,0.876) -- (-0.397,0.877) -- (-0.400,0.878) -- (-0.403,0.878) -- (-0.406,0.879) -- (-0.409,0.879) -- (-0.412,0.880) -- (-0.415,0.881) -- (-0.418,0.881) -- (-0.420,0.882) -- (-0.423,0.882) -- (-0.426,0.883) -- (-0.429,0.884) -- (-0.432,0.884) -- (-0.435,0.885) -- (-0.438,0.885) -- (-0.441,0.886) -- (-0.444,0.886) -- (-0.447,0.887) -- (-0.449,0.887) -- (-0.452,0.888) -- (-0.455,0.888) -- (-0.458,0.889) -- (-0.461,0.889) -- (-0.464,0.890) -- (-0.467,0.890) -- (-0.470,0.891) -- (-0.473,0.891) -- (-0.476,0.891) -- (-0.479,0.892) -- (-0.482,0.892) -- (-0.485,0.893) -- (-0.487,0.893) -- (-0.490,0.893) -- (-0.493,0.894) -- (-0.496,0.894) -- (-0.499,0.894) -- (-0.502,0.895) -- (-0.505,0.895) -- (-0.508,0.895) -- (-0.511,0.896) -- (-0.514,0.896) -- (-0.517,0.896) -- (-0.520,0.896) -- (-0.523,0.897) -- (-0.526,0.897) -- (-0.529,0.897) -- (-0.532,0.897) -- (-0.535,0.898) -- (-0.537,0.898) -- (-0.540,0.898) -- (-0.543,0.898) -- (-0.546,0.898) -- (-0.549,0.899) -- (-0.552,0.899) -- (-0.555,0.899) -- (-0.558,0.899) -- (-0.561,0.899) -- (-0.564,0.899) -- (-0.567,0.899) -- (-0.570,0.899) -- (-0.573,0.900) -- (-0.576,0.900) -- (-0.579,0.900) -- (-0.582,0.900) -- (-0.585,0.900) -- (-0.588,0.900) -- (-0.591,0.900) -- (-0.594,0.900) -- (-0.597,0.900) -- (-0.599,0.900) -- (-0.602,0.900) -- (-0.605,0.900) -- (-0.608,0.900) -- (-0.611,0.900) -- (-0.614,0.900) -- (-0.617,0.900) -- (-0.620,0.900) -- (-0.623,0.900) -- (-0.626,0.900) -- (-0.629,0.900) -- (-0.632,0.899) -- (-0.635,0.899) -- (-0.638,0.899) -- (-0.641,0.899) -- (-0.644,0.899) -- (-0.647,0.899) -- (-0.650,0.899) -- (-0.653,0.898) -- (-0.656,0.898) -- (-0.659,0.898) -- (-0.661,0.898) -- (-0.664,0.898) -- (-0.667,0.897) -- (-0.670,0.897) -- (-0.673,0.897) -- (-0.676,0.897) -- (-0.679,0.896) -- (-0.682,0.896) -- (-0.685,0.896) -- (-0.688,0.896) -- (-0.691,0.895) -- (-0.694,0.895) -- (-0.697,0.895) -- (-0.700,0.894) -- (-0.703,0.894) -- (-0.706,0.894) -- (-0.709,0.893) -- (-0.711,0.893) -- (-0.714,0.893) -- (-0.717,0.892) -- (-0.720,0.892) -- (-0.723,0.892) -- (-0.726,0.891) -- (-0.729,0.891) -- (-0.732,0.890) -- (-0.735,0.890) -- (-0.738,0.889) -- (-0.741,0.889) -- (-0.744,0.888) -- (-0.747,0.888) -- (-0.750,0.887) -- (-0.752,0.887) -- (-0.755,0.886) -- (-0.758,0.886) -- (-0.761,0.885) -- (-0.764,0.885) -- (-0.767,0.884) -- (-0.770,0.884) -- (-0.773,0.883) -- (-0.776,0.883) -- (-0.779,0.882) -- (-0.781,0.881) -- (-0.784,0.881) -- (-0.787,0.880) -- (-0.790,0.880) -- (-0.793,0.879) -- (-0.796,0.878) -- (-0.799,0.878) -- (-0.802,0.877) -- (-0.805,0.876) -- (-0.807,0.876) -- (-0.810,0.875) -- (-0.813,0.874) -- (-0.816,0.874) -- (-0.819,0.873) -- (-0.822,0.872) -- (-0.825,0.871) -- (-0.828,0.871) -- (-0.830,0.870) -- (-0.833,0.869) -- (-0.836,0.868) -- (-0.839,0.868) -- (-0.842,0.867) -- (-0.845,0.866) -- (-0.847,0.865) -- (-0.850,0.864) -- (-0.853,0.864) -- (-0.856,0.863) -- (-0.859,0.862) -- (-0.862,0.861) -- (-0.864,0.860) -- (-0.867,0.859) -- (-0.870,0.858) -- (-0.873,0.858) -- (-0.876,0.857) -- (-0.879,0.856) -- (-0.881,0.855) -- (-0.884,0.854) -- (-0.887,0.853) -- (-0.890,0.852) -- (-0.893,0.851) -- (-0.895,0.850) -- (-0.898,0.849) -- (-0.901,0.848) -- (-0.904,0.847) -- (-0.907,0.846) -- (-0.909,0.845) -- (-0.912,0.844) -- (-0.915,0.843) -- (-0.918,0.842) -- (-0.920,0.841) -- (-0.923,0.840) -- (-0.926,0.839) -- (-0.929,0.838) -- (-0.931,0.837) -- (-0.934,0.836) -- (-0.937,0.835) -- (-0.940,0.833) -- (-0.942,0.832) -- (-0.945,0.831) -- (-0.948,0.830) -- (-0.951,0.829) -- (-0.953,0.828) -- (-0.956,0.827) -- (-0.959,0.825) -- (-0.961,0.824) -- (-0.964,0.823) -- (-0.967,0.822) -- (-0.969,0.821) -- (-0.972,0.819) -- (-0.975,0.818) -- (-0.978,0.817) -- (-0.980,0.816) -- (-0.983,0.814) -- (-0.986,0.813) -- (-0.988,0.812) -- (-0.991,0.811) -- (-0.994,0.809) -- (-0.996,0.808) -- (-0.999,0.807) -- (-1.002,0.805) -- (-1.004,0.804) -- (-1.007,0.803) -- (-1.009,0.801) -- (-1.012,0.800) -- (-1.015,0.799) -- (-1.017,0.797) -- (-1.020,0.796) -- (-1.023,0.795) -- (-1.025,0.793) -- (-1.028,0.792) -- (-1.030,0.790) -- (-1.033,0.789) -- (-1.036,0.788) -- (-1.038,0.786) -- (-1.041,0.785) -- (-1.043,0.783) -- (-1.046,0.782) -- (-1.048,0.780) -- (-1.051,0.779) -- (-1.054,0.777) -- (-1.056,0.776) -- (-1.059,0.774) -- (-1.061,0.773) -- (-1.064,0.771) -- (-1.066,0.770) -- (-1.069,0.768) -- (-1.071,0.767) -- (-1.074,0.765) -- (-1.076,0.764) -- (-1.079,0.762) -- (-1.081,0.760) -- (-1.084,0.759) -- (-1.086,0.757) -- (-1.089,0.756) -- (-1.091,0.754) -- (-1.094,0.752) -- (-1.096,0.751) -- (-1.099,0.749) -- (-1.101,0.748) -- (-1.104,0.746) -- (-1.106,0.744) -- (-1.108,0.743) -- (-1.111,0.741) -- (-1.113,0.739) -- (-1.116,0.738) -- (-1.118,0.736) -- (-1.121,0.734) -- (-1.123,0.732) -- (-1.125,0.731) -- (-1.128,0.729) -- (-1.130,0.727) -- (-1.133,0.725) -- (-1.135,0.724) -- (-1.137,0.722) -- (-1.140,0.720) -- (-1.142,0.718) -- (-1.144,0.717) -- (-1.147,0.715) -- (-1.149,0.713) -- (-1.151,0.711) -- (-1.154,0.709) -- (-1.156,0.708) -- (-1.158,0.706) -- (-1.161,0.704) -- (-1.163,0.702) -- (-1.165,0.700) -- (-1.168,0.698) -- (-1.170,0.697) -- (-1.172,0.695) -- (-1.174,0.693) -- (-1.177,0.691) -- (-1.179,0.689) -- (-1.181,0.687) -- (-1.184,0.685) -- (-1.186,0.683) -- (-1.188,0.681) -- (-1.190,0.679) -- (-1.192,0.677) -- (-1.195,0.675) -- (-1.197,0.674) -- (-1.199,0.672) -- (-1.201,0.670) -- (-1.204,0.668) -- (-1.206,0.666) -- (-1.208,0.664) -- (-1.210,0.662) -- (-1.212,0.660) -- (-1.214,0.658) -- (-1.217,0.656) -- (-1.219,0.654) -- (-1.221,0.652) -- (-1.223,0.649) -- (-1.225,0.647) -- (-1.227,0.645) -- (-1.229,0.643) -- (-1.231,0.641) -- (-1.234,0.639) -- (-1.236,0.637) -- (-1.238,0.635) -- (-1.240,0.633) -- (-1.242,0.631) -- (-1.244,0.629) -- (-1.246,0.627) -- (-1.248,0.624) -- (-1.250,0.622) -- (-1.252,0.620) -- (-1.254,0.618) -- (-1.256,0.616) -- (-1.258,0.614) -- (-1.260,0.612) -- (-1.262,0.609) -- (-1.264,0.607) -- (-1.266,0.605) -- (-1.268,0.603) -- (-1.270,0.601) -- (-1.272,0.598) -- (-1.274,0.596) -- (-1.276,0.594) -- (-1.278,0.592) -- (-1.280,0.590) -- (-1.282,0.587) -- (-1.284,0.585) -- (-1.286,0.583) -- (-1.288,0.581) -- (-1.290,0.578) -- (-1.291,0.576) -- (-1.293,0.574) -- (-1.295,0.571) -- (-1.297,0.569) -- (-1.299,0.567) -- (-1.301,0.565) -- (-1.303,0.562) -- (-1.305,0.560) -- (-1.306,0.558) -- (-1.308,0.555) -- (-1.310,0.553) -- (-1.312,0.551) -- (-1.314,0.548) -- (-1.315,0.546) -- (-1.317,0.544) -- (-1.319,0.541) -- (-1.321,0.539) -- (-1.323,0.537) -- (-1.324,0.534) -- (-1.326,0.532) -- (-1.328,0.529) -- (-1.330,0.527) -- (-1.331,0.525) -- (-1.333,0.522) -- (-1.335,0.520) -- (-1.336,0.517) -- (-1.338,0.515) -- (-1.340,0.513) -- (-1.341,0.510) -- (-1.343,0.508) -- (-1.345,0.505) -- (-1.346,0.503) -- (-1.348,0.500) -- (-1.350,0.498) -- (-1.351,0.495) -- (-1.353,0.493) -- (-1.355,0.490) -- (-1.356,0.488) -- (-1.358,0.485) -- (-1.359,0.483) -- (-1.361,0.480) -- (-1.363,0.478) -- (-1.364,0.475) -- (-1.366,0.473) -- (-1.367,0.470) -- (-1.369,0.468) -- (-1.370,0.465) -- (-1.372,0.463) -- (-1.373,0.460) -- (-1.375,0.458) -- (-1.376,0.455) -- (-1.378,0.453) -- (-1.379,0.450) -- (-1.381,0.448) -- (-1.382,0.445) -- (-1.384,0.442) -- (-1.385,0.440) -- (-1.387,0.437) -- (-1.388,0.435) -- (-1.389,0.432) -- (-1.391,0.430) -- (-1.392,0.427) -- (-1.394,0.424) -- (-1.395,0.422) -- (-1.396,0.419) -- (-1.398,0.416) -- (-1.399,0.414) -- (-1.401,0.411) -- (-1.402,0.409) -- (-1.403,0.406) -- (-1.405,0.403) -- (-1.406,0.401) -- (-1.407,0.398) -- (-1.408,0.395) -- (-1.410,0.393) -- (-1.411,0.390) -- (-1.412,0.387) -- (-1.414,0.385) -- (-1.415,0.382) -- (-1.416,0.379) -- (-1.417,0.377) -- (-1.419,0.374) -- (-1.420,0.371) -- (-1.421,0.369) -- (-1.422,0.366) -- (-1.423,0.363) -- (-1.425,0.361) -- (-1.426,0.358) -- (-1.427,0.355) -- (-1.428,0.352) -- (-1.429,0.350) -- (-1.430,0.347) -- (-1.432,0.344) -- (-1.433,0.341) -- (-1.434,0.339) -- (-1.435,0.336) -- (-1.436,0.333) -- (-1.437,0.330) -- (-1.438,0.328) -- (-1.439,0.325) -- (-1.440,0.322) -- (-1.441,0.319) -- (-1.442,0.317) -- (-1.443,0.314) -- (-1.444,0.311) -- (-1.445,0.308) -- (-1.447,0.306) -- (-1.448,0.303) -- (-1.448,0.300) -- (-1.449,0.297) -- (-1.450,0.294) -- (-1.451,0.292) -- (-1.452,0.289) -- (-1.453,0.286) -- (-1.454,0.283) -- (-1.455,0.280) -- (-1.456,0.278) -- (-1.457,0.275) -- (-1.458,0.272) -- (-1.459,0.269) -- (-1.460,0.266) -- (-1.461,0.264) -- (-1.461,0.261) -- (-1.462,0.258) -- (-1.463,0.255) -- (-1.464,0.252) -- (-1.465,0.249) -- (-1.466,0.247) -- (-1.466,0.244) -- (-1.467,0.241) -- (-1.468,0.238) -- (-1.469,0.235) -- (-1.469,0.232) -- (-1.470,0.229) -- (-1.471,0.227) -- (-1.472,0.224) -- (-1.472,0.221) -- (-1.473,0.218) -- (-1.474,0.215) -- (-1.475,0.212) -- (-1.475,0.209) -- (-1.476,0.207) -- (-1.477,0.204) -- (-1.477,0.201) -- (-1.478,0.198) -- (-1.479,0.195) -- (-1.479,0.192) -- (-1.480,0.189) -- (-1.480,0.186) -- (-1.481,0.183) -- (-1.482,0.181) -- (-1.482,0.178) -- (-1.483,0.175) -- (-1.483,0.172) -- (-1.484,0.169) -- (-1.485,0.166) -- (-1.485,0.163) -- (-1.486,0.160) -- (-1.486,0.157) -- (-1.487,0.154) -- (-1.487,0.151) -- (-1.488,0.149) -- (-1.488,0.146) -- (-1.489,0.143) -- (-1.489,0.140) -- (-1.490,0.137) -- (-1.490,0.134) -- (-1.490,0.131) -- (-1.491,0.128) -- (-1.491,0.125) -- (-1.492,0.122) -- (-1.492,0.119) -- (-1.492,0.116) -- (-1.493,0.113) -- (-1.493,0.111) -- (-1.494,0.108) -- (-1.494,0.105) -- (-1.494,0.102) -- (-1.495,0.099) -- (-1.495,0.096) -- (-1.495,0.093) -- (-1.495,0.090) -- (-1.496,0.087) -- (-1.496,0.084) -- (-1.496,0.081) -- (-1.497,0.078) -- (-1.497,0.075) -- (-1.497,0.072) -- (-1.497,0.069) -- (-1.498,0.066) -- (-1.498,0.063) -- (-1.498,0.061) -- (-1.498,0.058) -- (-1.498,0.055) -- (-1.498,0.052) -- (-1.499,0.049) -- (-1.499,0.046) -- (-1.499,0.043) -- (-1.499,0.040) -- (-1.499,0.037) -- (-1.499,0.034) -- (-1.499,0.031) -- (-1.500,0.028) -- (-1.500,0.025) -- (-1.500,0.022) -- (-1.500,0.019) -- (-1.500,0.016) -- (-1.500,0.013) -- (-1.500,0.010) -- (-1.500,0.007) -- (-1.500,0.004) -- (-1.500,0.001) -- (-1.500,-0.001) -- (-1.500,-0.004) -- (-1.500,-0.007) -- (-1.500,-0.010) -- (-1.500,-0.013) -- (-1.500,-0.016) -- (-1.500,-0.019) -- (-1.500,-0.022) -- (-1.500,-0.025) -- (-1.500,-0.028) -- (-1.499,-0.031) -- (-1.499,-0.034) -- (-1.499,-0.037) -- (-1.499,-0.040) -- (-1.499,-0.043) -- (-1.499,-0.046) -- (-1.499,-0.049) -- (-1.498,-0.052) -- (-1.498,-0.055) -- (-1.498,-0.058) -- (-1.498,-0.061) -- (-1.498,-0.063) -- (-1.498,-0.066) -- (-1.497,-0.069) -- (-1.497,-0.072) -- (-1.497,-0.075) -- (-1.497,-0.078) -- (-1.496,-0.081) -- (-1.496,-0.084) -- (-1.496,-0.087) -- (-1.495,-0.090) -- (-1.495,-0.093) -- (-1.495,-0.096) -- (-1.495,-0.099) -- (-1.494,-0.102) -- (-1.494,-0.105) -- (-1.494,-0.108) -- (-1.493,-0.111) -- (-1.493,-0.113) -- (-1.492,-0.116) -- (-1.492,-0.119) -- (-1.492,-0.122) -- (-1.491,-0.125) -- (-1.491,-0.128) -- (-1.490,-0.131) -- (-1.490,-0.134) -- (-1.490,-0.137) -- (-1.489,-0.140) -- (-1.489,-0.143) -- (-1.488,-0.146) -- (-1.488,-0.149) -- (-1.487,-0.151) -- (-1.487,-0.154) -- (-1.486,-0.157) -- (-1.486,-0.160) -- (-1.485,-0.163) -- (-1.485,-0.166) -- (-1.484,-0.169) -- (-1.483,-0.172) -- (-1.483,-0.175) -- (-1.482,-0.178) -- (-1.482,-0.181) -- (-1.481,-0.183) -- (-1.480,-0.186) -- (-1.480,-0.189) -- (-1.479,-0.192) -- (-1.479,-0.195) -- (-1.478,-0.198) -- (-1.477,-0.201) -- (-1.477,-0.204) -- (-1.476,-0.207) -- (-1.475,-0.209) -- (-1.475,-0.212) -- (-1.474,-0.215) -- (-1.473,-0.218) -- (-1.472,-0.221) -- (-1.472,-0.224) -- (-1.471,-0.227) -- (-1.470,-0.229) -- (-1.469,-0.232) -- (-1.469,-0.235) -- (-1.468,-0.238) -- (-1.467,-0.241) -- (-1.466,-0.244) -- (-1.466,-0.247) -- (-1.465,-0.249) -- (-1.464,-0.252) -- (-1.463,-0.255) -- (-1.462,-0.258) -- (-1.461,-0.261) -- (-1.461,-0.264) -- (-1.460,-0.266) -- (-1.459,-0.269) -- (-1.458,-0.272) -- (-1.457,-0.275) -- (-1.456,-0.278) -- (-1.455,-0.280) -- (-1.454,-0.283) -- (-1.453,-0.286) -- (-1.452,-0.289) -- (-1.451,-0.292) -- (-1.450,-0.294) -- (-1.449,-0.297) -- (-1.448,-0.300) -- (-1.448,-0.303) -- (-1.447,-0.306) -- (-1.445,-0.308) -- (-1.444,-0.311) -- (-1.443,-0.314) -- (-1.442,-0.317) -- (-1.441,-0.319) -- (-1.440,-0.322) -- (-1.439,-0.325) -- (-1.438,-0.328) -- (-1.437,-0.330) -- (-1.436,-0.333) -- (-1.435,-0.336) -- (-1.434,-0.339) -- (-1.433,-0.341) -- (-1.432,-0.344) -- (-1.430,-0.347) -- (-1.429,-0.350) -- (-1.428,-0.352) -- (-1.427,-0.355) -- (-1.426,-0.358) -- (-1.425,-0.361) -- (-1.423,-0.363) -- (-1.422,-0.366) -- (-1.421,-0.369) -- (-1.420,-0.371) -- (-1.419,-0.374) -- (-1.417,-0.377) -- (-1.416,-0.379) -- (-1.415,-0.382) -- (-1.414,-0.385) -- (-1.412,-0.387) -- (-1.411,-0.390) -- (-1.410,-0.393) -- (-1.408,-0.395) -- (-1.407,-0.398) -- (-1.406,-0.401) -- (-1.405,-0.403) -- (-1.403,-0.406) -- (-1.402,-0.409) -- (-1.401,-0.411) -- (-1.399,-0.414) -- (-1.398,-0.416) -- (-1.396,-0.419) -- (-1.395,-0.422) -- (-1.394,-0.424) -- (-1.392,-0.427) -- (-1.391,-0.430) -- (-1.389,-0.432) -- (-1.388,-0.435) -- (-1.387,-0.437) -- (-1.385,-0.440) -- (-1.384,-0.442) -- (-1.382,-0.445) -- (-1.381,-0.448) -- (-1.379,-0.450) -- (-1.378,-0.453) -- (-1.376,-0.455) -- (-1.375,-0.458) -- (-1.373,-0.460) -- (-1.372,-0.463) -- (-1.370,-0.465) -- (-1.369,-0.468) -- (-1.367,-0.470) -- (-1.366,-0.473) -- (-1.364,-0.475) -- (-1.363,-0.478) -- (-1.361,-0.480) -- (-1.359,-0.483) -- (-1.358,-0.485) -- (-1.356,-0.488) -- (-1.355,-0.490) -- (-1.353,-0.493) -- (-1.351,-0.495) -- (-1.350,-0.498) -- (-1.348,-0.500) -- (-1.346,-0.503) -- (-1.345,-0.505) -- (-1.343,-0.508) -- (-1.341,-0.510) -- (-1.340,-0.513) -- (-1.338,-0.515) -- (-1.336,-0.517) -- (-1.335,-0.520) -- (-1.333,-0.522) -- (-1.331,-0.525) -- (-1.330,-0.527) -- (-1.328,-0.529) -- (-1.326,-0.532) -- (-1.324,-0.534) -- (-1.323,-0.537) -- (-1.321,-0.539) -- (-1.319,-0.541) -- (-1.317,-0.544) -- (-1.315,-0.546) -- (-1.314,-0.548) -- (-1.312,-0.551) -- (-1.310,-0.553) -- (-1.308,-0.555) -- (-1.306,-0.558) -- (-1.305,-0.560) -- (-1.303,-0.562) -- (-1.301,-0.565) -- (-1.299,-0.567) -- (-1.297,-0.569) -- (-1.295,-0.571) -- (-1.293,-0.574) -- (-1.291,-0.576) -- (-1.290,-0.578) -- (-1.288,-0.581) -- (-1.286,-0.583) -- (-1.284,-0.585) -- (-1.282,-0.587) -- (-1.280,-0.590) -- (-1.278,-0.592) -- (-1.276,-0.594) -- (-1.274,-0.596) -- (-1.272,-0.598) -- (-1.270,-0.601) -- (-1.268,-0.603) -- (-1.266,-0.605) -- (-1.264,-0.607) -- (-1.262,-0.609) -- (-1.260,-0.612) -- (-1.258,-0.614) -- (-1.256,-0.616) -- (-1.254,-0.618) -- (-1.252,-0.620) -- (-1.250,-0.622) -- (-1.248,-0.624) -- (-1.246,-0.627) -- (-1.244,-0.629) -- (-1.242,-0.631) -- (-1.240,-0.633) -- (-1.238,-0.635) -- (-1.236,-0.637) -- (-1.234,-0.639) -- (-1.231,-0.641) -- (-1.229,-0.643) -- (-1.227,-0.645) -- (-1.225,-0.647) -- (-1.223,-0.649) -- (-1.221,-0.652) -- (-1.219,-0.654) -- (-1.217,-0.656) -- (-1.214,-0.658) -- (-1.212,-0.660) -- (-1.210,-0.662) -- (-1.208,-0.664) -- (-1.206,-0.666) -- (-1.204,-0.668) -- (-1.201,-0.670) -- (-1.199,-0.672) -- (-1.197,-0.674) -- (-1.195,-0.675) -- (-1.192,-0.677) -- (-1.190,-0.679) -- (-1.188,-0.681) -- (-1.186,-0.683) -- (-1.184,-0.685) -- (-1.181,-0.687) -- (-1.179,-0.689) -- (-1.177,-0.691) -- (-1.174,-0.693) -- (-1.172,-0.695) -- (-1.170,-0.697) -- (-1.168,-0.698) -- (-1.165,-0.700) -- (-1.163,-0.702) -- (-1.161,-0.704) -- (-1.158,-0.706) -- (-1.156,-0.708) -- (-1.154,-0.709) -- (-1.151,-0.711) -- (-1.149,-0.713) -- (-1.147,-0.715) -- (-1.144,-0.717) -- (-1.142,-0.718) -- (-1.140,-0.720) -- (-1.137,-0.722) -- (-1.135,-0.724) -- (-1.133,-0.725) -- (-1.130,-0.727) -- (-1.128,-0.729) -- (-1.125,-0.731) -- (-1.123,-0.732) -- (-1.121,-0.734) -- (-1.118,-0.736) -- (-1.116,-0.738) -- (-1.113,-0.739) -- (-1.111,-0.741) -- (-1.108,-0.743) -- (-1.106,-0.744) -- (-1.104,-0.746) -- (-1.101,-0.748) -- (-1.099,-0.749) -- (-1.096,-0.751) -- (-1.094,-0.752) -- (-1.091,-0.754) -- (-1.089,-0.756) -- (-1.086,-0.757) -- (-1.084,-0.759) -- (-1.081,-0.760) -- (-1.079,-0.762) -- (-1.076,-0.764) -- (-1.074,-0.765) -- (-1.071,-0.767) -- (-1.069,-0.768) -- (-1.066,-0.770) -- (-1.064,-0.771) -- (-1.061,-0.773) -- (-1.059,-0.774) -- (-1.056,-0.776) -- (-1.054,-0.777) -- (-1.051,-0.779) -- (-1.048,-0.780) -- (-1.046,-0.782) -- (-1.043,-0.783) -- (-1.041,-0.785) -- (-1.038,-0.786) -- (-1.036,-0.788) -- (-1.033,-0.789) -- (-1.030,-0.790) -- (-1.028,-0.792) -- (-1.025,-0.793) -- (-1.023,-0.795) -- (-1.020,-0.796) -- (-1.017,-0.797) -- (-1.015,-0.799) -- (-1.012,-0.800) -- (-1.009,-0.801) -- (-1.007,-0.803) -- (-1.004,-0.804) -- (-1.002,-0.805) -- (-0.999,-0.807) -- (-0.996,-0.808) -- (-0.994,-0.809) -- (-0.991,-0.811) -- (-0.988,-0.812) -- (-0.986,-0.813) -- (-0.983,-0.814) -- (-0.980,-0.816) -- (-0.978,-0.817) -- (-0.975,-0.818) -- (-0.972,-0.819) -- (-0.969,-0.821) -- (-0.967,-0.822) -- (-0.964,-0.823) -- (-0.961,-0.824) -- (-0.959,-0.825) -- (-0.956,-0.827) -- (-0.953,-0.828) -- (-0.951,-0.829) -- (-0.948,-0.830) -- (-0.945,-0.831) -- (-0.942,-0.832) -- (-0.940,-0.833) -- (-0.937,-0.835) -- (-0.934,-0.836) -- (-0.931,-0.837) -- (-0.929,-0.838) -- (-0.926,-0.839) -- (-0.923,-0.840) -- (-0.920,-0.841) -- (-0.918,-0.842) -- (-0.915,-0.843) -- (-0.912,-0.844) -- (-0.909,-0.845) -- (-0.907,-0.846) -- (-0.904,-0.847) -- (-0.901,-0.848) -- (-0.898,-0.849) -- (-0.895,-0.850) -- (-0.893,-0.851) -- (-0.890,-0.852) -- (-0.887,-0.853) -- (-0.884,-0.854) -- (-0.881,-0.855) -- (-0.879,-0.856) -- (-0.876,-0.857) -- (-0.873,-0.858) -- (-0.870,-0.858) -- (-0.867,-0.859) -- (-0.864,-0.860) -- (-0.862,-0.861) -- (-0.859,-0.862) -- (-0.856,-0.863) -- (-0.853,-0.864) -- (-0.850,-0.864) -- (-0.847,-0.865) -- (-0.845,-0.866) -- (-0.842,-0.867) -- (-0.839,-0.868) -- (-0.836,-0.868) -- (-0.833,-0.869) -- (-0.830,-0.870) -- (-0.828,-0.871) -- (-0.825,-0.871) -- (-0.822,-0.872) -- (-0.819,-0.873) -- (-0.816,-0.874) -- (-0.813,-0.874) -- (-0.810,-0.875) -- (-0.807,-0.876) -- (-0.805,-0.876) -- (-0.802,-0.877) -- (-0.799,-0.878) -- (-0.796,-0.878) -- (-0.793,-0.879) -- (-0.790,-0.880) -- (-0.787,-0.880) -- (-0.784,-0.881) -- (-0.781,-0.881) -- (-0.779,-0.882) -- (-0.776,-0.883) -- (-0.773,-0.883) -- (-0.770,-0.884) -- (-0.767,-0.884) -- (-0.764,-0.885) -- (-0.761,-0.885) -- (-0.758,-0.886) -- (-0.755,-0.886) -- (-0.752,-0.887) -- (-0.750,-0.887) -- (-0.747,-0.888) -- (-0.744,-0.888) -- (-0.741,-0.889) -- (-0.738,-0.889) -- (-0.735,-0.890) -- (-0.732,-0.890) -- (-0.729,-0.891) -- (-0.726,-0.891) -- (-0.723,-0.892) -- (-0.720,-0.892) -- (-0.717,-0.892) -- (-0.714,-0.893) -- (-0.711,-0.893) -- (-0.709,-0.893) -- (-0.706,-0.894) -- (-0.703,-0.894) -- (-0.700,-0.894) -- (-0.697,-0.895) -- (-0.694,-0.895) -- (-0.691,-0.895) -- (-0.688,-0.896) -- (-0.685,-0.896) -- (-0.682,-0.896) -- (-0.679,-0.896) -- (-0.676,-0.897) -- (-0.673,-0.897) -- (-0.670,-0.897) -- (-0.667,-0.897) -- (-0.664,-0.898) -- (-0.661,-0.898) -- (-0.659,-0.898) -- (-0.656,-0.898) -- (-0.653,-0.898) -- (-0.650,-0.899) -- (-0.647,-0.899) -- (-0.644,-0.899) -- (-0.641,-0.899) -- (-0.638,-0.899) -- (-0.635,-0.899) -- (-0.632,-0.899) -- (-0.629,-0.900) -- (-0.626,-0.900) -- (-0.623,-0.900) -- (-0.620,-0.900) -- (-0.617,-0.900) -- (-0.614,-0.900) -- (-0.611,-0.900) -- (-0.608,-0.900) -- (-0.605,-0.900) -- (-0.602,-0.900) -- (-0.599,-0.900) -- (-0.597,-0.900) -- (-0.594,-0.900) -- (-0.591,-0.900) -- (-0.588,-0.900) -- (-0.585,-0.900) -- (-0.582,-0.900) -- (-0.579,-0.900) -- (-0.576,-0.900) -- (-0.573,-0.900) -- (-0.570,-0.899) -- (-0.567,-0.899) -- (-0.564,-0.899) -- (-0.561,-0.899) -- (-0.558,-0.899) -- (-0.555,-0.899) -- (-0.552,-0.899) -- (-0.549,-0.899) -- (-0.546,-0.898) -- (-0.543,-0.898) -- (-0.540,-0.898) -- (-0.537,-0.898) -- (-0.535,-0.898) -- (-0.532,-0.897) -- (-0.529,-0.897) -- (-0.526,-0.897) -- (-0.523,-0.897) -- (-0.520,-0.896) -- (-0.517,-0.896) -- (-0.514,-0.896) -- (-0.511,-0.896) -- (-0.508,-0.895) -- (-0.505,-0.895) -- (-0.502,-0.895) -- (-0.499,-0.894) -- (-0.496,-0.894) -- (-0.493,-0.894) -- (-0.490,-0.893) -- (-0.487,-0.893) -- (-0.485,-0.893) -- (-0.482,-0.892) -- (-0.479,-0.892) -- (-0.476,-0.891) -- (-0.473,-0.891) -- (-0.470,-0.891) -- (-0.467,-0.890) -- (-0.464,-0.890) -- (-0.461,-0.889) -- (-0.458,-0.889) -- (-0.455,-0.888) -- (-0.452,-0.888) -- (-0.449,-0.887) -- (-0.447,-0.887) -- (-0.444,-0.886) -- (-0.441,-0.886) -- (-0.438,-0.885) -- (-0.435,-0.885) -- (-0.432,-0.884) -- (-0.429,-0.884) -- (-0.426,-0.883) -- (-0.423,-0.882) -- (-0.420,-0.882) -- (-0.418,-0.881) -- (-0.415,-0.881) -- (-0.412,-0.880) -- (-0.409,-0.879) -- (-0.406,-0.879) -- (-0.403,-0.878) -- (-0.400,-0.878) -- (-0.397,-0.877) -- (-0.394,-0.876) -- (-0.392,-0.876) -- (-0.389,-0.875) -- (-0.386,-0.874) -- (-0.383,-0.873) -- (-0.380,-0.873) -- (-0.377,-0.872) -- (-0.374,-0.871) -- (-0.371,-0.870) -- (-0.369,-0.870) -- (-0.366,-0.869) -- (-0.363,-0.868) -- (-0.360,-0.867) -- (-0.357,-0.867) -- (-0.354,-0.866) -- (-0.352,-0.865) -- (-0.349,-0.864) -- (-0.346,-0.863) -- (-0.343,-0.863) -- (-0.340,-0.862) -- (-0.337,-0.861) -- (-0.335,-0.860) -- (-0.332,-0.859) -- (-0.329,-0.858) -- (-0.326,-0.857) -- (-0.323,-0.856) -- (-0.320,-0.855) -- (-0.318,-0.855) -- (-0.315,-0.854) -- (-0.312,-0.853) -- (-0.309,-0.852) -- (-0.306,-0.851) -- (-0.304,-0.850) -- (-0.301,-0.849) -- (-0.298,-0.848) -- (-0.295,-0.847) -- (-0.293,-0.846) -- (-0.290,-0.845) -- (-0.287,-0.844) -- (-0.284,-0.843) -- (-0.281,-0.842) -- (-0.279,-0.841) -- (-0.276,-0.840) -- (-0.273,-0.839) -- (-0.270,-0.837) -- (-0.268,-0.836) -- (-0.265,-0.835) -- (-0.262,-0.834) -- (-0.259,-0.833) -- (-0.256,-0.832) -- (-0.251,-0.830) -- (-0.246,-0.827) -- (-0.239,-0.825) -- (-0.233,-0.822) -- (-0.227,-0.820) -- (-0.220,-0.817) -- (-0.214,-0.814) -- (-0.208,-0.812) -- (-0.201,-0.809) -- (-0.195,-0.806) -- (-0.189,-0.804) -- (-0.182,-0.801) -- (-0.176,-0.798) -- (-0.170,-0.796) -- (-0.163,-0.793) -- (-0.157,-0.791) -- (-0.151,-0.788) -- (-0.144,-0.785) -- (-0.138,-0.783) -- (-0.132,-0.780) -- (-0.125,-0.777) -- (-0.119,-0.775) -- (-0.113,-0.772) -- (-0.106,-0.770) -- (-0.100,-0.767) -- (-0.093,-0.764) -- (-0.087,-0.762) -- (-0.081,-0.759) -- (-0.074,-0.756) -- (-0.068,-0.754) -- (-0.062,-0.751) -- (-0.055,-0.749) -- (-0.049,-0.746) -- (-0.043,-0.743) -- (-0.036,-0.741) -- (-0.030,-0.738) -- (-0.024,-0.735) -- (-0.017,-0.733) -- (-0.011,-0.730) -- (-0.005,-0.727) -- (0.002,-0.725) -- (0.008,-0.722) -- (0.014,-0.720) -- (0.021,-0.717) -- (0.027,-0.714) -- (0.033,-0.712) -- (0.040,-0.709) -- (0.046,-0.706) -- (0.052,-0.704) -- (0.059,-0.701) -- (0.065,-0.699) -- (0.071,-0.696) -- (0.078,-0.693) -- (0.084,-0.691) -- (0.091,-0.688) -- (0.097,-0.685) -- (0.103,-0.683) -- (0.110,-0.680) -- (0.116,-0.677) -- (0.122,-0.675) -- (0.129,-0.672) -- (0.135,-0.670) -- (0.141,-0.667) -- (0.148,-0.664) -- (0.154,-0.662) -- (0.160,-0.659) -- (0.167,-0.656) -- (0.173,-0.654) -- (0.179,-0.651) -- (0.186,-0.649) -- (0.192,-0.646) -- (0.198,-0.643) -- (0.205,-0.641) -- (0.211,-0.638) -- (0.217,-0.635) -- (0.224,-0.633) -- (0.230,-0.630) -- (0.236,-0.627) -- (0.243,-0.625) -- (0.249,-0.622) -- (0.256,-0.620) -- (0.262,-0.617) -- (0.268,-0.614) -- (0.275,-0.612) -- (0.281,-0.609) -- (0.287,-0.606) -- (0.294,-0.604) -- (0.300,-0.601) -- (0.306,-0.599) -- (0.313,-0.596) -- (0.319,-0.593) -- (0.325,-0.591) -- (0.332,-0.588) -- (0.338,-0.585) -- (0.344,-0.583) -- (0.351,-0.580) -- (0.357,-0.577) -- (0.363,-0.575) -- (0.370,-0.572) -- (0.376,-0.570) -- (0.382,-0.567) -- (0.389,-0.564) -- (0.395,-0.562) -- (0.401,-0.559) -- (0.408,-0.556) -- (0.414,-0.554) -- (0.421,-0.551) -- (0.427,-0.549) -- (0.433,-0.546) -- (0.440,-0.543) -- (0.446,-0.541) -- (0.452,-0.538) -- (0.459,-0.535) -- (0.465,-0.533) -- (0.471,-0.530) -- (0.478,-0.528) -- (0.484,-0.525) -- (0.490,-0.522) -- (0.497,-0.520) -- (0.503,-0.517) -- (0.509,-0.514) -- (0.516,-0.512) -- (0.522,-0.509) -- (0.528,-0.506) -- (0.535,-0.504) -- (0.541,-0.501) -- (0.547,-0.499) -- (0.554,-0.496) -- (0.560,-0.493) -- (0.566,-0.491) -- (0.573,-0.488) -- (0.579,-0.485) -- (0.586,-0.483) -- (0.592,-0.480) -- (0.598,-0.478) -- (0.605,-0.475) -- (0.611,-0.472) -- (0.617,-0.470) -- (0.624,-0.467) -- (0.630,-0.464) -- (0.636,-0.462) -- (0.643,-0.459) -- (0.649,-0.456) -- (0.655,-0.454) -- (0.662,-0.451) -- (0.668,-0.449) -- (0.674,-0.446) -- (0.681,-0.443) -- (0.687,-0.441) -- (0.693,-0.438) -- (0.700,-0.435) -- (0.706,-0.433) -- (0.712,-0.430) -- (0.719,-0.428) -- (0.725,-0.425) -- (0.731,-0.422) -- (0.738,-0.420) -- (0.744,-0.417) -- (0.751,-0.414) -- (0.757,-0.412) -- (0.763,-0.409) -- (0.770,-0.406) -- (0.776,-0.404) -- (0.782,-0.401) -- (0.789,-0.399) -- (0.795,-0.396) -- (0.801,-0.393) -- (0.808,-0.391) -- (0.814,-0.388) -- (0.820,-0.385) -- (0.827,-0.383) -- (0.833,-0.380) -- (0.839,-0.378) -- (0.846,-0.375) -- (0.852,-0.372) -- (0.858,-0.370) -- (0.865,-0.367) -- (0.871,-0.364) -- (0.877,-0.362) -- (0.884,-0.359) -- (0.890,-0.356) -- (0.896,-0.354) -- (0.903,-0.351) -- (0.909,-0.349) -- (0.916,-0.346) -- (0.922,-0.343) -- (0.928,-0.341) -- (0.935,-0.338) -- (0.941,-0.335) -- (0.947,-0.333) -- (0.954,-0.330) -- (0.960,-0.328) -- (0.966,-0.325) -- (0.973,-0.322) -- (0.979,-0.320) -- (0.985,-0.317) -- (0.992,-0.314) -- (0.998,-0.312) -- (1.004,-0.309) -- (1.011,-0.307) -- (1.017,-0.304) -- (1.023,-0.301) -- (1.030,-0.299) -- (1.036,-0.296) -- (1.042,-0.293) -- (1.049,-0.291) -- (1.055,-0.288) -- (1.061,-0.285) -- (1.068,-0.283) -- (1.074,-0.280) -- (1.081,-0.278) -- (1.087,-0.275) -- (1.093,-0.272) -- (1.100,-0.270) -- (1.106,-0.267) -- (1.112,-0.264) -- (1.119,-0.262) -- (1.125,-0.259) -- (1.131,-0.257) -- (1.138,-0.254) -- (1.144,-0.251) -- (1.150,-0.249) -- (1.157,-0.246) -- (1.163,-0.243) -- (1.169,-0.241) -- (1.176,-0.238) -- (1.182,-0.235) -- (1.188,-0.233) -- (1.195,-0.230) -- (1.201,-0.228) -- (1.207,-0.225) -- (1.214,-0.222) -- (1.220,-0.220) -- (1.226,-0.217) -- (1.233,-0.214) -- (1.239,-0.212) -- (1.245,-0.209) -- (1.252,-0.207) -- (1.258,-0.204) -- (1.265,-0.201) -- (1.271,-0.199) -- (1.277,-0.196) -- (1.284,-0.193) -- (1.290,-0.191) -- (1.296,-0.188) -- (1.303,-0.185) -- (1.309,-0.183) -- (1.315,-0.180) -- (1.322,-0.178) -- (1.328,-0.175) -- (1.334,-0.172) -- (1.341,-0.170) -- (1.347,-0.167) -- (1.353,-0.164) -- (1.360,-0.162) -- (1.366,-0.159) -- (1.372,-0.157) -- (1.379,-0.154) -- (1.385,-0.151) -- (1.391,-0.149) -- (1.398,-0.146) -- (1.404,-0.143) -- (1.410,-0.141) -- (1.417,-0.138) -- (1.423,-0.135) -- (1.430,-0.133) -- (1.436,-0.130) -- (1.442,-0.128) -- (1.449,-0.125) -- (1.455,-0.122) -- (1.461,-0.120) -- (1.468,-0.117) -- (1.474,-0.114) -- (1.480,-0.112) -- (1.487,-0.109) -- (1.493,-0.107) -- (1.499,-0.104) -- (1.506,-0.101) -- (1.512,-0.099) -- (1.518,-0.096) -- (1.525,-0.093) -- (1.531,-0.091) -- (1.537,-0.088) -- (1.544,-0.086) -- (1.550,-0.083) -- (1.556,-0.080) -- (1.563,-0.078) -- (1.569,-0.075) -- (1.575,-0.072) -- (1.582,-0.070) -- (1.588,-0.067) -- (1.595,-0.064) -- (1.601,-0.062) -- (1.607,-0.059) -- (1.614,-0.057) -- (1.620,-0.054) -- (1.626,-0.051) -- (1.633,-0.049) -- (1.639,-0.046) -- (1.645,-0.043) -- (1.652,-0.041) -- (1.658,-0.038) -- (1.664,-0.036) -- (1.671,-0.033) -- (1.677,-0.030) -- (1.683,-0.028) -- (1.690,-0.025) -- (1.696,-0.022) -- (1.702,-0.020) -- (1.709,-0.017) -- (1.715,-0.014) -- (1.721,-0.012) -- (1.728,-0.009) -- (1.734,-0.007) -- (1.740,-0.004) -- (1.744,-0.001) -- (1.744,0.001) -- (1.740,0.004) -- (1.734,0.007) -- (1.728,0.009) -- (1.721,0.012) -- (1.715,0.014) -- (1.709,0.017) -- (1.702,0.020) -- (1.696,0.022) -- (1.690,0.025) -- (1.683,0.028) -- (1.677,0.030) -- (1.671,0.033) -- (1.664,0.036) -- (1.658,0.038) -- (1.652,0.041) -- (1.645,0.043) -- (1.639,0.046) -- (1.633,0.049) -- (1.626,0.051) -- (1.620,0.054) -- (1.614,0.057) -- (1.607,0.059) -- (1.601,0.062) -- (1.595,0.064) -- (1.588,0.067) -- (1.582,0.070) -- (1.575,0.072) -- (1.569,0.075) -- (1.563,0.078) -- (1.556,0.080) -- (1.550,0.083) -- (1.544,0.086) -- (1.537,0.088) -- (1.531,0.091) -- (1.525,0.093) -- (1.518,0.096) -- (1.512,0.099) -- (1.506,0.101) -- (1.499,0.104) -- (1.493,0.107) -- (1.487,0.109) -- (1.480,0.112) -- (1.474,0.114) -- (1.468,0.117) -- (1.461,0.120) -- (1.455,0.122) -- (1.449,0.125) -- (1.442,0.128) -- (1.436,0.130) -- (1.430,0.133) -- (1.423,0.135) -- (1.417,0.138) -- (1.410,0.141) -- (1.404,0.143) -- (1.398,0.146) -- (1.391,0.149) -- (1.385,0.151) -- (1.379,0.154) -- (1.372,0.157) -- (1.366,0.159) -- (1.360,0.162) -- (1.353,0.164) -- (1.347,0.167) -- (1.341,0.170) -- (1.334,0.172) -- (1.328,0.175) -- (1.322,0.178) -- (1.315,0.180) -- (1.309,0.183) -- (1.303,0.185) -- (1.296,0.188) -- (1.290,0.191) -- (1.284,0.193) -- (1.277,0.196) -- (1.271,0.199) -- (1.265,0.201) -- (1.258,0.204) -- (1.252,0.207) -- (1.245,0.209) -- (1.239,0.212) -- (1.233,0.214) -- (1.226,0.217) -- (1.220,0.220) -- (1.214,0.222) -- (1.207,0.225) -- (1.201,0.228) -- (1.195,0.230) -- (1.188,0.233) -- (1.182,0.235) -- (1.176,0.238) -- (1.169,0.241) -- (1.163,0.243) -- (1.157,0.246) -- (1.150,0.249) -- (1.144,0.251) -- (1.138,0.254) -- (1.131,0.257) -- (1.125,0.259) -- (1.119,0.262) -- (1.112,0.264) -- (1.106,0.267) -- (1.100,0.270) -- (1.093,0.272) -- (1.087,0.275) -- (1.081,0.278) -- (1.074,0.280) -- (1.068,0.283) -- (1.061,0.285) -- (1.055,0.288) -- (1.049,0.291) -- (1.042,0.293) -- (1.036,0.296) -- (1.030,0.299) -- (1.023,0.301) -- (1.017,0.304) -- (1.011,0.307) -- (1.004,0.309) -- (0.998,0.312) -- (0.992,0.314) -- (0.985,0.317) -- (0.979,0.320) -- (0.973,0.322) -- (0.966,0.325) -- (0.960,0.328) -- (0.954,0.330) -- (0.947,0.333) -- (0.941,0.335) -- (0.935,0.338) -- (0.928,0.341) -- (0.922,0.343) -- (0.916,0.346) -- (0.909,0.349) -- (0.903,0.351) -- (0.896,0.354) -- (0.890,0.356) -- (0.884,0.359) -- (0.877,0.362) -- (0.871,0.364) -- (0.865,0.367) -- (0.858,0.370) -- (0.852,0.372) -- (0.846,0.375) -- (0.839,0.378) -- (0.833,0.380) -- (0.827,0.383) -- (0.820,0.385) -- (0.814,0.388) -- (0.808,0.391) -- (0.801,0.393) -- (0.795,0.396) -- (0.789,0.399) -- (0.782,0.401) -- (0.776,0.404) -- (0.770,0.406) -- (0.763,0.409) -- (0.757,0.412) -- (0.751,0.414) -- (0.744,0.417) -- (0.738,0.420) -- (0.731,0.422) -- (0.725,0.425) -- (0.719,0.428) -- (0.712,0.430) -- (0.706,0.433) -- (0.700,0.435) -- (0.693,0.438) -- (0.687,0.441) -- (0.681,0.443) -- (0.674,0.446) -- (0.668,0.449) -- (0.662,0.451) -- (0.655,0.454) -- (0.649,0.456) -- (0.643,0.459) -- (0.636,0.462) -- (0.630,0.464) -- (0.624,0.467) -- (0.617,0.470) -- (0.611,0.472) -- (0.605,0.475) -- (0.598,0.478) -- (0.592,0.480) -- (0.586,0.483) -- (0.579,0.485) -- (0.573,0.488) -- (0.566,0.491) -- (0.560,0.493) -- (0.554,0.496) -- (0.547,0.499) -- (0.541,0.501) -- (0.535,0.504) -- (0.528,0.506) -- (0.522,0.509) -- (0.516,0.512) -- (0.509,0.514) -- (0.503,0.517) -- (0.497,0.520) -- (0.490,0.522) -- (0.484,0.525) -- (0.478,0.528) -- (0.471,0.530) -- (0.465,0.533) -- (0.459,0.535) -- (0.452,0.538) -- (0.446,0.541) -- (0.440,0.543) -- (0.433,0.546) -- (0.427,0.549) -- (0.421,0.551) -- (0.414,0.554) -- (0.408,0.556) -- (0.401,0.559) -- (0.395,0.562) -- (0.389,0.564) -- (0.382,0.567) -- (0.376,0.570) -- (0.370,0.572) -- (0.363,0.575) -- (0.357,0.577) -- (0.351,0.580) -- (0.344,0.583) -- (0.338,0.585) -- (0.332,0.588) -- (0.325,0.591) -- (0.319,0.593) -- (0.313,0.596) -- (0.306,0.599) -- (0.300,0.601) -- (0.294,0.604) -- (0.287,0.606) -- (0.281,0.609) -- (0.275,0.612) -- (0.268,0.614) -- (0.262,0.617) -- (0.256,0.620) -- (0.249,0.622) -- (0.243,0.625) -- (0.236,0.627) -- (0.230,0.630) -- (0.224,0.633) -- (0.217,0.635) -- (0.211,0.638) -- (0.205,0.641) -- (0.198,0.643) -- (0.192,0.646) -- (0.186,0.649) -- (0.179,0.651) -- (0.173,0.654) -- (0.167,0.656) -- (0.160,0.659) -- (0.154,0.662) -- (0.148,0.664) -- (0.141,0.667) -- (0.135,0.670) -- (0.129,0.672) -- (0.122,0.675) -- (0.116,0.677) -- (0.110,0.680) -- (0.103,0.683) -- (0.097,0.685) -- (0.091,0.688) -- (0.084,0.691) -- (0.078,0.693) -- (0.071,0.696) -- (0.065,0.699) -- (0.059,0.701) -- (0.052,0.704) -- (0.046,0.706) -- (0.040,0.709) -- (0.033,0.712) -- (0.027,0.714) -- (0.021,0.717) -- (0.014,0.720) -- (0.008,0.722) -- (0.002,0.725) -- (-0.005,0.727) -- (-0.011,0.730) -- (-0.017,0.733) -- (-0.024,0.735) -- (-0.030,0.738) -- (-0.036,0.741) -- (-0.043,0.743) -- (-0.049,0.746) -- (-0.055,0.749) -- (-0.062,0.751) -- (-0.068,0.754) -- (-0.074,0.756) -- (-0.081,0.759) -- (-0.087,0.762) -- (-0.093,0.764) -- (-0.100,0.767) -- (-0.106,0.770) -- (-0.113,0.772) -- (-0.119,0.775) -- (-0.125,0.777) -- (-0.132,0.780) -- (-0.138,0.783) -- (-0.144,0.785) -- (-0.151,0.788) -- (-0.157,0.791) -- (-0.163,0.793) -- (-0.170,0.796) -- (-0.176,0.798) -- (-0.182,0.801) -- (-0.189,0.804) -- (-0.195,0.806) -- (-0.201,0.809) -- (-0.208,0.812) -- (-0.214,0.814) -- (-0.220,0.817) -- (-0.227,0.820) -- (-0.233,0.822) -- (-0.239,0.825) -- (-0.246,0.827) -- (-0.251,0.830) -- (-0.256,0.832) -- cycle}

    \fill[top color=azulvelo!80!white, bottom color=metal] \formagota;
    % Conductor con punta - densidad superficial
% Generado por scripts/generar_datos.py (no editar a mano).
\SG{-0.259}{0.833}{-0.262}{0.834}{0.183}
\SG{-0.262}{0.834}{-0.265}{0.835}{0.181}
\SG{-0.265}{0.835}{-0.268}{0.836}{0.179}
\SG{-0.268}{0.836}{-0.270}{0.837}{0.179}
\SG{-0.270}{0.837}{-0.273}{0.839}{0.181}
\SG{-0.273}{0.839}{-0.276}{0.840}{0.183}
\SG{-0.276}{0.840}{-0.279}{0.841}{0.185}
\SG{-0.279}{0.841}{-0.281}{0.842}{0.187}
\SG{-0.281}{0.842}{-0.284}{0.843}{0.190}
\SG{-0.284}{0.843}{-0.287}{0.844}{0.194}
\SG{-0.287}{0.844}{-0.290}{0.845}{0.197}
\SG{-0.290}{0.845}{-0.293}{0.846}{0.201}
\SG{-0.293}{0.846}{-0.295}{0.847}{0.203}
\SG{-0.295}{0.847}{-0.298}{0.848}{0.206}
\SG{-0.298}{0.848}{-0.301}{0.849}{0.203}
\SG{-0.301}{0.849}{-0.304}{0.850}{0.201}
\SG{-0.304}{0.850}{-0.306}{0.851}{0.200}
\SG{-0.306}{0.851}{-0.309}{0.852}{0.199}
\SG{-0.309}{0.852}{-0.312}{0.853}{0.199}
\SG{-0.312}{0.853}{-0.315}{0.854}{0.201}
\SG{-0.315}{0.854}{-0.318}{0.855}{0.204}
\SG{-0.318}{0.855}{-0.320}{0.855}{0.206}
\SG{-0.320}{0.855}{-0.323}{0.856}{0.209}
\SG{-0.323}{0.856}{-0.326}{0.857}{0.212}
\SG{-0.326}{0.857}{-0.329}{0.858}{0.216}
\SG{-0.329}{0.858}{-0.332}{0.859}{0.221}
\SG{-0.332}{0.859}{-0.335}{0.860}{0.225}
\SG{-0.335}{0.860}{-0.337}{0.861}{0.229}
\SG{-0.337}{0.861}{-0.340}{0.862}{0.233}
\SG{-0.340}{0.862}{-0.343}{0.863}{0.234}
\SG{-0.343}{0.863}{-0.346}{0.863}{0.232}
\SG{-0.346}{0.863}{-0.349}{0.864}{0.231}
\SG{-0.349}{0.864}{-0.352}{0.865}{0.231}
\SG{-0.352}{0.865}{-0.354}{0.866}{0.231}
\SG{-0.354}{0.866}{-0.357}{0.867}{0.232}
\SG{-0.357}{0.867}{-0.360}{0.867}{0.236}
\SG{-0.360}{0.867}{-0.363}{0.868}{0.239}
\SG{-0.363}{0.868}{-0.366}{0.869}{0.242}
\SG{-0.366}{0.869}{-0.369}{0.870}{0.246}
\SG{-0.369}{0.870}{-0.371}{0.870}{0.250}
\SG{-0.371}{0.870}{-0.374}{0.871}{0.257}
\SG{-0.374}{0.871}{-0.377}{0.872}{0.262}
\SG{-0.377}{0.872}{-0.380}{0.873}{0.266}
\SG{-0.380}{0.873}{-0.383}{0.873}{0.271}
\SG{-0.383}{0.873}{-0.386}{0.874}{0.273}
\SG{-0.386}{0.874}{-0.389}{0.875}{0.269}
\SG{-0.389}{0.875}{-0.392}{0.876}{0.266}
\SG{-0.392}{0.876}{-0.394}{0.876}{0.262}
\SG{-0.394}{0.876}{-0.397}{0.877}{0.259}
\SG{-0.397}{0.877}{-0.400}{0.878}{0.257}
\SG{-0.400}{0.878}{-0.403}{0.878}{0.253}
\SG{-0.403}{0.878}{-0.406}{0.879}{0.251}
\SG{-0.406}{0.879}{-0.409}{0.879}{0.248}
\SG{-0.409}{0.879}{-0.412}{0.880}{0.246}
\SG{-0.412}{0.880}{-0.415}{0.881}{0.244}
\SG{-0.415}{0.881}{-0.418}{0.881}{0.244}
\SG{-0.418}{0.881}{-0.420}{0.882}{0.245}
\SG{-0.420}{0.882}{-0.423}{0.882}{0.245}
\SG{-0.423}{0.882}{-0.426}{0.883}{0.246}
\SG{-0.426}{0.883}{-0.429}{0.884}{0.247}
\SG{-0.429}{0.884}{-0.432}{0.884}{0.252}
\SG{-0.432}{0.884}{-0.435}{0.885}{0.257}
\SG{-0.435}{0.885}{-0.438}{0.885}{0.262}
\SG{-0.438}{0.885}{-0.441}{0.886}{0.266}
\SG{-0.441}{0.886}{-0.444}{0.886}{0.271}
\SG{-0.444}{0.886}{-0.447}{0.887}{0.279}
\SG{-0.447}{0.887}{-0.449}{0.887}{0.286}
\SG{-0.449}{0.887}{-0.452}{0.888}{0.292}
\SG{-0.452}{0.888}{-0.455}{0.888}{0.299}
\SG{-0.455}{0.888}{-0.458}{0.889}{0.303}
\SG{-0.458}{0.889}{-0.461}{0.889}{0.302}
\SG{-0.461}{0.889}{-0.464}{0.890}{0.301}
\SG{-0.464}{0.890}{-0.467}{0.890}{0.299}
\SG{-0.467}{0.890}{-0.470}{0.891}{0.298}
\SG{-0.470}{0.891}{-0.473}{0.891}{0.297}
\SG{-0.473}{0.891}{-0.476}{0.891}{0.295}
\SG{-0.476}{0.891}{-0.479}{0.892}{0.293}
\SG{-0.479}{0.892}{-0.482}{0.892}{0.290}
\SG{-0.482}{0.892}{-0.485}{0.893}{0.289}
\SG{-0.485}{0.893}{-0.487}{0.893}{0.287}
\SG{-0.487}{0.893}{-0.490}{0.893}{0.285}
\SG{-0.490}{0.893}{-0.493}{0.894}{0.283}
\SG{-0.493}{0.894}{-0.496}{0.894}{0.281}
\SG{-0.496}{0.894}{-0.499}{0.894}{0.280}
\SG{-0.499}{0.894}{-0.502}{0.895}{0.278}
\SG{-0.502}{0.895}{-0.505}{0.895}{0.277}
\SG{-0.505}{0.895}{-0.508}{0.895}{0.276}
\SG{-0.508}{0.895}{-0.511}{0.896}{0.275}
\SG{-0.511}{0.896}{-0.514}{0.896}{0.273}
\SG{-0.514}{0.896}{-0.517}{0.896}{0.272}
\SG{-0.517}{0.896}{-0.520}{0.896}{0.271}
\SG{-0.520}{0.896}{-0.523}{0.897}{0.271}
\SG{-0.523}{0.897}{-0.526}{0.897}{0.270}
\SG{-0.526}{0.897}{-0.529}{0.897}{0.269}
\SG{-0.529}{0.897}{-0.532}{0.897}{0.268}
\SG{-0.532}{0.897}{-0.535}{0.898}{0.268}
\SG{-0.535}{0.898}{-0.537}{0.898}{0.267}
\SG{-0.537}{0.898}{-0.540}{0.898}{0.267}
\SG{-0.540}{0.898}{-0.543}{0.898}{0.266}
\SG{-0.543}{0.898}{-0.546}{0.898}{0.266}
\SG{-0.546}{0.898}{-0.549}{0.899}{0.266}
\SG{-0.549}{0.899}{-0.552}{0.899}{0.266}
\SG{-0.552}{0.899}{-0.555}{0.899}{0.265}
\SG{-0.555}{0.899}{-0.558}{0.899}{0.265}
\SG{-0.558}{0.899}{-0.561}{0.899}{0.265}
\SG{-0.561}{0.899}{-0.564}{0.899}{0.265}
\SG{-0.564}{0.899}{-0.567}{0.899}{0.265}
\SG{-0.567}{0.899}{-0.570}{0.899}{0.265}
\SG{-0.570}{0.899}{-0.573}{0.900}{0.265}
\SG{-0.573}{0.900}{-0.576}{0.900}{0.265}
\SG{-0.576}{0.900}{-0.579}{0.900}{0.266}
\SG{-0.579}{0.900}{-0.582}{0.900}{0.266}
\SG{-0.582}{0.900}{-0.585}{0.900}{0.266}
\SG{-0.585}{0.900}{-0.588}{0.900}{0.267}
\SG{-0.588}{0.900}{-0.591}{0.900}{0.267}
\SG{-0.591}{0.900}{-0.594}{0.900}{0.268}
\SG{-0.594}{0.900}{-0.597}{0.900}{0.268}
\SG{-0.597}{0.900}{-0.599}{0.900}{0.269}
\SG{-0.599}{0.900}{-0.602}{0.900}{0.269}
\SG{-0.602}{0.900}{-0.605}{0.900}{0.270}
\SG{-0.605}{0.900}{-0.608}{0.900}{0.270}
\SG{-0.608}{0.900}{-0.611}{0.900}{0.271}
\SG{-0.611}{0.900}{-0.614}{0.900}{0.272}
\SG{-0.614}{0.900}{-0.617}{0.900}{0.273}
\SG{-0.617}{0.900}{-0.620}{0.900}{0.273}
\SG{-0.620}{0.900}{-0.623}{0.900}{0.274}
\SG{-0.623}{0.900}{-0.626}{0.900}{0.275}
\SG{-0.626}{0.900}{-0.629}{0.900}{0.276}
\SG{-0.629}{0.900}{-0.632}{0.899}{0.277}
\SG{-0.632}{0.899}{-0.635}{0.899}{0.278}
\SG{-0.635}{0.899}{-0.638}{0.899}{0.280}
\SG{-0.638}{0.899}{-0.641}{0.899}{0.281}
\SG{-0.641}{0.899}{-0.644}{0.899}{0.282}
\SG{-0.644}{0.899}{-0.647}{0.899}{0.283}
\SG{-0.647}{0.899}{-0.650}{0.899}{0.284}
\SG{-0.650}{0.899}{-0.653}{0.898}{0.286}
\SG{-0.653}{0.898}{-0.656}{0.898}{0.287}
\SG{-0.656}{0.898}{-0.659}{0.898}{0.289}
\SG{-0.659}{0.898}{-0.661}{0.898}{0.290}
\SG{-0.661}{0.898}{-0.664}{0.898}{0.292}
\SG{-0.664}{0.898}{-0.667}{0.897}{0.294}
\SG{-0.667}{0.897}{-0.670}{0.897}{0.295}
\SG{-0.670}{0.897}{-0.673}{0.897}{0.297}
\SG{-0.673}{0.897}{-0.676}{0.897}{0.299}
\SG{-0.676}{0.897}{-0.679}{0.896}{0.301}
\SG{-0.679}{0.896}{-0.682}{0.896}{0.303}
\SG{-0.682}{0.896}{-0.685}{0.896}{0.305}
\SG{-0.685}{0.896}{-0.688}{0.896}{0.307}
\SG{-0.688}{0.896}{-0.691}{0.895}{0.310}
\SG{-0.691}{0.895}{-0.694}{0.895}{0.312}
\SG{-0.694}{0.895}{-0.697}{0.895}{0.315}
\SG{-0.697}{0.895}{-0.700}{0.894}{0.317}
\SG{-0.700}{0.894}{-0.703}{0.894}{0.320}
\SG{-0.703}{0.894}{-0.706}{0.894}{0.322}
\SG{-0.706}{0.894}{-0.709}{0.893}{0.325}
\SG{-0.709}{0.893}{-0.711}{0.893}{0.328}
\SG{-0.711}{0.893}{-0.714}{0.893}{0.331}
\SG{-0.714}{0.893}{-0.717}{0.892}{0.334}
\SG{-0.717}{0.892}{-0.720}{0.892}{0.338}
\SG{-0.720}{0.892}{-0.723}{0.892}{0.341}
\SG{-0.723}{0.892}{-0.726}{0.891}{0.344}
\SG{-0.726}{0.891}{-0.729}{0.891}{0.348}
\SG{-0.729}{0.891}{-0.732}{0.890}{0.351}
\SG{-0.732}{0.890}{-0.735}{0.890}{0.354}
\SG{-0.735}{0.890}{-0.738}{0.889}{0.357}
\SG{-0.738}{0.889}{-0.741}{0.889}{0.359}
\SG{-0.741}{0.889}{-0.744}{0.888}{0.362}
\SG{-0.744}{0.888}{-0.747}{0.888}{0.365}
\SG{-0.747}{0.888}{-0.750}{0.887}{0.364}
\SG{-0.750}{0.887}{-0.752}{0.887}{0.359}
\SG{-0.752}{0.887}{-0.755}{0.886}{0.353}
\SG{-0.755}{0.886}{-0.758}{0.886}{0.348}
\SG{-0.758}{0.886}{-0.761}{0.885}{0.342}
\SG{-0.761}{0.885}{-0.764}{0.885}{0.338}
\SG{-0.764}{0.885}{-0.767}{0.884}{0.334}
\SG{-0.767}{0.884}{-0.770}{0.884}{0.331}
\SG{-0.770}{0.884}{-0.773}{0.883}{0.328}
\SG{-0.773}{0.883}{-0.776}{0.883}{0.324}
\SG{-0.776}{0.883}{-0.779}{0.882}{0.323}
\SG{-0.779}{0.882}{-0.781}{0.881}{0.324}
\SG{-0.781}{0.881}{-0.784}{0.881}{0.325}
\SG{-0.784}{0.881}{-0.787}{0.880}{0.326}
\SG{-0.787}{0.880}{-0.790}{0.880}{0.327}
\SG{-0.790}{0.880}{-0.793}{0.879}{0.330}
\SG{-0.793}{0.879}{-0.796}{0.878}{0.334}
\SG{-0.796}{0.878}{-0.799}{0.878}{0.338}
\SG{-0.799}{0.878}{-0.802}{0.877}{0.342}
\SG{-0.802}{0.877}{-0.805}{0.876}{0.346}
\SG{-0.805}{0.876}{-0.807}{0.876}{0.351}
\SG{-0.807}{0.876}{-0.810}{0.875}{0.355}
\SG{-0.810}{0.875}{-0.813}{0.874}{0.360}
\SG{-0.813}{0.874}{-0.816}{0.874}{0.365}
\SG{-0.816}{0.874}{-0.819}{0.873}{0.370}
\SG{-0.819}{0.873}{-0.822}{0.872}{0.373}
\SG{-0.822}{0.872}{-0.825}{0.871}{0.370}
\SG{-0.825}{0.871}{-0.828}{0.871}{0.367}
\SG{-0.828}{0.871}{-0.830}{0.870}{0.363}
\SG{-0.830}{0.870}{-0.833}{0.869}{0.359}
\SG{-0.833}{0.869}{-0.836}{0.868}{0.355}
\SG{-0.836}{0.868}{-0.839}{0.868}{0.353}
\SG{-0.839}{0.868}{-0.842}{0.867}{0.352}
\SG{-0.842}{0.867}{-0.845}{0.866}{0.350}
\SG{-0.845}{0.866}{-0.847}{0.865}{0.349}
\SG{-0.847}{0.865}{-0.850}{0.864}{0.348}
\SG{-0.850}{0.864}{-0.853}{0.864}{0.350}
\SG{-0.853}{0.864}{-0.856}{0.863}{0.352}
\SG{-0.856}{0.863}{-0.859}{0.862}{0.355}
\SG{-0.859}{0.862}{-0.862}{0.861}{0.358}
\SG{-0.862}{0.861}{-0.864}{0.860}{0.360}
\SG{-0.864}{0.860}{-0.867}{0.859}{0.358}
\SG{-0.867}{0.859}{-0.870}{0.858}{0.356}
\SG{-0.870}{0.858}{-0.873}{0.858}{0.354}
\SG{-0.873}{0.858}{-0.876}{0.857}{0.351}
\SG{-0.876}{0.857}{-0.879}{0.856}{0.349}
\SG{-0.879}{0.856}{-0.881}{0.855}{0.348}
\SG{-0.881}{0.855}{-0.884}{0.854}{0.347}
\SG{-0.884}{0.854}{-0.887}{0.853}{0.346}
\SG{-0.887}{0.853}{-0.890}{0.852}{0.345}
\SG{-0.890}{0.852}{-0.893}{0.851}{0.345}
\SG{-0.893}{0.851}{-0.895}{0.850}{0.347}
\SG{-0.895}{0.850}{-0.898}{0.849}{0.350}
\SG{-0.898}{0.849}{-0.901}{0.848}{0.353}
\SG{-0.901}{0.848}{-0.904}{0.847}{0.357}
\SG{-0.904}{0.847}{-0.907}{0.846}{0.362}
\SG{-0.907}{0.846}{-0.909}{0.845}{0.363}
\SG{-0.909}{0.845}{-0.912}{0.844}{0.362}
\SG{-0.912}{0.844}{-0.915}{0.843}{0.361}
\SG{-0.915}{0.843}{-0.918}{0.842}{0.360}
\SG{-0.918}{0.842}{-0.920}{0.841}{0.358}
\SG{-0.920}{0.841}{-0.923}{0.840}{0.356}
\SG{-0.923}{0.840}{-0.926}{0.839}{0.356}
\SG{-0.926}{0.839}{-0.929}{0.838}{0.355}
\SG{-0.929}{0.838}{-0.931}{0.837}{0.355}
\SG{-0.931}{0.837}{-0.934}{0.836}{0.355}
\SG{-0.934}{0.836}{-0.937}{0.835}{0.354}
\SG{-0.937}{0.835}{-0.940}{0.833}{0.353}
\SG{-0.940}{0.833}{-0.942}{0.832}{0.351}
\SG{-0.942}{0.832}{-0.945}{0.831}{0.350}
\SG{-0.945}{0.831}{-0.948}{0.830}{0.348}
\SG{-0.948}{0.830}{-0.951}{0.829}{0.345}
\SG{-0.951}{0.829}{-0.953}{0.828}{0.345}
\SG{-0.953}{0.828}{-0.956}{0.827}{0.346}
\SG{-0.956}{0.827}{-0.959}{0.825}{0.346}
\SG{-0.959}{0.825}{-0.961}{0.824}{0.346}
\SG{-0.961}{0.824}{-0.964}{0.823}{0.346}
\SG{-0.964}{0.823}{-0.967}{0.822}{0.348}
\SG{-0.967}{0.822}{-0.969}{0.821}{0.352}
\SG{-0.969}{0.821}{-0.972}{0.819}{0.357}
\SG{-0.972}{0.819}{-0.975}{0.818}{0.362}
\SG{-0.975}{0.818}{-0.978}{0.817}{0.369}
\SG{-0.978}{0.817}{-0.980}{0.816}{0.375}
\SG{-0.980}{0.816}{-0.983}{0.814}{0.377}
\SG{-0.983}{0.814}{-0.986}{0.813}{0.378}
\SG{-0.986}{0.813}{-0.988}{0.812}{0.379}
\SG{-0.988}{0.812}{-0.991}{0.811}{0.380}
\SG{-0.991}{0.811}{-0.994}{0.809}{0.380}
\SG{-0.994}{0.809}{-0.996}{0.808}{0.381}
\SG{-0.996}{0.808}{-0.999}{0.807}{0.383}
\SG{-0.999}{0.807}{-1.002}{0.805}{0.385}
\SG{-1.002}{0.805}{-1.004}{0.804}{0.388}
\SG{-1.004}{0.804}{-1.007}{0.803}{0.391}
\SG{-1.007}{0.803}{-1.009}{0.801}{0.394}
\SG{-1.009}{0.801}{-1.012}{0.800}{0.393}
\SG{-1.012}{0.800}{-1.015}{0.799}{0.392}
\SG{-1.015}{0.799}{-1.017}{0.797}{0.391}
\SG{-1.017}{0.797}{-1.020}{0.796}{0.390}
\SG{-1.020}{0.796}{-1.023}{0.795}{0.389}
\SG{-1.023}{0.795}{-1.025}{0.793}{0.386}
\SG{-1.025}{0.793}{-1.028}{0.792}{0.383}
\SG{-1.028}{0.792}{-1.030}{0.790}{0.380}
\SG{-1.030}{0.790}{-1.033}{0.789}{0.377}
\SG{-1.033}{0.789}{-1.036}{0.788}{0.373}
\SG{-1.036}{0.788}{-1.038}{0.786}{0.368}
\SG{-1.038}{0.786}{-1.041}{0.785}{0.367}
\SG{-1.041}{0.785}{-1.043}{0.783}{0.366}
\SG{-1.043}{0.783}{-1.046}{0.782}{0.365}
\SG{-1.046}{0.782}{-1.048}{0.780}{0.363}
\SG{-1.048}{0.780}{-1.051}{0.779}{0.363}
\SG{-1.051}{0.779}{-1.054}{0.777}{0.363}
\SG{-1.054}{0.777}{-1.056}{0.776}{0.363}
\SG{-1.056}{0.776}{-1.059}{0.774}{0.363}
\SG{-1.059}{0.774}{-1.061}{0.773}{0.363}
\SG{-1.061}{0.773}{-1.064}{0.771}{0.363}
\SG{-1.064}{0.771}{-1.066}{0.770}{0.363}
\SG{-1.066}{0.770}{-1.069}{0.768}{0.365}
\SG{-1.069}{0.768}{-1.071}{0.767}{0.367}
\SG{-1.071}{0.767}{-1.074}{0.765}{0.369}
\SG{-1.074}{0.765}{-1.076}{0.764}{0.372}
\SG{-1.076}{0.764}{-1.079}{0.762}{0.377}
\SG{-1.079}{0.762}{-1.081}{0.760}{0.382}
\SG{-1.081}{0.760}{-1.084}{0.759}{0.383}
\SG{-1.084}{0.759}{-1.086}{0.757}{0.385}
\SG{-1.086}{0.757}{-1.089}{0.756}{0.386}
\SG{-1.089}{0.756}{-1.091}{0.754}{0.388}
\SG{-1.091}{0.754}{-1.094}{0.752}{0.390}
\SG{-1.094}{0.752}{-1.096}{0.751}{0.391}
\SG{-1.096}{0.751}{-1.099}{0.749}{0.391}
\SG{-1.099}{0.749}{-1.101}{0.748}{0.391}
\SG{-1.101}{0.748}{-1.104}{0.746}{0.390}
\SG{-1.104}{0.746}{-1.106}{0.744}{0.388}
\SG{-1.106}{0.744}{-1.108}{0.743}{0.386}
\SG{-1.108}{0.743}{-1.111}{0.741}{0.384}
\SG{-1.111}{0.741}{-1.113}{0.739}{0.382}
\SG{-1.113}{0.739}{-1.116}{0.738}{0.380}
\SG{-1.116}{0.738}{-1.118}{0.736}{0.377}
\SG{-1.118}{0.736}{-1.121}{0.734}{0.374}
\SG{-1.121}{0.734}{-1.123}{0.732}{0.371}
\SG{-1.123}{0.732}{-1.125}{0.731}{0.368}
\SG{-1.125}{0.731}{-1.128}{0.729}{0.367}
\SG{-1.128}{0.729}{-1.130}{0.727}{0.367}
\SG{-1.130}{0.727}{-1.133}{0.725}{0.368}
\SG{-1.133}{0.725}{-1.135}{0.724}{0.368}
\SG{-1.135}{0.724}{-1.137}{0.722}{0.369}
\SG{-1.137}{0.722}{-1.140}{0.720}{0.372}
\SG{-1.140}{0.720}{-1.142}{0.718}{0.373}
\SG{-1.142}{0.718}{-1.144}{0.717}{0.375}
\SG{-1.144}{0.717}{-1.147}{0.715}{0.377}
\SG{-1.147}{0.715}{-1.149}{0.713}{0.379}
\SG{-1.149}{0.713}{-1.151}{0.711}{0.381}
\SG{-1.151}{0.711}{-1.154}{0.709}{0.384}
\SG{-1.154}{0.709}{-1.156}{0.708}{0.386}
\SG{-1.156}{0.708}{-1.158}{0.706}{0.388}
\SG{-1.158}{0.706}{-1.161}{0.704}{0.390}
\SG{-1.161}{0.704}{-1.163}{0.702}{0.391}
\SG{-1.163}{0.702}{-1.165}{0.700}{0.391}
\SG{-1.165}{0.700}{-1.168}{0.698}{0.392}
\SG{-1.168}{0.698}{-1.170}{0.697}{0.392}
\SG{-1.170}{0.697}{-1.172}{0.695}{0.391}
\SG{-1.172}{0.695}{-1.174}{0.693}{0.391}
\SG{-1.174}{0.693}{-1.177}{0.691}{0.390}
\SG{-1.177}{0.691}{-1.179}{0.689}{0.388}
\SG{-1.179}{0.689}{-1.181}{0.687}{0.387}
\SG{-1.181}{0.687}{-1.184}{0.685}{0.386}
\SG{-1.184}{0.685}{-1.186}{0.683}{0.385}
\SG{-1.186}{0.683}{-1.188}{0.681}{0.384}
\SG{-1.188}{0.681}{-1.190}{0.679}{0.382}
\SG{-1.190}{0.679}{-1.192}{0.677}{0.380}
\SG{-1.192}{0.677}{-1.195}{0.675}{0.379}
\SG{-1.195}{0.675}{-1.197}{0.674}{0.377}
\SG{-1.197}{0.674}{-1.199}{0.672}{0.376}
\SG{-1.199}{0.672}{-1.201}{0.670}{0.375}
\SG{-1.201}{0.670}{-1.204}{0.668}{0.374}
\SG{-1.204}{0.668}{-1.206}{0.666}{0.373}
\SG{-1.206}{0.666}{-1.208}{0.664}{0.372}
\SG{-1.208}{0.664}{-1.210}{0.662}{0.371}
\SG{-1.210}{0.662}{-1.212}{0.660}{0.370}
\SG{-1.212}{0.660}{-1.214}{0.658}{0.370}
\SG{-1.214}{0.658}{-1.217}{0.656}{0.369}
\SG{-1.217}{0.656}{-1.219}{0.654}{0.369}
\SG{-1.219}{0.654}{-1.221}{0.652}{0.368}
\SG{-1.221}{0.652}{-1.223}{0.649}{0.368}
\SG{-1.223}{0.649}{-1.225}{0.647}{0.367}
\SG{-1.225}{0.647}{-1.227}{0.645}{0.367}
\SG{-1.227}{0.645}{-1.229}{0.643}{0.367}
\SG{-1.229}{0.643}{-1.231}{0.641}{0.367}
\SG{-1.231}{0.641}{-1.234}{0.639}{0.367}
\SG{-1.234}{0.639}{-1.236}{0.637}{0.367}
\SG{-1.236}{0.637}{-1.238}{0.635}{0.367}
\SG{-1.238}{0.635}{-1.240}{0.633}{0.367}
\SG{-1.240}{0.633}{-1.242}{0.631}{0.367}
\SG{-1.242}{0.631}{-1.244}{0.629}{0.368}
\SG{-1.244}{0.629}{-1.246}{0.627}{0.368}
\SG{-1.246}{0.627}{-1.248}{0.624}{0.368}
\SG{-1.248}{0.624}{-1.250}{0.622}{0.369}
\SG{-1.250}{0.622}{-1.252}{0.620}{0.370}
\SG{-1.252}{0.620}{-1.254}{0.618}{0.370}
\SG{-1.254}{0.618}{-1.256}{0.616}{0.371}
\SG{-1.256}{0.616}{-1.258}{0.614}{0.372}
\SG{-1.258}{0.614}{-1.260}{0.612}{0.373}
\SG{-1.260}{0.612}{-1.262}{0.609}{0.374}
\SG{-1.262}{0.609}{-1.264}{0.607}{0.375}
\SG{-1.264}{0.607}{-1.266}{0.605}{0.376}
\SG{-1.266}{0.605}{-1.268}{0.603}{0.377}
\SG{-1.268}{0.603}{-1.270}{0.601}{0.379}
\SG{-1.270}{0.601}{-1.272}{0.598}{0.380}
\SG{-1.272}{0.598}{-1.274}{0.596}{0.381}
\SG{-1.274}{0.596}{-1.276}{0.594}{0.382}
\SG{-1.276}{0.594}{-1.278}{0.592}{0.384}
\SG{-1.278}{0.592}{-1.280}{0.590}{0.386}
\SG{-1.280}{0.590}{-1.282}{0.587}{0.388}
\SG{-1.282}{0.587}{-1.284}{0.585}{0.390}
\SG{-1.284}{0.585}{-1.286}{0.583}{0.392}
\SG{-1.286}{0.583}{-1.288}{0.581}{0.393}
\SG{-1.288}{0.581}{-1.290}{0.578}{0.394}
\SG{-1.290}{0.578}{-1.291}{0.576}{0.396}
\SG{-1.291}{0.576}{-1.293}{0.574}{0.398}
\SG{-1.293}{0.574}{-1.295}{0.571}{0.399}
\SG{-1.295}{0.571}{-1.297}{0.569}{0.400}
\SG{-1.297}{0.569}{-1.299}{0.567}{0.401}
\SG{-1.299}{0.567}{-1.301}{0.565}{0.402}
\SG{-1.301}{0.565}{-1.303}{0.562}{0.402}
\SG{-1.303}{0.562}{-1.305}{0.560}{0.402}
\SG{-1.305}{0.560}{-1.306}{0.558}{0.402}
\SG{-1.306}{0.558}{-1.308}{0.555}{0.401}
\SG{-1.308}{0.555}{-1.310}{0.553}{0.400}
\SG{-1.310}{0.553}{-1.312}{0.551}{0.398}
\SG{-1.312}{0.551}{-1.314}{0.548}{0.396}
\SG{-1.314}{0.548}{-1.315}{0.546}{0.393}
\SG{-1.315}{0.546}{-1.317}{0.544}{0.391}
\SG{-1.317}{0.544}{-1.319}{0.541}{0.389}
\SG{-1.319}{0.541}{-1.321}{0.539}{0.388}
\SG{-1.321}{0.539}{-1.323}{0.537}{0.387}
\SG{-1.323}{0.537}{-1.324}{0.534}{0.385}
\SG{-1.324}{0.534}{-1.326}{0.532}{0.384}
\SG{-1.326}{0.532}{-1.328}{0.529}{0.383}
\SG{-1.328}{0.529}{-1.330}{0.527}{0.383}
\SG{-1.330}{0.527}{-1.331}{0.525}{0.383}
\SG{-1.331}{0.525}{-1.333}{0.522}{0.384}
\SG{-1.333}{0.522}{-1.335}{0.520}{0.385}
\SG{-1.335}{0.520}{-1.336}{0.517}{0.389}
\SG{-1.336}{0.517}{-1.338}{0.515}{0.392}
\SG{-1.338}{0.515}{-1.340}{0.513}{0.396}
\SG{-1.340}{0.513}{-1.341}{0.510}{0.399}
\SG{-1.341}{0.510}{-1.343}{0.508}{0.401}
\SG{-1.343}{0.508}{-1.345}{0.505}{0.403}
\SG{-1.345}{0.505}{-1.346}{0.503}{0.406}
\SG{-1.346}{0.503}{-1.348}{0.500}{0.408}
\SG{-1.348}{0.500}{-1.350}{0.498}{0.410}
\SG{-1.350}{0.498}{-1.351}{0.495}{0.411}
\SG{-1.351}{0.495}{-1.353}{0.493}{0.412}
\SG{-1.353}{0.493}{-1.355}{0.490}{0.411}
\SG{-1.355}{0.490}{-1.356}{0.488}{0.410}
\SG{-1.356}{0.488}{-1.358}{0.485}{0.408}
\SG{-1.358}{0.485}{-1.359}{0.483}{0.407}
\SG{-1.359}{0.483}{-1.361}{0.480}{0.406}
\SG{-1.361}{0.480}{-1.363}{0.478}{0.405}
\SG{-1.363}{0.478}{-1.364}{0.475}{0.402}
\SG{-1.364}{0.475}{-1.366}{0.473}{0.398}
\SG{-1.366}{0.473}{-1.367}{0.470}{0.394}
\SG{-1.367}{0.470}{-1.369}{0.468}{0.392}
\SG{-1.369}{0.468}{-1.370}{0.465}{0.390}
\SG{-1.370}{0.465}{-1.372}{0.463}{0.388}
\SG{-1.372}{0.463}{-1.373}{0.460}{0.388}
\SG{-1.373}{0.460}{-1.375}{0.458}{0.388}
\SG{-1.375}{0.458}{-1.376}{0.455}{0.388}
\SG{-1.376}{0.455}{-1.378}{0.453}{0.388}
\SG{-1.378}{0.453}{-1.379}{0.450}{0.389}
\SG{-1.379}{0.450}{-1.381}{0.448}{0.389}
\SG{-1.381}{0.448}{-1.382}{0.445}{0.389}
\SG{-1.382}{0.445}{-1.384}{0.442}{0.391}
\SG{-1.384}{0.442}{-1.385}{0.440}{0.393}
\SG{-1.385}{0.440}{-1.387}{0.437}{0.394}
\SG{-1.387}{0.437}{-1.388}{0.435}{0.395}
\SG{-1.388}{0.435}{-1.389}{0.432}{0.398}
\SG{-1.389}{0.432}{-1.391}{0.430}{0.403}
\SG{-1.391}{0.430}{-1.392}{0.427}{0.407}
\SG{-1.392}{0.427}{-1.394}{0.424}{0.411}
\SG{-1.394}{0.424}{-1.395}{0.422}{0.414}
\SG{-1.395}{0.422}{-1.396}{0.419}{0.417}
\SG{-1.396}{0.419}{-1.398}{0.416}{0.420}
\SG{-1.398}{0.416}{-1.399}{0.414}{0.422}
\SG{-1.399}{0.414}{-1.401}{0.411}{0.423}
\SG{-1.401}{0.411}{-1.402}{0.409}{0.424}
\SG{-1.402}{0.409}{-1.403}{0.406}{0.426}
\SG{-1.403}{0.406}{-1.405}{0.403}{0.426}
\SG{-1.405}{0.403}{-1.406}{0.401}{0.423}
\SG{-1.406}{0.401}{-1.407}{0.398}{0.420}
\SG{-1.407}{0.398}{-1.408}{0.395}{0.418}
\SG{-1.408}{0.395}{-1.410}{0.393}{0.416}
\SG{-1.410}{0.393}{-1.411}{0.390}{0.414}
\SG{-1.411}{0.390}{-1.412}{0.387}{0.415}
\SG{-1.412}{0.387}{-1.414}{0.385}{0.415}
\SG{-1.414}{0.385}{-1.415}{0.382}{0.414}
\SG{-1.415}{0.382}{-1.416}{0.379}{0.413}
\SG{-1.416}{0.379}{-1.417}{0.377}{0.412}
\SG{-1.417}{0.377}{-1.419}{0.374}{0.409}
\SG{-1.419}{0.374}{-1.420}{0.371}{0.403}
\SG{-1.420}{0.371}{-1.421}{0.369}{0.397}
\SG{-1.421}{0.369}{-1.422}{0.366}{0.392}
\SG{-1.422}{0.366}{-1.423}{0.363}{0.388}
\SG{-1.423}{0.363}{-1.425}{0.361}{0.385}
\SG{-1.425}{0.361}{-1.426}{0.358}{0.384}
\SG{-1.426}{0.358}{-1.427}{0.355}{0.385}
\SG{-1.427}{0.355}{-1.428}{0.352}{0.385}
\SG{-1.428}{0.352}{-1.429}{0.350}{0.385}
\SG{-1.429}{0.350}{-1.430}{0.347}{0.385}
\SG{-1.430}{0.347}{-1.432}{0.344}{0.387}
\SG{-1.432}{0.344}{-1.433}{0.341}{0.389}
\SG{-1.433}{0.341}{-1.434}{0.339}{0.392}
\SG{-1.434}{0.339}{-1.435}{0.336}{0.394}
\SG{-1.435}{0.336}{-1.436}{0.333}{0.396}
\SG{-1.436}{0.333}{-1.437}{0.330}{0.397}
\SG{-1.437}{0.330}{-1.438}{0.328}{0.397}
\SG{-1.438}{0.328}{-1.439}{0.325}{0.398}
\SG{-1.439}{0.325}{-1.440}{0.322}{0.399}
\SG{-1.440}{0.322}{-1.441}{0.319}{0.400}
\SG{-1.441}{0.319}{-1.442}{0.317}{0.401}
\SG{-1.442}{0.317}{-1.443}{0.314}{0.403}
\SG{-1.443}{0.314}{-1.444}{0.311}{0.405}
\SG{-1.444}{0.311}{-1.445}{0.308}{0.407}
\SG{-1.445}{0.308}{-1.447}{0.306}{0.409}
\SG{-1.447}{0.306}{-1.448}{0.303}{0.410}
\SG{-1.448}{0.303}{-1.448}{0.300}{0.407}
\SG{-1.448}{0.300}{-1.449}{0.297}{0.403}
\SG{-1.449}{0.297}{-1.450}{0.294}{0.399}
\SG{-1.450}{0.294}{-1.451}{0.292}{0.397}
\SG{-1.451}{0.292}{-1.452}{0.289}{0.394}
\SG{-1.452}{0.289}{-1.453}{0.286}{0.394}
\SG{-1.453}{0.286}{-1.454}{0.283}{0.395}
\SG{-1.454}{0.283}{-1.455}{0.280}{0.396}
\SG{-1.455}{0.280}{-1.456}{0.278}{0.397}
\SG{-1.456}{0.278}{-1.457}{0.275}{0.398}
\SG{-1.457}{0.275}{-1.458}{0.272}{0.400}
\SG{-1.458}{0.272}{-1.459}{0.269}{0.403}
\SG{-1.459}{0.269}{-1.460}{0.266}{0.406}
\SG{-1.460}{0.266}{-1.461}{0.264}{0.409}
\SG{-1.461}{0.264}{-1.461}{0.261}{0.411}
\SG{-1.461}{0.261}{-1.462}{0.258}{0.412}
\SG{-1.462}{0.258}{-1.463}{0.255}{0.409}
\SG{-1.463}{0.255}{-1.464}{0.252}{0.407}
\SG{-1.464}{0.252}{-1.465}{0.249}{0.405}
\SG{-1.465}{0.249}{-1.466}{0.247}{0.403}
\SG{-1.466}{0.247}{-1.466}{0.244}{0.402}
\SG{-1.466}{0.244}{-1.467}{0.241}{0.404}
\SG{-1.467}{0.241}{-1.468}{0.238}{0.406}
\SG{-1.468}{0.238}{-1.469}{0.235}{0.408}
\SG{-1.469}{0.235}{-1.469}{0.232}{0.410}
\SG{-1.469}{0.232}{-1.470}{0.229}{0.412}
\SG{-1.470}{0.229}{-1.471}{0.227}{0.418}
\SG{-1.471}{0.227}{-1.472}{0.224}{0.422}
\SG{-1.472}{0.224}{-1.472}{0.221}{0.426}
\SG{-1.472}{0.221}{-1.473}{0.218}{0.429}
\SG{-1.473}{0.218}{-1.474}{0.215}{0.432}
\SG{-1.474}{0.215}{-1.475}{0.212}{0.428}
\SG{-1.475}{0.212}{-1.475}{0.209}{0.423}
\SG{-1.475}{0.209}{-1.476}{0.207}{0.419}
\SG{-1.476}{0.207}{-1.477}{0.204}{0.414}
\SG{-1.477}{0.204}{-1.477}{0.201}{0.411}
\SG{-1.477}{0.201}{-1.478}{0.198}{0.406}
\SG{-1.478}{0.198}{-1.479}{0.195}{0.402}
\SG{-1.479}{0.195}{-1.479}{0.192}{0.399}
\SG{-1.479}{0.192}{-1.480}{0.189}{0.395}
\SG{-1.480}{0.189}{-1.480}{0.186}{0.392}
\SG{-1.480}{0.186}{-1.481}{0.183}{0.391}
\SG{-1.481}{0.183}{-1.482}{0.181}{0.390}
\SG{-1.482}{0.181}{-1.482}{0.178}{0.389}
\SG{-1.482}{0.178}{-1.483}{0.175}{0.389}
\SG{-1.483}{0.175}{-1.483}{0.172}{0.389}
\SG{-1.483}{0.172}{-1.484}{0.169}{0.392}
\SG{-1.484}{0.169}{-1.485}{0.166}{0.396}
\SG{-1.485}{0.166}{-1.485}{0.163}{0.400}
\SG{-1.485}{0.163}{-1.486}{0.160}{0.403}
\SG{-1.486}{0.160}{-1.486}{0.157}{0.407}
\SG{-1.486}{0.157}{-1.487}{0.154}{0.413}
\SG{-1.487}{0.154}{-1.487}{0.151}{0.419}
\SG{-1.487}{0.151}{-1.488}{0.149}{0.425}
\SG{-1.488}{0.149}{-1.488}{0.146}{0.430}
\SG{-1.488}{0.146}{-1.489}{0.143}{0.436}
\SG{-1.489}{0.143}{-1.489}{0.140}{0.434}
\SG{-1.489}{0.140}{-1.490}{0.137}{0.432}
\SG{-1.490}{0.137}{-1.490}{0.134}{0.430}
\SG{-1.490}{0.134}{-1.490}{0.131}{0.428}
\SG{-1.490}{0.131}{-1.491}{0.128}{0.426}
\SG{-1.491}{0.128}{-1.491}{0.125}{0.423}
\SG{-1.491}{0.125}{-1.492}{0.122}{0.420}
\SG{-1.492}{0.122}{-1.492}{0.119}{0.417}
\SG{-1.492}{0.119}{-1.492}{0.116}{0.414}
\SG{-1.492}{0.116}{-1.493}{0.113}{0.412}
\SG{-1.493}{0.113}{-1.493}{0.111}{0.409}
\SG{-1.493}{0.111}{-1.494}{0.108}{0.407}
\SG{-1.494}{0.108}{-1.494}{0.105}{0.404}
\SG{-1.494}{0.105}{-1.494}{0.102}{0.402}
\SG{-1.494}{0.102}{-1.495}{0.099}{0.399}
\SG{-1.495}{0.099}{-1.495}{0.096}{0.397}
\SG{-1.495}{0.096}{-1.495}{0.093}{0.395}
\SG{-1.495}{0.093}{-1.495}{0.090}{0.393}
\SG{-1.495}{0.090}{-1.496}{0.087}{0.391}
\SG{-1.496}{0.087}{-1.496}{0.084}{0.389}
\SG{-1.496}{0.084}{-1.496}{0.081}{0.388}
\SG{-1.496}{0.081}{-1.497}{0.078}{0.386}
\SG{-1.497}{0.078}{-1.497}{0.075}{0.384}
\SG{-1.497}{0.075}{-1.497}{0.072}{0.383}
\SG{-1.497}{0.072}{-1.497}{0.069}{0.381}
\SG{-1.497}{0.069}{-1.498}{0.066}{0.380}
\SG{-1.498}{0.066}{-1.498}{0.063}{0.379}
\SG{-1.498}{0.063}{-1.498}{0.061}{0.378}
\SG{-1.498}{0.061}{-1.498}{0.058}{0.377}
\SG{-1.498}{0.058}{-1.498}{0.055}{0.375}
\SG{-1.498}{0.055}{-1.498}{0.052}{0.375}
\SG{-1.498}{0.052}{-1.499}{0.049}{0.374}
\SG{-1.499}{0.049}{-1.499}{0.046}{0.373}
\SG{-1.499}{0.046}{-1.499}{0.043}{0.372}
\SG{-1.499}{0.043}{-1.499}{0.040}{0.371}
\SG{-1.499}{0.040}{-1.499}{0.037}{0.370}
\SG{-1.499}{0.037}{-1.499}{0.034}{0.370}
\SG{-1.499}{0.034}{-1.499}{0.031}{0.369}
\SG{-1.499}{0.031}{-1.500}{0.028}{0.369}
\SG{-1.500}{0.028}{-1.500}{0.025}{0.368}
\SG{-1.500}{0.025}{-1.500}{0.022}{0.368}
\SG{-1.500}{0.022}{-1.500}{0.019}{0.367}
\SG{-1.500}{0.019}{-1.500}{0.016}{0.367}
\SG{-1.500}{0.016}{-1.500}{0.013}{0.367}
\SG{-1.500}{0.013}{-1.500}{0.010}{0.366}
\SG{-1.500}{0.010}{-1.500}{0.007}{0.366}
\SG{-1.500}{0.007}{-1.500}{0.004}{0.366}
\SG{-1.500}{0.004}{-1.500}{0.001}{0.366}
\SG{-1.500}{0.001}{-1.500}{-0.001}{0.366}
\SG{-1.500}{-0.001}{-1.500}{-0.004}{0.366}
\SG{-1.500}{-0.004}{-1.500}{-0.007}{0.366}
\SG{-1.500}{-0.007}{-1.500}{-0.010}{0.366}
\SG{-1.500}{-0.010}{-1.500}{-0.013}{0.366}
\SG{-1.500}{-0.013}{-1.500}{-0.016}{0.366}
\SG{-1.500}{-0.016}{-1.500}{-0.019}{0.367}
\SG{-1.500}{-0.019}{-1.500}{-0.022}{0.367}
\SG{-1.500}{-0.022}{-1.500}{-0.025}{0.367}
\SG{-1.500}{-0.025}{-1.500}{-0.028}{0.368}
\SG{-1.500}{-0.028}{-1.499}{-0.031}{0.368}
\SG{-1.499}{-0.031}{-1.499}{-0.034}{0.369}
\SG{-1.499}{-0.034}{-1.499}{-0.037}{0.369}
\SG{-1.499}{-0.037}{-1.499}{-0.040}{0.370}
\SG{-1.499}{-0.040}{-1.499}{-0.043}{0.370}
\SG{-1.499}{-0.043}{-1.499}{-0.046}{0.371}
\SG{-1.499}{-0.046}{-1.499}{-0.049}{0.372}
\SG{-1.499}{-0.049}{-1.498}{-0.052}{0.373}
\SG{-1.498}{-0.052}{-1.498}{-0.055}{0.374}
\SG{-1.498}{-0.055}{-1.498}{-0.058}{0.374}
\SG{-1.498}{-0.058}{-1.498}{-0.061}{0.375}
\SG{-1.498}{-0.061}{-1.498}{-0.063}{0.376}
\SG{-1.498}{-0.063}{-1.498}{-0.066}{0.378}
\SG{-1.498}{-0.066}{-1.497}{-0.069}{0.379}
\SG{-1.497}{-0.069}{-1.497}{-0.072}{0.380}
\SG{-1.497}{-0.072}{-1.497}{-0.075}{0.381}
\SG{-1.497}{-0.075}{-1.497}{-0.078}{0.383}
\SG{-1.497}{-0.078}{-1.496}{-0.081}{0.384}
\SG{-1.496}{-0.081}{-1.496}{-0.084}{0.386}
\SG{-1.496}{-0.084}{-1.496}{-0.087}{0.388}
\SG{-1.496}{-0.087}{-1.495}{-0.090}{0.389}
\SG{-1.495}{-0.090}{-1.495}{-0.093}{0.391}
\SG{-1.495}{-0.093}{-1.495}{-0.096}{0.393}
\SG{-1.495}{-0.096}{-1.495}{-0.099}{0.395}
\SG{-1.495}{-0.099}{-1.494}{-0.102}{0.397}
\SG{-1.494}{-0.102}{-1.494}{-0.105}{0.399}
\SG{-1.494}{-0.105}{-1.494}{-0.108}{0.402}
\SG{-1.494}{-0.108}{-1.493}{-0.111}{0.404}
\SG{-1.493}{-0.111}{-1.493}{-0.113}{0.407}
\SG{-1.493}{-0.113}{-1.492}{-0.116}{0.409}
\SG{-1.492}{-0.116}{-1.492}{-0.119}{0.412}
\SG{-1.492}{-0.119}{-1.492}{-0.122}{0.414}
\SG{-1.492}{-0.122}{-1.491}{-0.125}{0.417}
\SG{-1.491}{-0.125}{-1.491}{-0.128}{0.420}
\SG{-1.491}{-0.128}{-1.490}{-0.131}{0.423}
\SG{-1.490}{-0.131}{-1.490}{-0.134}{0.426}
\SG{-1.490}{-0.134}{-1.490}{-0.137}{0.428}
\SG{-1.490}{-0.137}{-1.489}{-0.140}{0.430}
\SG{-1.489}{-0.140}{-1.489}{-0.143}{0.432}
\SG{-1.489}{-0.143}{-1.488}{-0.146}{0.435}
\SG{-1.488}{-0.146}{-1.488}{-0.149}{0.436}
\SG{-1.488}{-0.149}{-1.487}{-0.151}{0.430}
\SG{-1.487}{-0.151}{-1.487}{-0.154}{0.425}
\SG{-1.487}{-0.154}{-1.486}{-0.157}{0.419}
\SG{-1.486}{-0.157}{-1.486}{-0.160}{0.412}
\SG{-1.486}{-0.160}{-1.485}{-0.163}{0.407}
\SG{-1.485}{-0.163}{-1.485}{-0.166}{0.403}
\SG{-1.485}{-0.166}{-1.484}{-0.169}{0.399}
\SG{-1.484}{-0.169}{-1.483}{-0.172}{0.396}
\SG{-1.483}{-0.172}{-1.483}{-0.175}{0.392}
\SG{-1.483}{-0.175}{-1.482}{-0.178}{0.389}
\SG{-1.482}{-0.178}{-1.482}{-0.181}{0.389}
\SG{-1.482}{-0.181}{-1.481}{-0.183}{0.389}
\SG{-1.481}{-0.183}{-1.480}{-0.186}{0.390}
\SG{-1.480}{-0.186}{-1.480}{-0.189}{0.391}
\SG{-1.480}{-0.189}{-1.479}{-0.192}{0.392}
\SG{-1.479}{-0.192}{-1.479}{-0.195}{0.395}
\SG{-1.479}{-0.195}{-1.478}{-0.198}{0.399}
\SG{-1.478}{-0.198}{-1.477}{-0.201}{0.402}
\SG{-1.477}{-0.201}{-1.477}{-0.204}{0.406}
\SG{-1.477}{-0.204}{-1.476}{-0.207}{0.411}
\SG{-1.476}{-0.207}{-1.475}{-0.209}{0.414}
\SG{-1.475}{-0.209}{-1.475}{-0.212}{0.419}
\SG{-1.475}{-0.212}{-1.474}{-0.215}{0.423}
\SG{-1.474}{-0.215}{-1.473}{-0.218}{0.428}
\SG{-1.473}{-0.218}{-1.472}{-0.221}{0.432}
\SG{-1.472}{-0.221}{-1.472}{-0.224}{0.429}
\SG{-1.472}{-0.224}{-1.471}{-0.227}{0.426}
\SG{-1.471}{-0.227}{-1.470}{-0.229}{0.422}
\SG{-1.470}{-0.229}{-1.469}{-0.232}{0.418}
\SG{-1.469}{-0.232}{-1.469}{-0.235}{0.412}
\SG{-1.469}{-0.235}{-1.468}{-0.238}{0.410}
\SG{-1.468}{-0.238}{-1.467}{-0.241}{0.408}
\SG{-1.467}{-0.241}{-1.466}{-0.244}{0.406}
\SG{-1.466}{-0.244}{-1.466}{-0.247}{0.404}
\SG{-1.466}{-0.247}{-1.465}{-0.249}{0.402}
\SG{-1.465}{-0.249}{-1.464}{-0.252}{0.403}
\SG{-1.464}{-0.252}{-1.463}{-0.255}{0.405}
\SG{-1.463}{-0.255}{-1.462}{-0.258}{0.407}
\SG{-1.462}{-0.258}{-1.461}{-0.261}{0.410}
\SG{-1.461}{-0.261}{-1.461}{-0.264}{0.412}
\SG{-1.461}{-0.264}{-1.460}{-0.266}{0.411}
\SG{-1.460}{-0.266}{-1.459}{-0.269}{0.409}
\SG{-1.459}{-0.269}{-1.458}{-0.272}{0.406}
\SG{-1.458}{-0.272}{-1.457}{-0.275}{0.403}
\SG{-1.457}{-0.275}{-1.456}{-0.278}{0.400}
\SG{-1.456}{-0.278}{-1.455}{-0.280}{0.398}
\SG{-1.455}{-0.280}{-1.454}{-0.283}{0.397}
\SG{-1.454}{-0.283}{-1.453}{-0.286}{0.396}
\SG{-1.453}{-0.286}{-1.452}{-0.289}{0.395}
\SG{-1.452}{-0.289}{-1.451}{-0.292}{0.394}
\SG{-1.451}{-0.292}{-1.450}{-0.294}{0.394}
\SG{-1.450}{-0.294}{-1.449}{-0.297}{0.397}
\SG{-1.449}{-0.297}{-1.448}{-0.300}{0.399}
\SG{-1.448}{-0.300}{-1.448}{-0.303}{0.403}
\SG{-1.448}{-0.303}{-1.447}{-0.306}{0.407}
\SG{-1.447}{-0.306}{-1.445}{-0.308}{0.410}
\SG{-1.445}{-0.308}{-1.444}{-0.311}{0.409}
\SG{-1.444}{-0.311}{-1.443}{-0.314}{0.407}
\SG{-1.443}{-0.314}{-1.442}{-0.317}{0.405}
\SG{-1.442}{-0.317}{-1.441}{-0.319}{0.403}
\SG{-1.441}{-0.319}{-1.440}{-0.322}{0.401}
\SG{-1.440}{-0.322}{-1.439}{-0.325}{0.400}
\SG{-1.439}{-0.325}{-1.438}{-0.328}{0.399}
\SG{-1.438}{-0.328}{-1.437}{-0.330}{0.398}
\SG{-1.437}{-0.330}{-1.436}{-0.333}{0.397}
\SG{-1.436}{-0.333}{-1.435}{-0.336}{0.397}
\SG{-1.435}{-0.336}{-1.434}{-0.339}{0.396}
\SG{-1.434}{-0.339}{-1.433}{-0.341}{0.394}
\SG{-1.433}{-0.341}{-1.432}{-0.344}{0.392}
\SG{-1.432}{-0.344}{-1.430}{-0.347}{0.389}
\SG{-1.430}{-0.347}{-1.429}{-0.350}{0.386}
\SG{-1.429}{-0.350}{-1.428}{-0.352}{0.385}
\SG{-1.428}{-0.352}{-1.427}{-0.355}{0.385}
\SG{-1.427}{-0.355}{-1.426}{-0.358}{0.385}
\SG{-1.426}{-0.358}{-1.425}{-0.361}{0.385}
\SG{-1.425}{-0.361}{-1.423}{-0.363}{0.384}
\SG{-1.423}{-0.363}{-1.422}{-0.366}{0.385}
\SG{-1.422}{-0.366}{-1.421}{-0.369}{0.388}
\SG{-1.421}{-0.369}{-1.420}{-0.371}{0.392}
\SG{-1.420}{-0.371}{-1.419}{-0.374}{0.397}
\SG{-1.419}{-0.374}{-1.417}{-0.377}{0.403}
\SG{-1.417}{-0.377}{-1.416}{-0.379}{0.410}
\SG{-1.416}{-0.379}{-1.415}{-0.382}{0.412}
\SG{-1.415}{-0.382}{-1.414}{-0.385}{0.413}
\SG{-1.414}{-0.385}{-1.412}{-0.387}{0.414}
\SG{-1.412}{-0.387}{-1.411}{-0.390}{0.415}
\SG{-1.411}{-0.390}{-1.410}{-0.393}{0.415}
\SG{-1.410}{-0.393}{-1.408}{-0.395}{0.414}
\SG{-1.408}{-0.395}{-1.407}{-0.398}{0.416}
\SG{-1.407}{-0.398}{-1.406}{-0.401}{0.418}
\SG{-1.406}{-0.401}{-1.405}{-0.403}{0.420}
\SG{-1.405}{-0.403}{-1.403}{-0.406}{0.423}
\SG{-1.403}{-0.406}{-1.402}{-0.409}{0.426}
\SG{-1.402}{-0.409}{-1.401}{-0.411}{0.425}
\SG{-1.401}{-0.411}{-1.399}{-0.414}{0.424}
\SG{-1.399}{-0.414}{-1.398}{-0.416}{0.423}
\SG{-1.398}{-0.416}{-1.396}{-0.419}{0.422}
\SG{-1.396}{-0.419}{-1.395}{-0.422}{0.420}
\SG{-1.395}{-0.422}{-1.394}{-0.424}{0.417}
\SG{-1.394}{-0.424}{-1.392}{-0.427}{0.414}
\SG{-1.392}{-0.427}{-1.391}{-0.430}{0.411}
\SG{-1.391}{-0.430}{-1.389}{-0.432}{0.407}
\SG{-1.389}{-0.432}{-1.388}{-0.435}{0.403}
\SG{-1.388}{-0.435}{-1.387}{-0.437}{0.398}
\SG{-1.387}{-0.437}{-1.385}{-0.440}{0.395}
\SG{-1.385}{-0.440}{-1.384}{-0.442}{0.394}
\SG{-1.384}{-0.442}{-1.382}{-0.445}{0.393}
\SG{-1.382}{-0.445}{-1.381}{-0.448}{0.391}
\SG{-1.381}{-0.448}{-1.379}{-0.450}{0.389}
\SG{-1.379}{-0.450}{-1.378}{-0.453}{0.389}
\SG{-1.378}{-0.453}{-1.376}{-0.455}{0.389}
\SG{-1.376}{-0.455}{-1.375}{-0.458}{0.388}
\SG{-1.375}{-0.458}{-1.373}{-0.460}{0.388}
\SG{-1.373}{-0.460}{-1.372}{-0.463}{0.388}
\SG{-1.372}{-0.463}{-1.370}{-0.465}{0.388}
\SG{-1.370}{-0.465}{-1.369}{-0.468}{0.388}
\SG{-1.369}{-0.468}{-1.367}{-0.470}{0.390}
\SG{-1.367}{-0.470}{-1.366}{-0.473}{0.392}
\SG{-1.366}{-0.473}{-1.364}{-0.475}{0.394}
\SG{-1.364}{-0.475}{-1.363}{-0.478}{0.398}
\SG{-1.363}{-0.478}{-1.361}{-0.480}{0.402}
\SG{-1.361}{-0.480}{-1.359}{-0.483}{0.405}
\SG{-1.359}{-0.483}{-1.358}{-0.485}{0.406}
\SG{-1.358}{-0.485}{-1.356}{-0.488}{0.407}
\SG{-1.356}{-0.488}{-1.355}{-0.490}{0.408}
\SG{-1.355}{-0.490}{-1.353}{-0.493}{0.410}
\SG{-1.353}{-0.493}{-1.351}{-0.495}{0.412}
\SG{-1.351}{-0.495}{-1.350}{-0.498}{0.412}
\SG{-1.350}{-0.498}{-1.348}{-0.500}{0.411}
\SG{-1.348}{-0.500}{-1.346}{-0.503}{0.410}
\SG{-1.346}{-0.503}{-1.345}{-0.505}{0.408}
\SG{-1.345}{-0.505}{-1.343}{-0.508}{0.406}
\SG{-1.343}{-0.508}{-1.341}{-0.510}{0.403}
\SG{-1.341}{-0.510}{-1.340}{-0.513}{0.401}
\SG{-1.340}{-0.513}{-1.338}{-0.515}{0.399}
\SG{-1.338}{-0.515}{-1.336}{-0.517}{0.396}
\SG{-1.336}{-0.517}{-1.335}{-0.520}{0.392}
\SG{-1.335}{-0.520}{-1.333}{-0.522}{0.389}
\SG{-1.333}{-0.522}{-1.331}{-0.525}{0.385}
\SG{-1.331}{-0.525}{-1.330}{-0.527}{0.384}
\SG{-1.330}{-0.527}{-1.328}{-0.529}{0.383}
\SG{-1.328}{-0.529}{-1.326}{-0.532}{0.383}
\SG{-1.326}{-0.532}{-1.324}{-0.534}{0.383}
\SG{-1.324}{-0.534}{-1.323}{-0.537}{0.384}
\SG{-1.323}{-0.537}{-1.321}{-0.539}{0.385}
\SG{-1.321}{-0.539}{-1.319}{-0.541}{0.387}
\SG{-1.319}{-0.541}{-1.317}{-0.544}{0.388}
\SG{-1.317}{-0.544}{-1.315}{-0.546}{0.389}
\SG{-1.315}{-0.546}{-1.314}{-0.548}{0.391}
\SG{-1.314}{-0.548}{-1.312}{-0.551}{0.393}
\SG{-1.312}{-0.551}{-1.310}{-0.553}{0.396}
\SG{-1.310}{-0.553}{-1.308}{-0.555}{0.398}
\SG{-1.308}{-0.555}{-1.306}{-0.558}{0.400}
\SG{-1.306}{-0.558}{-1.305}{-0.560}{0.401}
\SG{-1.305}{-0.560}{-1.303}{-0.562}{0.402}
\SG{-1.303}{-0.562}{-1.301}{-0.565}{0.402}
\SG{-1.301}{-0.565}{-1.299}{-0.567}{0.402}
\SG{-1.299}{-0.567}{-1.297}{-0.569}{0.402}
\SG{-1.297}{-0.569}{-1.295}{-0.571}{0.401}
\SG{-1.295}{-0.571}{-1.293}{-0.574}{0.400}
\SG{-1.293}{-0.574}{-1.291}{-0.576}{0.399}
\SG{-1.291}{-0.576}{-1.290}{-0.578}{0.398}
\SG{-1.290}{-0.578}{-1.288}{-0.581}{0.396}
\SG{-1.288}{-0.581}{-1.286}{-0.583}{0.394}
\SG{-1.286}{-0.583}{-1.284}{-0.585}{0.393}
\SG{-1.284}{-0.585}{-1.282}{-0.587}{0.392}
\SG{-1.282}{-0.587}{-1.280}{-0.590}{0.390}
\SG{-1.280}{-0.590}{-1.278}{-0.592}{0.388}
\SG{-1.278}{-0.592}{-1.276}{-0.594}{0.386}
\SG{-1.276}{-0.594}{-1.274}{-0.596}{0.384}
\SG{-1.274}{-0.596}{-1.272}{-0.598}{0.382}
\SG{-1.272}{-0.598}{-1.270}{-0.601}{0.381}
\SG{-1.270}{-0.601}{-1.268}{-0.603}{0.380}
\SG{-1.268}{-0.603}{-1.266}{-0.605}{0.379}
\SG{-1.266}{-0.605}{-1.264}{-0.607}{0.377}
\SG{-1.264}{-0.607}{-1.262}{-0.609}{0.376}
\SG{-1.262}{-0.609}{-1.260}{-0.612}{0.375}
\SG{-1.260}{-0.612}{-1.258}{-0.614}{0.374}
\SG{-1.258}{-0.614}{-1.256}{-0.616}{0.373}
\SG{-1.256}{-0.616}{-1.254}{-0.618}{0.372}
\SG{-1.254}{-0.618}{-1.252}{-0.620}{0.371}
\SG{-1.252}{-0.620}{-1.250}{-0.622}{0.370}
\SG{-1.250}{-0.622}{-1.248}{-0.624}{0.370}
\SG{-1.248}{-0.624}{-1.246}{-0.627}{0.369}
\SG{-1.246}{-0.627}{-1.244}{-0.629}{0.368}
\SG{-1.244}{-0.629}{-1.242}{-0.631}{0.368}
\SG{-1.242}{-0.631}{-1.240}{-0.633}{0.368}
\SG{-1.240}{-0.633}{-1.238}{-0.635}{0.367}
\SG{-1.238}{-0.635}{-1.236}{-0.637}{0.367}
\SG{-1.236}{-0.637}{-1.234}{-0.639}{0.367}
\SG{-1.234}{-0.639}{-1.231}{-0.641}{0.367}
\SG{-1.231}{-0.641}{-1.229}{-0.643}{0.367}
\SG{-1.229}{-0.643}{-1.227}{-0.645}{0.367}
\SG{-1.227}{-0.645}{-1.225}{-0.647}{0.367}
\SG{-1.225}{-0.647}{-1.223}{-0.649}{0.367}
\SG{-1.223}{-0.649}{-1.221}{-0.652}{0.367}
\SG{-1.221}{-0.652}{-1.219}{-0.654}{0.368}
\SG{-1.219}{-0.654}{-1.217}{-0.656}{0.368}
\SG{-1.217}{-0.656}{-1.214}{-0.658}{0.369}
\SG{-1.214}{-0.658}{-1.212}{-0.660}{0.369}
\SG{-1.212}{-0.660}{-1.210}{-0.662}{0.370}
\SG{-1.210}{-0.662}{-1.208}{-0.664}{0.370}
\SG{-1.208}{-0.664}{-1.206}{-0.666}{0.371}
\SG{-1.206}{-0.666}{-1.204}{-0.668}{0.372}
\SG{-1.204}{-0.668}{-1.201}{-0.670}{0.373}
\SG{-1.201}{-0.670}{-1.199}{-0.672}{0.374}
\SG{-1.199}{-0.672}{-1.197}{-0.674}{0.375}
\SG{-1.197}{-0.674}{-1.195}{-0.675}{0.376}
\SG{-1.195}{-0.675}{-1.192}{-0.677}{0.377}
\SG{-1.192}{-0.677}{-1.190}{-0.679}{0.379}
\SG{-1.190}{-0.679}{-1.188}{-0.681}{0.380}
\SG{-1.188}{-0.681}{-1.186}{-0.683}{0.382}
\SG{-1.186}{-0.683}{-1.184}{-0.685}{0.384}
\SG{-1.184}{-0.685}{-1.181}{-0.687}{0.385}
\SG{-1.181}{-0.687}{-1.179}{-0.689}{0.385}
\SG{-1.179}{-0.689}{-1.177}{-0.691}{0.387}
\SG{-1.177}{-0.691}{-1.174}{-0.693}{0.388}
\SG{-1.174}{-0.693}{-1.172}{-0.695}{0.390}
\SG{-1.172}{-0.695}{-1.170}{-0.697}{0.391}
\SG{-1.170}{-0.697}{-1.168}{-0.698}{0.391}
\SG{-1.168}{-0.698}{-1.165}{-0.700}{0.391}
\SG{-1.165}{-0.700}{-1.163}{-0.702}{0.392}
\SG{-1.163}{-0.702}{-1.161}{-0.704}{0.391}
\SG{-1.161}{-0.704}{-1.158}{-0.706}{0.391}
\SG{-1.158}{-0.706}{-1.156}{-0.708}{0.390}
\SG{-1.156}{-0.708}{-1.154}{-0.709}{0.388}
\SG{-1.154}{-0.709}{-1.151}{-0.711}{0.386}
\SG{-1.151}{-0.711}{-1.149}{-0.713}{0.384}
\SG{-1.149}{-0.713}{-1.147}{-0.715}{0.381}
\SG{-1.147}{-0.715}{-1.144}{-0.717}{0.379}
\SG{-1.144}{-0.717}{-1.142}{-0.718}{0.377}
\SG{-1.142}{-0.718}{-1.140}{-0.720}{0.375}
\SG{-1.140}{-0.720}{-1.137}{-0.722}{0.373}
\SG{-1.137}{-0.722}{-1.135}{-0.724}{0.372}
\SG{-1.135}{-0.724}{-1.133}{-0.725}{0.369}
\SG{-1.133}{-0.725}{-1.130}{-0.727}{0.368}
\SG{-1.130}{-0.727}{-1.128}{-0.729}{0.368}
\SG{-1.128}{-0.729}{-1.125}{-0.731}{0.367}
\SG{-1.125}{-0.731}{-1.123}{-0.732}{0.367}
\SG{-1.123}{-0.732}{-1.121}{-0.734}{0.368}
\SG{-1.121}{-0.734}{-1.118}{-0.736}{0.371}
\SG{-1.118}{-0.736}{-1.116}{-0.738}{0.374}
\SG{-1.116}{-0.738}{-1.113}{-0.739}{0.377}
\SG{-1.113}{-0.739}{-1.111}{-0.741}{0.380}
\SG{-1.111}{-0.741}{-1.108}{-0.743}{0.382}
\SG{-1.108}{-0.743}{-1.106}{-0.744}{0.384}
\SG{-1.106}{-0.744}{-1.104}{-0.746}{0.386}
\SG{-1.104}{-0.746}{-1.101}{-0.748}{0.388}
\SG{-1.101}{-0.748}{-1.099}{-0.749}{0.390}
\SG{-1.099}{-0.749}{-1.096}{-0.751}{0.391}
\SG{-1.096}{-0.751}{-1.094}{-0.752}{0.391}
\SG{-1.094}{-0.752}{-1.091}{-0.754}{0.391}
\SG{-1.091}{-0.754}{-1.089}{-0.756}{0.390}
\SG{-1.089}{-0.756}{-1.086}{-0.757}{0.388}
\SG{-1.086}{-0.757}{-1.084}{-0.759}{0.386}
\SG{-1.084}{-0.759}{-1.081}{-0.760}{0.385}
\SG{-1.081}{-0.760}{-1.079}{-0.762}{0.383}
\SG{-1.079}{-0.762}{-1.076}{-0.764}{0.381}
\SG{-1.076}{-0.764}{-1.074}{-0.765}{0.376}
\SG{-1.074}{-0.765}{-1.071}{-0.767}{0.372}
\SG{-1.071}{-0.767}{-1.069}{-0.768}{0.369}
\SG{-1.069}{-0.768}{-1.066}{-0.770}{0.367}
\SG{-1.066}{-0.770}{-1.064}{-0.771}{0.365}
\SG{-1.064}{-0.771}{-1.061}{-0.773}{0.363}
\SG{-1.061}{-0.773}{-1.059}{-0.774}{0.363}
\SG{-1.059}{-0.774}{-1.056}{-0.776}{0.363}
\SG{-1.056}{-0.776}{-1.054}{-0.777}{0.363}
\SG{-1.054}{-0.777}{-1.051}{-0.779}{0.362}
\SG{-1.051}{-0.779}{-1.048}{-0.780}{0.363}
\SG{-1.048}{-0.780}{-1.046}{-0.782}{0.363}
\SG{-1.046}{-0.782}{-1.043}{-0.783}{0.363}
\SG{-1.043}{-0.783}{-1.041}{-0.785}{0.365}
\SG{-1.041}{-0.785}{-1.038}{-0.786}{0.366}
\SG{-1.038}{-0.786}{-1.036}{-0.788}{0.367}
\SG{-1.036}{-0.788}{-1.033}{-0.789}{0.368}
\SG{-1.033}{-0.789}{-1.030}{-0.790}{0.373}
\SG{-1.030}{-0.790}{-1.028}{-0.792}{0.377}
\SG{-1.028}{-0.792}{-1.025}{-0.793}{0.380}
\SG{-1.025}{-0.793}{-1.023}{-0.795}{0.383}
\SG{-1.023}{-0.795}{-1.020}{-0.796}{0.386}
\SG{-1.020}{-0.796}{-1.017}{-0.797}{0.389}
\SG{-1.017}{-0.797}{-1.015}{-0.799}{0.390}
\SG{-1.015}{-0.799}{-1.012}{-0.800}{0.391}
\SG{-1.012}{-0.800}{-1.009}{-0.801}{0.392}
\SG{-1.009}{-0.801}{-1.007}{-0.803}{0.393}
\SG{-1.007}{-0.803}{-1.004}{-0.804}{0.394}
\SG{-1.004}{-0.804}{-1.002}{-0.805}{0.391}
\SG{-1.002}{-0.805}{-0.999}{-0.807}{0.388}
\SG{-0.999}{-0.807}{-0.996}{-0.808}{0.385}
\SG{-0.996}{-0.808}{-0.994}{-0.809}{0.383}
\SG{-0.994}{-0.809}{-0.991}{-0.811}{0.381}
\SG{-0.991}{-0.811}{-0.988}{-0.812}{0.380}
\SG{-0.988}{-0.812}{-0.986}{-0.813}{0.380}
\SG{-0.986}{-0.813}{-0.983}{-0.814}{0.379}
\SG{-0.983}{-0.814}{-0.980}{-0.816}{0.378}
\SG{-0.980}{-0.816}{-0.978}{-0.817}{0.377}
\SG{-0.978}{-0.817}{-0.975}{-0.818}{0.375}
\SG{-0.975}{-0.818}{-0.972}{-0.819}{0.368}
\SG{-0.972}{-0.819}{-0.969}{-0.821}{0.362}
\SG{-0.969}{-0.821}{-0.967}{-0.822}{0.356}
\SG{-0.967}{-0.822}{-0.964}{-0.823}{0.352}
\SG{-0.964}{-0.823}{-0.961}{-0.824}{0.348}
\SG{-0.961}{-0.824}{-0.959}{-0.825}{0.346}
\SG{-0.959}{-0.825}{-0.956}{-0.827}{0.346}
\SG{-0.956}{-0.827}{-0.953}{-0.828}{0.346}
\SG{-0.953}{-0.828}{-0.951}{-0.829}{0.346}
\SG{-0.951}{-0.829}{-0.948}{-0.830}{0.345}
\SG{-0.948}{-0.830}{-0.945}{-0.831}{0.345}
\SG{-0.945}{-0.831}{-0.942}{-0.832}{0.348}
\SG{-0.942}{-0.832}{-0.940}{-0.833}{0.350}
\SG{-0.940}{-0.833}{-0.937}{-0.835}{0.352}
\SG{-0.937}{-0.835}{-0.934}{-0.836}{0.353}
\SG{-0.934}{-0.836}{-0.931}{-0.837}{0.354}
\SG{-0.931}{-0.837}{-0.929}{-0.838}{0.355}
\SG{-0.929}{-0.838}{-0.926}{-0.839}{0.355}
\SG{-0.926}{-0.839}{-0.923}{-0.840}{0.355}
\SG{-0.923}{-0.840}{-0.920}{-0.841}{0.356}
\SG{-0.920}{-0.841}{-0.918}{-0.842}{0.356}
\SG{-0.918}{-0.842}{-0.915}{-0.843}{0.358}
\SG{-0.915}{-0.843}{-0.912}{-0.844}{0.360}
\SG{-0.912}{-0.844}{-0.909}{-0.845}{0.361}
\SG{-0.909}{-0.845}{-0.907}{-0.846}{0.363}
\SG{-0.907}{-0.846}{-0.904}{-0.847}{0.363}
\SG{-0.904}{-0.847}{-0.901}{-0.848}{0.362}
\SG{-0.901}{-0.848}{-0.898}{-0.849}{0.357}
\SG{-0.898}{-0.849}{-0.895}{-0.850}{0.353}
\SG{-0.895}{-0.850}{-0.893}{-0.851}{0.350}
\SG{-0.893}{-0.851}{-0.890}{-0.852}{0.347}
\SG{-0.890}{-0.852}{-0.887}{-0.853}{0.345}
\SG{-0.887}{-0.853}{-0.884}{-0.854}{0.345}
\SG{-0.884}{-0.854}{-0.881}{-0.855}{0.346}
\SG{-0.881}{-0.855}{-0.879}{-0.856}{0.347}
\SG{-0.879}{-0.856}{-0.876}{-0.857}{0.348}
\SG{-0.876}{-0.857}{-0.873}{-0.858}{0.349}
\SG{-0.873}{-0.858}{-0.870}{-0.858}{0.352}
\SG{-0.870}{-0.858}{-0.867}{-0.859}{0.354}
\SG{-0.867}{-0.859}{-0.864}{-0.860}{0.356}
\SG{-0.864}{-0.860}{-0.862}{-0.861}{0.358}
\SG{-0.862}{-0.861}{-0.859}{-0.862}{0.360}
\SG{-0.859}{-0.862}{-0.856}{-0.863}{0.357}
\SG{-0.856}{-0.863}{-0.853}{-0.864}{0.354}
\SG{-0.853}{-0.864}{-0.850}{-0.864}{0.352}
\SG{-0.850}{-0.864}{-0.847}{-0.865}{0.350}
\SG{-0.847}{-0.865}{-0.845}{-0.866}{0.348}
\SG{-0.845}{-0.866}{-0.842}{-0.867}{0.349}
\SG{-0.842}{-0.867}{-0.839}{-0.868}{0.350}
\SG{-0.839}{-0.868}{-0.836}{-0.868}{0.352}
\SG{-0.836}{-0.868}{-0.833}{-0.869}{0.353}
\SG{-0.833}{-0.869}{-0.830}{-0.870}{0.355}
\SG{-0.830}{-0.870}{-0.828}{-0.871}{0.359}
\SG{-0.828}{-0.871}{-0.825}{-0.871}{0.364}
\SG{-0.825}{-0.871}{-0.822}{-0.872}{0.367}
\SG{-0.822}{-0.872}{-0.819}{-0.873}{0.370}
\SG{-0.819}{-0.873}{-0.816}{-0.874}{0.373}
\SG{-0.816}{-0.874}{-0.813}{-0.874}{0.370}
\SG{-0.813}{-0.874}{-0.810}{-0.875}{0.365}
\SG{-0.810}{-0.875}{-0.807}{-0.876}{0.360}
\SG{-0.807}{-0.876}{-0.805}{-0.876}{0.355}
\SG{-0.805}{-0.876}{-0.802}{-0.877}{0.351}
\SG{-0.802}{-0.877}{-0.799}{-0.878}{0.346}
\SG{-0.799}{-0.878}{-0.796}{-0.878}{0.342}
\SG{-0.796}{-0.878}{-0.793}{-0.879}{0.338}
\SG{-0.793}{-0.879}{-0.790}{-0.880}{0.334}
\SG{-0.790}{-0.880}{-0.787}{-0.880}{0.330}
\SG{-0.787}{-0.880}{-0.784}{-0.881}{0.327}
\SG{-0.784}{-0.881}{-0.781}{-0.881}{0.326}
\SG{-0.781}{-0.881}{-0.779}{-0.882}{0.325}
\SG{-0.779}{-0.882}{-0.776}{-0.883}{0.324}
\SG{-0.776}{-0.883}{-0.773}{-0.883}{0.323}
\SG{-0.773}{-0.883}{-0.770}{-0.884}{0.325}
\SG{-0.770}{-0.884}{-0.767}{-0.884}{0.328}
\SG{-0.767}{-0.884}{-0.764}{-0.885}{0.331}
\SG{-0.764}{-0.885}{-0.761}{-0.885}{0.335}
\SG{-0.761}{-0.885}{-0.758}{-0.886}{0.338}
\SG{-0.758}{-0.886}{-0.755}{-0.886}{0.342}
\SG{-0.755}{-0.886}{-0.752}{-0.887}{0.348}
\SG{-0.752}{-0.887}{-0.750}{-0.887}{0.354}
\SG{-0.750}{-0.887}{-0.747}{-0.888}{0.359}
\SG{-0.747}{-0.888}{-0.744}{-0.888}{0.364}
\SG{-0.744}{-0.888}{-0.741}{-0.889}{0.365}
\SG{-0.741}{-0.889}{-0.738}{-0.889}{0.362}
\SG{-0.738}{-0.889}{-0.735}{-0.890}{0.359}
\SG{-0.735}{-0.890}{-0.732}{-0.890}{0.357}
\SG{-0.732}{-0.890}{-0.729}{-0.891}{0.354}
\SG{-0.729}{-0.891}{-0.726}{-0.891}{0.351}
\SG{-0.726}{-0.891}{-0.723}{-0.892}{0.348}
\SG{-0.723}{-0.892}{-0.720}{-0.892}{0.344}
\SG{-0.720}{-0.892}{-0.717}{-0.892}{0.341}
\SG{-0.717}{-0.892}{-0.714}{-0.893}{0.338}
\SG{-0.714}{-0.893}{-0.711}{-0.893}{0.334}
\SG{-0.711}{-0.893}{-0.709}{-0.893}{0.331}
\SG{-0.709}{-0.893}{-0.706}{-0.894}{0.328}
\SG{-0.706}{-0.894}{-0.703}{-0.894}{0.325}
\SG{-0.703}{-0.894}{-0.700}{-0.894}{0.322}
\SG{-0.700}{-0.894}{-0.697}{-0.895}{0.320}
\SG{-0.697}{-0.895}{-0.694}{-0.895}{0.317}
\SG{-0.694}{-0.895}{-0.691}{-0.895}{0.315}
\SG{-0.691}{-0.895}{-0.688}{-0.896}{0.312}
\SG{-0.688}{-0.896}{-0.685}{-0.896}{0.310}
\SG{-0.685}{-0.896}{-0.682}{-0.896}{0.307}
\SG{-0.682}{-0.896}{-0.679}{-0.896}{0.305}
\SG{-0.679}{-0.896}{-0.676}{-0.897}{0.303}
\SG{-0.676}{-0.897}{-0.673}{-0.897}{0.301}
\SG{-0.673}{-0.897}{-0.670}{-0.897}{0.299}
\SG{-0.670}{-0.897}{-0.667}{-0.897}{0.297}
\SG{-0.667}{-0.897}{-0.664}{-0.898}{0.295}
\SG{-0.664}{-0.898}{-0.661}{-0.898}{0.294}
\SG{-0.661}{-0.898}{-0.659}{-0.898}{0.292}
\SG{-0.659}{-0.898}{-0.656}{-0.898}{0.290}
\SG{-0.656}{-0.898}{-0.653}{-0.898}{0.289}
\SG{-0.653}{-0.898}{-0.650}{-0.899}{0.287}
\SG{-0.650}{-0.899}{-0.647}{-0.899}{0.286}
\SG{-0.647}{-0.899}{-0.644}{-0.899}{0.284}
\SG{-0.644}{-0.899}{-0.641}{-0.899}{0.283}
\SG{-0.641}{-0.899}{-0.638}{-0.899}{0.282}
\SG{-0.638}{-0.899}{-0.635}{-0.899}{0.281}
\SG{-0.635}{-0.899}{-0.632}{-0.899}{0.280}
\SG{-0.632}{-0.899}{-0.629}{-0.900}{0.278}
\SG{-0.629}{-0.900}{-0.626}{-0.900}{0.277}
\SG{-0.626}{-0.900}{-0.623}{-0.900}{0.276}
\SG{-0.623}{-0.900}{-0.620}{-0.900}{0.275}
\SG{-0.620}{-0.900}{-0.617}{-0.900}{0.274}
\SG{-0.617}{-0.900}{-0.614}{-0.900}{0.273}
\SG{-0.614}{-0.900}{-0.611}{-0.900}{0.273}
\SG{-0.611}{-0.900}{-0.608}{-0.900}{0.272}
\SG{-0.608}{-0.900}{-0.605}{-0.900}{0.271}
\SG{-0.605}{-0.900}{-0.602}{-0.900}{0.270}
\SG{-0.602}{-0.900}{-0.599}{-0.900}{0.270}
\SG{-0.599}{-0.900}{-0.597}{-0.900}{0.269}
\SG{-0.597}{-0.900}{-0.594}{-0.900}{0.268}
\SG{-0.594}{-0.900}{-0.591}{-0.900}{0.268}
\SG{-0.591}{-0.900}{-0.588}{-0.900}{0.267}
\SG{-0.588}{-0.900}{-0.585}{-0.900}{0.267}
\SG{-0.585}{-0.900}{-0.582}{-0.900}{0.267}
\SG{-0.582}{-0.900}{-0.579}{-0.900}{0.266}
\SG{-0.579}{-0.900}{-0.576}{-0.900}{0.266}
\SG{-0.576}{-0.900}{-0.573}{-0.900}{0.266}
\SG{-0.573}{-0.900}{-0.570}{-0.899}{0.265}
\SG{-0.570}{-0.899}{-0.567}{-0.899}{0.265}
\SG{-0.567}{-0.899}{-0.564}{-0.899}{0.265}
\SG{-0.564}{-0.899}{-0.561}{-0.899}{0.265}
\SG{-0.561}{-0.899}{-0.558}{-0.899}{0.265}
\SG{-0.558}{-0.899}{-0.555}{-0.899}{0.265}
\SG{-0.555}{-0.899}{-0.552}{-0.899}{0.265}
\SG{-0.552}{-0.899}{-0.549}{-0.899}{0.265}
\SG{-0.549}{-0.899}{-0.546}{-0.898}{0.265}
\SG{-0.546}{-0.898}{-0.543}{-0.898}{0.266}
\SG{-0.543}{-0.898}{-0.540}{-0.898}{0.266}
\SG{-0.540}{-0.898}{-0.537}{-0.898}{0.266}
\SG{-0.537}{-0.898}{-0.535}{-0.898}{0.267}
\SG{-0.535}{-0.898}{-0.532}{-0.897}{0.267}
\SG{-0.532}{-0.897}{-0.529}{-0.897}{0.268}
\SG{-0.529}{-0.897}{-0.526}{-0.897}{0.268}
\SG{-0.526}{-0.897}{-0.523}{-0.897}{0.269}
\SG{-0.523}{-0.897}{-0.520}{-0.896}{0.270}
\SG{-0.520}{-0.896}{-0.517}{-0.896}{0.270}
\SG{-0.517}{-0.896}{-0.514}{-0.896}{0.271}
\SG{-0.514}{-0.896}{-0.511}{-0.896}{0.272}
\SG{-0.511}{-0.896}{-0.508}{-0.895}{0.273}
\SG{-0.508}{-0.895}{-0.505}{-0.895}{0.274}
\SG{-0.505}{-0.895}{-0.502}{-0.895}{0.276}
\SG{-0.502}{-0.895}{-0.499}{-0.894}{0.277}
\SG{-0.499}{-0.894}{-0.496}{-0.894}{0.278}
\SG{-0.496}{-0.894}{-0.493}{-0.894}{0.280}
\SG{-0.493}{-0.894}{-0.490}{-0.893}{0.281}
\SG{-0.490}{-0.893}{-0.487}{-0.893}{0.283}
\SG{-0.487}{-0.893}{-0.485}{-0.893}{0.285}
\SG{-0.485}{-0.893}{-0.482}{-0.892}{0.287}
\SG{-0.482}{-0.892}{-0.479}{-0.892}{0.288}
\SG{-0.479}{-0.892}{-0.476}{-0.891}{0.290}
\SG{-0.476}{-0.891}{-0.473}{-0.891}{0.292}
\SG{-0.473}{-0.891}{-0.470}{-0.891}{0.295}
\SG{-0.470}{-0.891}{-0.467}{-0.890}{0.297}
\SG{-0.467}{-0.890}{-0.464}{-0.890}{0.298}
\SG{-0.464}{-0.890}{-0.461}{-0.889}{0.299}
\SG{-0.461}{-0.889}{-0.458}{-0.889}{0.300}
\SG{-0.458}{-0.889}{-0.455}{-0.888}{0.302}
\SG{-0.455}{-0.888}{-0.452}{-0.888}{0.303}
\SG{-0.452}{-0.888}{-0.449}{-0.887}{0.298}
\SG{-0.449}{-0.887}{-0.447}{-0.887}{0.292}
\SG{-0.447}{-0.887}{-0.444}{-0.886}{0.285}
\SG{-0.444}{-0.886}{-0.441}{-0.886}{0.278}
\SG{-0.441}{-0.886}{-0.438}{-0.885}{0.271}
\SG{-0.438}{-0.885}{-0.435}{-0.885}{0.266}
\SG{-0.435}{-0.885}{-0.432}{-0.884}{0.262}
\SG{-0.432}{-0.884}{-0.429}{-0.884}{0.257}
\SG{-0.429}{-0.884}{-0.426}{-0.883}{0.252}
\SG{-0.426}{-0.883}{-0.423}{-0.882}{0.247}
\SG{-0.423}{-0.882}{-0.420}{-0.882}{0.246}
\SG{-0.420}{-0.882}{-0.418}{-0.881}{0.245}
\SG{-0.418}{-0.881}{-0.415}{-0.881}{0.245}
\SG{-0.415}{-0.881}{-0.412}{-0.880}{0.244}
\SG{-0.412}{-0.880}{-0.409}{-0.879}{0.244}
\SG{-0.409}{-0.879}{-0.406}{-0.879}{0.246}
\SG{-0.406}{-0.879}{-0.403}{-0.878}{0.248}
\SG{-0.403}{-0.878}{-0.400}{-0.878}{0.251}
\SG{-0.400}{-0.878}{-0.397}{-0.877}{0.253}
\SG{-0.397}{-0.877}{-0.394}{-0.876}{0.256}
\SG{-0.394}{-0.876}{-0.392}{-0.876}{0.259}
\SG{-0.392}{-0.876}{-0.389}{-0.875}{0.262}
\SG{-0.389}{-0.875}{-0.386}{-0.874}{0.266}
\SG{-0.386}{-0.874}{-0.383}{-0.873}{0.269}
\SG{-0.383}{-0.873}{-0.380}{-0.873}{0.273}
\SG{-0.380}{-0.873}{-0.377}{-0.872}{0.271}
\SG{-0.377}{-0.872}{-0.374}{-0.871}{0.266}
\SG{-0.374}{-0.871}{-0.371}{-0.870}{0.262}
\SG{-0.371}{-0.870}{-0.369}{-0.870}{0.256}
\SG{-0.369}{-0.870}{-0.366}{-0.869}{0.250}
\SG{-0.366}{-0.869}{-0.363}{-0.868}{0.246}
\SG{-0.363}{-0.868}{-0.360}{-0.867}{0.242}
\SG{-0.360}{-0.867}{-0.357}{-0.867}{0.239}
\SG{-0.357}{-0.867}{-0.354}{-0.866}{0.236}
\SG{-0.354}{-0.866}{-0.352}{-0.865}{0.232}
\SG{-0.352}{-0.865}{-0.349}{-0.864}{0.231}
\SG{-0.349}{-0.864}{-0.346}{-0.863}{0.231}
\SG{-0.346}{-0.863}{-0.343}{-0.863}{0.231}
\SG{-0.343}{-0.863}{-0.340}{-0.862}{0.232}
\SG{-0.340}{-0.862}{-0.337}{-0.861}{0.234}
\SG{-0.337}{-0.861}{-0.335}{-0.860}{0.233}
\SG{-0.335}{-0.860}{-0.332}{-0.859}{0.229}
\SG{-0.332}{-0.859}{-0.329}{-0.858}{0.225}
\SG{-0.329}{-0.858}{-0.326}{-0.857}{0.221}
\SG{-0.326}{-0.857}{-0.323}{-0.856}{0.216}
\SG{-0.323}{-0.856}{-0.320}{-0.855}{0.212}
\SG{-0.320}{-0.855}{-0.318}{-0.855}{0.209}
\SG{-0.318}{-0.855}{-0.315}{-0.854}{0.206}
\SG{-0.315}{-0.854}{-0.312}{-0.853}{0.203}
\SG{-0.312}{-0.853}{-0.309}{-0.852}{0.201}
\SG{-0.309}{-0.852}{-0.306}{-0.851}{0.199}
\SG{-0.306}{-0.851}{-0.304}{-0.850}{0.199}
\SG{-0.304}{-0.850}{-0.301}{-0.849}{0.200}
\SG{-0.301}{-0.849}{-0.298}{-0.848}{0.201}
\SG{-0.298}{-0.848}{-0.295}{-0.847}{0.203}
\SG{-0.295}{-0.847}{-0.293}{-0.846}{0.206}
\SG{-0.293}{-0.846}{-0.290}{-0.845}{0.203}
\SG{-0.290}{-0.845}{-0.287}{-0.844}{0.200}
\SG{-0.287}{-0.844}{-0.284}{-0.843}{0.197}
\SG{-0.284}{-0.843}{-0.281}{-0.842}{0.194}
\SG{-0.281}{-0.842}{-0.279}{-0.841}{0.190}
\SG{-0.279}{-0.841}{-0.276}{-0.840}{0.187}
\SG{-0.276}{-0.840}{-0.273}{-0.839}{0.185}
\SG{-0.273}{-0.839}{-0.270}{-0.837}{0.183}
\SG{-0.270}{-0.837}{-0.268}{-0.836}{0.181}
\SG{-0.268}{-0.836}{-0.265}{-0.835}{0.179}
\SG{-0.265}{-0.835}{-0.262}{-0.834}{0.179}
\SG{-0.262}{-0.834}{-0.259}{-0.833}{0.181}
\SG{-0.259}{-0.833}{-0.256}{-0.832}{0.183}
\SG{-0.256}{-0.832}{-0.251}{-0.830}{0.187}
\SG{-0.251}{-0.830}{-0.246}{-0.827}{0.193}
\SG{-0.246}{-0.827}{-0.239}{-0.825}{0.189}
\SG{-0.239}{-0.825}{-0.233}{-0.822}{0.183}
\SG{-0.233}{-0.822}{-0.227}{-0.820}{0.175}
\SG{-0.227}{-0.820}{-0.220}{-0.817}{0.171}
\SG{-0.220}{-0.817}{-0.214}{-0.814}{0.167}
\SG{-0.214}{-0.814}{-0.208}{-0.812}{0.158}
\SG{-0.208}{-0.812}{-0.201}{-0.809}{0.147}
\SG{-0.201}{-0.809}{-0.195}{-0.806}{0.138}
\SG{-0.195}{-0.806}{-0.189}{-0.804}{0.132}
\SG{-0.189}{-0.804}{-0.182}{-0.801}{0.127}
\SG{-0.182}{-0.801}{-0.176}{-0.798}{0.119}
\SG{-0.176}{-0.798}{-0.170}{-0.796}{0.110}
\SG{-0.170}{-0.796}{-0.163}{-0.793}{0.106}
\SG{-0.163}{-0.793}{-0.157}{-0.791}{0.102}
\SG{-0.157}{-0.791}{-0.151}{-0.788}{0.104}
\SG{-0.151}{-0.788}{-0.144}{-0.785}{0.112}
\SG{-0.144}{-0.785}{-0.138}{-0.783}{0.118}
\SG{-0.138}{-0.783}{-0.132}{-0.780}{0.114}
\SG{-0.132}{-0.780}{-0.125}{-0.777}{0.108}
\SG{-0.125}{-0.777}{-0.119}{-0.775}{0.106}
\SG{-0.119}{-0.775}{-0.113}{-0.772}{0.104}
\SG{-0.113}{-0.772}{-0.106}{-0.770}{0.099}
\SG{-0.106}{-0.770}{-0.100}{-0.767}{0.092}
\SG{-0.100}{-0.767}{-0.093}{-0.764}{0.085}
\SG{-0.093}{-0.764}{-0.087}{-0.762}{0.081}
\SG{-0.087}{-0.762}{-0.081}{-0.759}{0.079}
\SG{-0.081}{-0.759}{-0.074}{-0.756}{0.085}
\SG{-0.074}{-0.756}{-0.068}{-0.754}{0.095}
\SG{-0.068}{-0.754}{-0.062}{-0.751}{0.099}
\SG{-0.062}{-0.751}{-0.055}{-0.749}{0.097}
\SG{-0.055}{-0.749}{-0.049}{-0.746}{0.091}
\SG{-0.049}{-0.746}{-0.043}{-0.743}{0.089}
\SG{-0.043}{-0.743}{-0.036}{-0.741}{0.089}
\SG{-0.036}{-0.741}{-0.030}{-0.738}{0.084}
\SG{-0.030}{-0.738}{-0.024}{-0.735}{0.076}
\SG{-0.024}{-0.735}{-0.017}{-0.733}{0.068}
\SG{-0.017}{-0.733}{-0.011}{-0.730}{0.064}
\SG{-0.011}{-0.730}{-0.005}{-0.727}{0.061}
\SG{-0.005}{-0.727}{0.002}{-0.725}{0.055}
\SG{0.002}{-0.725}{0.008}{-0.722}{0.048}
\SG{0.008}{-0.722}{0.014}{-0.720}{0.045}
\SG{0.014}{-0.720}{0.021}{-0.717}{0.042}
\SG{0.021}{-0.717}{0.027}{-0.714}{0.043}
\SG{0.027}{-0.714}{0.033}{-0.712}{0.051}
\SG{0.033}{-0.712}{0.040}{-0.709}{0.062}
\SG{0.040}{-0.709}{0.046}{-0.706}{0.060}
\SG{0.046}{-0.706}{0.052}{-0.704}{0.057}
\SG{0.052}{-0.704}{0.059}{-0.701}{0.054}
\SG{0.059}{-0.701}{0.065}{-0.699}{0.053}
\SG{0.065}{-0.699}{0.071}{-0.696}{0.052}
\SG{0.071}{-0.696}{0.078}{-0.693}{0.046}
\SG{0.078}{-0.693}{0.084}{-0.691}{0.039}
\SG{0.084}{-0.691}{0.091}{-0.688}{0.037}
\SG{0.091}{-0.688}{0.097}{-0.685}{0.035}
\SG{0.097}{-0.685}{0.103}{-0.683}{0.038}
\SG{0.103}{-0.683}{0.110}{-0.680}{0.047}
\SG{0.110}{-0.680}{0.116}{-0.677}{0.058}
\SG{0.116}{-0.677}{0.122}{-0.675}{0.056}
\SG{0.122}{-0.675}{0.129}{-0.672}{0.053}
\SG{0.129}{-0.672}{0.135}{-0.670}{0.050}
\SG{0.135}{-0.670}{0.141}{-0.667}{0.050}
\SG{0.141}{-0.667}{0.148}{-0.664}{0.048}
\SG{0.148}{-0.664}{0.154}{-0.662}{0.042}
\SG{0.154}{-0.662}{0.160}{-0.659}{0.034}
\SG{0.160}{-0.659}{0.167}{-0.656}{0.030}
\SG{0.167}{-0.656}{0.173}{-0.654}{0.028}
\SG{0.173}{-0.654}{0.179}{-0.651}{0.024}
\SG{0.179}{-0.651}{0.186}{-0.649}{0.018}
\SG{0.186}{-0.649}{0.192}{-0.646}{0.013}
\SG{0.192}{-0.646}{0.198}{-0.643}{0.012}
\SG{0.198}{-0.643}{0.205}{-0.641}{0.011}
\SG{0.205}{-0.641}{0.211}{-0.638}{0.017}
\SG{0.211}{-0.638}{0.217}{-0.635}{0.028}
\SG{0.217}{-0.635}{0.224}{-0.633}{0.032}
\SG{0.224}{-0.633}{0.230}{-0.630}{0.030}
\SG{0.230}{-0.630}{0.236}{-0.627}{0.027}
\SG{0.236}{-0.627}{0.243}{-0.625}{0.026}
\SG{0.243}{-0.625}{0.249}{-0.622}{0.026}
\SG{0.249}{-0.622}{0.256}{-0.620}{0.022}
\SG{0.256}{-0.620}{0.262}{-0.617}{0.016}
\SG{0.262}{-0.617}{0.268}{-0.614}{0.013}
\SG{0.268}{-0.614}{0.275}{-0.612}{0.011}
\SG{0.275}{-0.612}{0.281}{-0.609}{0.011}
\SG{0.281}{-0.609}{0.287}{-0.606}{0.019}
\SG{0.287}{-0.606}{0.294}{-0.604}{0.031}
\SG{0.294}{-0.604}{0.300}{-0.601}{0.033}
\SG{0.300}{-0.601}{0.306}{-0.599}{0.031}
\SG{0.306}{-0.599}{0.313}{-0.596}{0.027}
\SG{0.313}{-0.596}{0.319}{-0.593}{0.027}
\SG{0.319}{-0.593}{0.325}{-0.591}{0.027}
\SG{0.325}{-0.591}{0.332}{-0.588}{0.022}
\SG{0.332}{-0.588}{0.338}{-0.585}{0.016}
\SG{0.338}{-0.585}{0.344}{-0.583}{0.012}
\SG{0.344}{-0.583}{0.351}{-0.580}{0.011}
\SG{0.351}{-0.580}{0.357}{-0.577}{0.014}
\SG{0.357}{-0.577}{0.363}{-0.575}{0.024}
\SG{0.363}{-0.575}{0.370}{-0.572}{0.038}
\SG{0.370}{-0.572}{0.376}{-0.570}{0.037}
\SG{0.376}{-0.570}{0.382}{-0.567}{0.035}
\SG{0.382}{-0.567}{0.389}{-0.564}{0.032}
\SG{0.389}{-0.564}{0.395}{-0.562}{0.032}
\SG{0.395}{-0.562}{0.401}{-0.559}{0.032}
\SG{0.401}{-0.559}{0.408}{-0.556}{0.027}
\SG{0.408}{-0.556}{0.414}{-0.554}{0.020}
\SG{0.414}{-0.554}{0.421}{-0.551}{0.015}
\SG{0.421}{-0.551}{0.427}{-0.549}{0.014}
\SG{0.427}{-0.549}{0.433}{-0.546}{0.011}
\SG{0.433}{-0.546}{0.440}{-0.543}{0.006}
\SG{0.440}{-0.543}{0.446}{-0.541}{0.001}
\SG{0.446}{-0.541}{0.452}{-0.538}{0.000}
\SG{0.452}{-0.538}{0.459}{-0.535}{0.000}
\SG{0.459}{-0.535}{0.465}{-0.533}{0.006}
\SG{0.465}{-0.533}{0.471}{-0.530}{0.017}
\SG{0.471}{-0.530}{0.478}{-0.528}{0.024}
\SG{0.478}{-0.528}{0.484}{-0.525}{0.023}
\SG{0.484}{-0.525}{0.490}{-0.522}{0.020}
\SG{0.490}{-0.522}{0.497}{-0.520}{0.020}
\SG{0.497}{-0.520}{0.503}{-0.517}{0.021}
\SG{0.503}{-0.517}{0.509}{-0.514}{0.019}
\SG{0.509}{-0.514}{0.516}{-0.512}{0.014}
\SG{0.516}{-0.512}{0.522}{-0.509}{0.010}
\SG{0.522}{-0.509}{0.528}{-0.506}{0.009}
\SG{0.528}{-0.506}{0.535}{-0.504}{0.009}
\SG{0.535}{-0.504}{0.541}{-0.501}{0.018}
\SG{0.541}{-0.501}{0.547}{-0.499}{0.030}
\SG{0.547}{-0.499}{0.554}{-0.496}{0.036}
\SG{0.554}{-0.496}{0.560}{-0.493}{0.036}
\SG{0.560}{-0.493}{0.566}{-0.491}{0.032}
\SG{0.566}{-0.491}{0.573}{-0.488}{0.032}
\SG{0.573}{-0.488}{0.579}{-0.485}{0.034}
\SG{0.579}{-0.485}{0.586}{-0.483}{0.030}
\SG{0.586}{-0.483}{0.592}{-0.480}{0.025}
\SG{0.592}{-0.480}{0.598}{-0.478}{0.019}
\SG{0.598}{-0.478}{0.605}{-0.475}{0.017}
\SG{0.605}{-0.475}{0.611}{-0.472}{0.016}
\SG{0.611}{-0.472}{0.617}{-0.470}{0.012}
\SG{0.617}{-0.470}{0.624}{-0.467}{0.007}
\SG{0.624}{-0.467}{0.630}{-0.464}{0.006}
\SG{0.630}{-0.464}{0.636}{-0.462}{0.006}
\SG{0.636}{-0.462}{0.643}{-0.459}{0.009}
\SG{0.643}{-0.459}{0.649}{-0.456}{0.019}
\SG{0.649}{-0.456}{0.655}{-0.454}{0.031}
\SG{0.655}{-0.454}{0.662}{-0.451}{0.031}
\SG{0.662}{-0.451}{0.668}{-0.449}{0.030}
\SG{0.668}{-0.449}{0.674}{-0.446}{0.028}
\SG{0.674}{-0.446}{0.681}{-0.443}{0.029}
\SG{0.681}{-0.443}{0.687}{-0.441}{0.029}
\SG{0.687}{-0.441}{0.693}{-0.438}{0.026}
\SG{0.693}{-0.438}{0.700}{-0.435}{0.020}
\SG{0.700}{-0.435}{0.706}{-0.433}{0.020}
\SG{0.706}{-0.433}{0.712}{-0.430}{0.020}
\SG{0.712}{-0.430}{0.719}{-0.428}{0.026}
\SG{0.719}{-0.428}{0.725}{-0.425}{0.037}
\SG{0.725}{-0.425}{0.731}{-0.422}{0.048}
\SG{0.731}{-0.422}{0.738}{-0.420}{0.048}
\SG{0.738}{-0.420}{0.744}{-0.417}{0.046}
\SG{0.744}{-0.417}{0.751}{-0.414}{0.045}
\SG{0.751}{-0.414}{0.757}{-0.412}{0.047}
\SG{0.757}{-0.412}{0.763}{-0.409}{0.046}
\SG{0.763}{-0.409}{0.770}{-0.406}{0.042}
\SG{0.770}{-0.406}{0.776}{-0.404}{0.036}
\SG{0.776}{-0.404}{0.782}{-0.401}{0.033}
\SG{0.782}{-0.401}{0.789}{-0.399}{0.033}
\SG{0.789}{-0.399}{0.795}{-0.396}{0.030}
\SG{0.795}{-0.396}{0.801}{-0.393}{0.026}
\SG{0.801}{-0.393}{0.808}{-0.391}{0.024}
\SG{0.808}{-0.391}{0.814}{-0.388}{0.024}
\SG{0.814}{-0.388}{0.820}{-0.385}{0.025}
\SG{0.820}{-0.385}{0.827}{-0.383}{0.034}
\SG{0.827}{-0.383}{0.833}{-0.380}{0.046}
\SG{0.833}{-0.380}{0.839}{-0.378}{0.050}
\SG{0.839}{-0.378}{0.846}{-0.375}{0.050}
\SG{0.846}{-0.375}{0.852}{-0.372}{0.048}
\SG{0.852}{-0.372}{0.858}{-0.370}{0.049}
\SG{0.858}{-0.370}{0.865}{-0.367}{0.051}
\SG{0.865}{-0.367}{0.871}{-0.364}{0.048}
\SG{0.871}{-0.364}{0.877}{-0.362}{0.044}
\SG{0.877}{-0.362}{0.884}{-0.359}{0.042}
\SG{0.884}{-0.359}{0.890}{-0.356}{0.042}
\SG{0.890}{-0.356}{0.896}{-0.354}{0.045}
\SG{0.896}{-0.354}{0.903}{-0.351}{0.055}
\SG{0.903}{-0.351}{0.909}{-0.349}{0.069}
\SG{0.909}{-0.349}{0.916}{-0.346}{0.071}
\SG{0.916}{-0.346}{0.922}{-0.343}{0.070}
\SG{0.922}{-0.343}{0.928}{-0.341}{0.068}
\SG{0.928}{-0.341}{0.935}{-0.338}{0.070}
\SG{0.935}{-0.338}{0.941}{-0.335}{0.071}
\SG{0.941}{-0.335}{0.947}{-0.333}{0.068}
\SG{0.947}{-0.333}{0.954}{-0.330}{0.064}
\SG{0.954}{-0.330}{0.960}{-0.328}{0.062}
\SG{0.960}{-0.328}{0.966}{-0.325}{0.063}
\SG{0.966}{-0.325}{0.973}{-0.322}{0.068}
\SG{0.973}{-0.322}{0.979}{-0.320}{0.080}
\SG{0.979}{-0.320}{0.985}{-0.317}{0.095}
\SG{0.985}{-0.317}{0.992}{-0.314}{0.095}
\SG{0.992}{-0.314}{0.998}{-0.312}{0.094}
\SG{0.998}{-0.312}{1.004}{-0.309}{0.094}
\SG{1.004}{-0.309}{1.011}{-0.307}{0.096}
\SG{1.011}{-0.307}{1.017}{-0.304}{0.097}
\SG{1.017}{-0.304}{1.023}{-0.301}{0.094}
\SG{1.023}{-0.301}{1.030}{-0.299}{0.087}
\SG{1.030}{-0.299}{1.036}{-0.296}{0.086}
\SG{1.036}{-0.296}{1.042}{-0.293}{0.086}
\SG{1.042}{-0.293}{1.049}{-0.291}{0.085}
\SG{1.049}{-0.291}{1.055}{-0.288}{0.082}
\SG{1.055}{-0.288}{1.061}{-0.285}{0.079}
\SG{1.061}{-0.285}{1.068}{-0.283}{0.080}
\SG{1.068}{-0.283}{1.074}{-0.280}{0.082}
\SG{1.074}{-0.280}{1.081}{-0.278}{0.091}
\SG{1.081}{-0.278}{1.087}{-0.275}{0.104}
\SG{1.087}{-0.275}{1.093}{-0.272}{0.112}
\SG{1.093}{-0.272}{1.100}{-0.270}{0.112}
\SG{1.100}{-0.270}{1.106}{-0.267}{0.112}
\SG{1.106}{-0.267}{1.112}{-0.264}{0.113}
\SG{1.112}{-0.264}{1.119}{-0.262}{0.116}
\SG{1.119}{-0.262}{1.125}{-0.259}{0.116}
\SG{1.125}{-0.259}{1.131}{-0.257}{0.113}
\SG{1.131}{-0.257}{1.138}{-0.254}{0.111}
\SG{1.138}{-0.254}{1.144}{-0.251}{0.112}
\SG{1.144}{-0.251}{1.150}{-0.249}{0.115}
\SG{1.150}{-0.249}{1.157}{-0.246}{0.126}
\SG{1.157}{-0.246}{1.163}{-0.243}{0.141}
\SG{1.163}{-0.243}{1.169}{-0.241}{0.147}
\SG{1.169}{-0.241}{1.176}{-0.238}{0.149}
\SG{1.176}{-0.238}{1.182}{-0.235}{0.148}
\SG{1.182}{-0.235}{1.188}{-0.233}{0.151}
\SG{1.188}{-0.233}{1.195}{-0.230}{0.155}
\SG{1.195}{-0.230}{1.201}{-0.228}{0.152}
\SG{1.201}{-0.228}{1.207}{-0.225}{0.149}
\SG{1.207}{-0.225}{1.214}{-0.222}{0.146}
\SG{1.214}{-0.222}{1.220}{-0.220}{0.146}
\SG{1.220}{-0.220}{1.226}{-0.217}{0.147}
\SG{1.226}{-0.217}{1.233}{-0.214}{0.146}
\SG{1.233}{-0.214}{1.239}{-0.212}{0.143}
\SG{1.239}{-0.212}{1.245}{-0.209}{0.145}
\SG{1.245}{-0.209}{1.252}{-0.207}{0.147}
\SG{1.252}{-0.207}{1.258}{-0.204}{0.154}
\SG{1.258}{-0.204}{1.265}{-0.201}{0.166}
\SG{1.265}{-0.201}{1.271}{-0.199}{0.180}
\SG{1.271}{-0.199}{1.277}{-0.196}{0.183}
\SG{1.277}{-0.196}{1.284}{-0.193}{0.184}
\SG{1.284}{-0.193}{1.290}{-0.191}{0.185}
\SG{1.290}{-0.191}{1.296}{-0.188}{0.189}
\SG{1.296}{-0.188}{1.303}{-0.185}{0.191}
\SG{1.303}{-0.185}{1.309}{-0.183}{0.190}
\SG{1.309}{-0.183}{1.315}{-0.180}{0.189}
\SG{1.315}{-0.180}{1.322}{-0.178}{0.191}
\SG{1.322}{-0.178}{1.328}{-0.175}{0.194}
\SG{1.328}{-0.175}{1.334}{-0.172}{0.204}
\SG{1.334}{-0.172}{1.341}{-0.170}{0.219}
\SG{1.341}{-0.170}{1.347}{-0.167}{0.232}
\SG{1.347}{-0.167}{1.353}{-0.164}{0.235}
\SG{1.353}{-0.164}{1.360}{-0.162}{0.235}
\SG{1.360}{-0.162}{1.366}{-0.159}{0.239}
\SG{1.366}{-0.159}{1.372}{-0.157}{0.244}
\SG{1.372}{-0.157}{1.379}{-0.154}{0.246}
\SG{1.379}{-0.154}{1.385}{-0.151}{0.245}
\SG{1.385}{-0.151}{1.391}{-0.149}{0.243}
\SG{1.391}{-0.149}{1.398}{-0.146}{0.244}
\SG{1.398}{-0.146}{1.404}{-0.143}{0.247}
\SG{1.404}{-0.143}{1.410}{-0.141}{0.248}
\SG{1.410}{-0.141}{1.417}{-0.138}{0.248}
\SG{1.417}{-0.138}{1.423}{-0.135}{0.250}
\SG{1.423}{-0.135}{1.430}{-0.133}{0.254}
\SG{1.430}{-0.133}{1.436}{-0.130}{0.260}
\SG{1.436}{-0.130}{1.442}{-0.128}{0.273}
\SG{1.442}{-0.128}{1.449}{-0.125}{0.290}
\SG{1.449}{-0.125}{1.455}{-0.122}{0.298}
\SG{1.455}{-0.122}{1.461}{-0.120}{0.302}
\SG{1.461}{-0.120}{1.468}{-0.117}{0.305}
\SG{1.468}{-0.117}{1.474}{-0.114}{0.311}
\SG{1.474}{-0.114}{1.480}{-0.112}{0.317}
\SG{1.480}{-0.112}{1.487}{-0.109}{0.319}
\SG{1.487}{-0.109}{1.493}{-0.107}{0.320}
\SG{1.493}{-0.107}{1.499}{-0.104}{0.324}
\SG{1.499}{-0.104}{1.506}{-0.101}{0.330}
\SG{1.506}{-0.101}{1.512}{-0.099}{0.339}
\SG{1.512}{-0.099}{1.518}{-0.096}{0.356}
\SG{1.518}{-0.096}{1.525}{-0.093}{0.376}
\SG{1.525}{-0.093}{1.531}{-0.091}{0.383}
\SG{1.531}{-0.091}{1.537}{-0.088}{0.388}
\SG{1.537}{-0.088}{1.544}{-0.086}{0.393}
\SG{1.544}{-0.086}{1.550}{-0.083}{0.401}
\SG{1.550}{-0.083}{1.556}{-0.080}{0.410}
\SG{1.556}{-0.080}{1.563}{-0.078}{0.414}
\SG{1.563}{-0.078}{1.569}{-0.075}{0.417}
\SG{1.569}{-0.075}{1.575}{-0.072}{0.424}
\SG{1.575}{-0.072}{1.582}{-0.070}{0.433}
\SG{1.582}{-0.070}{1.588}{-0.067}{0.448}
\SG{1.588}{-0.067}{1.595}{-0.064}{0.470}
\SG{1.595}{-0.064}{1.601}{-0.062}{0.493}
\SG{1.601}{-0.062}{1.607}{-0.059}{0.503}
\SG{1.607}{-0.059}{1.614}{-0.057}{0.511}
\SG{1.614}{-0.057}{1.620}{-0.054}{0.522}
\SG{1.620}{-0.054}{1.626}{-0.051}{0.536}
\SG{1.626}{-0.051}{1.633}{-0.049}{0.548}
\SG{1.633}{-0.049}{1.639}{-0.046}{0.556}
\SG{1.639}{-0.046}{1.645}{-0.043}{0.562}
\SG{1.645}{-0.043}{1.652}{-0.041}{0.575}
\SG{1.652}{-0.041}{1.658}{-0.038}{0.589}
\SG{1.658}{-0.038}{1.664}{-0.036}{0.603}
\SG{1.664}{-0.036}{1.671}{-0.033}{0.615}
\SG{1.671}{-0.033}{1.677}{-0.030}{0.629}
\SG{1.677}{-0.030}{1.683}{-0.028}{0.648}
\SG{1.683}{-0.028}{1.690}{-0.025}{0.670}
\SG{1.690}{-0.025}{1.696}{-0.022}{0.701}
\SG{1.696}{-0.022}{1.702}{-0.020}{0.737}
\SG{1.702}{-0.020}{1.709}{-0.017}{0.764}
\SG{1.709}{-0.017}{1.715}{-0.014}{0.783}
\SG{1.715}{-0.014}{1.721}{-0.012}{0.803}
\SG{1.721}{-0.012}{1.728}{-0.009}{0.824}
\SG{1.728}{-0.009}{1.734}{-0.007}{0.850}
\SG{1.734}{-0.007}{1.740}{-0.004}{0.830}
\SG{1.740}{-0.004}{1.744}{-0.001}{0.869}
\SG{1.744}{-0.001}{1.744}{0.001}{1.000}
\SG{1.744}{0.001}{1.740}{0.004}{0.854}
\SG{1.740}{0.004}{1.734}{0.007}{0.792}
\SG{1.734}{0.007}{1.728}{0.009}{0.830}
\SG{1.728}{0.009}{1.721}{0.012}{0.850}
\SG{1.721}{0.012}{1.715}{0.014}{0.824}
\SG{1.715}{0.014}{1.709}{0.017}{0.803}
\SG{1.709}{0.017}{1.702}{0.020}{0.783}
\SG{1.702}{0.020}{1.696}{0.022}{0.764}
\SG{1.696}{0.022}{1.690}{0.025}{0.737}
\SG{1.690}{0.025}{1.683}{0.028}{0.701}
\SG{1.683}{0.028}{1.677}{0.030}{0.670}
\SG{1.677}{0.030}{1.671}{0.033}{0.648}
\SG{1.671}{0.033}{1.664}{0.036}{0.629}
\SG{1.664}{0.036}{1.658}{0.038}{0.615}
\SG{1.658}{0.038}{1.652}{0.041}{0.603}
\SG{1.652}{0.041}{1.645}{0.043}{0.589}
\SG{1.645}{0.043}{1.639}{0.046}{0.575}
\SG{1.639}{0.046}{1.633}{0.049}{0.562}
\SG{1.633}{0.049}{1.626}{0.051}{0.556}
\SG{1.626}{0.051}{1.620}{0.054}{0.548}
\SG{1.620}{0.054}{1.614}{0.057}{0.536}
\SG{1.614}{0.057}{1.607}{0.059}{0.522}
\SG{1.607}{0.059}{1.601}{0.062}{0.511}
\SG{1.601}{0.062}{1.595}{0.064}{0.503}
\SG{1.595}{0.064}{1.588}{0.067}{0.493}
\SG{1.588}{0.067}{1.582}{0.070}{0.470}
\SG{1.582}{0.070}{1.575}{0.072}{0.448}
\SG{1.575}{0.072}{1.569}{0.075}{0.433}
\SG{1.569}{0.075}{1.563}{0.078}{0.424}
\SG{1.563}{0.078}{1.556}{0.080}{0.417}
\SG{1.556}{0.080}{1.550}{0.083}{0.414}
\SG{1.550}{0.083}{1.544}{0.086}{0.410}
\SG{1.544}{0.086}{1.537}{0.088}{0.401}
\SG{1.537}{0.088}{1.531}{0.091}{0.393}
\SG{1.531}{0.091}{1.525}{0.093}{0.388}
\SG{1.525}{0.093}{1.518}{0.096}{0.383}
\SG{1.518}{0.096}{1.512}{0.099}{0.376}
\SG{1.512}{0.099}{1.506}{0.101}{0.356}
\SG{1.506}{0.101}{1.499}{0.104}{0.339}
\SG{1.499}{0.104}{1.493}{0.107}{0.330}
\SG{1.493}{0.107}{1.487}{0.109}{0.324}
\SG{1.487}{0.109}{1.480}{0.112}{0.320}
\SG{1.480}{0.112}{1.474}{0.114}{0.319}
\SG{1.474}{0.114}{1.468}{0.117}{0.317}
\SG{1.468}{0.117}{1.461}{0.120}{0.311}
\SG{1.461}{0.120}{1.455}{0.122}{0.305}
\SG{1.455}{0.122}{1.449}{0.125}{0.302}
\SG{1.449}{0.125}{1.442}{0.128}{0.298}
\SG{1.442}{0.128}{1.436}{0.130}{0.290}
\SG{1.436}{0.130}{1.430}{0.133}{0.273}
\SG{1.430}{0.133}{1.423}{0.135}{0.260}
\SG{1.423}{0.135}{1.417}{0.138}{0.254}
\SG{1.417}{0.138}{1.410}{0.141}{0.250}
\SG{1.410}{0.141}{1.404}{0.143}{0.248}
\SG{1.404}{0.143}{1.398}{0.146}{0.248}
\SG{1.398}{0.146}{1.391}{0.149}{0.247}
\SG{1.391}{0.149}{1.385}{0.151}{0.244}
\SG{1.385}{0.151}{1.379}{0.154}{0.243}
\SG{1.379}{0.154}{1.372}{0.157}{0.245}
\SG{1.372}{0.157}{1.366}{0.159}{0.246}
\SG{1.366}{0.159}{1.360}{0.162}{0.244}
\SG{1.360}{0.162}{1.353}{0.164}{0.239}
\SG{1.353}{0.164}{1.347}{0.167}{0.235}
\SG{1.347}{0.167}{1.341}{0.170}{0.235}
\SG{1.341}{0.170}{1.334}{0.172}{0.232}
\SG{1.334}{0.172}{1.328}{0.175}{0.219}
\SG{1.328}{0.175}{1.322}{0.178}{0.204}
\SG{1.322}{0.178}{1.315}{0.180}{0.194}
\SG{1.315}{0.180}{1.309}{0.183}{0.191}
\SG{1.309}{0.183}{1.303}{0.185}{0.189}
\SG{1.303}{0.185}{1.296}{0.188}{0.190}
\SG{1.296}{0.188}{1.290}{0.191}{0.191}
\SG{1.290}{0.191}{1.284}{0.193}{0.189}
\SG{1.284}{0.193}{1.277}{0.196}{0.185}
\SG{1.277}{0.196}{1.271}{0.199}{0.184}
\SG{1.271}{0.199}{1.265}{0.201}{0.183}
\SG{1.265}{0.201}{1.258}{0.204}{0.180}
\SG{1.258}{0.204}{1.252}{0.207}{0.166}
\SG{1.252}{0.207}{1.245}{0.209}{0.154}
\SG{1.245}{0.209}{1.239}{0.212}{0.147}
\SG{1.239}{0.212}{1.233}{0.214}{0.145}
\SG{1.233}{0.214}{1.226}{0.217}{0.143}
\SG{1.226}{0.217}{1.220}{0.220}{0.146}
\SG{1.220}{0.220}{1.214}{0.222}{0.147}
\SG{1.214}{0.222}{1.207}{0.225}{0.146}
\SG{1.207}{0.225}{1.201}{0.228}{0.146}
\SG{1.201}{0.228}{1.195}{0.230}{0.149}
\SG{1.195}{0.230}{1.188}{0.233}{0.152}
\SG{1.188}{0.233}{1.182}{0.235}{0.155}
\SG{1.182}{0.235}{1.176}{0.238}{0.151}
\SG{1.176}{0.238}{1.169}{0.241}{0.148}
\SG{1.169}{0.241}{1.163}{0.243}{0.149}
\SG{1.163}{0.243}{1.157}{0.246}{0.147}
\SG{1.157}{0.246}{1.150}{0.249}{0.141}
\SG{1.150}{0.249}{1.144}{0.251}{0.126}
\SG{1.144}{0.251}{1.138}{0.254}{0.115}
\SG{1.138}{0.254}{1.131}{0.257}{0.112}
\SG{1.131}{0.257}{1.125}{0.259}{0.111}
\SG{1.125}{0.259}{1.119}{0.262}{0.113}
\SG{1.119}{0.262}{1.112}{0.264}{0.116}
\SG{1.112}{0.264}{1.106}{0.267}{0.116}
\SG{1.106}{0.267}{1.100}{0.270}{0.113}
\SG{1.100}{0.270}{1.093}{0.272}{0.112}
\SG{1.093}{0.272}{1.087}{0.275}{0.112}
\SG{1.087}{0.275}{1.081}{0.278}{0.112}
\SG{1.081}{0.278}{1.074}{0.280}{0.104}
\SG{1.074}{0.280}{1.068}{0.283}{0.091}
\SG{1.068}{0.283}{1.061}{0.285}{0.082}
\SG{1.061}{0.285}{1.055}{0.288}{0.080}
\SG{1.055}{0.288}{1.049}{0.291}{0.079}
\SG{1.049}{0.291}{1.042}{0.293}{0.082}
\SG{1.042}{0.293}{1.036}{0.296}{0.085}
\SG{1.036}{0.296}{1.030}{0.299}{0.086}
\SG{1.030}{0.299}{1.023}{0.301}{0.086}
\SG{1.023}{0.301}{1.017}{0.304}{0.087}
\SG{1.017}{0.304}{1.011}{0.307}{0.094}
\SG{1.011}{0.307}{1.004}{0.309}{0.097}
\SG{1.004}{0.309}{0.998}{0.312}{0.096}
\SG{0.998}{0.312}{0.992}{0.314}{0.094}
\SG{0.992}{0.314}{0.985}{0.317}{0.094}
\SG{0.985}{0.317}{0.979}{0.320}{0.095}
\SG{0.979}{0.320}{0.973}{0.322}{0.095}
\SG{0.973}{0.322}{0.966}{0.325}{0.080}
\SG{0.966}{0.325}{0.960}{0.328}{0.068}
\SG{0.960}{0.328}{0.954}{0.330}{0.063}
\SG{0.954}{0.330}{0.947}{0.333}{0.062}
\SG{0.947}{0.333}{0.941}{0.335}{0.064}
\SG{0.941}{0.335}{0.935}{0.338}{0.068}
\SG{0.935}{0.338}{0.928}{0.341}{0.071}
\SG{0.928}{0.341}{0.922}{0.343}{0.070}
\SG{0.922}{0.343}{0.916}{0.346}{0.068}
\SG{0.916}{0.346}{0.909}{0.349}{0.070}
\SG{0.909}{0.349}{0.903}{0.351}{0.071}
\SG{0.903}{0.351}{0.896}{0.354}{0.069}
\SG{0.896}{0.354}{0.890}{0.356}{0.055}
\SG{0.890}{0.356}{0.884}{0.359}{0.045}
\SG{0.884}{0.359}{0.877}{0.362}{0.042}
\SG{0.877}{0.362}{0.871}{0.364}{0.042}
\SG{0.871}{0.364}{0.865}{0.367}{0.044}
\SG{0.865}{0.367}{0.858}{0.370}{0.048}
\SG{0.858}{0.370}{0.852}{0.372}{0.051}
\SG{0.852}{0.372}{0.846}{0.375}{0.049}
\SG{0.846}{0.375}{0.839}{0.378}{0.048}
\SG{0.839}{0.378}{0.833}{0.380}{0.050}
\SG{0.833}{0.380}{0.827}{0.383}{0.050}
\SG{0.827}{0.383}{0.820}{0.385}{0.046}
\SG{0.820}{0.385}{0.814}{0.388}{0.034}
\SG{0.814}{0.388}{0.808}{0.391}{0.025}
\SG{0.808}{0.391}{0.801}{0.393}{0.024}
\SG{0.801}{0.393}{0.795}{0.396}{0.024}
\SG{0.795}{0.396}{0.789}{0.399}{0.026}
\SG{0.789}{0.399}{0.782}{0.401}{0.030}
\SG{0.782}{0.401}{0.776}{0.404}{0.033}
\SG{0.776}{0.404}{0.770}{0.406}{0.033}
\SG{0.770}{0.406}{0.763}{0.409}{0.036}
\SG{0.763}{0.409}{0.757}{0.412}{0.042}
\SG{0.757}{0.412}{0.751}{0.414}{0.046}
\SG{0.751}{0.414}{0.744}{0.417}{0.047}
\SG{0.744}{0.417}{0.738}{0.420}{0.045}
\SG{0.738}{0.420}{0.731}{0.422}{0.046}
\SG{0.731}{0.422}{0.725}{0.425}{0.048}
\SG{0.725}{0.425}{0.719}{0.428}{0.048}
\SG{0.719}{0.428}{0.712}{0.430}{0.037}
\SG{0.712}{0.430}{0.706}{0.433}{0.026}
\SG{0.706}{0.433}{0.700}{0.435}{0.020}
\SG{0.700}{0.435}{0.693}{0.438}{0.020}
\SG{0.693}{0.438}{0.687}{0.441}{0.020}
\SG{0.687}{0.441}{0.681}{0.443}{0.026}
\SG{0.681}{0.443}{0.674}{0.446}{0.029}
\SG{0.674}{0.446}{0.668}{0.449}{0.029}
\SG{0.668}{0.449}{0.662}{0.451}{0.028}
\SG{0.662}{0.451}{0.655}{0.454}{0.030}
\SG{0.655}{0.454}{0.649}{0.456}{0.031}
\SG{0.649}{0.456}{0.643}{0.459}{0.031}
\SG{0.643}{0.459}{0.636}{0.462}{0.019}
\SG{0.636}{0.462}{0.630}{0.464}{0.009}
\SG{0.630}{0.464}{0.624}{0.467}{0.006}
\SG{0.624}{0.467}{0.617}{0.470}{0.006}
\SG{0.617}{0.470}{0.611}{0.472}{0.007}
\SG{0.611}{0.472}{0.605}{0.475}{0.012}
\SG{0.605}{0.475}{0.598}{0.478}{0.016}
\SG{0.598}{0.478}{0.592}{0.480}{0.017}
\SG{0.592}{0.480}{0.586}{0.483}{0.019}
\SG{0.586}{0.483}{0.579}{0.485}{0.025}
\SG{0.579}{0.485}{0.573}{0.488}{0.030}
\SG{0.573}{0.488}{0.566}{0.491}{0.034}
\SG{0.566}{0.491}{0.560}{0.493}{0.032}
\SG{0.560}{0.493}{0.554}{0.496}{0.032}
\SG{0.554}{0.496}{0.547}{0.499}{0.036}
\SG{0.547}{0.499}{0.541}{0.501}{0.036}
\SG{0.541}{0.501}{0.535}{0.504}{0.030}
\SG{0.535}{0.504}{0.528}{0.506}{0.018}
\SG{0.528}{0.506}{0.522}{0.509}{0.009}
\SG{0.522}{0.509}{0.516}{0.512}{0.009}
\SG{0.516}{0.512}{0.509}{0.514}{0.010}
\SG{0.509}{0.514}{0.503}{0.517}{0.014}
\SG{0.503}{0.517}{0.497}{0.520}{0.019}
\SG{0.497}{0.520}{0.490}{0.522}{0.021}
\SG{0.490}{0.522}{0.484}{0.525}{0.020}
\SG{0.484}{0.525}{0.478}{0.528}{0.020}
\SG{0.478}{0.528}{0.471}{0.530}{0.023}
\SG{0.471}{0.530}{0.465}{0.533}{0.024}
\SG{0.465}{0.533}{0.459}{0.535}{0.017}
\SG{0.459}{0.535}{0.452}{0.538}{0.006}
\SG{0.452}{0.538}{0.446}{0.541}{0.000}
\SG{0.446}{0.541}{0.440}{0.543}{0.000}
\SG{0.440}{0.543}{0.433}{0.546}{0.001}
\SG{0.433}{0.546}{0.427}{0.549}{0.006}
\SG{0.427}{0.549}{0.421}{0.551}{0.011}
\SG{0.421}{0.551}{0.414}{0.554}{0.014}
\SG{0.414}{0.554}{0.408}{0.556}{0.015}
\SG{0.408}{0.556}{0.401}{0.559}{0.020}
\SG{0.401}{0.559}{0.395}{0.562}{0.027}
\SG{0.395}{0.562}{0.389}{0.564}{0.032}
\SG{0.389}{0.564}{0.382}{0.567}{0.032}
\SG{0.382}{0.567}{0.376}{0.570}{0.032}
\SG{0.376}{0.570}{0.370}{0.572}{0.035}
\SG{0.370}{0.572}{0.363}{0.575}{0.037}
\SG{0.363}{0.575}{0.357}{0.577}{0.038}
\SG{0.357}{0.577}{0.351}{0.580}{0.024}
\SG{0.351}{0.580}{0.344}{0.583}{0.014}
\SG{0.344}{0.583}{0.338}{0.585}{0.011}
\SG{0.338}{0.585}{0.332}{0.588}{0.012}
\SG{0.332}{0.588}{0.325}{0.591}{0.016}
\SG{0.325}{0.591}{0.319}{0.593}{0.022}
\SG{0.319}{0.593}{0.313}{0.596}{0.027}
\SG{0.313}{0.596}{0.306}{0.599}{0.027}
\SG{0.306}{0.599}{0.300}{0.601}{0.027}
\SG{0.300}{0.601}{0.294}{0.604}{0.031}
\SG{0.294}{0.604}{0.287}{0.606}{0.033}
\SG{0.287}{0.606}{0.281}{0.609}{0.031}
\SG{0.281}{0.609}{0.275}{0.612}{0.019}
\SG{0.275}{0.612}{0.268}{0.614}{0.011}
\SG{0.268}{0.614}{0.262}{0.617}{0.011}
\SG{0.262}{0.617}{0.256}{0.620}{0.013}
\SG{0.256}{0.620}{0.249}{0.622}{0.016}
\SG{0.249}{0.622}{0.243}{0.625}{0.022}
\SG{0.243}{0.625}{0.236}{0.627}{0.026}
\SG{0.236}{0.627}{0.230}{0.630}{0.026}
\SG{0.230}{0.630}{0.224}{0.633}{0.027}
\SG{0.224}{0.633}{0.217}{0.635}{0.030}
\SG{0.217}{0.635}{0.211}{0.638}{0.032}
\SG{0.211}{0.638}{0.205}{0.641}{0.028}
\SG{0.205}{0.641}{0.198}{0.643}{0.017}
\SG{0.198}{0.643}{0.192}{0.646}{0.011}
\SG{0.192}{0.646}{0.186}{0.649}{0.012}
\SG{0.186}{0.649}{0.179}{0.651}{0.013}
\SG{0.179}{0.651}{0.173}{0.654}{0.018}
\SG{0.173}{0.654}{0.167}{0.656}{0.024}
\SG{0.167}{0.656}{0.160}{0.659}{0.028}
\SG{0.160}{0.659}{0.154}{0.662}{0.030}
\SG{0.154}{0.662}{0.148}{0.664}{0.034}
\SG{0.148}{0.664}{0.141}{0.667}{0.042}
\SG{0.141}{0.667}{0.135}{0.670}{0.048}
\SG{0.135}{0.670}{0.129}{0.672}{0.050}
\SG{0.129}{0.672}{0.122}{0.675}{0.050}
\SG{0.122}{0.675}{0.116}{0.677}{0.053}
\SG{0.116}{0.677}{0.110}{0.680}{0.056}
\SG{0.110}{0.680}{0.103}{0.683}{0.058}
\SG{0.103}{0.683}{0.097}{0.685}{0.047}
\SG{0.097}{0.685}{0.091}{0.688}{0.038}
\SG{0.091}{0.688}{0.084}{0.691}{0.035}
\SG{0.084}{0.691}{0.078}{0.693}{0.037}
\SG{0.078}{0.693}{0.071}{0.696}{0.039}
\SG{0.071}{0.696}{0.065}{0.699}{0.046}
\SG{0.065}{0.699}{0.059}{0.701}{0.052}
\SG{0.059}{0.701}{0.052}{0.704}{0.053}
\SG{0.052}{0.704}{0.046}{0.706}{0.054}
\SG{0.046}{0.706}{0.040}{0.709}{0.057}
\SG{0.040}{0.709}{0.033}{0.712}{0.060}
\SG{0.033}{0.712}{0.027}{0.714}{0.062}
\SG{0.027}{0.714}{0.021}{0.717}{0.051}
\SG{0.021}{0.717}{0.014}{0.720}{0.043}
\SG{0.014}{0.720}{0.008}{0.722}{0.042}
\SG{0.008}{0.722}{0.002}{0.725}{0.045}
\SG{0.002}{0.725}{-0.005}{0.727}{0.048}
\SG{-0.005}{0.727}{-0.011}{0.730}{0.055}
\SG{-0.011}{0.730}{-0.017}{0.733}{0.061}
\SG{-0.017}{0.733}{-0.024}{0.735}{0.064}
\SG{-0.024}{0.735}{-0.030}{0.738}{0.068}
\SG{-0.030}{0.738}{-0.036}{0.741}{0.076}
\SG{-0.036}{0.741}{-0.043}{0.743}{0.084}
\SG{-0.043}{0.743}{-0.049}{0.746}{0.089}
\SG{-0.049}{0.746}{-0.055}{0.749}{0.089}
\SG{-0.055}{0.749}{-0.062}{0.751}{0.091}
\SG{-0.062}{0.751}{-0.068}{0.754}{0.097}
\SG{-0.068}{0.754}{-0.074}{0.756}{0.099}
\SG{-0.074}{0.756}{-0.081}{0.759}{0.095}
\SG{-0.081}{0.759}{-0.087}{0.762}{0.085}
\SG{-0.087}{0.762}{-0.093}{0.764}{0.079}
\SG{-0.093}{0.764}{-0.100}{0.767}{0.081}
\SG{-0.100}{0.767}{-0.106}{0.770}{0.085}
\SG{-0.106}{0.770}{-0.113}{0.772}{0.092}
\SG{-0.113}{0.772}{-0.119}{0.775}{0.099}
\SG{-0.119}{0.775}{-0.125}{0.777}{0.104}
\SG{-0.125}{0.777}{-0.132}{0.780}{0.106}
\SG{-0.132}{0.780}{-0.138}{0.783}{0.108}
\SG{-0.138}{0.783}{-0.144}{0.785}{0.114}
\SG{-0.144}{0.785}{-0.151}{0.788}{0.118}
\SG{-0.151}{0.788}{-0.157}{0.791}{0.112}
\SG{-0.157}{0.791}{-0.163}{0.793}{0.104}
\SG{-0.163}{0.793}{-0.170}{0.796}{0.102}
\SG{-0.170}{0.796}{-0.176}{0.798}{0.106}
\SG{-0.176}{0.798}{-0.182}{0.801}{0.110}
\SG{-0.182}{0.801}{-0.189}{0.804}{0.119}
\SG{-0.189}{0.804}{-0.195}{0.806}{0.127}
\SG{-0.195}{0.806}{-0.201}{0.809}{0.132}
\SG{-0.201}{0.809}{-0.208}{0.812}{0.138}
\SG{-0.208}{0.812}{-0.214}{0.814}{0.147}
\SG{-0.214}{0.814}{-0.220}{0.817}{0.158}
\SG{-0.220}{0.817}{-0.227}{0.820}{0.167}
\SG{-0.227}{0.820}{-0.233}{0.822}{0.171}
\SG{-0.233}{0.822}{-0.239}{0.825}{0.175}
\SG{-0.239}{0.825}{-0.246}{0.827}{0.183}
\SG{-0.246}{0.827}{-0.251}{0.830}{0.189}
\SG{-0.251}{0.830}{-0.256}{0.832}{0.193}
\SG{-0.256}{0.832}{-0.259}{0.833}{0.187}

    \node[font=\footnotesize, text=azulprof] at (-0.6,0) {$\phi=\phi_0$};
    \draw[azulclaro, line width=0.5pt] (-2.6,-2.0) rectangle (2.6,2.0);
    % barra de color de sigma
    \foreach \k in {0,...,19}{\mapcolor{\k/19}\fill[mapc] ({-2.0+0.2*\k},-2.35) rectangle ++(0.2,0.13);}
    \node[font=\scriptsize, anchor=east] at (-2.05,-2.29) {$\sigma$ menor};
    \node[font=\scriptsize, anchor=west] at (2.05,-2.29) {mayor};
    \node at (0,-2.85) {(a)};
  \end{scope}
  % ---------------------------------------------------------------- (b)
  \begin{scope}[xshift=7.15cm, yshift=-2.15cm]
    \begin{axis}[tema, width=6.6cm, height=5.3cm,
        xlabel={$x$}, ylabel={$\sigma/\sigma_{\min}$},
        xmin=-1.6, xmax=1.85, ymin=0.8, ymax=3.7,
        xtick={-1.5,-1,-0.5,0,0.5,1,1.5}, ytick={1,1.5,2,2.5,3,3.5}]
      \addplot[azulprof, line width=1pt] table[x=x, y=sigma] {figuras/datos/punta_sigma.dat};
      \draw[punteado] (axis cs:1.75,0.8) -- (axis cs:1.75,3.7);
      \node[font=\scriptsize, anchor=east, text=grispizarra] at (axis cs:1.68,3.45) {punta $\rightarrow$};
      \node[font=\scriptsize, anchor=west, text=grispizarra, fill=white, inner sep=1pt] at (axis cs:-1.5,2.2) {cuerpo redondeado};
    \end{axis}
    \node at (2.4,-1.15) {(b)};
  \end{scope}
\end{tikzpicture}
\caption{Efecto punta en un conductor cargado con forma de pera (sección de un conductor cilíndrico indefinido, resuelta numéricamente con
$\Lap\phi=0$). (a) Líneas de campo trazadas con el mismo flujo entre dos consecutivas, equipotenciales (trazos) y superficie coloreada según
$\sigma$: el campo y la carga se concentran en la punta. (b) Densidad superficial a lo largo del contorno: es mínima en la zona cóncava de
unión y crece bruscamente al acercarse a la punta.}
\label{fig:punta}
\end{figure}


\begin{observacion}[El pararrayos y las descargas en punta]\label{obs:pararrayos}
El efecto punta es el fundamento del \emph{pararrayos} de Franklin. Como el campo en la superficie es $E=\sigma/\eo$, una punta afilada conectada a
tierra alcanza campos enormes con potenciales moderados: el aire a su alrededor se ioniza (descarga en corona) y se crea un camino conductor
que guía la descarga de la nube hacia tierra a través del pararrayos, y no a través del edificio. El mismo razonamiento explica por qué los
conductores con aristas o puntas pierden carga antes, y por qué en los equipos de alta tensión todas las superficies se redondean
(\cref{sem:globo}).
\end{observacion}

% =========================================================================
\section{Presión electrostática}\label{sec:presion}
% =========================================================================
Si el conductor está cargado, su carga reside en la superficie y cada elemento de carga es repelido por todos los demás. Aparece así una
fuerza por unidad de área que tiende a expulsar la carga hacia el exterior: la \emph{presión electrostática}. Calcularla exige un poco de
cuidado, porque el campo en la superficie es discontinuo: vale $\sigma/\eo$ justo fuera y $0$ justo dentro. ¿Cuál de los dos actúa sobre la carga?

% -------------------------------------------------------------------------
\subsection{Un argumento ingenuo y su fallo}\label{sec:presion-ingenuo}
% -------------------------------------------------------------------------
Una primera idea sería multiplicar la carga de un elemento $\Delta S$ por el campo exterior:
\[
  \Pel\overset{?}{=}\frac{F}{\Delta S}=\frac{\Delta q\,E}{\Delta S}=\frac{\sigma\,\Delta S\,(\sigma/\eo)}{\Delta S}=\frac{\sigma^{2}}{\eo}.
\]
El razonamiento tiene un fallo: el campo $\sigma/\eo$ incluye el que crea la propia carga $\Delta q$, y \emph{una distribución de carga no ejerce
fuerza neta sobre sí misma} (por la tercera ley de Newton, las fuerzas internas se cancelan dos a dos), igual que nadie puede levantarse a sí
mismo tirando de sus propias manos hacia arriba. Hay que separar el campo en la contribución del elemento y la del resto de la carga.

% -------------------------------------------------------------------------
\subsection{El campo que actúa sobre la carga superficial}\label{sec:presion-campo}
% -------------------------------------------------------------------------
% =========================================================================
%  Figura: presión electrostática. (a) Elemento DeltaS de la superficie, campo
%  del disco y campo del resto en dos puntos 1 (fuera) y 2 (dentro).
%  (b) Componente normal de cada contribución frente a la coordenada xi
%  normal a la superficie: E_disco salta de -sigma/2eps0 a +sigma/2eps0,
%  E_resto es continuo y la suma pasa de 0 a sigma/eps0.
% =========================================================================
\begin{figure}[t]
\centering
\begin{tikzpicture}[font=\small]
  % ---------------------------------------------------------------- (a)
  \begin{scope}
    % superficie curvada del conductor (vista en perspectiva)
    \fill[top color=azulvelo!70!white, bottom color=metal!80!white, draw=azulprof, line width=0.6pt]
      (-3.0,-0.3) .. controls (-1.6,0.55) and (0.9,0.65) .. (3.0,0.05)
      -- (2.25,-1.05) .. controls (0.3,-0.55) and (-1.9,-0.75) .. (-3.6,-1.2) -- cycle;
    \node[font=\footnotesize, text=azulprof] at (-2.55,-0.95) {$S'$};
    % disco DeltaS
    \fill[cargapos!25] (0,-0.12) ellipse (0.62 and 0.17);
    \draw[cargapos, line width=0.6pt] (0,-0.12) ellipse (0.62 and 0.17);
    \node[font=\footnotesize] at (0.92,-0.32) {$\Delta S$};
    % campo del disco (a ambos lados) y del resto
    \draw[-{Stealth[length=4pt]}, verdeazul, line width=0.9pt] (0.12,-0.05) -- (0.12,1.15);
    \draw[-{Stealth[length=4pt]}, verdeazul, line width=0.9pt, densely dashed] (0.12,-0.2) -- (0.12,-1.4);
    \node[font=\footnotesize, text=verdeazul] at (0.72,0.95) {$\Edisc$};
    \draw[-{Stealth[length=4pt]}, cargapos, line width=0.9pt] (-0.22,-0.9) -- (-0.22,1.55);
    \node[font=\footnotesize, text=cargapos] at (-0.78,1.42) {$\Erest$};
    \draw[-{Stealth[length=3.5pt]}, azulprof, line width=0.6pt] (1.75,0.35) -- (1.75,0.95) node[right, font=\footnotesize] {$\un$};
    % puntos 1 y 2
    \begin{scope}[shift={(2.85,-1.95)}]
      \draw[azulacero, line width=1.3pt] (-0.55,0) -- (0.55,0);
      \fill[azulprof] (0,0.22) circle (1.4pt) node[right=3pt, font=\footnotesize] {1 (fuera)};
      \fill[azulprof] (0,-0.22) circle (1.4pt) node[right=3pt, font=\footnotesize] {2 (dentro)};
      \node[font=\footnotesize, anchor=east] at (-0.6,0) {$\Delta S$};
    \end{scope}
    \draw[->, grismedio, line width=0.4pt] (0.65,-0.18) .. controls (1.6,-0.7) .. (2.25,-1.85);
    \node at (0,-2.95) {(a)};
  \end{scope}
  % ---------------------------------------------------------------- (b)
  \begin{scope}[xshift=6.3cm, yshift=-2.1cm]
    \begin{axis}[tema, width=7.0cm, height=5.4cm, axis lines=left,
        xlabel={$\xi$ (normal a la superficie)}, ylabel={$E_n$}, xmin=-2.1, xmax=2.1, ymin=-0.75, ymax=1.25,
        xtick={0}, xticklabels={$0$}, ytick={-0.5,0,0.5,1},
        yticklabels={$-\frac{\sigma}{2\eo}$,$0$,$\frac{\sigma}{2\eo}$,$\frac{\sigma}{\eo}$},
        grid=none, legend columns=1,
        legend style={at={(0.5,1.04)}, anchor=south, font=\scriptsize}]
      \addplot[verdeazul, line width=1.1pt, densely dashed, const plot] coordinates {(-2,-0.5) (0,0.5) (2,0.5)};
      \addlegendentry{campo del disco $E_{\mathrm{disco}}$}
      \addplot[cargapos, line width=1.1pt] coordinates {(-2,0.52) (2,0.52)};
      \addlegendentry{campo del resto $E_{\mathrm{resto}}$}
      \addplot[azulprof, line width=1.4pt, const plot] coordinates {(-2,0.0) (0,1) (2,1)};
      \addlegendentry{campo total $E=E_{\mathrm{disco}}+E_{\mathrm{resto}}$}
      \draw[punteado] (axis cs:0,-0.75) -- (axis cs:0,1.25);
      \node[font=\scriptsize, text=grispizarra] at (axis cs:-1.05,-0.63) {conductor};
      \node[font=\scriptsize, text=grispizarra] at (axis cs:1.05,-0.63) {vacío};
    \end{axis}
    \node at (2.6,-1.25) {(b)};
  \end{scope}
\end{tikzpicture}
\caption{Presión electrostática. (a) El campo en un elemento $\Delta S$ de la superficie es la suma del campo del propio disco, discontinuo
($\pm\sigma/2\eo$ a cada lado), y del campo $\Erest$ del resto de la carga, continuo. (b) Componente normal de cada contribución frente a la
coordenada $\xi$ perpendicular a la superficie: $\Erest$ duplica el campo fuera y lo anula dentro, y es el único que actúa sobre la carga de
$\Delta S$.}
\label{fig:presion}
\end{figure}


\begin{lema}[Campo del resto de la carga]\label{lem:resto}
Sea $\Delta S$ un elemento de la superficie de un conductor en equilibrio, tan pequeño que $\sigma$ es constante en él, y sea $\Erest$ el campo
creado en $\Delta S$ por toda la carga del sistema \emph{excepto} la de $\Delta S$. Entonces
\begin{equation}\label{eq:campo-resto}
  \resalta{\Erest=\frac12\bigl(\vE_1+\vE_2\bigr)=\frac{\sigma}{2\eo}\,\un ,}
\end{equation}
donde $\vE_1$ y $\vE_2$ son los campos totales justo fuera y justo dentro del conductor, respectivamente.
\end{lema}
\begin{proof}
Descomponemos el campo total en dos puntos muy próximos a $\Delta S$, uno fuera (1) y otro dentro (2), como suma del campo del disco $\Delta S$ y
del campo del resto (\figpanel{fig:presion}{a}): $\vE=\Edisc+\Erest$.

\emph{El campo del resto es continuo} en $\Delta S$: es el campo que crea en un punto del vacío una distribución de carga (la de $S'-\Delta S$) que
no tiene carga en ese punto, de modo que no presenta ninguna discontinuidad allí; toma el mismo valor $\Erest$ en 1 y en 2.

\emph{El campo del disco es discontinuo.} En puntos suficientemente próximos a $\Delta S$, el disco se ve como un plano indefinido con densidad
$\sigma$, cuyo campo es $\sigma/(2\eo)$, perpendicular al plano y dirigido hacia fuera de él a ambos lados~\cite{TEMA1}:
\[
  \vE_{\mathrm{disco},1}=+\frac{\sigma}{2\eo}\,\un,\qquad\vE_{\mathrm{disco},2}=-\frac{\sigma}{2\eo}\,\un .
\]
Por tanto, usando \cref{eq:campo-superficie} y \cref{eq:campo-nulo},
\begin{alignat*}{2}
  \vE_1&=\Erest+\frac{\sigma}{2\eo}\,\un=\frac{\sigma}{\eo}\,\un &\qquad&\text{(punto 1, fuera)},\\
  \vE_2&=\Erest-\frac{\sigma}{2\eo}\,\un=\vc{0} &&\text{(punto 2, dentro)}.
\end{alignat*}
Sumando miembro a miembro ambas igualdades, $\vE_1+\vE_2=2\,\Erest$, de donde
\[
  \Erest=\frac12\bigl(\vE_1+\vE_2\bigr)=\frac12\Bigl(\frac{\sigma}{\eo}\,\un+\vc{0}\Bigr)=\frac{\sigma}{2\eo}\,\un .
\]
\end{proof}

\begin{observacion}[Lectura del resultado]\label{obs:resto}
El campo del resto de la carga vale exactamente la mitad del campo exterior, y es el \emph{promedio} de los campos a ambos lados de la superficie.
Su papel es doble (\figpanel{fig:presion}{b}): fuera, se suma al del disco y \emph{duplica} la intensidad hasta $\sigma/\eo$; dentro, se opone a él y lo
\emph{cancela} por completo. El equilibrio del conductor es, desde este punto de vista, un ajuste fino entre la carga de cada elemento y la de todos
los demás.
\end{observacion}

% -------------------------------------------------------------------------
\subsection{La presión electrostática}\label{sec:presion-valor}
% -------------------------------------------------------------------------
\begin{teorema}[Presión electrostática]\label{teo:presion}
Sobre la superficie de un conductor en equilibrio actúa una fuerza por unidad de área
\begin{equation}\label{eq:presion}
  \resalta{\Pel=\frac{\dif F}{\dif S}=\frac{\sigma^{2}}{2\eo}=\frac{\eo}{2}\,E^{2},}\qquad
  \frac{\dif\vF}{\dif S}=\frac{\sigma^{2}}{2\eo}\,\un ,
\end{equation}
dirigida siempre hacia el exterior del conductor, cualquiera que sea el signo de $\sigma$; aquí $E=\abs{\sigma}/\eo$ es el campo justo fuera de la superficie.
\end{teorema}
\begin{proof}
La fuerza sobre la carga $\Delta q=\sigma\,\Delta S$ del elemento la ejerce solo el resto de la carga, de modo que, por el \cref{lem:resto},
\[
  \Delta\vF=\Delta q\,\Erest=\sigma\,\Delta S\,\frac{\sigma}{2\eo}\,\un
  \quad\Longrightarrow\quad\Pel=\frac{\Delta F}{\Delta S}=\frac{\sigma^{2}}{2\eo} .
\]
Como $\sigma^{2}\ge0$, la fuerza tiene siempre el sentido de $\un$, hacia fuera. Por último, $\sigma^{2}/(2\eo)=\tfrac{\eo}{2}(\sigma/\eo)^{2}=\tfrac{\eo}{2}E^{2}$.
\end{proof}

\begin{observacion}[Presión y densidad de energía]\label{obs:presion-energia}
La expresión $\Pel=\tfrac{\eo}{2}E^{2}$ coincide con la densidad de energía del campo electrostático, $u=\tfrac{\eo}{2}E^{2}$, que estudiaremos más
adelante. No es casual: si un elemento de superficie se desplaza hacia fuera una distancia $\dif\xi$, el volumen $\dif S\,\dif\xi$, en el que antes había
un campo $E$, pasa a formar parte del metal y queda sin campo; el trabajo de la fuerza electrostática, $\Pel\,\dif S\,\dif\xi$, es igual a la energía
$u\,\dif S\,\dif\xi$ que desaparece.
\end{observacion}

\begin{ejemplo}[Fuerza entre los dos hemisferios de una esfera cargada]\label{ej:hemisferios}
Una esfera conductora de radio $R$ tiene carga $Q$. ¿Con qué fuerza se repelen sus dos hemisferios?

\emph{Resolución.} La carga se reparte uniformemente, $\sigma=Q/(4\pi R^{2})$, y la presión $\Pel=\sigma^{2}/(2\eo)$ es la misma en toda la superficie. Por
simetría, la fuerza sobre el hemisferio norte apunta según $\uz$, y solo contribuye la componente $\un\cdot\uz=\cos\theta$:
\begin{align*}
  F_z&=\iint_{\text{hemisferio}}\Pel\cos\theta\,\dsc
     =\frac{\sigma^{2}}{2\eo}\int_{0}^{2\pi}\!\!\int_{0}^{\pi/2}\cos\theta\,R^{2}\sen\theta\,\dif\theta\,\dif\varphi\\
     &=\frac{\sigma^{2}}{2\eo}\,2\pi R^{2}\,\frac12=\frac{\sigma^{2}\pi R^{2}}{2\eo}=\frac{Q^{2}}{32\pi\eo R^{2}} .
\end{align*}
El resultado tiene una lectura geométrica: es la presión por el área de la \emph{proyección} del hemisferio sobre el plano ecuatorial, $\pi R^{2}$.
\end{ejemplo}

\begin{observacion}[Órdenes de magnitud]\label{obs:presion-orden}
El campo máximo que soporta el aire seco antes de ionizarse es $\Erup\simeq3\cdot10^{6}\,\mathrm{V/m}$ (\cref{sem:globo}); a ese campo,
$\Pel=\tfrac{\eo}{2}\Erup^{2}\simeq40\,\mathrm{Pa}$, unas $2500$ veces menor que la presión atmosférica. La presión electrostática es pequeña,
pero perceptible en membranas ligeras: una pompa de jabón cargada se hincha ligeramente.
\end{observacion}

% =========================================================================
\section{Ecuación de Laplace, cavidades y capacidad}\label{sec:cavidades}
% =========================================================================
Las propiedades anteriores son locales. Para tratar conductores con cavidades, o sistemas de varios conductores, necesitamos una herramienta
global: la formulación del problema electrostático como un problema de contorno para el potencial, y los teoremas que garantizan que su
solución es única.

% -------------------------------------------------------------------------
\subsection{Ecuaciones de Poisson y Laplace}\label{sec:poisson}
% -------------------------------------------------------------------------
\begin{proposicion}[Ecuaciones de Poisson y Laplace]\label{prop:poisson}
En electrostática, el potencial satisface la \emph{ecuación de Poisson}
\begin{equation}\label{eq:poisson}
  \resalta{\Lap\phi=-\frac{\rho}{\eo},}
\end{equation}
y, en las regiones sin carga ($\rho=0$), la \emph{ecuación de Laplace}, $\Lap\phi=0$. Las funciones de clase $C^{2}$ que satisfacen la ecuación de
Laplace en un abierto se llaman \emph{armónicas}.
\end{proposicion}
\begin{proof}
Basta sustituir $\vE=-\nab\phi$ en la ley de Gauss diferencial, $\nab\cdot\vE=\rho/\eo$:
$\nab\cdot(-\nab\phi)=-\Lap\phi=\rho/\eo$~\cite{TEMA0,TEMA1}.
\end{proof}

En presencia de conductores, el espacio vacío $\Omega=\R^{3}\setminus\bigcup_i\Vc_i$ es una región en la que (si no hay cargas libres en el vacío)
$\phi$ es armónica, y sobre cuya frontera, las superficies $\Sc_i$ de los conductores, el potencial es constante. El problema electrostático se
convierte en un \emph{problema de contorno}. Dos propiedades de las funciones armónicas son las que hacen tratable este problema.

\begin{lema}[Principio del máximo]\label{lem:maximo}
Sea $\phi$ armónica en una región acotada $\Omega$ y continua en su cierre. Entonces $\phi$ alcanza su máximo y su mínimo sobre la frontera
$\partial\Omega$, y no tiene extremos locales en el interior salvo que sea constante. Si $\Omega$ es exterior a unos conductores, el resultado se
mantiene contando el infinito como parte de la frontera.
\end{lema}
\begin{proof}
Las funciones armónicas cumplen la \emph{propiedad del valor medio}: $\phi(\vc{r}_0)$ es igual al promedio de $\phi$ sobre cualquier esfera centrada
en $\vc{r}_0$ y contenida en $\Omega$~\cite{GRIFFITHS,KELLOGG}. Si $\phi$ tuviese un máximo local estricto en $\vc{r}_0$, en una esfera pequeña alrededor de
$\vc{r}_0$ todos los valores serían menores que $\phi(\vc{r}_0)$ y su promedio también, lo que es absurdo. El caso de un máximo no estricto se reduce
al anterior por un argumento de conexión~\cite{KELLOGG}.
\end{proof}

\begin{teorema}[Teoremas de unicidad]\label{teo:unicidad}
Sea $\Omega$ la región exterior a un conjunto de conductores $\Vc_1,\dots,\Vc_N$, con una densidad de carga $\rho$ dada en $\Omega$ y la condición
$\phi\to0$ en el infinito. La solución de la ecuación de Poisson en $\Omega$ es única si se conoce, para cada conductor:
\begin{enumerate}[label=\textup{(\alph*)}]
  \item su potencial $\phi_i$ (\emph{primer teorema de unicidad}, problema de Dirichlet), o bien
  \item su carga total $Q_i$ (\emph{segundo teorema de unicidad}),
\end{enumerate}
o una combinación de ambos datos (el potencial para unos conductores y la carga para los demás).
\end{teorema}
\begin{proof}
Sean $\phi_a$ y $\phi_b$ dos soluciones y $\psi=\phi_a-\phi_b$. Entonces $\Lap\psi=0$ en $\Omega$, $\psi\to0$ en el infinito y $\psi$ es constante, de
valor $\psi_i$, sobre cada $\Sc_i$. Aplicamos el teorema de la divergencia al campo $\psi\nab\psi$, cuya divergencia es
$\nab\cdot(\psi\nab\psi)=\norm{\nab\psi}^{2}+\psi\Lap\psi=\norm{\nab\psi}^{2}$, en la región comprendida entre los conductores y una esfera de radio
$R\to\infty$ (en la que el integrando decae como $R^{-3}$ y su contribución se anula):
\[
  \iiint_{\Omega}\norm{\nab\psi}^{2}\dV=-\sum_{i=1}^{N}\oiint_{\Sc_i}\psi\,\nab\psi\cdot\un\,\dsc
  =-\sum_{i=1}^{N}\psi_i\oiint_{\Sc_i}\nab\psi\cdot\un\,\dsc ,
\]
donde $\un$ es la normal exterior a cada conductor; el signo menos aparece porque la normal exterior a la región $\Omega$ sobre $\Sc_i$ es
$-\un$. Ahora bien: (a) si los potenciales coinciden, $\psi_i=0$; (b) si las cargas coinciden, como $Q_i=\eo\oiint_{\Sc_i}\vE\cdot\un\,\dsc
=-\eo\oiint_{\Sc_i}\nab\phi\cdot\un\,\dsc$ por \cref{eq:campo-superficie}, se tiene
$\oiint_{\Sc_i}\nab\psi\cdot\un\,\dsc=-\frac{1}{\eo}\bigl(Q_i^{(a)}-Q_i^{(b)}\bigr)=0$. En cualquiera de los dos casos, y en las combinaciones, cada sumando se
anula, de modo que $\int_{\Omega}\norm{\nab\psi}^{2}=0$, $\nab\psi=\vc{0}$ y $\psi$ es constante; como $\psi\to0$ en el infinito, $\psi\equiv0$.
\end{proof}

\begin{observacion}[El significado práctico de la unicidad]\label{obs:unicidad}
La unicidad legitima el método de «adivinar y comprobar»: cualquier función que cumpla la ecuación y las condiciones de contorno es \emph{la}
solución. Es la base del método de las imágenes, que usaremos en la \cref{fig:corteza}, y del principio de superposición para conductores,
que da lugar a los coeficientes de potencial (\cref{sec:sistemas}).
\end{observacion}

% -------------------------------------------------------------------------
\subsection{Conductores con cavidades}\label{sec:cavidades-casos}
% -------------------------------------------------------------------------
Consideramos un conductor con una cavidad interior (una región vacía completamente rodeada de metal). Su superficie tiene dos partes: la
\emph{cara exterior} y la \emph{cara interior}, que limita la cavidad.

\begin{teorema}[Cavidad vacía: jaula de Faraday]\label{teo:faraday}
Si una cavidad de un conductor en equilibrio no contiene carga, el campo en su interior es nulo y la cara interior no tiene carga en ningún
punto, cualesquiera que sean las cargas exteriores y la carga del conductor:
\begin{equation}\label{eq:faraday}
  \resalta{\Ecav=\vc{0},\qquad\sigma\big|_{\text{cara interior}}=0 .}
\end{equation}
\end{teorema}
\begin{proof}
En la cavidad $\phi$ es armónica, y en su frontera (la cara interior) vale la constante $\phi_{\Vc}$, el potencial del conductor. La función
constante $\phi\equiv\phi_{\Vc}$ cumple la ecuación y la condición de contorno; por el primer teorema de unicidad (o, directamente, por el principio
del máximo: el máximo y el mínimo en la cavidad se alcanzan en la frontera, donde valen lo mismo) es la única solución. Por tanto
$\Ecav=-\nab\phi=\vc{0}$, y por \cref{eq:campo-superficie} $\sigma=\eo\,\vE\cdot\un=0$ en la cara interior.

\emph{Argumento alternativo con líneas de campo.} Si hubiese campo en la cavidad, sus líneas tendrían que nacer y morir en la cara interior
(no hay cargas dentro de la cavidad), es decir, saldrían del conductor y volverían a él, lo que prohíbe el \cref{cor:lineas}.
\end{proof}

\begin{proposicion}[Regla de las caras]\label{prop:caras}
Sea un conductor de carga total $Q$ con una cavidad que contiene una carga total $q$. Entonces, en el equilibrio,
\begin{equation}\label{eq:caras}
  \resalta{\Qint=-q,\qquad\Qext=Q+q .}
\end{equation}
\end{proposicion}
\begin{proof}
Tomamos una superficie gaussiana $S_G$ contenida por completo en el metal y que rodee la cavidad (curvas a trazos de la \cref{fig:cavidades}).
Sobre ella $\vE=\vc{0}$, luego $\Qenc=\eo\oiint_{S_G}\vE\cdot\dS=0$. La carga encerrada es la de la cavidad más la de la cara interior, $q+\Qint=0$, de donde
$\Qint=-q$. Como el metal no tiene carga en volumen (\cref{teo:propiedades}), la carga del conductor se reparte entre sus dos caras,
$Q=\Qint+\Qext$, y $\Qext=Q-\Qint=Q+q$.
\end{proof}

% =========================================================================
%  Figura: los cuatro casos de conductor hueco (como en las figuras del tema).
%  Contornos en polares (centro común):
%    cara exterior  r(t) = 1.50 + 0.16 cos 3t + 0.10 sen 2t
%    gaussiana      r(t) = 1.06 + 0.08 cos 3t + 0.05 sen 2t
%    cara interior  r(t) = 0.62 + 0.08 cos(2t + 40)
% =========================================================================
\newcommand{\cavExt}{plot[domain=0:360, samples=90, smooth cycle, variable=\t]
  ({(1.50+0.16*cos(3*\t)+0.10*sin(2*\t))*cos(\t)},{(1.50+0.16*cos(3*\t)+0.10*sin(2*\t))*sin(\t)})}
\newcommand{\cavInt}{plot[domain=0:360, samples=70, smooth cycle, variable=\t]
  ({(0.62+0.08*cos(2*\t+40))*cos(\t)},{(0.62+0.08*cos(2*\t+40))*sin(\t)})}
\newcommand{\cavGauss}{plot[domain=0:360, samples=80, smooth cycle, variable=\t]
  ({(1.06+0.08*cos(3*\t)+0.05*sin(2*\t))*cos(\t)},{(1.06+0.08*cos(3*\t)+0.05*sin(2*\t))*sin(\t)})}
% signos sobre la cara exterior (#1 = 1) o interior (#1 = 2), en el lado del metal
\newcommand{\signosExt}{\foreach \t in {0,30,...,330}{
  \node[signopos, font=\tiny\bfseries] at ({(1.36+0.16*cos(3*\t)+0.10*sin(2*\t))*cos(\t)},{(1.36+0.16*cos(3*\t)+0.10*sin(2*\t))*sin(\t)}) {$+$};}}
\newcommand{\signosInt}{\foreach \t in {0,36,...,324}{
  \node[signoneg, font=\tiny\bfseries] at ({(0.76+0.08*cos(2*\t+40))*cos(\t)},{(0.76+0.08*cos(2*\t+40))*sin(\t)}) {$-$};}}
\newcommand{\panelCavidad}[4]{% #1 título, #2 signos ext (0/1), #3 carga en cavidad (0/1), #4 rótulos
  \fill[metal!55!white, even odd rule] \cavExt \cavInt;
  \fill[pattern=north east lines, pattern color=azulcielo!70, even odd rule] \cavExt \cavInt;
  \draw[azulprof, line width=0.6pt] \cavExt;
  \draw[azulprof, line width=0.6pt] \cavInt;
  \draw[gauss] \cavGauss;
  \ifnum#2=1 \signosExt\fi
  \ifnum#3=1 \signosInt \node[qpos, minimum size=8pt] at (0.05,-0.02) {$+$};
             \node[font=\scriptsize, text=cargapos] at (0.05,0.3) {$q$};\fi
  \node[font=\scriptsize, text=azulprof, fill=white, inner sep=1pt, rounded corners=1pt] at (0.15,-1.25) {$\vE=\vc{0}$};
  \node[anchor=west, align=left, font=\footnotesize] at (1.95,0) {#4};
  \node[anchor=south, font=\small] at (1.15,1.85) {#1};}

\begin{figure}[t]
\centering
\begin{tikzpicture}[x=0.95cm, y=0.95cm]
  \begin{scope}
    \panelCavidad{(a) Neutro, cavidad vacía}{0}{0}
      {$\Qint=0$\\ $\Qext=0$\\[2pt] $\Ecav=\vc{0}$\\ $\Eext=\vc{0}$}
  \end{scope}
  \begin{scope}[xshift=8.2cm]
    \panelCavidad{(b) Carga $Q$, cavidad vacía}{1}{0}
      {$\Qint=0$\\ $\Qext=Q$\\[2pt] $\Ecav=\vc{0}$\\ $\Eext\neq\vc{0}$}
  \end{scope}
  \begin{scope}[yshift=-4.4cm]
    \panelCavidad{(c) Neutro, carga $q$ en la cavidad}{1}{1}
      {$\Qint=-q$\\ $\Qext=+q$\\[2pt] $\Ecav\neq\vc{0}$\\ $\Eext\neq\vc{0}$}
  \end{scope}
  \begin{scope}[xshift=8.2cm, yshift=-4.4cm]
    \panelCavidad{(d) Carga $Q$, carga $q$ en la cavidad}{1}{1}
      {$\Qint=-q$\\ $\Qext=Q+q$\\[2pt] $\Ecav\neq\vc{0}$\\ $\Eext\neq\vc{0}$}
  \end{scope}
  % leyenda común
  \draw[gauss] (2.6,-7.05) -- (3.4,-7.05);
  \node[anchor=west, font=\footnotesize] at (3.5,-7.05) {superficie de Gauss (cerrada) dentro del metal};
\end{tikzpicture}
\caption{Conductores huecos en equilibrio. La superficie de Gauss trazada dentro del metal, donde $\vE=\vc{0}$, encierra carga nula: la cara
interior lleva siempre la carga opuesta a la de la cavidad, y la exterior el resto de la carga del conductor. Una cavidad vacía no tiene campo.}
\label{fig:cavidades}
\end{figure}


La \cref{tab:cavidades} resume los cuatro casos posibles, que la \cref{fig:cavidades} representa.

\begin{table}[ht]
\centering
\caption{Carga de cada cara y campo en la cavidad y en el exterior de un conductor hueco aislado, sin más cargas que las indicadas.}
\label{tab:cavidades}
\begin{tabular}{@{}lcccc@{}}
  \toprule
  Situación & Cara interior & Cara exterior & $\Ecav$ & $\Eext$\\
  \midrule
  Neutro, cavidad vacía                 & $0$  & $0$    & $\vc{0}$        & $\vc{0}$\\
  Carga $Q$, cavidad vacía              & $0$  & $Q$    & $\vc{0}$        & $\neq\vc{0}$\\
  Neutro, carga $q$ en la cavidad       & $-q$ & $+q$   & $\neq\vc{0}$    & $\neq\vc{0}$\\
  Carga $Q$, carga $q$ en la cavidad    & $-q$ & $Q+q$  & $\neq\vc{0}$    & $\neq\vc{0}$\\
  \bottomrule
\end{tabular}
\end{table}

\begin{teorema}[Apantallamiento en los dos sentidos]\label{teo:apantallamiento}
Sea un conductor hueco cerrado, con cargas en su cavidad y en el exterior.
\begin{enumerate}[label=\textup{(\alph*)}]
  \item El campo en la cavidad depende solo de las cargas de la cavidad y de la forma de la cara interior: no se ve afectado por las cargas
        exteriores (salvo por una constante aditiva en el potencial).
  \item El campo exterior depende solo de las cargas exteriores, de la forma de la cara exterior y de la carga total $\Qext=Q+q$ de esta cara: no
        depende de cómo estén colocadas las cargas dentro de la cavidad.
\end{enumerate}
\end{teorema}
\begin{proof}
(a) En la cavidad, $\phi-\phi_{\Vc}$ cumple la ecuación de Poisson con las cargas de la cavidad y se anula sobre la cara interior: por el primer
teorema de unicidad está determinada sin referencia alguna al exterior. (b) En el exterior, $\phi$ cumple la ecuación de Poisson con las cargas
exteriores, tiende a cero en el infinito, es constante sobre la cara exterior y la carga de esa cara es $\Qext=Q+q$ (\cref{prop:caras}): por el segundo
teorema de unicidad el problema exterior queda determinado, y en sus datos solo interviene el valor $q$, no la posición de la carga.
\end{proof}

La \cref{fig:corteza} ilustra el teorema con una corteza esférica gruesa, que puede resolverse exactamente por el método de las imágenes~\cite{JACKSON}.
Con una carga $q$ en el centro (a) todo tiene simetría esférica. Al desplazar la carga (b), las líneas de la cavidad se deforman, pero fuera de la
corteza el campo sigue siendo exactamente el de una carga $q$ en el centro: la cara exterior «no se entera». Al añadir una carga exterior (c), el
campo exterior se deforma, pero el de la cavidad es idéntico al de (b). Por último, sin carga en la cavidad (d), el interior queda completamente
apantallado: es la jaula de Faraday.

% =========================================================================
%  Figura «tipo Mathematica»: corteza esférica conductora gruesa (radios
%  R1 = 1 de la cavidad y R2 = 2 exterior), neutra, en el plano meridiano.
%  Solución exacta por imágenes:
%   - cavidad: q en s u_x  +  imagen -q R1/s en (R1^2/s) u_x;
%   - exterior: Qe en -D u_x, imagen -Qe R2/D en -(R2^2/D) u_x y la carga
%     q + Qe R2/D en el centro (la corteza aislada tiene Q_ext = q).
%  Datos: figuras/datos/corteza_{a,b,c,d}_*.
% =========================================================================
\newcommand{\panelCorteza}[2]{% #1 = a, b, c, d ; #2 = título
  \fill[azulfondo!60!white] (-5.4,-3.3) rectangle (3.4,3.3);
  \begin{scope}
    \clip (-5.4,-3.3) rectangle (3.4,3.3);
    \input{figuras/datos/corteza_#1_lc}
    \input{figuras/datos/corteza_#1_fl}
  \end{scope}
  \fill[metal!70!white, even odd rule] (0,0) circle (2) (0,0) circle (1);
  \fill[top color=azulvelo!30, bottom color=metalosc!45, opacity=0.6, even odd rule] (0,0) circle (2) (0,0) circle (1);
  \draw[azulprof, line width=0.55pt] (0,0) circle (2) (0,0) circle (1);
  \input{figuras/datos/corteza_#1_q}
  \draw[azulclaro, line width=0.5pt] (-5.4,-3.3) rectangle (3.4,3.3);
  \node[anchor=north, font=\small] at (-1.0,-3.4) {#2};}

\begin{figure}[t]
\centering
\begin{tikzpicture}[x=0.74cm, y=0.74cm]
  \tikzset{numLC/.style={azulmedio, line width=0.42pt},
           numFL/.style={-{Stealth[length=3pt, width=2.5pt]}, azulmedio, line width=0.42pt}}
  \begin{scope}
    \panelCorteza{a}{(a) Carga $q$ en el centro}
    \node[qpos, minimum size=7pt, font=\tiny\bfseries] at (0,0) {$+$};
  \end{scope}
  \begin{scope}[xshift=8.9cm]
    \panelCorteza{b}{(b) Carga $q$ desplazada}
    \node[qpos, minimum size=7pt, font=\tiny\bfseries] at (0.45,0) {$+$};
  \end{scope}
  \begin{scope}[yshift=-6.0cm]
    \panelCorteza{c}{(c) Carga $q$ desplazada y carga exterior $Q_e$}
    \node[qpos, minimum size=7pt, font=\tiny\bfseries] at (0.45,0) {$+$};
    \node[qpos, minimum size=8pt, font=\tiny\bfseries] at (-4,0) {$+$};
    \node[fondo, font=\scriptsize, text=cargapos] at (-4,0.52) {$Q_e$};
  \end{scope}
  \begin{scope}[xshift=8.9cm, yshift=-6.0cm]
    \panelCorteza{d}{(d) Cavidad vacía y carga exterior $Q_e$}
    \node[qpos, minimum size=8pt, font=\tiny\bfseries] at (-4,0) {$+$};
    \node[fondo, font=\scriptsize, text=cargapos] at (-4,0.52) {$Q_e$};
    \node[font=\scriptsize, text=azulprof, fill=white, inner sep=1pt, rounded corners=1pt] at (0,0) {$\Ecav=\vc{0}$};
  \end{scope}
\end{tikzpicture}
\caption{Corteza esférica conductora neutra (radios interior $R_1$ y exterior $R_2=2R_1$), resuelta exactamente por el método de las imágenes.
(a)~Carga $q$ en el centro: simetría esférica. (b)~Al desplazar $q$ cambia el campo de la cavidad, pero el exterior sigue siendo el de una carga
$q$ en el centro. (c)~Una carga exterior $Q_e$ deforma el campo exterior, pero no altera el de la cavidad. (d)~Sin carga en la cavidad, el interior
queda apantallado (jaula de Faraday). Se indican los signos de la carga inducida en cada cara.}
\label{fig:corteza}
\end{figure}


\begin{observacion}[Aplicaciones del apantallamiento]\label{obs:jaula}
Una caja metálica cerrada protege su interior de cualquier campo electrostático exterior: es el principio de las jaulas de Faraday, de los cables
coaxiales con malla exterior y de la protección de los ocupantes de un coche o un avión alcanzado por un rayo. En la práctica, una malla con aberturas
pequeñas frente a la escala de variación del campo funciona casi igual que una superficie continua. Recíprocamente, conectando el conductor a tierra
($\phi_{\Vc}=0$, $\Qext=0$) también se impide que las cargas de la cavidad creen campo fuera.
\end{observacion}

% -------------------------------------------------------------------------
\subsection{Capacidad de un conductor aislado}\label{sec:capacidad}
% -------------------------------------------------------------------------
Sea un conductor aislado, cargado con carga $Q$, en equilibrio. El potencial que crea en un punto exterior es, por la ley de Coulomb para
distribuciones superficiales~\cite{TEMA1},
\begin{equation}\label{eq:pot-conductor}
  \phi(\vc{r})=\frac{1}{4\pi\eo}\oiint_{\Sc}\frac{\sigma(\rp)}{\norm{\vc{r}-\rp}}\,\dSp ,
\end{equation}
y, si $\vc{r}\in\Sc$, se obtiene el potencial del propio conductor, $\phi_{\Vc}$.

\begin{teorema}[Capacidad]\label{teo:capacidad}
Para un conductor aislado, el cociente entre su carga y su potencial no depende de $Q$:
\begin{equation}\label{eq:capacidad}
  \resalta{C=\frac{Q}{\phi_{\Vc}}=\frac{Q'}{\phi'_{\Vc}}=\mathrm{cte},}
\end{equation}
y solo depende de la geometría del conductor (y del medio). Se llama \emph{capacidad} del conductor; su unidad en el SI es el
\emph{faradio}, $1\,\mathrm{F}=1\,\mathrm{C/V}$.
\end{teorema}
\begin{proof}
Sea $\sigma$ la densidad de equilibrio cuando la carga es $Q$, con potencial $\phi_{\Vc}$. Si la carga pasa a ser $Q'=\lambda Q$, la densidad
$\lambda\sigma$ tiene carga total $\lambda Q$ y, por la linealidad de \cref{eq:pot-conductor}, crea un potencial $\lambda\phi$, que es constante
($\lambda\phi_{\Vc}$) sobre el conductor y nulo en el infinito. Por el segundo teorema de unicidad es la solución de equilibrio: la geometría no cambia,
luego la forma en que se distribuye la carga tampoco, y solo se reescala. Así $\phi'_{\Vc}=\lambda\phi_{\Vc}$ y $Q'/\phi'_{\Vc}=Q/\phi_{\Vc}$.
\end{proof}

\begin{ejemplo}[Esfera conductora aislada]\label{ej:esfera}
Para una esfera de radio $R$, la carga se reparte uniformemente y el potencial es el de una carga puntual en el centro evaluado en $r=R$,
$\phi=\kc Q/R$. Así
\begin{equation}\label{eq:cap-esfera}
  C_{\text{esfera}}=4\pi\eo R .
\end{equation}
El faradio es una unidad enorme: la propia Tierra, tratada como una esfera conductora de radio $6{,}37\cdot10^{6}\,\mathrm{m}$, tiene una capacidad de solo
$C\simeq7{,}1\cdot10^{-4}\,\mathrm{F}\simeq709\,\mu\mathrm{F}$. En la práctica se usan sus submúltiplos: $1\,\mu\mathrm{F}=10^{-6}\,\mathrm{F}$,
$1\,\mathrm{nF}=10^{-9}\,\mathrm{F}$ y $1\,\mathrm{pF}=10^{-12}\,\mathrm{F}$.
\end{ejemplo}

% =========================================================================
\section{Sistemas de conductores}\label{sec:sistemas}
% =========================================================================
Si en lugar de un conductor tenemos un sistema de $N$ conductores finitos situados a distancias finitas, cada uno influye en todos los demás: la
carga de uno induce cargas en los otros, que a su vez modifican su potencial. Lo que nos interesa es la relación entre los potenciales
$\phi_i$ de los conductores y sus cargas $Q_i$. Veremos que es \emph{lineal} y que sus coeficientes, introducidos por Maxwell en su
\emph{Tratado}~\cite{MAXWELL}, dependen solo de la geometría.

% =========================================================================
%  Figura: geometría de un sistema de conductores. Punto fuente r'_j sobre el
%  conductor j, punto r_i sobre el conductor i y punto de observación r;
%  R_j = r - r'_j y R_ij = r_i - r'_j.
% =========================================================================
\begin{figure}[t]
\centering
\begin{tikzpicture}[font=\small, x=1cm, y=1cm]
  % ejes
  \draw[eje] (0,0) -- (0,1.6) node[above, font=\footnotesize] {$z$};
  \draw[eje] (0,0) -- (1.3,0) node[right, font=\footnotesize] {$y$};
  \draw[eje] (0,0) -- (-0.75,-0.6) node[below left=-2pt, font=\footnotesize] {$x$};
  \node[below=1pt, font=\footnotesize] at (0,0) {$O$};
  % conductor j
  \fill[top color=azulvelo!80!white, bottom color=metal, draw=azulprof, line width=0.6pt]
    (1.0,2.0) .. controls (0.8,2.9) and (1.9,3.5) .. (2.5,2.9)
    .. controls (3.1,2.4) and (2.6,1.4) .. (1.9,1.35) .. controls (1.4,1.3) and (1.1,1.5) .. cycle;
  \node[font=\footnotesize, text=azulprof] at (1.1,3.35) {$S_j'$};
  \node[font=\footnotesize] at (2.25,2.55) {$Q_j$};
  \node[font=\footnotesize, text=cargapos] at (3.55,3.35) {$\sigma_j(\rp_j)$};
  % conductor i
  \fill[top color=azulvelo!80!white, bottom color=metal, draw=azulprof, line width=0.6pt]
    (6.1,0.2) .. controls (5.7,0.9) and (6.6,1.4) .. (7.3,1.0)
    .. controls (7.9,0.6) and (7.6,-0.7) .. (7.0,-0.9) .. controls (6.5,-1.05) and (6.4,-0.4) .. cycle;
  \node[font=\footnotesize, text=azulprof] at (7.95,0.95) {$S_i'$};
  \node[font=\footnotesize] at (7.0,-0.35) {$Q_i$};
  % tercer conductor
  \fill[top color=azulvelo!80!white, bottom color=metal, draw=azulprof, line width=0.6pt]
    (7.2,3.0) .. controls (7.0,3.6) and (7.8,3.9) .. (8.2,3.5) .. controls (8.6,3.1) and (8.1,2.5) .. (7.7,2.6) .. controls (7.4,2.65) .. cycle;
  \node[font=\footnotesize] at (7.75,3.15) {$Q_N$};
  % vectores
  \coordinate (rpj) at (1.95,1.36);
  \coordinate (ri) at (6.25,0.45);
  \coordinate (r) at (5.3,2.4);
  \draw[vect, azulacero] (0,0) -- (rpj) node[midway, left=1pt, font=\footnotesize] {$\rp_j$};
  \draw[vect, azulacero] (0,0) -- (ri) node[pos=0.6, below=1pt, font=\footnotesize] {$\vc{r}_i$};
  \draw[vect, azulacero] (0,0) -- (r) node[pos=0.62, above=1pt, font=\footnotesize] {$\vc{r}$};
  \draw[vectrojo] (rpj) -- (r) node[pos=0.62, above=1pt, font=\footnotesize] {$\vc{R}_j$};
  \draw[vectrojo] (rpj) -- (ri) node[pos=0.72, below left=0pt, font=\footnotesize] {$\vc{R}_{ij}$};
  \fill[azulprof] (rpj) circle (1.3pt) (ri) circle (1.3pt) (r) circle (1.3pt);
\end{tikzpicture}
\caption{Sistema de conductores finitos a distancias finitas. El conductor $j$ tiene carga $Q_j$ repartida con densidad $\sigma_j$ sobre su superficie
$S_j'$; $\vc{R}_j=\vc{r}-\rp_j$ va de un punto fuente al punto de observación, y $\vc{R}_{ij}=\vc{r}_i-\rp_j$, de un punto fuente de $j$ a un punto del
conductor $i$.}
\label{fig:sistema-geom}
\end{figure}


% -------------------------------------------------------------------------
\subsection{Coeficientes de potencial}\label{sec:pij}
% -------------------------------------------------------------------------
\paso{La idea subyacente.} Supongamos que solo el conductor $j$ está cargado, $Q_j\neq0$ y $Q_k=0$ para $k\neq j$. Aparecen densidades de carga
$\sigma(\rp)$ en todos los conductores (inducidas en los neutros, con carga neta nula), tales que $\vE=\vc{0}$ en el interior de cada uno, y con ellas un
potencial $\phi(\vc{r})$ y un campo $\vE(\vc{r})$ en todo el espacio. Nos interesan los potenciales de los conductores, $\phi_i^{(j)}$, que son constantes.
Si multiplicamos $Q_j$ por un factor, la densidad de carga en todos los conductores se multiplica por el mismo factor (la distribución \emph{es
única}, \cref{teo:unicidad}, y es lineal en $Q_j$), y con ella el potencial y el campo. Por tanto $\phi_i^{(j)}=p_{ij}\,Q_j$, con un coeficiente de
proporcionalidad $p_{ij}$ que tiene que ver solo con la geometría del conjunto. Si ahora todas las cargas son no nulas, las distribuciones
inducidas por cada una son independientes y, por la linealidad de las ecuaciones, se superponen.

\begin{teorema}[Coeficientes de potencial]\label{teo:pij}
En un sistema de $N$ conductores, el potencial de cada uno depende linealmente de las cargas de todos, incluida la suya:
\begin{equation}\label{eq:pij}
  \resalta{\phi_i=\sum_{j=1}^{N}p_{ij}\,Q_j,\qquad i=1,\dots,N,}\qquad\text{es decir,}\qquad\vphi=\Pm\,\vQ ,
\end{equation}
donde los \emph{coeficientes de potencial} $p_{ij}$ dependen solo de la geometría del sistema (y del medio) y admiten la interpretación
\begin{equation}\label{eq:pij-def}
  p_{ij}=\left.\frac{\phi_i}{Q_j}\right|_{Q_k=0,\ k\neq j}:
\end{equation}
$p_{ij}$ es el potencial que adquiere el conductor $i$ por cada unidad de carga que ponemos en el conductor $j$, con todos los demás descargados.
Su unidad es $[p_{ij}]=\mathrm{V/C}=\mathrm{F^{-1}}$.
\end{teorema}
\begin{proof}
Para cada $j$, sea $u^{(j)}(\vc{r})$ el potencial del \emph{estado elemental} $j$: el conductor $j$ con carga unidad y los demás aislados y descargados,
con $u^{(j)}\to0$ en el infinito. Existe y es único (\cref{teo:unicidad}\,(b)). Definimos $p_{ij}$ como el valor (constante) de $u^{(j)}$ sobre el
conductor $i$. Dadas cargas arbitrarias $Q_1,\dots,Q_N$, consideremos
\[
  \phi(\vc{r})=\sum_{j=1}^{N}Q_j\,u^{(j)}(\vc{r}) .
\]
Esta función es armónica fuera de los conductores (combinación lineal de armónicas), se anula en el infinito, es constante sobre cada conductor
$i$, con valor $\sum_jQ_j\,p_{ij}$, y la carga que determina sobre el conductor $i$ es, por linealidad del flujo,
$\sum_jQ_j\,\delt{i}{j}=Q_i$. Cumple, por tanto, todas las condiciones del problema con cargas $Q_i$, y por el segundo teorema de unicidad es su
solución. Evaluándola sobre el conductor $i$ se obtiene \cref{eq:pij}. Como cada $u^{(j)}$ solo depende de la forma y posición de los conductores,
lo mismo ocurre con los $p_{ij}$; tomando $Q_k=0$ para $k\neq j$ en \cref{eq:pij} queda $\phi_i=p_{ij}Q_j$, que es \cref{eq:pij-def}.
\end{proof}

\begin{observacion}[Expresión integral]\label{obs:pij-integral}
Es instructivo escribir explícitamente la dependencia en la geometría (\cref{fig:sistema-geom}). El potencial en un punto cualquiera, en
particular sobre el conductor $i$, es la suma de las contribuciones de las cargas superficiales de todos los conductores,
\[
  \phi(\vc{r})=\frac{1}{4\pi\eo}\sum_{j=1}^{N}\oiint_{S_j'}\frac{\sigma_j(\rp_j)}{R_j}\,\dif S_j',\qquad
  \phi_i=\frac{1}{4\pi\eo}\sum_{j=1}^{N}\oiint_{S_j'}\frac{\sigma_j(\rp_j)}{R_{ij}}\,\dif S_j' ,
\]
con $R_{ij}=\norm{\vc{r}_i-\rp_j}$. Multiplicando por $\sigma_i$ e integrando sobre $S_i'$, y usando que $\phi_i$ es constante sobre $S_i'$,
$\oiint_{S_i'}\phi_i\,\sigma_i\,\dif S_i'=\phi_iQ_i$, de modo que
\[
  \phi_i=\sum_{j=1}^{N}\left[\frac{1}{4\pi\eo}\,\frac{1}{Q_iQ_j}\oiint_{S_i'}\oiint_{S_j'}\frac{\sigma_i\,\sigma_j}{R_{ij}}\,\dif S_j'\,\dif S_i'\right]Q_j ,
\]
y el corchete hace el papel de $p_{ij}$. Si $Q_i$ o $Q_j$ varían, $\sigma_i$ o $\sigma_j$ lo hacen en la misma proporción, y el corchete queda inalterado.
Esta lectura es exacta cuando se aplica al estado elemental $j$ (con $\sigma_i,\sigma_j$ las densidades de ese estado); el \cref{teo:pij} le da
forma rigurosa sin necesidad de conocer las densidades.
\end{observacion}

\begin{observacion}[Cómo se calculan en la práctica]\label{obs:pij-practica}
En general, los $p_{ij}$ \emph{no} se calculan con su expresión integral, sino a través de su interpretación \cref{eq:pij-def}: se carga un conductor,
se dejan los demás descargados y se calculan los potenciales de todos. Un sistema de $N$ conductores tiene $N^{2}$ coeficientes de potencial, aunque,
como vamos a ver, solo $N(N+1)/2$ son independientes.
\end{observacion}

\paso{Propiedades.} Para demostrarlas necesitamos un resultado clásico de Green~\cite{GREEN}.

\begin{lema}[Teorema de reciprocidad de Green]\label{lem:reciprocidad}
Sean dos estados de equilibrio del mismo sistema de conductores, con cargas $Q_i$ y potenciales $\phi_i$ en el primero, y $Q_i'$ y $\phi_i'$ en el segundo.
Entonces
\begin{equation}\label{eq:reciprocidad}
  \sum_{i=1}^{N}Q_i\,\phi_i'=\sum_{i=1}^{N}Q_i'\,\phi_i .
\end{equation}
\end{lema}
\begin{proof}
Sean $\sigma$ y $\sigma'$ las densidades superficiales de ambos estados. Como $\phi'$ es constante sobre cada conductor,
\[
  \sum_iQ_i\,\phi_i'=\sum_i\oiint_{\Sc_i}\sigma(\vc{r})\,\phi'(\vc{r})\,\dsc
  =\frac{1}{4\pi\eo}\oiint_{\Sc}\oiint_{\Sc}\frac{\sigma(\vc{r})\,\sigma'(\rp)}{\norm{\vc{r}-\rp}}\,\dSp\,\dsc ,
\]
donde $\Sc=\bigcup_i\Sc_i$ y hemos sustituido $\phi'$ por su expresión de Coulomb. El integrando es simétrico al intercambiar $(\vc{r},\sigma)$ con
$(\rp,\sigma')$, de modo que la integral doble es también igual a $\sum_iQ_i'\,\phi_i$.
\end{proof}

\begin{teorema}[Propiedades de los coeficientes de potencial]\label{teo:pij-prop}
\begin{enumerate}[label=\textup{(\alph*)}]
  \item Solo dependen de la geometría del conjunto de conductores y del medio.
  \item Son simétricos: $p_{ij}=p_{ji}$. Si una carga $Q$ en el conductor $j$ lleva al conductor $i$ a un potencial $\phi$, la misma carga $Q$ colocada en $i$
        lleva a $j$ al mismo potencial $\phi$.
  \item Son no negativos y el diagonal domina en su columna: $p_{jj}\ge p_{ij}\ge0$, y $p_{jj}>0$.
\end{enumerate}
\end{teorema}
\begin{proof}
(a) Es parte del \cref{teo:pij}.

(b) Aplicamos el \cref{lem:reciprocidad} al estado elemental $j$ ($Q_k=\delt{k}{j}$, $\phi_k=p_{kj}$) y al estado elemental $i$
($Q_k'=\delt{k}{i}$, $\phi_k'=p_{ki}$): $\sum_k\delt{k}{j}\,p_{ki}=\sum_k\delt{k}{i}\,p_{kj}$, es decir, $p_{ji}=p_{ij}$.

(c) En el estado elemental $j$, el conductor $j$ tiene carga positiva y los demás son neutros. Por el \cref{lem:maximo}, el máximo y el mínimo del
potencial se alcanzan sobre algún conductor o en el infinito, donde $\phi=0$. Sea $m$ el conductor (o el infinito) con el potencial \emph{máximo}. Ninguna
línea de campo puede llegar a él (vendría de un potencial mayor, \cref{cor:lineas}), luego $\sigma\ge0$ en toda su superficie. Si $m\neq j$ fuese un
conductor neutro, la carga total nula con $\sigma\ge0$ obligaría a que $\sigma\equiv0$, es decir, a que $\partial\phi/\partial n=0$ en toda su
superficie, lo que es imposible en el máximo de una función armónica no constante (lema de Hopf~\cite{KELLOGG}). Así, el máximo se alcanza en $j$
(quizá compartido con algún otro conductor, si está encerrado por $j$) y $p_{jj}\ge p_{ij}$. Análogamente, en el conductor con el potencial
\emph{mínimo} no puede nacer ninguna línea, luego $\sigma\le0$ en él, y no puede ser ni $j$, que tiene carga positiva, ni un conductor neutro, por
el mismo argumento: el mínimo se alcanza en el infinito, y $p_{ij}\ge0$. Por
último, de $j$ salen líneas, que acaban a menor potencial, luego $p_{jj}>0$.
\end{proof}

\begin{observacion}[Lectura física del signo]\label{obs:pij-signo}
Si todas las cargas son positivas, todos los potenciales son positivos. Las líneas de campo que salen del conductor cargado llegan a los demás o al
infinito; un conductor neutro que recibe líneas por una cara las reemite por la otra (\cref{fig:sistema-libre}), y como el potencial decrece a lo
largo de las líneas, $\phi_j>\phi_{i\neq j}>\phi_{\infty}=0$. El potencial de un conductor no es solo «suyo»: depende de la carga de todo el sistema,
incluida la propia.
\end{observacion}

% =========================================================================
%  Figura «tipo Mathematica»: estado «disco cargado, resto neutros» del
%  sistema de cuatro conductores (combinación de los estados elementales con
%  phi = P e_C). Líneas de campo con flujo fijo por línea, equipotenciales y
%  signo de la carga inducida. Datos: figuras/datos/sistema_libre_*.
% =========================================================================
\begin{figure}[t]
\centering
\begin{tikzpicture}[x=0.66cm, y=0.66cm, font=\small]
  % Contornos de los cuatro conductores
% Generado por scripts/generar_datos.py (no editar a mano).
\def\formaA{(-4.250,-2.950) -- (-3.050,-3.950) -- (0.220,-0.180) -- cycle}
\def\formaB{(-2.750,1.250) -- (-0.650,1.250) -- (-0.650,2.850) -- (-2.750,2.850) -- cycle}
\def\formaC{(1.330,0.350) -- (1.330,0.363) -- (1.329,0.375) -- (1.329,0.388) -- (1.327,0.400) -- (1.326,0.413) -- (1.324,0.425) -- (1.322,0.437) -- (1.320,0.450) -- (1.317,0.462) -- (1.314,0.474) -- (1.310,0.486) -- (1.307,0.498) -- (1.302,0.510) -- (1.298,0.522) -- (1.293,0.534) -- (1.289,0.545) -- (1.283,0.557) -- (1.278,0.568) -- (1.272,0.579) -- (1.266,0.590) -- (1.259,0.601) -- (1.253,0.611) -- (1.246,0.622) -- (1.238,0.632) -- (1.231,0.642) -- (1.223,0.652) -- (1.215,0.662) -- (1.207,0.671) -- (1.198,0.680) -- (1.189,0.689) -- (1.180,0.698) -- (1.171,0.707) -- (1.162,0.715) -- (1.152,0.723) -- (1.142,0.731) -- (1.132,0.738) -- (1.122,0.746) -- (1.111,0.753) -- (1.101,0.759) -- (1.090,0.766) -- (1.079,0.772) -- (1.068,0.778) -- (1.057,0.783) -- (1.045,0.789) -- (1.034,0.793) -- (1.022,0.798) -- (1.010,0.802) -- (0.998,0.807) -- (0.986,0.810) -- (0.974,0.814) -- (0.962,0.817) -- (0.950,0.820) -- (0.937,0.822) -- (0.925,0.824) -- (0.913,0.826) -- (0.900,0.827) -- (0.888,0.829) -- (0.875,0.829) -- (0.863,0.830) -- (0.850,0.830) -- (0.837,0.830) -- (0.825,0.829) -- (0.812,0.829) -- (0.800,0.827) -- (0.787,0.826) -- (0.775,0.824) -- (0.763,0.822) -- (0.750,0.820) -- (0.738,0.817) -- (0.726,0.814) -- (0.714,0.810) -- (0.702,0.807) -- (0.690,0.802) -- (0.678,0.798) -- (0.666,0.793) -- (0.655,0.789) -- (0.643,0.783) -- (0.632,0.778) -- (0.621,0.772) -- (0.610,0.766) -- (0.599,0.759) -- (0.589,0.753) -- (0.578,0.746) -- (0.568,0.738) -- (0.558,0.731) -- (0.548,0.723) -- (0.538,0.715) -- (0.529,0.707) -- (0.520,0.698) -- (0.511,0.689) -- (0.502,0.680) -- (0.493,0.671) -- (0.485,0.662) -- (0.477,0.652) -- (0.469,0.642) -- (0.462,0.632) -- (0.454,0.622) -- (0.447,0.611) -- (0.441,0.601) -- (0.434,0.590) -- (0.428,0.579) -- (0.422,0.568) -- (0.417,0.557) -- (0.411,0.545) -- (0.407,0.534) -- (0.402,0.522) -- (0.398,0.510) -- (0.393,0.498) -- (0.390,0.486) -- (0.386,0.474) -- (0.383,0.462) -- (0.380,0.450) -- (0.378,0.437) -- (0.376,0.425) -- (0.374,0.413) -- (0.373,0.400) -- (0.371,0.388) -- (0.371,0.375) -- (0.370,0.363) -- (0.370,0.350) -- (0.370,0.337) -- (0.371,0.325) -- (0.371,0.312) -- (0.373,0.300) -- (0.374,0.287) -- (0.376,0.275) -- (0.378,0.263) -- (0.380,0.250) -- (0.383,0.238) -- (0.386,0.226) -- (0.390,0.214) -- (0.393,0.202) -- (0.398,0.190) -- (0.402,0.178) -- (0.407,0.166) -- (0.411,0.155) -- (0.417,0.143) -- (0.422,0.132) -- (0.428,0.121) -- (0.434,0.110) -- (0.441,0.099) -- (0.447,0.089) -- (0.454,0.078) -- (0.462,0.068) -- (0.469,0.058) -- (0.477,0.048) -- (0.485,0.038) -- (0.493,0.029) -- (0.502,0.020) -- (0.511,0.011) -- (0.520,0.002) -- (0.529,-0.007) -- (0.538,-0.015) -- (0.548,-0.023) -- (0.558,-0.031) -- (0.568,-0.038) -- (0.578,-0.046) -- (0.589,-0.053) -- (0.599,-0.059) -- (0.610,-0.066) -- (0.621,-0.072) -- (0.632,-0.078) -- (0.643,-0.083) -- (0.655,-0.089) -- (0.666,-0.093) -- (0.678,-0.098) -- (0.690,-0.102) -- (0.702,-0.107) -- (0.714,-0.110) -- (0.726,-0.114) -- (0.738,-0.117) -- (0.750,-0.120) -- (0.763,-0.122) -- (0.775,-0.124) -- (0.787,-0.126) -- (0.800,-0.127) -- (0.812,-0.129) -- (0.825,-0.129) -- (0.837,-0.130) -- (0.850,-0.130) -- (0.863,-0.130) -- (0.875,-0.129) -- (0.888,-0.129) -- (0.900,-0.127) -- (0.913,-0.126) -- (0.925,-0.124) -- (0.937,-0.122) -- (0.950,-0.120) -- (0.962,-0.117) -- (0.974,-0.114) -- (0.986,-0.110) -- (0.998,-0.107) -- (1.010,-0.102) -- (1.022,-0.098) -- (1.034,-0.093) -- (1.045,-0.089) -- (1.057,-0.083) -- (1.068,-0.078) -- (1.079,-0.072) -- (1.090,-0.066) -- (1.101,-0.059) -- (1.111,-0.053) -- (1.122,-0.046) -- (1.132,-0.038) -- (1.142,-0.031) -- (1.152,-0.023) -- (1.162,-0.015) -- (1.171,-0.007) -- (1.180,0.002) -- (1.189,0.011) -- (1.198,0.020) -- (1.207,0.029) -- (1.215,0.038) -- (1.223,0.048) -- (1.231,0.058) -- (1.238,0.068) -- (1.246,0.078) -- (1.253,0.089) -- (1.259,0.099) -- (1.266,0.110) -- (1.272,0.121) -- (1.278,0.132) -- (1.283,0.143) -- (1.289,0.155) -- (1.293,0.166) -- (1.298,0.178) -- (1.302,0.190) -- (1.307,0.202) -- (1.310,0.214) -- (1.314,0.226) -- (1.317,0.238) -- (1.320,0.250) -- (1.322,0.263) -- (1.324,0.275) -- (1.326,0.287) -- (1.327,0.300) -- (1.329,0.312) -- (1.329,0.325) -- (1.330,0.337) -- cycle}
\def\formaD{(1.050,1.320) -- (3.550,1.320) -- (3.550,1.680) -- (1.050,1.680) -- cycle}

  %% potenciales del estado libre (phi_C = 1)
\expandafter\def\csname philibreA\endcsname{0{,}23}
\expandafter\def\csname philibreB\endcsname{0{,}20}
\expandafter\def\csname philibreC\endcsname{1{,}00}
\expandafter\def\csname philibreD\endcsname{0{,}41}

  \fill[azulfondo!60!white] (-5,-4.5) rectangle (5,4.5);
  \begin{scope}
    \clip (-5,-4.5) rectangle (5,4.5);
    \tikzset{numEQ/.style={azulcielo, line width=0.35pt, densely dashed}}
    % Disco cargado, resto neutros - equipotenciales
% Generado por scripts/generar_datos.py (no editar a mano).
\EQ{(-3.260,-4.241) (-3.220,-4.246) (-3.180,-4.249) (-3.140,-4.252) (-3.100,-4.253) (-3.060,-4.254) (-3.020,-4.253) (-2.980,-4.251) (-2.940,-4.248) (-2.900,-4.244) (-2.860,-4.239) (-2.820,-4.233) (-2.780,-4.226) (-2.740,-4.219) (-2.700,-4.211) (-2.660,-4.203) (-2.620,-4.194) (-2.580,-4.185) (-2.540,-4.176) (-2.500,-4.166) (-2.460,-4.156) (-2.420,-4.146) (-2.380,-4.135) (-2.340,-4.125) (-2.300,-4.115) (-2.260,-4.104) (-2.220,-4.094) (-2.180,-4.083) (-2.140,-4.072) (-2.100,-4.062) (-2.060,-4.051) (-2.020,-4.041) (-1.980,-4.030) (-1.941,-4.020) (-1.900,-4.009) (-1.860,-3.999) (-1.820,-3.989) (-1.780,-3.979) (-1.740,-3.969) (-1.700,-3.959) (-1.660,-3.949) (-1.620,-3.940) (-1.580,-3.930) (-1.540,-3.921) (-1.500,-3.912) (-1.460,-3.903) (-1.420,-3.894) (-1.380,-3.886) (-1.340,-3.877) (-1.300,-3.869) (-1.260,-3.861) (-1.220,-3.853) (-1.180,-3.846) (-1.140,-3.838) (-1.100,-3.831) (-1.060,-3.824) (-1.020,-3.818) (-0.980,-3.811) (-0.940,-3.805) (-0.900,-3.799) (-0.860,-3.793) (-0.820,-3.787) (-0.780,-3.782) (-0.740,-3.777) (-0.700,-3.772) (-0.660,-3.767) (-0.620,-3.763) (-0.580,-3.759) (-0.540,-3.755) (-0.500,-3.751) (-0.460,-3.747) (-0.420,-3.744) (-0.380,-3.741) (-0.340,-3.737) (-0.300,-3.735) (-0.260,-3.732) (-0.220,-3.729) (-0.180,-3.727) (-0.140,-3.724) (-0.100,-3.722) (-0.060,-3.720) (-0.020,-3.718) (0.020,-3.716) (0.060,-3.715) (0.100,-3.713) (0.140,-3.711) (0.180,-3.710) (0.220,-3.708) (0.260,-3.707) (0.300,-3.706) (0.340,-3.704) (0.380,-3.703) (0.420,-3.702) (0.460,-3.700) (0.500,-3.699) (0.540,-3.698) (0.580,-3.696) (0.620,-3.695) (0.660,-3.693) (0.700,-3.692) (0.740,-3.690) (0.780,-3.689) (0.820,-3.687) (0.860,-3.685) (0.900,-3.684) (0.940,-3.682) (0.980,-3.680) (1.020,-3.678) (1.060,-3.675) (1.100,-3.673) (1.140,-3.670) (1.180,-3.668) (1.220,-3.665) (1.260,-3.662) (1.300,-3.659) (1.340,-3.656) (1.380,-3.653) (1.420,-3.649) (1.460,-3.645) (1.500,-3.641) (1.540,-3.637) (1.580,-3.633) (1.620,-3.628) (1.660,-3.624) (1.700,-3.619) (1.740,-3.613) (1.780,-3.608) (1.820,-3.602) (1.860,-3.596) (1.900,-3.590) (1.940,-3.584) (1.980,-3.577) (2.020,-3.570) (2.060,-3.562) (2.100,-3.555) (2.140,-3.547) (2.180,-3.538) (2.220,-3.529) (2.260,-3.520) (2.300,-3.511) (2.340,-3.501) (2.380,-3.491) (2.420,-3.480) (2.460,-3.469) (2.500,-3.457) (2.540,-3.445) (2.580,-3.433) (2.618,-3.420) (2.660,-3.406) (2.700,-3.391) (2.740,-3.377) (2.780,-3.361) (2.820,-3.345) (2.860,-3.328) (2.900,-3.310) (2.940,-3.292) (2.980,-3.273) (3.020,-3.253) (3.060,-3.232) (3.100,-3.210) (3.140,-3.187) (3.180,-3.163) (3.217,-3.140) (3.260,-3.112) (3.300,-3.084) (3.334,-3.060) (3.380,-3.025) (3.412,-3.000) (3.460,-2.960) (3.500,-2.924) (3.540,-2.887) (3.580,-2.848) (3.620,-2.806) (3.660,-2.762) (3.696,-2.720) (3.729,-2.680) (3.760,-2.640) (3.790,-2.600) (3.818,-2.560) (3.846,-2.520) (3.872,-2.480) (3.896,-2.440) (3.920,-2.401) (3.940,-2.366) (3.960,-2.330) (3.980,-2.293) (4.000,-2.254) (4.020,-2.214) (4.040,-2.172) (4.060,-2.129) (4.080,-2.084) (4.099,-2.040) (4.115,-2.000) (4.131,-1.960) (4.146,-1.920) (4.161,-1.880) (4.175,-1.840) (4.189,-1.800) (4.203,-1.760) (4.216,-1.720) (4.228,-1.680) (4.241,-1.640) (4.253,-1.600) (4.264,-1.560) (4.275,-1.520) (4.286,-1.480) (4.297,-1.440) (4.307,-1.400) (4.317,-1.360) (4.327,-1.320) (4.336,-1.280) (4.345,-1.240) (4.354,-1.200) (4.363,-1.160) (4.371,-1.120) (4.379,-1.080) (4.387,-1.040) (4.395,-1.000) (4.403,-0.960) (4.410,-0.920) (4.417,-0.880) (4.424,-0.840) (4.431,-0.800) (4.438,-0.760) (4.444,-0.720) (4.451,-0.680) (4.457,-0.640) (4.463,-0.600) (4.469,-0.560) (4.475,-0.520) (4.481,-0.480) (4.486,-0.440) (4.492,-0.400) (4.497,-0.360) (4.502,-0.320) (4.507,-0.280) (4.512,-0.240) (4.517,-0.200) (4.522,-0.160) (4.527,-0.120) (4.532,-0.080) (4.536,-0.040) (4.541,-0.000) (4.545,0.040) (4.550,0.080) (4.554,0.120) (4.558,0.160) (4.562,0.200) (4.566,0.240) (4.570,0.280) (4.574,0.320) (4.578,0.360) (4.582,0.400) (4.585,0.440) (4.589,0.480) (4.592,0.520) (4.596,0.560) (4.599,0.600) (4.602,0.640) (4.605,0.680) (4.608,0.720) (4.611,0.760) (4.613,0.800) (4.616,0.840) (4.618,0.880) (4.620,0.920) (4.622,0.960) (4.624,1.000) (4.626,1.040) (4.627,1.080) (4.628,1.120) (4.629,1.160) (4.630,1.200) (4.630,1.240) (4.631,1.280) (4.630,1.320) (4.630,1.360) (4.629,1.400) (4.628,1.440) (4.627,1.480) (4.625,1.520) (4.623,1.560) (4.621,1.600) (4.618,1.640) (4.615,1.680) (4.611,1.720) (4.608,1.760) (4.603,1.800) (4.599,1.840) (4.593,1.880) (4.588,1.920) (4.582,1.960) (4.575,2.000) (4.569,2.040) (4.561,2.080) (4.553,2.120) (4.545,2.160) (4.536,2.200) (4.527,2.240) (4.517,2.280) (4.507,2.320) (4.496,2.360) (4.484,2.400) (4.472,2.440) (4.459,2.480) (4.445,2.520) (4.431,2.560) (4.416,2.600) (4.400,2.640) (4.384,2.680) (4.366,2.720) (4.348,2.760) (4.329,2.800) (4.308,2.840) (4.287,2.880) (4.264,2.920) (4.240,2.960) (4.214,3.000) (4.187,3.040) (4.160,3.078) (4.128,3.120) (4.100,3.154) (4.060,3.200) (4.023,3.240) (3.982,3.280) (3.940,3.319) (3.900,3.353) (3.866,3.380) (3.820,3.415) (3.783,3.440) (3.740,3.468) (3.700,3.493) (3.660,3.516) (3.620,3.537) (3.580,3.558) (3.540,3.577) (3.500,3.595) (3.460,3.612) (3.420,3.628) (3.380,3.643) (3.340,3.657) (3.300,3.671) (3.260,3.684) (3.220,3.696) (3.180,3.707) (3.140,3.718) (3.100,3.728) (3.060,3.738) (3.020,3.748) (2.980,3.756) (2.940,3.765) (2.900,3.773) (2.861,3.780) (2.820,3.787) (2.780,3.794) (2.740,3.800) (2.700,3.806) (2.660,3.812) (2.620,3.818) (2.580,3.823) (2.540,3.828) (2.500,3.832) (2.460,3.837) (2.420,3.841) (2.380,3.845) (2.340,3.849) (2.300,3.852) (2.260,3.855) (2.220,3.858) (2.180,3.861) (2.140,3.864) (2.100,3.867) (2.060,3.869) (2.020,3.871) (1.980,3.873) (1.940,3.875) (1.900,3.877) (1.860,3.878) (1.820,3.880) (1.780,3.881) (1.740,3.882) (1.700,3.883) (1.660,3.884) (1.620,3.885) (1.580,3.886) (1.540,3.886) (1.500,3.887) (1.460,3.887) (1.420,3.888) (1.380,3.888) (1.340,3.888) (1.300,3.888) (1.260,3.888) (1.220,3.887) (1.180,3.887) (1.140,3.886) (1.100,3.886) (1.060,3.885) (1.020,3.885) (0.980,3.884) (0.940,3.883) (0.900,3.882) (0.860,3.881) (0.820,3.880) (0.780,3.878) (0.740,3.877) (0.700,3.876) (0.660,3.874) (0.620,3.873) (0.580,3.871) (0.540,3.870) (0.500,3.868) (0.460,3.866) (0.420,3.865) (0.380,3.863) (0.340,3.861) (0.300,3.859) (0.260,3.858) (0.220,3.856) (0.180,3.854) (0.140,3.852) (0.100,3.850) (0.060,3.848) (0.020,3.846) (-0.020,3.845) (-0.060,3.843) (-0.100,3.841) (-0.140,3.840) (-0.180,3.838) (-0.220,3.836) (-0.260,3.835) (-0.300,3.833) (-0.340,3.832) (-0.380,3.830) (-0.420,3.829) (-0.460,3.828) (-0.500,3.827) (-0.540,3.826) (-0.580,3.825) (-0.620,3.824) (-0.660,3.823) (-0.700,3.822) (-0.740,3.821) (-0.780,3.820) (-0.820,3.819) (-0.860,3.818) (-0.900,3.817) (-0.940,3.817) (-0.980,3.816) (-1.020,3.815) (-1.060,3.814) (-1.100,3.813) (-1.140,3.813) (-1.180,3.812) (-1.220,3.811) (-1.260,3.810) (-1.300,3.809) (-1.340,3.808) (-1.380,3.807) (-1.420,3.806) (-1.460,3.804) (-1.500,3.803) (-1.540,3.802) (-1.580,3.800) (-1.620,3.798) (-1.660,3.797) (-1.700,3.795) (-1.740,3.793) (-1.780,3.791) (-1.820,3.788) (-1.860,3.786) (-1.900,3.783) (-1.940,3.781) (-1.980,3.777) (-2.020,3.774) (-2.060,3.771) (-2.100,3.767) (-2.140,3.763) (-2.180,3.758) (-2.220,3.754) (-2.260,3.749) (-2.300,3.743) (-2.340,3.737) (-2.380,3.731) (-2.420,3.725) (-2.460,3.717) (-2.500,3.710) (-2.540,3.701) (-2.580,3.693) (-2.620,3.683) (-2.660,3.673) (-2.700,3.662) (-2.740,3.650) (-2.780,3.638) (-2.820,3.624) (-2.860,3.610) (-2.900,3.595) (-2.940,3.578) (-2.980,3.561) (-3.020,3.542) (-3.060,3.522) (-3.100,3.501) (-3.136,3.480) (-3.180,3.453) (-3.220,3.426) (-3.257,3.400) (-3.300,3.368) (-3.334,3.340) (-3.379,3.300) (-3.420,3.261) (-3.460,3.220) (-3.496,3.180) (-3.529,3.140) (-3.560,3.100) (-3.589,3.060) (-3.616,3.020) (-3.640,2.982) (-3.665,2.940) (-3.687,2.900) (-3.708,2.860) (-3.727,2.820) (-3.745,2.780) (-3.762,2.740) (-3.779,2.700) (-3.794,2.660) (-3.808,2.620) (-3.821,2.580) (-3.834,2.540) (-3.846,2.500) (-3.857,2.460) (-3.868,2.420) (-3.877,2.380) (-3.887,2.340) (-3.895,2.300) (-3.903,2.260) (-3.911,2.220) (-3.918,2.180) (-3.925,2.140) (-3.931,2.100) (-3.937,2.060) (-3.942,2.020) (-3.948,1.980) (-3.952,1.940) (-3.957,1.900) (-3.961,1.860) (-3.964,1.820) (-3.968,1.780) (-3.971,1.740) (-3.974,1.700) (-3.977,1.660) (-3.979,1.620) (-3.981,1.580) (-3.983,1.540) (-3.985,1.500) (-3.987,1.460) (-3.988,1.420) (-3.990,1.380) (-3.991,1.340) (-3.992,1.300) (-3.993,1.260) (-3.994,1.220) (-3.994,1.180) (-3.995,1.140) (-3.996,1.100) (-3.996,1.060) (-3.997,1.020) (-3.997,0.980) (-3.998,0.940) (-3.999,0.900) (-3.999,0.860) (-4.000,0.820) (-4.000,0.780) (-4.001,0.740) (-4.002,0.700) (-4.003,0.660) (-4.004,0.620) (-4.005,0.580) (-4.006,0.540) (-4.007,0.500) (-4.009,0.460) (-4.010,0.420) (-4.012,0.380) (-4.014,0.340) (-4.016,0.300) (-4.018,0.260) (-4.020,0.220) (-4.023,0.180) (-4.026,0.140) (-4.029,0.100) (-4.032,0.060) (-4.035,0.020) (-4.039,-0.020) (-4.042,-0.060) (-4.046,-0.100) (-4.050,-0.140) (-4.055,-0.180) (-4.059,-0.220) (-4.064,-0.260) (-4.069,-0.300) (-4.074,-0.340) (-4.080,-0.380) (-4.085,-0.420) (-4.091,-0.460) (-4.097,-0.500) (-4.104,-0.540) (-4.110,-0.580) (-4.117,-0.620) (-4.124,-0.660) (-4.131,-0.700) (-4.139,-0.740) (-4.146,-0.780) (-4.154,-0.820) (-4.162,-0.860) (-4.170,-0.900) (-4.179,-0.940) (-4.187,-0.980) (-4.196,-1.020) (-4.205,-1.060) (-4.215,-1.100) (-4.224,-1.140) (-4.234,-1.180) (-4.243,-1.220) (-4.253,-1.260) (-4.264,-1.300) (-4.274,-1.340) (-4.284,-1.380) (-4.295,-1.420) (-4.306,-1.460) (-4.316,-1.500) (-4.327,-1.540) (-4.338,-1.580) (-4.350,-1.620) (-4.361,-1.660) (-4.372,-1.700) (-4.383,-1.740) (-4.395,-1.780) (-4.406,-1.820) (-4.418,-1.860) (-4.429,-1.900) (-4.441,-1.940) (-4.453,-1.980) (-4.464,-2.020) (-4.475,-2.060) (-4.487,-2.100) (-4.498,-2.140) (-4.509,-2.180) (-4.520,-2.220) (-4.531,-2.260) (-4.541,-2.300) (-4.552,-2.340) (-4.562,-2.380) (-4.571,-2.420) (-4.581,-2.460) (-4.590,-2.500) (-4.598,-2.540) (-4.606,-2.580) (-4.613,-2.620) (-4.620,-2.660) (-4.626,-2.700) (-4.631,-2.740) (-4.635,-2.780) (-4.638,-2.820) (-4.641,-2.860) (-4.642,-2.900) (-4.642,-2.940) (-4.641,-2.980) (-4.639,-3.020) (-4.636,-3.060) (-4.632,-3.100) (-4.626,-3.140) (-4.620,-3.180) (-4.612,-3.220) (-4.603,-3.260) (-4.592,-3.300) (-4.581,-3.340) (-4.568,-3.380) (-4.553,-3.420) (-4.538,-3.460) (-4.520,-3.500) (-4.501,-3.540) (-4.480,-3.580) (-4.460,-3.616) (-4.432,-3.660) (-4.404,-3.700) (-4.373,-3.740) (-4.340,-3.779) (-4.300,-3.820) (-4.260,-3.857) (-4.220,-3.891) (-4.181,-3.920) (-4.140,-3.949) (-4.100,-3.974) (-4.060,-3.997) (-4.020,-4.018) (-3.980,-4.038) (-3.940,-4.056) (-3.900,-4.073) (-3.860,-4.089) (-3.820,-4.104) (-3.780,-4.119) (-3.740,-4.132) (-3.700,-4.145) (-3.660,-4.157) (-3.620,-4.168) (-3.580,-4.178) (-3.540,-4.189) (-3.500,-4.198) (-3.460,-4.207) (-3.420,-4.215) (-3.380,-4.223) (-3.340,-4.230) (-3.300,-4.236) (-3.260,-4.241)}
\EQ{(-3.120,-4.042) (-3.080,-4.047) (-3.040,-4.046) (-3.000,-4.041) (-2.960,-4.031) (-2.920,-4.019) (-2.880,-4.003) (-2.840,-3.985) (-2.800,-3.966) (-2.760,-3.945) (-2.720,-3.923) (-2.680,-3.900) (-2.645,-3.880) (-2.600,-3.853) (-2.560,-3.829) (-2.520,-3.805) (-2.480,-3.781) (-2.440,-3.756) (-2.400,-3.731) (-2.360,-3.707) (-2.320,-3.682) (-2.285,-3.660) (-2.240,-3.632) (-2.200,-3.608) (-2.160,-3.583) (-2.122,-3.560) (-2.080,-3.534) (-2.040,-3.510) (-2.000,-3.486) (-1.960,-3.462) (-1.922,-3.440) (-1.880,-3.415) (-1.840,-3.392) (-1.800,-3.370) (-1.760,-3.347) (-1.720,-3.325) (-1.680,-3.303) (-1.640,-3.282) (-1.600,-3.261) (-1.560,-3.240) (-1.520,-3.220) (-1.480,-3.201) (-1.440,-3.182) (-1.400,-3.163) (-1.360,-3.145) (-1.320,-3.127) (-1.280,-3.110) (-1.240,-3.094) (-1.200,-3.078) (-1.160,-3.063) (-1.120,-3.049) (-1.080,-3.035) (-1.040,-3.022) (-1.000,-3.010) (-0.960,-2.998) (-0.920,-2.987) (-0.880,-2.977) (-0.840,-2.967) (-0.800,-2.958) (-0.760,-2.950) (-0.720,-2.942) (-0.680,-2.935) (-0.640,-2.929) (-0.600,-2.923) (-0.560,-2.918) (-0.520,-2.913) (-0.480,-2.909) (-0.440,-2.906) (-0.400,-2.903) (-0.360,-2.901) (-0.320,-2.899) (-0.280,-2.897) (-0.240,-2.896) (-0.200,-2.895) (-0.160,-2.894) (-0.120,-2.894) (-0.080,-2.894) (-0.040,-2.894) (-0.000,-2.895) (0.040,-2.895) (0.080,-2.896) (0.120,-2.897) (0.160,-2.898) (0.200,-2.899) (0.240,-2.900) (0.280,-2.901) (0.320,-2.902) (0.360,-2.903) (0.400,-2.904) (0.440,-2.905) (0.480,-2.905) (0.520,-2.906) (0.560,-2.907) (0.600,-2.907) (0.640,-2.907) (0.680,-2.907) (0.720,-2.907) (0.760,-2.907) (0.800,-2.907) (0.840,-2.906) (0.880,-2.905) (0.920,-2.904) (0.960,-2.902) (1.000,-2.901) (1.040,-2.899) (1.080,-2.896) (1.120,-2.894) (1.160,-2.891) (1.200,-2.888) (1.240,-2.884) (1.280,-2.880) (1.320,-2.876) (1.360,-2.871) (1.400,-2.866) (1.440,-2.861) (1.480,-2.855) (1.520,-2.849) (1.560,-2.843) (1.600,-2.836) (1.640,-2.828) (1.680,-2.820) (1.720,-2.812) (1.760,-2.803) (1.800,-2.794) (1.840,-2.784) (1.880,-2.774) (1.920,-2.763) (1.960,-2.752) (2.000,-2.740) (2.040,-2.728) (2.080,-2.715) (2.120,-2.701) (2.160,-2.687) (2.200,-2.672) (2.240,-2.657) (2.280,-2.640) (2.320,-2.624) (2.360,-2.606) (2.400,-2.588) (2.440,-2.568) (2.480,-2.548) (2.520,-2.527) (2.560,-2.506) (2.600,-2.483) (2.639,-2.460) (2.680,-2.435) (2.720,-2.409) (2.760,-2.382) (2.800,-2.355) (2.840,-2.326) (2.874,-2.300) (2.920,-2.264) (2.960,-2.231) (2.995,-2.200) (3.040,-2.160) (3.080,-2.122) (3.120,-2.083) (3.160,-2.042) (3.199,-2.000) (3.234,-1.960) (3.268,-1.920) (3.300,-1.881) (3.332,-1.840) (3.360,-1.803) (3.391,-1.760) (3.420,-1.720) (3.447,-1.680) (3.473,-1.640) (3.498,-1.600) (3.520,-1.563) (3.546,-1.520) (3.568,-1.480) (3.590,-1.440) (3.611,-1.400) (3.632,-1.360) (3.652,-1.320) (3.671,-1.280) (3.690,-1.240) (3.708,-1.200) (3.726,-1.160) (3.743,-1.120) (3.760,-1.080) (3.776,-1.040) (3.791,-1.000) (3.807,-0.960) (3.821,-0.920) (3.836,-0.880) (3.850,-0.840) (3.863,-0.800) (3.877,-0.760) (3.890,-0.720) (3.902,-0.680) (3.915,-0.640) (3.927,-0.600) (3.938,-0.560) (3.950,-0.520) (3.961,-0.480) (3.972,-0.440) (3.983,-0.400) (3.993,-0.360) (4.004,-0.320) (4.014,-0.280) (4.024,-0.240) (4.033,-0.200) (4.043,-0.160) (4.053,-0.120) (4.062,-0.080) (4.071,-0.040) (4.080,-0.000) (4.089,0.040) (4.098,0.080) (4.106,0.120) (4.115,0.160) (4.123,0.200) (4.132,0.240) (4.140,0.280) (4.148,0.320) (4.156,0.360) (4.164,0.400) (4.172,0.440) (4.180,0.480) (4.187,0.520) (4.195,0.560) (4.202,0.600) (4.209,0.640) (4.216,0.680) (4.223,0.720) (4.230,0.760) (4.236,0.800) (4.242,0.840) (4.248,0.880) (4.254,0.920) (4.259,0.960) (4.263,1.000) (4.268,1.040) (4.272,1.080) (4.275,1.120) (4.278,1.160) (4.281,1.200) (4.282,1.240) (4.284,1.280) (4.284,1.320) (4.284,1.360) (4.284,1.400) (4.282,1.440) (4.280,1.480) (4.277,1.520) (4.274,1.560) (4.269,1.600) (4.264,1.640) (4.258,1.680) (4.251,1.720) (4.243,1.760) (4.234,1.800) (4.224,1.840) (4.213,1.880) (4.201,1.920) (4.188,1.960) (4.174,2.000) (4.160,2.038) (4.143,2.080) (4.125,2.120) (4.106,2.160) (4.086,2.200) (4.065,2.240) (4.042,2.280) (4.020,2.316) (3.992,2.360) (3.964,2.400) (3.940,2.433) (3.904,2.480) (3.871,2.520) (3.840,2.555) (3.800,2.598) (3.760,2.638) (3.720,2.675) (3.680,2.710) (3.644,2.740) (3.600,2.775) (3.566,2.800) (3.520,2.832) (3.480,2.859) (3.446,2.880) (3.400,2.907) (3.360,2.930) (3.320,2.951) (3.280,2.971) (3.240,2.991) (3.200,3.009) (3.160,3.026) (3.120,3.042) (3.080,3.058) (3.040,3.073) (3.000,3.087) (2.961,3.100) (2.920,3.113) (2.880,3.125) (2.840,3.137) (2.800,3.148) (2.760,3.158) (2.720,3.168) (2.680,3.177) (2.640,3.186) (2.600,3.195) (2.560,3.203) (2.520,3.210) (2.480,3.217) (2.440,3.224) (2.400,3.231) (2.360,3.237) (2.320,3.242) (2.280,3.248) (2.240,3.253) (2.200,3.257) (2.160,3.262) (2.120,3.266) (2.080,3.270) (2.040,3.274) (2.000,3.277) (1.960,3.280) (1.920,3.283) (1.880,3.285) (1.840,3.288) (1.800,3.290) (1.760,3.292) (1.720,3.293) (1.680,3.295) (1.640,3.296) (1.600,3.297) (1.560,3.298) (1.520,3.298) (1.480,3.299) (1.440,3.299) (1.400,3.299) (1.360,3.299) (1.320,3.298) (1.280,3.297) (1.240,3.297) (1.200,3.296) (1.160,3.294) (1.120,3.293) (1.080,3.291) (1.040,3.290) (1.000,3.287) (0.960,3.285) (0.920,3.283) (0.880,3.280) (0.840,3.277) (0.800,3.274) (0.760,3.271) (0.720,3.268) (0.680,3.264) (0.640,3.260) (0.600,3.256) (0.560,3.252) (0.520,3.248) (0.480,3.243) (0.440,3.238) (0.400,3.234) (0.360,3.228) (0.320,3.223) (0.280,3.218) (0.240,3.212) (0.200,3.207) (0.160,3.201) (0.120,3.195) (0.080,3.189) (0.040,3.184) (-0.000,3.178) (-0.040,3.172) (-0.080,3.166) (-0.120,3.161) (-0.160,3.156) (-0.200,3.151) (-0.240,3.146) (-0.280,3.142) (-0.320,3.138) (-0.360,3.135) (-0.400,3.132) (-0.440,3.130) (-0.480,3.129) (-0.520,3.128) (-0.560,3.128) (-0.600,3.129) (-0.640,3.129) (-0.680,3.130) (-0.720,3.131) (-0.760,3.132) (-0.800,3.133) (-0.840,3.133) (-0.880,3.133) (-0.920,3.134) (-0.960,3.134) (-1.000,3.134) (-1.040,3.134) (-1.080,3.134) (-1.120,3.134) (-1.160,3.133) (-1.200,3.133) (-1.240,3.133) (-1.280,3.132) (-1.320,3.132) (-1.360,3.131) (-1.400,3.131) (-1.440,3.130) (-1.480,3.129) (-1.520,3.129) (-1.560,3.128) (-1.600,3.127) (-1.640,3.126) (-1.680,3.125) (-1.720,3.124) (-1.760,3.122) (-1.800,3.121) (-1.840,3.120) (-1.880,3.118) (-1.920,3.116) (-1.960,3.114) (-2.000,3.112) (-2.040,3.110) (-2.080,3.107) (-2.120,3.105) (-2.160,3.102) (-2.200,3.099) (-2.240,3.095) (-2.280,3.091) (-2.320,3.087) (-2.360,3.082) (-2.400,3.077) (-2.440,3.071) (-2.480,3.064) (-2.520,3.056) (-2.560,3.048) (-2.600,3.038) (-2.640,3.026) (-2.680,3.013) (-2.720,2.996) (-2.760,2.975) (-2.800,2.949) (-2.833,2.920) (-2.860,2.890) (-2.895,2.840) (-2.916,2.800) (-2.934,2.760) (-2.950,2.720) (-2.964,2.680) (-2.976,2.640) (-2.987,2.600) (-2.997,2.560) (-3.006,2.520) (-3.014,2.480) (-3.022,2.440) (-3.028,2.400) (-3.035,2.360) (-3.041,2.320) (-3.046,2.280) (-3.051,2.240) (-3.055,2.200) (-3.059,2.160) (-3.063,2.120) (-3.067,2.080) (-3.070,2.040) (-3.072,2.000) (-3.075,1.960) (-3.077,1.920) (-3.078,1.880) (-3.080,1.840) (-3.081,1.800) (-3.082,1.760) (-3.082,1.720) (-3.082,1.680) (-3.082,1.640) (-3.081,1.600) (-3.080,1.560) (-3.078,1.520) (-3.076,1.480) (-3.073,1.440) (-3.069,1.400) (-3.065,1.360) (-3.060,1.320) (-3.054,1.280) (-3.047,1.240) (-3.039,1.200) (-3.030,1.160) (-3.021,1.120) (-3.011,1.080) (-3.000,1.040) (-2.990,1.000) (-2.980,0.960) (-2.969,0.920) (-2.960,0.880) (-2.950,0.840) (-2.942,0.800) (-2.934,0.760) (-2.927,0.720) (-2.921,0.680) (-2.915,0.640) (-2.911,0.600) (-2.907,0.560) (-2.903,0.520) (-2.901,0.480) (-2.899,0.440) (-2.898,0.400) (-2.897,0.360) (-2.897,0.320) (-2.898,0.280) (-2.899,0.240) (-2.901,0.200) (-2.903,0.160) (-2.906,0.120) (-2.909,0.080) (-2.913,0.040) (-2.917,-0.000) (-2.921,-0.040) (-2.927,-0.080) (-2.932,-0.120) (-2.938,-0.160) (-2.945,-0.200) (-2.952,-0.240) (-2.960,-0.280) (-2.968,-0.320) (-2.976,-0.360) (-2.985,-0.400) (-2.995,-0.440) (-3.005,-0.480) (-3.016,-0.520) (-3.028,-0.560) (-3.040,-0.600) (-3.052,-0.640) (-3.065,-0.680) (-3.079,-0.720) (-3.093,-0.760) (-3.108,-0.800) (-3.124,-0.840) (-3.140,-0.880) (-3.156,-0.920) (-3.174,-0.960) (-3.192,-1.000) (-3.210,-1.040) (-3.229,-1.080) (-3.249,-1.120) (-3.269,-1.160) (-3.290,-1.200) (-3.312,-1.240) (-3.334,-1.280) (-3.356,-1.320) (-3.379,-1.360) (-3.400,-1.396) (-3.427,-1.440) (-3.451,-1.480) (-3.476,-1.520) (-3.500,-1.558) (-3.527,-1.600) (-3.553,-1.640) (-3.580,-1.680) (-3.607,-1.720) (-3.634,-1.760) (-3.660,-1.798) (-3.689,-1.840) (-3.718,-1.880) (-3.746,-1.920) (-3.775,-1.960) (-3.800,-1.995) (-3.832,-2.040) (-3.860,-2.078) (-3.890,-2.120) (-3.919,-2.160) (-3.948,-2.200) (-3.978,-2.240) (-4.007,-2.280) (-4.035,-2.320) (-4.060,-2.354) (-4.092,-2.400) (-4.120,-2.439) (-4.148,-2.480) (-4.175,-2.520) (-4.200,-2.558) (-4.226,-2.600) (-4.250,-2.640) (-4.273,-2.680) (-4.294,-2.720) (-4.313,-2.760) (-4.329,-2.800) (-4.342,-2.840) (-4.352,-2.880) (-4.358,-2.920) (-4.359,-2.960) (-4.356,-3.000) (-4.348,-3.040) (-4.337,-3.080) (-4.322,-3.120) (-4.304,-3.160) (-4.283,-3.200) (-4.260,-3.238) (-4.233,-3.280) (-4.204,-3.320) (-4.173,-3.360) (-4.140,-3.399) (-4.104,-3.440) (-4.066,-3.480) (-4.025,-3.520) (-3.981,-3.560) (-3.940,-3.596) (-3.900,-3.629) (-3.860,-3.660) (-3.820,-3.690) (-3.780,-3.718) (-3.740,-3.746) (-3.700,-3.772) (-3.660,-3.797) (-3.623,-3.820) (-3.580,-3.845) (-3.540,-3.868) (-3.500,-3.890) (-3.460,-3.911) (-3.420,-3.931) (-3.380,-3.950) (-3.340,-3.969) (-3.300,-3.986) (-3.260,-4.002) (-3.220,-4.017) (-3.180,-4.029) (-3.140,-4.039) (-3.120,-4.042)}
\EQ{(-0.047,-0.507) (-0.072,-0.540) (-0.100,-0.573) (-0.138,-0.620) (-0.160,-0.654) (-0.195,-0.700) (-0.220,-0.732) (-0.257,-0.780) (-0.280,-0.814) (-0.312,-0.860) (-0.340,-0.900) (-0.367,-0.940) (-0.392,-0.980) (-0.416,-1.020) (-0.440,-1.059) (-0.460,-1.094) (-0.480,-1.133) (-0.500,-1.175) (-0.519,-1.220) (-0.533,-1.260) (-0.545,-1.300) (-0.553,-1.340) (-0.557,-1.380) (-0.557,-1.420) (-0.553,-1.460) (-0.544,-1.500) (-0.530,-1.540) (-0.511,-1.580) (-0.486,-1.620) (-0.460,-1.655) (-0.420,-1.700) (-0.380,-1.740) (-0.340,-1.774) (-0.308,-1.800) (-0.260,-1.835) (-0.224,-1.860) (-0.180,-1.888) (-0.140,-1.912) (-0.100,-1.934) (-0.060,-1.955) (-0.020,-1.975) (0.020,-1.994) (0.060,-2.012) (0.100,-2.029) (0.140,-2.045) (0.178,-2.060) (0.220,-2.075) (0.260,-2.089) (0.300,-2.102) (0.340,-2.114) (0.380,-2.126) (0.420,-2.136) (0.460,-2.146) (0.500,-2.156) (0.540,-2.165) (0.580,-2.173) (0.620,-2.180) (0.660,-2.187) (0.700,-2.193) (0.740,-2.199) (0.780,-2.203) (0.820,-2.208) (0.860,-2.211) (0.900,-2.215) (0.940,-2.217) (0.980,-2.219) (1.020,-2.220) (1.060,-2.221) (1.100,-2.221) (1.140,-2.220) (1.180,-2.219) (1.220,-2.217) (1.260,-2.215) (1.300,-2.212) (1.340,-2.208) (1.380,-2.204) (1.420,-2.199) (1.460,-2.193) (1.500,-2.187) (1.540,-2.180) (1.580,-2.172) (1.620,-2.164) (1.660,-2.155) (1.700,-2.145) (1.740,-2.135) (1.780,-2.124) (1.820,-2.112) (1.860,-2.099) (1.900,-2.086) (1.940,-2.072) (1.980,-2.057) (2.020,-2.041) (2.060,-2.024) (2.100,-2.006) (2.140,-1.988) (2.180,-1.969) (2.220,-1.948) (2.260,-1.927) (2.300,-1.904) (2.340,-1.881) (2.374,-1.860) (2.420,-1.831) (2.460,-1.804) (2.494,-1.780) (2.540,-1.746) (2.574,-1.720) (2.620,-1.683) (2.660,-1.649) (2.693,-1.620) (2.737,-1.580) (2.778,-1.540) (2.818,-1.500) (2.856,-1.460) (2.892,-1.420) (2.920,-1.387) (2.959,-1.340) (2.991,-1.300) (3.020,-1.262) (3.051,-1.220) (3.080,-1.180) (3.107,-1.140) (3.133,-1.100) (3.158,-1.060) (3.180,-1.025) (3.207,-0.980) (3.229,-0.940) (3.252,-0.900) (3.273,-0.860) (3.294,-0.820) (3.314,-0.780) (3.333,-0.740) (3.352,-0.700) (3.370,-0.660) (3.388,-0.620) (3.405,-0.580) (3.422,-0.540) (3.438,-0.500) (3.454,-0.460) (3.470,-0.420) (3.485,-0.380) (3.499,-0.340) (3.514,-0.300) (3.528,-0.260) (3.542,-0.220) (3.555,-0.180) (3.569,-0.140) (3.582,-0.100) (3.595,-0.060) (3.608,-0.020) (3.620,0.019) (3.633,0.060) (3.645,0.100) (3.658,0.140) (3.670,0.180) (3.682,0.220) (3.695,0.260) (3.707,0.300) (3.719,0.340) (3.732,0.380) (3.744,0.420) (3.757,0.460) (3.770,0.500) (3.782,0.540) (3.795,0.580) (3.808,0.620) (3.820,0.658) (3.833,0.700) (3.846,0.740) (3.859,0.780) (3.871,0.820) (3.884,0.860) (3.896,0.900) (3.907,0.940) (3.918,0.980) (3.929,1.020) (3.939,1.060) (3.948,1.100) (3.956,1.140) (3.964,1.180) (3.970,1.220) (3.975,1.260) (3.979,1.300) (3.982,1.340) (3.984,1.380) (3.984,1.420) (3.983,1.460) (3.981,1.500) (3.977,1.540) (3.972,1.580) (3.965,1.620) (3.956,1.660) (3.946,1.700) (3.934,1.740) (3.920,1.779) (3.903,1.820) (3.885,1.860) (3.864,1.900) (3.841,1.940) (3.815,1.980) (3.787,2.020) (3.760,2.054) (3.720,2.100) (3.682,2.140) (3.640,2.180) (3.600,2.215) (3.569,2.240) (3.520,2.277) (3.487,2.300) (3.440,2.331) (3.400,2.356) (3.360,2.379) (3.322,2.400) (3.283,2.420) (3.243,2.440) (3.200,2.460) (3.160,2.477) (3.120,2.494) (3.080,2.510) (3.040,2.525) (3.000,2.539) (2.960,2.553) (2.920,2.566) (2.880,2.578) (2.840,2.590) (2.800,2.601) (2.760,2.611) (2.720,2.621) (2.680,2.631) (2.640,2.640) (2.600,2.648) (2.560,2.656) (2.520,2.664) (2.480,2.671) (2.440,2.678) (2.400,2.684) (2.360,2.690) (2.320,2.696) (2.280,2.702) (2.240,2.707) (2.200,2.711) (2.160,2.716) (2.120,2.720) (2.080,2.724) (2.040,2.727) (2.000,2.731) (1.960,2.734) (1.920,2.737) (1.880,2.739) (1.840,2.741) (1.800,2.743) (1.760,2.745) (1.720,2.747) (1.680,2.748) (1.640,2.749) (1.600,2.750) (1.560,2.750) (1.520,2.750) (1.480,2.750) (1.440,2.750) (1.400,2.750) (1.360,2.749) (1.320,2.748) (1.280,2.747) (1.240,2.745) (1.200,2.744) (1.160,2.741) (1.120,2.739) (1.080,2.736) (1.040,2.733) (1.000,2.730) (0.960,2.727) (0.920,2.723) (0.880,2.718) (0.840,2.713) (0.800,2.708) (0.760,2.703) (0.720,2.697) (0.680,2.690) (0.640,2.683) (0.600,2.676) (0.560,2.668) (0.520,2.659) (0.480,2.650) (0.440,2.640) (0.400,2.630) (0.360,2.619) (0.320,2.607) (0.280,2.594) (0.240,2.580) (0.200,2.565) (0.160,2.549) (0.120,2.532) (0.080,2.514) (0.040,2.494) (-0.000,2.472) (-0.040,2.448) (-0.080,2.422) (-0.120,2.392) (-0.159,2.360) (-0.200,2.321) (-0.238,2.280) (-0.270,2.240) (-0.299,2.200) (-0.320,2.166) (-0.340,2.131) (-0.360,2.093) (-0.380,2.050) (-0.400,2.003) (-0.416,1.960) (-0.430,1.920) (-0.443,1.880) (-0.455,1.840) (-0.467,1.800) (-0.477,1.760) (-0.487,1.720) (-0.497,1.680) (-0.506,1.640) (-0.514,1.600) (-0.523,1.560) (-0.531,1.520) (-0.539,1.480) (-0.548,1.440) (-0.556,1.400) (-0.566,1.360) (-0.577,1.320) (-0.590,1.280) (-0.610,1.240) (-0.640,1.206) (-0.677,1.180) (-0.713,1.160) (-0.753,1.140) (-0.794,1.120) (-0.836,1.100) (-0.877,1.080) (-0.916,1.060) (-0.954,1.040) (-0.989,1.020) (-1.040,0.989) (-1.080,0.963) (-1.113,0.940) (-1.160,0.904) (-1.200,0.870) (-1.232,0.840) (-1.260,0.811) (-1.288,0.780) (-1.320,0.740) (-1.349,0.700) (-1.374,0.660) (-1.397,0.620) (-1.416,0.580) (-1.434,0.540) (-1.449,0.500) (-1.462,0.460) (-1.473,0.420) (-1.483,0.380) (-1.491,0.340) (-1.497,0.300) (-1.501,0.260) (-1.505,0.220) (-1.507,0.180) (-1.507,0.140) (-1.507,0.100) (-1.505,0.060) (-1.502,0.020) (-1.497,-0.020) (-1.491,-0.060) (-1.484,-0.100) (-1.476,-0.140) (-1.466,-0.180) (-1.455,-0.220) (-1.443,-0.260) (-1.429,-0.300) (-1.413,-0.340) (-1.395,-0.380) (-1.375,-0.420) (-1.352,-0.460) (-1.325,-0.500) (-1.300,-0.534) (-1.260,-0.578) (-1.220,-0.614) (-1.183,-0.640) (-1.146,-0.660) (-1.100,-0.679) (-1.060,-0.690) (-1.020,-0.696) (-0.980,-0.699) (-0.940,-0.698) (-0.900,-0.695) (-0.860,-0.689) (-0.820,-0.681) (-0.780,-0.671) (-0.740,-0.660) (-0.700,-0.648) (-0.660,-0.634) (-0.620,-0.619) (-0.580,-0.602) (-0.540,-0.585) (-0.500,-0.568) (-0.460,-0.550) (-0.420,-0.530) (-0.380,-0.510) (-0.340,-0.492) (-0.300,-0.472) (-0.260,-0.451) (-0.220,-0.432) (-0.180,-0.409) (-0.140,-0.388) (-0.108,-0.368)}
\EQ{(0.021,-0.419) (-0.015,-0.460) (-0.040,-0.496)}
\EQ{(0.235,-0.185) (0.225,-0.240) (0.210,-0.280) (0.196,-0.320) (0.182,-0.360) (0.169,-0.400) (0.156,-0.440) (0.145,-0.480) (0.134,-0.520) (0.124,-0.560) (0.116,-0.600) (0.109,-0.640) (0.103,-0.680) (0.099,-0.720) (0.096,-0.760) (0.094,-0.800) (0.094,-0.840) (0.096,-0.880) (0.100,-0.920) (0.106,-0.960) (0.114,-1.000) (0.124,-1.040) (0.136,-1.080) (0.151,-1.120) (0.168,-1.160) (0.189,-1.200) (0.213,-1.240) (0.240,-1.279) (0.272,-1.320) (0.300,-1.351) (0.328,-1.380) (0.360,-1.409) (0.397,-1.440) (0.440,-1.472) (0.480,-1.499) (0.516,-1.520) (0.552,-1.540) (0.592,-1.560) (0.637,-1.580) (0.680,-1.597) (0.720,-1.612) (0.760,-1.625) (0.800,-1.636) (0.840,-1.646) (0.880,-1.655) (0.920,-1.663) (0.960,-1.670) (1.000,-1.675) (1.040,-1.680) (1.080,-1.683) (1.120,-1.686) (1.160,-1.687) (1.200,-1.687) (1.240,-1.687) (1.280,-1.685) (1.320,-1.683) (1.360,-1.679) (1.400,-1.675) (1.440,-1.669) (1.480,-1.663) (1.520,-1.656) (1.560,-1.648) (1.600,-1.638) (1.640,-1.628) (1.680,-1.617) (1.720,-1.605) (1.760,-1.592) (1.800,-1.578) (1.840,-1.562) (1.880,-1.546) (1.920,-1.529) (1.960,-1.510) (2.000,-1.490) (2.040,-1.469) (2.080,-1.447) (2.120,-1.424) (2.158,-1.400) (2.200,-1.373) (2.240,-1.345) (2.275,-1.320) (2.320,-1.285) (2.351,-1.260) (2.399,-1.220) (2.440,-1.183) (2.480,-1.145) (2.520,-1.105) (2.560,-1.062) (2.598,-1.020) (2.632,-0.980) (2.660,-0.946) (2.696,-0.900) (2.720,-0.867) (2.753,-0.820) (2.780,-0.781) (2.807,-0.740) (2.832,-0.700) (2.856,-0.660) (2.879,-0.620) (2.900,-0.581) (2.920,-0.544) (2.940,-0.505) (2.960,-0.465) (2.980,-0.423) (3.000,-0.380) (3.017,-0.340) (3.035,-0.300) (3.051,-0.260) (3.067,-0.220) (3.083,-0.180) (3.098,-0.140) (3.113,-0.100) (3.127,-0.060) (3.141,-0.020) (3.155,0.020) (3.169,0.060) (3.182,0.100) (3.195,0.140) (3.208,0.180) (3.222,0.220) (3.235,0.260) (3.248,0.300) (3.261,0.340) (3.275,0.380) (3.289,0.420) (3.303,0.460) (3.318,0.500) (3.333,0.540) (3.350,0.580) (3.366,0.620) (3.384,0.660) (3.403,0.700) (3.422,0.740) (3.443,0.780) (3.465,0.820) (3.488,0.860) (3.512,0.900) (3.537,0.940) (3.560,0.976) (3.588,1.020) (3.614,1.060) (3.638,1.100) (3.660,1.138) (3.680,1.176) (3.700,1.220) (3.715,1.260) (3.727,1.300) (3.736,1.340) (3.742,1.380) (3.746,1.420) (3.747,1.460) (3.746,1.500) (3.742,1.540) (3.737,1.580) (3.728,1.620) (3.717,1.660) (3.701,1.700) (3.681,1.740) (3.655,1.780) (3.621,1.820) (3.580,1.858) (3.540,1.889) (3.500,1.915) (3.460,1.938) (3.420,1.959) (3.380,1.977) (3.340,1.995) (3.300,2.011) (3.260,2.025) (3.220,2.039) (3.180,2.052) (3.140,2.064) (3.100,2.075) (3.060,2.086) (3.020,2.096) (2.980,2.106) (2.940,2.115) (2.900,2.124) (2.860,2.132) (2.820,2.139) (2.780,2.147) (2.740,2.154) (2.700,2.160) (2.660,2.166) (2.620,2.172) (2.580,2.178) (2.540,2.183) (2.500,2.188) (2.460,2.193) (2.420,2.197) (2.380,2.201) (2.340,2.205) (2.300,2.209) (2.260,2.213) (2.220,2.216) (2.180,2.219) (2.140,2.222) (2.100,2.225) (2.060,2.227) (2.020,2.229) (1.980,2.231) (1.940,2.233) (1.900,2.235) (1.860,2.237) (1.820,2.238) (1.780,2.239) (1.740,2.240) (1.700,2.241) (1.660,2.242) (1.620,2.243) (1.580,2.243) (1.540,2.243) (1.500,2.243) (1.460,2.243) (1.420,2.243) (1.380,2.243) (1.340,2.242) (1.300,2.241) (1.260,2.240) (1.220,2.239) (1.180,2.238) (1.140,2.236) (1.100,2.234) (1.060,2.232) (1.020,2.230) (0.980,2.227) (0.940,2.224) (0.900,2.220) (0.860,2.216) (0.820,2.212) (0.780,2.207) (0.740,2.201) (0.700,2.195) (0.660,2.188) (0.621,2.180) (0.580,2.171) (0.540,2.161) (0.500,2.151) (0.460,2.139) (0.420,2.126) (0.380,2.111) (0.340,2.095) (0.300,2.078) (0.263,2.060) (0.224,2.040) (0.189,2.020) (0.140,1.990) (0.100,1.962) (0.060,1.932) (0.022,1.900) (-0.020,1.862) (-0.060,1.822) (-0.097,1.780) (-0.130,1.740) (-0.160,1.701) (-0.189,1.660) (-0.215,1.620) (-0.240,1.580) (-0.264,1.540) (-0.286,1.500) (-0.307,1.460) (-0.328,1.420) (-0.348,1.380) (-0.368,1.340) (-0.389,1.300) (-0.410,1.260) (-0.431,1.220) (-0.454,1.180) (-0.479,1.140) (-0.505,1.100) (-0.533,1.060) (-0.560,1.022) (-0.591,0.980) (-0.619,0.940) (-0.647,0.900) (-0.673,0.860) (-0.696,0.820) (-0.718,0.780) (-0.738,0.740) (-0.755,0.700) (-0.770,0.660) (-0.783,0.620) (-0.793,0.580) (-0.801,0.540) (-0.807,0.500) (-0.811,0.460) (-0.812,0.420) (-0.811,0.380) (-0.808,0.340) (-0.801,0.300) (-0.793,0.260) (-0.781,0.220) (-0.767,0.180) (-0.749,0.140) (-0.727,0.100) (-0.701,0.060) (-0.670,0.020) (-0.640,-0.012) (-0.610,-0.040) (-0.560,-0.079) (-0.520,-0.104) (-0.480,-0.125) (-0.440,-0.143) (-0.400,-0.157) (-0.360,-0.169) (-0.320,-0.178) (-0.280,-0.186) (-0.240,-0.191) (-0.200,-0.194) (-0.160,-0.196) (-0.120,-0.196) (-0.080,-0.194) (-0.040,-0.192) (-0.000,-0.188) (0.040,-0.183) (0.080,-0.178) (0.120,-0.172) (0.160,-0.166) (0.200,-0.162) (0.216,-0.164)}
\EQ{(3.560,1.342) (3.565,1.400) (3.567,1.440) (3.567,1.480) (3.567,1.520) (3.566,1.560) (3.564,1.600) (3.560,1.640) (3.560,1.641)}
\EQ{(3.513,1.700) (3.460,1.709) (3.420,1.713) (3.380,1.717) (3.340,1.720) (3.300,1.723) (3.260,1.725) (3.220,1.728) (3.180,1.730) (3.140,1.732) (3.100,1.734) (3.060,1.735) (3.020,1.737) (2.980,1.738) (2.940,1.740) (2.900,1.741) (2.860,1.743) (2.820,1.744) (2.780,1.745) (2.740,1.746) (2.700,1.747) (2.660,1.748) (2.620,1.749) (2.580,1.750) (2.540,1.751) (2.500,1.751) (2.460,1.752) (2.420,1.753) (2.380,1.754) (2.340,1.754) (2.300,1.755) (2.260,1.755) (2.220,1.756) (2.180,1.756) (2.140,1.757) (2.100,1.757) (2.060,1.758) (2.020,1.758) (1.980,1.758) (1.940,1.759) (1.900,1.759) (1.860,1.759) (1.820,1.759) (1.780,1.760) (1.740,1.760) (1.700,1.760) (1.660,1.760) (1.620,1.760) (1.580,1.760) (1.540,1.760) (1.500,1.760) (1.460,1.760) (1.420,1.760) (1.380,1.760) (1.340,1.760) (1.300,1.761) (1.260,1.761) (1.220,1.761) (1.180,1.761) (1.140,1.762) (1.100,1.762) (1.060,1.764) (1.020,1.767) (0.980,1.772) (0.940,1.777) (0.900,1.783) (0.860,1.788) (0.820,1.792) (0.780,1.795) (0.740,1.797) (0.700,1.797) (0.660,1.795) (0.620,1.792) (0.580,1.787) (0.540,1.780) (0.500,1.772) (0.460,1.762) (0.420,1.750) (0.380,1.736) (0.340,1.720) (0.300,1.702) (0.260,1.682) (0.222,1.660) (0.180,1.634) (0.140,1.606) (0.106,1.580) (0.060,1.542) (0.020,1.505) (-0.020,1.464) (-0.059,1.420) (-0.093,1.380) (-0.120,1.345) (-0.153,1.300) (-0.180,1.261) (-0.207,1.220) (-0.232,1.180) (-0.257,1.140) (-0.280,1.101) (-0.304,1.060) (-0.326,1.020) (-0.348,0.980) (-0.368,0.940) (-0.387,0.900) (-0.405,0.860) (-0.422,0.820) (-0.437,0.780) (-0.450,0.740) (-0.461,0.700) (-0.471,0.660) (-0.478,0.620) (-0.483,0.580) (-0.486,0.540) (-0.487,0.500) (-0.485,0.460) (-0.481,0.420) (-0.474,0.380) (-0.464,0.340) (-0.452,0.300) (-0.436,0.260) (-0.416,0.220) (-0.392,0.180) (-0.364,0.140) (-0.329,0.100) (-0.300,0.071) (-0.263,0.040) (-0.220,0.009) (-0.180,-0.015) (-0.140,-0.035) (-0.100,-0.052) (-0.060,-0.067) (-0.020,-0.079) (0.020,-0.090) (0.060,-0.099) (0.100,-0.107) (0.140,-0.114) (0.180,-0.123) (0.220,-0.135) (0.260,-0.178) (0.269,-0.220) (0.270,-0.260) (0.269,-0.300) (0.268,-0.340) (0.268,-0.380) (0.268,-0.420) (0.270,-0.460) (0.272,-0.500) (0.277,-0.540) (0.282,-0.580) (0.290,-0.620) (0.299,-0.660) (0.310,-0.700) (0.324,-0.740) (0.340,-0.780) (0.358,-0.820) (0.380,-0.860) (0.405,-0.900) (0.433,-0.940) (0.460,-0.974) (0.500,-1.018) (0.540,-1.055) (0.580,-1.088) (0.620,-1.117) (0.655,-1.140) (0.700,-1.166) (0.740,-1.186) (0.780,-1.204) (0.820,-1.220) (0.860,-1.233) (0.900,-1.245) (0.940,-1.255) (0.980,-1.264) (1.020,-1.271) (1.060,-1.277) (1.100,-1.281) (1.140,-1.283) (1.180,-1.285) (1.220,-1.285) (1.260,-1.284) (1.300,-1.281) (1.340,-1.278) (1.380,-1.273) (1.420,-1.266) (1.460,-1.259) (1.500,-1.250) (1.540,-1.240) (1.580,-1.229) (1.620,-1.216) (1.660,-1.203) (1.700,-1.187) (1.740,-1.171) (1.780,-1.153) (1.820,-1.133) (1.860,-1.113) (1.900,-1.090) (1.940,-1.066) (1.980,-1.040) (2.020,-1.013) (2.060,-0.984) (2.100,-0.952) (2.139,-0.920) (2.180,-0.883) (2.220,-0.845) (2.260,-0.805) (2.300,-0.761) (2.335,-0.720) (2.368,-0.680) (2.398,-0.640) (2.427,-0.600) (2.455,-0.560) (2.480,-0.521) (2.505,-0.480) (2.528,-0.440) (2.550,-0.400) (2.571,-0.360) (2.591,-0.320) (2.610,-0.280) (2.628,-0.240) (2.645,-0.200) (2.660,-0.161) (2.676,-0.120) (2.690,-0.080) (2.703,-0.040) (2.716,-0.000) (2.728,0.040) (2.740,0.080) (2.751,0.120) (2.761,0.160) (2.770,0.200) (2.779,0.240) (2.788,0.280) (2.796,0.320) (2.804,0.360) (2.811,0.400) (2.818,0.440) (2.825,0.480) (2.832,0.520) (2.838,0.560) (2.844,0.600) (2.850,0.640) (2.857,0.680) (2.863,0.720) (2.870,0.760) (2.878,0.800) (2.886,0.840) (2.895,0.880) (2.906,0.920) (2.919,0.960) (2.935,1.000) (2.956,1.040) (2.980,1.077) (3.019,1.120) (3.060,1.152) (3.100,1.176) (3.140,1.195) (3.180,1.210) (3.220,1.224) (3.260,1.235) (3.300,1.245) (3.340,1.254) (3.380,1.263) (3.420,1.272) (3.460,1.280) (3.500,1.290) (3.530,1.300)}
\EQ{(0.980,-0.903) (1.020,-0.909) (1.060,-0.914) (1.100,-0.917) (1.140,-0.918) (1.180,-0.918) (1.220,-0.916) (1.260,-0.912) (1.300,-0.907) (1.340,-0.900) (1.380,-0.891) (1.420,-0.881) (1.460,-0.869) (1.500,-0.856) (1.540,-0.840) (1.580,-0.823) (1.620,-0.804) (1.660,-0.784) (1.700,-0.761) (1.740,-0.736) (1.780,-0.708) (1.818,-0.680) (1.860,-0.646) (1.900,-0.611) (1.932,-0.580) (1.971,-0.540) (2.000,-0.508) (2.040,-0.460) (2.070,-0.420) (2.098,-0.380) (2.120,-0.346) (2.148,-0.300) (2.170,-0.260) (2.190,-0.220) (2.209,-0.180) (2.226,-0.140) (2.242,-0.100) (2.256,-0.060) (2.269,-0.020) (2.281,0.020) (2.291,0.060) (2.300,0.100) (2.308,0.140) (2.314,0.180) (2.319,0.220) (2.323,0.260) (2.325,0.300) (2.326,0.340) (2.325,0.380) (2.323,0.420) (2.319,0.460) (2.314,0.500) (2.307,0.540) (2.298,0.580) (2.287,0.620) (2.273,0.660) (2.257,0.700) (2.238,0.740) (2.215,0.780) (2.189,0.820) (2.160,0.857) (2.121,0.900) (2.080,0.937) (2.040,0.969) (2.000,0.997) (1.962,1.020) (1.926,1.040) (1.886,1.060) (1.842,1.080) (1.800,1.097) (1.760,1.111) (1.720,1.125) (1.680,1.137) (1.640,1.148) (1.600,1.159) (1.560,1.168) (1.520,1.177) (1.480,1.185) (1.440,1.193) (1.400,1.200) (1.360,1.206) (1.320,1.213) (1.280,1.219) (1.240,1.224) (1.200,1.230) (1.160,1.237) (1.120,1.244) (1.080,1.253) (1.040,1.269) (1.000,1.295) (0.960,1.330) (0.925,1.360) (0.880,1.394) (0.840,1.417) (0.800,1.435) (0.760,1.448) (0.720,1.457) (0.680,1.463) (0.640,1.465) (0.600,1.465) (0.560,1.461) (0.520,1.455) (0.480,1.447) (0.440,1.436) (0.400,1.422) (0.360,1.406) (0.320,1.387) (0.280,1.365) (0.240,1.340) (0.200,1.312) (0.160,1.281) (0.120,1.246) (0.080,1.207) (0.040,1.164) (0.004,1.120) (-0.027,1.080) (-0.054,1.040) (-0.080,1.000) (-0.103,0.960) (-0.125,0.920) (-0.144,0.880) (-0.160,0.843) (-0.177,0.800) (-0.190,0.760) (-0.201,0.720) (-0.211,0.680) (-0.218,0.640) (-0.223,0.600) (-0.226,0.560) (-0.226,0.520) (-0.225,0.480) (-0.220,0.440) (-0.213,0.400) (-0.203,0.360) (-0.190,0.320) (-0.174,0.280) (-0.154,0.240) (-0.130,0.200) (-0.101,0.160) (-0.067,0.120) (-0.025,0.080) (0.020,0.043) (0.060,0.015) (0.100,-0.010) (0.140,-0.033) (0.180,-0.056) (0.220,-0.080) (0.260,-0.111) (0.299,-0.160) (0.319,-0.200) (0.333,-0.240) (0.346,-0.280) (0.358,-0.320) (0.371,-0.360) (0.386,-0.400) (0.402,-0.440) (0.420,-0.480) (0.440,-0.518) (0.465,-0.560) (0.492,-0.600) (0.520,-0.637) (0.559,-0.680) (0.600,-0.720) (0.640,-0.753) (0.678,-0.780) (0.720,-0.807) (0.760,-0.828) (0.800,-0.847) (0.840,-0.864) (0.880,-0.877) (0.920,-0.889) (0.960,-0.899) (0.980,-0.903)}
\EQ{(0.920,-0.640) (0.960,-0.648) (1.000,-0.653) (1.040,-0.656) (1.080,-0.657) (1.120,-0.655) (1.160,-0.652) (1.200,-0.647) (1.239,-0.640) (1.280,-0.631) (1.320,-0.619) (1.360,-0.606) (1.400,-0.591) (1.440,-0.573) (1.480,-0.553) (1.520,-0.530) (1.560,-0.505) (1.595,-0.480) (1.640,-0.444) (1.680,-0.409) (1.710,-0.380) (1.740,-0.348) (1.780,-0.300) (1.810,-0.260) (1.837,-0.220) (1.860,-0.183) (1.880,-0.147) (1.900,-0.107) (1.920,-0.061) (1.936,-0.020) (1.949,0.020) (1.961,0.060) (1.970,0.100) (1.978,0.140) (1.984,0.180) (1.988,0.220) (1.991,0.260) (1.991,0.300) (1.990,0.340) (1.987,0.380) (1.981,0.420) (1.974,0.460) (1.964,0.500) (1.952,0.540) (1.938,0.580) (1.920,0.620) (1.900,0.660) (1.877,0.700) (1.849,0.740) (1.820,0.777) (1.782,0.820) (1.740,0.860) (1.700,0.893) (1.664,0.920) (1.620,0.950) (1.580,0.974) (1.540,0.995) (1.500,1.015) (1.460,1.033) (1.420,1.050) (1.380,1.065) (1.340,1.079) (1.300,1.092) (1.260,1.104) (1.220,1.115) (1.180,1.126) (1.140,1.137) (1.100,1.148) (1.061,1.160) (1.020,1.173) (0.980,1.187) (0.940,1.202) (0.900,1.216) (0.860,1.229) (0.820,1.240) (0.780,1.248) (0.740,1.254) (0.700,1.257) (0.660,1.258) (0.620,1.255) (0.580,1.250) (0.540,1.242) (0.500,1.231) (0.460,1.218) (0.420,1.201) (0.380,1.181) (0.343,1.160) (0.300,1.132) (0.260,1.102) (0.220,1.067) (0.180,1.027) (0.140,0.981) (0.108,0.940) (0.081,0.900) (0.060,0.864) (0.040,0.827) (0.020,0.784) (0.003,0.740) (-0.011,0.700) (-0.021,0.660) (-0.030,0.620) (-0.036,0.580) (-0.039,0.540) (-0.040,0.500) (-0.038,0.460) (-0.034,0.420) (-0.027,0.380) (-0.017,0.340) (-0.004,0.300) (0.012,0.260) (0.032,0.220) (0.056,0.180) (0.080,0.146) (0.118,0.100) (0.156,0.060) (0.198,0.020) (0.240,-0.018) (0.280,-0.056) (0.320,-0.100) (0.350,-0.140) (0.376,-0.180) (0.400,-0.220) (0.424,-0.260) (0.448,-0.300) (0.475,-0.340) (0.500,-0.374) (0.539,-0.420) (0.577,-0.460) (0.620,-0.498) (0.660,-0.528) (0.700,-0.554) (0.740,-0.576) (0.780,-0.595) (0.820,-0.612) (0.860,-0.625) (0.900,-0.636) (0.920,-0.640)}
\EQ{(0.940,-0.462) (0.980,-0.465) (1.020,-0.465) (1.060,-0.463) (1.100,-0.459) (1.140,-0.452) (1.180,-0.443) (1.220,-0.431) (1.260,-0.417) (1.300,-0.400) (1.340,-0.380) (1.375,-0.360) (1.420,-0.331) (1.460,-0.301) (1.500,-0.266) (1.540,-0.226) (1.580,-0.180) (1.610,-0.140) (1.637,-0.100) (1.660,-0.060) (1.680,-0.020) (1.697,0.020) (1.712,0.060) (1.724,0.100) (1.734,0.140) (1.742,0.180) (1.747,0.220) (1.750,0.260) (1.751,0.300) (1.750,0.340) (1.747,0.380) (1.741,0.420) (1.733,0.460) (1.722,0.500) (1.708,0.540) (1.692,0.580) (1.673,0.620) (1.650,0.660) (1.623,0.700) (1.593,0.740) (1.560,0.777) (1.520,0.816) (1.480,0.850) (1.441,0.880) (1.400,0.908) (1.360,0.932) (1.320,0.953) (1.280,0.973) (1.240,0.991) (1.200,1.007) (1.160,1.022) (1.120,1.036) (1.080,1.049) (1.040,1.061) (1.000,1.072) (0.960,1.082) (0.920,1.091) (0.880,1.099) (0.840,1.105) (0.800,1.109) (0.760,1.111) (0.720,1.111) (0.680,1.108) (0.640,1.102) (0.600,1.094) (0.560,1.083) (0.520,1.068) (0.480,1.051) (0.440,1.030) (0.400,1.005) (0.365,0.980) (0.320,0.942) (0.280,0.902) (0.244,0.860) (0.220,0.827) (0.190,0.780) (0.168,0.740) (0.150,0.700) (0.135,0.660) (0.123,0.620) (0.113,0.580) (0.107,0.540) (0.103,0.500) (0.102,0.460) (0.104,0.420) (0.108,0.380) (0.115,0.340) (0.126,0.300) (0.139,0.260) (0.156,0.220) (0.177,0.180) (0.200,0.141) (0.228,0.100) (0.259,0.060) (0.293,0.020) (0.320,-0.012) (0.360,-0.059) (0.394,-0.100) (0.420,-0.133) (0.458,-0.180) (0.492,-0.220) (0.520,-0.251) (0.550,-0.280) (0.580,-0.307) (0.620,-0.339) (0.660,-0.366) (0.700,-0.389) (0.740,-0.409) (0.780,-0.425) (0.820,-0.439) (0.860,-0.449) (0.900,-0.457) (0.940,-0.462)}
\EQ{(0.900,-0.320) (0.940,-0.322) (0.980,-0.321) (1.020,-0.317) (1.060,-0.310) (1.100,-0.301) (1.140,-0.288) (1.180,-0.273) (1.220,-0.254) (1.260,-0.232) (1.300,-0.205) (1.333,-0.180) (1.377,-0.140) (1.415,-0.100) (1.440,-0.069) (1.474,-0.020) (1.497,0.020) (1.516,0.060) (1.533,0.100) (1.546,0.140) (1.556,0.180) (1.564,0.220) (1.569,0.260) (1.571,0.300) (1.571,0.340) (1.568,0.380) (1.563,0.420) (1.555,0.460) (1.544,0.500) (1.530,0.540) (1.512,0.580) (1.491,0.620) (1.467,0.660) (1.440,0.696) (1.403,0.740) (1.362,0.780) (1.320,0.816) (1.288,0.840) (1.240,0.871) (1.200,0.894) (1.160,0.913) (1.120,0.931) (1.080,0.946) (1.040,0.959) (1.000,0.971) (0.961,0.980) (0.920,0.988) (0.880,0.993) (0.840,0.996) (0.800,0.997) (0.760,0.996) (0.720,0.992) (0.680,0.985) (0.640,0.976) (0.600,0.963) (0.560,0.947) (0.520,0.927) (0.480,0.903) (0.440,0.874) (0.400,0.840) (0.361,0.800) (0.328,0.760) (0.301,0.720) (0.280,0.684) (0.260,0.643) (0.243,0.600) (0.231,0.560) (0.222,0.520) (0.216,0.480) (0.214,0.440) (0.214,0.400) (0.217,0.360) (0.223,0.320) (0.232,0.280) (0.244,0.240) (0.260,0.200) (0.279,0.160) (0.300,0.122) (0.327,0.080) (0.356,0.040) (0.387,-0.000) (0.420,-0.039) (0.456,-0.080) (0.494,-0.120) (0.536,-0.160) (0.580,-0.196) (0.614,-0.220) (0.660,-0.248) (0.700,-0.268) (0.740,-0.285) (0.780,-0.298) (0.820,-0.309) (0.860,-0.316) (0.900,-0.320)}

    % Disco cargado, resto neutros - líneas de campo
% Generado por scripts/generar_datos.py (no editar a mano).
\LC{(-4.174,-3.059) (-4.201,-3.082) (-4.231,-3.103) (-4.262,-3.121) (-4.294,-3.138) (-4.326,-3.153) (-4.360,-3.167) (-4.393,-3.180) (-4.427,-3.192) (-4.461,-3.204) (-4.496,-3.214) (-4.530,-3.224) (-4.565,-3.233) (-4.600,-3.242) (-4.635,-3.249) (-4.670,-3.257) (-4.706,-3.263) (-4.741,-3.269) (-4.777,-3.274) (-4.813,-3.278) (-4.849,-3.282) (-4.884,-3.285) (-4.920,-3.287) (-4.950,-3.288)}
\LC{(-3.897,-3.290) (-3.920,-3.317) (-3.944,-3.344) (-3.969,-3.371) (-3.993,-3.397) (-4.018,-3.423) (-4.044,-3.448) (-4.069,-3.474) (-4.095,-3.499) (-4.121,-3.523) (-4.147,-3.548) (-4.174,-3.572) (-4.201,-3.597) (-4.227,-3.621) (-4.254,-3.644) (-4.282,-3.668) (-4.309,-3.691) (-4.336,-3.715) (-4.364,-3.738) (-4.392,-3.760) (-4.420,-3.783) (-4.448,-3.805) (-4.477,-3.827) (-4.505,-3.849) (-4.534,-3.871) (-4.564,-3.892) (-4.593,-3.912) (-4.623,-3.932) (-4.653,-3.952) (-4.684,-3.970) (-4.715,-3.988) (-4.747,-4.005) (-4.780,-4.021) (-4.813,-4.035) (-4.847,-4.047) (-4.881,-4.058) (-4.916,-4.066) (-4.950,-4.071)}
\LC{(-3.565,-3.567) (-3.587,-3.595) (-3.609,-3.624) (-3.630,-3.653) (-3.651,-3.682) (-3.672,-3.711) (-3.692,-3.741) (-3.712,-3.771) (-3.731,-3.801) (-3.750,-3.832) (-3.768,-3.863) (-3.786,-3.894) (-3.803,-3.926) (-3.820,-3.958) (-3.835,-3.991) (-3.850,-4.023) (-3.865,-4.056) (-3.878,-4.090) (-3.891,-4.123) (-3.903,-4.157) (-3.914,-4.192) (-3.924,-4.226) (-3.933,-4.261) (-3.940,-4.296) (-3.947,-4.332) (-3.952,-4.367) (-3.956,-4.403) (-3.959,-4.439) (-3.960,-4.450)}
\LC{(-3.275,-3.808) (-3.297,-3.837) (-3.316,-3.867) (-3.334,-3.898) (-3.351,-3.930) (-3.367,-3.962) (-3.381,-3.995) (-3.394,-4.029) (-3.407,-4.063) (-3.418,-4.097) (-3.428,-4.132) (-3.437,-4.166) (-3.445,-4.201) (-3.452,-4.237) (-3.458,-4.272) (-3.463,-4.308) (-3.468,-4.344) (-3.471,-4.379) (-3.474,-4.415) (-3.475,-4.450)}
\LC{(-3.085,-3.967) (-3.091,-4.002) (-3.096,-4.038) (-3.099,-4.074) (-3.101,-4.109) (-3.103,-4.145) (-3.105,-4.181) (-3.106,-4.217) (-3.107,-4.253) (-3.108,-4.289) (-3.109,-4.325) (-3.109,-4.361) (-3.110,-4.397) (-3.110,-4.433) (-3.110,-4.450)}
\LC{(-2.877,-3.803) (-2.851,-3.829) (-2.828,-3.857) (-2.807,-3.886) (-2.788,-3.916) (-2.770,-3.948) (-2.754,-3.980) (-2.739,-4.013) (-2.726,-4.046) (-2.714,-4.080) (-2.703,-4.114) (-2.693,-4.149) (-2.685,-4.184) (-2.677,-4.219) (-2.671,-4.254) (-2.665,-4.290) (-2.660,-4.326) (-2.657,-4.362) (-2.654,-4.397) (-2.652,-4.433) (-2.651,-4.450)}
\LC{(-2.604,-3.489) (-2.577,-3.513) (-2.552,-3.539) (-2.528,-3.565) (-2.504,-3.593) (-2.481,-3.621) (-2.459,-3.649) (-2.438,-3.678) (-2.418,-3.708) (-2.399,-3.739) (-2.381,-3.770) (-2.363,-3.801) (-2.347,-3.833) (-2.331,-3.866) (-2.316,-3.898) (-2.302,-3.932) (-2.289,-3.965) (-2.277,-3.999) (-2.266,-4.033) (-2.256,-4.068) (-2.246,-4.102) (-2.237,-4.137) (-2.229,-4.172) (-2.222,-4.208) (-2.216,-4.243) (-2.211,-4.279) (-2.206,-4.314) (-2.202,-4.350) (-2.199,-4.386) (-2.197,-4.422) (-2.196,-4.450)}
\LC{(-2.163,-2.981) (-2.137,-3.006) (-2.112,-3.031) (-2.086,-3.057) (-2.062,-3.083) (-2.038,-3.110) (-2.014,-3.137) (-1.992,-3.165) (-1.969,-3.194) (-1.948,-3.222) (-1.927,-3.252) (-1.906,-3.281) (-1.887,-3.311) (-1.867,-3.342) (-1.849,-3.373) (-1.831,-3.404) (-1.814,-3.435) (-1.797,-3.467) (-1.781,-3.500) (-1.765,-3.532) (-1.750,-3.565) (-1.736,-3.598) (-1.722,-3.631) (-1.709,-3.665) (-1.697,-3.698) (-1.685,-3.732) (-1.673,-3.766) (-1.662,-3.801) (-1.652,-3.835) (-1.642,-3.870) (-1.633,-3.905) (-1.624,-3.940) (-1.616,-3.975) (-1.609,-4.010) (-1.602,-4.045) (-1.595,-4.081) (-1.589,-4.116) (-1.584,-4.152) (-1.579,-4.187) (-1.575,-4.223) (-1.571,-4.259) (-1.568,-4.295) (-1.565,-4.331) (-1.563,-4.367) (-1.561,-4.403) (-1.560,-4.439) (-1.560,-4.450)}
\LC{(-1.349,-2.043) (-1.324,-2.068) (-1.300,-2.095) (-1.277,-2.122) (-1.255,-2.151) (-1.234,-2.180) (-1.214,-2.210) (-1.195,-2.241) (-1.177,-2.272) (-1.159,-2.303) (-1.143,-2.335) (-1.127,-2.368) (-1.112,-2.400) (-1.097,-2.433) (-1.083,-2.466) (-1.069,-2.500) (-1.056,-2.533) (-1.043,-2.567) (-1.031,-2.601) (-1.019,-2.635) (-1.007,-2.669) (-0.996,-2.703) (-0.985,-2.737) (-0.974,-2.771) (-0.964,-2.806) (-0.953,-2.840) (-0.943,-2.875) (-0.933,-2.910) (-0.924,-2.944) (-0.915,-2.979) (-0.905,-3.014) (-0.897,-3.049) (-0.888,-3.084) (-0.879,-3.119) (-0.871,-3.154) (-0.863,-3.189) (-0.855,-3.224) (-0.847,-3.259) (-0.840,-3.294) (-0.833,-3.330) (-0.826,-3.365) (-0.819,-3.400) (-0.812,-3.436) (-0.806,-3.471) (-0.799,-3.506) (-0.793,-3.542) (-0.787,-3.577) (-0.782,-3.613) (-0.776,-3.649) (-0.771,-3.684) (-0.766,-3.720) (-0.761,-3.756) (-0.757,-3.791) (-0.752,-3.827) (-0.748,-3.863) (-0.744,-3.898) (-0.740,-3.934) (-0.737,-3.970) (-0.734,-4.006) (-0.731,-4.042) (-0.728,-4.078) (-0.725,-4.114) (-0.723,-4.150) (-0.721,-4.186) (-0.719,-4.221) (-0.717,-4.257) (-0.716,-4.293) (-0.714,-4.329) (-0.713,-4.365) (-0.713,-4.401) (-0.712,-4.437) (-0.712,-4.450)}
\LC{(-2.983,-2.123) (-3.003,-2.094) (-3.024,-2.064) (-3.045,-2.035) (-3.067,-2.007) (-3.090,-1.979) (-3.113,-1.951) (-3.136,-1.924) (-3.161,-1.897) (-3.185,-1.871) (-3.211,-1.846) (-3.236,-1.820) (-3.263,-1.796) (-3.289,-1.772) (-3.317,-1.748) (-3.344,-1.725) (-3.372,-1.703) (-3.401,-1.681) (-3.430,-1.659) (-3.459,-1.638) (-3.489,-1.618) (-3.518,-1.598) (-3.549,-1.578) (-3.579,-1.559) (-3.610,-1.541) (-3.642,-1.523) (-3.673,-1.506) (-3.705,-1.489) (-3.737,-1.473) (-3.769,-1.457) (-3.802,-1.441) (-3.835,-1.426) (-3.868,-1.412) (-3.901,-1.398) (-3.934,-1.385) (-3.968,-1.372) (-4.002,-1.359) (-4.036,-1.347) (-4.070,-1.336) (-4.104,-1.325) (-4.138,-1.314) (-4.173,-1.304) (-4.208,-1.294) (-4.242,-1.285) (-4.277,-1.276) (-4.312,-1.268) (-4.347,-1.260) (-4.383,-1.252) (-4.418,-1.245) (-4.453,-1.238) (-4.489,-1.232) (-4.524,-1.226) (-4.560,-1.221) (-4.595,-1.216) (-4.631,-1.212) (-4.667,-1.208) (-4.703,-1.204) (-4.739,-1.201) (-4.775,-1.198) (-4.810,-1.196) (-4.846,-1.194) (-4.882,-1.192) (-4.918,-1.191) (-4.950,-1.190)}
\LC{(-3.751,-2.600) (-3.772,-2.570) (-3.793,-2.541) (-3.816,-2.513) (-3.839,-2.486) (-3.864,-2.459) (-3.889,-2.434) (-3.915,-2.409) (-3.942,-2.385) (-3.969,-2.361) (-3.997,-2.339) (-4.026,-2.317) (-4.055,-2.297) (-4.085,-2.277) (-4.116,-2.258) (-4.147,-2.239) (-4.178,-2.222) (-4.210,-2.205) (-4.243,-2.190) (-4.276,-2.175) (-4.309,-2.161) (-4.342,-2.147) (-4.376,-2.135) (-4.410,-2.123) (-4.444,-2.112) (-4.479,-2.101) (-4.513,-2.092) (-4.548,-2.083) (-4.583,-2.075) (-4.618,-2.067) (-4.654,-2.061) (-4.689,-2.055) (-4.725,-2.049) (-4.761,-2.044) (-4.796,-2.040) (-4.832,-2.037) (-4.868,-2.034) (-4.904,-2.032) (-4.940,-2.031) (-4.950,-2.030)}
\LC{(-4.160,-2.853) (-4.184,-2.827) (-4.212,-2.803) (-4.241,-2.783) (-4.272,-2.765) (-4.305,-2.749) (-4.338,-2.734) (-4.371,-2.721) (-4.405,-2.710) (-4.440,-2.700) (-4.475,-2.690) (-4.510,-2.682) (-4.545,-2.675) (-4.580,-2.668) (-4.616,-2.662) (-4.651,-2.657) (-4.687,-2.652) (-4.723,-2.648) (-4.759,-2.644) (-4.795,-2.641) (-4.830,-2.639) (-4.866,-2.637) (-4.902,-2.636) (-4.938,-2.635) (-4.950,-2.634)}
\LC{(-0.858,2.885) (-0.858,2.921) (-0.857,2.957) (-0.857,2.993) (-0.856,3.029) (-0.855,3.065) (-0.855,3.101) (-0.854,3.137) (-0.854,3.173) (-0.854,3.209) (-0.854,3.245) (-0.854,3.281) (-0.854,3.317) (-0.854,3.353) (-0.855,3.389) (-0.855,3.425) (-0.856,3.461) (-0.856,3.497) (-0.857,3.533) (-0.858,3.569) (-0.858,3.605) (-0.859,3.641) (-0.860,3.677) (-0.861,3.713) (-0.861,3.749) (-0.862,3.785) (-0.863,3.821) (-0.863,3.857) (-0.864,3.893) (-0.865,3.929) (-0.865,3.965) (-0.866,4.001) (-0.867,4.037) (-0.867,4.073) (-0.868,4.109) (-0.868,4.145) (-0.869,4.181) (-0.869,4.217) (-0.869,4.253) (-0.870,4.289) (-0.870,4.325) (-0.870,4.361) (-0.870,4.397) (-0.871,4.433) (-0.871,4.450)}
\LC{(-1.716,2.885) (-1.716,2.921) (-1.717,2.957) (-1.717,2.993) (-1.718,3.029) (-1.719,3.065) (-1.720,3.101) (-1.721,3.137) (-1.722,3.173) (-1.723,3.209) (-1.725,3.245) (-1.726,3.281) (-1.728,3.317) (-1.729,3.353) (-1.731,3.389) (-1.733,3.425) (-1.735,3.461) (-1.737,3.497) (-1.739,3.533) (-1.741,3.568) (-1.743,3.604) (-1.745,3.640) (-1.746,3.676) (-1.748,3.712) (-1.750,3.748) (-1.752,3.784) (-1.754,3.820) (-1.756,3.856) (-1.758,3.892) (-1.760,3.928) (-1.761,3.964) (-1.763,4.000) (-1.764,4.036) (-1.766,4.072) (-1.767,4.108) (-1.768,4.144) (-1.769,4.180) (-1.770,4.216) (-1.771,4.252) (-1.772,4.288) (-1.773,4.324) (-1.773,4.360) (-1.774,4.396) (-1.774,4.432) (-1.774,4.450)}
\LC{(-2.485,2.885) (-2.488,2.921) (-2.491,2.957) (-2.496,2.992) (-2.502,3.028) (-2.508,3.063) (-2.515,3.099) (-2.523,3.134) (-2.532,3.169) (-2.541,3.204) (-2.550,3.238) (-2.560,3.273) (-2.570,3.308) (-2.580,3.342) (-2.590,3.377) (-2.600,3.411) (-2.610,3.446) (-2.621,3.480) (-2.631,3.515) (-2.641,3.549) (-2.651,3.584) (-2.660,3.619) (-2.670,3.653) (-2.679,3.688) (-2.688,3.723) (-2.697,3.758) (-2.706,3.793) (-2.714,3.828) (-2.722,3.863) (-2.730,3.898) (-2.737,3.933) (-2.744,3.969) (-2.751,4.004) (-2.757,4.040) (-2.763,4.075) (-2.768,4.111) (-2.773,4.146) (-2.778,4.182) (-2.782,4.218) (-2.785,4.254) (-2.788,4.290) (-2.791,4.325) (-2.793,4.361) (-2.795,4.397) (-2.796,4.433) (-2.797,4.450)}
\LC{(-2.785,2.594) (-2.821,2.597) (-2.857,2.601) (-2.892,2.607) (-2.927,2.615) (-2.962,2.623) (-2.997,2.633) (-3.031,2.644) (-3.065,2.655) (-3.099,2.668) (-3.133,2.681) (-3.166,2.695) (-3.199,2.709) (-3.232,2.724) (-3.265,2.739) (-3.297,2.754) (-3.329,2.770) (-3.361,2.787) (-3.393,2.803) (-3.425,2.820) (-3.457,2.837) (-3.488,2.855) (-3.520,2.872) (-3.551,2.890) (-3.583,2.907) (-3.614,2.925) (-3.645,2.943) (-3.676,2.961) (-3.707,2.980) (-3.738,2.998) (-3.769,3.016) (-3.800,3.034) (-3.831,3.053) (-3.862,3.071) (-3.893,3.089) (-3.924,3.108) (-3.955,3.126) (-3.986,3.144) (-4.018,3.162) (-4.049,3.180) (-4.080,3.198) (-4.111,3.215) (-4.143,3.233) (-4.174,3.250) (-4.206,3.267) (-4.238,3.284) (-4.270,3.301) (-4.302,3.317) (-4.334,3.333) (-4.366,3.349) (-4.399,3.364) (-4.432,3.379) (-4.465,3.393) (-4.498,3.407) (-4.532,3.421) (-4.565,3.433) (-4.599,3.445) (-4.633,3.457) (-4.668,3.467) (-4.703,3.477) (-4.737,3.485) (-4.773,3.493) (-4.808,3.500) (-4.843,3.506) (-4.879,3.510) (-4.915,3.514) (-4.950,3.516)}
\LC{(-2.785,1.698) (-2.821,1.698) (-2.857,1.698) (-2.893,1.697) (-2.929,1.697) (-2.965,1.697) (-3.001,1.697) (-3.037,1.697) (-3.073,1.697) (-3.109,1.697) (-3.145,1.697) (-3.181,1.697) (-3.217,1.697) (-3.253,1.698) (-3.289,1.699) (-3.325,1.700) (-3.361,1.701) (-3.397,1.702) (-3.433,1.703) (-3.469,1.704) (-3.505,1.706) (-3.541,1.708) (-3.577,1.710) (-3.613,1.712) (-3.649,1.714) (-3.685,1.716) (-3.720,1.718) (-3.756,1.721) (-3.792,1.723) (-3.828,1.726) (-3.864,1.728) (-3.900,1.731) (-3.936,1.734) (-3.972,1.736) (-4.008,1.739) (-4.044,1.742) (-4.080,1.745) (-4.115,1.747) (-4.151,1.750) (-4.187,1.753) (-4.223,1.755) (-4.259,1.758) (-4.295,1.760) (-4.331,1.763) (-4.367,1.765) (-4.403,1.767) (-4.439,1.769) (-4.475,1.772) (-4.511,1.773) (-4.547,1.775) (-4.582,1.777) (-4.618,1.779) (-4.654,1.780) (-4.690,1.781) (-4.726,1.783) (-4.762,1.784) (-4.798,1.785) (-4.834,1.785) (-4.870,1.786) (-4.906,1.787) (-4.942,1.787) (-4.950,1.787)}
\LC{(1.361,0.418) (1.396,0.423) (1.432,0.429) (1.467,0.435) (1.502,0.442) (1.538,0.450) (1.573,0.458) (1.608,0.466) (1.643,0.475) (1.677,0.485) (1.712,0.495) (1.746,0.505) (1.781,0.516) (1.815,0.527) (1.849,0.539) (1.883,0.551) (1.917,0.563) (1.950,0.576) (1.984,0.589) (2.017,0.603) (2.050,0.616) (2.084,0.630) (2.117,0.645) (2.150,0.659) (2.182,0.674) (2.215,0.689) (2.248,0.703) (2.281,0.718) (2.314,0.733) (2.347,0.748) (2.380,0.762) (2.413,0.776) (2.446,0.790) (2.480,0.803) (2.513,0.815) (2.548,0.826) (2.582,0.837) (2.617,0.845) (2.652,0.853) (2.688,0.858) (2.724,0.861) (2.760,0.862) (2.796,0.861) (2.832,0.857) (2.867,0.851) (2.902,0.843) (2.937,0.834) (2.971,0.823) (3.005,0.811) (3.039,0.798) (3.072,0.784) (3.105,0.770) (3.138,0.755) (3.171,0.740) (3.203,0.725) (3.236,0.710) (3.269,0.695) (3.301,0.680) (3.334,0.665) (3.367,0.650) (3.400,0.635) (3.433,0.621) (3.466,0.607) (3.499,0.593) (3.533,0.580) (3.566,0.567) (3.600,0.554) (3.634,0.542) (3.667,0.529) (3.702,0.518) (3.736,0.506) (3.770,0.495) (3.804,0.485) (3.839,0.474) (3.873,0.464) (3.908,0.455) (3.943,0.446) (3.978,0.437) (4.013,0.428) (4.048,0.420) (4.083,0.412) (4.118,0.404) (4.153,0.397) (4.189,0.390) (4.224,0.384) (4.260,0.377) (4.295,0.371) (4.331,0.366) (4.366,0.360) (4.402,0.355) (4.438,0.351) (4.473,0.346) (4.509,0.342) (4.545,0.338) (4.581,0.335) (4.616,0.331) (4.652,0.328) (4.688,0.326) (4.724,0.323) (4.760,0.321) (4.796,0.320) (4.832,0.318) (4.868,0.317) (4.904,0.316) (4.940,0.315) (4.950,0.315)}
\LC{(1.331,0.535) (1.364,0.549) (1.397,0.564) (1.429,0.579) (1.461,0.596) (1.493,0.613) (1.524,0.630) (1.555,0.649) (1.586,0.668) (1.615,0.688) (1.645,0.709) (1.674,0.731) (1.702,0.753) (1.729,0.776) (1.756,0.800) (1.782,0.825) (1.807,0.851) (1.832,0.877) (1.855,0.905) (1.877,0.933) (1.898,0.962) (1.918,0.992) (1.937,1.023) (1.954,1.055) (1.969,1.087) (1.982,1.121) (1.994,1.155) (2.004,1.190) (2.011,1.225) (2.016,1.260) (2.019,1.296) (2.020,1.320)}
\LC{(1.274,0.642) (1.303,0.664) (1.331,0.687) (1.358,0.710) (1.385,0.734) (1.411,0.759) (1.436,0.785) (1.460,0.811) (1.484,0.838) (1.507,0.867) (1.528,0.895) (1.549,0.925) (1.568,0.955) (1.586,0.987) (1.602,1.019) (1.617,1.051) (1.631,1.085) (1.643,1.119) (1.653,1.153) (1.661,1.188) (1.667,1.224) (1.672,1.259) (1.674,1.295) (1.675,1.331)}
\LC{(1.204,0.724) (1.228,0.751) (1.251,0.779) (1.273,0.807) (1.295,0.836) (1.315,0.865) (1.335,0.896) (1.353,0.926) (1.371,0.958) (1.387,0.990) (1.402,1.023) (1.416,1.056) (1.428,1.090) (1.438,1.124) (1.447,1.159) (1.455,1.194) (1.460,1.230) (1.464,1.266) (1.466,1.302) (1.466,1.326)}
\LC{(1.130,0.782) (1.149,0.813) (1.167,0.844) (1.185,0.875) (1.202,0.907) (1.218,0.939) (1.233,0.972) (1.247,1.005) (1.260,1.039) (1.272,1.072) (1.283,1.107) (1.292,1.142) (1.300,1.177) (1.307,1.212) (1.311,1.248) (1.314,1.284) (1.315,1.320) (1.315,1.332)}
\LC{(1.047,0.826) (1.060,0.860) (1.073,0.893) (1.085,0.927) (1.098,0.961) (1.110,0.994) (1.122,1.029) (1.133,1.063) (1.143,1.097) (1.153,1.132) (1.162,1.167) (1.170,1.202) (1.177,1.237) (1.181,1.273) (1.183,1.309) (1.183,1.321)}
\LC{(0.957,0.854) (0.965,0.889) (0.973,0.924) (0.981,0.959) (0.990,0.994) (0.998,1.029) (1.008,1.064) (1.017,1.098) (1.027,1.133) (1.038,1.167) (1.049,1.202) (1.061,1.236) (1.072,1.270) (1.081,1.305) (1.082,1.329)}
\LC{(0.863,0.865) (0.864,0.901) (0.867,0.937) (0.869,0.973) (0.873,1.008) (0.877,1.044) (0.882,1.080) (0.889,1.115) (0.896,1.150) (0.906,1.185) (0.917,1.219) (0.931,1.252) (0.948,1.284) (0.969,1.313) (0.995,1.338) (1.026,1.356) (1.062,1.363)}
\LC{(0.769,0.859) (0.764,0.894) (0.760,0.930) (0.756,0.966) (0.753,1.002) (0.751,1.038) (0.750,1.074) (0.749,1.110) (0.750,1.146) (0.752,1.182) (0.755,1.217) (0.760,1.253) (0.766,1.289) (0.773,1.324) (0.783,1.359) (0.794,1.393) (0.806,1.427) (0.820,1.460) (0.836,1.492) (0.853,1.524) (0.871,1.555) (0.890,1.586) (0.908,1.617) (0.926,1.648) (0.941,1.681) (0.952,1.715) (0.960,1.750) (0.965,1.786) (0.969,1.821) (0.972,1.857) (0.973,1.893) (0.974,1.929) (0.974,1.965) (0.974,2.001) (0.973,2.037) (0.972,2.073) (0.970,2.109) (0.968,2.145) (0.966,2.181) (0.963,2.217) (0.960,2.253) (0.957,2.289) (0.954,2.325) (0.951,2.360) (0.947,2.396) (0.944,2.432) (0.940,2.468) (0.936,2.504) (0.933,2.540) (0.929,2.575) (0.925,2.611) (0.921,2.647) (0.917,2.683) (0.914,2.719) (0.910,2.754) (0.906,2.790) (0.902,2.826) (0.899,2.862) (0.895,2.898) (0.892,2.933) (0.888,2.969) (0.885,3.005) (0.882,3.041) (0.879,3.077) (0.876,3.113) (0.873,3.149) (0.870,3.184) (0.867,3.220) (0.865,3.256) (0.862,3.292) (0.860,3.328) (0.857,3.364) (0.855,3.400) (0.853,3.436) (0.851,3.472) (0.849,3.508) (0.847,3.544) (0.845,3.580) (0.844,3.616) (0.842,3.652) (0.841,3.688) (0.839,3.724) (0.838,3.760) (0.837,3.795) (0.835,3.831) (0.834,3.867) (0.833,3.903) (0.832,3.939) (0.831,3.975) (0.831,4.011) (0.830,4.047) (0.829,4.083) (0.829,4.119) (0.828,4.155) (0.828,4.191) (0.827,4.227) (0.827,4.263) (0.826,4.299) (0.826,4.335) (0.826,4.371) (0.826,4.407) (0.826,4.443) (0.826,4.450)}
\LC{(0.665,0.831) (0.653,0.864) (0.641,0.898) (0.629,0.932) (0.618,0.967) (0.607,1.001) (0.596,1.035) (0.586,1.070) (0.576,1.105) (0.567,1.139) (0.558,1.174) (0.550,1.209) (0.541,1.244) (0.533,1.279) (0.526,1.314) (0.518,1.350) (0.510,1.385) (0.503,1.420) (0.495,1.455) (0.487,1.490) (0.479,1.525) (0.471,1.560) (0.462,1.595) (0.453,1.630) (0.443,1.665) (0.433,1.699) (0.422,1.734) (0.411,1.768) (0.399,1.802) (0.387,1.836) (0.375,1.870) (0.361,1.903) (0.348,1.936) (0.334,1.970) (0.320,2.003) (0.305,2.035) (0.290,2.068) (0.274,2.101) (0.259,2.133) (0.243,2.165) (0.227,2.198) (0.211,2.230) (0.194,2.262) (0.178,2.294) (0.161,2.326) (0.144,2.358) (0.128,2.390) (0.111,2.422) (0.095,2.454) (0.079,2.486) (0.063,2.518) (0.047,2.551) (0.032,2.583) (0.017,2.616) (0.003,2.649) (-0.011,2.682) (-0.024,2.716) (-0.037,2.749) (-0.049,2.784) (-0.060,2.818) (-0.070,2.852) (-0.079,2.887) (-0.088,2.922) (-0.096,2.957) (-0.103,2.992) (-0.110,3.028) (-0.116,3.063) (-0.122,3.099) (-0.127,3.134) (-0.132,3.170) (-0.137,3.206) (-0.141,3.241) (-0.145,3.277) (-0.148,3.313) (-0.152,3.349) (-0.155,3.385) (-0.158,3.421) (-0.161,3.456) (-0.163,3.492) (-0.166,3.528) (-0.168,3.564) (-0.170,3.600) (-0.172,3.636) (-0.174,3.672) (-0.176,3.708) (-0.178,3.744) (-0.180,3.780) (-0.181,3.816) (-0.183,3.852) (-0.184,3.888) (-0.185,3.924) (-0.187,3.960) (-0.188,3.996) (-0.189,4.032) (-0.190,4.068) (-0.190,4.104) (-0.191,4.140) (-0.192,4.176) (-0.193,4.212) (-0.193,4.248) (-0.194,4.284) (-0.194,4.320) (-0.194,4.356) (-0.195,4.392) (-0.195,4.428) (-0.195,4.450)}
\LC{(0.581,0.789) (0.561,0.819) (0.543,0.850) (0.524,0.881) (0.506,0.912) (0.487,0.943) (0.469,0.974) (0.451,1.005) (0.433,1.036) (0.415,1.068) (0.397,1.099) (0.379,1.130) (0.361,1.161) (0.343,1.192) (0.325,1.223) (0.307,1.254) (0.288,1.285) (0.269,1.316) (0.250,1.346) (0.230,1.376) (0.211,1.406) (0.190,1.436) (0.169,1.466) (0.148,1.495) (0.126,1.523) (0.104,1.552) (0.081,1.579) (0.058,1.607) (0.034,1.634) (0.009,1.660) (-0.016,1.686) (-0.041,1.711) (-0.068,1.735) (-0.095,1.759) (-0.122,1.782) (-0.151,1.805) (-0.180,1.826) (-0.209,1.847) (-0.239,1.866) (-0.270,1.885) (-0.302,1.902) (-0.334,1.919) (-0.366,1.934) (-0.400,1.948) (-0.433,1.960) (-0.468,1.971) (-0.503,1.980) (-0.538,1.987) (-0.573,1.993) (-0.609,1.996) (-0.645,1.998) (-0.657,1.998)}
\LC{(0.495,0.723) (0.470,0.749) (0.445,0.775) (0.420,0.801) (0.395,0.827) (0.370,0.852) (0.345,0.878) (0.319,0.904) (0.294,0.929) (0.268,0.955) (0.243,0.980) (0.217,1.005) (0.191,1.030) (0.165,1.055) (0.139,1.080) (0.113,1.105) (0.086,1.129) (0.060,1.153) (0.033,1.177) (0.005,1.200) (-0.022,1.224) (-0.050,1.247) (-0.078,1.269) (-0.106,1.291) (-0.135,1.313) (-0.165,1.333) (-0.194,1.354) (-0.224,1.374) (-0.255,1.392) (-0.286,1.411) (-0.318,1.428) (-0.350,1.444) (-0.383,1.459) (-0.416,1.472) (-0.450,1.485) (-0.484,1.495) (-0.519,1.504) (-0.554,1.511) (-0.590,1.516) (-0.626,1.519) (-0.662,1.520)}
\LC{(0.425,0.641) (0.396,0.662) (0.366,0.682) (0.336,0.702) (0.305,0.721) (0.275,0.740) (0.244,0.759) (0.214,0.778) (0.183,0.797) (0.152,0.815) (0.121,0.833) (0.090,0.851) (0.058,0.869) (0.027,0.887) (-0.004,0.905) (-0.036,0.922) (-0.067,0.940) (-0.098,0.957) (-0.130,0.975) (-0.161,0.993) (-0.193,1.010) (-0.224,1.028) (-0.256,1.045) (-0.287,1.063) (-0.318,1.081) (-0.349,1.099) (-0.380,1.117) (-0.411,1.135) (-0.442,1.154) (-0.473,1.173) (-0.504,1.191) (-0.534,1.210) (-0.565,1.229) (-0.596,1.248) (-0.628,1.264) (-0.651,1.272)}
\LC{(0.374,0.547) (0.340,0.559) (0.306,0.572) (0.273,0.584) (0.238,0.595) (0.204,0.606) (0.170,0.617) (0.135,0.627) (0.101,0.637) (0.066,0.646) (0.031,0.655) (-0.004,0.664) (-0.039,0.673) (-0.074,0.682) (-0.109,0.690) (-0.144,0.699) (-0.179,0.707) (-0.214,0.716) (-0.248,0.725) (-0.283,0.734) (-0.318,0.744) (-0.352,0.754) (-0.387,0.764) (-0.421,0.776) (-0.455,0.788) (-0.489,0.801) (-0.522,0.814) (-0.555,0.829) (-0.587,0.845) (-0.618,0.862) (-0.649,0.881) (-0.679,0.901) (-0.708,0.923) (-0.735,0.946) (-0.762,0.971) (-0.786,0.997) (-0.808,1.026) (-0.828,1.055) (-0.846,1.087) (-0.861,1.119) (-0.873,1.153) (-0.882,1.188) (-0.887,1.224) (-0.889,1.260)}
\LC{(0.346,0.457) (0.311,0.463) (0.275,0.469) (0.240,0.473) (0.204,0.477) (0.168,0.480) (0.132,0.482) (0.096,0.484) (0.060,0.485) (0.024,0.485) (-0.012,0.484) (-0.048,0.483) (-0.084,0.481) (-0.120,0.479) (-0.156,0.477) (-0.192,0.474) (-0.228,0.470) (-0.263,0.467) (-0.299,0.463) (-0.335,0.460) (-0.371,0.456) (-0.407,0.452) (-0.442,0.448) (-0.478,0.444) (-0.514,0.441) (-0.550,0.438) (-0.586,0.435) (-0.622,0.432) (-0.658,0.430) (-0.694,0.428) (-0.730,0.427) (-0.766,0.426) (-0.802,0.426) (-0.838,0.427) (-0.874,0.428) (-0.909,0.430) (-0.945,0.432) (-0.981,0.436) (-1.017,0.440) (-1.053,0.445) (-1.088,0.451) (-1.123,0.458) (-1.159,0.465) (-1.194,0.473) (-1.229,0.483) (-1.263,0.493) (-1.297,0.503) (-1.332,0.515) (-1.365,0.527) (-1.399,0.540) (-1.432,0.554) (-1.465,0.569) (-1.498,0.584) (-1.530,0.599) (-1.562,0.616) (-1.594,0.632) (-1.626,0.650) (-1.657,0.667) (-1.688,0.685) (-1.719,0.704) (-1.750,0.722) (-1.781,0.741) (-1.811,0.760) (-1.842,0.779) (-1.872,0.798) (-1.903,0.817) (-1.934,0.835) (-1.965,0.853) (-1.997,0.871) (-2.029,0.887) (-2.062,0.902) (-2.095,0.915) (-2.129,0.926) (-2.164,0.935) (-2.200,0.940) (-2.236,0.942) (-2.272,0.940) (-2.308,0.936) (-2.343,0.929) (-2.378,0.920) (-2.412,0.910) (-2.447,0.899) (-2.481,0.887) (-2.515,0.875) (-2.548,0.863) (-2.582,0.851) (-2.616,0.839) (-2.651,0.828) (-2.685,0.817) (-2.719,0.807) (-2.754,0.797) (-2.789,0.788) (-2.824,0.780) (-2.859,0.772) (-2.894,0.764) (-2.930,0.757) (-2.965,0.751) (-3.001,0.745) (-3.036,0.740) (-3.072,0.735) (-3.108,0.731) (-3.143,0.727) (-3.179,0.723) (-3.215,0.720) (-3.251,0.717) (-3.287,0.715) (-3.323,0.712) (-3.359,0.710) (-3.395,0.709) (-3.431,0.707) (-3.467,0.706) (-3.503,0.705) (-3.539,0.704) (-3.575,0.704) (-3.611,0.703) (-3.647,0.703) (-3.683,0.703) (-3.719,0.703) (-3.755,0.703) (-3.791,0.704) (-3.827,0.704) (-3.863,0.704) (-3.899,0.705) (-3.935,0.706) (-3.971,0.706) (-4.007,0.707) (-4.043,0.708) (-4.079,0.708) (-4.115,0.709) (-4.151,0.710) (-4.186,0.711) (-4.222,0.712) (-4.258,0.713) (-4.294,0.714) (-4.330,0.714) (-4.366,0.715) (-4.402,0.716) (-4.438,0.717) (-4.474,0.718) (-4.510,0.718) (-4.546,0.719) (-4.582,0.720) (-4.618,0.720) (-4.654,0.721) (-4.690,0.721) (-4.726,0.722) (-4.762,0.722) (-4.798,0.723) (-4.834,0.723) (-4.870,0.723) (-4.906,0.723) (-4.942,0.724) (-4.950,0.724)}
\LC{(0.335,0.363) (0.299,0.362) (0.263,0.359) (0.227,0.356) (0.192,0.351) (0.156,0.345) (0.121,0.338) (0.086,0.330) (0.051,0.320) (0.017,0.309) (-0.017,0.297) (-0.050,0.284) (-0.083,0.269) (-0.116,0.253) (-0.147,0.236) (-0.178,0.218) (-0.209,0.199) (-0.239,0.179) (-0.267,0.157) (-0.295,0.135) (-0.323,0.111) (-0.349,0.086) (-0.374,0.060) (-0.398,0.034) (-0.421,0.006) (-0.443,-0.023) (-0.463,-0.052) (-0.482,-0.083) (-0.500,-0.114) (-0.516,-0.146) (-0.530,-0.179) (-0.543,-0.213) (-0.554,-0.247) (-0.563,-0.282) (-0.569,-0.318) (-0.574,-0.353) (-0.576,-0.389) (-0.576,-0.425) (-0.573,-0.461) (-0.567,-0.497) (-0.559,-0.532) (-0.547,-0.566) (-0.534,-0.599) (-0.518,-0.632) (-0.513,-0.642)}
\LC{(0.339,0.282) (0.304,0.276) (0.269,0.267) (0.235,0.257) (0.201,0.245) (0.167,0.232) (0.134,0.217) (0.102,0.200) (0.072,0.182) (0.042,0.162) (0.013,0.140) (-0.014,0.116) (-0.039,0.091) (-0.063,0.064) (-0.085,0.035) (-0.104,0.004) (-0.121,-0.027) (-0.135,-0.060) (-0.147,-0.094) (-0.155,-0.129) (-0.161,-0.165) (-0.163,-0.201) (-0.161,-0.237) (-0.156,-0.272) (-0.148,-0.307) (-0.136,-0.341) (-0.119,-0.373) (-0.105,-0.393)}
\LC{(0.356,0.203) (0.323,0.190) (0.290,0.176) (0.258,0.160) (0.226,0.142) (0.196,0.122) (0.168,0.100) (0.141,0.076) (0.116,0.050) (0.093,0.022) (0.074,-0.008) (0.057,-0.040) (0.045,-0.074) (0.036,-0.109) (0.033,-0.144) (0.034,-0.180) (0.041,-0.216) (0.052,-0.250) (0.067,-0.283)}
\LC{(0.385,0.128) (0.354,0.110) (0.323,0.091) (0.294,0.071) (0.265,0.049) (0.239,0.025) (0.214,-0.001) (0.191,-0.029) (0.172,-0.060) (0.158,-0.093) (0.149,-0.128) (0.149,-0.164) (0.157,-0.199) (0.165,-0.221)}
\LC{(0.418,0.069) (0.389,0.048) (0.360,0.026) (0.332,0.004) (0.305,-0.020) (0.279,-0.044) (0.254,-0.071) (0.232,-0.099) (0.215,-0.130) (0.205,-0.165) (0.206,-0.189)}
\LC{(0.450,0.026) (0.423,0.002) (0.395,-0.021) (0.367,-0.044) (0.340,-0.067) (0.312,-0.090) (0.285,-0.114) (0.259,-0.138) (0.233,-0.164) (0.209,-0.190)}
\LC{(0.496,-0.024) (0.470,-0.049) (0.443,-0.073) (0.416,-0.097) (0.388,-0.119) (0.359,-0.140) (0.329,-0.160) (0.297,-0.177) (0.263,-0.189) (0.228,-0.193) (0.206,-0.195)}
\LC{(0.548,-0.067) (0.524,-0.094) (0.499,-0.120) (0.473,-0.145) (0.446,-0.168) (0.417,-0.190) (0.387,-0.209) (0.354,-0.225) (0.321,-0.238) (0.285,-0.245) (0.250,-0.245) (0.215,-0.236) (0.184,-0.218)}
\LC{(0.617,-0.109) (0.597,-0.139) (0.576,-0.169) (0.554,-0.197) (0.530,-0.223) (0.504,-0.248) (0.476,-0.271) (0.447,-0.292) (0.415,-0.310) (0.383,-0.325) (0.349,-0.337) (0.314,-0.344) (0.278,-0.347) (0.242,-0.345) (0.206,-0.338) (0.173,-0.326) (0.141,-0.309) (0.113,-0.287)}
\LC{(0.679,-0.136) (0.664,-0.168) (0.647,-0.200) (0.629,-0.231) (0.609,-0.261) (0.587,-0.290) (0.563,-0.317) (0.538,-0.343) (0.511,-0.367) (0.483,-0.388) (0.452,-0.408) (0.421,-0.425) (0.388,-0.440) (0.354,-0.451) (0.319,-0.459) (0.283,-0.464) (0.247,-0.466) (0.211,-0.464) (0.176,-0.458) (0.141,-0.448) (0.108,-0.435) (0.076,-0.417) (0.047,-0.395) (0.038,-0.388)}
\LC{(0.770,-0.159) (0.762,-0.194) (0.752,-0.228) (0.741,-0.263) (0.728,-0.296) (0.714,-0.329) (0.698,-0.362) (0.680,-0.393) (0.661,-0.424) (0.641,-0.453) (0.619,-0.482) (0.596,-0.509) (0.571,-0.535) (0.545,-0.560) (0.517,-0.583) (0.489,-0.605) (0.459,-0.625) (0.428,-0.644) (0.396,-0.660) (0.363,-0.675) (0.329,-0.688) (0.295,-0.698) (0.260,-0.707) (0.224,-0.713) (0.189,-0.717) (0.153,-0.718) (0.117,-0.717) (0.081,-0.714) (0.045,-0.708) (0.011,-0.699) (-0.024,-0.688) (-0.057,-0.674) (-0.089,-0.657) (-0.120,-0.639) (-0.149,-0.618) (-0.159,-0.611)}
\LC{(0.864,-0.165) (0.863,-0.201) (0.860,-0.237) (0.857,-0.273) (0.852,-0.308) (0.845,-0.344) (0.838,-0.379) (0.829,-0.414) (0.820,-0.448) (0.809,-0.483) (0.797,-0.517) (0.783,-0.550) (0.769,-0.583) (0.753,-0.616) (0.737,-0.647) (0.719,-0.679) (0.700,-0.710) (0.681,-0.740) (0.660,-0.769) (0.639,-0.798) (0.617,-0.827) (0.594,-0.854) (0.570,-0.881) (0.545,-0.907) (0.519,-0.933) (0.493,-0.958) (0.467,-0.982) (0.439,-1.005) (0.411,-1.027) (0.382,-1.049) (0.353,-1.070) (0.323,-1.091) (0.293,-1.110) (0.262,-1.129) (0.231,-1.147) (0.200,-1.164) (0.168,-1.181) (0.135,-1.196) (0.103,-1.211) (0.069,-1.225) (0.036,-1.239) (0.002,-1.252) (-0.032,-1.263) (-0.066,-1.274) (-0.100,-1.285) (-0.135,-1.294) (-0.170,-1.302) (-0.205,-1.310) (-0.241,-1.317) (-0.276,-1.322) (-0.312,-1.327) (-0.348,-1.330) (-0.384,-1.332) (-0.420,-1.333) (-0.456,-1.332) (-0.492,-1.330) (-0.527,-1.326) (-0.563,-1.319) (-0.598,-1.310) (-0.631,-1.297) (-0.664,-1.282) (-0.694,-1.263) (-0.713,-1.248)}
\LC{(0.971,-0.151) (0.978,-0.186) (0.985,-0.221) (0.990,-0.257) (0.994,-0.293) (0.998,-0.328) (1.000,-0.364) (1.002,-0.400) (1.003,-0.436) (1.003,-0.472) (1.002,-0.508) (1.001,-0.544) (0.999,-0.580) (0.996,-0.616) (0.992,-0.652) (0.988,-0.688) (0.983,-0.723) (0.977,-0.759) (0.971,-0.794) (0.964,-0.830) (0.957,-0.865) (0.949,-0.900) (0.941,-0.935) (0.932,-0.970) (0.923,-1.005) (0.914,-1.040) (0.904,-1.074) (0.894,-1.109) (0.884,-1.143) (0.873,-1.178) (0.862,-1.212) (0.851,-1.246) (0.840,-1.280) (0.828,-1.315) (0.817,-1.349) (0.805,-1.383) (0.793,-1.417) (0.782,-1.451) (0.770,-1.485) (0.758,-1.519) (0.746,-1.553) (0.735,-1.587) (0.723,-1.621) (0.712,-1.655) (0.701,-1.690) (0.690,-1.724) (0.679,-1.758) (0.668,-1.793) (0.658,-1.827) (0.648,-1.862) (0.638,-1.896) (0.628,-1.931) (0.619,-1.966) (0.610,-2.001) (0.602,-2.036) (0.594,-2.071) (0.586,-2.106) (0.578,-2.141) (0.571,-2.176) (0.565,-2.212) (0.558,-2.247) (0.552,-2.283) (0.547,-2.318) (0.541,-2.354) (0.537,-2.389) (0.532,-2.425) (0.528,-2.461) (0.524,-2.497) (0.521,-2.533) (0.517,-2.568) (0.514,-2.604) (0.512,-2.640) (0.510,-2.676) (0.508,-2.712) (0.506,-2.748) (0.504,-2.784) (0.503,-2.820) (0.502,-2.856) (0.501,-2.892) (0.501,-2.928) (0.500,-2.964) (0.500,-3.000) (0.500,-3.036) (0.500,-3.072) (0.501,-3.108) (0.501,-3.144) (0.501,-3.180) (0.502,-3.216) (0.503,-3.252) (0.503,-3.288) (0.504,-3.324) (0.505,-3.360) (0.506,-3.396) (0.507,-3.432) (0.509,-3.468) (0.510,-3.504) (0.511,-3.540) (0.512,-3.576) (0.513,-3.612) (0.514,-3.648) (0.516,-3.684) (0.517,-3.720) (0.518,-3.756) (0.519,-3.792) (0.520,-3.828) (0.521,-3.864) (0.523,-3.900) (0.524,-3.936) (0.525,-3.972) (0.526,-4.008) (0.526,-4.044) (0.527,-4.080) (0.528,-4.116) (0.529,-4.152) (0.529,-4.188) (0.530,-4.224) (0.531,-4.260) (0.531,-4.296) (0.531,-4.332) (0.532,-4.368) (0.532,-4.404) (0.532,-4.440) (0.532,-4.450)}
\LC{(1.084,-0.109) (1.100,-0.141) (1.115,-0.174) (1.130,-0.207) (1.144,-0.240) (1.157,-0.273) (1.170,-0.307) (1.182,-0.341) (1.194,-0.375) (1.205,-0.409) (1.216,-0.443) (1.226,-0.478) (1.236,-0.512) (1.246,-0.547) (1.255,-0.582) (1.264,-0.617) (1.272,-0.652) (1.280,-0.687) (1.288,-0.722) (1.295,-0.757) (1.302,-0.793) (1.309,-0.828) (1.315,-0.863) (1.322,-0.899) (1.328,-0.934) (1.334,-0.970) (1.340,-1.005) (1.345,-1.041) (1.351,-1.076) (1.356,-1.112) (1.362,-1.148) (1.367,-1.183) (1.372,-1.219) (1.377,-1.255) (1.382,-1.290) (1.387,-1.326) (1.392,-1.362) (1.397,-1.397) (1.401,-1.433) (1.406,-1.469) (1.411,-1.504) (1.416,-1.540) (1.421,-1.576) (1.426,-1.611) (1.431,-1.647) (1.436,-1.682) (1.441,-1.718) (1.447,-1.754) (1.452,-1.789) (1.457,-1.825) (1.462,-1.861) (1.468,-1.896) (1.473,-1.932) (1.479,-1.967) (1.485,-2.003) (1.490,-2.038) (1.496,-2.074) (1.502,-2.109) (1.508,-2.145) (1.514,-2.180) (1.520,-2.216) (1.526,-2.251) (1.533,-2.287) (1.539,-2.322) (1.545,-2.358) (1.552,-2.393) (1.558,-2.429) (1.565,-2.464) (1.571,-2.499) (1.578,-2.535) (1.584,-2.570) (1.591,-2.605) (1.598,-2.641) (1.604,-2.676) (1.611,-2.712) (1.618,-2.747) (1.624,-2.782) (1.631,-2.818) (1.638,-2.853) (1.644,-2.888) (1.651,-2.924) (1.658,-2.959) (1.664,-2.995) (1.671,-3.030) (1.677,-3.065) (1.684,-3.101) (1.690,-3.136) (1.696,-3.172) (1.702,-3.207) (1.709,-3.243) (1.715,-3.278) (1.721,-3.314) (1.726,-3.349) (1.732,-3.385) (1.738,-3.420) (1.743,-3.456) (1.749,-3.491) (1.754,-3.527) (1.759,-3.563) (1.764,-3.598) (1.769,-3.634) (1.774,-3.670) (1.778,-3.705) (1.783,-3.741) (1.787,-3.777) (1.791,-3.813) (1.795,-3.848) (1.799,-3.884) (1.802,-3.920) (1.805,-3.956) (1.809,-3.992) (1.812,-4.028) (1.814,-4.064) (1.817,-4.100) (1.819,-4.135) (1.821,-4.171) (1.823,-4.207) (1.825,-4.243) (1.826,-4.279) (1.828,-4.315) (1.829,-4.351) (1.829,-4.387) (1.830,-4.423) (1.830,-4.450)}
\LC{(1.185,-0.041) (1.208,-0.069) (1.231,-0.096) (1.254,-0.125) (1.276,-0.153) (1.297,-0.182) (1.319,-0.210) (1.340,-0.240) (1.361,-0.269) (1.382,-0.298) (1.402,-0.328) (1.423,-0.358) (1.442,-0.388) (1.462,-0.418) (1.482,-0.448) (1.501,-0.478) (1.520,-0.509) (1.539,-0.540) (1.558,-0.570) (1.576,-0.601) (1.595,-0.632) (1.613,-0.663) (1.631,-0.694) (1.649,-0.725) (1.667,-0.757) (1.685,-0.788) (1.702,-0.819) (1.720,-0.851) (1.738,-0.882) (1.755,-0.914) (1.772,-0.945) (1.790,-0.977) (1.807,-1.008) (1.824,-1.040) (1.842,-1.072) (1.859,-1.103) (1.876,-1.135) (1.893,-1.166) (1.910,-1.198) (1.928,-1.230) (1.945,-1.261) (1.962,-1.293) (1.979,-1.325) (1.996,-1.356) (2.014,-1.388) (2.031,-1.419) (2.048,-1.451) (2.066,-1.482) (2.083,-1.514) (2.100,-1.545) (2.118,-1.577) (2.135,-1.608) (2.153,-1.640) (2.170,-1.671) (2.188,-1.703) (2.206,-1.734) (2.223,-1.765) (2.241,-1.797) (2.259,-1.828) (2.277,-1.859) (2.295,-1.890) (2.313,-1.922) (2.331,-1.953) (2.349,-1.984) (2.367,-2.015) (2.385,-2.046) (2.403,-2.077) (2.421,-2.108) (2.439,-2.140) (2.458,-2.171) (2.476,-2.202) (2.494,-2.233) (2.513,-2.264) (2.531,-2.295) (2.549,-2.325) (2.568,-2.356) (2.586,-2.387) (2.605,-2.418) (2.623,-2.449) (2.642,-2.480) (2.660,-2.511) (2.678,-2.542) (2.697,-2.573) (2.715,-2.604) (2.734,-2.635) (2.752,-2.665) (2.771,-2.696) (2.789,-2.727) (2.808,-2.758) (2.826,-2.789) (2.844,-2.820) (2.862,-2.851) (2.881,-2.882) (2.899,-2.913) (2.917,-2.945) (2.935,-2.976) (2.953,-3.007) (2.971,-3.038) (2.989,-3.069) (3.006,-3.101) (3.024,-3.132) (3.041,-3.164) (3.059,-3.195) (3.076,-3.227) (3.093,-3.258) (3.110,-3.290) (3.127,-3.322) (3.143,-3.354) (3.160,-3.386) (3.176,-3.418) (3.192,-3.450) (3.208,-3.483) (3.224,-3.515) (3.239,-3.548) (3.254,-3.580) (3.269,-3.613) (3.283,-3.646) (3.297,-3.679) (3.311,-3.713) (3.325,-3.746) (3.338,-3.779) (3.351,-3.813) (3.363,-3.847) (3.375,-3.881) (3.386,-3.915) (3.397,-3.949) (3.407,-3.984) (3.417,-4.019) (3.426,-4.053) (3.435,-4.088) (3.443,-4.123) (3.450,-4.159) (3.457,-4.194) (3.463,-4.229) (3.469,-4.265) (3.473,-4.301) (3.477,-4.337) (3.480,-4.372) (3.483,-4.408) (3.484,-4.444) (3.485,-4.450)}
\LC{(1.275,0.059) (1.304,0.038) (1.333,0.017) (1.362,-0.005) (1.391,-0.026) (1.420,-0.047) (1.449,-0.069) (1.478,-0.090) (1.507,-0.112) (1.535,-0.134) (1.564,-0.156) (1.592,-0.178) (1.620,-0.201) (1.648,-0.223) (1.676,-0.245) (1.704,-0.268) (1.732,-0.291) (1.760,-0.314) (1.788,-0.337) (1.815,-0.360) (1.843,-0.383) (1.870,-0.407) (1.898,-0.430) (1.925,-0.453) (1.952,-0.477) (1.979,-0.501) (2.006,-0.525) (2.033,-0.548) (2.060,-0.572) (2.087,-0.596) (2.114,-0.620) (2.141,-0.644) (2.167,-0.668) (2.194,-0.693) (2.221,-0.717) (2.247,-0.741) (2.274,-0.765) (2.301,-0.789) (2.327,-0.814) (2.354,-0.838) (2.380,-0.862) (2.407,-0.886) (2.434,-0.911) (2.460,-0.935) (2.487,-0.959) (2.514,-0.983) (2.540,-1.007) (2.567,-1.031) (2.594,-1.055) (2.621,-1.079) (2.648,-1.103) (2.675,-1.127) (2.702,-1.151) (2.729,-1.175) (2.756,-1.198) (2.783,-1.222) (2.811,-1.245) (2.838,-1.269) (2.865,-1.292) (2.893,-1.315) (2.920,-1.338) (2.948,-1.361) (2.976,-1.384) (3.004,-1.407) (3.032,-1.429) (3.060,-1.452) (3.088,-1.474) (3.116,-1.497) (3.145,-1.519) (3.173,-1.541) (3.202,-1.562) (3.231,-1.584) (3.260,-1.606) (3.289,-1.627) (3.318,-1.648) (3.347,-1.669) (3.377,-1.690) (3.406,-1.710) (3.436,-1.731) (3.466,-1.751) (3.495,-1.771) (3.526,-1.791) (3.556,-1.810) (3.586,-1.830) (3.617,-1.849) (3.647,-1.867) (3.678,-1.886) (3.709,-1.904) (3.740,-1.922) (3.772,-1.940) (3.803,-1.958) (3.835,-1.975) (3.867,-1.992) (3.899,-2.008) (3.931,-2.024) (3.963,-2.040) (3.996,-2.056) (4.028,-2.071) (4.061,-2.086) (4.094,-2.100) (4.127,-2.114) (4.160,-2.128) (4.194,-2.141) (4.228,-2.154) (4.261,-2.166) (4.295,-2.178) (4.330,-2.189) (4.364,-2.200) (4.398,-2.211) (4.433,-2.220) (4.468,-2.230) (4.503,-2.239) (4.538,-2.247) (4.573,-2.255) (4.608,-2.262) (4.644,-2.268) (4.679,-2.274) (4.715,-2.280) (4.750,-2.285) (4.786,-2.289) (4.822,-2.292) (4.858,-2.295) (4.894,-2.297) (4.930,-2.299) (4.950,-2.300)}
\LC{(1.331,0.166) (1.365,0.153) (1.398,0.140) (1.432,0.127) (1.465,0.114) (1.499,0.101) (1.533,0.088) (1.566,0.076) (1.600,0.063) (1.633,0.049) (1.667,0.036) (1.700,0.023) (1.734,0.010) (1.767,-0.003) (1.801,-0.017) (1.834,-0.030) (1.867,-0.044) (1.901,-0.057) (1.934,-0.071) (1.967,-0.085) (2.000,-0.099) (2.033,-0.113) (2.066,-0.128) (2.099,-0.142) (2.132,-0.157) (2.165,-0.171) (2.198,-0.186) (2.231,-0.201) (2.263,-0.217) (2.296,-0.232) (2.328,-0.248) (2.361,-0.263) (2.393,-0.279) (2.425,-0.295) (2.458,-0.311) (2.490,-0.327) (2.522,-0.343) (2.554,-0.360) (2.586,-0.376) (2.618,-0.393) (2.650,-0.409) (2.682,-0.426) (2.714,-0.442) (2.746,-0.459) (2.778,-0.476) (2.809,-0.493) (2.841,-0.510) (2.873,-0.526) (2.905,-0.543) (2.937,-0.560) (2.969,-0.577) (3.000,-0.593) (3.032,-0.610) (3.064,-0.627) (3.096,-0.643) (3.128,-0.660) (3.160,-0.676) (3.193,-0.692) (3.225,-0.708) (3.257,-0.724) (3.289,-0.740) (3.322,-0.756) (3.354,-0.772) (3.387,-0.787) (3.419,-0.803) (3.452,-0.818) (3.484,-0.833) (3.517,-0.848) (3.550,-0.862) (3.583,-0.876) (3.616,-0.891) (3.650,-0.905) (3.683,-0.918) (3.716,-0.932) (3.750,-0.945) (3.783,-0.958) (3.817,-0.971) (3.851,-0.983) (3.885,-0.995) (3.919,-1.007) (3.953,-1.019) (3.987,-1.030) (4.021,-1.041) (4.055,-1.052) (4.090,-1.062) (4.124,-1.072) (4.159,-1.082) (4.194,-1.092) (4.229,-1.101) (4.264,-1.109) (4.299,-1.118) (4.334,-1.125) (4.369,-1.133) (4.404,-1.140) (4.440,-1.147) (4.475,-1.153) (4.511,-1.159) (4.546,-1.165) (4.582,-1.170) (4.617,-1.175) (4.653,-1.179) (4.689,-1.183) (4.725,-1.187) (4.761,-1.190) (4.797,-1.193) (4.832,-1.195) (4.868,-1.197) (4.904,-1.198) (4.940,-1.199) (4.950,-1.199)}
\LC{(1.362,0.297) (1.398,0.294) (1.434,0.291) (1.470,0.288) (1.506,0.286) (1.542,0.283) (1.578,0.281) (1.614,0.279) (1.650,0.278) (1.686,0.276) (1.722,0.275) (1.757,0.273) (1.793,0.272) (1.829,0.271) (1.865,0.270) (1.901,0.269) (1.937,0.268) (1.973,0.267) (2.009,0.266) (2.045,0.265) (2.081,0.264) (2.117,0.262) (2.153,0.261) (2.189,0.259) (2.225,0.257) (2.261,0.255) (2.297,0.252) (2.333,0.249) (2.369,0.246) (2.405,0.243) (2.440,0.239) (2.476,0.234) (2.512,0.229) (2.547,0.224) (2.583,0.218) (2.618,0.212) (2.654,0.205) (2.689,0.198) (2.724,0.190) (2.759,0.182) (2.794,0.173) (2.829,0.164) (2.864,0.155) (2.898,0.145) (2.933,0.135) (2.968,0.125) (3.002,0.114) (3.036,0.104) (3.071,0.093) (3.105,0.081) (3.139,0.070) (3.173,0.059) (3.207,0.047) (3.241,0.035) (3.275,0.024) (3.309,0.012) (3.343,0.000) (3.377,-0.011) (3.412,-0.023) (3.446,-0.035) (3.480,-0.046) (3.514,-0.057) (3.548,-0.069) (3.582,-0.080) (3.617,-0.091) (3.651,-0.102) (3.685,-0.113) (3.720,-0.123) (3.754,-0.134) (3.789,-0.144) (3.823,-0.154) (3.858,-0.164) (3.892,-0.173) (3.927,-0.183) (3.962,-0.192) (3.997,-0.201) (4.032,-0.209) (4.067,-0.218) (4.102,-0.226) (4.137,-0.233) (4.172,-0.241) (4.208,-0.248) (4.243,-0.255) (4.278,-0.262) (4.314,-0.268) (4.349,-0.274) (4.385,-0.280) (4.420,-0.286) (4.456,-0.291) (4.492,-0.296) (4.527,-0.300) (4.563,-0.305) (4.599,-0.308) (4.635,-0.312) (4.671,-0.315) (4.706,-0.318) (4.742,-0.321) (4.778,-0.323) (4.814,-0.325) (4.850,-0.326) (4.886,-0.328) (4.922,-0.329) (4.950,-0.329)}
\LC{(3.366,1.285) (3.373,1.250) (3.384,1.215) (3.399,1.183) (3.417,1.152) (3.438,1.123) (3.462,1.095) (3.487,1.070) (3.514,1.046) (3.543,1.024) (3.572,1.004) (3.603,0.985) (3.634,0.967) (3.666,0.951) (3.699,0.935) (3.732,0.921) (3.765,0.907) (3.799,0.895) (3.833,0.883) (3.867,0.872) (3.902,0.861) (3.936,0.851) (3.971,0.842) (4.006,0.834) (4.041,0.825) (4.076,0.818) (4.111,0.811) (4.147,0.804) (4.182,0.797) (4.218,0.791) (4.253,0.786) (4.289,0.781) (4.325,0.776) (4.360,0.771) (4.396,0.767) (4.432,0.763) (4.468,0.759) (4.503,0.756) (4.539,0.753) (4.575,0.750) (4.611,0.747) (4.647,0.745) (4.683,0.743) (4.719,0.741) (4.755,0.739) (4.791,0.738) (4.827,0.737) (4.863,0.736) (4.899,0.735) (4.935,0.735) (4.950,0.734)}
\LC{(3.585,1.402) (3.621,1.398) (3.657,1.394) (3.692,1.389) (3.728,1.385) (3.764,1.381) (3.800,1.377) (3.835,1.374) (3.871,1.371) (3.907,1.369) (3.943,1.367) (3.979,1.365) (4.015,1.364) (4.051,1.363) (4.087,1.362) (4.123,1.362) (4.159,1.362) (4.195,1.362) (4.231,1.362) (4.267,1.362) (4.303,1.362) (4.339,1.363) (4.375,1.363) (4.411,1.364) (4.447,1.364) (4.483,1.365) (4.519,1.366) (4.555,1.366) (4.591,1.367) (4.627,1.367) (4.663,1.368) (4.699,1.368) (4.735,1.369) (4.771,1.369) (4.807,1.370) (4.843,1.370) (4.879,1.370) (4.915,1.370) (4.950,1.371)}
\LC{(3.585,1.615) (3.620,1.621) (3.655,1.630) (3.690,1.640) (3.724,1.650) (3.759,1.661) (3.793,1.672) (3.827,1.683) (3.862,1.694) (3.896,1.704) (3.931,1.715) (3.965,1.725) (4.000,1.735) (4.034,1.745) (4.069,1.755) (4.104,1.764) (4.138,1.773) (4.173,1.782) (4.208,1.790) (4.243,1.799) (4.279,1.807) (4.314,1.814) (4.349,1.821) (4.384,1.828) (4.420,1.835) (4.455,1.841) (4.491,1.847) (4.526,1.853) (4.562,1.858) (4.598,1.862) (4.633,1.867) (4.669,1.871) (4.705,1.874) (4.741,1.878) (4.777,1.880) (4.813,1.883) (4.848,1.885) (4.884,1.886) (4.920,1.887) (4.950,1.888)}
\LC{(3.422,1.715) (3.428,1.751) (3.437,1.785) (3.449,1.819) (3.463,1.852) (3.479,1.884) (3.497,1.916) (3.516,1.946) (3.537,1.976) (3.558,2.005) (3.581,2.033) (3.604,2.060) (3.628,2.087) (3.652,2.114) (3.677,2.140) (3.703,2.165) (3.729,2.190) (3.755,2.214) (3.782,2.238) (3.809,2.262) (3.837,2.285) (3.865,2.307) (3.893,2.329) (3.922,2.351) (3.951,2.372) (3.981,2.393) (4.011,2.413) (4.041,2.433) (4.071,2.452) (4.102,2.471) (4.133,2.489) (4.164,2.507) (4.196,2.524) (4.228,2.541) (4.260,2.557) (4.292,2.573) (4.325,2.588) (4.358,2.602) (4.391,2.616) (4.425,2.629) (4.458,2.642) (4.492,2.654) (4.527,2.665) (4.561,2.675) (4.596,2.685) (4.630,2.694) (4.665,2.703) (4.701,2.710) (4.736,2.717) (4.772,2.723) (4.807,2.728) (4.843,2.732) (4.879,2.735) (4.915,2.738) (4.950,2.739)}
\LC{(3.069,1.715) (3.071,1.751) (3.073,1.787) (3.077,1.823) (3.081,1.858) (3.087,1.894) (3.093,1.929) (3.100,1.965) (3.108,2.000) (3.117,2.035) (3.126,2.069) (3.137,2.104) (3.148,2.138) (3.160,2.172) (3.172,2.206) (3.185,2.240) (3.198,2.273) (3.212,2.306) (3.227,2.339) (3.242,2.372) (3.257,2.404) (3.273,2.437) (3.289,2.469) (3.306,2.501) (3.323,2.532) (3.341,2.564) (3.358,2.595) (3.376,2.626) (3.395,2.657) (3.413,2.688) (3.432,2.719) (3.452,2.749) (3.471,2.779) (3.491,2.810) (3.511,2.840) (3.531,2.869) (3.551,2.899) (3.572,2.929) (3.592,2.958) (3.613,2.987) (3.634,3.017) (3.656,3.046) (3.677,3.075) (3.699,3.103) (3.720,3.132) (3.742,3.161) (3.764,3.189) (3.786,3.218) (3.809,3.246) (3.831,3.274) (3.854,3.302) (3.876,3.330) (3.899,3.358) (3.922,3.386) (3.945,3.413) (3.968,3.441) (3.991,3.469) (4.014,3.496) (4.038,3.524) (4.061,3.551) (4.085,3.578) (4.108,3.605) (4.132,3.633) (4.155,3.660) (4.179,3.687) (4.203,3.714) (4.227,3.741) (4.250,3.768) (4.274,3.795) (4.298,3.822) (4.322,3.849) (4.346,3.876) (4.370,3.903) (4.394,3.930) (4.417,3.957) (4.441,3.984) (4.465,4.011) (4.488,4.038) (4.512,4.065) (4.535,4.093) (4.558,4.120) (4.581,4.148) (4.604,4.176) (4.626,4.204) (4.648,4.233) (4.669,4.262) (4.689,4.292) (4.709,4.322) (4.727,4.354) (4.742,4.386) (4.755,4.420) (4.764,4.450)}
\LC{(2.580,1.715) (2.581,1.751) (2.582,1.787) (2.583,1.823) (2.585,1.859) (2.587,1.895) (2.590,1.931) (2.593,1.967) (2.596,2.002) (2.600,2.038) (2.604,2.074) (2.608,2.110) (2.613,2.146) (2.618,2.181) (2.623,2.217) (2.629,2.252) (2.634,2.288) (2.641,2.323) (2.647,2.359) (2.654,2.394) (2.661,2.429) (2.668,2.465) (2.676,2.500) (2.683,2.535) (2.691,2.570) (2.700,2.605) (2.708,2.640) (2.716,2.675) (2.725,2.710) (2.734,2.745) (2.743,2.780) (2.752,2.815) (2.761,2.850) (2.770,2.884) (2.780,2.919) (2.789,2.954) (2.799,2.989) (2.808,3.023) (2.818,3.058) (2.828,3.093) (2.837,3.127) (2.847,3.162) (2.857,3.197) (2.867,3.231) (2.876,3.266) (2.886,3.301) (2.896,3.335) (2.905,3.370) (2.914,3.405) (2.924,3.440) (2.933,3.474) (2.942,3.509) (2.951,3.544) (2.960,3.579) (2.969,3.614) (2.977,3.649) (2.985,3.684) (2.993,3.719) (3.001,3.754) (3.009,3.789) (3.016,3.825) (3.023,3.860) (3.030,3.895) (3.037,3.931) (3.043,3.966) (3.049,4.002) (3.054,4.037) (3.059,4.073) (3.064,4.108) (3.069,4.144) (3.073,4.180) (3.076,4.216) (3.080,4.252) (3.082,4.287) (3.085,4.323) (3.087,4.359) (3.088,4.395) (3.089,4.431) (3.090,4.450)}
\LC{(1.996,1.715) (1.996,1.751) (1.996,1.787) (1.997,1.823) (1.997,1.859) (1.998,1.895) (1.999,1.931) (2.000,1.967) (2.001,2.003) (2.002,2.039) (2.003,2.075) (2.005,2.111) (2.007,2.147) (2.008,2.183) (2.010,2.219) (2.012,2.255) (2.014,2.291) (2.016,2.327) (2.019,2.362) (2.021,2.398) (2.023,2.434) (2.026,2.470) (2.029,2.506) (2.031,2.542) (2.034,2.578) (2.037,2.614) (2.040,2.650) (2.043,2.686) (2.046,2.721) (2.049,2.757) (2.052,2.793) (2.056,2.829) (2.059,2.865) (2.062,2.901) (2.066,2.937) (2.069,2.972) (2.072,3.008) (2.076,3.044) (2.079,3.080) (2.083,3.116) (2.086,3.152) (2.090,3.187) (2.093,3.223) (2.097,3.259) (2.100,3.295) (2.104,3.331) (2.107,3.367) (2.110,3.402) (2.114,3.438) (2.117,3.474) (2.120,3.510) (2.123,3.546) (2.126,3.582) (2.129,3.618) (2.132,3.653) (2.135,3.689) (2.138,3.725) (2.141,3.761) (2.143,3.797) (2.146,3.833) (2.148,3.869) (2.151,3.905) (2.153,3.941) (2.155,3.977) (2.157,4.013) (2.159,4.049) (2.161,4.084) (2.162,4.120) (2.164,4.156) (2.165,4.192) (2.166,4.228) (2.167,4.264) (2.168,4.300) (2.169,4.336) (2.169,4.372) (2.170,4.408) (2.170,4.444) (2.170,4.450)}
\LC{(1.362,1.715) (1.363,1.751) (1.363,1.787) (1.363,1.823) (1.363,1.859) (1.363,1.895) (1.363,1.931) (1.363,1.967) (1.363,2.003) (1.362,2.039) (1.362,2.075) (1.362,2.111) (1.362,2.147) (1.361,2.183) (1.361,2.219) (1.360,2.255) (1.360,2.291) (1.359,2.327) (1.358,2.363) (1.358,2.399) (1.357,2.435) (1.356,2.471) (1.355,2.507) (1.354,2.543) (1.354,2.579) (1.353,2.615) (1.352,2.651) (1.351,2.687) (1.350,2.723) (1.349,2.759) (1.349,2.795) (1.348,2.831) (1.347,2.867) (1.346,2.903) (1.346,2.939) (1.345,2.975) (1.344,3.011) (1.344,3.047) (1.343,3.083) (1.342,3.119) (1.342,3.155) (1.341,3.191) (1.341,3.227) (1.340,3.263) (1.340,3.299) (1.340,3.335) (1.339,3.371) (1.339,3.407) (1.339,3.443) (1.338,3.479) (1.338,3.515) (1.338,3.551) (1.338,3.587) (1.338,3.623) (1.338,3.659) (1.337,3.695) (1.337,3.731) (1.337,3.767) (1.337,3.803) (1.337,3.839) (1.337,3.875) (1.337,3.911) (1.337,3.947) (1.337,3.983) (1.337,4.019) (1.337,4.055) (1.337,4.091) (1.337,4.127) (1.337,4.163) (1.337,4.199) (1.337,4.235) (1.337,4.271) (1.337,4.307) (1.337,4.343) (1.337,4.379) (1.337,4.415) (1.337,4.450)}

    % Disco cargado, resto neutros - flechas
% Generado por scripts/generar_datos.py (no editar a mano).
\FL{-4.553}{-3.230}{-165.4}
\FL{-4.383}{-3.753}{-140.8}
\FL{-3.830}{-3.980}{-115.9}
\FL{-3.424}{-4.120}{-105.9}
\FL{-3.106}{-4.217}{-91.9}
\FL{-2.703}{-4.114}{-73.5}
\FL{-2.302}{-3.932}{-68.0}
\FL{-1.713}{-3.653}{-68.7}
\FL{-0.860}{-3.200}{-77.3}
\FL{-3.857}{-1.417}{156.6}
\FL{-4.287}{-2.170}{156.6}
\FL{-4.533}{-2.677}{168.3}
\FL{-0.860}{3.677}{91.2}
\FL{-1.746}{3.676}{93.1}
\FL{-2.673}{3.665}{105.1}
\FL{-3.862}{3.071}{149.4}
\FL{-3.876}{1.729}{175.8}
\FL{3.149}{0.750}{-24.5}
\FL{1.791}{0.834}{45.2}
\FL{1.562}{0.945}{58.1}
\FL{1.392}{1.001}{65.0}
\FL{1.264}{1.050}{70.2}
\FL{1.136}{1.074}{72.6}
\FL{1.017}{1.098}{74.1}
\FL{0.896}{1.150}{76.3}
\FL{0.923}{2.635}{96.1}
\FL{0.027}{2.594}{114.5}
\FL{0.097}{1.561}{129.1}
\FL{-0.022}{1.224}{140.2}
\FL{-0.109}{0.963}{150.9}
\FL{-0.375}{0.761}{162.7}
\FL{-2.260}{0.941}{-176.9}
\FL{-0.366}{0.069}{-134.0}
\FL{-0.070}{0.054}{-127.6}
\FL{0.093}{0.022}{-126.0}
\FL{0.206}{-0.010}{-129.4}
\FL{0.387}{-0.209}{-150.3}
\FL{0.394}{-0.321}{-156.3}
\FL{0.410}{-0.430}{-155.5}
\FL{0.417}{-0.649}{-152.0}
\FL{0.323}{-1.091}{-146.4}
\FL{0.554}{-2.271}{-99.5}
\FL{1.531}{-2.275}{-80.0}
\FL{2.470}{-2.191}{-59.5}
\FL{2.976}{-1.384}{-39.3}
\FL{3.107}{-0.649}{-27.2}
\FL{3.173}{0.059}{-18.6}
\FL{4.076}{0.818}{-11.9}
\FL{4.267}{1.362}{0.5}
\FL{4.267}{1.804}{12.6}
\FL{4.041}{2.433}{32.8}
\FL{3.757}{3.180}{52.2}
\FL{2.828}{3.093}{74.4}
\FL{2.081}{3.092}{84.5}
\FL{1.343}{3.083}{90.9}

  \end{scope}
  \foreach \f in {\formaA,\formaB,\formaD}{\fill[top color=azulvelo!80!white, bottom color=metal, draw=azulprof, line width=0.55pt] \f;}
  \fill[top color=azulcielo, bottom color=azulmedio, draw=azulprof, line width=0.55pt] \formaC;
  % Disco cargado, resto neutros - cargas
% Generado por scripts/generar_datos.py (no editar a mano).
\CP{-3.977}{-3.386}
\CP{-3.165}{-4.063}
\CP{-2.068}{-3.062}
\CP{-3.817}{-2.494}
\CN{0.084}{-0.581}
\CN{0.336}{-0.291}
\CN{-0.150}{-0.221}
\CP{-1.508}{3.010}
\CP{-2.910}{2.402}
\CN{-0.810}{1.090}
\CN{-0.490}{1.714}
\CP{1.447}{0.581}
\CP{1.197}{0.888}
\CP{0.865}{0.990}
\CP{0.514}{0.895}
\CP{0.258}{0.593}
\CP{0.216}{0.264}
\CP{0.314}{-0.000}
\CP{0.476}{-0.169}
\CP{0.752}{-0.282}
\CP{1.142}{-0.219}
\CP{1.448}{0.123}
\CP{3.710}{1.369}
\CP{3.270}{1.840}
\CP{1.891}{1.840}
\CN{1.170}{1.160}
\CN{1.531}{1.160}
\CN{2.452}{1.160}

  \draw[azulprof, line width=0.7pt] (-5,-4.5) rectangle (5,4.5);
  % rótulos
  \node[font=\footnotesize, fill=white, inner sep=1.5pt, rounded corners=1pt, align=center] at (-3.55,-1.55) {A\\[-1pt]\scriptsize$\phi=\csname philibreA\endcsname$};
  \node[font=\footnotesize, inner sep=1pt, align=center] at (-1.7,2.05) {B\\[-1pt]\scriptsize$\phi=\csname philibreB\endcsname$};
  \node[font=\footnotesize, fill=white, inner sep=1.5pt, rounded corners=1pt, align=center] at (2.05,-0.6) {C\ ($Q>0$)\\[-1pt]\scriptsize$\phi=\csname philibreC\endcsname$};
  \node[font=\footnotesize, fill=white, inner sep=1.5pt, rounded corners=1pt, align=center] at (2.3,2.45) {D\\[-1pt]\scriptsize$\phi=\csname philibreD\endcsname$};
  \node[font=\scriptsize, text=grispizarra, anchor=north east] at (4.9,-3.8) {caja a tierra, $\phi=0$};
\end{tikzpicture}
\caption{Estado «un conductor cargado, los demás neutros», que define la columna $\mathrm{C}$ de los coeficientes de potencial (sección de conductores
cilíndricos, calculado numéricamente; potenciales normalizados a $\phi_{\mathrm{C}}=1$). Las líneas salen del disco cargado; los conductores neutros
reciben líneas por la cara próxima, donde se induce carga negativa, y las reemiten por la opuesta. Todos los potenciales son positivos y el del
conductor cargado es el mayor: $\phi_{\mathrm{C}}>\phi_{i}>0$.}
\label{fig:sistema-libre}
\end{figure}


% -------------------------------------------------------------------------
\subsection{Ejemplo: dos esferas concéntricas}\label{sec:ej-concentricas}
% -------------------------------------------------------------------------
\begin{ejemplo}[Esfera y cáscara esférica concéntricas]\label{ej:concentricas}
Una esfera conductora de radio $R_1$ (conductor~1) está rodeada por una cáscara esférica conductora delgada y concéntrica de radio $R_2>R_1$
(conductor~2). Calculemos sus coeficientes de potencial.

\emph{Estado $Q_1\neq0$, $Q_2=0$.} La cáscara neutra adquiere $-Q_1$ en su cara interior y $+Q_1$ en la exterior, ambas en el mismo radio, de modo que
no modifica el campo. Por Gauss, el campo es el de una carga $Q_1$ en el centro para $r>R_1$, y el potencial (\figpanel{fig:concentricas}{b})
\[
  \phi(r)=\begin{cases}\dfrac{Q_1}{4\pi\eo r}, & r\ge R_1,\\[8pt] \dfrac{Q_1}{4\pi\eo R_1}, & r\le R_1,\end{cases}
  \qquad\Longrightarrow\qquad p_{11}=\left.\frac{\phi_1}{Q_1}\right|_{Q_2=0}=\frac{1}{4\pi\eo R_1},\quad
  p_{21}=\left.\frac{\phi_2}{Q_1}\right|_{Q_2=0}=\frac{1}{4\pi\eo R_2}.
\]
\emph{Estado $Q_2\neq0$, $Q_1=0$.} La carga de la cáscara crea campo solo fuera de ella, y dentro el potencial es constante (\figpanel{fig:concentricas}{c}):
\[
  \phi(r)=\begin{cases}\dfrac{Q_2}{4\pi\eo r}, & r\ge R_2,\\[8pt] \dfrac{Q_2}{4\pi\eo R_2}, & r\le R_2,\end{cases}
  \qquad\Longrightarrow\qquad p_{12}=\left.\frac{\phi_1}{Q_2}\right|_{Q_1=0}=\frac{1}{4\pi\eo R_2}=p_{22}.
\]
En resumen,
\begin{equation}\label{eq:p-concentricas}
  \Pm=\frac{1}{4\pi\eo}\begin{pmatrix}1/R_1 & 1/R_2\\ 1/R_2 & 1/R_2\end{pmatrix},
\end{equation}
que es simétrica ($p_{12}=p_{21}$), de elementos positivos, y con $p_{11}>p_{21}$ y $p_{22}=p_{12}$: la igualdad del \cref{teo:pij-prop}\,(c) se alcanza
porque la esfera interior está encerrada por la cáscara y comparte su potencial cuando solo esta está cargada.
\end{ejemplo}

% =========================================================================
%  Figura: esfera (R1) y cáscara (R2 = 2 R1) concéntricas. Potencial de los
%  dos estados elementales, normalizado a kQ/R1:
%  (a) Q1 = Q, Q2 = 0:  phi = 1 (r <= 1),  1/r (r >= 1)   (r en unidades de R1)
%  (b) Q2 = Q, Q1 = 0:  phi = 1/2 (r <= 2), 1/r (r >= 2)
% =========================================================================
\begin{figure}[t]
\centering
\begin{tikzpicture}[font=\small]
  % esquema
  \begin{scope}[x=0.75cm, y=0.75cm]
    \draw[azulprof, line width=0.9pt, fill=azulfondo] (0,0) circle (2);
    \shade[ball color=azulvelo] (0,0) circle (1);
    \draw[azulprof, line width=0.6pt] (0,0) circle (1);
    \draw[-{Stealth[length=4pt]}, azulprof, line width=0.5pt] (0,0) -- (-0.71,0.71) node[pos=0.55, below left=-2pt, font=\footnotesize] {$R_1$};
    \draw[-{Stealth[length=4pt]}, azulprof, line width=0.5pt] (0,0) -- (1.73,-1.0) node[pos=0.7, above right=-2pt, font=\footnotesize] {$R_2$};
    \node[font=\footnotesize] at (0.25,0.3) {$1$};
    \node[font=\footnotesize] at (1.55,1.55) {$2$};
    \node at (0,-4.5) {(a)};
  \end{scope}
  % estado Q1
  \begin{scope}[xshift=2.3cm, yshift=-2.05cm]
    \begin{axis}[tema, width=6.0cm, height=4.6cm, xmin=0, xmax=4.2, ymin=0, ymax=1.15,
        xlabel={$r/R_1$}, ylabel={$\phi\,/\,(\kc Q/R_1)$}, xtick={0,1,2,3,4}, ytick={0,0.5,1},
        title={\footnotesize $Q_1=Q,\ Q_2=0$}, title style={yshift=-4pt}]
      \addplot[azulprof, line width=1.1pt, domain=0:1] {1};
      \addplot[azulprof, line width=1.1pt, domain=1:4.2, samples=60] {1/x};
      \draw[punteado] (axis cs:1,0) -- (axis cs:1,1.15) (axis cs:2,0) -- (axis cs:2,1.15);
      \fill[cargapos] (axis cs:1,1) circle (1.6pt) node[above right=-1pt, font=\scriptsize] {$p_{11}$};
      \fill[cargapos] (axis cs:2,0.5) circle (1.6pt) node[above right=-1pt, font=\scriptsize] {$p_{21}$};
    \end{axis}
    \node at (2.3,-1.35) {(b)};
  \end{scope}
  % estado Q2
  \begin{scope}[xshift=8.7cm, yshift=-2.05cm]
    \begin{axis}[tema, width=6.0cm, height=4.6cm, xmin=0, xmax=4.2, ymin=0, ymax=1.15,
        xlabel={$r/R_1$}, ylabel={$\phi\,/\,(\kc Q/R_1)$}, xtick={0,1,2,3,4}, ytick={0,0.5,1},
        title={\footnotesize $Q_1=0,\ Q_2=Q$}, title style={yshift=-4pt}]
      \addplot[azulacero, line width=1.1pt, domain=0:2] {0.5};
      \addplot[azulacero, line width=1.1pt, domain=2:4.2, samples=60] {1/x};
      \draw[punteado] (axis cs:1,0) -- (axis cs:1,1.15) (axis cs:2,0) -- (axis cs:2,1.15);
      \fill[cargapos] (axis cs:1,0.5) circle (1.6pt) node[above=1pt, font=\scriptsize] {$p_{12}$};
      \fill[cargapos] (axis cs:2,0.5) circle (1.6pt) node[above right=-1pt, font=\scriptsize] {$p_{22}$};
    \end{axis}
    \node at (2.3,-1.35) {(c)};
  \end{scope}
\end{tikzpicture}
\caption{Coeficientes de potencial de una esfera y una cáscara concéntricas ($R_2=2R_1$). (a) Geometría. (b) Con la esfera cargada y la cáscara
neutra, el potencial decae como $1/r$ desde $R_1$. (c) Con la cáscara cargada, el interior queda a potencial constante, $\kc Q/R_2$, de modo que
$p_{12}=p_{22}$.}
\label{fig:concentricas}
\end{figure}


El ejemplo se generaliza de inmediato, y da dos reglas muy útiles para los seminarios.

\begin{proposicion}[Reglas prácticas para esferas]\label{prop:reglas}
\begin{enumerate}[label=\textup{(\alph*)}]
  \item \emph{Esferas concéntricas delgadas} de radios $R_1<R_2<\dots<R_N$:
        \begin{equation}\label{eq:regla-concentricas}
          \resalta{p_{ij}=\frac{1}{4\pi\eo}\,\frac{1}{\max(R_i,R_j)} .}
        \end{equation}
  \item \emph{Esferas lejanas} de radios $R_i$ con distancias entre centros $d_{ij}\gg R_i,R_j$: al orden más bajo en $R/d$,
        \begin{equation}\label{eq:regla-lejanas}
          \resalta{p_{ii}\simeq\frac{1}{4\pi\eo R_i},\qquad p_{ij}\simeq\frac{1}{4\pi\eo\,d_{ij}}\quad(i\neq j).}
        \end{equation}
\end{enumerate}
\end{proposicion}
\begin{proof}
(a) En el estado elemental $j$, las cáscaras neutras llevan cargas opuestas en sus dos caras (en el mismo radio) y no modifican el campo; por Gauss,
este vale $\kc/r^{2}$ fuera de $R_j$ y es nulo dentro. El potencial es $\kc/r$ para $r\ge R_j$ y constante, $\kc/R_j$, para $r\le R_j$. El conductor $i$,
en el radio $R_i$, tiene potencial $\kc/R_i$ si $R_i\ge R_j$ y $\kc/R_j$ si $R_i\le R_j$, es decir, $\kc/\max(R_i,R_j)$.

(b) En el estado elemental $j$, la carga de la esfera $j$ apenas se ve perturbada por las demás, que están muy lejos: se reparte casi uniformemente y
su potencial es el de una esfera aislada, $\kc/R_j$. Vista desde la esfera $i$, a distancia $d_{ij}$, es una carga puntual que crea un potencial
$\kc/d_{ij}$; la esfera $i$, neutra, solo adquiere una distribución inducida de carga neta nula (dipolar), cuya contribución a los potenciales es de
orden superior en $R/d$.
\end{proof}

La regla (a) se recuerda así: cada carga «se ve» desde su propio radio hacia dentro con el potencial de su superficie, y desde fuera como una
carga puntual. Con tres esferas,
\[
  \Pm=\kc\begin{pmatrix}1/R_1 & 1/R_2 & 1/R_3\\ 1/R_2 & 1/R_2 & 1/R_3\\ 1/R_3 & 1/R_3 & 1/R_3\end{pmatrix}:
\]
cada fila contiene el inverso del radio propio y los inversos de los radios mayores.

% -------------------------------------------------------------------------
\subsection{Coeficientes de capacidad e influencia}\label{sec:cij}
% -------------------------------------------------------------------------
Normalmente lo que se puede fijar es el potencial de un conductor (conectándolo a tierra o a una batería), y lo que se pregunta es qué carga adquiere.
Este es el fundamento de los coeficientes de capacidad e influencia: hay que invertir \cref{eq:pij}.

\begin{lema}[La matriz de coeficientes de potencial es definida positiva]\label{lem:definida}
Para todo vector de cargas $\vQ\neq\vc{0}$,
\begin{equation}\label{eq:definida}
  \vQ\Tt\Pm\,\vQ=\sum_{i,j}Q_i\,p_{ij}\,Q_j=\sum_iQ_i\,\phi_i=\eo\iiint_{\Omega}\norm{\vE}^{2}\dV>0 .
\end{equation}
En consecuencia, $\Pm$ es invertible, $\det\Pm>0$, y su inversa es también simétrica y definida positiva.
\end{lema}
\begin{proof}
Con las notaciones de la demostración del \cref{teo:unicidad}, aplicando el teorema de la divergencia a $\phi\nab\phi$ en la región exterior $\Omega$,
en la que $\Lap\phi=0$,
\[
  \eo\iiint_{\Omega}\norm{\nab\phi}^{2}\dV=-\eo\sum_i\phi_i\oiint_{\Sc_i}\nab\phi\cdot\un\,\dsc=\sum_i\phi_i\,Q_i .
\]
El primer miembro es positivo salvo que $\vE=-\nab\phi$ sea idénticamente nulo, en cuyo caso todas las cargas (los flujos) son nulas. Una matriz simétrica
y definida positiva tiene todos sus autovalores positivos: su determinante es positivo y su inversa es simétrica y definida positiva. (La cantidad
$\tfrac12\sum_iQ_i\phi_i$ resultará ser la energía electrostática del sistema.)
\end{proof}

\begin{definicion}[Coeficientes de capacidad e influencia]\label{def:cij}
Se llama matriz de \emph{coeficientes de capacidad e influencia} a $\Cm=(c_{ij})=\Pm^{-1}$, de modo que
\begin{equation}\label{eq:cij}
  \resalta{Q_i=\sum_{j=1}^{N}c_{ij}\,\phi_j,\qquad\vQ=\Cm\,\vphi,}\qquad
  c_{ij}=\left.\frac{Q_i}{\phi_j}\right|_{\phi_k=0,\ k\neq j},\qquad[c_{ij}]=\mathrm{C/V}=\mathrm{F}.
\end{equation}
Los elementos diagonales $c_{ii}$ se llaman \emph{coeficientes de capacidad}, y los no diagonales $c_{ij}$ ($i\neq j$), \emph{coeficientes de influencia}
(o de inducción).
\end{definicion}

La columna $j$ de $\Cm$ tiene una lectura inmediata: son las cargas de todos los conductores cuando el $j$ está a potencial unidad y todos los demás están
a tierra. Así, $c_{jj}$ es la carga del propio conductor $j$ en ese estado, y $c_{ij}$ la carga que aparece en el conductor $i$ (a tierra) por influencia
del $j$.

\begin{teorema}[Propiedades de los coeficientes de capacidad e influencia]\label{teo:cij-prop}
\begin{enumerate}[label=\textup{(\alph*)}]
  \item Solo dependen de la geometría y del medio, y son simétricos: $c_{ij}=c_{ji}$.
  \item Los coeficientes de capacidad son positivos, $c_{jj}>0$, y los de influencia negativos o nulos, $c_{ij}\le0$ para $i\neq j$.
  \item $c_{jj}$ \emph{no} es la capacidad del conductor $j$ con los demás descargados, $1/p_{jj}$: es la carga de $j$ por unidad de potencial cuando
        todos los demás están a tierra, y $c_{jj}\ge1/p_{jj}$.
\end{enumerate}
\end{teorema}
\begin{proof}
(a) $\Cm=\Pm^{-1}$, y la inversa de una matriz simétrica es simétrica. Que $\Cm$ solo depende de la geometría se hereda de $\Pm$.

(b) Consideremos el estado $\phi_j=1$, $\phi_k=0$ ($k\neq j$), con $\phi=0$ también en el infinito. Por el principio del máximo, $0\le\phi\le1$ en todo
el espacio. En la superficie de $j$, el potencial máximo, ninguna línea de campo puede llegar (vendría de un potencial mayor): $\sigma\ge0$ en todo $\Sc_j$,
y no es idénticamente nula (si lo fuese, $\phi$ sería constante e igual a $1$ en todo el espacio, en contra de $\phi(\infty)=0$). Luego $c_{jj}=Q_j>0$.
En la superficie de un conductor $i\neq j$, al potencial mínimo $0$, ninguna línea puede nacer (iría a un potencial menor que el mínimo): $\sigma\le0$,
y $c_{ij}=Q_i\le0$.

\emph{Argumento con los coeficientes de potencial.} Con solo dos conductores, $\phi_i=p_{ii}Q_i+p_{ij}Q_j=0$ da $Q_i=-(p_{ij}/p_{ii})\,Q_j$; como
$p_{ij},p_{ii}>0$, si $\phi_j>0$ (y por tanto $Q_j>0$), $Q_i<0$: la carga inducida tiene signo contrario.

(c) Con los demás descargados, $\phi_j=p_{jj}Q_j$, y el cociente $Q_j/\phi_j$ vale $1/p_{jj}$. Para la desigualdad, sea $\bm{e}_j$ el
$j$-ésimo vector de la base canónica. Como $\Pm$ es simétrica y definida positiva (\cref{lem:definida}), tiene raíz cuadrada $\Pm^{1/2}$, y por la
desigualdad de Cauchy--Schwarz
\[
  1=\bigl(\bm{e}_j\Tt\bm{e}_j\bigr)^{2}=\bigl[(\Pm^{1/2}\bm{e}_j)\Tt(\Pm^{-1/2}\bm{e}_j)\bigr]^{2}
  \le\bigl(\bm{e}_j\Tt\Pm\,\bm{e}_j\bigr)\bigl(\bm{e}_j\Tt\Pm^{-1}\bm{e}_j\bigr)=p_{jj}\,c_{jj} .
\]
Físicamente: los conductores a tierra «tiran» de carga hacia el conductor $j$ más que los conductores aislados.
\end{proof}

% =========================================================================
%  Figura: influencia total. Un tubo delgado de líneas de campo une el
%  conductor j (phi_j = 1) con el conductor i (a tierra). La superficie
%  gaussiana cerrada S = S_lat ∪ B_1 ∪ B_2 (a trazos) está formada por la
%  pared del tubo (E tangente, E ⊥ dS) y por dos tapas dentro de los
%  conductores (E = 0). Estilo minimalista: bordes marcados e interiores con
%  un velo radial sutil (estilo «conductorfino»).
% =========================================================================
\begin{figure}[t]
\centering
\begin{tikzpicture}[font=\small, x=1cm, y=1cm]
  % --- conductores --------------------------------------------------------------
  \path[conductorfino]
    (-3.7,1.35) .. controls (-2.6,1.65) and (-1.6,1.2) .. (-1.55,0.3)
    .. controls (-1.5,-0.6) and (-2.4,-1.45) .. (-3.5,-1.25) .. controls (-4.35,-1.05) and (-4.45,1.05) .. cycle;
  \path[conductorfino]
    (2.0,0.95) .. controls (2.6,1.55) and (3.85,1.25) .. (3.95,0.2)
    .. controls (4.05,-0.85) and (3.0,-1.35) .. (2.2,-0.95) .. controls (1.6,-0.65) and (1.6,0.5) .. cycle;
  \node[font=\footnotesize] at (-3.2,0.75) {$i$};
  \node[font=\footnotesize] at (3.3,0.75) {$j$};
  \node[font=\footnotesize, anchor=west] at (3.35,1.45) {$\phi_j=1$};
  \draw[azulacero, line width=0.6pt] (-3.3,-1.27) -- (-3.3,-1.6);
  \pic at (-3.3,-1.6) {tierra};
  \node[font=\footnotesize, anchor=west] at (-3.0,-1.82) {$\phi_i=0$};
  % --- tubo: interior con degradado tenue ---------------------------------------
  \fill[top color=azulvelo!55, bottom color=azulvelo!55, middle color=white]
    (-1.86,0.28) .. controls (-0.6,0.36) and (0.7,0.36) .. (2.0,0.28)
    -- (2.0,-0.28) .. controls (0.7,-0.36) and (-0.6,-0.36) .. (-1.86,-0.28) -- cycle;
  % líneas de campo (de j a i)
  \draw[lcampo] (1.64,0) -- (-1.53,0);
  \draw[azulacero!70, line width=0.4pt] (1.64,0.16) .. controls (0.6,0.2) and (-0.5,0.2) .. (-1.53,0.16);
  \draw[azulacero!70, line width=0.4pt] (1.64,-0.16) .. controls (0.6,-0.2) and (-0.5,-0.2) .. (-1.53,-0.16);
  % --- superficie gaussiana: pared lateral y tapas --------------------------------
  \draw[gauss, line width=0.6pt] (-1.86,0.28) .. controls (-0.6,0.36) and (0.7,0.36) .. (2.0,0.28);
  \draw[gauss, line width=0.6pt] (-1.86,-0.28) .. controls (-0.6,-0.36) and (0.7,-0.36) .. (2.0,-0.28);
  \draw[gauss, line width=0.6pt] (-1.86,0) ellipse (0.07 and 0.28);
  \draw[gauss, line width=0.6pt] (2.0,0) ellipse (0.07 and 0.28);
  % sección transversal del tubo (da la sensación de cilindro)
  \draw[cargapos!60, line width=0.4pt] (0.05,0.34) arc (90:270:0.07 and 0.34);
  \draw[cargapos!60, line width=0.4pt, densely dotted] (0.05,-0.34) arc (-90:90:0.07 and 0.34);
  % --- cargas en los extremos -------------------------------------------------
  \node[signoneg, font=\tiny\bfseries] at (-1.68,0.12) {$-$};
  \node[signoneg, font=\tiny\bfseries] at (-1.68,-0.12) {$-$};
  \node[signopos, font=\tiny\bfseries] at (1.8,0.12) {$+$};
  \node[signopos, font=\tiny\bfseries] at (1.8,-0.12) {$+$};
  % --- rótulos de la superficie --------------------------------------------------
  \draw[grismedio, line width=0.35pt] (-1.86,-0.3) -- (-2.25,-0.75);
  \node[fondo, font=\footnotesize, text=cargapos!85!black] at (-2.4,-0.88) {$B_1$};
  \draw[grismedio, line width=0.35pt] (2.0,-0.3) -- (2.4,-0.75);
  \node[fondo, font=\footnotesize, text=cargapos!85!black] at (2.55,-0.88) {$B_2$};
  \draw[grismedio, line width=0.35pt] (0.9,-0.33) -- (1.15,-0.75);
  \node[fondo, font=\footnotesize, text=cargapos!85!black] at (1.25,-0.9) {$S_{\mathrm{lat}}$};
  \node[font=\scriptsize, text=azulprof] at (-2.75,0.15) {$\vE=\vc{0}$};
  \node[font=\scriptsize, text=azulprof] at (3.05,0.15) {$\vE=\vc{0}$};
  % --- E y dS en un punto de la pared --------------------------------------------
  \coordinate (P) at (-0.5,0.34);
  \draw[-{Stealth[length=3.5pt]}, verdeazul, line width=0.7pt] (P) -- ++(0,0.55);
  \node[fondo, font=\footnotesize, text=verdeazul, anchor=west] at ($(P)+(0.04,0.45)$) {$\dS$};
  \draw[-{Stealth[length=3.5pt]}, azulprof, line width=0.7pt] (P) -- ++(-0.55,0);
  \node[fondo, font=\footnotesize, text=azulprof, anchor=south] at ($(P)+(-0.42,0.04)$) {$\vE$};
  \fill[azulprof] (P) circle (0.9pt);
  % --- leyenda --------------------------------------------------------------------
  \draw[gauss, line width=0.6pt] (-1.3,-1.95) -- (-0.7,-1.95);
  \node[font=\scriptsize, anchor=west] at (-0.6,-1.95) {superficie gaussiana $S=S_{\mathrm{lat}}\cup B_1\cup B_2$};
\end{tikzpicture}
\caption{Influencia total. La superficie cerrada $S$ formada por un tubo de líneas de campo (pared lateral $S_{\mathrm{lat}}$, donde $\vE\perp\dS$) y por
dos tapas $B_1$ y $B_2$ dentro de los conductores (donde $\vE=\vc{0}$) tiene flujo nulo: la carga positiva de la que nace el tubo en $j$ es igual y
opuesta a la negativa en la que muere en $i$.}
\label{fig:tubo}
\end{figure}


% -------------------------------------------------------------------------
\subsection{Influencia total}\label{sec:influencia}
% -------------------------------------------------------------------------
Pongamos el conductor $j$ a potencial $\phi_j=1$ y todos los demás a tierra. Las líneas de campo, que van siempre hacia potenciales menores, salen de
$j$ ($Q_j>0$) y pueden acabar en los otros conductores o irse al infinito.

\begin{teorema}[Desigualdad de influencia]\label{teo:influencia}
En el estado $\phi_j>0$, $\phi_k=0$ ($k\neq j$), la carga inducida en los demás conductores no puede superar, en valor absoluto, a la del conductor $j$:
\begin{equation}\label{eq:influencia}
  Q_j\ge-\sum_{i\neq j}Q_i ,\qquad\text{equivalentemente}\qquad\resalta{c_{jj}\ge-\sum_{i\neq j}c_{ij}\quad\Longleftrightarrow\quad\sum_{i=1}^{N}c_{ij}\ge0 .}
\end{equation}
La igualdad se da si, y solo si, todas las líneas de campo que salen de $j$ terminan en los otros conductores y ninguna llega al infinito; se dice entonces
que el sistema está en \emph{influencia total}. El caso típico es el de un conductor encerrado dentro de otro.
\end{teorema}
\begin{proof}
Consideremos primero dos conductores con todas las líneas de $j$ terminando en $i$. Elegimos una superficie cerrada $S$ formada por el tubo de líneas de
campo que sale de $j$ (su pared lateral, $S_{\mathrm{lat}}$) y cerrada por dos tapas $B_1$ y $B_2$ dentro de los conductores $i$ y $j$
(\cref{fig:tubo}). Por la ley de Gauss,
\[
  \PhiE=\oiint_{S}\vE\cdot\dS=\underbrace{\iint_{S_{\mathrm{lat}}}\vE\cdot\dS}_{0\ (\vE\perp\dS)}+\underbrace{\iint_{B_1}\vE\cdot\dS}_{0\ (\vE=\vc{0})}
  +\underbrace{\iint_{B_2}\vE\cdot\dS}_{0\ (\vE=\vc{0})}=0=\frac{\Qenc}{\eo},
\]
porque en la pared lateral el campo es tangente a la superficie y en las tapas, interiores a los conductores, es nulo. La carga encerrada es la que hay en
los extremos del tubo, la de salida en $j$ y la de llegada en $i$, luego $Q_j=-Q_i$. Con $N$ conductores el argumento se aplica a cada tubo elemental: la
carga positiva en el origen de cada tubo es igual y opuesta a la negativa en su extremo. Si todas las líneas que salen de $j$ terminan en los demás
conductores, $Q_j=-\sum_{i\neq j}Q_i$; si alguna se va al infinito, la carga que llega a los demás es menor y $Q_j>-\sum_{i\neq j}Q_i$. Además, como los
demás están al potencial mínimo, no emiten líneas (\cref{teo:cij-prop}), así que no hay otras contribuciones. Dividiendo por $\phi_j$ y usando
$Q_i=c_{ij}\phi_j$ se obtiene la forma con los coeficientes.
\end{proof}

\begin{observacion}[Una comprobación numérica]\label{obs:numerico}
La \cref{fig:sistema-cij} muestra los cuatro estados elementales de un sistema de cuatro conductores (una cuña, un rectángulo, un disco y una barra),
calculados numéricamente en una caja a tierra, que hace el papel del infinito. Integrando el flujo alrededor de cada conductor se obtiene la matriz de
la \cref{tab:sistema-coef}. Se comprueban, dentro de la precisión numérica, todas las propiedades demostradas: simetría, $c_{jj}>0$, $c_{ij}<0$ y
$\sum_ic_{ij}>0$ (parte de las líneas termina en la caja), así como $p_{ij}>0$ y $p_{jj}\ge p_{ij}$ para la matriz inversa. Los coeficientes de influencia
más intensos son los de los pares más próximos (cuña--disco y disco--barra), y los más débiles, los de los pares lejanos y apantallados (cuña--barra).
\end{observacion}

% =========================================================================
%  Figura «tipo Mathematica»: los cuatro estados elementales phi_j = 1,
%  phi_k = 0 (k != j) de un sistema de cuatro conductores en una caja a tierra.
%  Fila superior: potencial (escala azulpot de 0 a 1) con equipotenciales.
%  Fila inferior: módulo del campo (escala azulcampo, logarítmica) con las
%  mismas equipotenciales. Mapas: figuras/datos/sistema_{pot,campo}_j.png;
%  contornos: figuras/datos/sistema_eq_j.tex y sistema_formas.tex.
% =========================================================================
\pgfdeclarehorizontalshading{azulpotsh}{100bp}{rgb(0bp)=(0.961,0.973,0.988);
  rgb(20bp)=(0.863,0.906,0.961); rgb(40bp)=(0.616,0.725,0.867); rgb(60bp)=(0.227,0.431,0.647);
  rgb(80bp)=(0.122,0.306,0.537); rgb(100bp)=(0.043,0.145,0.271)}
\pgfdeclarehorizontalshading{azulcamposh}{100bp}{rgb(0bp)=(0.043,0.145,0.271);
  rgb(20bp)=(0.122,0.306,0.537); rgb(40bp)=(0.227,0.525,0.710); rgb(60bp)=(0.498,0.765,0.851);
  rgb(80bp)=(0.812,0.918,0.945); rgb(100bp)=(1,1,1)}
\newcommand{\panelSistema}[3]{% #1 = j (1..4), #2 = pot / campo, #3 = conductor activo (A..D)
  \node[anchor=south west, inner sep=0] at (-5,-4.5)
    {\includegraphics[width=3.7cm, height=3.33cm]{figuras/datos/sistema_#2_#1.png}};
  \begin{scope}
    \clip (-5,-4.5) rectangle (5,4.5);
    \input{figuras/datos/sistema_eq_#1}
  \end{scope}
  \foreach \f in {\formaA,\formaB,\formaC,\formaD}{\fill[white, draw=azulprof, line width=0.4pt] \f;}
  \fill[azulprof, draw=azulprof, line width=0.4pt] \csname forma#3\endcsname;
  \draw[azulprof, line width=0.5pt] (-5,-4.5) rectangle (5,4.5);}

\begin{figure}[t]
\centering
\begin{tikzpicture}[x=0.37cm, y=0.37cm, font=\footnotesize]
  % Contornos de los cuatro conductores
% Generado por scripts/generar_datos.py (no editar a mano).
\def\formaA{(-4.250,-2.950) -- (-3.050,-3.950) -- (0.220,-0.180) -- cycle}
\def\formaB{(-2.750,1.250) -- (-0.650,1.250) -- (-0.650,2.850) -- (-2.750,2.850) -- cycle}
\def\formaC{(1.330,0.350) -- (1.330,0.363) -- (1.329,0.375) -- (1.329,0.388) -- (1.327,0.400) -- (1.326,0.413) -- (1.324,0.425) -- (1.322,0.437) -- (1.320,0.450) -- (1.317,0.462) -- (1.314,0.474) -- (1.310,0.486) -- (1.307,0.498) -- (1.302,0.510) -- (1.298,0.522) -- (1.293,0.534) -- (1.289,0.545) -- (1.283,0.557) -- (1.278,0.568) -- (1.272,0.579) -- (1.266,0.590) -- (1.259,0.601) -- (1.253,0.611) -- (1.246,0.622) -- (1.238,0.632) -- (1.231,0.642) -- (1.223,0.652) -- (1.215,0.662) -- (1.207,0.671) -- (1.198,0.680) -- (1.189,0.689) -- (1.180,0.698) -- (1.171,0.707) -- (1.162,0.715) -- (1.152,0.723) -- (1.142,0.731) -- (1.132,0.738) -- (1.122,0.746) -- (1.111,0.753) -- (1.101,0.759) -- (1.090,0.766) -- (1.079,0.772) -- (1.068,0.778) -- (1.057,0.783) -- (1.045,0.789) -- (1.034,0.793) -- (1.022,0.798) -- (1.010,0.802) -- (0.998,0.807) -- (0.986,0.810) -- (0.974,0.814) -- (0.962,0.817) -- (0.950,0.820) -- (0.937,0.822) -- (0.925,0.824) -- (0.913,0.826) -- (0.900,0.827) -- (0.888,0.829) -- (0.875,0.829) -- (0.863,0.830) -- (0.850,0.830) -- (0.837,0.830) -- (0.825,0.829) -- (0.812,0.829) -- (0.800,0.827) -- (0.787,0.826) -- (0.775,0.824) -- (0.763,0.822) -- (0.750,0.820) -- (0.738,0.817) -- (0.726,0.814) -- (0.714,0.810) -- (0.702,0.807) -- (0.690,0.802) -- (0.678,0.798) -- (0.666,0.793) -- (0.655,0.789) -- (0.643,0.783) -- (0.632,0.778) -- (0.621,0.772) -- (0.610,0.766) -- (0.599,0.759) -- (0.589,0.753) -- (0.578,0.746) -- (0.568,0.738) -- (0.558,0.731) -- (0.548,0.723) -- (0.538,0.715) -- (0.529,0.707) -- (0.520,0.698) -- (0.511,0.689) -- (0.502,0.680) -- (0.493,0.671) -- (0.485,0.662) -- (0.477,0.652) -- (0.469,0.642) -- (0.462,0.632) -- (0.454,0.622) -- (0.447,0.611) -- (0.441,0.601) -- (0.434,0.590) -- (0.428,0.579) -- (0.422,0.568) -- (0.417,0.557) -- (0.411,0.545) -- (0.407,0.534) -- (0.402,0.522) -- (0.398,0.510) -- (0.393,0.498) -- (0.390,0.486) -- (0.386,0.474) -- (0.383,0.462) -- (0.380,0.450) -- (0.378,0.437) -- (0.376,0.425) -- (0.374,0.413) -- (0.373,0.400) -- (0.371,0.388) -- (0.371,0.375) -- (0.370,0.363) -- (0.370,0.350) -- (0.370,0.337) -- (0.371,0.325) -- (0.371,0.312) -- (0.373,0.300) -- (0.374,0.287) -- (0.376,0.275) -- (0.378,0.263) -- (0.380,0.250) -- (0.383,0.238) -- (0.386,0.226) -- (0.390,0.214) -- (0.393,0.202) -- (0.398,0.190) -- (0.402,0.178) -- (0.407,0.166) -- (0.411,0.155) -- (0.417,0.143) -- (0.422,0.132) -- (0.428,0.121) -- (0.434,0.110) -- (0.441,0.099) -- (0.447,0.089) -- (0.454,0.078) -- (0.462,0.068) -- (0.469,0.058) -- (0.477,0.048) -- (0.485,0.038) -- (0.493,0.029) -- (0.502,0.020) -- (0.511,0.011) -- (0.520,0.002) -- (0.529,-0.007) -- (0.538,-0.015) -- (0.548,-0.023) -- (0.558,-0.031) -- (0.568,-0.038) -- (0.578,-0.046) -- (0.589,-0.053) -- (0.599,-0.059) -- (0.610,-0.066) -- (0.621,-0.072) -- (0.632,-0.078) -- (0.643,-0.083) -- (0.655,-0.089) -- (0.666,-0.093) -- (0.678,-0.098) -- (0.690,-0.102) -- (0.702,-0.107) -- (0.714,-0.110) -- (0.726,-0.114) -- (0.738,-0.117) -- (0.750,-0.120) -- (0.763,-0.122) -- (0.775,-0.124) -- (0.787,-0.126) -- (0.800,-0.127) -- (0.812,-0.129) -- (0.825,-0.129) -- (0.837,-0.130) -- (0.850,-0.130) -- (0.863,-0.130) -- (0.875,-0.129) -- (0.888,-0.129) -- (0.900,-0.127) -- (0.913,-0.126) -- (0.925,-0.124) -- (0.937,-0.122) -- (0.950,-0.120) -- (0.962,-0.117) -- (0.974,-0.114) -- (0.986,-0.110) -- (0.998,-0.107) -- (1.010,-0.102) -- (1.022,-0.098) -- (1.034,-0.093) -- (1.045,-0.089) -- (1.057,-0.083) -- (1.068,-0.078) -- (1.079,-0.072) -- (1.090,-0.066) -- (1.101,-0.059) -- (1.111,-0.053) -- (1.122,-0.046) -- (1.132,-0.038) -- (1.142,-0.031) -- (1.152,-0.023) -- (1.162,-0.015) -- (1.171,-0.007) -- (1.180,0.002) -- (1.189,0.011) -- (1.198,0.020) -- (1.207,0.029) -- (1.215,0.038) -- (1.223,0.048) -- (1.231,0.058) -- (1.238,0.068) -- (1.246,0.078) -- (1.253,0.089) -- (1.259,0.099) -- (1.266,0.110) -- (1.272,0.121) -- (1.278,0.132) -- (1.283,0.143) -- (1.289,0.155) -- (1.293,0.166) -- (1.298,0.178) -- (1.302,0.190) -- (1.307,0.202) -- (1.310,0.214) -- (1.314,0.226) -- (1.317,0.238) -- (1.320,0.250) -- (1.322,0.263) -- (1.324,0.275) -- (1.326,0.287) -- (1.327,0.300) -- (1.329,0.312) -- (1.329,0.325) -- (1.330,0.337) -- cycle}
\def\formaD{(1.050,1.320) -- (3.550,1.320) -- (3.550,1.680) -- (1.050,1.680) -- cycle}

  \tikzset{numEQ/.style={azulprof!80, line width=0.25pt}}
  \foreach \j/\c/\l in {1/A/a,2/B/b,3/C/c,4/D/d}{
    \begin{scope}[xshift={(\j-1)*3.95cm}]
      \panelSistema{\j}{pot}{\c}
      \node[anchor=south] at (0,4.55) {(\l) $\phi_{\mathrm{\c}}=1$};
    \end{scope}
    \begin{scope}[xshift={(\j-1)*3.95cm}, yshift=-3.75cm]
      \tikzset{numEQ/.style={white!85!azulprof, line width=0.25pt}}
      \panelSistema{\j}{campo}{\c}
    \end{scope}}
  % barras de color
  \begin{scope}[yshift=-5.9cm]
    \shade[shading=azulpotsh] (-5,-0.2) rectangle (14.2,0.35);
    \draw[azulclaro, line width=0.3pt] (-5,-0.2) rectangle (14.2,0.35);
    \node[anchor=north] at (-5,-0.25) {$0$};
    \node[anchor=north] at (14.2,-0.25) {$1$};
    \node[anchor=north] at (4.6,-0.25) {potencial $\phi$};
    \shade[shading=azulcamposh] (17.8,-0.2) rectangle (37.03,0.35);
    \draw[azulclaro, line width=0.3pt] (17.8,-0.2) rectangle (37.03,0.35);
    \node[anchor=north] at (17.8,-0.25) {débil};
    \node[anchor=north] at (37.03,-0.25) {intenso};
    \node[anchor=north] at (27.4,-0.25) {$\norm{\vE}$ (escala logarítmica)};
  \end{scope}
\end{tikzpicture}
\caption{Estados elementales que definen los coeficientes de capacidad e influencia de un sistema de cuatro conductores (A~cuña, B~rectángulo, C~disco
y D~barra) dentro de una caja a tierra, calculados numéricamente. En cada columna, el conductor $j$ (relleno oscuro) está a potencial unidad y los demás a
tierra: la carga de cada conductor en ese estado es la columna $j$ de $\Cm$. Arriba, el potencial; abajo, el módulo del campo, que es máximo junto a las
aristas y en los huecos estrechos entre conductores próximos.}
\label{fig:sistema-cij}
\end{figure}


% =========================================================================
%  Tabla: coeficientes de capacidad e influencia y de potencial del sistema
%  de cuatro conductores (valores numéricos de figuras/datos/sistema_coef.tex,
%  por unidad de longitud, en unidades de eps0 y de 1/eps0).
% =========================================================================
% Sistema de cuatro conductores - coeficientes
% Generado por scripts/generar_datos.py (no editar a mano).
\expandafter\def\csname cijAA\endcsname{9{,}83}
\expandafter\def\csname pijAA\endcsname{0{,}1177}
\expandafter\def\csname cijAB\endcsname{-1{,}10}
\expandafter\def\csname pijAB\endcsname{0{,}0312}
\expandafter\def\csname cijAC\endcsname{-2{,}02}
\expandafter\def\csname pijAC\endcsname{0{,}0593}
\expandafter\def\csname cijAD\endcsname{-0{,}11}
\expandafter\def\csname pijAD\endcsname{0{,}0286}
\expandafter\def\csname sumcA\endcsname{6{,}60}
\expandafter\def\csname cijBA\endcsname{-1{,}10}
\expandafter\def\csname pijBA\endcsname{0{,}0312}
\expandafter\def\csname cijBB\endcsname{5{,}95}
\expandafter\def\csname pijBB\endcsname{0{,}1832}
\expandafter\def\csname cijBC\endcsname{-0{,}66}
\expandafter\def\csname pijBC\endcsname{0{,}0498}
\expandafter\def\csname cijBD\endcsname{-0{,}59}
\expandafter\def\csname pijBD\endcsname{0{,}0393}
\expandafter\def\csname sumcB\endcsname{3{,}60}
\expandafter\def\csname cijCA\endcsname{-2{,}02}
\expandafter\def\csname pijCA\endcsname{0{,}0593}
\expandafter\def\csname cijCB\endcsname{-0{,}66}
\expandafter\def\csname pijCB\endcsname{0{,}0498}
\expandafter\def\csname cijCC\endcsname{5{,}41}
\expandafter\def\csname pijCC\endcsname{0{,}2555}
\expandafter\def\csname cijCD\endcsname{-2{,}18}
\expandafter\def\csname pijCD\endcsname{0{,}1056}
\expandafter\def\csname sumcC\endcsname{0{,}55}
\expandafter\def\csname cijDA\endcsname{-0{,}11}
\expandafter\def\csname pijDA\endcsname{0{,}0286}
\expandafter\def\csname cijDB\endcsname{-0{,}59}
\expandafter\def\csname pijDB\endcsname{0{,}0393}
\expandafter\def\csname cijDC\endcsname{-2{,}18}
\expandafter\def\csname pijDC\endcsname{0{,}1056}
\expandafter\def\csname cijDD\endcsname{5{,}62}
\expandafter\def\csname pijDD\endcsname{0{,}2237}
\expandafter\def\csname sumcD\endcsname{2{,}73}

\newcommand{\cc}[2]{\csname cij#1#2\endcsname}
\newcommand{\pp}[2]{\csname pij#1#2\endcsname}
\begin{table}[t]
\centering
\caption{Coeficientes del sistema de la \cref{fig:sistema-cij}, calculados numéricamente (por unidad de longitud):
A~cuña, B~rectángulo, C~disco y D~barra. A la izquierda, $c_{ij}/\eo$ y la suma de cada columna; a la derecha, $\eo\,p_{ij}$, la matriz inversa.}
\label{tab:sistema-coef}
\small
\begin{tabular}{@{}crrrr@{}}
  \toprule
  $c_{ij}/\eo$ & A & B & C & D\\
  \midrule
  A & \textbf{\cc AA} & \cc AB & \cc AC & \cc AD\\
  B & \cc BA & \textbf{\cc BB} & \cc BC & \cc BD\\
  C & \cc CA & \cc CB & \textbf{\cc CC} & \cc CD\\
  D & \cc DA & \cc DB & \cc DC & \textbf{\cc DD}\\
  \midrule
  $\sum_i c_{ij}/\eo$ & \csname sumcA\endcsname & \csname sumcB\endcsname & \csname sumcC\endcsname & \csname sumcD\endcsname\\
  \bottomrule
\end{tabular}
\qquad
\begin{tabular}{@{}cllll@{}}
  \toprule
  $\eo\,p_{ij}$ & A & B & C & D\\
  \midrule
  A & \textbf{\pp AA} & \pp AB & \pp AC & \pp AD\\
  B & \pp BA & \textbf{\pp BB} & \pp BC & \pp BD\\
  C & \pp CA & \pp CB & \textbf{\pp CC} & \pp CD\\
  D & \pp DA & \pp DB & \pp DC & \textbf{\pp DD}\\
  \bottomrule
\end{tabular}
\end{table}


% =========================================================================
\section{Condensadores}\label{sec:condensadores}
% =========================================================================
El sistema de dos conductores en influencia total tiene un nombre propio y una importancia práctica enorme: es el condensador, el dispositivo que
almacena carga (y, como veremos, energía) con una diferencia de potencial dada.

% -------------------------------------------------------------------------
\subsection{Definición y capacidad}\label{sec:condensador-def}
% -------------------------------------------------------------------------
\begin{definicion}[Condensador]\label{def:condensador}
Un \emph{condensador} es un sistema de dos conductores en influencia total cargados con cargas de igual módulo y signo contrario, $Q_1=Q$ y
$Q_2=-Q$ ($Q>0$). Cada uno de los conductores se llama \emph{placa} o \emph{armadura}. Su \emph{capacidad} es el cociente entre el valor absoluto de la carga
común y la diferencia de potencial entre las armaduras:
\begin{equation}\label{eq:cap-condensador}
  \resalta{C=\frac{Q}{\phi_1-\phi_2}=\frac{Q}{\DV} .}
\end{equation}
\end{definicion}

\begin{teorema}[Capacidad en función de los coeficientes]\label{teo:cap-coef}
La capacidad de un condensador solo depende de la geometría de las armaduras y del medio, es positiva, y vale
\begin{equation}\label{eq:cap-coef}
  \resalta{C=\frac{1}{p_{11}+p_{22}-2p_{12}}=\frac{c_{11}c_{22}-c_{12}^{2}}{c_{11}+c_{22}+2c_{12}} .}
\end{equation}
\end{teorema}
\begin{proof}
Haciendo $Q_1=Q$ y $Q_2=-Q$ en las ecuaciones de los coeficientes de potencial \cref{eq:pij},
\[
  \left.\begin{aligned}
    \phi_1&=p_{11}Q_1+p_{12}Q_2=(p_{11}-p_{12})\,Q\\
    \phi_2&=p_{21}Q_1+p_{22}Q_2=(p_{21}-p_{22})\,Q
  \end{aligned}\right\}
  \;\Longrightarrow\;\phi_1-\phi_2=\bigl(p_{11}+p_{22}-2p_{12}\bigr)Q ,
\]
donde hemos usado $p_{21}=p_{12}$. Dividiendo, $C=Q/(\phi_1-\phi_2)=(p_{11}+p_{22}-2p_{12})^{-1}$, que solo depende de la geometría. Es positiva porque
$p_{11}+p_{22}-2p_{12}=\bm{v}\Tt\Pm\,\bm{v}$ con $\bm{v}=(1,-1)\Tt$, y $\Pm$ es definida positiva (\cref{lem:definida}). Para la segunda forma basta sustituir
los elementos de la inversa de $\Cm$: con $D=\det\Cm=c_{11}c_{22}-c_{12}^{2}$, $p_{11}=c_{22}/D$, $p_{22}=c_{11}/D$ y $p_{12}=-c_{12}/D$, de modo que
$p_{11}+p_{22}-2p_{12}=(c_{11}+c_{22}+2c_{12})/D$.
\end{proof}

\begin{corolario}[Conductor encerrado]\label{cor:encerrado}
Si la armadura~1 está encerrada por la armadura~2, entonces
\begin{equation}\label{eq:cap-encerrado}
  C=c_{11}=-c_{12}=-c_{21} .
\end{equation}
En particular, con la armadura exterior a tierra ($\phi_2=0$) y la interior a potencial $V$, se tiene $Q_1=c_{11}V$ y $Q_2=c_{21}V=-Q_1$: es exactamente un
condensador cargado a $\DV=V$, y $C=\tfrac12(c_{11}-c_{21})$.
\end{corolario}
\begin{proof}
Con $\phi_2=0$ y $\phi_1=V$, \cref{eq:cij} da $Q_1=c_{11}V$ y $Q_2=c_{21}V$. Al estar 1 encerrado por 2, todas sus líneas acaban en 2 (influencia total,
\cref{teo:influencia}): $Q_2=-Q_1$, es decir, $c_{21}=-c_{11}$. Entonces $C=Q_1/(\phi_1-\phi_2)=c_{11}$. Restando las dos ecuaciones,
$2Q=(c_{11}-c_{21})V$, que es la última forma. (Esta última expresión \emph{no} es general: solo vale cuando $\phi_2=0$ y hay influencia total.)
\end{proof}

\begin{observacion}[Condensador «generalizado»]\label{obs:generalizado}
La fórmula~\eqref{eq:cap-coef} tiene sentido aunque los dos conductores no estén en influencia total (por ejemplo, dos esferas próximas): es la capacidad
que se mide al transferir una carga $Q$ de un conductor al otro. En un condensador de verdad, en cambio, la influencia total garantiza que el campo
fuera de las armaduras es nulo y que la capacidad no depende de los objetos del entorno.
\end{observacion}

% -------------------------------------------------------------------------
\subsection{Ejemplo: condensador esférico}\label{sec:ej-esferico}
% -------------------------------------------------------------------------
\begin{ejemplo}[Condensador esférico de armaduras delgadas]\label{ej:esferico}
Una esfera conductora de radio $a$ (armadura~1) está rodeada por una cáscara esférica delgada y concéntrica de radio $b>a$ (armadura~2). Calculemos sus
coeficientes y su capacidad.

\emph{Coeficientes de potencial.} Por el \cref{ej:concentricas} (o la regla~\eqref{eq:regla-concentricas}),
\[
  \Pm=\frac{1}{4\pi\eo}\begin{pmatrix}1/a & 1/b\\ 1/b & 1/b\end{pmatrix},\qquad
  C=\frac{1}{p_{11}+p_{22}-2p_{12}}=4\pi\eo\Bigl(\frac1a+\frac1b-\frac2b\Bigr)^{-1}=4\pi\eo\Bigl(\frac1a-\frac1b\Bigr)^{-1},
\]
\begin{equation}\label{eq:cap-esferico}
  \resalta{C=4\pi\eo\,\frac{ab}{b-a} .}
\end{equation}

\emph{Coeficientes de capacidad e influencia.} El determinante es
\[
  \det\Pm=\Bigl(\frac{1}{4\pi\eo}\Bigr)^{2}\Bigl(\frac1a\,\frac1b-\frac1b\,\frac1b\Bigr)=\Bigl(\frac{1}{4\pi\eo}\Bigr)^{2}\frac{b-a}{ab^{2}} ,
\]
y la inversa de una matriz $2\times2$ es $\Pm^{-1}=\frac{1}{\det\Pm}\bigl(\begin{smallmatrix}p_{22}&-p_{12}\\-p_{12}&p_{11}\end{smallmatrix}\bigr)$:
\[
  \Cm=(4\pi\eo)^{2}\,\frac{ab^{2}}{b-a}\,\frac{1}{4\pi\eo}\begin{pmatrix}1/b & -1/b\\ -1/b & 1/a\end{pmatrix}
  =4\pi\eo\,\frac{ab}{b-a}\begin{pmatrix}1 & -1\\ -1 & b/a\end{pmatrix}.
\]
Se comprueban todas las propiedades: $\Cm$ es simétrica, $c_{11},c_{22}>0$ y $c_{12}<0$. Además, $c_{11}=-c_{21}=C$: la esfera interior está en influencia
total (\cref{cor:encerrado}). En cambio, $c_{22}+c_{12}=4\pi\eo\,\frac{b^{2}-ab}{b-a}=4\pi\eo\,b>0$: parte de las líneas que salen de la cáscara se van al
infinito, y lo que «sobra» es justamente la capacidad de una esfera aislada de radio $b$.

\emph{Placas próximas.} Si $b=a+d$ con $d\ll a$,
$C=4\pi\eo\,\frac{a(a+d)}{d}\simeq4\pi\eo\,\frac{a^{2}}{d}=\eo\,\frac{A}{d}$ con $A=4\pi a^{2}$: el condensador esférico se comporta como uno plano de
área $A$ y separación $d$ (\cref{sem:plano}).
\end{ejemplo}

% =========================================================================
%  Figura: condensador esférico (b = 2a).
%  (a) Sección: +Q en la esfera, -Q en la cara interior de la cáscara, campo
%      radial entre armaduras y gaussiana a < r < b.
%  (b) E/E0 frente a r/a, con E0 = kQ/a^2: el área sombreada es DeltaV/phi0.
%  (c) phi/phi0 con la armadura exterior a tierra, phi0 = kQ/a.
% =========================================================================
\begin{figure}[t]
\centering
\begin{tikzpicture}[font=\small]
  % ---------------------------------------------------------------- (a)
  \begin{scope}[xshift=7.6cm, x=0.85cm, y=0.85cm]
    \foreach \a in {15,45,...,345}{
      \draw[azulacero, line width=0.55pt, -{Stealth[length=3.5pt]}] (\a:1.08) -- (\a:1.6);
      \draw[azulacero, line width=0.55pt] (\a:1.6) -- (\a:1.9);}
    \draw[azulprof, line width=1pt] (0,0) circle (2);
    \shade[ball color=azulvelo] (0,0) circle (1);
    \draw[azulprof, line width=0.6pt] (0,0) circle (1);
    \foreach \a in {0,45,...,315}{\node[signopos, font=\tiny\bfseries] at (\a:0.82) {$+$};
                                 \node[signoneg, font=\tiny\bfseries] at (\a+22.5:1.86) {$-$};}
    \draw[gauss] (0,0) circle (1.45);
    \node[font=\footnotesize, text=cargapos] at (0,0.3) {$+Q$};
    \node[font=\footnotesize] at (0,-0.25) {$a$};
    \node[font=\footnotesize, text=carganeg] at (2.25,1.55) {$-Q$};
    \node[font=\footnotesize] at (2.05,-1.5) {$b$};
    \node[font=\footnotesize, text=grispizarra, align=center] at (0,-2.55) {$\vE=\vc{0}$ fuera de $b$};
    \node at (0,-3.2) {(a)};
  \end{scope}
  % ---------------------------------------------------------------- (b)
  \begin{scope}[xshift=1.2cm, yshift=-7.0cm]
    \begin{axis}[tema, width=6.6cm, height=4.6cm, xmin=0, xmax=3.2, ymin=0, ymax=1.15,
        xlabel={$r/a$}, ylabel={$E/E_0$}, xtick={0,1,2,3}, ytick={0,0.25,0.5,1}]
      \addplot[draw=none, fill=azulvelo, domain=1:2, samples=40] {1/x^2} \closedcycle;
      \addplot[azulprof, line width=1.1pt, domain=1:2, samples=40] {1/x^2};
      \addplot[azulprof, line width=1.1pt] coordinates {(0,0) (1,0)};
      \addplot[azulprof, line width=1.1pt] coordinates {(2,0) (3.2,0)};
      \draw[punteado] (axis cs:1,0) -- (axis cs:1,1) (axis cs:2,0) -- (axis cs:2,0.25);
      \node[font=\scriptsize] at (axis cs:1.45,0.18) {$\DV$};
      \node[font=\scriptsize, text=grispizarra] at (axis cs:0.5,0.12) {$E=0$};
      \node[font=\scriptsize, text=grispizarra] at (axis cs:2.6,0.12) {$E=0$};
    \end{axis}
    \node at (2.6,-1.35) {(b)};
  \end{scope}
  % ---------------------------------------------------------------- (c)
  \begin{scope}[xshift=9.0cm, yshift=-7.0cm]
    \begin{axis}[tema, width=6.6cm, height=4.6cm, xmin=0, xmax=3.2, ymin=0, ymax=0.62,
        xlabel={$r/a$}, ylabel={$\phi/\phi_0$}, xtick={0,1,2,3}, ytick={0,0.25,0.5}]
      \addplot[azulacero, line width=1.1pt] coordinates {(0,0.5) (1,0.5)};
      \addplot[azulacero, line width=1.1pt, domain=1:2, samples=40] {1/x-0.5};
      \addplot[azulacero, line width=1.1pt] coordinates {(2,0) (3.2,0)};
      \draw[punteado] (axis cs:1,0) -- (axis cs:1,0.5) (axis cs:2,0) -- (axis cs:2,0.6);
      \draw[{Stealth[length=3pt]}-{Stealth[length=3pt]}, cargapos] (axis cs:2.35,0) -- (axis cs:2.35,0.5);
      \draw[punteado] (axis cs:1,0.5) -- (axis cs:2.45,0.5);
      \node[font=\scriptsize, text=cargapos, anchor=west] at (axis cs:2.4,0.27) {$\DV$};
    \end{axis}
    \node at (2.6,-1.35) {(c)};
  \end{scope}
\end{tikzpicture}
\caption{Condensador esférico con $b=2a$. (a) Sección: toda la carga $-Q$ de la armadura exterior está en su cara interior, y la gaussiana entre armaduras
solo encierra $+Q$. (b) Campo, con $E_0=\kc Q/a^{2}$: es nulo dentro de la esfera y fuera de la cáscara, y el área bajo la curva entre $a$ y $b$ es
$\DV/\phi_0$. (c) Potencial con la armadura exterior a tierra, $\phi_0=\kc Q/a$: constante en la esfera y nulo desde la cáscara hacia fuera.}
\label{fig:condensador-esferico}
\end{figure}


\begin{ejemplo}[El mismo sistema, en otras situaciones]\label{ej:esferico-situaciones}
Los coeficientes resuelven cualquier situación del sistema, no solo la del condensador.

\emph{Esfera interior a tierra.} Si $\phi_1=0$ y $\phi_2\neq0$, $Q_2=c_{22}\phi_2$ y $Q_1=c_{12}\phi_2$, de modo que
\[
  \frac{Q_1}{Q_2}=\frac{c_{12}}{c_{22}}=-\frac{a}{b} :
\]
el sistema \emph{no} está en influencia total respecto de la cáscara, y la esfera interior adquiere menos carga (en valor absoluto) que la exterior. Si
$b=2a$, $Q_1=-Q_2/2$: la cara interior de la cáscara lleva $+Q_2/2$ y la exterior $+Q_2/2$.

\emph{Cargas dadas.} Si conocemos $Q_1$ y $Q_2$, los potenciales salen directamente con los coeficientes de potencial,
\[
  \phi_1=p_{11}Q_1+p_{12}Q_2=\frac{1}{4\pi\eo}\Bigl(\frac{Q_1}{a}+\frac{Q_2}{b}\Bigr),\qquad
  \phi_2=p_{21}Q_1+p_{22}Q_2=\frac{1}{4\pi\eo}\,\frac{Q_1+Q_2}{b} ,
\]
mientras que con los coeficientes de capacidad habría que resolver un sistema lineal. Es útil comprobarlo resolviendo el problema «como en el Tema~1»:
por Gauss,
\[
  \vE(\vc{r})=\begin{cases}\dfrac{Q_1+Q_2}{4\pi\eo r^{2}}\,\ur, & r>b,\\[8pt] \dfrac{Q_1}{4\pi\eo r^{2}}\,\ur, & a<r<b,\\[8pt] \vc{0}, & r<a,\end{cases}
\]
con un salto de $\vE$ en $r=b$ debido a la carga de la cáscara. El potencial es continuo, con $\phi(\infty)=0$: para $r\ge b$,
$\phi=\frac{Q_1+Q_2}{4\pi\eo r}$; para $a\le r\le b$,
\[
  \phi(r)-\phi(b)=-\int_{b}^{r}\vE\cdot\dif\vc{r}'=\int_{r}^{b}\frac{Q_1}{4\pi\eo r'^{2}}\,\dif r'=\frac{Q_1}{4\pi\eo}\Bigl(\frac1r-\frac1b\Bigr)
  \;\Longrightarrow\;\phi(r)=\frac{1}{4\pi\eo}\Bigl(\frac{Q_1}{r}+\frac{Q_2}{b}\Bigr),
\]
y para $r\le a$, $\phi=\frac{1}{4\pi\eo}\bigl(\frac{Q_1}{a}+\frac{Q_2}{b}\bigr)$, que coincide con $\phi_1$. Con los coeficientes de potencial el cálculo es inmediato; con los de
capacidad, para pasar de cargas a potenciales, habría que invertir la matriz.
\end{ejemplo}

% -------------------------------------------------------------------------
\subsection{Ejemplo: condensador esférico con armadura exterior gruesa}\label{sec:ej-grueso}
% -------------------------------------------------------------------------
\begin{ejemplo}[Esfera y corteza esférica gruesa]\label{ej:grueso}
El conductor~1 es una esfera de radio $a$; el conductor~2, una corteza esférica concéntrica de radio interior $b$ y exterior $c$. Escribimos
$\kc=1/(4\pi\eo)$.

\emph{Caso A: $Q_1=Q$, $Q_2=0$.} Por influencia total, la cara interior de la corteza adquiere $-Q$ (Gauss con una superficie dentro del metal) y, por
conservación de la carga, la exterior $+Q$. Fuera de la corteza el campo es el de una carga $Q$ en el centro, y $\phi(\vc{r})=-\int_{\infty}^{\vc{r}}\vE\cdot\dl$
da, sobre la corteza y sobre la esfera,
\[
  \phi_2=\int_{c}^{\infty}\frac{\kc Q}{r^{2}}\,\dif r=\frac{\kc Q}{c},\qquad
  \phi_1=\phi_2+\int_{a}^{b}\frac{\kc Q}{r^{2}}\,\dif r=\frac{\kc Q}{c}+\kc Q\Bigl(\frac1a-\frac1b\Bigr),
\]
luego $p_{11}=\kc\bigl(\frac1a-\frac1b+\frac1c\bigr)$ y $p_{21}=\kc/c$.

\emph{Caso B: $Q_1=0$, $Q_2=Q$.} La carga se reparte en la cara exterior de la corteza; dentro el campo es nulo y la esfera interior queda al mismo potencial
que la corteza, $\phi_1=\phi_2=\kc Q/c$. Así $p_{12}=p_{22}=\kc/c$, y
\[
  \Pm=\frac{1}{4\pi\eo}\begin{pmatrix}\dfrac1a-\dfrac1b+\dfrac1c & \dfrac1c\\[8pt] \dfrac1c & \dfrac1c\end{pmatrix}.
\]

\emph{Capacidad del condensador.}
\[
  C=\frac{Q}{\phi_1-\phi_2}=\frac{1}{p_{11}+p_{22}-2p_{12}}=\frac{4\pi\eo}{\frac1a-\frac1b+\frac1c+\frac1c-\frac2c}=4\pi\eo\,\frac{ab}{b-a} ,
\]
la misma que con armaduras delgadas: el espesor de la armadura exterior no influye, porque entre las armaduras el campo solo depende de $Q$. Con $d=b-a$,
$C=\eo\,A_{\mathrm{ef}}/d$ con $A_{\mathrm{ef}}=4\pi ab=\sqrt{4\pi a^{2}\cdot4\pi b^{2}}$, la media geométrica de las áreas de las armaduras.

\emph{Inversión de la matriz.} Con $\det\Pm=p_{11}p_{22}-p_{12}^{2}=\kc^{2}\,\frac{b-a}{abc}$ y la matriz adjunta,
$\Pm^{-1}=\frac{1}{\det\Pm}\,[\adj\Pm]\Tt=\frac{1}{\det\Pm}\bigl(\begin{smallmatrix}p_{22}&-p_{12}\\-p_{21}&p_{11}\end{smallmatrix}\bigr)$. Como $\Pm$ es simétrica, la
matriz de cofactores también lo es, y trasponer no cambia nada: por eso $\Cm$ sale simétrica automáticamente. Resulta
\[
  c_{11}=\frac{p_{22}}{\det\Pm}=4\pi\eo\,\frac{ab}{b-a},\qquad c_{12}=c_{21}=-\frac{p_{12}}{\det\Pm}=-4\pi\eo\,\frac{ab}{b-a},
\]
\[
  c_{22}=\frac{p_{11}}{\det\Pm}=4\pi\eo\Bigl(c+\frac{ab}{b-a}\Bigr).
\]
\emph{Interpretación.} $c_{11}>0$ y $c_{12}<0$; $c_{11}=-c_{21}$ es la igualdad de la influencia total, pues todas las líneas de la esfera interior acaban en la
corteza; en cambio $c_{22}>-c_{12}$, porque parte de las líneas de la corteza van al infinito: $c_{22}=C+4\pi\eo c$, la capacidad del condensador más la de
una esfera aislada de radio $c$.
\end{ejemplo}

% =========================================================================
%  Figura: esfera (a) y corteza gruesa (b, c), con a = 1, b = 2, c = 2,5.
%  Potencial normalizado a kQ/a en los dos estados elementales:
%  Caso A (Q1 = Q, Q2 = 0): 1/x (x >= c); 1/c (b <= x <= c);
%                           1/c + 1/x - 1/b (a <= x <= b); 1/c + 1/a - 1/b (x <= a)
%  Caso B (Q1 = 0, Q2 = Q): 1/x (x >= c); 1/c (x <= c)
% =========================================================================
\begin{figure}[t]
\centering
\begin{tikzpicture}[font=\small]
  % ---------------------------------------------------------------- (a)
  \begin{scope}[x=0.72cm, y=0.72cm]
    \fill[metal!65!white, even odd rule] (0,0) circle (2.5) (0,0) circle (2);
    \fill[pattern=north east lines, pattern color=azulcielo!70, even odd rule] (0,0) circle (2.5) (0,0) circle (2);
    \draw[azulprof, line width=0.6pt] (0,0) circle (2.5) (0,0) circle (2);
    \shade[ball color=azulvelo] (0,0) circle (1);
    \draw[azulprof, line width=0.6pt] (0,0) circle (1);
    \draw[-{Stealth[length=3.5pt]}, azulprof, line width=0.5pt] (0,0) -- (90:1) node[pos=0.5, right=-1pt, font=\scriptsize] {$a$};
    \draw[-{Stealth[length=3.5pt]}, azulprof, line width=0.5pt] (0,0) -- (210:2) node[pos=0.6, above left=-2pt, font=\scriptsize] {$b$};
    \draw[-{Stealth[length=3.5pt]}, azulprof, line width=0.5pt] (0,0) -- (-30:2.5) node[pos=0.8, above right=-2pt, font=\scriptsize] {$c$};
    \node[font=\footnotesize] at (0,-0.35) {$1$};
    \node[font=\footnotesize, fill=white, inner sep=1pt] at (0,2.25) {$2$};
    \node at (0,-3.45) {(a)};
  \end{scope}
  % ---------------------------------------------------------------- (b)
  \begin{scope}[xshift=4.7cm, yshift=-2.1cm]
    \begin{axis}[tema, width=6.4cm, height=4.6cm, xmin=0, xmax=4.5, ymin=0, ymax=1.02,
        xlabel={$r/a$}, ylabel={$\phi\,/\,(\kc Q/a)$}, xtick={0,1,2,2.5,4}, xticklabels={$0$,$1$,$2$,$2{,}5$,$4$},
        ytick={0,0.4,0.9}, legend style={at={(0.98,0.97)}, anchor=north east}]
      \fill[metal!45!white] (axis cs:2,0) rectangle (axis cs:2.5,1.02);
      \addplot[azulprof, line width=1.1pt] coordinates {(0,0.9) (1,0.9)};
      \addlegendentry{caso A: $Q_1=Q$}
      \addplot[azulprof, line width=1.1pt, domain=1:2, samples=40, forget plot] {0.4+1/x-0.5};
      \addplot[azulprof, line width=1.1pt, forget plot] coordinates {(2,0.4) (2.5,0.4)};
      \addplot[azulprof, line width=1.1pt, domain=2.5:4.5, samples=40, forget plot] {1/x};
      \addplot[azulcielo, line width=1.1pt, densely dashed] coordinates {(0,0.4) (2.5,0.4)};
      \addlegendentry{caso B: $Q_2=Q$}
      \addplot[azulcielo, line width=1.1pt, densely dashed, domain=2.5:4.5, samples=40, forget plot] {1/x};
      \draw[{Stealth[length=3pt]}-{Stealth[length=3pt]}, cargapos] (axis cs:0.5,0.4) -- (axis cs:0.5,0.9)
        node[midway, right, font=\scriptsize] {$\phi_1-\phi_2$};
    \end{axis}
    \node at (2.45,-1.35) {(b)};
  \end{scope}
\end{tikzpicture}
\caption{Condensador esférico con armadura exterior gruesa ($a=1$, $b=2$, $c=2{,}5$). (a) Geometría: la corteza, entre $b$ y $c$, es un volumen
equipotencial. (b) Potencial en los dos estados que definen los coeficientes de potencial: con la esfera cargada (caso A) el potencial cae entre $a$ y $b$, y
la diferencia $\phi_1-\phi_2$ solo depende de $a$ y $b$; con la corteza cargada (caso B) el interior queda a potencial constante $\kc Q/c$.}
\label{fig:corteza-gruesa}
\end{figure}


% -------------------------------------------------------------------------
\subsection{Los tres condensadores canónicos}\label{sec:canonicos}
% -------------------------------------------------------------------------
Para geometrías indefinidas (cilindros o planos infinitos) no se puede tomar $\phi(\infty)=0$, porque el potencial de una distribución infinita de carga
diverge; los coeficientes de potencial solo pueden definirse respecto de una referencia arbitraria (\cref{sem:cilindrico}). En un condensador, sin
embargo, esto no supone ningún problema: el campo fuera de las armaduras es nulo y solo interesa la diferencia de potencial, que es finita. La
\cref{tab:canonicos} reúne los tres casos clásicos, que se deducen en detalle en los seminarios (\cref{sem:esferico,sem:cilindrico,sem:plano}).

\begin{table}[ht]
\centering
\caption{Los tres condensadores canónicos. El campo entre las armaduras se obtiene con la ley de Gauss, la diferencia de potencial integrándolo de la
armadura positiva a la negativa, y la capacidad como $C=Q/\DV$.}
\label{tab:canonicos}
\begin{tabular}{@{}lccc@{}}
  \toprule
  Condensador & Campo entre armaduras & $\DV$ & Capacidad\\
  \midrule
  Esférico (radios $a<b$) & $\dfrac{Q}{4\pi\eo r^{2}}$ & $\dfrac{Q}{4\pi\eo}\Bigl(\dfrac1a-\dfrac1b\Bigr)$ & $4\pi\eo\,\dfrac{ab}{b-a}$\\[12pt]
  Cilíndrico (radios $a<b$, longitud $\ell$) & $\dfrac{\lambda}{2\pi\eo r}$ & $\dfrac{\lambda}{2\pi\eo}\ln\dfrac{b}{a}$ & $\dfrac{2\pi\eo\,\ell}{\ln(b/a)}$\\[12pt]
  Plano (área $A$, separación $d$) & $\dfrac{Q}{\eo A}$ & $\dfrac{Qd}{\eo A}$ & $\eo\,\dfrac{A}{d}$\\[8pt]
  \bottomrule
\end{tabular}
\end{table}

\begin{observacion}[Un límite común]\label{obs:limite-plano}
Cuando la separación entre armaduras es pequeña frente a sus radios de curvatura, los tres casos tienden a la misma expresión $C\simeq\eo A/d$: localmente,
dos superficies muy próximas se ven planas. Para el cilíndrico, con $b=a+d$, $\ln(1+d/a)\simeq d/a$ y $C/\ell\simeq2\pi\eo a/d=\eo(A/\ell)/d$.
\end{observacion}

% =========================================================================
\section{Seminarios prácticos}\label{sec:seminarios}
% =========================================================================
Esta sección resuelve, con todo detalle, los siete problemas de la relación del tema~\cite{PROBLEMAS-T3}, más un ejemplo modelo del método de los
coeficientes de potencial~\cite{GUIA-T3}. Cada problema sigue la misma estructura que en los temas anteriores: el \emph{enunciado}, la \emph{idea} que guía
la resolución, la \emph{resolución} paso a paso y un comentario sobre la \emph{interpretación} del resultado y sus comprobaciones. Salvo que se diga lo
contrario, $\phi(\infty)=0$ y $\kc=1/(4\pi\eo)$.

% -------------------------------------------------------------------------
\subsection{Cómo plantear un problema de conductores}\label{sem:metodo}
% -------------------------------------------------------------------------
Todo problema de conductores se reduce a un sistema lineal: \emph{cada conductor aporta una ecuación y una incógnita}, porque de cada uno se conoce o su
carga o su potencial. La primera tarea es, pues, traducir el enunciado (\cref{tab:traduccion}); la segunda, escoger el método (\cref{tab:mapa}).

\begin{table}[ht]
\centering
\caption{Traducción de la situación física a ecuaciones.}
\label{tab:traduccion}
\small
\begin{tabularx}{\linewidth}{@{}>{\raggedright\arraybackslash}p{4.1cm}>{\raggedright\arraybackslash}X>{\raggedright\arraybackslash}p{4.3cm}@{}}
  \toprule
  Situación & Dato & Incógnita\\
  \midrule
  Aislado & $Q_i$ conocida (se conserva) & $\phi_i$\\
  Conectado a tierra & $\phi_i=0$ & $Q_i$ (fluye carga desde o hacia tierra)\\
  Unido a $j$ por un hilo fino & $\phi_i=\phi_j$ y $Q_i+Q_j$ conservada & el reparto $Q_i$, $Q_j$\\
  Batería $V_0$ entre $i$ y $j$ & $\phi_i-\phi_j=V_0$ y $Q_i+Q_j$ conservada & $Q_i$ (y $Q_j=-Q_i$ si eran neutros)\\
  Se conecta y se desconecta & tras desconectar, la carga queda congelada & la carga de cada etapa\\
  \bottomrule
\end{tabularx}
\end{table}

\begin{table}[ht]
\centering
\caption{Mapa de decisión: qué método conviene según lo que pide el enunciado.}
\label{tab:mapa}
\small
\begin{tabularx}{\linewidth}{@{}>{\raggedright\arraybackslash}X>{\raggedright\arraybackslash}p{5.6cm}c@{}}
  \toprule
  Lo que pide el enunciado & Método & Problemas\\
  \midrule
  Potencial o campo de un conductor aislado; ruptura dieléctrica; efecto punta & M1. Conductor aislado & 3.1\\
  Esferas o cilindros concéntricos con cargas, tierras o baterías & M2. Cadena de Gauss y potenciales & 3.2, 3.4\\
  «Calcula los coeficientes de capacidad e influencia» & M3. Coeficientes (caminos A y B) & 3.3\\
  Placas planas paralelas; cargas en cada cara & M4. Superposición de planos & 3.5\\
  Esferas pequeñas y lejanas ($d\gg R$); tierras sucesivas & M5. Esferas lejanas con $p_{ij}$ aproximados & 3.6, 3.7\\
  \bottomrule
\end{tabularx}
\end{table}

\begin{receta}[Los cinco métodos]
\begin{enumerate}[label=\textbf{M\arabic*.}, leftmargin=2.2em]
  \item \emph{Conductor aislado.} $\phi=\kc Q/R$, $E_{\mathrm{sup}}=\sigma/\eo=\phi/R$ (esfera).
  \item \emph{Concéntricos.} (1) Cargas en cada cara con la regla de las caras (\cref{prop:caras}); (2) Gauss en cada hueco: $E=\kc\Qenc/r^{2}$ (o
        $\lambda_{\mathrm{enc}}/2\pi\eo r$), con $\Qenc$ la suma de las cargas netas interiores, y $\vE=\vc{0}$ dentro del metal; (3) potencial del más
        externo, $\phi_N=\kc Q_{\mathrm{total}}/R_N$; (4) cadena hacia dentro, $\phi_k-\phi_{k+1}=\kc\Qenc\bigl(\frac{1}{R_k}-\frac{1}{R_{k+1}}\bigr)$;
        (5) imponer tierras, baterías y aislamientos.
  \item \emph{Coeficientes.} Camino A (directo): $\phi_j=1$, los demás a tierra; las cargas obtenidas son la columna $j$ de $\Cm$. Camino B: $Q_j=1$,
        los demás descargados; los potenciales son la columna $j$ de $\Pm$; luego se invierte.
  \item \emph{Placas paralelas.} Cada cara es un plano que crea $\pm\sigma/2\eo$; se impone conservación de carga y $\vE=\vc{0}$ dentro de cada placa.
  \item \emph{Esferas lejanas.} $p_{ii}\simeq\kc/R_i$, $p_{ij}\simeq\kc/d_{ij}$; en procesos por etapas, solo cambia la carga de la esfera conectada.
\end{enumerate}
\end{receta}

% =========================================================================
\subsection{Problema 3.1. Globo conductor y campo de ruptura}\label{sem:globo}
% =========================================================================
\begin{problema}[Potencial máximo de un globo]\label{prob:globo}
Supongamos un globo esférico de radio $5\,\mathrm{cm}$ recubierto por una fina película conductora. Si el campo electrostático de ruptura en la atmósfera
es de $3\cdot10^{6}\,\mathrm{V/m}$, calcula el máximo potencial al que puede conectarse el globo.
\end{problema}

\begin{resolucion}
\paso{Idea.} La \emph{ruptura dieléctrica} es el fenómeno por el que un aislante (aquí el aire) se ioniza y se convierte en un medio conductor cuando el
campo supera un valor crítico $\Erup$: el globo empezaría a descargarse a través del aire. Como el campo es máximo justo en la superficie del conductor, la
condición de ruptura hay que imponerla allí. Es el método M1.

\paso{Paso 1: campo y potencial del globo.} La película conductora es una esfera conductora de radio $R$. Con carga $Q$, la carga se reparte
uniformemente y, por Gauss, el campo exterior es el de una carga puntual en el centro:
\[
  \vE(\vc{r})=\frac{1}{4\pi\eo}\,\frac{Q}{r^{2}}\,\ur\quad(r>R),\qquad E(R)=\frac{1}{4\pi\eo}\,\frac{Q}{R^{2}},\qquad
  \phi=\frac{1}{4\pi\eo}\,\frac{Q}{R} .
\]
\paso{Paso 2: relación entre campo superficial y potencial.} Dividiendo $E(R)$ entre $\phi$ desaparecen la carga y la constante:
\[
  \frac{E(R)}{\phi}=\frac{\kc Q/R^{2}}{\kc Q/R}=\frac{R}{R^{2}}=\frac1R\quad\Longrightarrow\quad\resalta{\phi=E(R)\,R .}
\]
\paso{Paso 3: condición de ruptura.} Como $E(r)$ decrece fuera del globo, el campo más intenso está en la superficie. Si $E(R)$ supera $\Erup$, el aire se
ioniza y el globo se descarga. Por tanto,
\[
  \phi_{\max}=\Erup\,R=3\cdot10^{6}\,\frac{\mathrm{V}}{\mathrm{m}}\cdot5\cdot10^{-2}\,\mathrm{m}=1{,}5\cdot10^{5}\,\mathrm{V}
  \qquad\Longrightarrow\qquad\boxed{\;\phi_{\max}=150\,\mathrm{kV}\;}
\]
\end{resolucion}

% =========================================================================
%  Figura (Problema 3.1): globo conductor en el umbral de ruptura.
%  (a) E/E_rup y phi/phi_max frente a r/R: E = (R/r)^2, phi = R/r fuera.
%  (b) phi_max = E_rup R frente a R (escala logarítmica), E_rup = 3e6 V/m.
% =========================================================================
\begin{figure}[t]
\centering
\begin{tikzpicture}[font=\small]
  \begin{axis}[tema, name=a, width=7.0cm, height=5.0cm, xmin=0, xmax=4.2, ymin=0, ymax=1.25,
      xlabel={$r/R$}, ylabel={magnitud normalizada}, xtick={0,1,2,3,4}, ytick={0,0.5,1},
      legend style={at={(0.97,0.72)}, anchor=north east}]
    \fill[azulvelo!60] (axis cs:0,0) rectangle (axis cs:1,1.25);
    \node[font=\scriptsize, text=azulprof, rotate=90] at (axis cs:0.5,0.62) {globo};
    \addplot[azulprof, line width=1.2pt, domain=1:4.2, samples=60] {1/x^2};
    \addlegendentry{$E/\Erup$}
    \addplot[azulacero, line width=1.2pt, densely dashed, domain=1:4.2, samples=60] {1/x};
    \addlegendentry{$\phi/\phi_{\max}$}
    \addplot[azulprof, line width=1.2pt, forget plot] coordinates {(0,0) (1,0)};
    \addplot[azulacero, line width=1.2pt, densely dashed, forget plot] coordinates {(0,1) (1,1)};
    \draw[cargapos, line width=0.7pt, dashdotted] (axis cs:0,1) -- (axis cs:4.2,1);
    \node[font=\scriptsize, text=cargapos, anchor=south east] at (axis cs:4.15,1.01) {umbral de ruptura};
  \end{axis}
  \node[font=\small, anchor=north] at ($(a.outer south)+(0,-0.05)$) {(a)};
  \begin{axis}[tema, name=b, at={($(a.east)+(1.6cm,0)$)}, anchor=west, width=7.0cm, height=5.0cm,
      xmode=log, ymode=log, xmin=1e-4, xmax=1, ymin=1e2, ymax=1e7,
      xlabel={$R$ (m)}, ylabel={$\phi_{\max}$ (V)}, log basis x=10, log basis y=10]
    \addplot[azulprof, line width=1.2pt, domain=1e-4:1, samples=20] {3e6*x};
    \fill[cargapos] (axis cs:0.05,1.5e5) circle (2pt);
    \node[font=\scriptsize, anchor=north west, text=cargapos] at (axis cs:0.055,1.5e5) {$150\,\mathrm{kV}$};
    \node[font=\scriptsize, anchor=north west, text=grispizarra] at (axis cs:1.3e-4,5e6) {$\phi_{\max}=\Erup R$};
  \end{axis}
  \node[font=\small, anchor=north] at ($(b.outer south)+(0,-0.05)$) {(b)};
\end{tikzpicture}
\caption{Problema 3.1. (a) Campo y potencial de una esfera conductora cargada justo en el umbral de ruptura: el campo es máximo en la superficie y decae como
$1/r^{2}$, de modo que la ruptura empieza allí. (b) Potencial máximo frente al radio: cuanto más pequeño es el conductor, antes se descarga (efecto punta). El
punto marca el globo del enunciado.}
\label{fig:sem-globo}
\end{figure}


\begin{observacion}[Interpretación y órdenes de magnitud]\label{obs:sem-globo}
La carga máxima es $Q_{\max}=4\pi\eo R^{2}\Erup=\frac{(5\cdot10^{-2})^{2}\cdot3\cdot10^{6}}{8{,}988\cdot10^{9}}\,\mathrm{C}\simeq0{,}83\,\mu\mathrm{C}$, con densidad
$\sigma_{\max}=\eo\Erup\simeq2{,}7\cdot10^{-5}\,\mathrm{C/m^{2}}$ y presión electrostática $\Pel=\eo\Erup^{2}/2\simeq40\,\mathrm{Pa}$. La relación
$\phi_{\max}=\Erup R$ explica por qué los conductores pequeños o con puntas se descargan antes (\figpanel{fig:sem-globo}{b}): a igual potencial, cuanto menor es
$R$, mayor es el campo en la superficie. Una esfera de $1\,\mathrm{mm}$ no soporta más de $3\,\mathrm{kV}$; por eso los generadores de Van de Graaff tienen
cúpulas grandes y lisas, y los pararrayos, puntas finas (\cref{obs:pararrayos}).
\end{observacion}

% =========================================================================
\subsection{Problema 3.2. Cuatro esferas concéntricas con una batería}\label{sem:cuatro}
% =========================================================================
\begin{problema}[Cuatro superficies esféricas concéntricas]\label{prob:cuatro}
Sean cuatro superficies conductoras esféricas concéntricas $S_1$, $S_2$, $S_3$ y $S_4$ de radios $R_1<R_2<R_3<R_4$. Suponiendo que son inicialmente neutras y
aplicando una diferencia de potencial $V_0$ entre $S_1$ y $S_4$ ($V_0=\phi_1-\phi_4$), calcula los potenciales de las esferas $S_2$ y $S_3$ así como la carga
que adquiere $S_1$.
\end{problema}

\begin{resolucion}
\paso{Idea.} La batería, conectada entre $S_1$ y $S_4$ (a través de pequeños orificios en $S_2$ y $S_3$, cuyo efecto se desprecia), traslada carga de una a otra:
como estaban neutras, $Q_4=-Q_1$. Las esferas $S_2$ y $S_3$ siguen aisladas y neutras. Lo resolvemos de dos formas: con coeficientes de potencial (M3, camino B)
y con la cadena de Gauss (M2).

\paso{Método 1: coeficientes de potencial.} Por la regla \eqref{eq:regla-concentricas}, $p_{ij}=\kc/\max(R_i,R_j)$. Si ponemos una carga $Q$ en la esfera $S_j$,
para la esfera $S_i$ hay dos posibilidades: si está dentro ($R_i<R_j$), está en el interior de la capa cargada y su potencial es $\kc Q/R_j$; si está fuera
($R_i>R_j$), ve la carga como puntual y su potencial es $\kc Q/R_i$. Con $Q_2=Q_3=0$ y $Q_4=-Q_1$:
\[
  \phi_4=p_{41}Q_1+p_{44}Q_4=\frac{\kc}{R_4}\bigl(Q_1+Q_4\bigr)=0 ,
\]
así que la esfera exterior queda al mismo potencial que el infinito (la carga total es nula). La condición de la batería da
\[
  V_0=\phi_1-\phi_4=\bigl(p_{11}-p_{14}\bigr)Q_1=\frac{Q_1}{4\pi\eo}\Bigl(\frac{1}{R_1}-\frac{1}{R_4}\Bigr)=\frac{Q_1}{4\pi\eo}\,\frac{R_4-R_1}{R_1R_4}
  \quad\Longrightarrow\quad\resalta{Q_1=4\pi\eo V_0\,\frac{R_1R_4}{R_4-R_1}.}
\]
Para las esferas intermedias,
\[
  \phi_2=p_{21}Q_1+p_{24}Q_4=\bigl(p_{21}-p_{24}\bigr)Q_1=\frac{Q_1}{4\pi\eo}\Bigl(\frac{1}{R_2}-\frac{1}{R_4}\Bigr)=V_0\,\frac{R_1}{R_2}\,\frac{R_4-R_2}{R_4-R_1},
\]
\[
  \phi_3=\bigl(p_{31}-p_{34}\bigr)Q_1=\frac{Q_1}{4\pi\eo}\Bigl(\frac{1}{R_3}-\frac{1}{R_4}\Bigr)=V_0\,\frac{R_1}{R_3}\,\frac{R_4-R_3}{R_4-R_1}.
\]

\paso{Método 2: cadena de Gauss.} Las esferas neutras $S_2$ y $S_3$ llevan $\mp Q_1$ en sus caras interior y exterior, ambas en el mismo radio, y no modifican el
campo. Por Gauss con esferas concéntricas,
\[
  \vE(\vc{r})=\begin{cases}\vc{0}, & r>R_4\ (\Qenc=Q_1+Q_4=0),\\[3pt] \dfrac{Q_1}{4\pi\eo r^{2}}\,\ur, & R_1<r<R_4,\\[8pt] \vc{0}, & r<R_1 ,\end{cases}
\]
y, como $\phi(r)-\phi(R_4)=\int_r^{R_4}E\,\dif r'$ con $\phi(R_4)=0$,
\[
  \phi(r)=\begin{cases}0, & r\ge R_4,\\[3pt] \dfrac{Q_1}{4\pi\eo}\Bigl(\dfrac1r-\dfrac{1}{R_4}\Bigr), & R_1\le r\le R_4,\\[8pt]
  \dfrac{Q_1}{4\pi\eo}\Bigl(\dfrac{1}{R_1}-\dfrac{1}{R_4}\Bigr)=V_0, & r\le R_1 .\end{cases}
\]
Evaluando en $R_2$ y $R_3$ se recuperan $\phi_2$ y $\phi_3$, y la última línea reproduce $Q_1$.
\end{resolucion}

% =========================================================================
%  Figura (Problema 3.2): cuatro esferas concéntricas con una batería entre
%  S1 y S4. (b) phi(r)/V0 con R_i = i: phi/V0 = (4/r - 1)/3 en [1, 4].
% =========================================================================
\begin{figure}[t]
\centering
\begin{tikzpicture}[font=\small]
  % ---------------------------------------------------------------- (a)
  \begin{scope}[x=1cm, y=1cm]
    \foreach \r/\n in {1.75/4,1.3/3,0.85/2}{\draw[azulprof, line width=0.7pt] (0,0) circle (\r);}
    \shade[ball color=azulvelo] (0,0) circle (0.4);
    \draw[azulprof, line width=0.6pt] (0,0) circle (0.4);
    \foreach \r/\n in {0.4/1,0.85/2,1.3/3,1.75/4}{
      \node[font=\scriptsize, inner sep=0.5pt] at (125:{\r+0.22}) {$S_{\n}$};}
    % hilo por los orificios hasta la batería
    \draw[azulprof, line width=0.6pt] (0.4,0) -- (2.35,0) -- (2.35,-0.35);
    \draw[azulprof, line width=1.1pt] (2.15,-0.35) -- (2.55,-0.35);
    \draw[azulprof, line width=0.6pt] (2.25,-0.47) -- (2.45,-0.47);
    \draw[azulprof, line width=0.6pt] (2.35,-0.47) -- (2.35,-1.0) -- (1.42,-1.0);
    \node[font=\footnotesize, anchor=west] at (2.6,-0.42) {$V_0$};
    \node[font=\scriptsize, anchor=west, text=cargapos] at (2.6,0.0) {$+$};
    \foreach \a in {60,180,240,300}{\node[signopos, font=\tiny\bfseries] at (\a:0.3) {$+$};}
    \foreach \a in {70,200,250}{\node[signoneg, font=\tiny\bfseries] at (\a:1.65) {$-$};}
    \node[font=\footnotesize, text=cargapos] at (0,0.0) {$Q_1$};
    \node[font=\footnotesize, text=carganeg] at (-1.62,-1.62) {$-Q_1$};
    \node at (0,-2.5) {(a)};
  \end{scope}
  % ---------------------------------------------------------------- (b)
  \begin{scope}[xshift=4.4cm, yshift=-2.1cm]
    \begin{axis}[tema, width=8.4cm, height=4.9cm, xmin=0, xmax=5.2, ymin=0, ymax=1.15,
        xlabel={$r$ (con $R_i=i$)}, ylabel={$\phi/V_0$}, xtick={0,1,2,3,4,5},
        xticklabels={$0$,$R_1$,$R_2$,$R_3$,$R_4$,}, ytick={0,0.111,0.333,1}, yticklabels={$0$,$\frac19$,$\frac13$,$1$}]
      \addplot[azulprof, line width=1.2pt] coordinates {(0,1) (1,1)};
      \addplot[azulprof, line width=1.2pt, domain=1:4, samples=60] {(4/x-1)/3};
      \addplot[azulprof, line width=1.2pt] coordinates {(4,0) (5.2,0)};
      \foreach \x in {1,2,3,4}{\edef\tmp{\noexpand\draw[punteado] (axis cs:\x,0) -- (axis cs:\x,1.15);}\tmp}
      \fill[cargapos] (axis cs:2,0.3333) circle (1.8pt) node[above right=-1pt, font=\scriptsize] {$\phi_2$};
      \fill[cargapos] (axis cs:3,0.1111) circle (1.8pt) node[above right=-1pt, font=\scriptsize] {$\phi_3$};
      \fill[cargapos] (axis cs:1,1) circle (1.8pt) node[above right=-1pt, font=\scriptsize] {$\phi_1=V_0$};
      \fill[cargapos] (axis cs:4,0) circle (1.8pt) node[above right=-1pt, font=\scriptsize] {$\phi_4=0$};
    \end{axis}
    \node at (3.4,-1.35) {(b)};
  \end{scope}
\end{tikzpicture}
\caption{Problema 3.2. (a) Cuatro esferas concéntricas; la batería conecta $S_1$ y $S_4$ a través de pequeños orificios en $S_2$ y $S_3$. (b) Potencial con
$R_i=i$: entre $R_1$ y $R_4$ es el de la carga $Q_1$ y fuera de $S_4$ es nulo; las esferas neutras $S_2$ y $S_3$ solo toman el valor del potencial en su radio.}
\label{fig:sem-cuatro}
\end{figure}


\begin{observacion}[Interpretación]\label{obs:sem-cuatro}
Las esferas $S_1$ y $S_4$ están en influencia total: forman un condensador esférico de capacidad $C=Q_1/V_0=4\pi\eo R_1R_4/(R_4-R_1)$, que no depende de
la presencia de $S_2$ y $S_3$. Es un error frecuente pensar que estas modifican el campo: al ser neutras y delgadas, solo adoptan el potencial del radio en el que
están, $V_0>\phi_2>\phi_3>0$ (\figpanel{fig:sem-cuatro}{b}).
\end{observacion}

% =========================================================================
\subsection{Problema 3.3. Condensador esférico: coeficientes}\label{sem:esferico}
% =========================================================================
\begin{problema}[Coeficientes de un condensador esférico]\label{prob:esferico}
Sea un condensador esférico formado por una esfera conductora de radio $a$ rodeada concéntricamente por otra esfera conductora de radio $b$. Calcula los
coeficientes de capacidad e influencia y la capacidad del condensador.
\end{problema}

\begin{resolucion}
\paso{Idea.} Hay dos caminos (método M3). Con conductores concéntricos y cargas fijas, el camino B (coeficientes de potencial y luego inversión) es muy
rápido; el camino A (potenciales impuestos) sirve de comprobación.

\paso{Camino B: coeficientes de potencial.} \emph{Estado $Q_1=1$, $Q_2=0$:} la esfera exterior queda neutra, con $-1$ y $+1$ en sus caras; fuera de la esfera
interior el campo es el de una carga unidad en el centro, luego $\phi_1=\kc/a$ y $\phi_2=\kc/b$: $p_{11}=\kc/a$, $p_{21}=\kc/b$. \emph{Estado $Q_1=0$, $Q_2=1$:} el
campo dentro de la esfera exterior es nulo, y $\phi_1=\phi_2=\kc/b$: $p_{12}=p_{22}=\kc/b$. Así,
\[
  \Pm=\kc\begin{pmatrix}1/a & 1/b\\ 1/b & 1/b\end{pmatrix},\qquad\det\Pm=\kc^{2}\Bigl(\frac{1}{ab}-\frac{1}{b^{2}}\Bigr)=\kc^{2}\,\frac{b-a}{ab^{2}},
\]
\[
  \Cm=\Pm^{-1}=\frac{1}{\det\Pm}\begin{pmatrix}p_{22} & -p_{12}\\ -p_{12} & p_{11}\end{pmatrix}
  =\frac{ab^{2}}{\kc(b-a)}\begin{pmatrix}1/b & -1/b\\ -1/b & 1/a\end{pmatrix}.
\]

\paso{Camino A: potenciales impuestos.} \emph{Estado $\phi_1=1$, $\phi_2=0$:} entre las esferas el campo es el de $Q_1$ en el centro, y
$\phi_1-\phi_2=\kc Q_1\bigl(\frac1a-\frac1b\bigr)=1$, de donde $c_{11}=Q_1=4\pi\eo\frac{ab}{b-a}$. La exterior está a tierra, así que fuera el potencial es nulo y la
carga total también: $c_{21}=Q_2=-Q_1=-c_{11}$. \emph{Estado $\phi_1=0$, $\phi_2=1$:} ahora $\phi_1-\phi_2=\kc Q_1\bigl(\frac1a-\frac1b\bigr)=-1$ da
$c_{12}=Q_1=-4\pi\eo\frac{ab}{b-a}$; y $\phi_2=\kc(Q_1+Q_2)/b=1$ da $Q_1+Q_2=4\pi\eo b$, es decir, $c_{22}=4\pi\eo b-c_{12}$ (\cref{fig:sem-esferico}).
Ambos caminos coinciden:
\[
  \resalta{c_{11}=4\pi\eo\,\frac{ab}{b-a},\qquad c_{12}=c_{21}=-4\pi\eo\,\frac{ab}{b-a},\qquad c_{22}=4\pi\eo\,\frac{b^{2}}{b-a},\qquad C=4\pi\eo\,\frac{ab}{b-a}.}
\]
La capacidad sale de \cref{eq:cap-coef}, $C=1/(p_{11}+p_{22}-2p_{12})=1/[\kc(1/a-1/b)]$, o del \cref{cor:encerrado}, $C=c_{11}$.
\end{resolucion}

% =========================================================================
%  Figura (Problema 3.3): los dos estados del camino A con a = 1, b = 2.
%  phi1 = 1, phi2 = 0: phi = 1 (r <= 1), 2/r - 1 (1..2), 0 (r >= 2)
%  phi1 = 0, phi2 = 1: phi = 0 (r <= 1), 2 - 2/r (1..2), 2/r (r >= 2)
% =========================================================================
\begin{figure}[t]
\centering
\begin{tikzpicture}[font=\small]
  \begin{axis}[tema, name=a, width=7.0cm, height=4.8cm, xmin=0, xmax=4.2, ymin=0, ymax=1.18,
      xlabel={$r/a$}, ylabel={$\phi$}, xtick={0,1,2,3,4}, ytick={0,0.5,1},
      title={\footnotesize $\phi_1=1,\ \phi_2=0$: columna 1 de $\Cm$}]
    \fill[azulvelo!60] (axis cs:0,0) rectangle (axis cs:1,1.18);
    \addplot[azulprof, line width=1.2pt] coordinates {(0,1) (1,1)};
    \addplot[azulprof, line width=1.2pt, domain=1:2, samples=40] {2/x-1};
    \addplot[azulprof, line width=1.2pt] coordinates {(2,0) (4.2,0)};
    \draw[punteado] (axis cs:2,0) -- (axis cs:2,1.18);
    \node[font=\scriptsize, anchor=west, align=left] at (axis cs:2.2,0.65) {$Q_1=c_{11}>0$\\ $Q_2=c_{21}=-c_{11}$};
  \end{axis}
  \node[anchor=north] at ($(a.outer south)+(0,-0.05)$) {(a)};
  \begin{axis}[tema, name=b, at={($(a.east)+(1.4cm,0)$)}, anchor=west, width=7.0cm, height=4.8cm,
      xmin=0, xmax=4.2, ymin=0, ymax=1.18,
      xlabel={$r/a$}, ylabel={$\phi$}, xtick={0,1,2,3,4}, ytick={0,0.5,1},
      title={\footnotesize $\phi_1=0,\ \phi_2=1$: columna 2 de $\Cm$}]
    \fill[azulvelo!60] (axis cs:0,0) rectangle (axis cs:1,1.18);
    \addplot[azulacero, line width=1.2pt] coordinates {(0,0) (1,0)};
    \addplot[azulacero, line width=1.2pt, domain=1:2, samples=40] {2-2/x};
    \addplot[azulacero, line width=1.2pt, domain=2:4.2, samples=40] {2/x};
    \draw[punteado] (axis cs:2,0) -- (axis cs:2,1.18);
    \node[font=\scriptsize, anchor=west, align=left] at (axis cs:2.45,0.38) {$Q_1=c_{12}<0$\\ $Q_2=c_{22}>0$};
  \end{axis}
  \node[anchor=north] at ($(b.outer south)+(0,-0.05)$) {(b)};
\end{tikzpicture}
\caption{Problema 3.3, camino A ($b=2a$). En cada estado se fija el potencial de un conductor a la unidad y el otro a tierra; las cargas que aparecen son una
columna de la matriz de coeficientes de capacidad e influencia. (a) La esfera interior a potencial unidad: todo su flujo acaba en la exterior. (b) La esfera
exterior a potencial unidad: parte de su flujo va al infinito y parte induce carga negativa en la interior, a tierra.}
\label{fig:sem-esferico}
\end{figure}


\begin{observacion}[Comprobaciones obligatorias]\label{obs:sem-esferico}
(i) \emph{Simetría:} $c_{12}=c_{21}$. (ii) \emph{Signos:} $c_{11},c_{22}>0$ y $c_{12}<0$. (iii) \emph{Influencia:} $c_{11}=-c_{21}$, porque la esfera interior está
en influencia total; y $c_{22}+c_{12}=4\pi\eo b>0$, porque parte de las líneas de la exterior van al infinito. (iv) \emph{Límites:} si $b=a+d$ con $d\ll a$,
$C\simeq\eo\,4\pi a^{2}/d$ (condensador plano); si $b\to\infty$, $C\to4\pi\eo a$, la capacidad de la esfera aislada. (v) \emph{Unidades:} $\eo$ por una
longitud, es decir, faradios.
\end{observacion}

% =========================================================================
\subsection{Problema 3.4. Condensador cilíndrico}\label{sem:cilindrico}
% =========================================================================
\begin{problema}[Capacidad por unidad de longitud]\label{prob:cilindrico}
Un condensador cilíndrico se compone de dos cilindros conductores coaxiales de radio interno $a$ y externo $b$. Calcula la capacidad por unidad de longitud
del condensador.
\end{problema}

\begin{resolucion}
\paso{Idea.} Trabajamos con cilindros de longitud $\ell\gg b$ y despreciamos los efectos de borde, de modo que hay simetría cilíndrica. Cargamos las
armaduras con $\pm Q$, es decir, $\pm\lambda=\pm Q/\ell$ por unidad de longitud, y seguimos los cinco pasos del método directo (M2).

\paso{Paso 1: campo entre armaduras.} Por la simetría, $\vE=E(r)\,\urho$. Con una superficie gaussiana cilíndrica coaxial de radio $a<r<b$ y longitud $\ell$, las
tapas no aportan flujo y solo se encierra la carga de la armadura interior:
\[
  2\pi r\ell\,E=\frac{\lambda\ell}{\eo}\quad\Longrightarrow\quad\vE=\frac{\lambda}{2\pi\eo r}\,\urho\qquad(a<r<b),
\]
y el campo es nulo dentro del cilindro interior y fuera del exterior (la carga encerrada es $\lambda-\lambda=0$).

\paso{Paso 2: diferencia de potencial.} Integrando de la armadura positiva a la negativa,
\[
  \DV=\phi_a-\phi_b=-\int_{b}^{a}\vE\cdot\dl=\int_{a}^{b}\frac{\lambda}{2\pi\eo r}\,\dif r=\frac{\lambda}{2\pi\eo}\,\ln\frac{b}{a}.
\]
\paso{Paso 3: capacidad.}
\[
  C=\frac{Q}{\DV}=\frac{\lambda\ell}{\frac{\lambda}{2\pi\eo}\ln(b/a)}\quad\Longrightarrow\quad\resalta{\frac{C}{\ell}=\frac{2\pi\eo}{\ln(b/a)} .}
\]

\paso{Con coeficientes de potencial.} Para un cilindro indefinido no se puede tomar $\phi(\infty)=0$: el potencial de un hilo cargado diverge
logarítmicamente. Tomamos el origen de potenciales a una distancia arbitraria $r_0$ del eje. Un cilindro de radio $R$ con carga $Q$ en una longitud $\ell$
crea entonces (\figpanel{fig:sem-cilindrico}{b})
\[
  \phi(r)=\begin{cases}\dfrac{Q}{2\pi\eo\ell}\,\ln\dfrac{r_0}{r}, & r\ge R,\\[8pt] \dfrac{Q}{2\pi\eo\ell}\,\ln\dfrac{r_0}{R}, & r\le R,\end{cases}
\]
y, razonando como en el caso esférico,
\[
  \Pm=\frac{1}{2\pi\eo\ell}\begin{pmatrix}\ln(r_0/a) & \ln(r_0/b)\\ \ln(r_0/b) & \ln(r_0/b)\end{pmatrix}.
\]
Como $p_{22}=p_{12}$, $p_{11}+p_{22}-2p_{12}=p_{11}-p_{12}$ y
\[
  C=\bigl(p_{11}-p_{12}\bigr)^{-1}=2\pi\eo\ell\Bigl[\ln\frac{r_0}{a}-\ln\frac{r_0}{b}\Bigr]^{-1}=\frac{2\pi\eo\ell}{\ln(b/a)} ,
\]
que no depende de $r_0$, como debe ser. Elegir $r_0=b$ simplifica mucho el problema, porque entonces $\ln(r_0/b)=\ln1=0$.
\end{resolucion}

% =========================================================================
%  Figura (Problema 3.4): condensador cilíndrico.
%  (a) Sección transversal: +lambda en r = a, -lambda en r = b, gaussiana.
%  (b) Potencial de un cilindro cargado respecto de una referencia r0:
%      phi ∝ ln(r0/R) (r <= R), ln(r0/r) (r >= R), con R = 1 y r0 = 3.
% =========================================================================
\begin{figure}[t]
\centering
\begin{tikzpicture}[font=\small]
  % ---------------------------------------------------------------- (a)
  \begin{scope}[x=0.95cm, y=0.95cm]
    \foreach \a in {0,30,...,330}{
      \draw[azulacero, line width=0.55pt, -{Stealth[length=3.5pt]}] (\a:0.78) -- (\a:1.35);
      \draw[azulacero, line width=0.55pt] (\a:1.35) -- (\a:1.85);}
    \draw[azulprof, line width=1.1pt] (0,0) circle (1.95);
    \shade[ball color=azulvelo] (0,0) circle (0.7);
    \draw[azulprof, line width=0.6pt] (0,0) circle (0.7);
    \draw[gauss] (0,0) circle (1.3);
    \foreach \a in {0,45,...,315}{\node[signopos, font=\tiny\bfseries] at (\a:0.55) {$+$};
                                 \node[signoneg, font=\tiny\bfseries] at (\a+22.5:1.82) {$-$};}
    \node[font=\scriptsize, text=cargapos] at (0,0) {$+\lambda$};
    \node[font=\scriptsize, text=carganeg] at (1.9,1.9) {$-\lambda$};
    \draw[-{Stealth[length=3.5pt]}, azulprof, line width=0.5pt] (0,0) -- (-125:0.7);
    \node[font=\scriptsize] at (-0.55,-0.25) {$a$};
    \draw[-{Stealth[length=3.5pt]}, azulprof, line width=0.5pt] (0,0) -- (-55:1.95);
    \node[font=\scriptsize] at (1.0,-1.25) {$b$};
    \node[font=\scriptsize, text=cargapos!85] at (-1.05,1.15) {$r$};
    \node at (0,-2.6) {(a)};
  \end{scope}
  % ---------------------------------------------------------------- (b)
  \begin{scope}[xshift=3.95cm, yshift=-2.15cm]
    \begin{axis}[tema, width=8.0cm, height=4.9cm, xmin=0, xmax=6.2, ymin=-0.75, ymax=1.35,
        axis x line=middle, axis y line=left,
        xlabel={$r$}, ylabel={$\phi\,/\,\frac{Q}{2\pi\eo\ell}$}, xtick={1,3}, xticklabels={$R$,$r_0$}, ytick={-0.5,0,0.5,1},
        xlabel style={anchor=west}]
      \addplot[azulprof, line width=1.2pt] coordinates {(0,1.0986) (1,1.0986)};
      \addplot[azulprof, line width=1.2pt, domain=1:6.2, samples=60] {ln(3/x)};
      \draw[punteado] (axis cs:1,-0.75) -- (axis cs:1,1.35);
      \fill[cargapos] (axis cs:3,0) circle (1.8pt);
      \node[font=\scriptsize, anchor=south west, text=grispizarra] at (axis cs:3.05,0.05) {referencia $\phi(r_0)=0$};
      \node[font=\scriptsize, anchor=west, text=grispizarra] at (axis cs:1.15,1.12) {$\ln(r_0/R)$};
    \end{axis}
    \node at (3.2,-1.35) {(b)};
  \end{scope}
\end{tikzpicture}
\caption{Problema 3.4. (a) Sección transversal del condensador cilíndrico (longitud $\ell$): la gaussiana coaxial de radio $a<r<b$ solo encierra la carga de la
armadura interior. (b) Potencial de un cilindro cargado de radio $R$: como diverge en el infinito, se refiere a una distancia arbitraria $r_0$; las diferencias
de potencial, y por tanto la capacidad, no dependen de esa elección.}
\label{fig:sem-cilindrico}
\end{figure}


\begin{observacion}[Interpretación]\label{obs:sem-cilindrico}
Con placas próximas, $b=a+d$ y $d\ll a$, $\ln(1+d/a)\simeq d/a$ y $C/\ell\simeq2\pi\eo a/d=\eo(A/\ell)/d$, con $A/\ell=2\pi a$: de nuevo el condensador plano. Los
coeficientes de potencial respecto de $r_0$ ya no cumplen todas las propiedades del \cref{teo:pij-prop} (si $r_0<b$, $p_{22}<0$), porque estas usaban
$\phi(\infty)=0$; las diferencias de potencial y la capacidad, en cambio, son independientes de la referencia.
\end{observacion}

% =========================================================================
\subsection{Problema 3.5. Placas planas paralelas}\label{sem:plano}
% =========================================================================
\begin{problema}[Condensador plano y reparto de carga por caras]\label{prob:plano}
Sean dos conductores planos paralelos de superficie $A$ y separados una distancia $d$. Despreciando el efecto de bordes, calcula la capacidad del
condensador. Si la carga de cada conductor es $Q_1$ y $Q_2$, calcula las cargas en cada una de las caras de los conductores.
\end{problema}

\begin{resolucion}
\paso{Idea.} Si las dimensiones de las placas son grandes frente a $d$, cada cara se comporta como un plano indefinido cargado, que crea un campo
$\sigma/2\eo$ a cada lado, dirigido hacia fuera del plano si $\sigma>0$ (método M4). Con dos placas hay cuatro caras, con cargas $q_1,q_2,q_3,q_4$ de
izquierda a derecha (\cref{fig:sem-plano}).

\paso{Paso 1: condensador ($Q_1=Q$, $Q_2=-Q$).} Las cargas $\pm Q$ se sitúan en las caras enfrentadas, con densidades $\pm\sigma=\pm Q/A$. El campo entre las
placas puede obtenerse de dos formas. \emph{Por superposición:} entre las placas los campos de ambas se suman, $\sigma/2\eo+\sigma/2\eo=\sigma/\eo$, y fuera se
anulan. \emph{Por Gauss}, con una caja que tiene una tapa dentro del metal ($\vE=\vc{0}$) y otra entre las placas:
\[
  E\,A_{\mathrm{caja}}=\frac{\sigma A_{\mathrm{caja}}}{\eo}\quad\Longrightarrow\quad E=\frac{\sigma}{\eo}=\frac{Q}{\eo A}.
\]
El campo es uniforme, de modo que $\DV=E\,d=\frac{Qd}{\eo A}$ y
\[
  \resalta{C=\frac{Q}{\DV}=\eo\,\frac{A}{d} .}
\]

\paso{Paso 2: cargas cualesquiera $Q_1$ y $Q_2$.} Tomamos como positivo el sentido hacia la derecha. Un punto interior a la placa~1 está a la derecha de la cara
$q_1$ y a la izquierda de $q_2$, $q_3$ y $q_4$; uno interior a la placa~2, a la derecha de $q_1,q_2,q_3$ y a la izquierda de $q_4$. Imponiendo $\vE=\vc{0}$ en
ambos y la conservación de la carga en cada placa:
\[
  \begin{aligned}
    q_1-q_2-q_3-q_4&=0 &&\text{(campo nulo en la placa 1)},\\
    q_1+q_2+q_3-q_4&=0 &&\text{(campo nulo en la placa 2)},\\
    q_1+q_2&=Q_1,\qquad q_3+q_4=Q_2 &&\text{(conservación de la carga)}.
  \end{aligned}
\]
Restando las dos primeras, $2(q_2+q_3)=0$, es decir, $q_3=-q_2$; sumándolas, $2(q_1-q_4)=0$, es decir, $q_1=q_4$. Con la conservación,
$q_1+q_2=Q_1$ y $q_1-q_2=Q_2$, de donde
\[
  \resalta{q_1=q_4=\frac{Q_1+Q_2}{2}\ \ \text{(caras exteriores)},\qquad q_2=-q_3=\frac{Q_1-Q_2}{2}\ \ \text{(caras interiores)}.}
\]
\end{resolucion}

% =========================================================================
%  Figura (Problema 3.5): placas planas paralelas.
%  (a) Condensador ideal: +sigma y -sigma en las caras enfrentadas, campo
%      sigma/eps0 entre las placas y nulo fuera; caja de Gauss.
%  (b) Cargas cualesquiera: cuatro caras q1..q4 y campo en cada región.
% =========================================================================
\begin{figure}[t]
\centering
\begin{tikzpicture}[font=\small]
  % ---------------------------------------------------------------- (a)
  \begin{scope}
    \fill[top color=azulvelo!80!white, bottom color=metal, draw=azulprof, line width=0.6pt] (-1.75,-1.7) rectangle (-1.4,1.7);
    \fill[top color=azulvelo!80!white, bottom color=metal, draw=azulprof, line width=0.6pt] (1.4,-1.7) rectangle (1.75,1.7);
    \foreach \y in {-1.5,-1.1,...,1.6}{\node[signopos, font=\tiny\bfseries] at (-1.28,\y) {$+$};
                                       \node[signoneg, font=\tiny\bfseries] at (1.28,\y) {$-$};}
    \foreach \y in {-1.35,-0.85,0.85,1.35}{\draw[lcampo] (-1.1,\y) -- (1.1,\y);}
    \node[font=\footnotesize, fill=white, inner sep=1.5pt] at (0.25,-0.2) {$E=\dfrac{\sigma}{\eo}$};
    \draw[gaussazul] (-1.62,0.15) rectangle (-0.75,0.55);
    \node[font=\scriptsize, text=azulacero, anchor=west] at (-0.7,0.35) {gaussiana};
    \node[font=\footnotesize] at (-2.35,-0.9) {$E=0$};
    \node[font=\footnotesize] at (2.35,-0.9) {$E=0$};
    \node[font=\footnotesize, text=cargapos] at (-1.58,-2.0) {$+Q$};
    \node[font=\footnotesize, text=carganeg] at (1.58,-2.0) {$-Q$};
    \draw[{Stealth[length=3pt]}-{Stealth[length=3pt]}, grispizarra] (-1.4,1.95) -- (1.4,1.95) node[midway, above, font=\footnotesize] {$d$};
    \node at (0,-3.15) {(a)};
  \end{scope}
  % ---------------------------------------------------------------- (b)
  \begin{scope}[xshift=8.6cm]
    \fill[top color=azulvelo!80!white, bottom color=metal, draw=azulprof, line width=0.6pt] (-2.2,-1.7) rectangle (-1.6,1.7);
    \fill[top color=azulvelo!80!white, bottom color=metal, draw=azulprof, line width=0.6pt] (1.6,-1.7) rectangle (2.2,1.7);
    \foreach \x/\q in {-2.2/q_1,-1.6/q_2,1.6/q_3,2.2/q_4}{
      \draw[cargapos, line width=1.3pt] (\x,-1.7) -- (\x,1.7);
      \node[font=\footnotesize, anchor=south] at (\x,1.72) {$\q$};}
    \node[font=\scriptsize, rotate=90] at (-1.9,0) {$\vE=\vc{0}$};
    \node[font=\scriptsize, rotate=90] at (1.9,0) {$\vE=\vc{0}$};
    \node[font=\footnotesize] at (-1.9,-2.0) {$Q_1$};
    \node[font=\footnotesize] at (1.9,-2.0) {$Q_2$};
    \draw[-{Stealth[length=4pt]}, azulacero, line width=0.8pt] (-2.45,1.0) -- (-3.25,1.0);
    \node[font=\scriptsize, anchor=north] at (-2.85,0.85) {$\dfrac{Q_1+Q_2}{2\eo A}$};
    \draw[-{Stealth[length=4pt]}, azulacero, line width=0.8pt] (-0.4,1.0) -- (0.4,1.0);
    \node[font=\scriptsize, anchor=north] at (0,0.85) {$\dfrac{Q_1-Q_2}{2\eo A}$};
    \draw[-{Stealth[length=4pt]}, azulacero, line width=0.8pt] (2.45,1.0) -- (3.25,1.0);
    \node[font=\scriptsize, anchor=north] at (2.85,0.85) {$\dfrac{Q_1+Q_2}{2\eo A}$};
    \node[font=\scriptsize, align=center] at (0,-2.45) {$q_1=q_4=\tfrac12(Q_1+Q_2)$\\[3pt] $q_2=-q_3=\tfrac12(Q_1-Q_2)$};
    \node at (0,-3.15) {(b)};
  \end{scope}
\end{tikzpicture}
\caption{Problema 3.5. (a) Condensador plano ideal: la carga está en las caras enfrentadas y el campo, uniforme, solo existe entre las placas; la caja de Gauss
con una tapa dentro del metal da $E=\sigma/\eo$. (b) Placas con cargas cualesquiera: las caras exteriores se reparten la carga total y las interiores llevan
cargas opuestas; se indica el campo en cada región (para $Q_1>Q_2$ y $Q_1+Q_2>0$).}
\label{fig:sem-plano}
\end{figure}


\begin{observacion}[Interpretación]\label{obs:sem-plano}
Las caras exteriores comparten la carga total a partes iguales, y las enfrentadas llevan cargas opuestas: solo estas forman el condensador. El campo es
$\frac{Q_1-Q_2}{2\eo A}$ entre las placas y $\frac{Q_1+Q_2}{2\eo A}$ fuera, hacia fuera a cada lado (\cref{tab:sem-plano}). Con $Q_1=Q$ y $Q_2=-Q$ se recupera el
condensador ideal: caras exteriores descargadas y campo nulo fuera. Si, por ejemplo, $Q_2=0$, cada cara de la placa neutra lleva $\mp Q_1/2$: la placa neutra
«parte» el campo pero no lo anula.
\end{observacion}

\begin{table}[ht]
\centering
\caption{Campo de cada placa del condensador plano ideal ($\pm\sigma$) en las tres regiones (positivo hacia la derecha).}
\label{tab:sem-plano}
\begin{tabular}{@{}lccc@{}}
  \toprule
  & Izquierda & Entre placas & Derecha\\
  \midrule
  Placa $+\sigma$ & $-\dfrac{\sigma}{2\eo}$ & $+\dfrac{\sigma}{2\eo}$ & $+\dfrac{\sigma}{2\eo}$\\[8pt]
  Placa $-\sigma$ & $+\dfrac{\sigma}{2\eo}$ & $+\dfrac{\sigma}{2\eo}$ & $-\dfrac{\sigma}{2\eo}$\\[6pt]
  \midrule
  Total & $0$ & $\dfrac{\sigma}{\eo}$ & $0$\\[6pt]
  \bottomrule
\end{tabular}
\end{table}

% =========================================================================
\subsection{Problema 3.6. Tres esferas alineadas}\label{sem:alineadas}
% =========================================================================
\begin{problema}[Tres esferas alineadas, laterales a tierra]\label{prob:alineadas}
Tres esferas conductoras de radio $R$ se sitúan alineadas con sus centros separados entre sí $d$ ($d\gg R$). Las laterales se conectan a tierra y a la central
se le suministra una carga $Q$. Halla el potencial de esta última.
\end{problema}

\begin{resolucion}
\paso{Idea.} Como $d\gg R$, aplicamos el método M5: las esferas se ven entre sí como cargas puntuales y los coeficientes de potencial son los de la regla
\eqref{eq:regla-lejanas}. Numeramos 1 (izquierda), 2 (central) y 3 (derecha), con $d_{12}=d_{23}=d$ y $d_{13}=2d$.

\paso{Paso 1: matriz de coeficientes de potencial.}
\[
  \Pm=\kc\begin{pmatrix}1/R & 1/d & 1/(2d)\\ 1/d & 1/R & 1/d\\ 1/(2d) & 1/d & 1/R\end{pmatrix}.
\]
\paso{Paso 2: clasificación.} Las esferas 1 y 3 tienen como dato $\phi=0$ y como incógnita su carga; la 2 tiene como dato $Q_2=Q$ y como incógnita $\phi_2$. Por
simetría, $Q_1=Q_3\equiv q$.

\paso{Paso 3: esferas a tierra.}
\[
  0=\phi_1=\kc\Bigl(\frac{q}{R}+\frac{Q}{d}+\frac{q}{2d}\Bigr)\quad\Longrightarrow\quad q\,\frac{2d+R}{2dR}=-\frac{Q}{d}
  \quad\Longrightarrow\quad q=-\frac{2QR}{2d+R}.
\]
\paso{Paso 4: potencial buscado.}
\[
  \phi_2=\kc\Bigl(\frac{Q}{R}+\frac{2q}{d}\Bigr)=\frac{Q}{4\pi\eo}\Bigl[\frac1R-\frac{4R}{d(2d+R)}\Bigr]
  \quad\Longrightarrow\quad\resalta{\phi_2=\frac{Q}{4\pi\eo}\,\frac{2d^{2}+Rd-4R^{2}}{Rd\,(2d+R)} .}
\]
\end{resolucion}

% =========================================================================
%  Figura (Problema 3.6): tres esferas alineadas, laterales a tierra.
%  (b) phi2/(kQ/R) frente a d/R: 1 - 4/(x(2x+1)) y la aproximación 1 - 2/x^2.
% =========================================================================
\begin{figure}[t]
\centering
\begin{tikzpicture}[font=\small]
  % ---------------------------------------------------------------- (a)
  \begin{scope}
    \foreach \x/\n in {-2.2/1,0/2,2.2/3}{
      \shade[ball color=azulvelo] (\x,0) circle (0.38);
      \draw[azulprof, line width=0.6pt] (\x,0) circle (0.38);
      \node[font=\footnotesize] at (\x,0) {$\n$};}
    \foreach \x in {-2.2,2.2}{\draw[azulprof, line width=0.6pt] (\x,-0.38) -- (\x,-0.8); \pic at (\x,-0.8) {tierra};
                              \node[font=\footnotesize, text=carganeg] at (\x,0.65) {$q$};}
    \node[font=\footnotesize, text=cargapos] at (0,0.65) {$Q$};
    \draw[{Stealth[length=3pt]}-{Stealth[length=3pt]}, grispizarra] (-2.2,-1.45) -- (0,-1.45) node[midway, below, font=\footnotesize] {$d$};
    \draw[{Stealth[length=3pt]}-{Stealth[length=3pt]}, grispizarra] (0,-1.45) -- (2.2,-1.45) node[midway, below, font=\footnotesize] {$d$};
    \draw[-{Stealth[length=3pt]}, azulprof, line width=0.5pt] (0,0) ++(-130:0.02) -- ++(-130:0.36);
    \node[font=\scriptsize] at (-0.5,-0.55) {$R$};
    \node[font=\footnotesize, text=grispizarra] at (0,1.2) {$\phi_1=\phi_3=0$};
    \node at (0,-2.45) {(a)};
  \end{scope}
  % ---------------------------------------------------------------- (b)
  \begin{scope}[xshift=4.0cm, yshift=-2.0cm]
    \begin{axis}[tema, width=7.6cm, height=4.6cm, xmin=3, xmax=20, ymin=0.8, ymax=1.03,
        xlabel={$d/R$}, ylabel={$\phi_2\,/\,(\kc Q/R)$}, ytick={0.8,0.85,0.9,0.95,1},
        legend style={at={(0.97,0.05)}, anchor=south east}]
      \addplot[grismedio, line width=0.8pt, densely dashed, domain=3:20] {1};
      \addlegendentry{esfera aislada}
      \addplot[azulprof, line width=1.2pt, domain=3:20, samples=80] {1-4/(x*(2*x+1))};
      \addlegendentry{resultado}
      \addplot[azulcielo, line width=1.1pt, densely dotted, domain=3:20, samples=80] {1-2/x^2};
      \addlegendentry{$1-2R^{2}/d^{2}$}
    \end{axis}
    \node at (2.95,-1.35) {(b)};
  \end{scope}
\end{tikzpicture}
\caption{Problema 3.6. (a) Esfera central con carga $Q$ entre dos esferas a tierra. (b) Potencial de la esfera central normalizado al de una esfera aislada: las
cargas negativas inducidas en las laterales lo rebajan en una cantidad de orden $(R/d)^{2}$.}
\label{fig:sem-alineadas}
\end{figure}


\begin{observacion}[Interpretación y orden de la aproximación]\label{obs:sem-alineadas}
Las esferas a tierra adquieren carga negativa, $q\simeq-QR/d$, que rebaja el potencial de la central respecto del de una esfera aislada, $\kc Q/R$:
\[
  \phi_2=\frac{Q}{4\pi\eo R}\Bigl[1-\frac{4R^{2}}{d(2d+R)}\Bigr]\simeq\frac{Q}{4\pi\eo R}\Bigl(1-\frac{2R^{2}}{d^{2}}\Bigr).
\]
La corrección es pequeña, de orden $(R/d)^{2}$ (\figpanel{fig:sem-alineadas}{b}). El resultado solo tiene sentido al orden más bajo en $R/d$: los términos de orden
superior (como el $R^{3}/d^{3}$ que aparece al desarrollar la expresión exacta) no son significativos dentro de la aproximación de esferas lejanas, porque se
han despreciado efectos del mismo orden (la deformación de la carga de cada esfera por las vecinas).
\end{observacion}

% =========================================================================
\subsection{Problema 3.7. Tres esferas en triángulo, a tierra por etapas}\label{sem:triangulo}
% =========================================================================
\begin{problema}[Tres esferas a tierra una a una]\label{prob:triangulo}
Sean tres esferas conductoras de radio $a$ situadas en los vértices de un triángulo equilátero de lado $b$ ($b\gg a$). Las tres esferas tienen carga $q$.
Sucesivamente, una a una se llevan a potencial de tierra mediante una conexión con un hilo conductor muy delgado hasta que se alcanza el equilibrio
electrostático, desconectándolas posteriormente. Calcula las cargas $q_1$, $q_2$ y $q_3$ de las esferas al final del proceso.
\end{problema}

\begin{resolucion}
\paso{Idea.} Método M5 por etapas (\cref{fig:sem-triangulo}). En cada etapa, solo la esfera conectada a tierra tiene su carga como incógnita; las demás están
aisladas y conservan la carga que tenían, incluidas las ya desconectadas: \emph{lo que se congela es la carga, no el potencial}.

\paso{Paso 1: matriz de coeficientes de potencial.} Como $b\gg a$, cada esfera se ve desde las otras como una carga puntual. El coeficiente propio es el
potencial de una esfera aislada con carga unidad, $p_{ii}\simeq\kc/a$, y el mutuo, el que crea una carga unidad a distancia $b$, $p_{ij}\simeq\kc/b$; por la
simetría del triángulo, los tres mutuos son iguales:
\[
  \Pm=\kc\begin{pmatrix}1/a & 1/b & 1/b\\ 1/b & 1/a & 1/b\\ 1/b & 1/b & 1/a\end{pmatrix},\qquad
  \phi_i=\kc\Bigl(\frac{Q_i}{a}+\frac{Q_j+Q_l}{b}\Bigr),
\]
con $j$ y $l$ las otras dos esferas.

\paso{Paso 2: clasificación, etapa por etapa.} En la etapa 1 la esfera 1 está a tierra y $Q_2=Q_3=q$; en la etapa 2, la esfera 2 está a tierra, $Q_1=q_1$
(congelada) y $Q_3=q$; en la etapa 3, la esfera 3 está a tierra y $Q_1=q_1$, $Q_2=q_2$ (congeladas).

\paso{Paso 3: una ecuación por etapa.} Solo se usa la ecuación de la esfera a tierra:
\[
  \text{Etapa 1:}\quad0=\phi_1=\kc\Bigl(\frac{q_1}{a}+\frac{q+q}{b}\Bigr)\quad\Longrightarrow\quad q_1=-\frac{2qa}{b},
\]
\[
  \text{Etapa 2:}\quad0=\phi_2=\kc\Bigl(\frac{q_2}{a}+\frac{q_1+q}{b}\Bigr)\quad\Longrightarrow\quad
  q_2=-\frac{a}{b}\bigl(q+q_1\bigr)=-\frac{qa}{b}\Bigl(1-\frac{2a}{b}\Bigr)=-\frac{qa\,(b-2a)}{b^{2}},
\]
\[
  \text{Etapa 3:}\quad0=\phi_3=\kc\Bigl(\frac{q_3}{a}+\frac{q_1+q_2}{b}\Bigr)\quad\Longrightarrow\quad
  q_3=-\frac{a}{b}\bigl(q_1+q_2\bigr)=\frac{a}{b}\,\frac{qa}{b^{2}}\bigl[2b+(b-2a)\bigr]=\frac{qa^{2}(3b-2a)}{b^{3}} .
\]
En resumen,
\[
  \resalta{q_1=-\frac{2qa}{b},\qquad q_2=-\frac{qa\,(b-2a)}{b^{2}},\qquad q_3=\frac{qa^{2}(3b-2a)}{b^{3}} .}
\]
\end{resolucion}

% =========================================================================
%  Figura (Problema 3.7): proceso por etapas y cargas en cada etapa para
%  a/b = 0,1 (q = 1): inicial (1, 1, 1); etapa 1 (-0.2, 1, 1);
%  etapa 2 (-0.2, -0.08, 1); etapa 3 (-0.2, -0.08, 0.028).
% =========================================================================
\newcommand{\triEtapa}[5]{% #1 título, #2 #3 #4 rótulos de carga, #5 esfera a tierra (0 = ninguna)
  \coordinate (v1) at (-0.8,0.45); \coordinate (v2) at (0.8,0.45); \coordinate (v3) at (0,-0.94);
  \draw[grismedio, line width=0.4pt, densely dashed] (v1) -- (v2) -- (v3) -- cycle;
  \foreach \v/\n in {v1/1,v2/2,v3/3}{
    \ifnum\n=#5 \shade[ball color=azulcielo] (\v) circle (0.2); \else \shade[ball color=azulvelo] (\v) circle (0.2); \fi
    \draw[azulprof, line width=0.5pt] (\v) circle (0.2);
    \node[font=\tiny] at (\v) {\n};}
  \node[font=\scriptsize, anchor=south] at ($(v1)+(0,0.22)$) {#2};
  \node[font=\scriptsize, anchor=south] at ($(v2)+(0,0.22)$) {#3};
  \node[font=\scriptsize, anchor=west] at ($(v3)+(0.24,0)$) {#4};
  \ifnum#5=1 \draw[azulprof, line width=0.5pt] (v1) ++(-0.2,0) -- ++(-0.3,0) -- ++(0,-0.25); \pic at ($(v1)+(-0.5,-0.25)$) {tierra};\fi
  \ifnum#5=2 \draw[azulprof, line width=0.5pt] (v2) ++(0.2,0) -- ++(0.3,0) -- ++(0,-0.25); \pic at ($(v2)+(0.5,-0.25)$) {tierra};\fi
  \ifnum#5=3 \draw[azulprof, line width=0.5pt] (v3) ++(-0.2,0) -- ++(-0.3,0) -- ++(0,-0.25); \pic at ($(v3)+(-0.5,-0.25)$) {tierra};\fi
  \node[font=\footnotesize, anchor=north] at (0,-1.45) {#1};}

\begin{figure}[t]
\centering
\begin{tikzpicture}[font=\small]
  \begin{scope}[xshift=0cm]\triEtapa{Inicial}{$q$}{$q$}{$q$}{0}\end{scope}
  \begin{scope}[xshift=3.6cm]\triEtapa{Etapa 1: $\phi_1=0$}{$q_1$}{$q$}{$q$}{1}\end{scope}
  \begin{scope}[xshift=7.2cm]\triEtapa{Etapa 2: $\phi_2=0$}{$q_1$}{$q_2$}{$q$}{2}\end{scope}
  \begin{scope}[xshift=10.8cm]\triEtapa{Etapa 3: $\phi_3=0$}{$q_1$}{$q_2$}{$q_3$}{3}\end{scope}
  \node[font=\scriptsize, text=grispizarra] at (5.4,1.45) {lado $b$, radio $a$};
  \node at (5.4,-2.15) {(a)};
  \begin{scope}[yshift=-6.4cm, xshift=1.7cm]
    \begin{axis}[tema, ybar, bar width=7pt, width=11.8cm, height=4.6cm, ymin=-0.3, ymax=1.15,
        symbolic x coords={Inicial,Etapa 1,Etapa 2,Etapa 3}, xtick=data, ylabel={$Q_i/q$},
        enlarge x limits=0.15, axis x line*=bottom, axis y line*=left, ytick={-0.2,0,0.5,1},
        legend style={at={(0.5,1.03)}, anchor=south, legend columns=3}]
      \addplot[fill=azulprof, draw=azulprof] coordinates {(Inicial,1) (Etapa 1,-0.2) (Etapa 2,-0.2) (Etapa 3,-0.2)};
      \addlegendentry{esfera 1}
      \addplot[fill=azulacero, draw=azulacero] coordinates {(Inicial,1) (Etapa 1,1) (Etapa 2,-0.08) (Etapa 3,-0.08)};
      \addlegendentry{esfera 2}
      \addplot[fill=azulcielo, draw=azulcielo] coordinates {(Inicial,1) (Etapa 1,1) (Etapa 2,1) (Etapa 3,0.028)};
      \addlegendentry{esfera 3}
    \end{axis}
    \node at (5.2,-1.2) {(b)};
  \end{scope}
\end{tikzpicture}
\caption{Problema 3.7. (a) Esquema del proceso: en cada etapa se conecta a tierra una esfera (sombreada) y después se desconecta; las demás conservan su carga.
(b) Cargas de las tres esferas tras cada etapa para $a/b=0{,}1$: cada esfera conectada a tierra pierde casi toda su carga, y la última acaba con una pequeña
carga positiva porque sus vecinas ya son negativas.}
\label{fig:sem-triangulo}
\end{figure}


\begin{observacion}[Comprobaciones e interpretación]\label{obs:sem-triangulo}
\emph{Patrón común:} la esfera que se conecta a tierra adquiere $-a/b$ veces la suma de las cargas actuales de las otras dos, $Q_i^{\mathrm{final}}=-\frac{a}{b}(Q_j+Q_l)$.
\emph{Signos:} $q_1<0$ porque las otras dos son positivas; $q_2<0$ porque $q+q_1=q(1-2a/b)>0$ al ser $b\gg a$; y $q_3>0$ porque las otras dos son ya negativas.
\emph{Tamaños:} $\abs{q_1},\abs{q_2}\sim qa/b$ y $q_3\sim qa^{2}/b^{2}$: cada esfera conectada a tierra pierde casi toda su carga, y lo que le queda es lo que
inducen las vecinas lejanas. \emph{Lo que se congela es la carga:} al terminar la etapa 2, la esfera 1 ya no está a potencial cero,
\[
  \phi_1=\kc\Bigl(\frac{q_1}{a}+\frac{q_2+q}{b}\Bigr)=\kc\,\frac{q_2-q}{b}\neq0 :
\]
al conectar la esfera 2 su carga cambia y eso modifica el potencial de la 1, pero como la 1 está aislada su carga no puede variar. Por último, los términos como
$2a/b$ dentro de $q_2$ y $q_3$ provienen de propagar la aproximación de una etapa a la siguiente, y coinciden con la solución de la relación.
\end{observacion}

% =========================================================================
\subsection{Ejemplo modelo: tres esferas concéntricas, la interior a tierra}\label{sem:modelo}
% =========================================================================
\begin{problema}[Método de los coeficientes de potencial]\label{prob:modelo}
Tres esferas conductoras concéntricas de radios $a$, $b=2a$ y $c=3a$ tienen cargas $q_a=q$, $q_b=2q$ y $q_c=3q$. ¿Cuánto varía el potencial de la esfera
externa cuando la interna se conecta a tierra?
\end{problema}

\begin{resolucion}
\paso{Paso 1: matriz.} Con la regla \eqref{eq:regla-concentricas},
\[
  \Pm=\frac{\kc}{a}\begin{pmatrix}1 & 1/2 & 1/3\\ 1/2 & 1/2 & 1/3\\ 1/3 & 1/3 & 1/3\end{pmatrix}.
\]
\paso{Paso 2: estado inicial.} Todas aisladas, con cargas conocidas:
$\phi_c=p_{ca}q_a+p_{cb}q_b+p_{cc}q_c=\frac{\kc}{3a}(q+2q+3q)=\frac{2\kc q}{a}$.

\paso{Paso 3: clasificación tras conectar.} Para $a$, el dato es $\phi_a=0$ y la incógnita $q_a'$; las esferas $b$ y $c$ siguen aisladas, con $q_b=2q$ y
$q_c=3q$.

\paso{Paso 4: ecuación de la esfera a tierra.}
\[
  0=\phi_a=p_{aa}q_a'+p_{ab}q_b+p_{ac}q_c=\kc\Bigl(\frac{q_a'}{a}+\frac{2q}{2a}+\frac{3q}{3a}\Bigr)\quad\Longrightarrow\quad q_a'=-2q .
\]
Por el hilo pasa a tierra una carga $q_a-q_a'=3q$.

\paso{Paso 5: variación del potencial.} Como $\Pm$ solo depende de la geometría, si en un proceso cambian algunas cargas,
$\Delta\phi_i=\sum_jp_{ij}\,\Delta Q_j$, y solo hay que sumar sobre los conductores cuya carga cambia:
\[
  \Delta\phi_c=p_{ca}\bigl(q_a'-q_a\bigr)=\frac{\kc}{3a}(-2q-q)\quad\Longrightarrow\quad\resalta{\Delta\phi_c=-\frac{q}{4\pi\eo a}.}
\]
\end{resolucion}

\begin{observacion}[Comprobación con Gauss]\label{obs:sem-modelo}
Fuera de todo, la carga total pasa de $6q$ a $3q$, así que $\phi_c=\kc Q_{\mathrm{tot}}/c$ pasa de $\kc\,6q/3a=2\kc q/a$ a $\kc\,3q/3a=\kc q/a$: coincide. También se
comprueba $\phi_a'=0$ con la cadena de Gauss (\figpanel{fig:sem-modelo}{b}):
$\phi_a'=\phi_c'+\kc q_a'\bigl(\frac1a-\frac1b\bigr)+\kc(q_a'+q_b)\bigl(\frac1b-\frac1c\bigr)=\frac{\kc q}{a}-\frac{2\kc q}{a}\cdot\frac12+0=0$.
\emph{Trampa habitual:} al conectar a tierra cambia la carga del conductor conectado, no la de los aislados; pero cambian \emph{todos} los potenciales.
\end{observacion}

% =========================================================================
%  Figura (ejemplo modelo): tres esferas concéntricas a, 2a, 3a con cargas
%  q, 2q, 3q. Potencial (en unidades de kq/a) antes y después de conectar la
%  interior a tierra:
%   antes:   3 (r<=1), 2 + 1/r (1..2), 1 + 3/r (2..3), 6/r (r>=3)
%   después: 0 (r<=1), 2 - 2/r (1..2), 1 (2..3), 3/r (r>=3)
% =========================================================================
\begin{figure}[t]
\centering
\begin{tikzpicture}[font=\small]
  % ---------------------------------------------------------------- (a)
  \begin{scope}[x=0.62cm, y=0.62cm]
    \draw[azulprof, line width=0.8pt] (0,0) circle (3);
    \draw[azulprof, line width=0.8pt] (0,0) circle (2);
    \shade[ball color=azulvelo] (0,0) circle (1);
    \draw[azulprof, line width=0.6pt] (0,0) circle (1);
    \node[font=\scriptsize] at (0,0.25) {$q$};
    \node[font=\scriptsize] at (0,-0.3) {$a$};
    \node[font=\scriptsize, fill=white, inner sep=0.5pt] at (90:2) {$2q$};
    \node[font=\scriptsize, fill=white, inner sep=0.5pt] at (90:3) {$3q$};
    \draw[azulprof, line width=0.5pt] (1,0) -- (3.6,0) -- (3.6,-0.5);
    \pic at (3.6,-0.5) {tierra};
    \node[fondo, font=\scriptsize, anchor=south, text=grispizarra] at (2.3,0.05) {hilo};
    \node at (0,-4.1) {(a)};
  \end{scope}
  % ---------------------------------------------------------------- (b)
  \begin{scope}[xshift=4.4cm, yshift=-2.25cm]
    \begin{axis}[tema, width=8.4cm, height=5.0cm, xmin=0, xmax=5, ymin=0, ymax=3.3,
        xlabel={$r/a$}, ylabel={$\phi\,/\,(\kc q/a)$}, xtick={0,1,2,3,4,5}, ytick={0,1,2,3},
        legend style={at={(0.99,0.97)}, anchor=north east}]
      \addplot[azulprof, line width=1.2pt] coordinates {(0,3) (1,3)};
      \addlegendentry{antes}
      \addplot[azulprof, line width=1.2pt, domain=1:2, samples=30, forget plot] {2+1/x};
      \addplot[azulprof, line width=1.2pt, domain=2:3, samples=30, forget plot] {1+3/x};
      \addplot[azulprof, line width=1.2pt, domain=3:5, samples=30, forget plot] {6/x};
      \addplot[azulcielo, line width=1.2pt, densely dashed] coordinates {(0,0) (1,0)};
      \addlegendentry{interior a tierra}
      \addplot[azulcielo, line width=1.2pt, densely dashed, domain=1:2, samples=30, forget plot] {2-2/x};
      \addplot[azulcielo, line width=1.2pt, densely dashed, forget plot] coordinates {(2,1) (3,1)};
      \addplot[azulcielo, line width=1.2pt, densely dashed, domain=3:5, samples=30, forget plot] {3/x};
      \foreach \x in {1,2,3}{\edef\tmp{\noexpand\draw[punteado] (axis cs:\x,0) -- (axis cs:\x,3.3);}\tmp}
      \draw[{Stealth[length=3pt]}-{Stealth[length=3pt]}, cargapos] (axis cs:3.12,1) -- (axis cs:3.12,2)
        node[midway, right, font=\scriptsize] {$\Delta\phi_c$};
    \end{axis}
    \node at (3.6,-1.4) {(b)};
  \end{scope}
\end{tikzpicture}
\caption{Ejemplo modelo. (a) Tres esferas concéntricas de radios $a$, $2a$ y $3a$; la interior se conecta a tierra. (b) Potencial antes y después de la conexión:
la esfera interior pasa de $q$ a $-2q$, el potencial de la externa baja de $2\kc q/a$ a $\kc q/a$ y, entre $2a$ y $3a$, el campo se anula porque la carga
encerrada es $-2q+2q=0$.}
\label{fig:sem-modelo}
\end{figure}


% -------------------------------------------------------------------------
\subsection{Resumen de resultados y lista de control}\label{sem:resumen}
% -------------------------------------------------------------------------
\begin{table}[ht]
\centering
\caption{Resultados de los seminarios del Tema 3.}
\label{tab:resultados}
\small
\begin{tabular}{@{}clc@{}}
  \toprule
  Problema & Resultado & Método\\
  \midrule
  3.1 & $\phi_{\max}=\Erup R=150\,\mathrm{kV}$ & M1\\[3pt]
  3.2 & $Q_1=4\pi\eo V_0\dfrac{R_1R_4}{R_4-R_1}$,\quad $\phi_{2,3}=\dfrac{Q_1}{4\pi\eo}\Bigl(\dfrac{1}{R_{2,3}}-\dfrac{1}{R_4}\Bigr)$ & M2, M3\\[10pt]
  3.3 & $c_{11}=-c_{12}=4\pi\eo\dfrac{ab}{b-a}$,\quad $c_{22}=4\pi\eo\dfrac{b^{2}}{b-a}$,\quad $C=4\pi\eo\dfrac{ab}{b-a}$ & M3\\[10pt]
  3.4 & $C/\ell=2\pi\eo/\ln(b/a)$ & M2\\[3pt]
  3.5 & $C=\eo A/d$;\quad caras: $q_1=q_4=\tfrac12(Q_1+Q_2)$,\ $q_2=-q_3=\tfrac12(Q_1-Q_2)$ & M4\\[4pt]
  3.6 & $\phi_2=\dfrac{Q}{4\pi\eo}\,\dfrac{2d^{2}+Rd-4R^{2}}{Rd\,(2d+R)}$ & M5\\[10pt]
  3.7 & $q_1=-\dfrac{2qa}{b}$,\quad $q_2=-\dfrac{qa(b-2a)}{b^{2}}$,\quad $q_3=\dfrac{qa^{2}(3b-2a)}{b^{3}}$ & M5\\[10pt]
  Modelo & $\Delta\phi_c=-q/(4\pi\eo a)$ & M3\\
  \bottomrule
\end{tabular}
\end{table}

\needspace{20\baselineskip}
\begin{receta}[Lista de control]
\begin{itemize}
  \item ¿He asignado la carga de cada cara con la regla de las caras?
  \item ¿He distinguido \emph{a tierra} ($\phi=0$, pero $Q\neq0$ en general) de \emph{descargado} ($Q=0$, pero $\phi\neq0$ en general)?
  \item En los $p_{ij}$, ¿los demás conductores están descargados? En los $c_{ij}$, ¿están a tierra?
  \item ¿Son simétricas las matrices? ¿Tienen los signos correctos ($p_{ij}>0$, $c_{ii}>0$, $c_{ij}<0$)?
  \item ¿Se cumple $c_{jj}\ge-\sum_{i\neq j}c_{ij}$, con igualdad en influencia total?
  \item ¿Coincide la capacidad calculada directamente con $1/(p_{11}+p_{22}-2p_{12})$?
  \item En procesos por etapas, ¿he congelado la carga de los conductores ya desconectados?
  \item Un hilo fino solo iguala potenciales: ¿he despreciado su carga y su influencia?
\end{itemize}
\end{receta}

% =========================================================================
\appendix
\section{Formulario del Tema 3}\label{sec:apendice}
% =========================================================================
La \cref{tab:formulario} reúne las expresiones centrales del tema, con la referencia al enunciado en que se demuestran.

\begingroup
\centering
\captionof{table}{Formulario del Tema 3. Conductores en equilibrio electrostático ($\kc=1/4\pi\eo$, $\phi(\infty)=0$).}
\label{tab:formulario}
\small
\renewcommand{\arraystretch}{1.55}
\begin{tabularx}{\linewidth}{@{}>{\raggedright\arraybackslash}p{4.6cm}>{\centering\arraybackslash}X>{\raggedleft\arraybackslash}p{2.1cm}@{}}
  \toprule
  Magnitud o resultado & Expresión & Referencia\\
  \midrule
  \multicolumn{3}{@{}l}{\textsc{\color{azulmedio}Conductor en equilibrio}}\\
  Interior & $\vE=\vc{0}$,\quad $\phi=\mathrm{cte}$,\quad $\rho=0$ & \cref{teo:propiedades}\\
  Campo junto a la superficie & $\vE=\dfrac{\sigma}{\eo}\,\un$ & \cref{eq:campo-superficie}\\
  Esferas unidas por un hilo & $\sigma_1R_1=\sigma_2R_2$ & \cref{prop:punta}\\
  Campo del resto de la carga & $\Erest=\dfrac{\sigma}{2\eo}\,\un=\tfrac12(\vE_1+\vE_2)$ & \cref{lem:resto}\\
  Presión electrostática & $\Pel=\dfrac{\sigma^{2}}{2\eo}=\dfrac{\eo}{2}E^{2}$ (hacia fuera) & \cref{teo:presion}\\
  Ecuaciones de Poisson y Laplace & $\Lap\phi=-\rho/\eo$,\quad $\Lap\phi=0$ & \cref{prop:poisson}\\
  Cavidad vacía & $\Ecav=\vc{0}$ & \cref{teo:faraday}\\
  Regla de las caras & $\Qint=-q$,\quad $\Qext=Q+q$ & \cref{prop:caras}\\
  Capacidad de un conductor aislado & $C=Q/\phi$,\quad $C_{\text{esfera}}=4\pi\eo R$ & \cref{teo:capacidad}\\
  \midrule
  \multicolumn{3}{@{}l}{\textsc{\color{azulmedio}Sistemas de conductores}}\\
  Coeficientes de potencial & $\phi_i=\sum_jp_{ij}Q_j$,\quad $p_{ij}=\dfrac{\phi_i}{Q_j}\Big|_{Q_{k\neq j}=0}$ & \cref{teo:pij}\\
  Propiedades & $p_{ij}=p_{ji}$,\quad $p_{jj}\ge p_{ij}\ge0$ & \cref{teo:pij-prop}\\
  Esferas concéntricas delgadas & $p_{ij}=\kc/\max(R_i,R_j)$ & \cref{prop:reglas}\\
  Esferas lejanas & $p_{ii}\simeq\kc/R_i$,\quad $p_{ij}\simeq\kc/d_{ij}$ & \cref{prop:reglas}\\
  Coeficientes de capacidad e influencia & $Q_i=\sum_jc_{ij}\phi_j$,\quad $\Cm=\Pm^{-1}$,\quad $c_{ij}=\dfrac{Q_i}{\phi_j}\Big|_{\phi_{k\neq j}=0}$ & \cref{def:cij}\\
  Propiedades & $c_{ij}=c_{ji}$,\quad $c_{jj}>0$,\quad $c_{ij}\le0$ & \cref{teo:cij-prop}\\
  Desigualdad de influencia & $c_{jj}\ge-\sum_{i\neq j}c_{ij}$ (igualdad: influencia total) & \cref{teo:influencia}\\
  Inversa $2\times2$ & $\Cm=\dfrac{1}{p_{11}p_{22}-p_{12}^{2}}\begin{pmatrix}p_{22}&-p_{12}\\-p_{12}&p_{11}\end{pmatrix}$ & \cref{ej:esferico}\\
  \midrule
  \multicolumn{3}{@{}l}{\textsc{\color{azulmedio}Condensadores}}\\
  Capacidad & $C=\dfrac{Q}{\DV}=\dfrac{1}{p_{11}+p_{22}-2p_{12}}=\dfrac{c_{11}c_{22}-c_{12}^{2}}{c_{11}+c_{22}+2c_{12}}$ & \cref{teo:cap-coef}\\
  Armadura encerrada & $C=c_{11}=-c_{12}$ & \cref{cor:encerrado}\\
  Esférico, cilíndrico, plano & $4\pi\eo\dfrac{ab}{b-a}$,\quad $\dfrac{2\pi\eo\ell}{\ln(b/a)}$,\quad $\eo\dfrac{A}{d}$ & \cref{tab:canonicos}\\
  \bottomrule
\end{tabularx}
\par\endgroup


% =========================================================================
%                              BIBLIOGRAFÍA
% =========================================================================
\FloatBarrier
\printbibliography[title={Referencias}, heading=bibintoc]

\end{document}
